\beginnumbering
\eledsection*{Adhyāya 21}
\pstart \ledsidenote{\passref{4.21.1}}\edtext{\skt{अथातः}}{\lemma{\skt{अथातः}}\Afootnote{A H K\textsuperscript{pc}; \skt{अथात}\ K\textsuperscript{ac}}}\skt{ शूक\-दो\-ष\-चि\-कि\-त्सितं }\edtext{\skt{व्याख्यास्यामः}}{\lemma{\skt{व्याख्यास्यामः}}\Afootnote{A H; \skt{व्याव्याख्यास्यामः~॥}\ K}}\skt{~॥}\pend

\pstart\versepar \ledsidenote{\passref{4.21.3}}\edtext{\skt{विलिख्य}}{\lemma{\skt{विलिख्य}}\Afootnote{K H; \skt{संलिख्य}\ A}}\skt{ सर्षपीं सम्य\-क्क\-षा\-यै\-र\-व\-चू\-र्ण्ण\-येत्~। }\pend
\pstart\versepar \skt{कषायेष्वेव }\edtext{\skt{कुर्वीत}}{\lemma{\skt{कुर्वीत}}\Afootnote{K H; \emph{om.}\ A}}\skt{ तैलञ्च }\edtext{\skt{व्रणरोपणं}}{\lemma{\skt{व्रणरोपणं}}\Afootnote{K H; \skt{कुर्वीत व्रणरोपणम्॥}\ A}}\skt{~॥}\pend

\pstart\versepar \ledsidenote{\passref{4.21.4}}\skt{अष्ठीलिकां }\edtext{\edtext{\skt{जलौ\-का\-भि\-र्ग्रा\-ह\-ये\-द्ब\-हुशो}}{\lemma{\skt{जलौकाभिर्}}\Afootnote{K H; \skt{जलौकोभिर्}\ A}}}{\lemma{\skt{ग्राहयेद्बहुशो}}\Afootnote{K H; \skt{ग्राहयेच्च पुनः}\ A}}\skt{ भिषक्~। }\pend
\pstart\versepar \skt{तथा चानुपशाम्यन्तीं कफ\-ग्र\-न्थि\-व\-दु\-द्ध\-रेत्~।}\edtext{}{\lemma{\skt{भिषक्~। तथा}}\Afootnote{K H; \skt{पुनः~। तथा}\ A}}\pend

\pstart\versepar \ledsidenote{\passref{4.21.5}}\edtext{\skt{स्वेदयेद्ग्रथितां}}{\lemma{\skt{स्वेदयेद्ग्रथितां}}\Afootnote{K H; \skt{स्वेदयेद्ग्रथितं}\ A}}\skt{ }\edtext{\skt{शश्वन्नाडीं}}{\lemma{\skt{नाडीं}}\Afootnote{K H; \skt{नाडी}\ A}}\skt{ स्वेदेन बुद्धिमान्~। }\pend
\pstart\versepar \skt{सुखो\-ष्णै\-रु\-प\-ना\-हैश्च }\edtext{\skt{सुस्वि\-न्नै\-रु\-प\-ना\-ह\-येत्}}{\lemma{\skt{सुस्विन्नैर्}}\Afootnote{K H; \skt{सुस्निग्धैर्}\ A}}\skt{~॥}\pend

\pstart\versepar \ledsidenote{\passref{4.21.6}}\skt{कुम्भीकां }\edtext{\skt{पाकमापन्नां}}{\lemma{\skt{पाकमापन्नां}}\Afootnote{\skt{पाकमापन्नां}\ A K H}}\skt{ }\edtext{\skt{च्छिन्द्याच्छोध्य च}}{\lemma{\skt{च्छिन्द्याच्छोध्य च}}\Afootnote{K H; \skt{भिन्द्याच्छुद्धां तु}\ A}}\skt{ रोपयेत्~॥ }\pend
\pstart\versepar \skt{तैलेन }\edtext{\skt{लोध्र\-त्रि\-फ\-ला\-ति\-न्दु\-का\-म्र\-कृ\-तेन}}{\lemma{\skt{लोध्र\appbrk{}त्रि\appbrk{}फ\appbrk{}ला\appbrk{}ति\appbrk{}न्दु\appbrk{}का\appbrk{}म्र\appbrk{}कृ\appbrk{}तेन}}\Afootnote{K H; \skt{त्रिफ\appbrk{}ला\appbrk{}लो\appbrk{}ध्र\appbrk{}ति\appbrk{}न्दु\appbrk{}का\appbrk{}म्रा\appbrk{}त\appbrk{}केन}\ A}}\skt{ }\edtext{\skt{वा}}{\lemma{\skt{वा}}\Afootnote{K H; \skt{तु~॥}\ A}}\skt{~॥}\pend

\pstart\versepar \ledsidenote{\passref{4.21.7}}\edtext{\edtext{\skt{जला\-यु\-का\-भि\-र\-ल\-जी\-ङ्ग्रा\-ह\-ये\-त्स्वे\-द\-यीत}}{\lemma{\skt{जलायुकाभिर्}}\Afootnote{K H; \skt{ग्राहयित्वा जलौकोभिर्}\ A}}\skt{ च}}{\lemma{\skt{ग्राह\appbrk{}ये\appbrk{}त्स्वे\appbrk{}द\appbrk{}यीत च}}\Afootnote{K H; \skt{सेचयेत्ततः~।}\ A}}\skt{~। }\pend
\pstart\versepar \skt{कषायैस्तेषु सिद्धञ्च तैलं रोपणमिष्यते~॥}\pend

\pstart\versepar \ledsidenote{\passref{4.21.8}}\skt{बलातैलेन }\edtext{\skt{सिद्धेन}}{\lemma{\skt{सिद्धेन}}\Afootnote{K H; \skt{कोष्णेन}\ A}}\skt{ मृदि\-त\-म्प\-रि\-षे\-च\-येत्~। }\pend
\pstart\versepar \skt{सर्पिः }\edtext{\skt{स्निग्धैश्च}}{\lemma{\skt{स्निग्धैश्च}}\Afootnote{K H; \emph{om.}\ A}}\skt{ }\edtext{\skt{मधुरैः}}{\lemma{\skt{मधुरैः}}\Afootnote{K H; \skt{धुरैः}\ A}}\skt{ }\edtext{\skt{सुखो\-ष्णै\-रु\-प\-ना\-ह\-येत्}}{\lemma{\skt{सुखोष्णैर्}}\Afootnote{K H; \skt{सर्पिषा स्निग्धैः सुखोष्णैर्}\ A}}\skt{~॥}\edtext{}{\lemma{\skt{परिषेचयेत्~। सर्पिः}}\Afootnote{K H; \skt{परिषेचयेत्~। भ}\ A}}\pend

\pstart\versepar \ledsidenote{\passref{4.21.9}}\edtext{\skt{सम्मूढपिटकां}}{\lemma{\skt{सम्मूढपिटकां}}\Afootnote{K H; \skt{संमूढपिडकां}\ A}}\skt{ क्षिप्रं }\edtext{\skt{जलौ\-का\-भि\-रु\-पा\-च\-रेत्}}{\lemma{\skt{जलौकाभिर्}}\Afootnote{K H; \skt{जलौकोभिर्}\ A}}\skt{~। }\pend
\pstart\versepar \skt{भित्वा }\edtext{\skt{पर्यागतं}}{\lemma{\skt{पर्यागतं}}\Afootnote{K; \skt{पर्यागतां}\ A H}}\skt{ चापि }\edtext{\skt{लेप\-ये\-न्म\-धु\-स\-र्पिषा}}{\lemma{\skt{लेप\appbrk{}ये\appbrk{}न्म\appbrk{}धु\appbrk{}स\appbrk{}र्पिषा}}\Afootnote{H; \skt{लेप\appbrk{}ये\appbrk{}त्क्षौ\appbrk{}द्र\appbrk{}स\appbrk{}र्पिषा~॥}\ A; \skt{लेपये मधुसर्पिषा~॥}\ K}}\skt{~॥}\pend

\pstart\versepar \ledsidenote{\passref{4.21.10}}\skt{अवमन्थे गते पाकं भिन्ने तैलं विधीयते~। }\pend
\pstart\versepar \skt{}\edtext{\skt{धवा\-ज\-क\-र्ण्णा\-श्व\-क\-र्ण्ण\-स\-ल्ल\-की\-क्वा\-थ\-सा\-धितं}}{\lemma{\skt{धवा\appbrk{}ज\appbrk{}क\appbrk{}र्ण्णा\appbrk{}श्व\appbrk{}क\appbrk{}र्ण्ण\appbrk{}स\appbrk{}ल्ल\appbrk{}की\appbrk{}क्वा\appbrk{}थ\appbrk{}सा\appbrk{}धितं}}\Afootnote{K H; \skt{धवा\appbrk{}श्व\appbrk{}क\appbrk{}र्ण\appbrk{}प\appbrk{}त्त\appbrk{}ङ्ग\appbrk{}स\appbrk{}ल्ल\appbrk{}की\appbrk{}ति\appbrk{}न्दु\appbrk{}की\appbrk{}कृ\appbrk{}तम्~॥}\ A}}\skt{~॥}\pend

\pstart \ledsidenote{\passref{4.21.11}}\skt{क्रिया\-म्पु\-ष्क\-रि\-कायां तु शीतां सर्वां }\edtext{\skt{प्रयोजयेत्}}{\lemma{\skt{प्रयोजयेत्}}\Afootnote{K H; \skt{प्रयोजयेत्~। जलौ\appbrk{}को\appbrk{}भि\appbrk{}र्ह\appbrk{}रे\appbrk{}च्चा\appbrk{}सृ\appbrk{}क्स\appbrk{}र्पिषा चावसेचयेत्~॥}\ A}}\skt{~॥}\pend

\pstart\versepar \ledsidenote{\passref{4.21.12}}\edtext{\skt{स्पर्शहार्यां}}{\lemma{\skt{स्पर्शहार्यां}}\Afootnote{K H; \skt{स्पर्शहान्यां}\ A}}\skt{ हरेद्रक्तं }\edtext{\skt{मधुरैः}}{\lemma{\skt{मधुरैः}}\Afootnote{K H; \skt{प्रदि\appbrk{}ह्या\appbrk{}न्म\appbrk{}धु\appbrk{}रैर्}\ A}}\skt{ }\edtext{\skt{प्रदिहेदपि}}{\lemma{\skt{प्रदिहेदपि}}\Afootnote{K H; \skt{अपि~।}\ A}}\skt{~। }\pend
\pstart\versepar \skt{क्षीरे\-क्षु\-र\-स\-स\-र्पिर्भिः सेचयेच्च सुशीतलैः~॥}\pend

\pstart\versepar \ledsidenote{\passref{4.21.13}}\edtext{\skt{उत्तमाख्यां तु}}{\lemma{\skt{उत्तमाख्यां तु}}\Afootnote{K H; \skt{पिड\appbrk{}का\appbrk{}मु\appbrk{}त्त\appbrk{}माख्यां}\ A}}\skt{ }\edtext{\skt{पिटकां}}{\lemma{\skt{पिटकां}}\Afootnote{K H; \skt{च}\ A}}\skt{ बडि\-से\-नो\-द्ध\-रे\-द्भि\-षक्~। }\pend
\pstart\versepar \skt{उद्धृत्य मधुसंयुक्तैः कषा\-यै\-र\-व\-चू\-र्ण्ण\-येत्~॥}\edtext{}{\lemma{\skt{बडि\appbrk{}से\appbrk{}नो\appbrk{}द्ध\appbrk{}रे\appbrk{}द्भि\appbrk{}षक्~। उद्धृत्य}}\Afootnote{K H; \skt{बडि\appbrk{}शे\appbrk{}नो\appbrk{}द्ध\appbrk{}रे\appbrk{}द्भि\appbrk{}षक्~। उद्धृत्य}\ A}}\pend

\pstart\versepar \ledsidenote{\passref{4.21.14}}\skt{रसक्रिया विधातव्या लिखिते }\edtext{\skt{शतपोशनके}}{\lemma{\skt{शतपोशनके}}\Afootnote{K H\textsuperscript{pc}; \skt{शतपोनके}\ A H\textsuperscript{ac}}}\skt{~। }\pend
\pstart\versepar \skt{पृथ\-क्प\-र्ण्या\-दि\-सि\-द्धञ्च देय\-न्तै\-ल\-म\-न\-न्तरं~॥}\pend

\pstart\versepar \ledsidenote{\passref{4.21.15}}\skt{क्रियां }\edtext{\skt{प्रयुञ्ज्यात्कुशलं}}{\lemma{\skt{प्रयुञ्ज्यात्कुशलं}}\Afootnote{K H; \skt{कुर्या\appbrk{}द्भि\appbrk{}ष\appbrk{}क्प्रा\appbrk{}ज्ञस्}\ A}}\skt{ }\edtext{\skt{त्वक्पाकस्य}}{\lemma{\skt{त्वक्पाकस्य}}\Afootnote{A K\textsuperscript{pc} H; \skt{त्वक्प्कस्य}\ K\textsuperscript{ac}}}\skt{ विसर्पवत्~॥ }\pend
\pstart\versepar \skt{रक्त\-वि\-द्र\-धि\-व\-च्चापि }\edtext{\skt{क्रियां}}{\lemma{\skt{क्रियां}}\Afootnote{K H; \skt{क्रिया}\ A}}\skt{ शोणितजे }\edtext{\skt{र्बुदे}}{\lemma{\skt{र्बुदे}}\Afootnote{K; \skt{ऽर्बुदे}\ A H}}\skt{~।}\pend

\pstart\versepar \ledsidenote{\passref{4.21.16}}\skt{कषा\-य\-क\-ल्क\-स\-र्पींषि }\edtext{\skt{तैल\-चू\-र्ण्ण\-र\-स\-क्रियां}}{\lemma{\skt{तैल\appbrk{}चू\appbrk{}र्ण्ण\appbrk{}र\appbrk{}स\appbrk{}क्रियां}}\Afootnote{K H; \skt{तैलं चूर्णं रसक्रियाम्~।}\ A}}\skt{~॥ }\pend
\pstart\versepar \skt{}\edtext{\skt{शोधने}}{\lemma{\skt{शोधने}}\Afootnote{K H; \skt{शोधनं}\ A}}\skt{ }\edtext{\skt{रोपणे}}{\lemma{\skt{रोपणे}}\Afootnote{K H; \skt{रोपणं}\ A}}\skt{ चैव वीक्ष्य वीक्ष्या\-व\-चा\-र\-येत्~।}\pend

\pstart\versepar \ledsidenote{\passref{4.21.17}}\skt{हितञ्च सर्पिषः पानं पथ्यञ्चापि }\edtext{\skt{विरेचनं}}{\lemma{\skt{विरेचनं}}\Afootnote{K H; \skt{चिरेचनम्~।}\ A}}\skt{~। }\pend
\pstart\versepar \skt{हितः शोणितमोक्षश्च यच्चापि लघु भोजनं~॥}\pend

\pstart\versepar \ledsidenote{\passref{4.21.18}}\skt{अर्बु\-द\-म्मां\-स\-पा\-कञ्च विद्रधिं तिलकालकं~। }\pend
\pstart\versepar \skt{}\edtext{\skt{प्रत्याख्याय हि}}{\lemma{\skt{प्रत्याख्याय हि}}\Afootnote{K H; \skt{प्रत्याख्याय प्र}\ A}}\skt{ कुर्वीत }\edtext{\edtext{\edtext{\skt{भिष\-क्ते\-षा\-म्प्र\-ति\-क्रि\-या\-मिति}}{\lemma{\skt{भिषक्तेषाम्}}\Afootnote{K H; \skt{भिषक्सम्य}\ A}}}{\lemma{\skt{प्रतिक्रियाम्}}\Afootnote{K H; \skt{क्प्रतिक्रियाम्~॥}\ A}}}{\lemma{\skt{इति}}\Afootnote{K H; \emph{om.}\ A}}\skt{~॥}\pend

\pstart \ledsidenote{\passref{4.21.trailer}}\edtext{\skt{चिचिकित्सास्थाने}}{\lemma{\skt{चिचिकित्सास्थाने}}\Afootnote{\skt{इति सुश्रु\appbrk{}त\appbrk{}सं\appbrk{}हि\appbrk{}तायां चिकित्सास्थाने शूक\appbrk{}रो\appbrk{}ग\appbrk{}चि\appbrk{}कि\appbrk{}त्सितं नामै\appbrk{}क\appbrk{}विं\appbrk{}शोऽ\appbrk{}ध्यायः~॥२१~॥}\ A; \skt{चिचिकित्सास्थाने २१~॥~॥ ०~॥}\ K; \skt{चिकित्सा २१~॥~॥}\ H}}\skt{ २१~॥~॥ ०~॥}\pend

\endnumbering
\beginnumbering
\eledsection*{Adhyāya 22}
\pstart \ledsidenote{\passref{4.22.1}}\skt{अथातो }\edtext{\skt{मुखरोगाणां}}{\lemma{\skt{मुखरोगाणां}}\Afootnote{K H; \skt{मुखरोग}\ A}}\skt{ चिकित्सितं }\edtext{\skt{व्याख्यास्यामः}}{\lemma{\skt{व्याख्यास्यामः}}\Afootnote{A H; \emph{om.}\ K}}\skt{~॥}\pend

\pstart\versepar \ledsidenote{\passref{4.22.3}}\skt{चतुर्विधेन स्नेहेन }\edtext{\skt{मधु\-च्छि\-ष्ट\-यु\-तेन वा}}{\lemma{\skt{मधु\appbrk{}च्छि\appbrk{}ष्ट\appbrk{}यु\appbrk{}तेन वा}}\Afootnote{K; \skt{मधू\appbrk{}च्छि\appbrk{}ष्ट\appbrk{}यु\appbrk{}तेन च~।}\ A; \skt{मधू\appbrk{}च्छि\appbrk{}ष्ट\appbrk{}यु\appbrk{}तेन वा~।}\ H}}\skt{~। }\pend
\pstart\versepar \skt{}\edtext{\skt{वाति\-के\-भ्य\-ञ्ज\-नकं}}{\lemma{\skt{वाति\appbrk{}के\appbrk{}भ्य\appbrk{}ञ्ज\appbrk{}नकं}}\Afootnote{K; \skt{वातजेऽभ्यञ्जनं}\ A; \skt{वातिकेऽभ्यञ्जनं}\ H}}\skt{ }\edtext{\edtext{\skt{कार्यं\-ना\-डी\-स्वे\-द\-म्त\-थापि}}{\lemma{\skt{कार्यं\appbrk{}ना\appbrk{}डी\appbrk{}स्वे\appbrk{}दम्}}\Afootnote{K H; \skt{दन्}\ A}}}{\lemma{\skt{तथापि}}\Afootnote{K; \skt{तवैद}\ A; \emph{om.}\ H}}\skt{ }\edtext{\skt{च}}{\lemma{\skt{च}}\Afootnote{K; \skt{र्भे दन्तमूलानि शोधयेत्~। ततः क्सारं प्रयुञ्जीत क्रियाः सर्वाश्च शीतलाः~॥}\ A; \skt{तथापि च~॥}\ H}}\skt{~।}\pend

\pstart\versepar \ledsidenote{\passref{4.22.4}}\edtext{\skt{ओष्ठप्रकोपे}}{\lemma{\skt{ओष्ठप्रकोपे}}\Afootnote{K H; \skt{विद\appbrk{}ध्या\appbrk{}दो\appbrk{}ष्ठ\appbrk{}कोपे}\ A}}\skt{ }\edtext{\skt{भिषजा}}{\lemma{\skt{भिषजा}}\Afootnote{K H; \skt{तु}\ A}}\skt{ साल्वणं }\edtext{\skt{चोपनाहनं}}{\lemma{\skt{चोपनाहनं}}\Afootnote{K H; \skt{चोपनाहने।}\ A}}\skt{~। }\pend
\pstart\versepar \skt{मस्तिष्के चैव }\edtext{\skt{नस्ये च}}{\lemma{\skt{नस्ये च}}\Afootnote{A K H\textsuperscript{pc}; \skt{नस्ये}\ H\textsuperscript{ac}}}\skt{ तैलं }\edtext{\skt{वातहरै}}{\lemma{\skt{वातहरै}}\Afootnote{K; \skt{वातहरं}\ A; \skt{वातहरैः}\ H}}\skt{ }\edtext{\skt{स्रितं}}{\lemma{\skt{स्रितं}}\Afootnote{K; \skt{हितम्॥}\ A; \skt{श्रितं~॥}\ H}}\skt{~।}\pend

\pstart\versepar \ledsidenote{\passref{4.22.5}}\skt{श्रीवेष्टकं सर्जरसं }\edtext{\skt{गुग्गुलुं}}{\lemma{\skt{गुग्गुलुं}}\Afootnote{K H; \emph{om.}\ A}}\skt{ }\edtext{\skt{सुरदारु च}}{\lemma{\skt{सुरदारु च}}\Afootnote{K H; \skt{सुरदारु स}\ A}}\skt{~। }\pend
\pstart\versepar \skt{}\edtext{\edtext{\skt{यष्टी\-म\-धु\-क\-चू\-र्ण्णञ्च}}{\lemma{\skt{यष्टी\appbrk{}म\appbrk{}धु\appbrk{}क\appbrk{}चू\appbrk{}र्ण्णञ्}}\Afootnote{K; \skt{गुग्गुलु। यष्टी\appbrk{}म\appbrk{}धु\appbrk{}क\appbrk{}चूर्णं तु}\ A; \skt{यष्टि\appbrk{}म\appbrk{}धु\appbrk{}क\appbrk{}चू\appbrk{}र्ण्णञ्}\ H}}}{\lemma{\skt{च}}\Afootnote{K H; \skt{विदध्यात्प्र}\ A}}\skt{ }\edtext{\skt{हितमत्र}}{\lemma{\skt{हितमत्र}}\Afootnote{H; \skt{तिसारणम्॥}\ A; \skt{हितमत्र}\ K}}\skt{ }\edtext{\skt{प्रलेपने}}{\lemma{\skt{प्रलेपने}}\Afootnote{K H; \emph{om.}\ A}}\skt{~।}\pend

\pstart\versepar \ledsidenote{\passref{4.22.6}}\skt{पित्त\-र\-क्ता\-भि\-घा\-तोत्थं }\edtext{\skt{जलो\-का\-भि\-रु\-पा\-च\-रेत्}}{\lemma{\skt{जलोकाभिर्}}\Afootnote{K; \skt{जलौकोभिर्}\ A; \skt{जलौकाभिर्}\ H}}\skt{~। }\pend
\pstart\versepar \skt{पित्त\-वि\-द्र\-धि\-व\-च्चापि क्रियां कुर्यादशेषतः~।}\pend

\pstart\versepar \ledsidenote{\passref{4.22.7}}\skt{शिरोविरेचनं }\edtext{\skt{धूमं}}{\lemma{\skt{धूमं}}\Afootnote{K H; \skt{धूमः}\ A}}\skt{ }\edtext{\skt{स्वेदं}}{\lemma{\skt{स्वेदं}}\Afootnote{K H; \skt{स्वेदः}\ A}}\skt{ }\edtext{\skt{कवडधारणं}}{\lemma{\skt{कवडधारणं}}\Afootnote{K H; \skt{कवल एव च।}\ A}}\skt{~। }\pend
\pstart\versepar \skt{}\edtext{\edtext{\skt{हृत\-र\-क्त\-प्र\-यो\-क्त\-व्य\-मो\-ष्ठ\-कोपे}}{\lemma{\skt{हृत\appbrk{}र\appbrk{}क्त\appbrk{}प्र\appbrk{}यो\appbrk{}क्त\appbrk{}व्यम्}}\Afootnote{K; \skt{हृते रक्ते प्रयोक्तव्य}\ A; \skt{हृतरक्ते प्रयोक्तव्यम्}\ H}}}{\lemma{\skt{ओष्ठकोपे}}\Afootnote{K H; \skt{मोष्ठकोपे}\ A}}\skt{ कफात्मके~।}\pend

\pstart\versepar \ledsidenote{\passref{4.22.8}}\edtext{\skt{त्रिकटुं}}{\lemma{\skt{त्रिकटुं}}\Afootnote{K H; \skt{त्र्यूषणं}\ A}}\skt{ }\edtext{\skt{स्वर्जिकाक्षारं}}{\lemma{\skt{स्वर्जिकाक्षारं}}\Afootnote{K H; \skt{स्वर्जिकाक्षारो}\ A}}\skt{ }\edtext{\skt{क्षारञ्च}}{\lemma{\skt{क्षारञ्च}}\Afootnote{K H; \emph{om.}\ A}}\skt{ }\edtext{\skt{यवशूकजं}}{\lemma{\skt{यवशूकजं}}\Afootnote{K H; \skt{यवक्षारो विडं तथा।}\ A}}\skt{~। }\pend
\pstart\versepar \skt{}\edtext{\skt{क्षौद्रयुक्तं}}{\lemma{\skt{क्षौद्रयुक्तं}}\Afootnote{A K; \skt{क्षौद्रयुक्तम}\ H}}\skt{ }\edtext{\edtext{\skt{विधा\-त\-व्य\-मे\-तच्च}}{\lemma{\skt{विधातव्यम्}}\Afootnote{K H; \skt{विधातव्य}\ A}}}{\lemma{\skt{एतच्च}}\Afootnote{K H; \skt{मेतच्च}\ A}}\skt{ प्रतिसारणं~।}\pend

\pstart\versepar \ledsidenote{\passref{4.22.9}}\skt{मेदोजे स्वेदिते भिन्ने }\edtext{\skt{शोधिते}}{\lemma{\skt{शोधिते}}\Afootnote{A K; \skt{शोणिते}\ H}}\skt{ ज्वलनो हितः~। }\pend
\pstart\versepar \skt{}\edtext{\skt{प्रियं\-गु\-लो\-ध्र\-त्रि\-फला}}{\lemma{\skt{प्रियं\appbrk{}गु\appbrk{}लो\appbrk{}ध्र\appbrk{}त्रि\appbrk{}फला}}\Afootnote{K H; \skt{प्रियङ्गुत्रिफला}\ A}}\skt{ }\edtext{\skt{सक्षौद्रं}}{\lemma{\skt{सक्षौद्रं}}\Afootnote{K H; \skt{लोध्रं सक्षौद्रं}\ A}}\skt{ प्रतिसारणं~।}\pend

\pstart\versepar \ledsidenote{\passref{4.22.10}}\edtext{\skt{एत\-दो\-ष्ठ\-प्र\-को\-पानां}}{\lemma{\skt{एत\appbrk{}दो\appbrk{}ष्ठ\appbrk{}प्र\appbrk{}को\appbrk{}पानां}}\Afootnote{K H; \skt{एत\appbrk{}दो\appbrk{}ष्ठ\appbrk{}प्र\appbrk{}को\appbrk{}पानां}\ A}}\skt{ साध्यानां कर्म कीर्त्तितं~॥ }\pend
\pstart\versepar \skt{दन्त\-मू\-ल\-ग\-तानां तु }\edtext{\skt{रोगाणां}}{\lemma{\skt{रोगाणां}}\Afootnote{A K; \skt{रोगानाङ्}\ H}}\skt{ कर्म }\edtext{\skt{उच्यते}}{\lemma{\skt{उच्यते}}\Afootnote{K H; \skt{वक्ष्यते॥}\ A}}\skt{~।}\pend

\pstart\versepar \ledsidenote{\passref{4.22.11}}\skt{शीतादे }\edtext{\skt{हृतरक्तस्य}}{\lemma{\skt{हृतरक्तस्य}}\Afootnote{K H; \skt{हृतरक्ते तु}\ A}}\skt{ }\edtext{\skt{तोय\-ना\-ग\-र\-स\-र्ष\-पान्}}{\lemma{\skt{तोय\appbrk{}ना\appbrk{}ग\appbrk{}र\appbrk{}स\appbrk{}र्ष\appbrk{}पान्}}\Afootnote{K; \skt{तोये नाग\appbrk{}र\appbrk{}स\appbrk{}र्ष\appbrk{}पान्}\ A H}}\skt{~॥ }\pend
\pstart\versepar \skt{}\edtext{\skt{निःक्वाथ्य}}{\lemma{\skt{निःक्वाथ्य}}\Afootnote{K H; \skt{निष्क्वाथ्य}\ A}}\skt{ }\edtext{\skt{त्रिफलाञ्चापि}}{\lemma{\skt{चापि}}\Afootnote{K H; \emph{om.}\ A}}\skt{ }\edtext{\skt{कुर्या\-द्ग\-ण्डू\-ष\-धा\-रणं}}{\lemma{\skt{कुर्या\appbrk{}द्ग\appbrk{}ण्डू\appbrk{}ष\appbrk{}धा\appbrk{}रणं}}\Afootnote{K H; \skt{मुस्तं गण्डूषः सरसाञ्जनः॥}\ A}}\skt{~।}\pend

\pstart\versepar \ledsidenote{\passref{4.22.12}}\skt{प्रियङ्गवश्च }\edtext{\skt{मुस्ता च}}{\lemma{\skt{मुस्ता च}}\Afootnote{K; \skt{मुस्तं च}\ A; \skt{मुस्ताश्च}\ H}}\skt{ त्रिफला च प्रलेपनं~। }\pend
\pstart\versepar \skt{नस्यञ्च }\edtext{\skt{सिद्धं}}{\lemma{\skt{सिद्धं}}\Afootnote{K; \emph{om.}\ A; \skt{सिद्धिन्}\ H}}\skt{ }\edtext{\skt{त्रिफ\-ला\-म\-धु\-को\-त्प\-ल\-प\-द्मकैः}}{\lemma{\skt{त्रिफ\appbrk{}ला\appbrk{}म\appbrk{}धु\appbrk{}को\appbrk{}त्प\appbrk{}ल\appbrk{}प\appbrk{}द्मकैः}}\Afootnote{K H; \skt{त्रिफलासिद्धं मधु\appbrk{}को\appbrk{}त्प\appbrk{}ल\appbrk{}प\appbrk{}द्मकैः॥}\ A}}\skt{~॥}\pend

\pstart\versepar \ledsidenote{\passref{4.22.13}}\skt{दन्तपुप्पुटके कार्यं तरुणे रक्तमोक्षणं~। }\pend
\pstart\versepar \skt{स }\edtext{\skt{पञ्चलवणं}}{\lemma{\skt{पञ्चलवणं}}\Afootnote{K H; \skt{पञ्चलवणः}\ A}}\skt{ }\edtext{\skt{क्षारं}}{\lemma{\skt{क्षारं}}\Afootnote{K H; \skt{क्षारः}\ A}}\skt{ }\edtext{\skt{सक्षौद्रं}}{\lemma{\skt{सक्षौद्रं}}\Afootnote{K H; \skt{सक्षौद्रः}\ A}}\skt{ प्रतिसारणं}\pend

\pstart\versepar \ledsidenote{\passref{4.22.14}}\edtext{\skt{शिरोविरेकश्च}}{\lemma{\skt{शिरोविरेकश्च}}\Afootnote{K H; \skt{हितः शिरोविरेकश्च}\ A}}\skt{ }\edtext{\skt{हितो}}{\lemma{\skt{हितो}}\Afootnote{K H; \emph{om.}\ A}}\skt{ नस्यं }\edtext{\skt{स्निग्धश्च}}{\lemma{\skt{स्निग्धश्च}}\Afootnote{K; \skt{स्निग्धं च}\ A H}}\skt{ भोजनं~॥ }\pend
\pstart\versepar \skt{विस्राविते दन्तवेष्टे }\edtext{\skt{व्रणां तु}}{\lemma{\skt{व्रणां तु}}\Afootnote{A K; \skt{व्रणन्तु}\ H}}\skt{ प्रतिसारयेत्~।}\pend

\pstart\versepar \ledsidenote{\passref{4.22.15}}\edtext{\skt{लोध्र\-प\-त्र\-ङ्ग\-म\-धु\-क\-ल\-क्षा\-चू\-र्ण्णै\-र्म\-धू\-त्तरैः}}{\lemma{\skt{लोध्र\appbrk{}प\appbrk{}त्र\appbrk{}ङ्ग\appbrk{}म\appbrk{}धु\appbrk{}क\appbrk{}ल\appbrk{}क्षा\appbrk{}चू\appbrk{}र्ण्णै\appbrk{}र्म\appbrk{}धू\appbrk{}त्तरैः}}\Afootnote{\skt{रोध्र\appbrk{}प\appbrk{}त्त\appbrk{}ङ्ग\appbrk{}य\appbrk{}ष्ट्या\appbrk{}ह्व\appbrk{}ला\appbrk{}क्षा\appbrk{}चू\appbrk{}र्णै\appbrk{}र्म\appbrk{}धू\appbrk{}त्तरैः।}\ A; \skt{लोध्र\appbrk{}प\appbrk{}ङ्ग\appbrk{}म\appbrk{}धु\appbrk{}क\appbrk{}ल\appbrk{}क्षा\appbrk{}चू\appbrk{}र्ण्णै\appbrk{}र्म\appbrk{}धू\appbrk{}त्तरैः~।}\ K\textsuperscript{ac}; \skt{लोध्र\appbrk{}प\appbrk{}त्र\appbrk{}ङ्ग\appbrk{}म\appbrk{}धु\appbrk{}क\appbrk{}ल\appbrk{}क्षा\appbrk{}चू\appbrk{}र्ण्णै\appbrk{}र्म\appbrk{}धू\appbrk{}त्तरैः~।}\ K\textsuperscript{pc}; \skt{लोध्र\appbrk{}प\appbrk{}त्र\appbrk{}ङ्ग\appbrk{}म\appbrk{}धु\appbrk{}क\appbrk{}ला\appbrk{}क्षा\appbrk{}चू\appbrk{}र्ण्णै\appbrk{}र्म्म\appbrk{}धू\appbrk{}त्तरैः~।}\ H}}\skt{~। }\pend
\pstart\versepar \skt{गण्डूषे क्षीरिणो योज्याः सक्षौ\-द्र\-घृ\-त\-श\-र्कराः~॥}\pend

\pstart\versepar \ledsidenote{\passref{4.22.16}}\skt{काकोल्यादौ दशक्षीरसिद्धं सर्पिश्च }\edtext{\skt{नस्ततः}}{\lemma{\skt{नस्ततः}}\Afootnote{K H; \skt{नस्यतः।}\ A}}\skt{~॥ }\pend
\pstart\versepar \skt{}\edtext{\skt{सौषि\-रे\-हृ\-त\-रक्ते तु}}{\lemma{\skt{सौषि\appbrk{}रे\appbrk{}हृ\appbrk{}त\appbrk{}रक्ते तु}}\Afootnote{K; \skt{शौषिरे हृतरक्ते तु}\ A; \skt{सौशिरे हृतरक्ते तु}\ H}}\skt{ }\edtext{\skt{लोध्र\-मु\-स्त\-र\-सा\-ञ्जनैः}}{\lemma{\skt{लोध्र\appbrk{}मु\appbrk{}स्त\appbrk{}र\appbrk{}सा\appbrk{}ञ्जनैः}}\Afootnote{K H; \skt{रोध्र\appbrk{}मु\appbrk{}स्त\appbrk{}र\appbrk{}सा\appbrk{}ञ्जनैः॥}\ A}}\skt{~॥}\pend

\pstart\versepar \ledsidenote{\passref{4.22.18}}\skt{क्षीरे }\edtext{\skt{दशगुणा}}{\lemma{\skt{दशगुणा}}\Afootnote{K; \skt{दशगुणे}\ A H}}\skt{ }\edtext{\skt{सर्पिः}}{\lemma{\skt{सर्पिः}}\Afootnote{K H; \skt{सिद्धं}\ A}}\skt{ }\edtext{\skt{सिद्धं}}{\lemma{\skt{सिद्धं}}\Afootnote{K H; \skt{सर्पि}\ A}}\skt{ }\edtext{\skt{नस्ये च}}{\lemma{\skt{नस्ये च}}\Afootnote{K H; \skt{र्नस्ये च}\ A}}\skt{ पूजितं~॥ }\pend
\pstart\versepar \skt{क्रियां परिदरे }\edtext{\skt{कुर्यात्शीतादोक्तं}}{\lemma{\skt{कुर्यात्शीतादोक्तं}}\Afootnote{K H; \skt{कुयाच्छीताकोक्तां}\ A}}\skt{ विचक्षणः~।}\pend

\pstart\versepar \ledsidenote{\passref{4.22.19}}\edtext{\skt{संशोध्योभयत}}{\lemma{\skt{संशोध्योभयत}}\Afootnote{K; \skt{संशोध्योभयतः}\ A H}}\skt{ कार्यं }\edtext{\skt{शिरश्चोपकुशे}}{\lemma{\skt{शिरश्चोपकुशे}}\Afootnote{A K; \skt{शिरश्चोपकुले}\ H}}\skt{ }\edtext{\skt{हितः}}{\lemma{\skt{हितः}}\Afootnote{K H; \skt{तथा।}\ A}}\skt{~। }\pend
\pstart\versepar \skt{}\edtext{\edtext{\skt{काको\-द\-म्व\-रि\-का\-गो\-जी\-प\-त्रै\-र्वि\-स्रा\-व\-ये\-द\-सृक्}}{\lemma{\skt{काको\appbrk{}द\appbrk{}म्व\appbrk{}रि\appbrk{}का\appbrk{}गो\appbrk{}जी\appbrk{}प\appbrk{}त्रैर्}}\Afootnote{K; \skt{काको\appbrk{}दु\appbrk{}म्ब\appbrk{}रि\appbrk{}का\appbrk{}गो\appbrk{}जी\appbrk{}पत्रै}\ A; \skt{काको\appbrk{}दु\appbrk{}म्व\appbrk{}रि\appbrk{}का\appbrk{}गो\appbrk{}जी\appbrk{}प\appbrk{}त्त्रैर्}\ H}}}{\lemma{\skt{विस्रा\appbrk{}व\appbrk{}ये\appbrk{}द\appbrk{}सृक्}}\Afootnote{K H; \skt{र्विस्रा\appbrk{}व\appbrk{}ये\appbrk{}द\appbrk{}सृक्॥}\ A}}\skt{~।}\pend

\pstart\versepar \ledsidenote{\passref{4.22.20}}\skt{क्षौद्रयुक्तैः }\edtext{\skt{सलवणैः}}{\lemma{\skt{सलवणैः}}\Afootnote{K H; \skt{च लवणैः}\ A}}\skt{ सव्योषैः प्रतिसारयेत्~। }\pend
\pstart\versepar \skt{पिप्पल्यः सर्षपा }\edtext{\skt{श्वेताः}}{\lemma{\skt{श्वेताः}}\Afootnote{K; \skt{ञ्श्वेतान्}\ A; \skt{स्वेताः}\ H\textsuperscript{ac} H\textsuperscript{pc}}}\skt{ नागरं नैचुलं फलं~।}\edtext{}{\lemma{\skt{प्रतिसारयेत्~। पिप्पल्यः}}\Afootnote{K H; \skt{प्रतिसारयेत्। पिप्पलीः}\ A}}\pend

\pstart\versepar \ledsidenote{\passref{4.22.21}}\skt{सुखोदकेन संसृज्य }\edtext{\skt{कवडन्तस्य}}{\lemma{\skt{कवडन्तस्य}}\Afootnote{K H; \skt{कवलं चापि}\ A}}\skt{ योजयेत्~॥ }\pend
\pstart\versepar \skt{घृतं मधुरकैः सिद्धं हितं }\edtext{\skt{कवडनस्ययोः}}{\lemma{\skt{कवडनस्ययोः}}\Afootnote{K H; \skt{कवलनस्ययोः॥}\ A}}\skt{~।}\edtext{}{\lemma{\skt{योजयेत्~॥ घृतं}}\Afootnote{K H; \skt{धारयेत्। घृतं}\ A}}\pend

\pstart\versepar \ledsidenote{\passref{4.22.22}}\skt{शस्त्रेण }\edtext{\skt{दन्त\-वै\-द\-र्भ\-द\-न्त\-मू\-लानि}}{\lemma{\skt{दन्त\appbrk{}वै\appbrk{}द\appbrk{}र्भ\appbrk{}द\appbrk{}न्त\appbrk{}मू\appbrk{}लानि}}\Afootnote{K H; \skt{दन्तवैदर्भे दन्तमूलानि}\ A}}\skt{ शोधयेत्~। }\pend
\pstart\versepar \skt{ततःक्षारं प्रयुञ्जीत क्रियाः सर्वाश्च शीतलाः~॥}\edtext{}{\lemma{\skt{शोधयेत्~। ततःक्षारं}}\Afootnote{K H; \skt{शोधयेत्। ततः क्सारं}\ A}}\pend

\pstart\versepar \ledsidenote{\passref{4.22.23}}\edtext{\skt{उद्धृ\-त्या\-धि\-क\-मां\-सन्तु}}{\lemma{\skt{उद्धृ\appbrk{}त्या\appbrk{}धि\appbrk{}क\appbrk{}मां\appbrk{}सन्तु}}\Afootnote{K; \skt{उद्धृ\appbrk{}त्या\appbrk{}धि\appbrk{}क\appbrk{}दन्तं तु}\ A H}}\skt{ ततो }\edtext{\skt{ग्निम\-व\-चा\-र\-येत्}}{\lemma{\skt{ग्निम्}}\Afootnote{K; \skt{ऽग्नि}\ A; \skt{ऽग्निम्}\ H}}\skt{~। }\pend
\pstart\versepar \skt{क्रिमि\-द\-न्त\-क\-व\-च्चात्र }\edtext{\skt{विधिः}}{\lemma{\skt{विधिः}}\Afootnote{A K; \skt{विधि}\ H}}\skt{ कार्यो विजानता~।}\edtext{}{\lemma{\skt{अवचारयेत्~। क्रिमि\appbrk{}द\appbrk{}न्त\appbrk{}क\appbrk{}व\appbrk{}च्चात्र}}\Afootnote{K; \skt{मवचारयेत्~। कृमि\appbrk{}द\appbrk{}न्त\appbrk{}क\appbrk{}व\appbrk{}च्चापि}\ A; \skt{अवचारयेत्~। क्रिमि\appbrk{}द\appbrk{}न्त\appbrk{}क\appbrk{}व\appbrk{}च्चात्र}\ H}}\pend

\pstart\versepar \ledsidenote{\passref{4.22.24}}\edtext{\skt{च्छित्वाधिमांसं}}{\lemma{\skt{च्छित्वाधिमांसं}}\Afootnote{K; \skt{छित्त्वाऽधिमांसं}\ A; \skt{छिन्नाधिमांसं}\ H}}\skt{ }\edtext{\edtext{\skt{सक्षौ\-द्रै\-रि\-मै\-श्चूर्ण्णैः}}{\lemma{\skt{सक्षौद्रैरिमैश्}}\Afootnote{\skt{सक्षौद्रैरेभिश्}\ A; \skt{सक्षौद्रैरिमैश्}\ K; \skt{सक्षौद्रैरिमै}\ H}}}{\lemma{\skt{चूर्ण्णैः}}\Afootnote{A K; \skt{श्चूर्ण्णैः}\ H}}\skt{ समाचरेत्~। }\pend
\pstart\versepar \skt{वचा\-ते\-ज\-व\-ती\-पा\-ठा\-स्व\-र्जि\-का\-या\-व\-शू\-कजैः~।}\edtext{}{\lemma{\skt{समाचरेत्~। वचा\appbrk{}ते\appbrk{}ज\appbrk{}व\appbrk{}ती\appbrk{}पा\appbrk{}ठा\appbrk{}स्व\appbrk{}र्जि\appbrk{}का\appbrk{}या\appbrk{}व\appbrk{}शू\appbrk{}कजैः}}\Afootnote{K; \skt{उपाचरेत्~। वचा\appbrk{}ते\appbrk{}जो\appbrk{}व\appbrk{}ती\appbrk{}पा\appbrk{}ठा\appbrk{}स्व\appbrk{}र्जि\appbrk{}का\appbrk{}या\appbrk{}व\appbrk{}शू\appbrk{}कजैः~॥}\ A; \skt{समाचरेत्~। वचा\appbrk{}ते\appbrk{}जो\appbrk{}व\appbrk{}ती\appbrk{}पा\appbrk{}ठा\appbrk{}स्व\appbrk{}र्ज्जि\appbrk{}का\appbrk{}या\appbrk{}व\appbrk{}शू\appbrk{}कजैः~॥}\ H}}\pend

\pstart\versepar \ledsidenote{\passref{4.22.25}}\edtext{\skt{क्षौद्रद्वितीया}}{\lemma{\skt{क्षौद्रद्वितीया}}\Afootnote{K; \skt{क्षौद्रद्वितीयाः}\ A H}}\skt{ पिप्पल्यः }\edtext{\skt{कव\-ड\-श्चा\-त्र\-की\-र्त्तितः}}{\lemma{\skt{कवडश्}}\Afootnote{K H; \skt{कवलश्}\ A}}\skt{~। }\pend
\pstart\versepar \skt{}\edtext{\skt{पटो\-ल\-नि\-म्व\-त्रि\-फ\-ला\-क\-षा\-या\-श्चात्र}}{\lemma{\skt{पटो\appbrk{}ल\appbrk{}नि\appbrk{}म्व\appbrk{}त्रि\appbrk{}फ\appbrk{}ला\appbrk{}क\appbrk{}षा\appbrk{}याश्}}\Afootnote{K; \skt{पटो\appbrk{}ल\appbrk{}त्रि\appbrk{}फ\appbrk{}ला\appbrk{}नि\appbrk{}म्ब\appbrk{}क\appbrk{}षा\appbrk{}यश्}\ A; \skt{पटो\appbrk{}ल\appbrk{}नि\appbrk{}म्ब\appbrk{}त्रि\appbrk{}फ\appbrk{}ला\appbrk{}क\appbrk{}षा\appbrk{}याश्}\ H}}\skt{ }\edtext{\skt{धावनं}}{\lemma{\skt{धावनं}}\Afootnote{K H; \skt{धावने~।}\ A}}\skt{~। }\pend
\pstart\versepar \skt{}\edtext{\skt{सिरो}}{\lemma{\skt{सिरो}}\Afootnote{K; \skt{हितः}\ A; \emph{om.}\ H}}\skt{ }\edtext{\skt{विरेकश्च}}{\lemma{\skt{विरेकश्च}}\Afootnote{K; \skt{शिरोविरेकश्च}\ A H}}\skt{ }\edtext{\skt{हितो}}{\lemma{\skt{हितो}}\Afootnote{K H; \emph{om.}\ A}}\skt{ }\edtext{\skt{धूमो\-वै\-रे\-च\-नश्च}}{\lemma{\skt{धूमो\appbrk{}वै\appbrk{}रे\appbrk{}च\appbrk{}नश्च}}\Afootnote{K; \skt{धूमो वैरेचनश्च}\ A; \skt{धूमो वैरेचनश्च}\ H}}\skt{ यः~॥}\pend

\pstart \ledsidenote{\passref{4.22.26ab}}\skt{सामान्यं कर्म नाडीनां }\edtext{\skt{विशेषं}}{\lemma{\skt{विशेषं}}\Afootnote{A K; \skt{विशेषश्}\ H}}\skt{ चात्र मे }\edtext{\skt{श्रुणु}}{\lemma{\skt{श्रुणु}}\Afootnote{K; \skt{शृणु}\ A H}}\skt{~॥}\pend

\pstart\versepar \ledsidenote{\passref{4.22.27}}\edtext{\skt{यद्द\-न्त\-म\-धि\-जा\-येत}}{\lemma{\skt{यद्द\appbrk{}न्त\appbrk{}म\appbrk{}धि\appbrk{}जा\appbrk{}येत}}\Afootnote{\skt{यं दन्त\appbrk{}म\appbrk{}भि\appbrk{}जा\appbrk{}येत}\ A; \skt{यद्द\appbrk{}न्त\appbrk{}म\appbrk{}धि\appbrk{}जा\appbrk{}येत}\ K H}}\skt{ }\edtext{\skt{नाडीदन्तं}}{\lemma{\skt{नाडीदन्तं}}\Afootnote{K H; \skt{नाडी तं दन्}\ A}}\skt{ तमुद्धरेत्~। }\pend
\pstart\versepar \skt{छित्वा मांसानि शस्त्रेण यदि नोपरिजो भवेत्~।}\edtext{}{\lemma{\skt{तमुद्धरेत्~। छित्वा}}\Afootnote{A; \skt{तमुद्धरेत्~। छित्वा}\ K H}}\pend

\pstart\versepar \ledsidenote{\passref{4.22.28}}\skt{शोधयित्वा दहेच्चापि क्षारेण ज्वलनेन वा~। }\pend
\pstart\versepar \skt{}\edtext{\skt{उपेक्षिते}}{\lemma{\skt{उपेक्षिते}}\Afootnote{K H; \skt{भिनत्त्युपेक्षिते}\ A}}\skt{ }\edtext{\skt{तु}}{\lemma{\skt{तु}}\Afootnote{K H; \skt{दन्ते}\ A}}\skt{ }\edtext{\skt{दशने}}{\lemma{\skt{दशने}}\Afootnote{K H; \skt{हनुकास्थि}\ A}}\skt{ }\edtext{\skt{गतिर्हन्वस्थि}}{\lemma{\skt{गतिर्हन्वस्थि}}\Afootnote{K; \skt{गतिर्ध्रुवम्~॥}\ A; \skt{गतिहन्वस्थि}\ H}}\skt{ }\edtext{\skt{दारयेत्}}{\lemma{\skt{दारयेत्}}\Afootnote{K H; \emph{om.}\ A}}\skt{~।}\pend

\pstart\versepar \ledsidenote{\passref{4.22.29}}\edtext{\skt{तस्मात्समूलं}}{\lemma{\skt{तस्मात्समूलं}}\Afootnote{K; \skt{समूलं}\ A; \skt{तस्मात्समूल}\ H}}\skt{ }\edtext{\skt{दश\-न\-मु\-द्ध\-रे\-द्भ\-ग्न\-मस्थि च}}{\lemma{\skt{उद्ध\appbrk{}रे\appbrk{}द्भ\appbrk{}ग्न\appbrk{}मस्थि च}}\Afootnote{K H; \skt{तस्मा\appbrk{}दु\appbrk{}द्ध\appbrk{}रे\appbrk{}द्भ\appbrk{}ग्न\appbrk{}म\appbrk{}स्थि\appbrk{}रम्~।}\ A}}\skt{~। }\pend
\pstart\versepar \skt{उद्धृते }\edtext{\skt{तूत्तरे}}{\lemma{\skt{तूत्तरे}}\Afootnote{A K; \skt{तूत्तमे}\ H}}\skt{ दन्ते }\edtext{\skt{शोणितं}}{\lemma{\skt{शोणितं}}\Afootnote{K H; \skt{च स्था}\ A}}\skt{ }\edtext{\skt{संप्रसिच्यते}}{\lemma{\skt{संप्रसिच्यते}}\Afootnote{K H; \skt{नं विद\appbrk{}हे\appbrk{}त्सु\appbrk{}षि\appbrk{}रस्य च~। ततो विदा\appbrk{}री\appbrk{}य\appbrk{}ष्ट्या\appbrk{}ह्व\appbrk{}शृ\appbrk{}ङ्गा\appbrk{}ट\appbrk{}क\appbrk{}क\appbrk{}से\appbrk{}रुकैः~॥}\ A}}\skt{~।}\pend

\pstart \ledsidenote{\passref{4.22.30ab}}\edtext{\skt{रक्ता\-ति\-यो\-गा\-त्पू\-र्वोक्ता}}{\lemma{\skt{रक्ता\appbrk{}ति\appbrk{}यो\appbrk{}गा\appbrk{}त्पू\appbrk{}र्वोक्ता}}\Afootnote{K H; \skt{रक्ता\appbrk{}ति\appbrk{}यो\appbrk{}गा\appbrk{}त्पू\appbrk{}र्वेक्ता}\ A}}\skt{ रोगा घोरा भवन्ति }\edtext{\skt{च}}{\lemma{\skt{च}}\Afootnote{K H; \skt{हि।}\ A}}\skt{~।}\pend

\pstart \ledsidenote{\passref{4.22.31ab}}\edtext{\skt{चल\-म\-प्यु\-त्त\-म\-न्द\-न्त\-मतो}}{\lemma{\skt{उत्तमन्}}\Afootnote{K H; \skt{उत्तरं}\ A}}\skt{ नाप\-ह\-रे\-द्भि\-षक्~।}\pend

\pstart\versepar \ledsidenote{\passref{4.22.32}}\edtext{\skt{कषा\-य\-जा\-ती\-म\-द\-न\-क\-टु\-क\-स्वा\-दु\-क\-ण्टकैः}}{\lemma{\skt{कषा\appbrk{}य\appbrk{}जा\appbrk{}ती\appbrk{}म\appbrk{}द\appbrk{}न\appbrk{}क\appbrk{}टु\appbrk{}क\appbrk{}स्वा\appbrk{}दु\appbrk{}क\appbrk{}ण्टकैः}}\Afootnote{K H; \skt{कषायं जाति\appbrk{}म\appbrk{}द\appbrk{}न\appbrk{}क\appbrk{}टु\appbrk{}क\appbrk{}स्वा\appbrk{}दु\appbrk{}क\appbrk{}ण्टकैः।}\ A}}\skt{~। }\pend
\pstart\versepar \skt{}\edtext{\skt{लोध्र\-म\-ञ्जि\-ष्ठ\-ख\-दि\-र\-य\-ष्ट्या\-ह्वै\-श्चापि}}{\lemma{\skt{लोध्र\appbrk{}म\appbrk{}ञ्जि\appbrk{}ष्ठ\appbrk{}ख\appbrk{}दि\appbrk{}र\appbrk{}य\appbrk{}ष्ट्या\appbrk{}ह्वैश्}}\Afootnote{K H; \skt{यष्ट्या\appbrk{}ह्व\appbrk{}रो\appbrk{}ध्र\appbrk{}म\appbrk{}ञ्जि\appbrk{}ष्ठा\appbrk{}ख\appbrk{}दि\appbrk{}रैश्}\ A}}\skt{ यत्कृतं~।}\pend

\pstart\versepar \ledsidenote{\passref{4.22.33}}\skt{तैलं }\edtext{\skt{संशोधने}}{\lemma{\skt{संशोधने}}\Afootnote{K H; \skt{संशोधनं}\ A}}\skt{ }\edtext{\skt{तद्धि\-ह\-न्या\-द्द\-न्ति\-गतां}}{\lemma{\skt{तद्धि\appbrk{}ह\appbrk{}न्या\appbrk{}द्द\appbrk{}न्ति\appbrk{}गतां}}\Afootnote{K; \skt{तद्धि हन्याद्दन्तगतां}\ A H}}\skt{ गतिं~॥ }\pend
\pstart\versepar \skt{कीर्त्तिता }\edtext{\skt{दन्तमूलेषु}}{\lemma{\skt{दन्तमूलेषु}}\Afootnote{K H; \skt{दन्तमूले तु}\ A}}\skt{ क्रिया दन्तेषु वक्ष्यते~।}\pend

\pstart\versepar \ledsidenote{\passref{4.22.34}}\edtext{\skt{सुखोष्णः}}{\lemma{\skt{सुखोष्णः}}\Afootnote{K H; \emph{om.}\ A}}\skt{ }\edtext{\skt{स्नेहकवडाः}}{\lemma{\skt{स्नेहकवडाः}}\Afootnote{K H; \skt{स्नेहानां कवलाः}\ A}}\skt{ }\edtext{\skt{सर्पि\-ष\-स्त्रि\-वृ\-तस्य वा}}{\lemma{\skt{सर्पि\appbrk{}ष\appbrk{}स्त्रि\appbrk{}वृ\appbrk{}तस्य वा}}\Afootnote{K H; \skt{कोष्णाः सर्पि\appbrk{}ष\appbrk{}स्त्रै\appbrk{}वृ\appbrk{}तस्य वा।}\ A}}\skt{~। }\pend
\pstart\versepar \skt{}\edtext{\skt{निर्यू\-ह\-श्चा\-नि\-ल\-घ्नानां}}{\lemma{\skt{निर्यूहश्}}\Afootnote{K H; \skt{निर्यूहाश्}\ A}}\skt{ दन्त\-ह\-र्ष\-प्र\-म\-र्दनाः~।}\pend

\pstart\versepar \ledsidenote{\passref{4.22.35}}\skt{स्नैहिकश्च हितो धूमो नस्यं }\edtext{\edtext{\skt{स्नैहिकमेव}}{\lemma{\skt{स्नैहिकम्}}\Afootnote{K H; \skt{स्निग्धं च}\ A}}\skt{ च}}{\lemma{\skt{एव च}}\Afootnote{K H; \skt{भोजनम्।}\ A}}\skt{~। }\pend
\pstart\versepar \skt{}\edtext{\skt{रसा}}{\lemma{\skt{रसा}}\Afootnote{K H; \skt{रसो}\ A}}\skt{ }\edtext{\skt{रसै}}{\lemma{\skt{रसै}}\Afootnote{K; \skt{रस}\ A; \skt{रसे}\ H}}\skt{ यवाग्वश्च }\edtext{\skt{क्षीरसन्तानिक}}{\lemma{\skt{क्षीरसन्तानिक}}\Afootnote{K; \skt{क्षीरं सन्तानिका}\ A; \skt{क्षीरसन्तानिका}\ H}}\skt{ घृतं~॥}\pend

\pstart \ledsidenote{\passref{4.22.36cd}}\skt{अहिं\-स\-न्द\-न्त\-मू\-लानि शर्क\-रा\-मु\-द्ध\-रे\-द्भि\-षक्~।}\pend

\pstart\versepar \ledsidenote{\passref{4.22.37}}\edtext{\edtext{\edtext{\skt{लाक्षा\-चू\-र्ण्णै\-र्म\-धु\-यु\-तै\-स्त\-तस्तं}}{\lemma{\skt{लाक्षाचूर्ण्णैर्}}\Afootnote{K H; \skt{लाक्षाचूर्णै}\ A}}}{\lemma{\skt{मधुयुतैस्}}\Afootnote{K H; \skt{र्मधुयुतैस्}\ A}}}{\lemma{\skt{ततस्तं}}\Afootnote{K H; \skt{ततस्ताः}\ A}}\skt{ प्रतिसारयेत्~। }\pend
\pstart\versepar \skt{दन्त }\edtext{\skt{हर्षक्रियाश्चापि}}{\lemma{\skt{हर्षक्रियाश्}}\Afootnote{K; \skt{हर्षक्रियां}\ A H}}\skt{ कुर्या\-न्नि\-र\-व\-शे\-षतः~।}\pend

\pstart\versepar \ledsidenote{\passref{4.22.38}}\edtext{\skt{कपालकः}}{\lemma{\skt{कपालकः}}\Afootnote{K H; \skt{कपालिका}\ A}}\skt{ }\edtext{\skt{कृच्छ्रतमः}}{\lemma{\skt{कृच्छ्रतमः}}\Afootnote{K H; \skt{कृच्छ्रतमा}\ A}}\skt{ }\edtext{\skt{तत्राप्येषा}}{\lemma{\skt{तत्राप्येषा}}\Afootnote{K H; \skt{तत्राप्येषा}\ A}}\skt{ क्रिया हिता~। }\pend
\pstart\versepar \skt{}\edtext{\skt{जायेद्विस्रावणैः}}{\lemma{\skt{जायेद्विस्रावणैः}}\Afootnote{K; \skt{जयेद्विस्रावणैः}\ A H}}\skt{ }\edtext{\skt{स्विन्नमचलं}}{\lemma{\skt{स्विन्नमचलं}}\Afootnote{K; \skt{स्विन्नमचलं}\ A H}}\skt{ }\edtext{\skt{क्रिमिन्दकं}}{\lemma{\skt{क्रिमिन्दकं}}\Afootnote{K; \skt{कृमिदन्तकम्॥}\ A; \skt{क्रिमिदन्तकं~॥}\ H}}\skt{~।}\pend

\pstart\versepar \ledsidenote{\passref{4.22.39}}\edtext{\skt{तथा\-व\-पी\-डै\-र्वा\-तघ्नैः}}{\lemma{\skt{तथावपीडैर्}}\Afootnote{K H; \skt{तथा+\appbrk{}अ\appbrk{}व\appbrk{}पी\appbrk{}डैर्}\ A}}\skt{ स्नेह\-ग\-ण्डू\-ष\-धा\-रणैः~। }\pend
\pstart\versepar \skt{}\edtext{\skt{भद्र\-दा\-र्वा\-दि\-व\-र्षा\-भू\-र्लेपैः}}{\lemma{\skt{भद्र\appbrk{}दा\appbrk{}र्वा\appbrk{}दि\appbrk{}व\appbrk{}र्षा\appbrk{}भूर्}}\Afootnote{K H; \skt{भद्र\appbrk{}दा\appbrk{}र्वा\appbrk{}दि\appbrk{}व\appbrk{}र्षाभू}\ A}}\skt{ स्निग्धैश्च भोजनैः~।}\pend

\pstart\versepar \ledsidenote{\passref{4.22.40}}\edtext{\edtext{\skt{चलमुद्धृत्य}}{\lemma{\skt{चलम्}}\Afootnote{K H; \skt{चल}\ A}}\skt{ वा}}{\lemma{\skt{उद्धृत्य वा}}\Afootnote{K H; \skt{मुद्धृत्य च}\ A}}\skt{ स्थानं }\edtext{\skt{दहेत}}{\lemma{\skt{दहेत}}\Afootnote{K H; \skt{विदहे}\ A}}\skt{ }\edtext{\skt{सुषिरस्य वा}}{\lemma{\skt{सुषिरस्य वा}}\Afootnote{K H; \skt{त्सुषिरस्य च।}\ A}}\skt{~। }\pend
\pstart\versepar \skt{ततो विदारी }\edtext{\skt{यष्ट्या\-ह्व\-शृ\-ङ्गा\-ट\-क\-क\-से\-रुभिः}}{\lemma{\skt{यष्ट्या\appbrk{}ह्व\appbrk{}शृ\appbrk{}ङ्गा\appbrk{}ट\appbrk{}क\appbrk{}क\appbrk{}से\appbrk{}रुभिः}}\Afootnote{K; \skt{यष्ट्या\appbrk{}ह्व\appbrk{}शृ\appbrk{}ङ्गा\appbrk{}ट\appbrk{}क\appbrk{}क\appbrk{}से\appbrk{}रुकैः॥}\ A; \skt{यष्ट्या\appbrk{}ह्व\appbrk{}शृ\appbrk{}ङ्गा\appbrk{}ट\appbrk{}क\appbrk{}क\appbrk{}शे\appbrk{}रुभिः~॥}\ H}}\skt{~।}\pend

\pstart\versepar \ledsidenote{\passref{4.22.41}}\skt{तैलं }\edtext{\skt{दश\-गु\-ण\-क्षी\-र\-सिद्धं}}{\lemma{\skt{दश\appbrk{}गु\appbrk{}ण\appbrk{}क्षी\appbrk{}र\appbrk{}सिद्धं}}\Afootnote{K; \skt{दशगुणे क्षीरे सिद्धं}\ A; \skt{दशगुणं क्षीरं सिद्धं}\ H}}\skt{ }\edtext{\skt{नस्य च}}{\lemma{\skt{नस्य च}}\Afootnote{K; \skt{नस्ये}\ A; \skt{नस्ये च}\ H}}\skt{ }\edtext{\skt{पूजितं}}{\lemma{\skt{पूजितं}}\Afootnote{K H; \skt{हितं}\ A}}\skt{~। }\pend
\pstart\versepar \skt{}\edtext{\skt{हनु\-र्मो\-क्ष\-स\-मु\-द्दिष्टां}}{\lemma{\skt{हनु\appbrk{}र्मो\appbrk{}क्ष\appbrk{}स\appbrk{}मु\appbrk{}द्दिष्टां}}\Afootnote{K; \skt{भवेत्~। हनोमोक्षे समुद्दिष्टां}\ A; \skt{हनु\appbrk{}र्म्मो\appbrk{}क्ष\appbrk{}स\appbrk{}मु\appbrk{}द्धिष्टां}\ H}}\skt{ कुर्या\-च्चा\-र्दि\-त\-व\-त्क्रियां~।}\pend

\pstart\versepar \ledsidenote{\passref{4.22.42}}\edtext{\skt{फलान्यम्व्लानि}}{\lemma{\skt{फलान्यम्व्लानि}}\Afootnote{K; \skt{फलान्यम्लानि}\ A H}}\skt{ }\edtext{\skt{शीताम्वु}}{\lemma{\skt{शीताम्वु}}\Afootnote{K; \skt{शीताम्बु}\ A H}}\skt{ }\edtext{\skt{रूक्षान्नं}}{\lemma{\skt{रूक्षान्नं}}\Afootnote{A K; \skt{रूक्षान्न}\ H}}\skt{ दन्तधावनं~। }\pend
\pstart\versepar \skt{}\edtext{\skt{तथातिकठिनां}}{\lemma{\skt{तथातिकठिनां}}\Afootnote{K H; \skt{तथाऽतिकठिना}\ A}}\skt{ }\edtext{\skt{भक्षां}}{\lemma{\skt{भक्षां}}\Afootnote{K H; \skt{न्भक्ष्यान्}\ A}}\skt{ दन्तरोगी विवर्जयेत्~॥}\pend

\pstart\versepar \ledsidenote{\passref{4.22.43}}\skt{साध्यानां }\edtext{\edtext{\skt{दन्त\-रो\-गा\-णा\-ञ्चि\-कि\-त्सि\-त\-मु\-दा\-हृतं}}{\lemma{\skt{दन्त\appbrk{}रो\appbrk{}गा\appbrk{}णा\appbrk{}ञ्चि\appbrk{}कि\appbrk{}त्सि\appbrk{}तम्}}\Afootnote{K; \skt{दन्तरोगाणां चिकित्सित}\ A; \skt{दन्त\appbrk{}रो\appbrk{}गा\appbrk{}ना\appbrk{}ञ्चि\appbrk{}कि\appbrk{}त्सि\appbrk{}तम्}\ H}}}{\lemma{\skt{उदाहृतं}}\Afootnote{K H; \skt{मुदीरितम्~।}\ A}}\skt{~। }\pend
\pstart\versepar \skt{जिह्वागतानां साध्यानां कर्म वक्ष्यामि }\edtext{\skt{सिद्धयेत्}}{\lemma{\skt{सिद्धयेत्}}\Afootnote{K; \skt{सिद्धये}\ A H}}\skt{~॥}\pend

\pstart\versepar \ledsidenote{\passref{4.22.44}}\edtext{\skt{औष्ठ\-प्र\-को\-पे\-नि\-लजे}}{\lemma{\skt{औष्ठ\appbrk{}प्र\appbrk{}को\appbrk{}पे\appbrk{}नि\appbrk{}लजे}}\Afootnote{K; \skt{औष्ठ\appbrk{}प्र\appbrk{}को\appbrk{}पेऽ\appbrk{}नि\appbrk{}लजे}\ A H}}\skt{ यदुक्तं प्राक्चिकित्सितं~। }\pend
\pstart\versepar \skt{कण्ट\-के\-ष्व\-नि\-लो\-त्थेषु तत्कार्यं भिषजा भवेत्~।}\pend

\pstart\versepar \ledsidenote{\passref{4.22.45}}\skt{पित्तजेषु विघृष्टेषु }\edtext{\skt{निसृते}}{\lemma{\skt{निसृते}}\Afootnote{K H; \skt{निःसृते}\ A}}\skt{ दुष्टशोणिते~। }\pend
\pstart\versepar \skt{}\edtext{\skt{प्रति\-सा\-र\-ण\-ग\-ण्डूषा}}{\lemma{\skt{प्रति\appbrk{}सा\appbrk{}र\appbrk{}ण\appbrk{}ग\appbrk{}ण्डूषा}}\Afootnote{K H; \skt{प्रति\appbrk{}सा\appbrk{}र\appbrk{}ण\appbrk{}ग\appbrk{}ण्डूषं}\ A}}\skt{ नस्यञ्च मधुरं हितं~।}\pend

\pstart\versepar \ledsidenote{\passref{4.22.46}}\skt{कण्टकेषु कफोत्थेषु }\edtext{\edtext{\skt{लिखितेष्वसृजः}}{\lemma{\skt{लिखितेष्व्}}\Afootnote{A K; \skt{लिखिते}\ H}}}{\lemma{\skt{असृजः}}\Afootnote{A K; \skt{ष्वसृजः}\ H}}\skt{ क्षये~। }\pend
\pstart\versepar \skt{}\edtext{\edtext{\skt{पिप्प\-ल्या\-दि\-र्म\-धु\-युतः}}{\lemma{\skt{पिप्पल्यादिर्}}\Afootnote{K H; \skt{पिप्पल्यादि}\ A}}}{\lemma{\skt{मधुयुतः}}\Afootnote{K H; \skt{र्मधुयुतः}\ A}}\skt{ कार्यस्तु }\edtext{\skt{प्रतिसारणं}}{\lemma{\skt{प्रतिसारणं}}\Afootnote{K H; \skt{प्रतिसारणे~॥}\ A}}\skt{~।}\pend

\pstart\versepar \ledsidenote{\passref{4.22.47}}\edtext{\skt{कृह्णीयात्कवलं}}{\lemma{\skt{कृह्णीयात्कवलं}}\Afootnote{K\textsuperscript{pc}; \skt{गृह्णी\appbrk{}या\appbrk{}त्क\appbrk{}व\appbrk{}लांश्}\ A; \skt{कृह्णीया कवलं}\ K\textsuperscript{ac}; \skt{गृह्णी\appbrk{}या\appbrk{}त्क\appbrk{}व\appbrk{}लम्}\ H}}\skt{ }\edtext{\skt{चापि}}{\lemma{\skt{चापि}}\Afootnote{A K; \skt{वापि}\ H}}\skt{ गौर\-स\-र्ष\-प\-सै\-न्धवैः~। }\pend
\pstart\versepar \skt{}\edtext{\skt{पटो\-ल\-नि\-म्व\-वा\-र्त्ता\-कू\-क्षा\-र\-यू\-षैश्च}}{\lemma{\skt{पटो\appbrk{}ल\appbrk{}नि\appbrk{}म्व\appbrk{}वा\appbrk{}र्त्ता\appbrk{}कू\appbrk{}क्षा\appbrk{}र\appbrk{}यू\appbrk{}षैश्च}}\Afootnote{K; \skt{पटो\appbrk{}ल\appbrk{}नि\appbrk{}म्ब\appbrk{}वा\appbrk{}र्ता\appbrk{}कु\appbrk{}क्षा\appbrk{}र\appbrk{}यू\appbrk{}षैश्च}\ A H}}\skt{ भोजयेत्~।}\pend

\pstart\versepar \ledsidenote{\passref{4.22.48}}\skt{उपजिह्वान्तु संलिख्य क्षारेण प्रतिसारयेत्~। }\pend
\pstart\versepar \skt{शिरो\-वि\-रे\-क\-ग\-ण्डूष }\edtext{\skt{धूमै\-श्चै\-ना\-मु\-पा\-च\-रेत्}}{\lemma{\skt{धूमै\appbrk{}श्चै\appbrk{}ना\appbrk{}मु\appbrk{}पा\appbrk{}च\appbrk{}रेत्}}\Afootnote{K; \skt{धूमै\appbrk{}श्चै\appbrk{}न\appbrk{}मु\appbrk{}पा\appbrk{}च\appbrk{}रेत्~॥}\ A; \skt{धूमै\appbrk{}श्चै\appbrk{}ना\appbrk{}मु\appbrk{}पा\appbrk{}च\appbrk{}रेत्~॥}\ H}}\skt{~॥}\pend

\pstart\versepar \ledsidenote{\passref{4.22.49}}\skt{जिह्वागतानां कर्मोक्तं तालव्यानां प्रवक्ष्यते~। }\pend
\pstart\versepar \skt{}\edtext{\edtext{\skt{अङ्गु\-ष्ठां\-गु\-लि\-सं\-दं\-शै\-र्गृ\-हीत्वा}}{\lemma{\skt{अङ्गु\appbrk{}ष्ठां\appbrk{}गु\appbrk{}लि\appbrk{}सं\appbrk{}दं\appbrk{}शैर्}}\Afootnote{K H; \skt{अङ्गु\appbrk{}ष्ठा\appbrk{}ङ्गु\appbrk{}लि\appbrk{}सं\appbrk{}दंशे}\ A}}}{\lemma{\skt{गृहीत्वा}}\Afootnote{K H; \skt{नाकृष्य}\ A}}\skt{ गलशुण्डिकां}\pend

\pstart \ledsidenote{\passref{4.22.50ab}}\edtext{\skt{च्छेद\-ये\-त्म\-ण्ड\-ला\-ग्रेण}}{\lemma{\skt{च्छेद\appbrk{}ये\appbrk{}त्म\appbrk{}ण्ड\appbrk{}ला\appbrk{}ग्रेण}}\Afootnote{K H; \skt{छेद\appbrk{}ये\appbrk{}न्म\appbrk{}ण्ड\appbrk{}ला\appbrk{}ग्रेण}\ A}}\skt{ जिह्वोपरि तु }\edtext{\skt{यत्स्थितं}}{\lemma{\skt{यत्स्थितं}}\Afootnote{K\textsuperscript{pc} H; \skt{संस्थिताम्~।}\ A; \skt{त्स्थितं}\ K\textsuperscript{ac}}}\pend

\pstart\versepar \ledsidenote{\passref{4.22.51}}\edtext{\skt{अत्या\-दा\-ना\-त्स्र\-वे\-द्रक्तं}}{\lemma{\skt{अत्या\appbrk{}दा\appbrk{}ना\appbrk{}त्स्र\appbrk{}वे\appbrk{}द्रक्तं}}\Afootnote{K\textsuperscript{pc} H; \skt{अत्या\appbrk{}दा\appbrk{}ना\appbrk{}त्स्र\appbrk{}वे\appbrk{}द्रक्तं}\ A; \skt{अत्या\appbrk{}ना\appbrk{}त्स्र\appbrk{}वे\appbrk{}द्रक्तं}\ K\textsuperscript{ac}}}\skt{ तन्निमित्तं }\edtext{\skt{म्रियेत सः}}{\lemma{\skt{म्रियेत सः}}\Afootnote{K H; \skt{म्रियेत च~।}\ A}}\skt{~। }\pend
\pstart\versepar \skt{}\edtext{\skt{हीन\-च्छे\-दा\-द्भ\-वे\-च्छोफो}}{\lemma{\skt{हीन\appbrk{}च्छे\appbrk{}दा\appbrk{}द्भ\appbrk{}वे\appbrk{}च्छोफो}}\Afootnote{K H; \skt{हीन\appbrk{}च्छे\appbrk{}दा\appbrk{}द्भ\appbrk{}वे\appbrk{}च्छोफो}\ A}}\skt{ लाला }\edtext{\skt{स्रावो}}{\lemma{\skt{स्रावो}}\Afootnote{K H; \skt{निद्रा}\ A}}\skt{ }\edtext{\skt{भ्रमस्तमः}}{\lemma{\skt{भ्रमस्तमः}}\Afootnote{\skt{भ्रमस्तमः}\ A K H}}\skt{~।}\pend

\pstart\versepar \ledsidenote{\passref{4.22.52}}\skt{तस्माद्वैद्यः प्रयत्नेन दृष्टकर्मा }\edtext{\skt{विशारदः}}{\lemma{\skt{विशारदः}}\Afootnote{A K; \skt{विसारदः~॥}\ H}}\skt{~। }\pend
\pstart\versepar \skt{}\edtext{\skt{गलशुण्डिकां}}{\lemma{\skt{गलशुण्डिकां}}\Afootnote{K H; \skt{गलशुण्डीं}\ A}}\skt{ }\edtext{\skt{च्छेदयित्वा}}{\lemma{\skt{च्छेदयित्वा}}\Afootnote{K H; \skt{तु सञ्छिद्य}\ A}}\skt{ }\edtext{\skt{कुर्यात्प्राप्तमिदं}}{\lemma{\skt{कुर्यात्प्राप्तमिदं}}\Afootnote{K H; \skt{कुर्यात्प्राप्तमिमं}\ A}}\skt{ क्रमं~।}\pend

\pstart\versepar \ledsidenote{\passref{4.22.53}}\skt{मरि\-चा\-ति\-वि\-षा\-पा\-ठा\-व\-चा\-कु\-ष्ठ\-कु\-ट\-न्नटैः~। }\pend
\pstart\versepar \skt{क्षौद्रयुक्तैः }\edtext{\skt{सल\-व\-णै\-स्त\-तस्तां}}{\lemma{\skt{ततस्तां}}\Afootnote{A K; \skt{तततां}\ H}}\skt{ प्रतिसारयेत्~॥}\pend

\pstart\versepar \ledsidenote{\passref{4.22.54}}\edtext{\edtext{\skt{वचा\-म\-ति\-वि\-षा\-म्पाठं}}{\lemma{\skt{वचा\appbrk{}म\appbrk{}ति\appbrk{}वि\appbrk{}षाम्}}\Afootnote{A K; \skt{वचा\appbrk{}म\appbrk{}ति\appbrk{}वि\appbrk{}षाम्}\ H}}}{\lemma{\skt{पाठं}}\Afootnote{K H; \skt{पाठां}\ A}}\skt{ रास्नां कटुकरोहिणीं~। }\pend
\pstart\versepar \skt{}\edtext{\skt{निःक्वा\-थ्य\-पि\-चु\-म\-र्दञ्च}}{\lemma{\skt{निःक्वा\appbrk{}थ्य\appbrk{}पि\appbrk{}चु\appbrk{}म\appbrk{}र्दञ्च}}\Afootnote{K H; \skt{निष्क्वाथ्य पिचुमन्दं च}\ A}}\skt{ }\edtext{\skt{कवडन्तत्र}}{\lemma{\skt{कवडन्}}\Afootnote{K H; \skt{कवलं}\ A}}\skt{ योजयेत्~।}\pend

\pstart\versepar \ledsidenote{\passref{4.22.55}}\edtext{\skt{इङ्गु\-दी\-कि\-णि\-ही\-द\-न्ती\-स\-रलं}}{\lemma{\skt{इङ्गु\appbrk{}दी\appbrk{}कि\appbrk{}णि\appbrk{}ही\appbrk{}द\appbrk{}न्ती\appbrk{}स\appbrk{}रलं}}\Afootnote{K H; \skt{इङ्गु\appbrk{}दी\appbrk{}कि\appbrk{}णि\appbrk{}ही\appbrk{}द\appbrk{}न्ती\appbrk{}स\appbrk{}रला}\ A}}\skt{ }\edtext{\skt{देवदारु च}}{\lemma{\skt{देवदारु च}}\Afootnote{K H; \skt{सुरदारुभिः~।}\ A}}\skt{~। }\pend
\pstart\versepar \skt{}\edtext{\edtext{\edtext{\skt{पञ्चा\-ङ्गी\-ङ्का\-र\-ये\-द्व\-र्त्ति\-मे\-तै\-र्ग\-न्धो\-त्तरां}}{\lemma{\skt{पञ्चाङ्गीङ्}}\Afootnote{A K; \skt{पञ्चाङ्गी}\ H}}}{\lemma{\skt{कारयेद्वर्त्तिम्}}\Afootnote{K H; \skt{कारयेत्पिष्टैर्}\ A}}}{\lemma{\skt{एतैर्}}\Afootnote{K H; \skt{वर्तिं}\ A}}\skt{ शुभां~।}\pend

\pstart\versepar \ledsidenote{\passref{4.22.56}}\edtext{\skt{तस्या}}{\lemma{\skt{तस्या}}\Afootnote{K H; \skt{ततो}\ A}}\skt{ धूमं }\edtext{\skt{पिवेज्जन्तु}}{\lemma{\skt{पिवेज्जन्तु}}\Afootnote{K; \skt{पिबेज्जन्तु}\ A; \skt{पिबेज्जन्तुर्}\ H}}\skt{ }\edtext{\skt{द्विरह्नः}}{\lemma{\skt{द्विरह्नः}}\Afootnote{K H; \skt{र्द्विरह्नः}\ A}}\skt{ कफनाशनं~। }\pend
\pstart\versepar \skt{क्षारसिद्धेषु मुद्गेषु }\edtext{\edtext{\skt{यूषश्चाप्यशने}}{\lemma{\skt{चाप्य्}}\Afootnote{A H; \skt{चा}\ K}}}{\lemma{\skt{अशने}}\Afootnote{A H; \skt{प्यशने}\ K}}\skt{ हितः~॥}\pend

\pstart\versepar \ledsidenote{\passref{4.22.57}}\edtext{\skt{तुण्डिकेर्यध्रुषो}}{\lemma{\skt{तुण्डिकेर्यध्रुषो}}\Afootnote{K; \skt{तुण्डिकेर्यध्रुषे}\ A H}}\skt{ कूर्मे संघाते }\edtext{\skt{तालुपुप्पुटे}}{\lemma{\skt{तालुपुप्पुटे}}\Afootnote{A K; \skt{तुलुपुप्पुटे~॥}\ H}}\skt{~। }\pend
\pstart\versepar \skt{एष एव विधिः }\edtext{\skt{क्रायो}}{\lemma{\skt{क्रायो}}\Afootnote{K; \skt{कार्यो}\ A H}}\skt{ विशेषः शस्त्रकर्मणि~।}\pend

\pstart\versepar \ledsidenote{\passref{4.22.58}}\skt{तालुपाके तु कर्तव्यं }\edtext{\skt{विधाणं}}{\lemma{\skt{विधाणं}}\Afootnote{K; \skt{विधानं}\ A H}}\skt{ पित्तनाशनं~। }\pend
\pstart\versepar \skt{स्नेहस्वेदौ तालुशोषे विधि\-श्चा\-नि\-ल\-ना\-शनः~।}\pend

\pstart\versepar \ledsidenote{\passref{4.22.59}}\skt{कीर्त्तितं तालु जानान्तु कण्ठ्यानां कर्म वक्ष्यते~॥ }\pend
\pstart\versepar \skt{साध्यानां रोहिणीनां तु हितं शोणितमोक्षणं~।}\pend

\pstart\versepar \ledsidenote{\passref{4.22.60}}\skt{च्छर्दनं }\edtext{\skt{धूमपानञ्च}}{\lemma{\skt{धूमपानञ्च}}\Afootnote{A K; \skt{धूम्रपानञ्च}\ H}}\skt{ गण्डूषो }\edtext{\skt{नस्तकर्म च}}{\lemma{\skt{नस्तकर्म च}}\Afootnote{K H; \skt{नस्यकर्म च~।}\ A}}\skt{~। }\pend
\pstart\versepar \skt{}\edtext{\skt{वातिकां तु}}{\lemma{\skt{वातिकां तु}}\Afootnote{K H; \skt{वातिकीं तु}\ A}}\skt{ }\edtext{\skt{गते}}{\lemma{\skt{गते}}\Afootnote{K H; \skt{हृते}\ A}}\skt{ रक्ते लवणैः प्रतिसारयेत्~।}\pend

\pstart\versepar \ledsidenote{\passref{4.22.61}}\edtext{\skt{सुखोष्णां}}{\lemma{\skt{सुखोष्णां}}\Afootnote{K H; \skt{सुखोष्णा}\ A}}\skt{ }\edtext{\skt{स्नेहकवलां}}{\lemma{\skt{स्नेहकवलां}}\Afootnote{K H; \skt{न्स्नेहगण्डूषान्}\ A}}\skt{ }\edtext{\edtext{\skt{धार\-ये\-च्चा\-प्य\-भी\-क्ष्णशः}}{\lemma{\skt{धारयेच्चाप्य्}}\Afootnote{A H; \skt{धारयेच्चा}\ K}}}{\lemma{\skt{अभीक्ष्णशः}}\Afootnote{A H; \skt{प्यभीक्ष्णशः~।}\ K}}\skt{~। }\pend
\pstart\versepar \skt{पत्त\-ङ्ग\-श\-र्क\-रा\-क्षौद्रैः }\edtext{\skt{पैत्तिकां}}{\lemma{\skt{पैत्तिकां}}\Afootnote{K; \skt{पैत्तिकीं}\ A; \skt{पैत्तिकं}\ H}}\skt{ प्रतिसारयेत्~।}\pend

\pstart\versepar \ledsidenote{\passref{4.22.62}}\edtext{\skt{द्राक्षा\-प\-रू\-ष\-क\-क्वाथो}}{\lemma{\skt{द्राक्षा\appbrk{}प\appbrk{}रू\appbrk{}ष\appbrk{}क\appbrk{}क्वाथो}}\Afootnote{A K; \skt{द्राक्षा\appbrk{}फ\appbrk{}रू\appbrk{}ष\appbrk{}क\appbrk{}क्वाथो}\ H}}\skt{ }\edtext{\skt{हितश्च}}{\lemma{\skt{हितश्च}}\Afootnote{A K; \skt{हितञ्च}\ H}}\skt{ }\edtext{\skt{कवडग्रहे}}{\lemma{\skt{कवडग्रहे}}\Afootnote{K H; \skt{कवलग्रहे~।}\ A}}\skt{~। }\pend
\pstart\versepar \skt{}\edtext{\skt{अगा\-र\-धू\-म\-म\-धुकैः}}{\lemma{\skt{अगा\appbrk{}र\appbrk{}धू\appbrk{}म\appbrk{}म\appbrk{}धुकैः}}\Afootnote{K H; \skt{अगा\appbrk{}र\appbrk{}धू\appbrk{}म\appbrk{}क\appbrk{}टुकैः}\ A}}\skt{ }\edtext{\skt{कफजां}}{\lemma{\skt{कफजां}}\Afootnote{K; \skt{श्लैष्मिकीं}\ A; \skt{कफजम्}\ H}}\skt{ प्रतिसारयेत्~।}\pend

\pstart\versepar \ledsidenote{\passref{4.22.63}}\edtext{\skt{श्वेता}}{\lemma{\skt{श्वेता}}\Afootnote{A K; \skt{स्वेता}\ H}}\skt{ विडङ्गदन्तीषु तैलं सिद्धं ससैन्धवं~। }\pend
\pstart\versepar \skt{नस्यकर्मणि }\edtext{\skt{दातव्यं}}{\lemma{\skt{दातव्यं}}\Afootnote{K H; \skt{योक्तव्यं}\ A}}\skt{ }\edtext{\skt{कवडण्च}}{\lemma{\skt{कवडण्च}}\Afootnote{K H; \skt{तथा कव}\ A}}\skt{ }\edtext{\skt{कफोच्छ्रये}}{\lemma{\skt{कफोच्छ्रये}}\Afootnote{K H; \skt{लधारणे~॥}\ A}}\skt{~।}\pend

\pstart\versepar \ledsidenote{\passref{4.22.64}}\skt{पित्त\-व\-त्सा\-ध\-ये\-द्वैद्यो रोहिणीं रक्तसम्भवां~। }\pend
\pstart\versepar \skt{विस्राव्य }\edtext{\skt{कण्ठसालूकं}}{\lemma{\skt{कण्ठसालूकं}}\Afootnote{K H; \skt{कण्ठशालूकं}\ A}}\skt{ साध\-ये\-त्तु\-ण्डि\-के\-रि\-वत्~।}\pend

\pstart\versepar \ledsidenote{\passref{4.22.65}}\skt{एककालं यवान्नञ्च भुञ्जीत स्निग्धमल्पशः~। }\pend
\pstart\versepar \skt{}\edtext{\skt{उप\-जि\-ह्वि\-क\-व\-च्चापि}}{\lemma{\skt{उप\appbrk{}जि\appbrk{}ह्वि\appbrk{}क\appbrk{}व\appbrk{}च्चापि}}\Afootnote{A K H\textsuperscript{pc}; \skt{उप\appbrk{}जि\appbrk{}ह्वि\appbrk{}व\appbrk{}च्चापि}\ H\textsuperscript{ac}}}\skt{ }\edtext{\skt{शोधयेदपि}}{\lemma{\skt{शोधयेदपि}}\Afootnote{K; \skt{साधयेदधि}\ A; \skt{साधयेदपि}\ H}}\skt{ जिह्विकां~।}\pend

\pstart\versepar \ledsidenote{\passref{4.22.66}}\skt{एकवृन्दन्तु }\edtext{\skt{विस्राव्य}}{\lemma{\skt{विस्राव्य}}\Afootnote{A K; \skt{विस्राव्यम्}\ H}}\skt{ विधिं }\edtext{\skt{शोध\-ण\-मा\-च\-रेत्}}{\lemma{\skt{शोधणम्}}\Afootnote{K; \skt{शोधन}\ A; \skt{शोधनम्}\ H}}\skt{~॥ }\pend
\pstart\versepar \skt{गिलायुकापि यो }\edtext{\skt{व्याधिस्तञ्च}}{\lemma{\skt{व्याधिस्तञ्च}}\Afootnote{A K; \skt{व्याधिस्तञ्च}\ H}}\skt{ शस्त्रेण साधयेत्~।}\edtext{}{\lemma{\skt{आचरेत्~॥ गिलायुकापि यो}}\Afootnote{K H; \skt{माचरेत्~। गिलायुश्चापि यो}\ A}}\pend

\pstart\versepar \ledsidenote{\passref{4.22.67}}\skt{अमर्मस्थं }\edtext{\skt{सुसंपक्वं}}{\lemma{\skt{सुसंपक्वं}}\Afootnote{K H; \skt{सुपक्वं च}\ A}}\skt{ भेद\-ये\-द्ग\-ल\-वि\-द्रधिं~। }\pend
\pstart\versepar \skt{}\edtext{\edtext{\skt{वातां\-त्स\-र्व\-स\-र\-ञ्चू\-र्णै\-र्ला\-वणैः}}{\lemma{\skt{वातां\appbrk{}त्स\appbrk{}र्व\appbrk{}स\appbrk{}रञ्}}\Afootnote{K\textsuperscript{pc}; \skt{वातात्सर्वसरं}\ A; \skt{वातां\appbrk{}त्स\appbrk{}र्व\appbrk{}र\appbrk{}स\appbrk{}रञ्}\ K\textsuperscript{ac}; \skt{वातात्सर्व्वसरां}\ H}}}{\lemma{\skt{लावणैः}}\Afootnote{K; \skt{लवणैः}\ A H}}\skt{ प्रतिसारयेत्~।}\pend

\pstart\versepar \ledsidenote{\passref{4.22.68}}\skt{तैलं वातहरैः सिद्धं हितं }\edtext{\skt{कवडनस्तयोः}}{\lemma{\skt{कवडनस्तयोः}}\Afootnote{K H; \skt{कवलनस्ययोः~।}\ A}}\skt{~। }\pend
\pstart\versepar \skt{}\edtext{\skt{ततोस्मै}}{\lemma{\skt{ततोस्मै}}\Afootnote{K; \skt{ततोऽस्मै}\ A H}}\skt{ }\edtext{\edtext{\skt{स्नैहि\-क\-न्धू\-म\-मिदं}}{\lemma{\skt{धूमम्}}\Afootnote{K H; \skt{धूम}\ A}}}{\lemma{\skt{इदं}}\Afootnote{K H; \skt{मिमं}\ A}}\skt{ }\edtext{\skt{दद्याद्विचक्षणः}}{\lemma{\skt{दद्याद्विचक्षणः}}\Afootnote{K H; \skt{दद्याद्विचक्षणः~॥}\ A}}\skt{~॥}\pend

\pstart\versepar \ledsidenote{\passref{4.22.69}}\edtext{\skt{साल\-रा\-जा\-द\-नै\-र\-ण्ड\-चा\-रे\-ङ्गु\-दि\-म\-धू\-कजाः}}{\lemma{\skt{साल\appbrk{}रा\appbrk{}जा\appbrk{}द\appbrk{}नै\appbrk{}र\appbrk{}ण्ड\appbrk{}चा\appbrk{}रे\appbrk{}ङ्गु\appbrk{}दि\appbrk{}म\appbrk{}धू\appbrk{}कजाः}}\Afootnote{K; \skt{शाल\appbrk{}रा\appbrk{}जा\appbrk{}द\appbrk{}नै\appbrk{}र\appbrk{}ण्ड\appbrk{}सा\appbrk{}रै\appbrk{}ङ्गु\appbrk{}द\appbrk{}म\appbrk{}धू\appbrk{}कजाः~।}\ A; \skt{शाल\appbrk{}रा\appbrk{}जा\appbrk{}द\appbrk{}नै\appbrk{}र\appbrk{}ण्ड\appbrk{}चा\appbrk{}रे\appbrk{}ङ्गु\appbrk{}दि\appbrk{}म\appbrk{}धू\appbrk{}कजाः~॥\textbackslash\{\}}\ H}}\skt{~। }\pend
\pstart\versepar \skt{मज्जानो }\edtext{\skt{गुग्गु\-लु\-ध्या\-म\-मां\-सी\-का\-ला\-नु\-सा\-रिवा}}{\lemma{\skt{गुग्गु\appbrk{}लु\appbrk{}ध्या\appbrk{}म\appbrk{}मां\appbrk{}सी\appbrk{}का\appbrk{}ला\appbrk{}नु\appbrk{}सा\appbrk{}रिवा}}\Afootnote{K; \skt{गुग्गु\appbrk{}लु\appbrk{}ध्या\appbrk{}म\appbrk{}मां\appbrk{}सी\appbrk{}का\appbrk{}ला\appbrk{}नु\appbrk{}सा\appbrk{}रिवाः~।}\ A; \skt{गुग्गु\appbrk{}लु\appbrk{}ध्या\appbrk{}म\appbrk{}मा\appbrk{}न्सी\appbrk{}का\appbrk{}ला\appbrk{}नु\appbrk{}शा\appbrk{}रिवा~।}\ H}}\skt{~। }\pend
\pstart\versepar \skt{श्रीस\-र्ज\-र\-स\-शै\-ले\-य\-म\-धू\-च्छि\-ष्टानि चाहरेत्~।}\pend

\pstart\versepar \ledsidenote{\passref{4.22.70}}\skt{तत्सर्वं सुकृतं चूर्ण्णं स्नेहेनालोड्य युक्तितः~। }\pend
\pstart\versepar \skt{}\edtext{\skt{टुण्टू\-क\-वृ\-त्ता\-त्स\-क्षौ\-द्रा\-त्म\-तिमां}}{\lemma{\skt{टुण्टू\appbrk{}क\appbrk{}वृ\appbrk{}त्ता\appbrk{}त्स\appbrk{}क्षौ\appbrk{}द्रा\appbrk{}त्म\appbrk{}तिमां}}\Afootnote{K H; \skt{टिण्डूकवृन्तं सक्षौद्रं मतिमां}\ A}}\skt{ }\edtext{\skt{तेन}}{\lemma{\skt{तेन}}\Afootnote{K H; \skt{स्तेन}\ A}}\skt{ लेपयेत्~।}\pend

\pstart\versepar \ledsidenote{\passref{4.22.71}}\skt{एष सर्वसरे धूमः }\edtext{\skt{प्रश\-स्त\-स्नै\-हि\-को\-त्तमः}}{\lemma{\skt{प्रश\appbrk{}स्त\appbrk{}स्नै\appbrk{}हि\appbrk{}को\appbrk{}त्तमः}}\Afootnote{K H; \skt{प्रशस्तः स्नैहिको मतः~।}\ A}}\skt{~। }\pend
\pstart\versepar \skt{कफघ्नो }\edtext{\skt{मारुघ्नश्च}}{\lemma{\skt{मारुघ्नश्च}}\Afootnote{K; \skt{मारुतघ्नश्च}\ A H}}\skt{ मुख\-रो\-ग\-वि\-ना\-शनः~।}\pend

\pstart\versepar \ledsidenote{\passref{4.22.72}}\skt{पित्तात्मके सर्वसरे शुद्धकायस्य देहिनः~। }\pend
\pstart\versepar \skt{सर्वः पित्तहरः कार्यो विधि\-र्म\-धु\-र\-शी\-तलः~।}\pend

\pstart\versepar \ledsidenote{\passref{4.22.73}}\edtext{\skt{प्रति\-सा\-र\-ण\-ग\-ण्डूषा}}{\lemma{\skt{प्रति\appbrk{}सा\appbrk{}र\appbrk{}ण\appbrk{}ग\appbrk{}ण्डूषा}}\Afootnote{K H\textsuperscript{pc}; \skt{प्रति\appbrk{}सा\appbrk{}र\appbrk{}ण\appbrk{}ग\appbrk{}ण्डूषौ}\ A; \skt{प्रति\appbrk{}सा\appbrk{}र\appbrk{}ग\appbrk{}ण्डूषा}\ H\textsuperscript{ac}}}\skt{ }\edtext{\skt{धूमं}}{\lemma{\skt{धूमं}}\Afootnote{K H; \skt{धूमः}\ A}}\skt{ संशोधनानि च~॥ }\pend
\pstart\versepar \skt{कफात्मके सर्वसरे }\edtext{\skt{क्रमां}}{\lemma{\skt{क्रमां}}\Afootnote{K H; \skt{विधिं}\ A}}\skt{ कुर्यात्कफापहं~।}\pend

\pstart\versepar \ledsidenote{\passref{4.22.74}}\edtext{\skt{पिवेदतिविषां}}{\lemma{\skt{पिवेदतिविषां}}\Afootnote{K; \skt{पिबेदतिविषां}\ A; \skt{पिबे\appbrk{}द\appbrk{}ति\appbrk{}वि\appbrk{}षाम्}\ H}}\skt{ }\edtext{\skt{पाठं}}{\lemma{\skt{पाठं}}\Afootnote{K H; \skt{पाठां}\ A}}\skt{ }\edtext{\skt{मुस्तञ्च}}{\lemma{\skt{मुस्तञ्च}}\Afootnote{A K; \skt{मुस्ताञ्च}\ H}}\skt{ }\edtext{\skt{च}}{\lemma{\skt{च}}\Afootnote{\skt{सुरदारु च}\ A H K}}\skt{ }\pend
\pstart\versepar \skt{}\edtext{\skt{रोहिणी}}{\lemma{\skt{रोहिणी}}\Afootnote{K H; \skt{रोहिणीं}\ A}}\skt{ }\edtext{\skt{कटुकाख्यान्तु}}{\lemma{\skt{कटुकाख्यान्तु}}\Afootnote{K H; \skt{कटुकाख्यां च}\ A}}\skt{ }\edtext{\skt{कटजस्य}}{\lemma{\skt{कटजस्य}}\Afootnote{K; \skt{कुटजस्य}\ A H}}\skt{ फलानि च~।}\pend

\pstart\versepar \ledsidenote{\passref{4.22.75}}\skt{गवां मूत्रेण मनुजो }\edtext{\skt{भागां}}{\lemma{\skt{भागां}}\Afootnote{K H; \skt{भागैर्}\ A}}\skt{ }\edtext{\skt{धरणसंमितां}}{\lemma{\skt{धरणसंमितां}}\Afootnote{K H; \skt{धरणसंमितैः~।}\ A}}\skt{~। }\pend
\pstart\versepar \skt{एष }\edtext{\skt{सर्वां}}{\lemma{\skt{सर्वां}}\Afootnote{K H; \skt{सर्वा}\ A}}\skt{ }\edtext{\skt{कफकृतां}}{\lemma{\skt{कफकृतां}}\Afootnote{K H; \skt{न्कफकृता}\ A}}\skt{ }\edtext{\skt{रोगा\-न्यो\-गो\-प\-क\-र्षति}}{\lemma{\skt{रोगा\appbrk{}न्यो\appbrk{}गो\appbrk{}प\appbrk{}क\appbrk{}र्षति}}\Afootnote{K; \skt{न्रोगा\appbrk{}न्यो\appbrk{}गोऽ\appbrk{}प\appbrk{}क\appbrk{}र्षति~॥}\ A; \skt{रोगा\appbrk{}न्यो\appbrk{}गोऽ\appbrk{}प\appbrk{}क\appbrk{}र्षति~॥}\ H}}\skt{~।}\pend

\pstart\versepar \ledsidenote{\passref{4.22.76}}\edtext{\skt{क्षीरे\-क्षु\-र\-स\-गो\-मू\-त्र\-द\-धि\-स्त्व\-म्ल\-का\-ञ्जिकैः}}{\lemma{\skt{क्षीरे\appbrk{}क्षु\appbrk{}र\appbrk{}स\appbrk{}गो\appbrk{}मू\appbrk{}त्र\appbrk{}द\appbrk{}धि\appbrk{}स्त्व\appbrk{}म्ल\appbrk{}का\appbrk{}ञ्जिकैः}}\Afootnote{K; \skt{क्षीरे\appbrk{}क्षु\appbrk{}र\appbrk{}स\appbrk{}गो\appbrk{}मू\appbrk{}त्र\appbrk{}द\appbrk{}धि\appbrk{}म\appbrk{}स्त्व\appbrk{}म्ल\appbrk{}का\appbrk{}ञ्जिकैः~।}\ A H}}\skt{~। }\pend
\pstart\versepar \skt{}\edtext{\skt{विद\-ध्या\-त्क\-व\-लां\-वीक्ष्य}}{\lemma{\skt{विद\appbrk{}ध्या\appbrk{}त्क\appbrk{}व\appbrk{}लां\appbrk{}वीक्ष्य}}\Afootnote{K H; \skt{विद\appbrk{}ध्या\appbrk{}त्क\appbrk{}व\appbrk{}ला\appbrk{}न्वीक्ष्य}\ A}}\skt{ दोषं तैलघृतैरपि~।}\pend

\pstart\versepar \ledsidenote{\passref{4.22.77}}\edtext{\skt{रोगानां}}{\lemma{\skt{रोगानां}}\Afootnote{K; \skt{रोगाणां}\ A H}}\skt{ मुखजातानां साध्यानाङ्कर्म कीर्त्तितं~॥ }\pend
\pstart\versepar \skt{}\edtext{\skt{मपि}}{\lemma{\skt{अपि}}\Afootnote{\skt{असाध्या अपि}\ A; \skt{असाध्यामपि}\ K H}}\skt{ }\edtext{\skt{वक्ष्यामि}}{\lemma{\skt{वक्ष्यामि}}\Afootnote{K H; \skt{वक्ष्यन्ते}\ A}}\skt{ रोगा ये }\edtext{\skt{यत्र}}{\lemma{\skt{यत्र}}\Afootnote{K H; \skt{तत्र}\ A}}\skt{ कीर्त्तिताः~।}\edtext{}{\lemma{\skt{कीर्त्तितं~॥ म्}}\Afootnote{\skt{कीर्तितम्}\ A K H}}\pend

\pstart\versepar \ledsidenote{\passref{4.22.78}}\edtext{\skt{औष्ठप्रकोपो}}{\lemma{\skt{औष्ठप्रकोपो}}\Afootnote{K H; \skt{औष्ठप्रकोपे}\ A}}\skt{ }\edtext{\skt{वर्ज्या स्युर्मां\-स\-र\-क्त\-त्रि\-दो\-षजाः}}{\lemma{\skt{वर्ज्या स्युर्}}\Afootnote{K; \skt{वर्ज्याः स्युर्}\ A H}}\skt{~। }\pend
\pstart\versepar \skt{}\edtext{\skt{दन्तमूलेषु}}{\lemma{\skt{दन्तमूलेषु}}\Afootnote{A K; \skt{दन्तमूले तु}\ H\textsuperscript{ac}; \skt{दन्तमूलेषु तु}\ H\textsuperscript{pc}}}\skt{ }\edtext{\skt{वर्ज्यौ तु}}{\lemma{\skt{वर्ज्यौ तु}}\Afootnote{A K; \skt{वर्ज्य तु}\ H}}\skt{ त्रिलिं\-ग\-ग\-ति\-सौ\-षिरौ~।}\pend

\pstart\versepar \ledsidenote{\passref{4.22.79}}\skt{दन्तेषु च न सिध्यन्ति श्याव\-दा\-ल\-न\-भ\-ञ्जनाः~। }\pend
\pstart\versepar \skt{}\edtext{\skt{जिह्वा\-ग\-ते\-ष्व\-लां\-सन्तु}}{\lemma{\skt{अलांसन्तु}}\Afootnote{K; \skt{अलासस्तु}\ A; \skt{अला सन्तु}\ H}}\skt{ }\edtext{\skt{तालव्येष्वर्वुदं}}{\lemma{\skt{अर्वुदं}}\Afootnote{K H; \skt{अर्बुदं}\ A}}\skt{ तथा~।}\pend

\pstart\versepar \ledsidenote{\passref{4.22.80}}\skt{स्वरघ्नो }\edtext{\skt{वलयं}}{\lemma{\skt{वलयं}}\Afootnote{K H; \skt{वलयो}\ A}}\skt{ }\edtext{\skt{वृन्दं}}{\lemma{\skt{वृन्दं}}\Afootnote{K H; \skt{वृन्दो}\ A}}\skt{ }\edtext{\skt{वला\-स\-स\-ह\-चा\-रिणां}}{\lemma{\skt{वला\appbrk{}स\appbrk{}स\appbrk{}ह\appbrk{}चा\appbrk{}रिणां}}\Afootnote{K H; \skt{विदार्यलस एव च~।}\ A}}\skt{~। }\pend
\pstart\versepar \skt{}\edtext{\skt{गलौ\-घ\-मां\-स\-ता\-नश्च}}{\lemma{\skt{गलौ\appbrk{}घ\appbrk{}मां\appbrk{}स\appbrk{}ता\appbrk{}नश्च}}\Afootnote{K; \skt{गलौष्ठो मांसतानश्च}\ A; \skt{गलौ\appbrk{}घ\appbrk{}म्मा\appbrk{}न्स\appbrk{}ता\appbrk{}नञ्च}\ H}}\skt{ शतघ्नी }\edtext{\skt{रोहिणीतले}}{\lemma{\skt{रोहिणीतले}}\Afootnote{K; \skt{रोहिणी च या~॥}\ A; \skt{रोहिणीगले~॥}\ H}}\skt{~।}\pend

\pstart\versepar \ledsidenote{\passref{4.22.81}}\skt{असाध्याः कीर्त्तिता ह्येते रोगा नव दशैव च~। }\pend
\pstart\versepar \skt{}\edtext{\skt{तेषुचापि}}{\lemma{\skt{तेषुचापि}}\Afootnote{K H; \skt{तेषां चापि}\ A}}\skt{ }\edtext{\skt{क्रियां}}{\lemma{\skt{क्रियां}}\Afootnote{A K; \skt{क्रिया}\ H}}\skt{ वैद्यः प्रत्याख्याय }\edtext{\skt{समाचरेदिति}}{\lemma{\skt{समाचरेदिति}}\Afootnote{K; \skt{समाचरेत्~॥}\ A; \skt{समाचरेदिति~॥}\ H}}\skt{~॥}\pend

\pstart \ledsidenote{\passref{4.22.trailer}}\edtext{\skt{चि}}{\lemma{\skt{चि}}\Afootnote{K; \skt{इति सुश्रु\appbrk{}त\appbrk{}सं\appbrk{}हि\appbrk{}तायां चि}\ A; \emph{om.}\ H}}\skt{ }\edtext{\skt{थ२~॥ {\symfont ❈}}}{\lemma{\skt{थ२~॥ {\symfont ❈}}}\Afootnote{K; \skt{कित्सास्थाने मुख\appbrk{}रो\appbrk{}ग\appbrk{}चि\appbrk{}कि\appbrk{}त्सितं नाम द्वाविंशोऽध्यायः~॥२२~॥}\ A; \skt{चिकित्सा थं द्वि~॥~॥}\ H}}\skt{~॥}\pend

\endnumbering
\beginnumbering
\eledsection*{Adhyāya 23}
\pstart \ledsidenote{\passref{4.23.1}}\skt{अथातः }\edtext{\skt{शोफचिकित्सितं}}{\lemma{\skt{शोफचिकित्सितं}}\Afootnote{K H; \skt{शोफानां चिकित्सितं}\ A}}\skt{ }\edtext{\skt{व्याख्यास्यामः}}{\lemma{\skt{व्याख्यास्यामः}}\Afootnote{A H; \emph{om.}\ K}}\pend

\pstart \ledsidenote{\passref{4.23.3}}\edtext{\skt{तत्र}}{\lemma{\skt{तत्र}}\Afootnote{K H; \emph{om.}\ A}}\skt{ }\edtext{\skt{षड्वि\-धो\-व\-य\-व\-स\-मुत्थः}}{\lemma{\skt{षड्वि\appbrk{}धो\appbrk{}व\appbrk{}य\appbrk{}व\appbrk{}स\appbrk{}मुत्थः}}\Afootnote{K H; \skt{षड्वि\appbrk{}धोऽ\appbrk{}व\appbrk{}य\appbrk{}व\appbrk{}स\appbrk{}मुत्थः}\ A}}\skt{ }\edtext{\skt{शोफोभिहितो}}{\lemma{\skt{शोफोभिहितो}}\Afootnote{K; \skt{शोफोऽभिहितो}\ A H}}\skt{ लक्षणतः }\edtext{\skt{प्रतिकारतश्च}}{\lemma{\skt{प्रतिकारतश्च}}\Afootnote{K H; \skt{प्रतीकारतश्च;}\ A}}\skt{~। }\edtext{\skt{सर्वसरश्च}}{\lemma{\skt{सर्वसरश्च}}\Afootnote{K H; \skt{सर्वसरस्तु}\ A}}\skt{ पञ्चविधः~। }\edtext{\skt{तद्यथा}}{\lemma{\skt{तद्यथा}}\Afootnote{A K; \skt{तद्यथा}\ H}}\skt{~। }\edtext{\skt{वात\-पि\-त्त\-क\-फ\-सं\-नि\-पा\-त\-वि\-ष\-न्नि\-मित्तः}}{\lemma{\skt{वात\appbrk{}पि\appbrk{}त्त\appbrk{}क\appbrk{}फ\appbrk{}सं\appbrk{}नि\appbrk{}पा\appbrk{}त\appbrk{}वि\appbrk{}ष\appbrk{}न्नि\appbrk{}मित्तः}}\Afootnote{K; \skt{वात\appbrk{}पि\appbrk{}त्त\appbrk{}श्ले\appbrk{}ष्म\appbrk{}स\appbrk{}न्नि\appbrk{}पा\appbrk{}त\appbrk{}वि\appbrk{}ष\appbrk{}नि\appbrk{}मित्तः~॥}\ A; \skt{वात\appbrk{}पि\appbrk{}त्त\appbrk{}क\appbrk{}फ\appbrk{}स\appbrk{}न्नि\appbrk{}पा\appbrk{}त\appbrk{}वि\appbrk{}ष\appbrk{}नि\appbrk{}मित्तः~।}\ H}}\skt{~।}\pend

\pstart \ledsidenote{\passref{4.23.4}}\edtext{\skt{तत्रा\-प\-त\-र्पि\-त\-स्या\-भ्य\-र्थ\-म\-ध्व\-ग\-म\-ना\-द\-ति\-मा\-त्र\-म\-भ्य\-ह\-रतो वा}}{\lemma{\skt{तत्रा\appbrk{}प\appbrk{}त\appbrk{}र्पि\appbrk{}त\appbrk{}स्या\appbrk{}भ्य\appbrk{}र्थ\appbrk{}म\appbrk{}ध्व\appbrk{}ग\appbrk{}म\appbrk{}ना\appbrk{}द\appbrk{}ति\appbrk{}मा\appbrk{}त्र\appbrk{}म\appbrk{}भ्य\appbrk{}ह\appbrk{}रतो वा}}\Afootnote{K; \skt{तत्रा\appbrk{}प\appbrk{}त\appbrk{}र्पि\appbrk{}त\appbrk{}स्या\appbrk{}ध्व\appbrk{}ग\appbrk{}म\appbrk{}ना\appbrk{}द\appbrk{}ति\appbrk{}मा\appbrk{}त्र\appbrk{}म\appbrk{}भ्य\appbrk{}व\appbrk{}ह\appbrk{}रतो वा}\ A; \skt{तत्रापि तर्प्पि\appbrk{}त\appbrk{}स्या\appbrk{}र्थ\appbrk{}म\appbrk{}ध्व\appbrk{}ग\appbrk{}म\appbrk{}ना\appbrk{}द\appbrk{}ति\appbrk{}मा\appbrk{}त्र\appbrk{}म\appbrk{}भ्य\appbrk{}व\appbrk{}ह\appbrk{}रतो वा}\ H}}\skt{ }\edtext{\skt{पिष्टा\-न्न\-ह\-रि\-त\-शा\-क\-ल\-व\-णानि}}{\lemma{\skt{पिष्टा\appbrk{}न्न\appbrk{}ह\appbrk{}रि\appbrk{}त\appbrk{}शा\appbrk{}क\appbrk{}ल\appbrk{}व\appbrk{}णानि}}\Afootnote{K H; \skt{पिष्टा\appbrk{}न्न\appbrk{}ह\appbrk{}रि\appbrk{}त\appbrk{}क\appbrk{}शा\appbrk{}क\appbrk{}ल\appbrk{}व\appbrk{}णानि}\ A}}\skt{~। क्षीणस्य }\edtext{\edtext{\skt{वाति\-मा\-त्र\-म\-म्ल\-मु\-प\-से\-वतो}}{\lemma{\skt{वाति\appbrk{}मा\appbrk{}त्र\appbrk{}म\appbrk{}म्लम्}}\Afootnote{K; \skt{वाऽति\appbrk{}मा\appbrk{}त्र\appbrk{}म\appbrk{}म्लम्}\ A; \skt{चाति\appbrk{}मा\appbrk{}त्त्र\appbrk{}म\appbrk{}म्लम्}\ H}}}{\lemma{\skt{उपसेवतो}}\Afootnote{K H; \skt{उप\appbrk{}से\appbrk{}व\appbrk{}मा\appbrk{}नस्य}\ A}}\skt{ }\edtext{\edtext{\skt{मृत्प\-क्त\-लो\-ष्ट\-क\-ट\-श\-र्क\-रा\-मु\-पौ\-द\-क\-मां\-स\-से\-वितो}}{\lemma{\skt{मृत्प\appbrk{}क्त\appbrk{}लो\appbrk{}ष्ट\appbrk{}क\appbrk{}ट\appbrk{}श\appbrk{}र्क\appbrk{}राम्}}\Afootnote{K; \skt{मृत्प\appbrk{}क्व\appbrk{}लो\appbrk{}ष्ट\appbrk{}क\appbrk{}ट\appbrk{}श\appbrk{}र्करा}\ A; \skt{मृत्प\appbrk{}ङ्क\appbrk{}लो\appbrk{}ष्ट\appbrk{}क\appbrk{}ट\appbrk{}श\appbrk{}र्क्करा}\ H}}\skt{ वा}}{\lemma{\skt{उपौ\appbrk{}द\appbrk{}क\appbrk{}मां\appbrk{}स\appbrk{}से\appbrk{}वितो वा}}\Afootnote{K; \skt{नूपौ\appbrk{}द\appbrk{}क\appbrk{}मां\appbrk{}स\appbrk{}से\appbrk{}व\appbrk{}ना\appbrk{}द\appbrk{}जी\appbrk{}र्णिनो वा ग्राम्य\appbrk{}ध\appbrk{}र्म\appbrk{}से\appbrk{}वना}\ A; \skt{नूपो\appbrk{}द\appbrk{}क\appbrk{}मां\appbrk{}स\appbrk{}से\appbrk{}विनो वा}\ H}}\skt{ }\edtext{\skt{विरुद्धाहारस्य}}{\lemma{\skt{विरुद्धाहारस्य}}\Afootnote{K; \skt{द्विरु\appbrk{}द्धा\appbrk{}हा\appbrk{}र\appbrk{}सेव}\ A; \skt{विरुद्धाहार}\ H}}\skt{ }\edtext{\skt{हस्त्य\-श्व\-क्षो\-भ\-ना\-दा\-या\-सिता}}{\lemma{\skt{हस्त्य\appbrk{}श्व\appbrk{}क्षो\appbrk{}भ\appbrk{}ना\appbrk{}दा\appbrk{}या\appbrk{}सिता}}\Afootnote{K; \skt{नात्वा हस्त्य\appbrk{}श्वो\appbrk{}ष्ट्र\appbrk{}र\appbrk{}थ\appbrk{}प\appbrk{}दा\appbrk{}ति\appbrk{}स\appbrk{}ङ्क्षो\appbrk{}भ\appbrk{}णा\appbrk{}दयो सितस्य}\ A; \skt{स्य हस्त्य\appbrk{}श्व\appbrk{}क्षो\appbrk{}भ\appbrk{}ना\appbrk{}दा\appbrk{}या\appbrk{}सिता}\ H}}\skt{ दोषा }\edtext{\skt{धातूं}}{\lemma{\skt{धातूं}}\Afootnote{K; \skt{धातू}\ A H}}\skt{ }\edtext{\skt{प्रदूषयित्वा}}{\lemma{\skt{प्रदूषयित्वा}}\Afootnote{K; \skt{न्प्रदूष्य श्व}\ A; \skt{न्प्रदूषयित्वा}\ H}}\skt{ }\edtext{\skt{शोफ\-मा\-पा\-द\-य\-न्त्य\-खिले}}{\lemma{\skt{शोफ\appbrk{}मा\appbrk{}पा\appbrk{}द\appbrk{}य\appbrk{}न्त्य\appbrk{}खिले}}\Afootnote{K; \skt{यथु\appbrk{}मा\appbrk{}पा\appbrk{}द\appbrk{}य\appbrk{}न्त्य\appbrk{}खिले}\ A; \skt{शोफ\appbrk{}मा\appbrk{}पा\appbrk{}द\appbrk{}यन्ति अखिल}\ H}}\skt{ शरीरे~।}\pend

\pstart \ledsidenote{\passref{4.23.5}}\edtext{\skt{तत्र}}{\lemma{\skt{तत्र}}\Afootnote{K; \skt{तत्र वात\appbrk{}श्व\appbrk{}य\appbrk{}थु\appbrk{}र\appbrk{}रुणः कृष्णो वा मृदु\appbrk{}र\appbrk{}न\appbrk{}व\appbrk{}स्थि\appbrk{}ता\appbrk{}स्तो\appbrk{}दा\appbrk{}द\appbrk{}य\appbrk{}श्चात्र वेदनाविशेषाः; पित्तश्वयथुः पीतः सरक्तो वा मृदुः शीघ्रा\appbrk{}नु\appbrk{}सा\appbrk{}र्यू\appbrk{}षा\appbrk{}द\appbrk{}य\appbrk{}श्चात्र}\ A; \skt{तत्र वातपित्त}\ H}}\skt{ }\edtext{\skt{वात\-पि\-त्त\-क\-फ\-स\-न्नि\-पा\-त\-श्व\-य\-थू\-ना\-म\-व\-य\-व\-स\-मुत्थे}}{\lemma{\skt{वात\appbrk{}पि\appbrk{}त्त\appbrk{}क\appbrk{}फ\appbrk{}स\appbrk{}न्नि\appbrk{}पा\appbrk{}त\appbrk{}श्व\appbrk{}य\appbrk{}थू\appbrk{}ना\appbrk{}म\appbrk{}व\appbrk{}य\appbrk{}व\appbrk{}स\appbrk{}मुत्थे}}\Afootnote{K; \skt{वेदनाविशेषाः; श्लेष्मश्वयथुः पाण्डुः शुक्लो वा स्निग्धः कठिनः शीतो मन्दानुसारी, कण्ड्वादयश्चात्र वेदनाविशेषाः; सन्नि\appbrk{}पा\appbrk{}त\appbrk{}श्व\appbrk{}यथुः सर्ववर्णवे}\ A; \skt{कफ\appbrk{}स\appbrk{}न्नि\appbrk{}पा\appbrk{}त\appbrk{}श्व\appbrk{}य\appbrk{}थू\appbrk{}ना\appbrk{}म\appbrk{}व\appbrk{}य\appbrk{}व\appbrk{}स\appbrk{}मुत्थे}\ H}}\skt{ }\edtext{\skt{शोफलक्षणानि}}{\lemma{\skt{शोफलक्षणानि}}\Afootnote{K H; \skt{द}\ A}}\skt{ }\edtext{\skt{व्याख्यातानि}}{\lemma{\skt{व्याख्यातानि}}\Afootnote{K H; \skt{नः;}\ A}}\skt{~। }\edtext{\skt{विषनिमित्तस्तु}}{\lemma{\skt{विषनिमित्तस्तु}}\Afootnote{K H; \skt{विषनिमित्तस्तु}\ A}}\skt{ }\edtext{\skt{गरोपयोगाद्दु ष्टतो\-य\-से\-व\-ना\-त्प्र\-कु\-थि\-तो\-द\-का\-व\-से\-का\-त्स\-वि\-ष\-स\-त्व\-स\-ञ्चू\-र्ण्ण\-नाद्वा}}{\lemma{\skt{गरोपयोगाद्दु ष्टतो\appbrk{}य\appbrk{}से\appbrk{}व\appbrk{}ना\appbrk{}त्प्र\appbrk{}कु\appbrk{}थि\appbrk{}तो\appbrk{}द\appbrk{}का\appbrk{}व\appbrk{}से\appbrk{}का\appbrk{}त्स\appbrk{}वि\appbrk{}ष\appbrk{}स\appbrk{}त्व\appbrk{}स\appbrk{}ञ्चू\appbrk{}र्ण्ण\appbrk{}नाद्वा}}\Afootnote{K; \skt{गरो\appbrk{}प\appbrk{}यो\appbrk{}गा\appbrk{}द्दु\appbrk{}ष्ट\appbrk{}तो\appbrk{}य\appbrk{}से\appbrk{}व\appbrk{}ना\appbrk{}त्प्र\appbrk{}कु\appbrk{}थि\appbrk{}तो\appbrk{}द\appbrk{}का\appbrk{}व\appbrk{}गा\appbrk{}ह\appbrk{}ना\appbrk{}त्स\appbrk{}वि\appbrk{}ष\appbrk{}स\appbrk{}त्त्व\appbrk{}दि\appbrk{}ग्ध\appbrk{}चू\appbrk{}र्णा\appbrk{}व\appbrk{}चू\appbrk{}र्ण\appbrk{}नाद्वा}\ A; \skt{गरो\appbrk{}प\appbrk{}यो\appbrk{}गा\appbrk{}द्दु\appbrk{}ष्ट\appbrk{}तो\appbrk{}य\appbrk{}से\appbrk{}व\appbrk{}ना\appbrk{}त्प्र\appbrk{}कु\appbrk{}थि\appbrk{}तो\appbrk{}द\appbrk{}का\appbrk{}व\appbrk{}से\appbrk{}का\appbrk{}त्स\appbrk{}वि\appbrk{}ष\appbrk{}स\appbrk{}त्व\appbrk{}स\appbrk{}ञ्चू\appbrk{}र्ण्ण\appbrk{}नाद्वा~।}\ H\textsuperscript{ac}; \skt{गरो\appbrk{}प\appbrk{}यो\appbrk{}गा\appbrk{}द्दु\appbrk{}ष्ट\appbrk{}तो\appbrk{}य\appbrk{}से\appbrk{}व\appbrk{}ना\appbrk{}त्प्र\appbrk{}कु\appbrk{}पि\appbrk{}थि\appbrk{}तो\appbrk{}द\appbrk{}का\appbrk{}व\appbrk{}से\appbrk{}का\appbrk{}त्स\appbrk{}वि\appbrk{}ष\appbrk{}स\appbrk{}त्व\appbrk{}स\appbrk{}ञ्चू\appbrk{}र्ण्ण\appbrk{}नाद्वा~।}\ H\textsuperscript{pc}}}\skt{~। }\edtext{\skt{सवि\-ष\-मू\-त्र\-पु\-री\-श\-शु\-क्र\-स्पृ\-ष्टानां वा}}{\lemma{\skt{सवि\appbrk{}ष\appbrk{}मू\appbrk{}त्र\appbrk{}पु\appbrk{}री\appbrk{}श\appbrk{}शु\appbrk{}क्र\appbrk{}स्पृ\appbrk{}ष्टानां वा}}\Afootnote{K; \skt{सवि\appbrk{}ष\appbrk{}मू\appbrk{}त्र\appbrk{}पु\appbrk{}री\appbrk{}ष\appbrk{}शु\appbrk{}क्र\appbrk{}स्पृ\appbrk{}ष्टानां वा}\ A H}}\skt{ तृणकाष्ठादीनां }\edtext{\skt{संस्पर्शात्~॥ स तु}}{\lemma{\skt{संस्पर्शात्~॥ स तु}}\Afootnote{K H; \skt{संस्पर्शनात्, स तु}\ A}}\skt{ मृदुः }\edtext{\skt{क्षिप्रोत्थाम्च}}{\lemma{\skt{क्षिप्रोत्थाम्च}}\Afootnote{K; \skt{क्सिप्रो\appbrk{}त्था\appbrk{}नोऽ\appbrk{}व\appbrk{}लम्बी चलो}\ A; \skt{क्षिप्रो\appbrk{}त्था\appbrk{}न\appbrk{}श्च\appbrk{}ला\appbrk{}व\appbrk{}लम्बी}\ H}}\skt{ }\edtext{\skt{लावलम्वी}}{\lemma{\skt{लावलम्वी}}\Afootnote{K; \skt{ऽचलो वा}\ A; \emph{om.}\ H}}\skt{ }\edtext{\skt{दाह\-पा\-क\-प्रा\-यश्च}}{\lemma{\skt{दाह\appbrk{}पा\appbrk{}क\appbrk{}प्रा\appbrk{}यश्च}}\Afootnote{K; \skt{दाह\appbrk{}पा\appbrk{}क\appbrk{}रा\appbrk{}ग\appbrk{}प्रा\appbrk{}यश्च}\ A; \skt{दाह\appbrk{}पा\appbrk{}क\appbrk{}प्र\appbrk{}यश्च}\ H}}\skt{ भवति~॥}\pend

\pstart\versepar \ledsidenote{\passref{4.23.6}}\edtext{\skt{दोषाः}}{\lemma{\skt{दोषाः}}\Afootnote{K H; \skt{भवन्ति चात्र दोषाः}\ A}}\skt{ }\edtext{\skt{श्वयथुमूर्ध्वं हि}}{\lemma{\skt{श्वयथुमूर्ध्वं हि}}\Afootnote{A K; \skt{श्वयथुमूर्द्धं हि}\ H}}\skt{ }\edtext{\skt{कुर्व\-न्त्या\-मा\-श\-या\-श्रिताः}}{\lemma{\skt{कुर्व\appbrk{}न्त्या\appbrk{}मा\appbrk{}श\appbrk{}या\appbrk{}श्रिताः}}\Afootnote{K; \skt{कुर्व\appbrk{}न्त्या\appbrk{}मा\appbrk{}श\appbrk{}य\appbrk{}स्थिताः~।}\ A; \skt{कुर्व्व\appbrk{}न्त\appbrk{}या\appbrk{}मा\appbrk{}या\appbrk{}स्रिताः~।}\ H\textsuperscript{ac}; \skt{कुर्व्व\appbrk{}न्त\appbrk{}या\appbrk{}मा\appbrk{}स\appbrk{}या\appbrk{}स्रिताः~।}\ H\textsuperscript{pc}}}\skt{~। }\pend
\pstart\versepar \skt{पक्वाशयस्था }\edtext{\skt{मध्ये तु}}{\lemma{\skt{मध्ये तु}}\Afootnote{K H; \skt{मध्ये च}\ A}}\skt{ }\edtext{\skt{वर्च\-स्था\-न\-ग\-ता\-स्त्वधः}}{\lemma{\skt{वर्च\appbrk{}स्था\appbrk{}न\appbrk{}ग\appbrk{}ता\appbrk{}स्त्वधः}}\Afootnote{K H; \skt{वर्चः\appbrk{}स्था\appbrk{}न\appbrk{}ग\appbrk{}ता\appbrk{}स्त्वधः~॥}\ A}}\skt{~।}\pend

\pstart\versepar \ledsidenote{\passref{4.23.7}}\skt{कृत्स्नं देहमनुप्राप्ताः~। कुर्युः सर्वसरन्तथा~॥ }\pend
\pstart\versepar \skt{}\edtext{\skt{यो}}{\lemma{\skt{यो}}\Afootnote{K H; \skt{श्वयथुर्}\ A}}\skt{ मध्यदेशे }\edtext{\skt{श्वयथुः स}}{\lemma{\skt{श्वयथुः स}}\Afootnote{K; \skt{यः स}\ A; \skt{स्वयथुः स}\ H}}\skt{ कष्टः सर्वगश्च यः~।}\pend

\pstart\versepar \ledsidenote{\passref{4.23.8}}\edtext{\skt{अर्धा\-ङ्गे\-रि\-ष्ट\-भूतं}}{\lemma{\skt{अर्धा\appbrk{}ङ्गे\appbrk{}रि\appbrk{}ष्ट\appbrk{}भूतं}}\Afootnote{K H; \skt{अर्धा\appbrk{}ङ्गेऽ\appbrk{}रि\appbrk{}ष्ट\appbrk{}भू\appbrk{}तश्}\ A}}\skt{ }\edtext{\skt{स्याद्यश्चोर्ध्वं}}{\lemma{\skt{स्याद्यश्चोर्ध्वं}}\Afootnote{K; \skt{च यश्चोर्ध्वं}\ A; \skt{स्याद्यश्चोर्द्धम्}\ H}}\skt{ परिसर्पति~। }\pend
\pstart\versepar \skt{श्वासः }\edtext{\skt{पिपासाच्छर्दिश्च}}{\lemma{\skt{पिपासाच्छर्दिश्च}}\Afootnote{K H; \skt{पिपासा}\ A}}\skt{ }\edtext{\skt{दौर्वल्यं}}{\lemma{\skt{दौर्वल्यं}}\Afootnote{K; \skt{दौर्बल्यं}\ A H}}\skt{ }\edtext{\skt{ज्वर}}{\lemma{\skt{ज्वर}}\Afootnote{K H; \skt{ज्वरश्}\ A}}\skt{ }\edtext{\skt{एव}}{\lemma{\skt{एव}}\Afootnote{K H; \emph{om.}\ A}}\skt{ }\edtext{\skt{च}}{\lemma{\skt{च}}\Afootnote{K H; \skt{च्छर्दिररोचकः~॥}\ A}}\skt{~।}\pend

\pstart\versepar \ledsidenote{\passref{4.23.9}}\edtext{\skt{यस्य }\edtext{\skt{चान्नेरुचिर्नास्ति}}{\lemma{\skt{चान्नेरुचिर्नास्ति}}\Afootnote{K\textsuperscript{pc}; \skt{चान्नेरुर्नास्ति}\ K\textsuperscript{ac}; \skt{चान्ने रुचिर्न्नास्ति}\ H}}\skt{ श्वयथुं तं विवर्जयेत्}}{\lemma{\skt{… वि}}\Afootnote{\emph{om.}\ A}}\skt{~। }\pend
\pstart\versepar \skt{सामान्यतो विशेषाच्च तेषां वक्ष्यामि भेषजं~।}\edtext{}{\lemma{\skt{विवर्जयेत्~। सामान्यतो}}\Afootnote{K H; \skt{हिक्का\appbrk{}ती\appbrk{}सा\appbrk{}र\appbrk{}का\appbrk{}साश्च शूनं सङ्क्षपयन्ति हि~। सामान्यतो}\ A}}\pend

\pstart \ledsidenote{\passref{4.23.10}}\edtext{\skt{तत्र}}{\lemma{\skt{तत्र}}\Afootnote{K H; \emph{om.}\ A}}\skt{ शोफिनः सर्व }\edtext{\skt{एव}}{\lemma{\skt{एव}}\Afootnote{A K; \skt{ए}\ H}}\skt{ }\edtext{\skt{परिहरेयुः~।अ\-म्ल\-ल\-व\-ण\-म\-द्य\-गु\-ड\-द\-धि\-व\-सा\-तै\-ल\-पि\-ष्ट\-गु\-रूणि}}{\lemma{\skt{परिहरेयुः~।अ\appbrk{}म्ल\appbrk{}ल\appbrk{}व\appbrk{}ण\appbrk{}म\appbrk{}द्य\appbrk{}गु\appbrk{}ड\appbrk{}द\appbrk{}धि\appbrk{}व\appbrk{}सा\appbrk{}तै\appbrk{}ल\appbrk{}पि\appbrk{}ष्ट\appbrk{}गु\appbrk{}रूणि}}\Afootnote{K; \skt{परि\appbrk{}ह\appbrk{}रे\appbrk{}यु\appbrk{}र\appbrk{}म्ल\appbrk{}ल\appbrk{}व\appbrk{}ण\appbrk{}द\appbrk{}धि\appbrk{}गु\appbrk{}ड\appbrk{}व\appbrk{}सा\appbrk{}प\appbrk{}य\appbrk{}स्तै\appbrk{}ल\appbrk{}घृ\appbrk{}त\appbrk{}पि\appbrk{}ष्ट\appbrk{}म\appbrk{}य\appbrk{}गु\appbrk{}रूणि~॥}\ A; \skt{परिहरेयुः~। अम्ल\appbrk{}ल\appbrk{}व\appbrk{}ण\appbrk{}म\appbrk{}धु\appbrk{}गु\appbrk{}ड\appbrk{}द\appbrk{}धि\appbrk{}व\appbrk{}सा\appbrk{}तै\appbrk{}ल\appbrk{}पि\appbrk{}ष्ट\appbrk{}गु\appbrk{}रूणि}\ H}}\skt{ }\edtext{\skt{विशेषेण}}{\lemma{\skt{विशेषेण}}\Afootnote{K H; \emph{om.}\ A}}\skt{~॥}\pend

\pstart \ledsidenote{\passref{4.23.11}}\skt{तत्र वातश्वयथौ }\edtext{\skt{त्रिवृ\-त्तै\-र\-ण्ड\-तैलं वा}}{\lemma{\skt{त्रिवृ\appbrk{}त्तै\appbrk{}र\appbrk{}ण्ड\appbrk{}तैलं वा}}\Afootnote{K; \skt{त्रैवृ\appbrk{}त\appbrk{}मे\appbrk{}र\appbrk{}ण्ड\appbrk{}तैलं वा}\ A; \skt{तृवृ\appbrk{}त\appbrk{}मे\appbrk{}र\appbrk{}ण्ड\appbrk{}तै\appbrk{}लम्वा}\ H}}\skt{ }\edtext{\skt{मास\-म\-र्द्ध\-मा\-सम्वा}}{\lemma{\skt{मास\appbrk{}म\appbrk{}र्द्ध\appbrk{}मा\appbrk{}सम्वा}}\Afootnote{\skt{मासमर्धमासं वा}\ A K\textsuperscript{pc} H; \skt{मासम}\ K\textsuperscript{ac}}}\skt{ पाययेत्~। न्यग्रोधादि }\edtext{\skt{सिद्धं}}{\lemma{\skt{सिद्धं}}\Afootnote{K H; \skt{ककषायसिद्धं}\ A}}\skt{ सर्पिः~। पित्तश्वयथौ~। आर\-ग्व\-धा\-दि\-सिद्धं }\edtext{\skt{तैलं}}{\lemma{\skt{तैलं}}\Afootnote{K H; \skt{सर्पिः}\ A}}\skt{ श्लेष्मश्वयथौ~। }\edtext{\skt{सन्निपाश्वयथौ}}{\lemma{\skt{सन्निपाश्वयथौ}}\Afootnote{K; \skt{सन्नि\appbrk{}पा\appbrk{}त\appbrk{}श्व\appbrk{}यथौ}\ A H}}\skt{ }\edtext{\skt{स्नुहा\-क्षी\-र\-पा\-त्र\-तुल्यं}}{\lemma{\skt{स्नुहा\appbrk{}क्षी\appbrk{}र\appbrk{}पा\appbrk{}त्र\appbrk{}तुल्यं}}\Afootnote{K H; \skt{स्नुहीक्षीरपात्रं द्वा}\ A}}\skt{ दश\-भि\-र\-म्ल\-पात्रैः~। प्रतिसंसृज्य दन्ती द्रवन्ती प्र }\edtext{\skt{तिवायं}}{\lemma{\skt{तिवायं}}\Afootnote{K H; \skt{तीवापं}\ A}}\skt{ सर्पिः पाचयित्वा }\edtext{\skt{पाययेत्~। विष\-नि\-मि\-त्त\-मा\-क\-ल्पेषु}}{\lemma{\skt{पाययेत्~। विष\appbrk{}नि\appbrk{}मि\appbrk{}त्त\appbrk{}मा\appbrk{}क\appbrk{}ल्पेषु}}\Afootnote{K; \skt{पाययेत्, विषनिमित्तेषु कल्पेषु}\ A; \skt{पाययेत्~। विषनिमित्तस्य कल्पेषु}\ H}}\skt{ प्रतीकारः~॥}\pend

\pstart \ledsidenote{\passref{4.23.12}}\skt{अत }\edtext{\skt{ऊर्ध्वं}}{\lemma{\skt{ऊर्ध्वं}}\Afootnote{A K; \skt{ऊर्द्धन्}\ H}}\skt{ }\edtext{\skt{समान्यान्योगा}}{\lemma{\skt{समान्यान्योगा}}\Afootnote{K; \skt{सामान्यचिकि}\ A; \skt{सामान्यान्योगा}\ H}}\skt{ }\edtext{\skt{वक्ष्यामः~।ति\-ल्व\-क\-घृ\-त\-च\-तु\-र्थानि}}{\lemma{\skt{वक्ष्यामः~।ति\appbrk{}ल्व\appbrk{}क\appbrk{}घृ\appbrk{}त\appbrk{}च\appbrk{}तु\appbrk{}र्थानि}}\Afootnote{K; \skt{त्सित\appbrk{}मु\appbrk{}प\appbrk{}दे\appbrk{}क्ष्यामः तिल्व\appbrk{}क\appbrk{}घृ\appbrk{}त\appbrk{}च\appbrk{}तु\appbrk{}र्थानि}\ A; \skt{न्वक्ष्यामः~। तिल्व\appbrk{}क\appbrk{}घृ\appbrk{}त\appbrk{}च\appbrk{}तु\appbrk{}र्थानि}\ H}}\skt{ }\edtext{\skt{यान्युक्तानि}}{\lemma{\skt{यान्युक्तानि}}\Afootnote{K; \skt{यान्युक्ता}\ A H}}\skt{ }\edtext{\skt{उदरेषु}}{\lemma{\skt{उदरेषु}}\Afootnote{K; \skt{न्युदरेषु}\ A H}}\skt{ }\edtext{\edtext{\skt{ततो\-न्य\-त\-म\-मु\-प\-यु\-ज्य\-मानं}}{\lemma{\skt{ततोन्यतमम्}}\Afootnote{K; \skt{ततोऽन्यतम}\ A; \skt{ततोऽन्यतमम्}\ H}}}{\lemma{\skt{उपयुज्यमानं}}\Afootnote{K H; \skt{नुपयुज्यमानं}\ A}}\skt{ }\edtext{\skt{श्वय\-थू\-न\-प\-हन्ति}}{\lemma{\skt{श्वय\appbrk{}थू\appbrk{}न\appbrk{}प\appbrk{}हन्ति}}\Afootnote{K; \skt{श्वय\appbrk{}थु\appbrk{}म\appbrk{}प\appbrk{}हन्ति,}\ A; \skt{श्वयथूनु पहन्ति~।}\ H}}\skt{~। }\edtext{\skt{मूर्त्त\-व\-र्त्ति\-क्रि\-याम्वा}}{\lemma{\skt{मूर्त्त\appbrk{}व\appbrk{}र्त्ति\appbrk{}क्रि\appbrk{}याम्वा}}\Afootnote{K; \skt{मूत्रवर्तिक्रियां वा}\ A H}}\skt{ सेवेत~। नवायसं }\edtext{\skt{वाह\-र\-ह\-र्म\-धुना}}{\lemma{\skt{वाह\appbrk{}र\appbrk{}ह\appbrk{}र्म\appbrk{}धुना}}\Afootnote{K H; \skt{वाऽह\appbrk{}र\appbrk{}ह\appbrk{}र्म\appbrk{}धुना,}\ A}}\skt{~। }\edtext{\skt{विड\-ङ्गा\-ति\-वि\-षा\-कु\-ट\-ज\-फ\-ल\-भ\-द्र\-दा\-रु\-ना\-ग\-र\-म\-रि\-चानां वा}}{\lemma{\skt{विड\appbrk{}ङ्गा\appbrk{}ति\appbrk{}वि\appbrk{}षा\appbrk{}कु\appbrk{}ट\appbrk{}ज\appbrk{}फ\appbrk{}ल\appbrk{}भ\appbrk{}द्र\appbrk{}दा\appbrk{}रु\appbrk{}ना\appbrk{}ग\appbrk{}र\appbrk{}म\appbrk{}रि\appbrk{}चानां वा}}\Afootnote{K H; \skt{विड\appbrk{}ङ्गा\appbrk{}ति\appbrk{}वि\appbrk{}षा\appbrk{}कु\appbrk{}ट\appbrk{}ज\appbrk{}फ\appbrk{}ल\appbrk{}भ\appbrk{}द्र\appbrk{}दा\appbrk{}रु\appbrk{}ना\appbrk{}ग\appbrk{}र\appbrk{}म\appbrk{}रि\appbrk{}च\appbrk{}चूर्णं वा}\ A}}\skt{ }\edtext{\skt{धर\-ण\-मु\-ष्णा\-म्वुना}}{\lemma{\skt{धर\appbrk{}ण\appbrk{}मु\appbrk{}ष्णा\appbrk{}म्वुना}}\Afootnote{K; \skt{धर\appbrk{}ण\appbrk{}मु\appbrk{}ष्णा\appbrk{}म्बुना,}\ A; \skt{धरणुष्णाम्बुना~।}\ H}}\skt{~। }\edtext{\skt{त्रिक\-टु\-क\-क्षा\-र\-य\-श्चू\-र्णानि}}{\lemma{\skt{त्रिक\appbrk{}टु\appbrk{}क\appbrk{}क्षा\appbrk{}र\appbrk{}य\appbrk{}श्चू\appbrk{}र्णानि}}\Afootnote{K H; \skt{त्रिक\appbrk{}टु\appbrk{}क्षा\appbrk{}रा\appbrk{}य\appbrk{}श्चू\appbrk{}र्णानि वा}\ A}}\skt{ }\edtext{\skt{त्रिफलाक्वाथेन}}{\lemma{\skt{त्रिफलाक्वाथेन}}\Afootnote{K H; \skt{त्रिफ\appbrk{}ला\appbrk{}क\appbrk{}षा\appbrk{}येण,}\ A}}\skt{~। मूत्रम्वा तुल्यक्षीरं~। हरीतकीम्वा तुल्य\-गु\-डा\-मु\-प\-यु\-ञ्जीत~। देवदारुशुण्ठीं वा~। }\edtext{\edtext{\skt{गुग्गुलुम्वा}}{\lemma{\skt{गुग्गुलुम्}}\Afootnote{K H; \skt{गुग्गुलुं वा}\ A}}}{\lemma{\skt{वा}}\Afootnote{K H; \skt{मूत्रेण वर्ष\appbrk{}भू\appbrk{}क\appbrk{}षा\appbrk{}या\appbrk{}नु\appbrk{}पानं वा, तुल्यगुडं वा}\ A}}\skt{ मूत्रेण~। }\edtext{\skt{वर्षा\-भू\-क\-षा\-या\-नु\-पा\-नम्वा}}{\lemma{\skt{वर्षा\appbrk{}भू\appbrk{}क\appbrk{}षा\appbrk{}या\appbrk{}नु\appbrk{}पा\appbrk{}नम्वा}}\Afootnote{K H; \skt{वर्ष\appbrk{}भू\appbrk{}क\appbrk{}षा\appbrk{}या\appbrk{}नु\appbrk{}पानं वा,}\ A}}\skt{ तुल्यगुडं }\edtext{\skt{शृङ्गवेरं}}{\lemma{\skt{शृङ्गवेरं}}\Afootnote{K H; \skt{शृङ्गवेरं वा, वर्षाभूकषायं मूलकल्कं वा सशृङ्गवेरं}\ A}}\skt{~। }\edtext{\skt{व्योषं}}{\lemma{\skt{व्योषं}}\Afootnote{K H; \skt{पयोऽ\appbrk{}नु\appbrk{}पा\appbrk{}न\appbrk{}म\appbrk{}ह\appbrk{}र\appbrk{}ह\appbrk{}र्मासं, व्योष}\ A}}\skt{ वर्षा\-भू\-क\-षा\-य\-सि\-द्धेन वा~। }\edtext{\skt{सर्पि\-षा\-मु\-द्गो\-लुम्वां}}{\lemma{\skt{सर्पि\appbrk{}षा\appbrk{}मु\appbrk{}द्गो\appbrk{}लुम्वां}}\Afootnote{K; \skt{सर्पिषा मुद्गोलुम्बा}\ A; \skt{सर्प्पि\appbrk{}षा\appbrk{}मु\appbrk{}द्ग\appbrk{}स्तु\appbrk{}म्वाम्}\ H}}\skt{ }\edtext{\skt{भक्षयेत्~। पिप्प\-ली\-मू\-ल\-चि\-त्र\-क\-च\-व्या\-म\-यू\-र\-क\-व\-र्षा\-भू\-सिद्धं वा}}{\lemma{\skt{भक्षयेत्~। पिप्प\appbrk{}ली\appbrk{}मू\appbrk{}ल\appbrk{}चि\appbrk{}त्र\appbrk{}क\appbrk{}च\appbrk{}व्या\appbrk{}म\appbrk{}यू\appbrk{}र\appbrk{}क\appbrk{}व\appbrk{}र्षा\appbrk{}भू\appbrk{}सिद्धं वा}}\Afootnote{K H; \skt{न्भक्षयेत्, पिप्प\appbrk{}ली\appbrk{}पि\appbrk{}प्प\appbrk{}ली\appbrk{}मू\appbrk{}ल\appbrk{}च\appbrk{}व्य\appbrk{}चि\appbrk{}त्र\appbrk{}क\appbrk{}म\appbrk{}यू\appbrk{}र\appbrk{}क\appbrk{}व\appbrk{}र्षा\appbrk{}भू\appbrk{}सिद्धं वा}\ A}}\skt{ क्षीरं }\edtext{\skt{पिवेत्~। सहौ\-ष\-ध\-मु\-र\-ङ्गी\-मू\-ल\-सि\-द्धम्वा}}{\lemma{\skt{पिवेत्~। सहौ\appbrk{}ष\appbrk{}ध\appbrk{}मु\appbrk{}र\appbrk{}ङ्गी\appbrk{}मू\appbrk{}ल\appbrk{}सि\appbrk{}द्धम्वा}}\Afootnote{K; \skt{पिबेत्, सहौ\appbrk{}ष\appbrk{}ध\appbrk{}मु\appbrk{}र\appbrk{}ङ्गी\appbrk{}मू\appbrk{}ल\appbrk{}सिद्धं वा,}\ A; \skt{पिबेत्~। महौ\appbrk{}ष\appbrk{}ध\appbrk{}मु\appbrk{}रु\appbrk{}ङ्गी\appbrk{}मू\appbrk{}ल\appbrk{}सि\appbrk{}द्धम्वा~।}\ H}}\skt{~। }\edtext{\edtext{\skt{त्रिक\-चु\-कै\-र\-ण्ड\-श्या\-मा\-मू\-ल\-सि\-द्धम्वा}}{\lemma{\skt{त्रिकचुकैर्}}\Afootnote{K; \skt{त्रिकटुकैर्}\ A; \skt{त्रिकटुकै}\ H}}}{\lemma{\skt{अण्ड\appbrk{}श्या\appbrk{}मा\appbrk{}मू\appbrk{}ल\appbrk{}सि\appbrk{}द्धम्वा}}\Afootnote{A K; \skt{रण्ड\appbrk{}श्या\appbrk{}मा\appbrk{}मू\appbrk{}ल\appbrk{}सि\appbrk{}द्धम्वा~।}\ H}}\skt{~। }\edtext{\skt{वर्षा\-भू\-शृं\-ग\-वे\-र\-दे\-व\-दा\-रु\-सि\-द्धम्वा}}{\lemma{\skt{वर्षा\appbrk{}भू\appbrk{}शृं\appbrk{}ग\appbrk{}वे\appbrk{}र\appbrk{}दे\appbrk{}व\appbrk{}दा\appbrk{}रु\appbrk{}सि\appbrk{}द्धम्वा}}\Afootnote{K H; \skt{वर्षा\appbrk{}भू\appbrk{}शृ\appbrk{}ङ्ग\appbrk{}वे\appbrk{}र\appbrk{}स\appbrk{}हा\appbrk{}दे\appbrk{}व\appbrk{}दा\appbrk{}रु\appbrk{}सिद्धं वा,}\ A}}\skt{~। }\edtext{\skt{तथा\-ला\-वु\-वि\-भी\-त\-क\-फ\-ल\-क\-ल्कम्वा}}{\lemma{\skt{तथा\appbrk{}ला\appbrk{}वु\appbrk{}वि\appbrk{}भी\appbrk{}त\appbrk{}क\appbrk{}फ\appbrk{}ल\appbrk{}क\appbrk{}ल्कम्वा}}\Afootnote{K; \skt{तथाऽ\appbrk{}ला\appbrk{}बु\appbrk{}बि\appbrk{}भी\appbrk{}त\appbrk{}क\appbrk{}फ\appbrk{}ल\appbrk{}कल्कं वा}\ A; \skt{तथा\appbrk{}ला\appbrk{}बु\appbrk{}वि\appbrk{}भी\appbrk{}त\appbrk{}क\appbrk{}फ\appbrk{}ल\appbrk{}क\appbrk{}ल्कम्वा~।}\ H}}\skt{~। }\edtext{\skt{पिप्प\-ली\-म\-रि\-च\-शृ\-ङ्ग\-वे\-रा\-नु\-सि\-द्धेन च}}{\lemma{\skt{पिप्प\appbrk{}ली\appbrk{}म\appbrk{}रि\appbrk{}च\appbrk{}शृ\appbrk{}ङ्ग\appbrk{}वे\appbrk{}रा\appbrk{}नु\appbrk{}सि\appbrk{}द्धेन च}}\Afootnote{K; \skt{तण्डुलाम्बुना; क्षीर\appbrk{}पि\appbrk{}प्प\appbrk{}ली\appbrk{}म\appbrk{}रि\appbrk{}च\appbrk{}शृ\appbrk{}ङ्ग\appbrk{}वे\appbrk{}रा\appbrk{}नु\appbrk{}सि\appbrk{}द्धेन च}\ A; \skt{पिप्प\appbrk{}ली\appbrk{}म\appbrk{}रि\appbrk{}च\appbrk{}शृं\appbrk{}ग\appbrk{}वे\appbrk{}रा\appbrk{}नु\appbrk{}सि\appbrk{}द्धेन}\ H}}\skt{ }\edtext{\skt{मुद्ग\-यू\-षे\-णा\-ल\-व\-णे\-नाल्पं}}{\lemma{\skt{मुद्ग\appbrk{}यू\appbrk{}षे\appbrk{}णा\appbrk{}ल\appbrk{}व\appbrk{}णे\appbrk{}नाल्पं}}\Afootnote{K; \skt{मुद्ग\appbrk{}यू\appbrk{}षे\appbrk{}णा\appbrk{}ल\appbrk{}व\appbrk{}णे\appbrk{}नाल्प}\ A; \skt{मुद्गयुषेण लवणेनाल्पं}\ H}}\skt{ }\edtext{\skt{स्नेहं}}{\lemma{\skt{स्नेहं}}\Afootnote{K H; \skt{स्नेहेन}\ A}}\skt{ }\edtext{\skt{भोजयेत्~। यवगोधूमान्नं}}{\lemma{\skt{भोजयेत्~। यवगोधूमान्नं}}\Afootnote{K H; \skt{भोजयेद्यवान्नं गोधूमान्नं वा;}\ A}}\skt{~। }\edtext{\skt{वृक्ष\-का\-र्क\-न\-क्त\-मा\-ल\-नि\-म्व\-व\-र्षा\-भू\-क्वा\-थश्च}}{\lemma{\skt{वृक्ष\appbrk{}का\appbrk{}र्क\appbrk{}न\appbrk{}क्त\appbrk{}मा\appbrk{}ल\appbrk{}नि\appbrk{}म्व\appbrk{}व\appbrk{}र्षा\appbrk{}भू\appbrk{}क्वा\appbrk{}थश्च}}\Afootnote{K; \skt{वृक्ष\appbrk{}का\appbrk{}र्क\appbrk{}न\appbrk{}क्त\appbrk{}मा\appbrk{}ल\appbrk{}नि\appbrk{}म्ब\appbrk{}व\appbrk{}र्षा\appbrk{}भू\appbrk{}क्वा\appbrk{}थैश्च}\ A; \skt{वृक्ष\appbrk{}का\appbrk{}र्क्क\appbrk{}न\appbrk{}क्त\appbrk{}मा\appbrk{}ल\appbrk{}नि\appbrk{}म्ब\appbrk{}व\appbrk{}र्षा\appbrk{}भू\appbrk{}क्वा\appbrk{}थश्च}\ H}}\skt{ परिषेकः~। }\edtext{\skt{सर्ष\-प\-सौ\-व\-र्च\-ल\-सै\-न्ध\-व\-शार्ङ्गे}}{\lemma{\skt{सर्ष\appbrk{}प\appbrk{}सौ\appbrk{}व\appbrk{}र्च\appbrk{}ल\appbrk{}सै\appbrk{}न्ध\appbrk{}व\appbrk{}शार्ङ्गे}}\Afootnote{K; \skt{सर्ष\appbrk{}प\appbrk{}सु\appbrk{}व\appbrk{}र्च\appbrk{}ला\appbrk{}सै\appbrk{}न्ध\appbrk{}व\appbrk{}शार्ङ्गे}\ A; \skt{सर्षपसौव र्च्चल\appbrk{}सै\appbrk{}न्ध\appbrk{}व\appbrk{}शार्ङ्ग}\ H}}\skt{ ष्टाभिश्च प्रदेहः कार्यः~। यथादोषञ्च }\edtext{\skt{तीक्ष्णा\-नि\-व\-म\-न\-वि\-रे\-च\-ना\-स्था\-प\-ना\-नु\-वा\-स\-नानि}}{\lemma{\skt{तीक्ष्णा\appbrk{}नि\appbrk{}व\appbrk{}म\appbrk{}न\appbrk{}वि\appbrk{}रे\appbrk{}च\appbrk{}ना\appbrk{}स्था\appbrk{}प\appbrk{}ना\appbrk{}नु\appbrk{}वा\appbrk{}स\appbrk{}नानि}}\Afootnote{K H; \skt{वम\appbrk{}न\appbrk{}वि\appbrk{}रे\appbrk{}च\appbrk{}ना\appbrk{}स्था\appbrk{}प\appbrk{}नानि ती}\ A}}\skt{ }\edtext{\skt{सत\-त\-मु\-प\-से\-वेत}}{\lemma{\skt{सततम्}}\Afootnote{K H; \skt{क्ष्णान्यजस्रम्}\ A}}\skt{~। }\edtext{\skt{स्वेदोपनाहौ च}}{\lemma{\skt{स्वेदोपनाहौ च}}\Afootnote{K H; \skt{स्नेह\appbrk{}स्वे\appbrk{}दो\appbrk{}प\appbrk{}ना\appbrk{}हांश्च;}\ A}}\skt{~। }\edtext{\edtext{\skt{सिराश्चाभीक्ष्णं}}{\lemma{\skt{सिराश्}}\Afootnote{K; \skt{सिराभिश्}\ A; \skt{सिरा}\ H}}}{\lemma{\skt{चाभीक्ष्णं}}\Afootnote{A K; \skt{च्चाभीक्ष्णं}\ H}}\skt{ }\edtext{\skt{शोणि\-त\-म\-व\-से\-च\-येत्~॥ अन्य\-त्र\-पा\-ण्डु\-शो\-फा\-दिति}}{\lemma{\skt{शोणि\appbrk{}त\appbrk{}म\appbrk{}व\appbrk{}से\appbrk{}च\appbrk{}येत्~॥ अन्य\appbrk{}त्र\appbrk{}पा\appbrk{}ण्डु\appbrk{}शो\appbrk{}फा\appbrk{}दिति}}\Afootnote{K H; \skt{शोणि\appbrk{}त\appbrk{}म\appbrk{}व\appbrk{}से\appbrk{}च\appbrk{}ये\appbrk{}द\appbrk{}न्य\appbrk{}त्रो\appbrk{}प\appbrk{}द्र\appbrk{}व\appbrk{}शो\appbrk{}फा\appbrk{}दिति~॥}\ A}}\skt{~॥}\pend

\pstart\versepar \ledsidenote{\passref{4.23.13}}\edtext{\skt{पिष्टान्नमम्व्लं}}{\lemma{\skt{पिष्टान्नमम्व्लं}}\Afootnote{K; \skt{भवति चात्र पिष्टान्नमम्लं}\ A; \skt{पिष्टान्नमम्लं}\ H}}\skt{ लवणानि मद्यं मृदं }\edtext{\skt{दिवा\-स्व\-प्न\-म\-जा\-ङ्ग\-लञ्च}}{\lemma{\skt{दिवा\appbrk{}स्व\appbrk{}प्न\appbrk{}म\appbrk{}जा\appbrk{}ङ्ग\appbrk{}लञ्च}}\Afootnote{K; \skt{दिवा\appbrk{}स्व\appbrk{}प्न\appbrk{}म\appbrk{}जा\appbrk{}ङ्गलं च~।}\ A; \skt{दिवा\appbrk{}स्व\appbrk{}प्न\appbrk{}म\appbrk{}ज\appbrk{}ङ्ग\appbrk{}लञ्च~।}\ H}}\skt{~। }\pend
\pstart\versepar \skt{स्त्रियो }\edtext{\skt{गुडं}}{\lemma{\skt{गुडं}}\Afootnote{K H; \skt{घृतं}\ A}}\skt{ }\edtext{\skt{तैलमथो}}{\lemma{\skt{तैलमथो}}\Afootnote{K; \skt{तैलपयो}\ A; \skt{तेलमथो}\ H}}\skt{ गुरूणि शोफं }\edtext{\skt{जिघांसु}}{\lemma{\skt{जिघांसु}}\Afootnote{K; \skt{जिघांसुः}\ A H}}\skt{ }\edtext{\skt{परिवर्जयीत}}{\lemma{\skt{परिवर्जयीत}}\Afootnote{K H; \skt{परिवर्जयेत्तु~॥}\ A}}\skt{~॥}\pend

\pstart \ledsidenote{\passref{4.23.trailer}}\edtext{\skt{चि२३॥}}{\lemma{\skt{चि}}\Afootnote{K; \skt{इति सुश्रु\appbrk{}त\appbrk{}सं\appbrk{}हि\appbrk{}तायां चि}\ A; \emph{om.}\ H}}\skt{ }\edtext{\skt{ओ}}{\lemma{\skt{ओ}}\Afootnote{K; \skt{कित्सास्थाने सोथचिकित्सितं नाम त्रयोविंशोऽध्यायः~॥२३~॥}\ A; \skt{॥ चिकित्साथं त्रि~॥}\ H}}\skt{~॥}\pend

\endnumbering
\beginnumbering
\eledsection*{Adhyāya 24}
\pstart \ledsidenote{\passref{4.24.1}}\edtext{\skt{अथा\-तोऽ\-ना\-ग\-ता\-बा\-ध\-चि\-कि\-त्सितं}}{\lemma{\skt{ऽना\appbrk{}ग\appbrk{}ता\appbrk{}बा\appbrk{}ध\appbrk{}चि\appbrk{}कि\appbrk{}त्सितं}}\Afootnote{\skt{ऽना\appbrk{}ग\appbrk{}ता\appbrk{}बा\appbrk{}धा\appbrk{}प्र\appbrk{}ति\appbrk{}षेधं}\ A; \skt{नाग\appbrk{}ता\appbrk{}वा\appbrk{}ध\appbrk{}चि\appbrk{}कि\appbrk{}त्सितं}\ K H}}\skt{ }\edtext{\skt{व्याख्यास्यामः}}{\lemma{\skt{व्याख्यास्यामः}}\Afootnote{A H; \skt{व्या व्याख्यास्यामः~।}\ K}}\skt{~।}\pend

\pstart\versepar \ledsidenote{\passref{4.24.3}}\edtext{\skt{उत्था\-यो\-त्था\-या\-त्म\-वता}}{\lemma{\skt{उत्था\appbrk{}यो\appbrk{}त्था\appbrk{}या\appbrk{}त्म\appbrk{}वता}}\Afootnote{K H; \skt{उत्थायोत्थाय सततं}\ A}}\skt{ }\edtext{\edtext{\skt{स्वस्थे\-ना\-रो\-ग\-मि\-च्छता}}{\lemma{\skt{स्वस्थेनारोगम्}}\Afootnote{K H; \skt{स्वस्थेनारोग्यम्}\ A}}}{\lemma{\skt{इच्छता}}\Afootnote{K H; \skt{इच्छता~। धीमता}\ A}}\skt{~। }\pend
\pstart\versepar \skt{}\edtext{\edtext{\skt{यद\-नु\-ष्ठे\-य\-म\-स्कन्नं}}{\lemma{\skt{यदनुष्ठेयम्}}\Afootnote{A K; \skt{यदनूष्ठेयम्}\ H}}}{\lemma{\skt{अस्कन्नं}}\Afootnote{K H; \emph{om.}\ A}}\skt{ }\edtext{\skt{तत्सर्वं}}{\lemma{\skt{तत्सर्वं}}\Afootnote{A K; \skt{तत्सर्व्व}\ H}}\skt{ }\edtext{\skt{संप्रवक्ष्यते}}{\lemma{\skt{संप्रवक्ष्यते}}\Afootnote{A K; \skt{न्सम्प्रवक्ष्यते~॥}\ H}}\skt{~॥१॥}\pend

\pstart\versepar \ledsidenote{\passref{4.24.4}}\edtext{\skt{तत्रादितो}}{\lemma{\skt{तत्रादितो}}\Afootnote{K H; \skt{तत्रादौ दन्}\ A}}\skt{ }\edtext{\skt{दन्तकाष्ठं}}{\lemma{\skt{दन्तकाष्ठं}}\Afootnote{K H\textsuperscript{pc}; \skt{तपवनं}\ A; \skt{दन्तकाष्ठं म्}\ H\textsuperscript{ac}}}\skt{ }\edtext{\skt{द्वाद\-शां\-गु\-ल\-मा\-यतं}}{\lemma{\skt{द्वाद\appbrk{}शां\appbrk{}गु\appbrk{}ल\appbrk{}मा\appbrk{}यतं}}\Afootnote{A K; \skt{वाद\appbrk{}शा\appbrk{}ङ्गु\appbrk{}ल\appbrk{}मा\appbrk{}यतं~।}\ H\textsuperscript{ac}; \skt{द्वाद\appbrk{}शा\appbrk{}ङ्गु\appbrk{}ल\appbrk{}मा\appbrk{}यतं~।}\ H\textsuperscript{pc}}}\skt{~। }\pend
\pstart\versepar \skt{}\edtext{\edtext{\edtext{\skt{कनी\-नि\-का\-प\-रि\-णा\-ह\-मृ\-ज्व\-ग्र\-न्थि\-म\-था\-व्रणं}}{\lemma{\skt{कनी\appbrk{}नि\appbrk{}का\appbrk{}प\appbrk{}रि\appbrk{}णा\appbrk{}हम्}}\Afootnote{K H; \skt{कनि\appbrk{}ष्ठि\appbrk{}का\appbrk{}प\appbrk{}री\appbrk{}णा\appbrk{}हम्}\ A}}}{\lemma{\skt{ऋज्वग्रन्थिम्}}\Afootnote{K H; \skt{ऋज्वग्रन्थि}\ A}}}{\lemma{\skt{अथाव्रणं}}\Afootnote{K H; \skt{तमव्रणम्~॥}\ A}}\skt{~॥२॥}\pend

\pstart \ledsidenote{\passref{4.24.5}}\edtext{\edtext{\skt{अयुग्मग्रन्थि यं}}{\lemma{\skt{अयुग्मग्रन्थि यं}}\Afootnote{\skt{अयुग्मग्रन्थि य}\ A K; \skt{अयुग्मग्रन्}\ H}}\skt{ }\edtext{\skt{सक्षीरञ्च}}{\lemma{\skt{सक्षीरञ्च}}\Afootnote{K; \skt{च्चापि प्रत्य}\ A}}\skt{ }\edtext{\skt{प्रशस्यते}}{\lemma{\skt{प्रशस्यते}}\Afootnote{K; \skt{ग्रं शस्तभूमिजम्~। अवेक्ष्यर्तुं च दोषं च रसं वीर्यं च योजयेत्~॥}\ A}}\skt{~।}}{\lemma{\skt{थि… प्रशस्यते}}\Afootnote{\emph{om.}\ H}}\pend

\pstart\versepar \ledsidenote{\passref{4.24.6}}\edtext{\skt{कषा\-य\-क\-टु\-तिक्तं वा}}{\lemma{\skt{कषा\appbrk{}य\appbrk{}क\appbrk{}टु\appbrk{}तिक्तं वा}}\Afootnote{K H; \skt{कषायं मधुरं तिक्तं}\ A}}\skt{ }\edtext{\skt{तत्त\-त्खा\-दे\-त्प्रा\-त\-रु\-त्थितः}}{\lemma{\skt{तत्त\appbrk{}त्खा\appbrk{}दे\appbrk{}त्प्रा\appbrk{}तर्}}\Afootnote{\skt{कटुकं प्रातर्}\ A; \skt{तत\appbrk{}त्खा\appbrk{}दे\appbrk{}त्प्रा\appbrk{}तर्}\ K; \skt{तत्खादेत्प्रातर्}\ H}}\skt{~॥३॥ }\pend
\pstart\versepar \skt{}\edtext{\skt{मधु\-त्रि\-क\-टु\-का\-क्तम्वा }\edtext{\skt{गोमूत्रेणापि}}{\lemma{\skt{गोमूत्रेणापि}}\Afootnote{K; \skt{गोमूत्रेणाति}\ H}}\skt{ भावितं~।}}{\lemma{\skt{मधु\appbrk{}त्रि\appbrk{}क\appbrk{}टु\appbrk{}का\appbrk{}क्तम्… भावितं}}\Afootnote{\emph{om.}\ A}}\pend

\pstart \ledsidenote{\passref{4.24.9}}\edtext{\skt{तद्दौ\-र्ग्ग\-न्ध्यो\-प\-देहौ तु}}{\lemma{\skt{तद्दौ\appbrk{}र्ग्ग\appbrk{}न्ध्यो\appbrk{}प\appbrk{}देहौ तु}}\Afootnote{K H\textsuperscript{pc}; \skt{दन्त\appbrk{}शो\appbrk{}ध\appbrk{}न\appbrk{}चू\appbrk{}र्णेन दन्त\appbrk{}मां\appbrk{}सा\appbrk{}न्य\appbrk{}बा\appbrk{}ध\appbrk{}यन्~। तद्दौ\appbrk{}र्ग\appbrk{}न्ध्यो\appbrk{}प\appbrk{}देहौ तु}\ A; \skt{तद्दौ\appbrk{}र्ग्ग\appbrk{}न्ध्यो\appbrk{}प\appbrk{}देहौ}\ H\textsuperscript{ac}}}\skt{ श्लेष्माणं चापकर्षति~॥४॥}\pend

\pstart\versepar \ledsidenote{\passref{4.24.10}}\edtext{\edtext{\skt{वैश\-द्य\-म\-न्नेऽ\-भि\-रुचिं}}{\lemma{\skt{अन्ने}}\Afootnote{K H; \skt{अन्ना}\ A}}}{\lemma{\skt{ऽभिरुचिं}}\Afootnote{H; \skt{भिरुचिं}\ A; \skt{भिरुचिं}\ K}}\skt{ सौमनस्यं }\edtext{\skt{करोति च}}{\lemma{\skt{करोति च}}\Afootnote{K\textsuperscript{pc} H\textsuperscript{pc}; \skt{करोति}\ A K\textsuperscript{ac} H\textsuperscript{ac}}}\skt{~। }\pend
\pstart\versepar \skt{स्निग्धां }\edtext{\skt{शुक्लांश्च}}{\lemma{\skt{शुक्लांश्च}}\Afootnote{K\textsuperscript{pc}; \skt{शुष्कांश्च}\ H\textsuperscript{pc}}}\skt{ दशनां जिह्वायापि च लाघवं~॥५॥ }\pend
\pstart\versepar \skt{न }\edtext{\skt{चापि\-ग\-ल\-ता\-ल्वौ\-ष्ठ\-जि\-ह्वा\-रो\-ग\-स\-मु\-द्भवः}}{\lemma{\skt{चापि\appbrk{}ग\appbrk{}ल\appbrk{}ता\appbrk{}ल्वौ\appbrk{}ष्ठ\appbrk{}जि\appbrk{}ह्वा\appbrk{}रो\appbrk{}ग\appbrk{}स\appbrk{}मु\appbrk{}द्भवः}}\Afootnote{K; \skt{खादे\appbrk{}द्ग\appbrk{}ल\appbrk{}ता\appbrk{}ल्वो\appbrk{}ष्ठ\appbrk{}जि\appbrk{}ह्वा\appbrk{}रो\appbrk{}ग\appbrk{}स\appbrk{}मु\appbrk{}द्भवे~॥}\ A; \skt{चाति\appbrk{}ग\appbrk{}ल\appbrk{}ता\appbrk{}ल्वो\appbrk{}ष्ठ\appbrk{}जि\appbrk{}ह्वा\appbrk{}रो\appbrk{}ग\appbrk{}स\appbrk{}मु\appbrk{}द्भवः~।}\ H}}\skt{~।}\edtext{}{\lemma{\skt{च… च}}\Afootnote{\emph{om.}\ A}}\edtext{}{\lemma{\skt{च… च}}\Afootnote{\emph{om.}\ K\textsuperscript{ac}}}\edtext{}{\lemma{\skt{च… च}}\Afootnote{\emph{om.}\ H\textsuperscript{ac}}}\edtext{}{\lemma{\skt{लाघवं~॥५॥ न}}\Afootnote{K\textsuperscript{pc} H\textsuperscript{pc}; \skt{न}\ A K\textsuperscript{ac} H\textsuperscript{ac}}}\pend

\pstart\versepar \ledsidenote{\passref{4.24.11}}\edtext{\skt{न}}{\lemma{\skt{न}}\Afootnote{K H; \skt{अथास्यपाके श्वासे च}\ A}}\skt{ }\edtext{\skt{खादेत्मुखशोषी तु}}{\lemma{\skt{खादेत्मुखशोषी तु}}\Afootnote{K H; \skt{का}\ A}}\skt{ }\edtext{\skt{श्वास\-हि\-क्का\-व\-मीषु च}}{\lemma{\skt{श्वास\appbrk{}हि\appbrk{}क्का\appbrk{}व\appbrk{}मीषु च}}\Afootnote{K H; \skt{सहिक्कावमीषु च~।}\ A}}\skt{~॥६॥ }\pend
\pstart\versepar \skt{}\edtext{\edtext{\skt{दुर्ब\-लोऽ\-जी\-र्ण्ण\-भ\-क्तश्च}}{\lemma{\skt{दुर्बलो}}\Afootnote{A; \skt{दुर्वलो}\ K H}}}{\lemma{\skt{ऽजीर्ण्णभक्तश्च}}\Afootnote{A; \skt{जीर्ण्णभक्तश्च}\ K H}}\skt{ मूर्च्छार्त्तो }\edtext{\skt{मदपीडितः}}{\lemma{\skt{मदपीडितः}}\Afootnote{K H; \skt{मदपीषितः~॥}\ A}}\skt{~।}\pend

\pstart\versepar \ledsidenote{\passref{4.24.12}}\edtext{\edtext{\skt{शिरो\-रु\-जां\-र्त्त\-स्तृ\-षितः}}{\lemma{\skt{शिरोरुजांर्त्तस्}}\Afootnote{K; \skt{शिरोरुजार्त}\ A; \skt{शिरोरुजार्त्तस्}\ H}}}{\lemma{\skt{तृषितः}}\Afootnote{K H; \skt{स्तृषितः}\ A}}\skt{ }\edtext{\skt{श्रान्त}}{\lemma{\skt{श्रान्त}}\Afootnote{K H\textsuperscript{ac}; \skt{श्रान्तः}\ A; \skt{श्रान्तयाश्रान्त}\ H\textsuperscript{pc}}}\skt{ पानक्लमान्वितः~। }\pend
\pstart\versepar \skt{अर्दिती कर्ण्णशूली च दन्तरोगी च मानवः~॥७॥}\pend

\pstart\versepar \ledsidenote{\passref{4.24.13}}\skt{जिह्वानिर्लेखनं रौप्यं सौवर्ण्णं }\edtext{\skt{वार्क्षमेव वा}}{\lemma{\skt{एव वा}}\Afootnote{K; \skt{एव च~।}\ A H}}\skt{~। }\pend
\pstart\versepar \skt{}\edtext{\skt{तन्मलापहरं}}{\lemma{\skt{तन्मलापहरं}}\Afootnote{H; \skt{तन्मलापहरं}\ A K}}\skt{ }\edtext{\skt{शस्तं}}{\lemma{\skt{शस्तं}}\Afootnote{A K H\textsuperscript{ac}; \skt{शनस्तं}\ H\textsuperscript{pc}}}\skt{ मृदुश्लक्ष्णं दशांगुलं~॥८॥}\pend

\pstart\versepar \ledsidenote{\passref{4.24.14}}\edtext{\skt{मुखवैरस्य}}{\lemma{\skt{मुखवैरस्य}}\Afootnote{K H; \skt{मुखवैरस्य}\ A}}\skt{ }\edtext{\skt{वैगन्ध्यं}}{\lemma{\skt{वैगन्ध्यं}}\Afootnote{K; \skt{दौर्गन्ध्य}\ A; \skt{वैगन्ध्य}\ H}}\skt{ }\edtext{\skt{शोफजाड्यापहं}}{\lemma{\skt{शोफजाड्यापहं}}\Afootnote{K H; \skt{शोफजाड्यहरं}\ A}}\skt{ सुखं~। }\pend
\pstart\versepar \skt{}\edtext{\skt{दृढीकरञ्च}}{\lemma{\skt{दृढीकरञ्च}}\Afootnote{K H; \skt{दन्तदार्ढ्यक}\ A}}\skt{ }\edtext{\skt{दन्तानां}}{\lemma{\skt{दन्तानां}}\Afootnote{K H; \skt{रं रुच्यं}\ A}}\skt{ स्नेह\-ग\-ण्डू\-ष\-धा\-रणं~॥९॥}\pend

\pstart \ledsidenote{\passref{4.24.21}}\edtext{\skt{ताम्बूलपत्रं}}{\lemma{\skt{ताम्बूलपत्रं}}\Afootnote{H; \skt{कर्पू\appbrk{}र\appbrk{}जा\appbrk{}ती\appbrk{}क\appbrk{}क्को\appbrk{}ल\appbrk{}ल\appbrk{}व\appbrk{}ङ्ग\appbrk{}क\appbrk{}टु\appbrk{}का\appbrk{}ह्वयैः~। सचूर्णपूगैः सहितं पत्रं ताम्बुल}\ A; \skt{ताम्वूलपत्रं}\ K}}\skt{ }\edtext{\skt{सक्षारं}}{\lemma{\skt{सक्षारं}}\Afootnote{K H; \skt{जं}\ A}}\skt{ }\edtext{\skt{पूगं}}{\lemma{\skt{पूगं}}\Afootnote{K H; \skt{शुभम्~॥}\ A}}\skt{ }\edtext{\skt{हृद्यं}}{\lemma{\skt{हृद्यं}}\Afootnote{K H; \emph{om.}\ A}}\skt{ }\edtext{\skt{कफापहं}}{\lemma{\skt{कफापहं}}\Afootnote{K H; \emph{om.}\ A}}\skt{~।}\pend

\pstart \ledsidenote{\passref{4.24.22}}\edtext{\skt{मुख\-वै\-श\-द्य\-सौ\-ग\-न्ध्य\-का\-न्ति\-सौ\-ष्ठ\-व\-का\-रकं}}{\lemma{\skt{मुख\appbrk{}वै\appbrk{}श\appbrk{}द्य\appbrk{}सौ\appbrk{}ग\appbrk{}न्ध्य\appbrk{}का\appbrk{}न्ति\appbrk{}सौ\appbrk{}ष्ठ\appbrk{}व\appbrk{}का\appbrk{}रकं}}\Afootnote{K; \skt{मुख\appbrk{}वै\appbrk{}श\appbrk{}द्य\appbrk{}सौ\appbrk{}ग\appbrk{}न्ध्य\appbrk{}का\appbrk{}न्ति\appbrk{}सौ\appbrk{}ष्ठ\appbrk{}व\appbrk{}का\appbrk{}र\appbrk{}कम्~। हनु\appbrk{}द\appbrk{}न्त\appbrk{}स्व\appbrk{}र\appbrk{}म\appbrk{}ल\appbrk{}जि\appbrk{}ह्वे\appbrk{}न्द्रि\appbrk{}य\appbrk{}वि\appbrk{}शो\appbrk{}ध\appbrk{}नम्~॥}\ A; \skt{मुख\appbrk{}वै\appbrk{}ष\appbrk{}द्य\appbrk{}सौ\appbrk{}ग\appbrk{}न्ध्य\appbrk{}का\appbrk{}न्ति\appbrk{}सौ\appbrk{}ष्ठ\appbrk{}व\appbrk{}का\appbrk{}रकं~॥}\ H}}\skt{~॥१०॥}\pend

\pstart \ledsidenote{\passref{4.24.23}}\edtext{\skt{पथ्यं}}{\lemma{\skt{पथ्यं}}\Afootnote{K H; \skt{प्रसेकशमनं हृद्यं गला\appbrk{}म\appbrk{}य\appbrk{}वि\appbrk{}ना\appbrk{}श\appbrk{}नम्~। पथ्यं}\ A}}\skt{ सुप्तोत्थिते भुक्ते }\edtext{\skt{वान्ते}}{\lemma{\skt{वान्ते}}\Afootnote{K H; \skt{स्नाते}\ A}}\skt{ }\edtext{\skt{स्रान्ते च}}{\lemma{\skt{स्रान्ते च}}\Afootnote{K H; \skt{वान्ते च}\ A}}\skt{ }\edtext{\skt{मानवे}}{\lemma{\skt{मानवे}}\Afootnote{A K; \skt{मानवः~।}\ H}}\skt{~।}\pend

\pstart \ledsidenote{\passref{4.24.24}}\edtext{\edtext{\skt{रूक्ष\-दु\-र्ब\-ल\-म\-त्ता\-ना\-म\-हितं}}{\lemma{\skt{रूक्ष\appbrk{}दु\appbrk{}र्ब\appbrk{}ल\appbrk{}म\appbrk{}त्ता\appbrk{}नाम्}}\Afootnote{\skt{रक्त\appbrk{}पि\appbrk{}त्त\appbrk{}क्ष\appbrk{}त\appbrk{}क्षी\appbrk{}ण\appbrk{}तृ\appbrk{}ष्णा\appbrk{}मू\appbrk{}र्च्छा\appbrk{}प\appbrk{}री\appbrk{}ति\appbrk{}नाम्~। रूक्ष\appbrk{}दु\appbrk{}र्ब\appbrk{}ल\appbrk{}म\appbrk{}र्त्यानां}\ A; \skt{रूक्ष\appbrk{}दु\appbrk{}र्व\appbrk{}ल\appbrk{}म\appbrk{}त्ता\appbrk{}नाम्}\ K; \skt{रूक्ष\appbrk{}दु\appbrk{}र्व्व\appbrk{}ल\appbrk{}म\appbrk{}र्त्यानां}\ H}}}{\lemma{\skt{अहितं}}\Afootnote{K; \skt{न हितं}\ A; \skt{महितञ्}\ H}}\skt{ चास्यशोषिणां~॥११॥}\pend

\pstart\versepar \ledsidenote{\passref{4.24.25}}\edtext{\skt{केश\-भू\-मि\-ग\-ता\-न्रोगां शिरोऽ\-भ्य\-ङ्गोऽ\-प\-क\-र्षति}}{\lemma{\skt{केश\appbrk{}भू\appbrk{}मि\appbrk{}ग\appbrk{}ता\appbrk{}न्रोगां शिरोऽ\appbrk{}भ्य\appbrk{}ङ्गोऽ\appbrk{}प\appbrk{}क\appbrk{}र्षति}}\Afootnote{\skt{शिरोगतांस्तथा रोगा\appbrk{}ञ्छि\appbrk{}रो\appbrk{}भ\appbrk{}ङ्गोऽ\appbrk{}प\appbrk{}क\appbrk{}र्षति~।}\ A; \skt{केश\appbrk{}भू\appbrk{}मि\appbrk{}ग\appbrk{}ता\appbrk{}न्रोगां शिरो\appbrk{}भ्य\appbrk{}ङ्गो\appbrk{}प\appbrk{}क\appbrk{}र्षति~।}\ K; \skt{केश\appbrk{}भू\appbrk{}मि\appbrk{}ग\appbrk{}ता\appbrk{}न्रो\appbrk{}गा\appbrk{}न्शि\appbrk{}रो\appbrk{}भ्य\appbrk{}ङ्गो\appbrk{}प\appbrk{}क\appbrk{}र्षति~।}\ H}}\skt{~। }\pend
\pstart\versepar \skt{केशानां }\edtext{\skt{मृदुता}}{\lemma{\skt{मृदुता}}\Afootnote{K H\textsuperscript{ac}; \skt{मार्दवं दैर्घ्यं बहुत्वं}\ A; \skt{मृदुता स्निग्धं दृ}\ H\textsuperscript{pc}}}\skt{ }\edtext{\skt{स्निग्धं}}{\lemma{\skt{स्निग्धं}}\Afootnote{K H\textsuperscript{ac}; \skt{स्निग्ध}\ A; \emph{om.}\ H\textsuperscript{pc}}}\skt{ }\edtext{\skt{रुक्प्रशान्तिं}}{\lemma{\skt{रुक्प्रशान्तिं}}\Afootnote{K H\textsuperscript{pc}; \skt{कृष्णताम्~॥}\ A; \skt{दृक्प्रशान्तिङ्}\ H\textsuperscript{ac}}}\skt{ }\edtext{\skt{करोति}}{\lemma{\skt{करोति}}\Afootnote{K H; \emph{om.}\ A}}\skt{ }\edtext{\skt{च}}{\lemma{\skt{च}}\Afootnote{K H; \emph{om.}\ A}}\skt{~॥१२॥}\pend

\pstart\versepar \ledsidenote{\passref{4.24.26}}\skt{करोति }\edtext{\edtext{\skt{शिरसस्तृप्तिं}}{\lemma{\skt{शिरसस्}}\Afootnote{K H; \skt{शिरस}\ A}}}{\lemma{\skt{तृप्तिं}}\Afootnote{K H; \skt{स्तृप्तिं}\ A}}\skt{ }\edtext{\skt{केशानां}}{\lemma{\skt{केशानां}}\Afootnote{K H; \emph{om.}\ A}}\skt{ }\edtext{\skt{दृढतामपि}}{\lemma{\skt{दृढताम्}}\Afootnote{K H; \skt{सुत्वक्कम्}\ A}}\skt{~। }\pend
\pstart\versepar \skt{}\edtext{\edtext{\skt{सन्त\-र्प\-ण\-श्चे\-न्द्रि\-याणां}}{\lemma{\skt{सन्तर्पणश्}}\Afootnote{K; \skt{चाननम्~। सन्तर्पणं}\ A; \skt{सन्तर्पण}\ H}}}{\lemma{\skt{चेन्द्रियाणां}}\Afootnote{A K; \skt{श्चेन्द्रियाणां}\ H}}\skt{ }\edtext{\skt{दन्तोऽभ्यङ्गस्तु}}{\lemma{\skt{दन्तोऽभ्यङ्गस्तु}}\Afootnote{\skt{शिरसः प्रति}\ A; \skt{दन्तोभ्यङ्गस्तु}\ K; \skt{दलेभ्यङ्गस्तु}\ H}}\skt{ }\edtext{\skt{मूर्ध्नि}}{\lemma{\skt{मूर्ध्नि}}\Afootnote{K H; \skt{पूर}\ A}}\skt{ }\edtext{\skt{तु}}{\lemma{\skt{तु}}\Afootnote{K H; \skt{णम्~॥}\ A}}\skt{~॥१३॥}\pend

\pstart\versepar \ledsidenote{\passref{4.24.29}}\skt{केशप्रसाधनी केश्या रजोजन्तु मलापहा~। }\pend
\pstart\versepar \skt{हनु\-म\-न्या\-शि\-रः\-क\-र्ण्ण\-शू\-लघ्नं कर्ण्णपूरणं~॥१४॥}\pend

\pstart\versepar \ledsidenote{\passref{4.24.30}}\skt{अभ्यङ्गो मार्दवकरः कफ\-वा\-त\-नि\-रो\-धनः~। }\pend
\pstart\versepar \skt{}\edtext{\skt{धाऊ\-ना\-म्पु\-ष्टि\-ज\-ननो}}{\lemma{\skt{धाऊनाम्}}\Afootnote{K; \skt{धातूनां}\ A H}}\skt{ }\edtext{\skt{मृजा\-व\-र्ण्ण\-ब\-ल\-प्रदः}}{\lemma{\skt{मृजा\appbrk{}व\appbrk{}र्ण्ण\appbrk{}ब\appbrk{}ल\appbrk{}प्रदः}}\Afootnote{A; \skt{मृजा\appbrk{}व\appbrk{}र्ण्ण\appbrk{}व\appbrk{}ल\appbrk{}प्रदः}\ K H}}\skt{~॥१५॥}\pend

\pstart\versepar \ledsidenote{\passref{4.24.31}}\skt{सेकः }\edtext{\skt{श्रमघ्नोनिलहा}}{\lemma{\skt{श्रमघ्नोनिलहा}}\Afootnote{K; \skt{श्रमघ्नोऽनिलहृ}\ A; \skt{श्रमघ्नोऽनिलहा}\ H}}\skt{ }\edtext{\skt{भग्न\-स\-न्धि\-प्र\-सा\-धकः}}{\lemma{\skt{भग्न\appbrk{}स\appbrk{}न्धि\appbrk{}प्र\appbrk{}सा\appbrk{}धकः}}\Afootnote{K H; \skt{द्भग्न\appbrk{}स\appbrk{}न्धि\appbrk{}प्र\appbrk{}सा\appbrk{}धकः~।}\ A}}\skt{~। }\pend
\pstart\versepar \skt{क्षता\-ग्नि\-द\-ग्धा\-भि\-ह\-त\-वि\-घृ\-ष्टानां }\edtext{\skt{रुजापहः}}{\lemma{\skt{रुजापहः}}\Afootnote{A K; \skt{रुजापहाः~।}\ H}}\skt{~॥१६॥}\pend

\pstart\versepar \ledsidenote{\passref{4.24.32}}\edtext{\skt{सिक्त\-स्या\-भि\-प्र\-व\-र्धन्ते}}{\lemma{\skt{सिक्त\appbrk{}स्या\appbrk{}भि\appbrk{}प्र\appbrk{}व\appbrk{}र्धन्ते}}\Afootnote{K H; \skt{जलसिक्तस्य वर्धन्ते}\ A}}\skt{ यथा }\edtext{\edtext{\skt{पुष्पोऽङ्कुरस्तरोः}}{\lemma{\skt{पुष्पो}}\Afootnote{K; \skt{मूले}\ A; \skt{पुष्पोङ्}\ H}}}{\lemma{\skt{ऽङ्कुरस्तरोः}}\Afootnote{\skt{ऽङ्कुरास्तरोः~।}\ A; \skt{ङ्कुरस्तरोः~।}\ K; \skt{करस्तरोः~।}\ H}}\skt{~। }\pend
\pstart\versepar \skt{तथा }\edtext{\skt{धातुविवृद्धिर्हि}}{\lemma{\skt{धातुविवृद्धिर्हि}}\Afootnote{\skt{धातुविवृद्धिर्हि}\ A K; \skt{धातुर्विवृद्धिर्हि}\ H}}\skt{ स्नेहसिक्तस्य जायते~॥१७॥}\pend

\pstart\versepar \ledsidenote{\passref{4.24.33}}\edtext{\edtext{\skt{शिरा\-मु\-खै\-र्लो\-म\-कू\-पै\-र्ध\-म\-नी\-भिश्च}}{\lemma{\skt{शिरामुखैर्}}\Afootnote{K H; \skt{सिरामुखै}\ A}}}{\lemma{\skt{लोमकूपैर्}}\Afootnote{K H; \skt{रोमकूपैर्}\ A}}\skt{ तर्पयन्~। }\pend
\pstart\versepar \skt{शरीरे }\edtext{\skt{बलमाधत्ते}}{\lemma{\skt{बलम्}}\Afootnote{A; \skt{वलम्}\ K H}}\skt{ }\edtext{\skt{स्नेहो वै}}{\lemma{\skt{स्नेहो वै}}\Afootnote{K H; \skt{युक्तः स्नेहो}\ A}}\skt{ }\edtext{\skt{साधु}}{\lemma{\skt{साधु}}\Afootnote{K H; \emph{om.}\ A}}\skt{ }\edtext{\skt{योजितः}}{\lemma{\skt{योजितः}}\Afootnote{K H; \skt{ऽवगाहने~॥}\ A}}\skt{~॥१८॥}\edtext{}{\lemma{\skt{तर्पयन्~। शरीरे}}\Afootnote{K H; \skt{तर्पयम्~। शरीर}\ A}}\pend

\pstart\versepar \ledsidenote{\passref{4.24.34}}\skt{तत्र प्रकृ\-ति\-सा\-त्म्य\-र्तु\-दे\-श\-दो\-ष\-वि\-का\-र\-वित्~। }\pend
\pstart\versepar \skt{तैलं घृतम्वा }\edtext{\skt{मति\-मा\-न्यु\-ञ्जा\-द\-भ्य\-ङ्ग\-से\-कयोः}}{\lemma{\skt{मति\appbrk{}मा\appbrk{}न्यु\appbrk{}ञ्जा\appbrk{}द\appbrk{}भ्य\appbrk{}ङ्ग\appbrk{}से\appbrk{}कयोः}}\Afootnote{K; \skt{मति\appbrk{}मा\appbrk{}न्यु\appbrk{}ञ्ज्या\appbrk{}द\appbrk{}भ्य\appbrk{}ङ्ग\appbrk{}से\appbrk{}कयोः~॥}\ A; \skt{मति\appbrk{}मा\appbrk{}न्युं\appbrk{}ज्या\appbrk{}द\appbrk{}भ्य\appbrk{}ङ्ग\appbrk{}से\appbrk{}कयोः~॥}\ H}}\skt{~॥१९॥}\edtext{}{\lemma{\skt{प्रकृ\appbrk{}ति\appbrk{}सा\appbrk{}त्म्य\appbrk{}र्तु\appbrk{}दे\appbrk{}श\appbrk{}दो\appbrk{}ष\appbrk{}वि\appbrk{}का\appbrk{}र\appbrk{}वित्~। तैलं}}\Afootnote{A K; \skt{प्रकृतिसात्म्यर्त्तु देश\appbrk{}दो\appbrk{}ष\appbrk{}वि\appbrk{}का\appbrk{}र\appbrk{}वित्~। तैलं}\ H}}\pend

\pstart\versepar \ledsidenote{\passref{4.24.35}}\skt{केवलं }\edtext{\skt{त्वामदोषेषु न}}{\lemma{\skt{त्वामदोषेषु न}}\Afootnote{K H; \skt{सामदोषेषु न}\ A}}\skt{ कथञ्चन योजयेत्~। }\pend
\pstart\versepar \skt{तरु\-ण\-ज्व\-र्य\-जीर्ण्णी च }\edtext{\skt{नाभ्यक्तव्यो}}{\lemma{\skt{नाभ्यक्तव्यो}}\Afootnote{K H; \skt{नाभ्यक्तव्यौ}\ A}}\skt{ }\edtext{\skt{कदाचनः}}{\lemma{\skt{कदाचनः}}\Afootnote{K; \skt{कथञ्चन~॥}\ A; \skt{कदाचन~॥}\ H}}\skt{~॥२०॥}\pend

\pstart\versepar \ledsidenote{\passref{4.24.36}}\skt{तथा }\edtext{\skt{विरक्तो}}{\lemma{\skt{विरक्तो}}\Afootnote{K; \skt{विरिक्तो}\ A H}}\skt{ वान्तश्च निरूढो यश्च मानवः~। }\pend
\pstart\versepar \skt{पूर्वयोः कृच्छ्रता }\edtext{\skt{व्याध्यो\-र\-सा\-ध्य\-त्व\-म\-थापि}}{\lemma{\skt{व्याध्योर्}}\Afootnote{K H; \skt{व्याधेर्}\ A}}\skt{ वा~।}\pend

\pstart \ledsidenote{\passref{4.24.37}}\skt{शेषाणान्तदहः }\edtext{\edtext{\skt{प्रोक्ता\-स्त्व\-ग्नि\-सा\-दा\-दयो}}{\lemma{\skt{प्रोक्तास्}}\Afootnote{K; \skt{प्रोक्ता}\ A; \skt{प्रोक्त}\ H}}}{\lemma{\skt{त्वग्निसादादयो}}\Afootnote{K; \skt{अग्निमान्द्यादयो}\ A; \skt{स्त्वग्निसादयो}\ H\textsuperscript{ac}; \skt{स्त्वग्निसादादयो}\ H\textsuperscript{pc}}}\skt{ }\edtext{\skt{गदाः}}{\lemma{\skt{गदाः}}\Afootnote{K; \skt{गदाः~। सन्त\appbrk{}र्प\appbrk{}ण\appbrk{}स\appbrk{}मु\appbrk{}त्थानां रोगाणां नैव कारयेत्~॥}\ A; \skt{गताः~।}\ H}}\skt{~॥२१॥~।}\pend

\pstart\versepar \ledsidenote{\passref{4.24.39}}\edtext{\skt{कान्तिः}}{\lemma{\skt{कान्तिः}}\Afootnote{K; \emph{om.}\ A; \skt{कान्ति}\ H}}\skt{ शरीरोपचयो }\edtext{\skt{गात्राणां}}{\lemma{\skt{गात्राणां}}\Afootnote{K H; \skt{कान्तिर्गात्राणां}\ A}}\skt{ सुविभक्तता~। }\pend
\pstart\versepar \skt{}\edtext{\skt{मृजाथ}}{\lemma{\skt{मृजाथ}}\Afootnote{K H; \skt{दीप्ता\appbrk{}ग्नि\appbrk{}त्व\appbrk{}म\appbrk{}ना\appbrk{}लस्यं स्थिरत्वं लाघवं मृजा~॥}\ A}}\skt{ }\edtext{\edtext{\edtext{\skt{स्थिर\-का\-य\-त्व\-म\-ग्नि\-दी\-प्ति\-र\-रो\-गता}}{\lemma{\skt{स्थिरकायत्वम्}}\Afootnote{K H; \emph{om.}\ A}}}{\lemma{\skt{अग्निदीप्तिर्}}\Afootnote{K H; \emph{om.}\ A}}}{\lemma{\skt{अरोगता}}\Afootnote{K H; \emph{om.}\ A}}\skt{~॥२२॥}\pend

\pstart\versepar \ledsidenote{\passref{4.24.40}}\edtext{\skt{श्रम\-क्ल\-म\-पि\-पा\-सानां}}{\lemma{\skt{श्रम\appbrk{}क्ल\appbrk{}म\appbrk{}पि\appbrk{}पा\appbrk{}सानां}}\Afootnote{H; \skt{श्रम\appbrk{}क्ल\appbrk{}म\appbrk{}पि\appbrk{}पा\appbrk{}सो\appbrk{}ष्ण\appbrk{}शी\appbrk{}ता\appbrk{}दीनां}\ A; \skt{श्रम\appbrk{}त्क्ल\appbrk{}म\appbrk{}पि\appbrk{}पा\appbrk{}सानां}\ K}}\skt{ }\edtext{\edtext{\skt{सहि\-ष्णु\-त्व\-म\-सी\-दनं}}{\lemma{\skt{सहिष्णुत्वम्}}\Afootnote{K H; \skt{सहिष्णुता~।}\ A}}}{\lemma{\skt{असीदनं}}\Afootnote{K H; \emph{om.}\ A}}\skt{~। }\pend
\pstart\versepar \skt{आरोग्यञ्चापि परमं }\edtext{\skt{व्याया\-मे\-नो\-प\-जा\-यते}}{\lemma{\skt{व्याया\appbrk{}मे\appbrk{}नो\appbrk{}प\appbrk{}जा\appbrk{}यते}}\Afootnote{K H; \skt{व्याया\appbrk{}मा\appbrk{}दु\appbrk{}प\appbrk{}जा\appbrk{}यते~॥}\ A}}\skt{~॥२३॥}\pend

\pstart\versepar \ledsidenote{\passref{4.24.41}}\edtext{\skt{न च}}{\lemma{\skt{न च}}\Afootnote{K H; \skt{न चा}\ A}}\skt{ }\edtext{\skt{व्याया\-म\-स\-दृ\-श\-म\-न्य\-स्थू\-ला\-प\-क\-र्षणं}}{\lemma{\skt{व्याया\appbrk{}म\appbrk{}स\appbrk{}दृ\appbrk{}शम्}}\Afootnote{K H; \skt{स्ति सदृशं}\ A}}\skt{~। }\pend
\pstart\versepar \skt{न च व्यायामिनं }\edtext{\edtext{\edtext{\skt{मर्त्य\-म\-र्द\-य\-न्त्य\-रयो}}{\lemma{\skt{मर्त्यम्}}\Afootnote{H; \skt{मर्त्य}\ A; \emph{om.}\ K}}}{\lemma{\skt{अर्दयन्त्य्}}\Afootnote{\skt{मर्दयन्त्य्}\ A; \skt{यन्}\ K; \skt{मर्द्दयन्}\ H}}}{\lemma{\skt{अरयो}}\Afootnote{A; \skt{त्यरयो}\ K H}}\skt{ }\edtext{\skt{बलात्}}{\lemma{\skt{बलात्}}\Afootnote{A; \skt{वलात्~।}\ K H}}\skt{~॥२४॥}\edtext{}{\lemma{\skt{अन्य\appbrk{}स्थू\appbrk{}ला\appbrk{}प\appbrk{}क\appbrk{}र्षणं~। न च}}\Afootnote{K H; \skt{तेन किञ्चि\appbrk{}त्स्थौ\appbrk{}ल्या\appbrk{}प\appbrk{}क\appbrk{}र्ष\appbrk{}णम्~। न च}\ A}}\pend

\pstart\versepar \ledsidenote{\passref{4.24.42}}\skt{न चैनं }\edtext{\skt{सहसाक्रम्य}}{\lemma{\skt{सहसाक्रम्य}}\Afootnote{K H; \skt{सहसाऽऽआक्रम्य}\ A}}\skt{ जरा समधिरोहति~। }\pend
\pstart\versepar \skt{स्थिरी }\edtext{\skt{भवति}}{\lemma{\skt{भवति}}\Afootnote{A K; \skt{भवन्ति}\ H}}\skt{ }\edtext{\skt{मांसश्च}}{\lemma{\skt{मांसश्च}}\Afootnote{K; \skt{मांसं च}\ A H}}\skt{ }\edtext{\edtext{\skt{व्याया\-मा\-भि\-र\-तस्य}}{\lemma{\skt{व्यायामाभिर्}}\Afootnote{A H; \skt{व्यायामाभि}\ K}}}{\lemma{\skt{अतस्य}}\Afootnote{A H; \skt{रतस्य}\ K}}\skt{ }\edtext{\skt{वै}}{\lemma{\skt{वै}}\Afootnote{K H; \skt{च~॥}\ A}}\skt{~॥२५॥}\pend

\pstart \ledsidenote{\passref{4.24.44}}\edtext{\edtext{\skt{व्यायामङ्कुर्वतो}}{\lemma{\skt{व्यायामङ्}}\Afootnote{K H; \skt{वयो\appbrk{}रू\appbrk{}प\appbrk{}गु\appbrk{}णै\appbrk{}र्ही\appbrk{}न\appbrk{}मपि कुर्या\appbrk{}त्सु\appbrk{}द\appbrk{}र्श\appbrk{}नम्~। व्यायामं}\ A}}}{\lemma{\skt{कुर्वतो}}\Afootnote{A K; \skt{कुर्व्वतो}\ H}}\skt{ नित्यं विरुद्धमपि भोजनं~।}\pend

\pstart \ledsidenote{\passref{4.24.45}}\edtext{\skt{विदग्धमपि}}{\lemma{\skt{अपि}}\Afootnote{K H; \skt{अवि}\ A}}\skt{ दग्धम्वा निर्दोषं }\edtext{\skt{परिपच्यते}}{\lemma{\skt{परिपच्यते}}\Afootnote{K H; \skt{परिपच्यते~। व्यायामो हि सदा पथ्यो बलिनां स्निग्धभोजिनाम्~॥}\ A}}\skt{~॥२६॥}\pend

\pstart \ledsidenote{\passref{4.24.50}}\edtext{\skt{रक्तपित्ती}}{\lemma{\skt{रक्तपित्ती}}\Afootnote{K H; \skt{कास\appbrk{}शो\appbrk{}ष\appbrk{}ज्व\appbrk{}र\appbrk{}श्वासा अति\appbrk{}व्या\appbrk{}या\appbrk{}म\appbrk{}सं\appbrk{}भवाः~। रक्तपित्ती}\ A}}\skt{ }\edtext{\skt{कृशः}}{\lemma{\skt{कृशः}}\Afootnote{A K; \skt{कृषः}\ H}}\skt{ शोषी }\edtext{\skt{कासी}}{\lemma{\skt{कासी}}\Afootnote{K H; \skt{श्वास}\ A}}\skt{ }\edtext{\skt{श्वासी}}{\lemma{\skt{श्वासी}}\Afootnote{K H; \skt{कास}\ A}}\skt{ क्षतातुरः~।}\pend

\pstart\versepar \ledsidenote{\passref{4.24.51}}\edtext{\skt{भुक्तवांस्त्रीषु च}}{\lemma{\skt{भुक्तवांस्त्रीषु च}}\Afootnote{\skt{भुक्तवान्स्त्रीषु च}\ A H; \skt{भुक्तवांस्त्रीषु च}\ K}}\skt{ }\edtext{\skt{क्षीणो}}{\lemma{\skt{क्षीणो}}\Afootnote{K H; \skt{क्षी}\ A}}\skt{ }\edtext{\skt{व्यायामं}}{\lemma{\skt{व्यायामं}}\Afootnote{K; \skt{व्यायायम्}\ H}}\skt{ परिवर्जयेत्~॥२७॥ }\pend
\pstart\versepar \skt{व्यायामो हि }\edtext{\skt{समपथ्यो}}{\lemma{\skt{समपथ्यो}}\Afootnote{K\textsuperscript{pc}; \skt{समम्पथ्यो}\ H\textsuperscript{pc}}}\skt{ }\edtext{\skt{बलिनां}}{\lemma{\skt{बलिनां}}\Afootnote{\skt{वलिनां}\ K\textsuperscript{pc} H\textsuperscript{pc}}}\skt{ स्निग्धभोजिनां~। }\pend
\pstart\versepar \skt{स च शीते वसन्ते च तेषां पथ्यतमः स्मृतः~॥२८॥ }\pend
\pstart\versepar \skt{सर्वेष्वृतुषु सर्वैर्हि }\edtext{\skt{मर्त्यै\-रा\-त्म\-हि\-तै\-षिभिः}}{\lemma{\skt{मर्त्यै\appbrk{}रा\appbrk{}त्म\appbrk{}हि\appbrk{}तै\appbrk{}षिभिः}}\Afootnote{K\textsuperscript{pc}; \skt{मर्त्यै\appbrk{}न्यै\appbrk{}रा\appbrk{}त्म\appbrk{}हि\appbrk{}ते\appbrk{}षिभिः}\ H\textsuperscript{pc}}}\skt{~। }\pend
\pstart\versepar \skt{}\edtext{\skt{शक्त्य\-वे\-नो\-प\-क\-र्त्तव्यो}}{\lemma{\skt{शक्त्य\appbrk{}वे\appbrk{}नो\appbrk{}प\appbrk{}क\appbrk{}र्त्तव्यो}}\Afootnote{K\textsuperscript{pc}; \skt{शक्त्य\appbrk{}र्द्ध\appbrk{}नो\appbrk{}प\appbrk{}क\appbrk{}र्त्तव्यो}\ H\textsuperscript{pc}}}\skt{ व्यायमो }\edtext{\skt{हन्त्यतोन्यथा}}{\lemma{\skt{हन्त्यतोन्यथा}}\Afootnote{\skt{हंन्यतो न्यथा~।}\ K\textsuperscript{pc} H\textsuperscript{pc}}}\skt{~॥२९॥ }\pend
\pstart\versepar \skt{हृदि स्थानस्थितो वायुर्यदा नासां प्रपद्यते~। }\pend
\pstart\versepar \skt{व्यायामं कुर्वतो }\edtext{\skt{जन्तो\-स्त\-च्छ\-क्त्य\-वस्य}}{\lemma{\skt{तच्छक्त्यवस्य}}\Afootnote{K\textsuperscript{pc}; \skt{तच्छक्त्यर्द्धस्य}\ H\textsuperscript{pc}}}\skt{ लक्षणं~॥३०॥ }\pend
\pstart\versepar \skt{}\edtext{\skt{वयो बलं}}{\lemma{\skt{वयो बलं}}\Afootnote{\skt{वयो वलं}\ K\textsuperscript{pc}; \skt{वयोवलं}\ H\textsuperscript{pc}}}\skt{ }\edtext{\skt{देहबलं}}{\lemma{\skt{देहबलं}}\Afootnote{\skt{देहवलं}\ K\textsuperscript{pc} H\textsuperscript{pc}}}\skt{ }\edtext{\skt{देश\-का\-ल\-म\-थापि च}}{\lemma{\skt{देश\appbrk{}का\appbrk{}ल\appbrk{}म\appbrk{}थापि च}}\Afootnote{H\textsuperscript{pc}; \skt{देश\appbrk{}का\appbrk{}ल\appbrk{}म\appbrk{}थापि च~।}\ K\textsuperscript{pc}}}\skt{~। }\pend
\pstart\versepar \skt{समीक्ष्य कुर्याद्व्यायामं युक्त्या शक्त्या च बुद्धिमान्~॥३१॥ }\pend
\pstart\versepar \skt{क्षय\-स्तृ\-ष्णा\-रु\-चि\-च्छ\-र्दि\-र\-क्त\-पि\-त्त\-श्र\-म\-क्लमाः~। }\pend
\pstart\versepar \skt{}\edtext{\skt{कास\-शो\-ष\-ज्व\-र\-श्वा\-सा\-श्चा\-ति\-व्या\-या\-म\-स\-म्भवाः}}{\lemma{\skt{चाति\appbrk{}व्या\appbrk{}या\appbrk{}म\appbrk{}स\appbrk{}म्भवाः}}\Afootnote{K\textsuperscript{pc}; \skt{च वर्जयेत्~। उ}\ A; \skt{चाति\appbrk{}व्या\appbrk{}य\appbrk{}म\appbrk{}स\appbrk{}म्भवाः~।}\ H\textsuperscript{pc}}}\skt{~॥३२॥ }\pend
\pstart\versepar \skt{}\edtext{\skt{उद्वर्त्तनं}}{\lemma{\skt{उद्वर्त्तनं}}\Afootnote{K\textsuperscript{pc} H\textsuperscript{pc}; \skt{द्वर्तनं}\ A; \skt{तुद्वर्त्तनं}\ K\textsuperscript{ac} H\textsuperscript{ac}}}\skt{ वातहरं }\edtext{\skt{कफमेदो}}{\lemma{\skt{कफमेदो}}\Afootnote{K H; \skt{कफमेदो}\ A}}\skt{ विलापनं~।}\edtext{}{\lemma{\skt{णो… कासशोष}}\Afootnote{\emph{om.}\ A}}\edtext{}{\lemma{\skt{परिवर्जयेत्~॥२७॥ व्यायामो हि}}\Afootnote{K\textsuperscript{pc} H\textsuperscript{pc}; \skt{परिवर्जये}\ K\textsuperscript{ac} H\textsuperscript{ac}}}\edtext{}{\lemma{\skt{त्… उ}}\Afootnote{\emph{om.}\ K\textsuperscript{ac}}}\edtext{}{\lemma{\skt{त्… उ}}\Afootnote{\emph{om.}\ H\textsuperscript{ac}}}\edtext{}{\lemma{\skt{बुद्धिमान्~॥३१॥ क्षय\appbrk{}स्तृ\appbrk{}ष्णा\appbrk{}रु\appbrk{}चि\appbrk{}च्छ\appbrk{}र्दि\appbrk{}र\appbrk{}क्त\appbrk{}पि\appbrk{}त्त\appbrk{}श्र\appbrk{}म\appbrk{}क्लमाः}}\Afootnote{\skt{वुद्धि\appbrk{}मा\appbrk{}न्क्ष\appbrk{}य\appbrk{}स्तृष्णा रुचिच्छर्दि रक्त\appbrk{}पि\appbrk{}त्त\appbrk{}श्र\appbrk{}म\appbrk{}क्लमाः}\ K\textsuperscript{pc} H\textsuperscript{pc}}}\pend

\pstart\versepar \ledsidenote{\passref{4.24.52}}\edtext{\edtext{\skt{स्थिरी\-क\-र\-ण\-म\-ङ्गानां}}{\lemma{\skt{स्थिरीकरणम्}}\Afootnote{A K; \skt{स्थिरीकरण}\ H}}}{\lemma{\skt{अङ्गानां}}\Afootnote{A K; \skt{मङ्गानां}\ H}}\skt{ }\edtext{\skt{त्वक्प्रसादकरं}}{\lemma{\skt{त्वक्प्रसादकरं}}\Afootnote{K; \skt{त्वक्प्रस्थादकरं}\ A; \skt{त्वक्प्रसाधकरं}\ H}}\skt{ }\edtext{\skt{तथा}}{\lemma{\skt{तथा}}\Afootnote{K H; \skt{परम्~।}\ A}}\skt{~॥३३॥ }\pend
\pstart\versepar \skt{सिरा\-मु\-ख\-वि\-वि\-क्तत्वं }\edtext{\skt{त्वक्थस्याग्रेश्च}}{\lemma{\skt{त्वक्थस्याग्रेश्च}}\Afootnote{K; \skt{त्वक्स्थस्याग्नेश्च}\ A; \skt{तव्क्थस्याग्नेश्च}\ H}}\skt{ }\edtext{\skt{दीपनं}}{\lemma{\skt{दीपनं}}\Afootnote{K H; \skt{तेजनम्~॥}\ A}}\skt{~।}\pend

\pstart \ledsidenote{\passref{4.24.53}}\edtext{\skt{उद्घ\-र्ष\-णो\-द\-नाभ्यां}}{\lemma{\skt{उद्घ\appbrk{}र्ष\appbrk{}णो\appbrk{}द\appbrk{}नाभ्यां}}\Afootnote{\skt{उद्घ\appbrk{}र्ष\appbrk{}णो\appbrk{}त्सा\appbrk{}द\appbrk{}नाभ्यां}\ A K H}}\skt{ }\edtext{\skt{भवे\-या\-ता\-म\-सं\-शयं}}{\lemma{\skt{भवे\appbrk{}या\appbrk{}ता\appbrk{}म\appbrk{}सं\appbrk{}शयं}}\Afootnote{K\textsuperscript{pc} H; \skt{जाये\appbrk{}या\appbrk{}ता\appbrk{}म\appbrk{}सं\appbrk{}श\appbrk{}यम्~। उत्सा\appbrk{}द\appbrk{}ना\appbrk{}द्भ\appbrk{}वे\appbrk{}त्स्त्रीणां विशे\appbrk{}षा\appbrk{}त्का\appbrk{}न्ति\appbrk{}म\appbrk{}द्वपुः~॥}\ A; \skt{भवे\appbrk{}या\appbrk{}ता\appbrk{}य\appbrk{}म\appbrk{}सं\appbrk{}शयं~।}\ K\textsuperscript{ac}}}\skt{~।}\pend

\pstart\versepar \ledsidenote{\passref{4.24.58}}\skt{तन्द्रा\-पा\-प्मो\-प\-श\-मनं तुष्टिदं }\edtext{\skt{पुष्टिवर्धनं}}{\lemma{\skt{पुष्टिवर्धनं}}\Afootnote{K H; \skt{पुंस्त्ववर्धनम्।}\ A}}\skt{~। }\pend
\pstart\versepar \skt{रक्त\-प्र\-सा\-द\-न\-ञ्चापि स्नानमग्नेश्च दीपनं~।}\pend

\pstart\versepar \ledsidenote{\passref{4.24.62}}\edtext{\skt{तच्चा\-ति\-सा\-र\-ज्व\-रि\-त\-क\-र्ण्ण\-शू\-ला\-र्दि\-तो\-दिषु}}{\lemma{\skt{तच्चा\appbrk{}ति\appbrk{}सा\appbrk{}र\appbrk{}ज्व\appbrk{}रि\appbrk{}त\appbrk{}क\appbrk{}र्ण्ण\appbrk{}शू\appbrk{}ला\appbrk{}र्दि\appbrk{}तो\appbrk{}दिषु}}\Afootnote{\skt{तच्चा\appbrk{}ति\appbrk{}सा\appbrk{}र\appbrk{}ज्व\appbrk{}रि\appbrk{}त\appbrk{}क\appbrk{}र्ण\appbrk{}शू\appbrk{}ला\appbrk{}नि\appbrk{}ला\appbrk{}र्तिषु~।}\ A; \skt{तच्चा\appbrk{}ति\appbrk{}सा\appbrk{}र\appbrk{}ज्व\appbrk{}रि\appbrk{}त\appbrk{}क\appbrk{}र्ण्ण\appbrk{}शू\appbrk{}ला\appbrk{}र्दि\appbrk{}तो\appbrk{}दिषु}\ K H}}\skt{~। }\pend
\pstart\versepar \skt{}\edtext{\skt{आध्मा\-ता\-रो\-च\-का\-जी\-र्ण्ण\-भु\-क्त\-वत्सु च}}{\lemma{\skt{आध्मा\appbrk{}ता\appbrk{}रो\appbrk{}च\appbrk{}का\appbrk{}जी\appbrk{}र्ण्ण\appbrk{}भु\appbrk{}क्त\appbrk{}वत्सु च}}\Afootnote{K H; \skt{आध्मा\appbrk{}ना\appbrk{}रो\appbrk{}च\appbrk{}का\appbrk{}जी\appbrk{}र्ण\appbrk{}भु\appbrk{}क्त\appbrk{}वत्सु च}\ A}}\skt{ गर्हितं~।}\pend

\pstart\versepar \ledsidenote{\passref{4.24.63}}\edtext{\skt{तन्द्रा\-पा\-प्मो\-प\-श\-मनं}}{\lemma{\skt{तन्द्रा\appbrk{}पा\appbrk{}प्मो\appbrk{}प\appbrk{}श\appbrk{}मनं}}\Afootnote{K H; \skt{सौभाग्यदं वर्णकरं}\ A}}\skt{ }\edtext{\skt{प्रीत्योजो}}{\lemma{\skt{प्रीत्योजो}}\Afootnote{K H; \skt{प्रीत्योजो}\ A}}\skt{ }\edtext{\skt{बलवर्धनं}}{\lemma{\skt{बलवर्धनं}}\Afootnote{A; \skt{वलवर्धनं}\ K H}}\skt{~। }\pend
\pstart\versepar \skt{}\edtext{\skt{स्वेद\-दौ\-र्ग\-न्ध्य\-वै\-व\-र्ण्य\-श्र\-म\-घ्न\-म\-नु\-ले\-पनं}}{\lemma{\skt{स्वेद\appbrk{}दौ\appbrk{}र्ग\appbrk{}न्ध्य\appbrk{}वै\appbrk{}व\appbrk{}र्ण्य\appbrk{}श्र\appbrk{}म\appbrk{}घ्न\appbrk{}म\appbrk{}नु\appbrk{}ले\appbrk{}पनं}}\Afootnote{\skt{स्वेद\appbrk{}दौ\appbrk{}र्ग\appbrk{}न्ध्य\appbrk{}वै\appbrk{}व\appbrk{}र्ण्य\appbrk{}श्र\appbrk{}म\appbrk{}घ्न\appbrk{}म\appbrk{}नु\appbrk{}ले\appbrk{}प\appbrk{}नम्}\ A K H}}\skt{~।}\pend

\pstart \ledsidenote{\passref{4.24.64}}\edtext{\skt{रक्षोघ्नमथ}}{\lemma{\skt{रक्षोघ्नमथ}}\Afootnote{K; \skt{स्नानं येषां निषिद्धं तु तेषा\appbrk{}म\appbrk{}प्य\appbrk{}नु\appbrk{}ले\appbrk{}प\appbrk{}नम्~। रक्षोघ्नमथ}\ A; \skt{रक्षोघ्नमथ}\ H}}\skt{ चौजस्यं }\edtext{\skt{सौभा\-ग्यं\-क\-र\-मु\-त्तमं}}{\lemma{\skt{सौभाग्यंकरम्}}\Afootnote{K; \skt{सौभाग्यकरम्}\ A H}}\skt{~।}\pend

\pstart\versepar \ledsidenote{\passref{4.24.65}}\edtext{\skt{सुम\-नो\-म्व\-र\-र\-त्ना\-ना\-न्धा\-रणं}}{\lemma{\skt{सुम\appbrk{}नो\appbrk{}म्व\appbrk{}र\appbrk{}र\appbrk{}त्ना\appbrk{}नान्}}\Afootnote{K H; \skt{सुम\appbrk{}नो\appbrk{}म्ब\appbrk{}र\appbrk{}र\appbrk{}त्नानां}\ A}}\skt{ प्रीतिवर्धनं~। }\pend
\pstart\versepar \skt{}\edtext{\skt{मुखालेपं}}{\lemma{\skt{मुखालेपं}}\Afootnote{K; \skt{मुखालेपा}\ A H}}\skt{ }\edtext{\skt{दृढंचक्षुः}}{\lemma{\skt{दृढंचक्षुः}}\Afootnote{K; \skt{द्दृढं चक्षुः}\ A; \skt{द्दृढञ्चक्षुः}\ H}}\skt{ }\edtext{\skt{पीन\-ग\-ण्ड\-न्त\-था\-ननं}}{\lemma{\skt{पीन\appbrk{}ग\appbrk{}ण्ड\appbrk{}न्त\appbrk{}था\appbrk{}ननं}}\Afootnote{K H; \skt{पीनगण्डं तथाऽऽआननम्~॥}\ A}}\skt{~।}\pend

\pstart\versepar \ledsidenote{\passref{4.24.66}}\edtext{\skt{कान्त\-म\-व्य\-ङ्ग\-पि\-डकं}}{\lemma{\skt{कान्त\appbrk{}म\appbrk{}व्य\appbrk{}ङ्ग\appbrk{}पि\appbrk{}डकं}}\Afootnote{\skt{अव्यङ्गपिडकं}\ A; \skt{कान्त\appbrk{}म\appbrk{}व्य\appbrk{}ङ्ग\appbrk{}पि\appbrk{}डकं}\ K H}}\skt{ }\edtext{\skt{भव\-त्य\-म्वु\-ज\-स\-न्निभं}}{\lemma{\skt{भव\appbrk{}त्य\appbrk{}म्वु\appbrk{}ज\appbrk{}स\appbrk{}न्निभं}}\Afootnote{\skt{कान्तं भव\appbrk{}त्य\appbrk{}म्बु\appbrk{}ज\appbrk{}स\appbrk{}न्नि\appbrk{}भम्~।}\ A; \skt{भव\appbrk{}त्य\appbrk{}म्वु\appbrk{}ज\appbrk{}स\appbrk{}न्निभं~।}\ K H}}\skt{~। }\pend
\pstart\versepar \skt{}\edtext{\skt{पक्ष्मलं}}{\lemma{\skt{पक्ष्मलं}}\Afootnote{K H; \skt{पक्स्मलं}\ A}}\skt{ विशदं }\edtext{\skt{कान्त\-म\-म\-लो\-ज्व\-ल\-म\-ण्डलं}}{\lemma{\skt{अम\appbrk{}लो\appbrk{}ज्व\appbrk{}ल\appbrk{}म\appbrk{}ण्डलं}}\Afootnote{K H; \skt{अम\appbrk{}लो\appbrk{}ज्ज्व\appbrk{}ल\appbrk{}म\appbrk{}ण्ड\appbrk{}लम्~॥}\ A}}\skt{~।}\pend

\pstart\versepar \ledsidenote{\passref{4.24.67}}\skt{नेत्र\-म\-ञ्ज\-न\-सं\-यो\-गा\-द्भ\-वे\-च्चा\-म\-ल\-ता\-रकं~॥ }\pend
\pstart\versepar \skt{यशस्यं स्वर्ग्यमायुष्यं }\edtext{\skt{धन\-धा\-न्या\-भि\-व\-र्धनं}}{\lemma{\skt{धन\appbrk{}धा\appbrk{}न्या\appbrk{}भि\appbrk{}व\appbrk{}र्धनं}}\Afootnote{K H; \skt{धन\appbrk{}धा\appbrk{}न्य\appbrk{}वि\appbrk{}व\appbrk{}र्ध\appbrk{}नम्~॥}\ A}}\skt{~।}\pend

\pstart\versepar \ledsidenote{\passref{4.24.68}}\edtext{\skt{गुरु\-दे\-व\-त\-वि\-प्रा\-णा\-म्पू\-जनं}}{\lemma{\skt{गुरु\appbrk{}दे\appbrk{}व\appbrk{}त\appbrk{}वि\appbrk{}प्रा\appbrk{}णाम्}}\Afootnote{K; \skt{देव\appbrk{}ता\appbrk{}ति\appbrk{}थि\appbrk{}वि\appbrk{}प्राणां}\ A; \skt{गुरु\appbrk{}दै\appbrk{}व\appbrk{}त\appbrk{}वि\appbrk{}प्राणां}\ H}}\skt{ गोत्रवर्धनं~। }\pend
\pstart\versepar \skt{आहारः प्रीणनः सद्यो }\edtext{\skt{बल\-कृ\-द्दे\-ह\-धा\-रणंः}}{\lemma{\skt{बल\appbrk{}कृ\appbrk{}द्दे\appbrk{}ह\appbrk{}धा\appbrk{}रणंः}}\Afootnote{\skt{बल\appbrk{}कृ\appbrk{}द्दे\appbrk{}ह\appbrk{}धा\appbrk{}रकः~॥}\ A; \skt{वल\appbrk{}कृ\appbrk{}द्दे\appbrk{}ह\appbrk{}धा\appbrk{}रणंः~।}\ K; \skt{वल\appbrk{}कृ\appbrk{}द्दे\appbrk{}ह\appbrk{}धा\appbrk{}रणं~॥}\ H}}\skt{~।}\pend

\pstart\versepar \ledsidenote{\passref{4.24.69}}\edtext{\skt{स्मृत्यायुः}}{\lemma{\skt{स्मृत्यायुः}}\Afootnote{K H; \skt{आयु\appbrk{}स्ते\appbrk{}जः\appbrk{}स\appbrk{}मु\appbrk{}त्सा\appbrk{}हस्मृ}\ A}}\skt{ }\edtext{\skt{पक्कितेजोजः}}{\lemma{\skt{पक्कितेजोजः}}\Afootnote{K; \emph{om.}\ A; \skt{पक्तितेजोजः}\ H}}\skt{ }\edtext{\skt{सत्वो\-त्सा\-हा\-भि\-व\-र्धनः}}{\lemma{\skt{सत्वो\appbrk{}त्सा\appbrk{}हा\appbrk{}भि\appbrk{}व\appbrk{}र्धनः}}\Afootnote{K; \skt{त्योजो\appbrk{}ग्नि\appbrk{}वि\appbrk{}व\appbrk{}र्धनः~।}\ A; \skt{सत्वो\appbrk{}त्सा\appbrk{}हा\appbrk{}भि\appbrk{}व\appbrk{}र्द्धनाः~।}\ H}}\skt{~। }\pend
\pstart\versepar \skt{पादप्रक्षालनं }\edtext{\skt{पाद\-रो\-ग\-श्र\-म\-म\-ला\-पहं}}{\lemma{\skt{पाद\appbrk{}रो\appbrk{}ग\appbrk{}श्र\appbrk{}म\appbrk{}म\appbrk{}ला\appbrk{}पहं}}\Afootnote{K H; \skt{पाद\appbrk{}म\appbrk{}ल\appbrk{}रो\appbrk{}ग\appbrk{}श्र\appbrk{}मा\appbrk{}प\appbrk{}हम्~॥}\ A}}\skt{~।}\pend

\pstart\versepar \ledsidenote{\passref{4.24.70}}\edtext{\skt{चक्षु}}{\lemma{\skt{चक्षु}}\Afootnote{K; \skt{चक्षुः}\ A H}}\skt{ }\edtext{\skt{प्रसादजननं}}{\lemma{\skt{प्रसादजननं}}\Afootnote{K H; \skt{प्रसादनं वृष्यं}\ A}}\skt{ रक्षोघ्नं प्रीतिवर्धनं~। }\pend
\pstart\versepar \skt{निद्राकरो }\edtext{\skt{देहसुखं}}{\lemma{\skt{देहसुखं}}\Afootnote{K H; \skt{देहसुखश्}\ A}}\skt{ चक्षुष्यः }\edtext{\skt{पादरोगहा}}{\lemma{\skt{पादरोगहा}}\Afootnote{K H; \skt{श्रमसुप्तिनुत्~॥}\ A}}\skt{~।}\pend

\pstart\versepar \ledsidenote{\passref{4.24.71}}\edtext{\skt{पाद\-त्व\-ङ्मृ\-दु\-कर्त्ता च}}{\lemma{\skt{मृदुकर्त्ता च}}\Afootnote{K H; \skt{मृदुकारी च}\ A}}\skt{ पादाभ्यङ्गः }\edtext{\skt{प्रशस्यते}}{\lemma{\skt{प्रशस्यते}}\Afootnote{K H; \skt{सदा हितः~।}\ A}}\skt{~। }\pend
\pstart\versepar \skt{पादरोगहरं }\edtext{\skt{वृष्यमायुष्य}}{\lemma{\skt{आयुष्य}}\Afootnote{\emph{om.}\ A; \skt{आयुष्यञ्}\ K H}}\skt{ }\edtext{\skt{तं}}{\lemma{\skt{तं}}\Afootnote{\skt{रक्षोघ्नं प्रीतिवर्धनम्~॥}\ A; \skt{चक्षुषे हि तं}\ K H}}\skt{~।}\pend

\pstart\versepar \ledsidenote{\passref{4.24.72}}\edtext{\skt{सुख\-प्र\-चा\-र\-मौ\-जस्यं}}{\lemma{\skt{औजस्यं}}\Afootnote{K H; \skt{ओजस्यं}\ A}}\skt{ सदा\-पा\-द\-त्र\-धा\-र\-णम्~। }\pend
\pstart\versepar \skt{}\edtext{\edtext{\skt{अना\-रो\-ग्य\-म\-ना\-यु\-ष्य\-मि\-न्द्रि\-य\-घ्न\-म\-दृ\-ष्टि\-कृत्}}{\lemma{\skt{इन्द्रियघ्नम्}}\Afootnote{K H; \skt{चक्षुषोरु}\ A}}}{\lemma{\skt{अदृष्टिकृत्}}\Afootnote{K H; \skt{पघातकृत्~॥}\ A}}\skt{~।}\pend

\pstart\versepar \ledsidenote{\passref{4.24.73}}\edtext{\skt{पादा\-भ्या\-म\-नु\-पा\-नभ्यां}}{\lemma{\skt{अनुपानभ्यां}}\Afootnote{K; \skt{अनुपानद्भ्यां}\ A H}}\skt{ }\edtext{\skt{नृणाञ्चङ्क्रमणं}}{\lemma{\skt{नृणाञ्}}\Afootnote{K H; \skt{सदा}\ A}}\skt{ }\edtext{\skt{सदा}}{\lemma{\skt{सदा}}\Afootnote{K H; \skt{नृणाम्~।}\ A}}\skt{~। }\pend
\pstart\versepar \skt{पाप्मोपशमनं केश\-न\-ख\-रो\-मा\-प\-मा\-र्जनं~।}\pend

\pstart \ledsidenote{\passref{4.24.74}}\edtext{\edtext{\skt{हर्ष\-ला\-घ\-व\-सौ\-भा\-ग्य\-मु\-त्सा\-ह\-क\-र\-व\-र्धनं}}{\lemma{\skt{हर्ष\appbrk{}ला\appbrk{}घ\appbrk{}व\appbrk{}सौ\appbrk{}भा\appbrk{}ग्यम्}}\Afootnote{K; \skt{हर्ष\appbrk{}ला\appbrk{}घ\appbrk{}व\appbrk{}सौ\appbrk{}भाग्य}\ A H}}}{\lemma{\skt{उत्सा\appbrk{}ह\appbrk{}क\appbrk{}र\appbrk{}व\appbrk{}र्धनं}}\Afootnote{K; \skt{कर\appbrk{}मु\appbrk{}त्सा\appbrk{}ह\appbrk{}व\appbrk{}र्ध\appbrk{}नम्~। बाणवारं मृजा\appbrk{}व\appbrk{}र्ण\appbrk{}ते\appbrk{}जो\appbrk{}ब\appbrk{}ल\appbrk{}वि\appbrk{}व\appbrk{}र्ध\appbrk{}नम्~॥}\ A; \skt{कर\appbrk{}मु\appbrk{}त्सा\appbrk{}ह\appbrk{}व\appbrk{}र्द्धनं~।}\ H}}\skt{~।}\pend

\pstart\versepar \ledsidenote{\passref{4.24.75}}\skt{पवित्रं }\edtext{\skt{केश्यमुष्णीशं}}{\lemma{\skt{केश्यमुष्णीशं}}\Afootnote{K; \skt{केश्यमुष्णीषं}\ A; \skt{केश्यमुष्णीशं}\ H}}\skt{ }\edtext{\skt{वाता\-त\-प\-र\-जो\-नुदं}}{\lemma{\skt{वाता\appbrk{}त\appbrk{}प\appbrk{}र\appbrk{}जो\appbrk{}नुदं}}\Afootnote{K H; \skt{वाता\appbrk{}त\appbrk{}प\appbrk{}र\appbrk{}जो\appbrk{}प\appbrk{}हम्~।}\ A}}\skt{~। }\pend
\pstart\versepar \skt{}\edtext{\skt{छत्रं वातातपरुजो}}{\lemma{\skt{छत्रं वातातपरुजो}}\Afootnote{K\textsuperscript{pc} H\textsuperscript{pc}; \skt{वर्षा\appbrk{}नि\appbrk{}ल\appbrk{}र\appbrk{}जो\appbrk{}घर्म}\ A; \emph{om.}\ K\textsuperscript{ac} H\textsuperscript{ac}}}\skt{ }\edtext{\skt{मारुतापहरं}}{\lemma{\skt{मारुतापहरं}}\Afootnote{K\textsuperscript{pc} H\textsuperscript{pc}; \skt{हिमादीनां निवा}\ A; \emph{om.}\ K\textsuperscript{ac} H\textsuperscript{ac}}}\skt{ }\edtext{\skt{शुभं}}{\lemma{\skt{शुभं}}\Afootnote{K\textsuperscript{pc} H\textsuperscript{pc}; \skt{रणम्~॥}\ A; \emph{om.}\ K\textsuperscript{ac} H\textsuperscript{ac}}}\skt{~।}\pend

\pstart \ledsidenote{\passref{4.24.77}}\edtext{\skt{सत्वो\-त्सा\-ह\-ब\-ल\-स्थै\-र्य\-धै\-र्य\-वी\-र्या\-भि\-व\-र्धनं}}{\lemma{\skt{सत्वो\appbrk{}त्सा\appbrk{}ह\appbrk{}ब\appbrk{}ल\appbrk{}स्थै\appbrk{}र्य\appbrk{}धै\appbrk{}र्य\appbrk{}वी\appbrk{}र्या\appbrk{}भि\appbrk{}व\appbrk{}र्धनं}}\Afootnote{\skt{श्रम\appbrk{}स्ख\appbrk{}ल\appbrk{}न\appbrk{}दो\appbrk{}षघ्नं स्थविरे च प्रशस्यते~। सत्त्वो\appbrk{}त्सा\appbrk{}ह\appbrk{}ब\appbrk{}ल\appbrk{}स्थैर्य धैर्य\appbrk{}वी\appbrk{}र्य\appbrk{}वि\appbrk{}व\appbrk{}र्ध\appbrk{}नम्~॥}\ A; \skt{सत्वो\appbrk{}त्सा\appbrk{}ह\appbrk{}व\appbrk{}ल\appbrk{}स्थै\appbrk{}र्य\appbrk{}धै\appbrk{}र्य\appbrk{}वी\appbrk{}र्या\appbrk{}भि\appbrk{}व\appbrk{}र्धनं}\ K H}}\skt{~।}\pend

\pstart\versepar \ledsidenote{\passref{4.24.78}}\skt{अव\-ष्ट\-म्भ\-क\-र\-ञ्चापि }\edtext{\skt{रक्षोघ्नं}}{\lemma{\skt{रक्षोघ्नं}}\Afootnote{K\textsuperscript{pc} H; \skt{भयघ्नं}\ A; \skt{रक्षोष्णंघ्नं}\ K\textsuperscript{ac}}}\skt{ दण्डधारणं~। }\pend
\pstart\versepar \skt{}\edtext{\skt{आस्या\-व\-र्ण्ण\-श्ले\-ष्म\-मेदः}}{\lemma{\skt{आस्या\appbrk{}व\appbrk{}र्ण्ण\appbrk{}श्ले\appbrk{}ष्म\appbrk{}मेदः}}\Afootnote{K H; \skt{आस्या वर्णकफस्थौल्य}\ A}}\skt{ सौकुमार्यकरी सुखा~।}\pend

\pstart \ledsidenote{\passref{4.24.79}}\edtext{\skt{अध्वा\-व\-र्ण्ण\-श्ले\-ष्म\-मेदः}}{\lemma{\skt{अध्वा\appbrk{}व\appbrk{}र्ण्ण\appbrk{}श्ले\appbrk{}ष्म\appbrk{}मेदः}}\Afootnote{K; \skt{अध्वा वर्णकफस्थौल्य}\ A; \skt{अध्वा\appbrk{}व\appbrk{}र्ण्ण\appbrk{}श्ले\appbrk{}ष\appbrk{}मेदः}\ H}}\skt{ }\edtext{\skt{सौकु\-मा\-र्य\-वि\-ना\-शनः}}{\lemma{\skt{सौकु\appbrk{}मा\appbrk{}र्य\appbrk{}वि\appbrk{}ना\appbrk{}शनः}}\Afootnote{K H; \skt{सौकु\appbrk{}मा\appbrk{}र्य\appbrk{}वि\appbrk{}ना\appbrk{}शनः~। अत्यध्वा विप\appbrk{}री\appbrk{}तो+\appbrk{}आ\appbrk{}स्मा\appbrk{}ज्ज\appbrk{}रा\appbrk{}दौ\appbrk{}र्ब\appbrk{}ल्य\appbrk{}कृच्च सः~॥}\ A}}\skt{~।}\pend

\pstart\versepar \ledsidenote{\passref{4.24.80}}\skt{यत्तु }\edtext{\skt{चंक्र\-म\-ण\-न्ना\-ति\-दे\-ह\-पी\-डा\-करं}}{\lemma{\skt{नाति\appbrk{}दे\appbrk{}ह\appbrk{}पी\appbrk{}डा\appbrk{}करं}}\Afootnote{A K; \skt{नादि\appbrk{}दे\appbrk{}ह\appbrk{}पी\appbrk{}ड\appbrk{}क\appbrk{}रम्}\ H}}\skt{ भवेत्~। }\pend
\pstart\versepar \skt{}\edtext{\skt{तदा\-यु\-र्ब\-ल\-मे\-धा\-ग्नि\-प्र\-द\-मि\-न्द्रि\-य\-बो\-धणं}}{\lemma{\skt{इन्द्रियबोधणं}}\Afootnote{\skt{मिन्द्रि\appbrk{}य\appbrk{}बो\appbrk{}ध\appbrk{}नम्~॥}\ A; \skt{इन्द्रियवोधणं~।}\ K; \skt{मिन्द्रियवोधनं~।}\ H}}\skt{~।}\edtext{}{\lemma{\skt{भवेत्~। तदा\appbrk{}यु\appbrk{}र्ब\appbrk{}ल\appbrk{}मे\appbrk{}धा\appbrk{}ग्नि\appbrk{}प्र\appbrk{}दम्}}\Afootnote{\skt{भवेत्~। तदा\appbrk{}यु\appbrk{}र्ब\appbrk{}ल\appbrk{}मे\appbrk{}धा\appbrk{}ग्नि\appbrk{}प्रद}\ A; \skt{भवेत्~। तदा\appbrk{}यु\appbrk{}र्व\appbrk{}ल\appbrk{}मे\appbrk{}धा\appbrk{}ग्नि\appbrk{}प्र\appbrk{}दम्}\ K; \skt{भवेत्~॥ तदा\appbrk{}यु\appbrk{}र्व्व\appbrk{}ल\appbrk{}मे\appbrk{}धा\appbrk{}ग्नि\appbrk{}प्रद}\ H}}\pend

\pstart\versepar \ledsidenote{\passref{4.24.81}}\skt{श्रमानिलहरं वृष्यं पुष्टि\-नि\-द्रा\-धृ\-ति\-प्रदं~। }\pend
\pstart\versepar \skt{सुखं }\edtext{\skt{शय्या\-स\-न\-न्त\-स्मा\-द्वि\-प\-री\-त\-म\-तो\-न्यथा}}{\lemma{\skt{तस्मा\appbrk{}द्वि\appbrk{}प\appbrk{}री\appbrk{}त\appbrk{}म\appbrk{}तो\appbrk{}न्यथा}}\Afootnote{H; \skt{दुःखं विपरीतगुणं मतम्~॥}\ A; \skt{तस्मा\appbrk{}द्वि\appbrk{}प\appbrk{}री\appbrk{}त\appbrk{}मतो न्यथा~।}\ K}}\skt{~।}\pend

\pstart\versepar \ledsidenote{\passref{4.24.82}}\edtext{\skt{वाल\-व्य\-ज\-न\-मौ\-जस्यं}}{\lemma{\skt{वाल\appbrk{}व्य\appbrk{}ज\appbrk{}न\appbrk{}मौ\appbrk{}जस्यं}}\Afootnote{K H; \skt{बाल\appbrk{}व्य\appbrk{}ज\appbrk{}न\appbrk{}मौ\appbrk{}जस्यं}\ A}}\skt{ मक्षि\-का\-दी\-न\-पो\-हति~। }\pend
\pstart\versepar \skt{}\edtext{\skt{स्वेद\-दा\-ह\-तृ\-षा\-मू\-र्च्छा\-श्र\-मघ्नो}}{\lemma{\skt{स्वेद\appbrk{}दा\appbrk{}ह\appbrk{}तृ\appbrk{}षा\appbrk{}मू\appbrk{}र्च्छा\appbrk{}श्र\appbrk{}मघ्नो}}\Afootnote{K H; \skt{शोष\appbrk{}दा\appbrk{}ह\appbrk{}श्र\appbrk{}म\appbrk{}स्वे\appbrk{}द\appbrk{}मू\appbrk{}र्च्छाघ्नो}\ A}}\skt{ व्यजनानिलः~।}\pend

\pstart\versepar \ledsidenote{\passref{4.24.83}}\skt{प्रीतिनिद्राकरं वृष्यं कफ\-वा\-त\-श्र\-मा\-पहं~। }\pend
\pstart\versepar \skt{सम्वाहनं }\edtext{\skt{मांस\-रं\-स\-त्व\-क्प्र\-सा\-द\-करं}}{\lemma{\skt{मांस\appbrk{}रं\appbrk{}स\appbrk{}त्व\appbrk{}क्प्र\appbrk{}सा\appbrk{}द\appbrk{}करं}}\Afootnote{K; \skt{मांस\appbrk{}र\appbrk{}क्त\appbrk{}त्व\appbrk{}क्प्र\appbrk{}सा\appbrk{}द\appbrk{}करं}\ A; \skt{मांसरसं त्वक्प्र\appbrk{}सा\appbrk{}द\appbrk{}क\appbrk{}रन्}\ H}}\skt{ }\edtext{\skt{तथा}}{\lemma{\skt{तथा}}\Afootnote{K H; \skt{सुखम्~॥}\ A}}\skt{~।}\pend

\pstart\versepar \ledsidenote{\passref{4.24.84}}\skt{प्रवातं रौक्ष्य\-वै\-व\-र्ण्ण्स्त\-म्भ\-कृ\-ह\-प\-क्ति\-नुत्~। }\pend
\pstart\versepar \skt{}\edtext{\skt{स्वेद\-मू\-र्च्छा\-पि\-पा\-स\-घ्न\-म\-प्र\-वा\-त\-म\-तो\-न्यथा}}{\lemma{\skt{अप्र\appbrk{}वा\appbrk{}त\appbrk{}म\appbrk{}तो\appbrk{}न्यथा}}\Afootnote{K; \skt{मप्र\appbrk{}वा\appbrk{}त\appbrk{}म\appbrk{}तोऽ\appbrk{}न्यथा~॥}\ A; \skt{अप्र\appbrk{}वा\appbrk{}त\appbrk{}म\appbrk{}तोऽ\appbrk{}न्यथा~।}\ H}}\skt{~।}\edtext{}{\lemma{\skt{रौक्ष्य\appbrk{}वै\appbrk{}व\appbrk{}र्ण्ण्स्त\appbrk{}म्भ\appbrk{}कृ\appbrk{}ह\appbrk{}प\appbrk{}क्ति\appbrk{}नुत्~। स्वेद\appbrk{}मू\appbrk{}र्च्छा\appbrk{}पि\appbrk{}पा\appbrk{}स\appbrk{}घ्नम्}}\Afootnote{K; \skt{रौक्ष्य\appbrk{}वै\appbrk{}व\appbrk{}र्ण्य\appbrk{}स्त\appbrk{}म्भ\appbrk{}कृ\appbrk{}द्दा\appbrk{}ह\appbrk{}प\appbrk{}क्ति\appbrk{}नुत्~। स्वेद\appbrk{}मू\appbrk{}र्च्छा\appbrk{}पि\appbrk{}पा\appbrk{}साघ्न}\ A; \skt{रौक्ष्य\appbrk{}वै\appbrk{}व\appbrk{}र्ण्ण्य\appbrk{}स्त\appbrk{}म्भ\appbrk{}कृ\appbrk{}द्दा\appbrk{}ह\appbrk{}प\appbrk{}क्ति\appbrk{}नुत्~॥ स्वेद\appbrk{}मू\appbrk{}र्च्छा\appbrk{}पि\appbrk{}पा\appbrk{}स\appbrk{}घ्नम्}\ H}}\pend

\pstart\versepar \ledsidenote{\passref{4.24.85}}\skt{सुखं }\edtext{\skt{प्रवातं}}{\lemma{\skt{प्रवातं}}\Afootnote{K H; \skt{वातं}\ A}}\skt{ }\edtext{\skt{सेवेत}}{\lemma{\skt{सेवेत}}\Afootnote{K H; \skt{प्रसेवेत}\ A}}\skt{ ग्रीष्मे शरदि }\edtext{\skt{चान्तरा}}{\lemma{\skt{चान्तरा}}\Afootnote{K H; \skt{मानवः~।}\ A}}\skt{~। }\pend
\pstart\versepar \skt{}\edtext{\skt{निवातं}}{\lemma{\skt{निवातं}}\Afootnote{\skt{निवातं ह्य्}\ A K; \skt{निरावातं ह्य्}\ H}}\skt{ }\edtext{\skt{ह्यायुषे}}{\lemma{\skt{ह्यायुषे}}\Afootnote{\skt{आयुषे}\ A K H}}\skt{ }\edtext{\skt{सेव्यमारोग्याय च}}{\lemma{\skt{सेव्यमारोग्याय च}}\Afootnote{\skt{सेव्यमारोग्याय च}\ A K H}}\skt{ सर्वदा~।}\pend

\pstart\versepar \ledsidenote{\passref{4.24.86}}\edtext{\skt{आत\-प\-पि\-त्त\-र\-क्ताग्निः}}{\lemma{\skt{आत\appbrk{}प\appbrk{}पि\appbrk{}त्त\appbrk{}र\appbrk{}क्ताग्निः}}\Afootnote{K\textsuperscript{pc}; \skt{आतपः पित्ततृष्णाग्नि}\ A; \skt{आत\appbrk{}पः\appbrk{}पि\appbrk{}त्त\appbrk{}र\appbrk{}क्ताग्निः}\ K\textsuperscript{ac} H}}\skt{ }\edtext{\skt{स्वेद\-मू\-र्च्छा\-तृ\-षा\-पहं}}{\lemma{\skt{स्वेद\appbrk{}मू\appbrk{}र्च्छा\appbrk{}तृ\appbrk{}षा\appbrk{}पहं}}\Afootnote{K; \skt{स्वेद\appbrk{}मू\appbrk{}र्च्छा\appbrk{}भ्र\appbrk{}मा\appbrk{}स्रकृ}\ A; \skt{स्वेद\appbrk{}मू\appbrk{}र्च्छा\appbrk{}कृ\appbrk{}षा\appbrk{}पहः~॥}\ H\textsuperscript{ac}; \skt{स्वेद\appbrk{}मू\appbrk{}र्च्छा\appbrk{}हृ\appbrk{}षा\appbrk{}व\appbrk{}हः\appbrk{}कृ\appbrk{}षा\appbrk{}पहः~॥}\ H\textsuperscript{pc}}}\skt{~। }\pend
\pstart\versepar \skt{}\edtext{\skt{दाह\-वै\-व\-र्ण्ण्य\-ज\-न\-न\-स्था\-या\-त्वे\-ता\-न्व्य\-पो\-हति}}{\lemma{\skt{दाह\appbrk{}वै\appbrk{}व\appbrk{}र्ण्ण्य\appbrk{}ज\appbrk{}न\appbrk{}न\appbrk{}स्था\appbrk{}या\appbrk{}त्वे\appbrk{}ता\appbrk{}न्व्य\appbrk{}पो\appbrk{}हति}}\Afootnote{K H; \skt{त्~। दाह\appbrk{}वै\appbrk{}व\appbrk{}र्ण्य\appbrk{}कारी च छाया चैतानपोहति~॥}\ A}}\skt{~।}\pend

\pstart\versepar \ledsidenote{\passref{4.24.87}}\edtext{\skt{अग्नि\-र्वा\-त\-क\-फ\-स्तम्भी}}{\lemma{\skt{अग्नि\appbrk{}र्वा\appbrk{}त\appbrk{}क\appbrk{}फ\appbrk{}स्तम्भी}}\Afootnote{\skt{अग्नि\appbrk{}र्वा\appbrk{}त\appbrk{}क\appbrk{}फ\appbrk{}स्तम्भ}\ A; \skt{अग्नि\appbrk{}र्वा\appbrk{}त\appbrk{}क\appbrk{}फ\appbrk{}स्तम्भी}\ K H}}\skt{ }\edtext{\skt{शीत\-वे\-प\-थु\-ना\-शनः}}{\lemma{\skt{शीत\appbrk{}वे\appbrk{}प\appbrk{}थु\appbrk{}ना\appbrk{}शनः}}\Afootnote{A K\textsuperscript{pc} H; \skt{शीत\appbrk{}वे\appbrk{}प\appbrk{}ना\appbrk{}शनः~।}\ K\textsuperscript{ac}}}\skt{~। }\pend
\pstart\versepar \skt{}\edtext{\skt{आमा\-भि\-ष्य\-न्द\-श\-मनो}}{\lemma{\skt{आमा\appbrk{}भि\appbrk{}ष्य\appbrk{}न्द\appbrk{}श\appbrk{}मनो}}\Afootnote{K H; \skt{आमा\appbrk{}भि\appbrk{}ष्य\appbrk{}न्द\appbrk{}ज\appbrk{}रशो}\ A}}\skt{ रक्त\-पि\-त्त\-प्र\-दू\-षणः~।}\pend

\pstart\versepar \ledsidenote{\passref{4.24.88}}\edtext{\skt{पुष्टि\-व\-र्ण्ण\-व\-लो\-त्सा\-हा\-म\-ग्नि\-दी\-प्ति\-र\-त\-न्द्रितां}}{\lemma{\skt{पुष्टि\appbrk{}व\appbrk{}र्ण्ण\appbrk{}व\appbrk{}लो\appbrk{}त्सा\appbrk{}हा\appbrk{}म\appbrk{}ग्नि\appbrk{}दी\appbrk{}प्ति\appbrk{}र\appbrk{}त\appbrk{}न्द्रितां}}\Afootnote{\skt{पुष्टि\appbrk{}व\appbrk{}र्ण\appbrk{}ब\appbrk{}लो\appbrk{}त्सा\appbrk{}ह\appbrk{}म\appbrk{}ग्नि\appbrk{}दी\appbrk{}प्ति\appbrk{}म\appbrk{}त\appbrk{}न्द्रि\appbrk{}ताम्~।}\ A; \skt{पुष्टि\appbrk{}व\appbrk{}र्ण्ण\appbrk{}व\appbrk{}लो\appbrk{}त्सा\appbrk{}हा\appbrk{}म\appbrk{}ग्नि\appbrk{}दी\appbrk{}प्ति\appbrk{}र\appbrk{}त\appbrk{}न्द्रितां~।}\ K; \skt{पुष्टि\appbrk{}व\appbrk{}र्ण्ण\appbrk{}व\appbrk{}लो\appbrk{}त्सा\appbrk{}ह\appbrk{}म\appbrk{}ग्नि\appbrk{}दी\appbrk{}प्ति\appbrk{}र\appbrk{}त\appbrk{}न्द्रितां~॥}\ H}}\skt{~। }\pend
\pstart\versepar \skt{करोति धातुसाम्यञ्च }\edtext{\skt{नित्रा\-का\-ल\-नि\-षे\-विता}}{\lemma{\skt{नित्रा\appbrk{}का\appbrk{}ल\appbrk{}नि\appbrk{}षे\appbrk{}विता}}\Afootnote{K; \skt{निद्रा काले निषेविता~॥}\ A; \skt{निद्रा\appbrk{}का\appbrk{}ल\appbrk{}नि\appbrk{}षे\appbrk{}विताः~॥}\ H}}\skt{~॥}\pend

\pstart \ledsidenote{\passref{4.24.89}}\edtext{\skt{तत्रादितो}}{\lemma{\skt{तत्रादितो}}\Afootnote{K H; \skt{तत्रादित}\ A}}\skt{ }\edtext{\skt{नीचनखरोम्णा}}{\lemma{\skt{नीचनखरोम्णा}}\Afootnote{K; \skt{एव नीचनखरोम्णा}\ A; \skt{नीचनखरोम्}\ H}}\skt{ }\edtext{\skt{शुचिवाससा}}{\lemma{\skt{शुचिवाससा}}\Afootnote{K; \skt{शुचिना शुक्लवाससा}\ A; \skt{ना शुचिवाससा}\ H}}\skt{ }\edtext{\skt{लघू\-ष्णी\-ष\-च्छ\-त्रो\-पा\-न\-त्के\-ना\-द\-र\-का\-रिणा}}{\lemma{\skt{लघू\appbrk{}ष्णी\appbrk{}ष\appbrk{}च्छ\appbrk{}त्रो\appbrk{}पा\appbrk{}न\appbrk{}त्के\appbrk{}ना\appbrk{}द\appbrk{}र\appbrk{}का\appbrk{}रिणा}}\Afootnote{K H; \skt{लघू\appbrk{}ष्णी\appbrk{}ष\appbrk{}च्छ\appbrk{}त्रो\appbrk{}पा\appbrk{}न\appbrk{}त्केन}\ A}}\skt{~। दण्डपाणिना }\edtext{\skt{काल\-हि\-त\-मि\-त\-म\-धु\-र\-पू\-र्वा\-भि\-भा\-षिणा}}{\lemma{\skt{काल\appbrk{}हि\appbrk{}त\appbrk{}मि\appbrk{}त\appbrk{}म\appbrk{}धु\appbrk{}र\appbrk{}पू\appbrk{}र्वा\appbrk{}भि\appbrk{}भा\appbrk{}षिणा}}\Afootnote{K H; \skt{काले हित\appbrk{}मि\appbrk{}त\appbrk{}म\appbrk{}धु\appbrk{}र\appbrk{}पू\appbrk{}र्वा\appbrk{}भि\appbrk{}भा\appbrk{}षिणा}\ A}}\skt{ }\edtext{\skt{वन्धुभूतेन}}{\lemma{\skt{वन्धुभूतेन}}\Afootnote{K H; \skt{बन्धुभूतेन}\ A}}\skt{ भूतानां }\edtext{\skt{गुरु\-वृ\-द्धा\-नु\-म\-ते\-ना\-न\-न्य\-म\-नसा}}{\lemma{\skt{गुरु\appbrk{}वृ\appbrk{}द्धा\appbrk{}नु\appbrk{}म\appbrk{}ते\appbrk{}ना\appbrk{}न\appbrk{}न्य\appbrk{}म\appbrk{}नसा}}\Afootnote{K H; \skt{गुरु\appbrk{}वृ\appbrk{}द्धा\appbrk{}नु\appbrk{}म\appbrk{}तेन सुस\appbrk{}हा\appbrk{}ये\appbrk{}ना\appbrk{}न\appbrk{}न्य\appbrk{}म\appbrk{}नसा}\ A}}\skt{ }\edtext{\skt{सुस\-हा\-ये\-नो\-प\-क\-र्त्तव्यं}}{\lemma{\skt{सुस\appbrk{}हा\appbrk{}ये\appbrk{}नो\appbrk{}प\appbrk{}क\appbrk{}र्त्तव्यं}}\Afootnote{K H; \skt{खलू\appbrk{}प\appbrk{}च\appbrk{}रि\appbrk{}तव्यं,}\ A}}\skt{~। }\edtext{\skt{तदपि न}}{\lemma{\skt{तदपि न}}\Afootnote{H; \skt{तदपि न}\ A K}}\skt{ रात्रौ~। न }\edtext{\skt{केशा\-स्थि\-क\-ण्ट\-का\-श्म\-तु\-ष\-भ\-स्मो\-त्क\-र\-क\-पा\-ला\-ङ्गा\-रा\-मे\-ध्य\-स्था\-न\-न\-व\-लि\-भू\-मिषु~। न}}{\lemma{\skt{केशा\appbrk{}स्थि\appbrk{}क\appbrk{}ण्ट\appbrk{}का\appbrk{}श्म\appbrk{}तु\appbrk{}ष\appbrk{}भ\appbrk{}स्मो\appbrk{}त्क\appbrk{}र\appbrk{}क\appbrk{}पा\appbrk{}ला\appbrk{}ङ्गा\appbrk{}रा\appbrk{}मे\appbrk{}ध्य\appbrk{}स्था\appbrk{}न\appbrk{}न\appbrk{}व\appbrk{}लि\appbrk{}भू\appbrk{}मिषु~। न}}\Afootnote{\skt{केशा\appbrk{}स्थि\appbrk{}क\appbrk{}ण्ट\appbrk{}का\appbrk{}श्म\appbrk{}तु\appbrk{}ष\appbrk{}भ\appbrk{}स्मो\appbrk{}त्क\appbrk{}र\appbrk{}क\appbrk{}पा\appbrk{}ला\appbrk{}ङ्गा\appbrk{}रा\appbrk{}मे\appbrk{}ध्य\appbrk{}स्ना\appbrk{}न\appbrk{}ब\appbrk{}लि\appbrk{}भू\appbrk{}मिषु, न}\ A; \skt{केशा\appbrk{}स्थि\appbrk{}क\appbrk{}ण्ट\appbrk{}का\appbrk{}श्म\appbrk{}तु\appbrk{}ष\appbrk{}भ\appbrk{}स्मो\appbrk{}त्क\appbrk{}र\appbrk{}क\appbrk{}पा\appbrk{}ला\appbrk{}ङ्गा\appbrk{}रा\appbrk{}मे\appbrk{}ध्य\appbrk{}स्था\appbrk{}न\appbrk{}व\appbrk{}लि\appbrk{}भू\appbrk{}मिषु~। न}\ K; \skt{केशा\appbrk{}स्शि\appbrk{}क\appbrk{}क\appbrk{}ण्ट\appbrk{}का\appbrk{}श्म\appbrk{}तु\appbrk{}ष\appbrk{}भ\appbrk{}स्मो\appbrk{}त्क\appbrk{}र\appbrk{}क\appbrk{}पा\appbrk{}ला\appbrk{}ङ्गा\appbrk{}रा\appbrk{}मे\appbrk{}ध्य\appbrk{}स्था\appbrk{}न\appbrk{}व\appbrk{}लि\appbrk{}भू\appbrk{}मिषु~। न}\ H}}\skt{ }\edtext{\skt{विष\-मे\-न्द्र\-की\-ल\-च\-तु\-ष्प\-थ\-स्व\-भ्रा\-णा\-मु\-प\-रि\-ष्टात्}}{\lemma{\skt{विष\appbrk{}मे\appbrk{}न्द्र\appbrk{}की\appbrk{}ल\appbrk{}च\appbrk{}तु\appbrk{}ष्प\appbrk{}थ\appbrk{}स्व\appbrk{}भ्रा\appbrk{}णाम्}}\Afootnote{K H; \skt{विष\appbrk{}मे\appbrk{}न्द्र\appbrk{}की\appbrk{}ल\appbrk{}च\appbrk{}तु\appbrk{}ष्प\appbrk{}थ\appbrk{}श्व\appbrk{}भ्रा\appbrk{}णाम्}\ A}}\skt{~।}\pend

\pstart \ledsidenote{\passref{4.24.90}}\edtext{\skt{न राज\-द्वि\-ष्ट\-प\-रु\-ष\-पै\-शू\-न्या\-नृ\-तानि}}{\lemma{\skt{न राज\appbrk{}द्वि\appbrk{}ष्ट\appbrk{}प\appbrk{}रु\appbrk{}ष\appbrk{}पै\appbrk{}शू\appbrk{}न्या\appbrk{}नृ\appbrk{}तानि}}\Afootnote{K; \skt{न राज\appbrk{}द्वि\appbrk{}ष्ट\appbrk{}प\appbrk{}रु\appbrk{}ष\appbrk{}पै\appbrk{}शु\appbrk{}न्या\appbrk{}नृ\appbrk{}तानि}\ A; \skt{न राज\appbrk{}द्वि\appbrk{}ष्ट\appbrk{}पु\appbrk{}रु\appbrk{}ष\appbrk{}पै\appbrk{}शू\appbrk{}न्या\appbrk{}नृ\appbrk{}तानि}\ H}}\skt{ }\edtext{\skt{वदेत्~। देव\-व्रा\-ह्म\-ण\-गु\-रु\-वृ\-द्ध\-पि\-तृ\-प\-रि\-वा\-दांश्च~। न}}{\lemma{\skt{वदेत्~। देव\appbrk{}व्रा\appbrk{}ह्म\appbrk{}ण\appbrk{}गु\appbrk{}रु\appbrk{}वृ\appbrk{}द्ध\appbrk{}पि\appbrk{}तृ\appbrk{}प\appbrk{}रि\appbrk{}वा\appbrk{}दांश्च~। न}}\Afootnote{K H; \skt{वदेत्, न देव\appbrk{}ब्रा\appbrk{}ह्म\appbrk{}ण\appbrk{}पि\appbrk{}तृ\appbrk{}प\appbrk{}रि\appbrk{}वा\appbrk{}दांश्च; न}\ A}}\skt{ }\edtext{\skt{नरे\-न्द्र\-द्वि\-ष्टो\-न्म\-त्त\-प\-ति\-ति\-क्षुच्च}}{\lemma{\skt{नरे\appbrk{}न्द्र\appbrk{}द्वि\appbrk{}ष्टो\appbrk{}न्म\appbrk{}त्त\appbrk{}प\appbrk{}ति\appbrk{}ति\appbrk{}क्षुच्च}}\Afootnote{K; \skt{नरे\appbrk{}न्द्र\appbrk{}द्वि\appbrk{}ष्टो\appbrk{}न्म\appbrk{}त्त\appbrk{}प\appbrk{}ति\appbrk{}त\appbrk{}क्षुद्र}\ A; \skt{नरे\appbrk{}न्द्र\appbrk{}द्वि\appbrk{}ष्टो\appbrk{}न्म\appbrk{}त्त\appbrk{}प\appbrk{}ति\appbrk{}त\appbrk{}क्षुच्च}\ H}}\skt{ }\edtext{\skt{भीतानुपासीत}}{\lemma{\skt{भीतानुपासीत}}\Afootnote{K H; \skt{नीचानुपासीत~॥}\ A}}\skt{~।}\pend

\pstart \ledsidenote{\passref{4.24.91}}\edtext{\skt{वृक्ष\-प\-र्व\-त\-प्र\-पा\-त\-वि\-ष\-म\-व\-ल्मी\-क\-दु\-ष्ट\-वा\-जिनि}}{\lemma{\skt{वृक्ष\appbrk{}प\appbrk{}र्व\appbrk{}त\appbrk{}प्र\appbrk{}पा\appbrk{}त\appbrk{}वि\appbrk{}ष\appbrk{}म\appbrk{}व\appbrk{}ल्मी\appbrk{}क\appbrk{}दु\appbrk{}ष्ट\appbrk{}वा\appbrk{}जिनि}}\Afootnote{K; \skt{वृक्ष\appbrk{}प\appbrk{}र्व\appbrk{}त\appbrk{}प्र\appbrk{}पा\appbrk{}त\appbrk{}वि\appbrk{}ष\appbrk{}म\appbrk{}व\appbrk{}ल्मी\appbrk{}क\appbrk{}दु\appbrk{}ष्ट\appbrk{}वा\appbrk{}जि\appbrk{}कु\appbrk{}ञ्ज\appbrk{}रा\appbrk{}द्य\appbrk{}धि\appbrk{}रो\appbrk{}ह\appbrk{}णानि}\ A; \skt{वृक्ष\appbrk{}प\appbrk{}र्व्व\appbrk{}त\appbrk{}प्र\appbrk{}पा\appbrk{}त\appbrk{}वि\appbrk{}ष\appbrk{}म\appbrk{}व\appbrk{}ल्मी\appbrk{}क\appbrk{}दु\appbrk{}ष्ट\appbrk{}वा\appbrk{}जि\appbrk{}कु\appbrk{}ञ्ज\appbrk{}रा\appbrk{}रो\appbrk{}ह\appbrk{}णानि}\ H}}\skt{ }\edtext{\skt{परिहरेत्~। पूर्ण्ण\-न\-दी\-स\-मु\-द्रा\-वि\-दि\-त\-प\-ल्व\-ल\-स्व\-भ्र\-कू\-पा\-व\-त\-र\-णानि च}}{\lemma{\skt{परिहरेत्~। पूर्ण्ण\appbrk{}न\appbrk{}दी\appbrk{}स\appbrk{}मु\appbrk{}द्रा\appbrk{}वि\appbrk{}दि\appbrk{}त\appbrk{}प\appbrk{}ल्व\appbrk{}ल\appbrk{}स्व\appbrk{}भ्र\appbrk{}कू\appbrk{}पा\appbrk{}व\appbrk{}त\appbrk{}र\appbrk{}णानि च}}\Afootnote{K H; \skt{परिहरेत्, पूर्ण\appbrk{}न\appbrk{}दी\appbrk{}स\appbrk{}मु\appbrk{}द्रा\appbrk{}वि\appbrk{}दि\appbrk{}त\appbrk{}प\appbrk{}ल्व\appbrk{}ल\appbrk{}श्व\appbrk{}भ्र\appbrk{}कू\appbrk{}पा\appbrk{}व\appbrk{}त\appbrk{}र\appbrk{}णानि,}\ A}}\skt{ }\edtext{\skt{भिन्न}}{\lemma{\skt{भिन्न}}\Afootnote{A K; \skt{भिन्ना}\ H}}\skt{ }\edtext{\skt{शून्या\-गा\-र\-श्म\-शा\-न\-वि\-ज\-ना\-र\-ण्य\-ग्रा\-म\-से\-वाश्च}}{\lemma{\skt{शून्या\appbrk{}गा\appbrk{}र\appbrk{}श्म\appbrk{}शा\appbrk{}न\appbrk{}वि\appbrk{}ज\appbrk{}ना\appbrk{}र\appbrk{}ण्य\appbrk{}ग्रा\appbrk{}म\appbrk{}से\appbrk{}वाश्च}}\Afootnote{K H; \skt{शून्या\appbrk{}गा\appbrk{}र\appbrk{}श्म\appbrk{}शा\appbrk{}न\appbrk{}वि\appbrk{}ज\appbrk{}ना\appbrk{}र\appbrk{}ण्य\appbrk{}वा\appbrk{}सा\appbrk{}ग्नि\appbrk{}सं\appbrk{}भ्र\appbrk{}म\appbrk{}व्या\appbrk{}ल\appbrk{}भु\appbrk{}ज\appbrk{}ङ्ग\appbrk{}की\appbrk{}ट\appbrk{}से\appbrk{}वाश्च,}\ A}}\skt{~। }\edtext{\skt{ग्राम\-घा\-त\-क\-ल\-ह\-श\-स्त्र\-स\-न्नि\-पा\-ता\-ग्नि\-स\-म्भ्र\-म\-व्या\-ड\-स\-रीसृ}}{\lemma{\skt{ग्राम\appbrk{}घा\appbrk{}त\appbrk{}क\appbrk{}ल\appbrk{}ह\appbrk{}श\appbrk{}स्त्र\appbrk{}स\appbrk{}न्नि\appbrk{}पा\appbrk{}ता\appbrk{}ग्नि\appbrk{}स\appbrk{}म्भ्र\appbrk{}म\appbrk{}व्या\appbrk{}ड\appbrk{}स\appbrk{}रीसृ}}\Afootnote{K H; \skt{ग्रामा\appbrk{}घा\appbrk{}त\appbrk{}क\appbrk{}ल\appbrk{}ह\appbrk{}श\appbrk{}स्त्र\appbrk{}स\appbrk{}न्नि\appbrk{}पा\appbrk{}त\appbrk{}व्या\appbrk{}ल\appbrk{}स\appbrk{}रीसृ}\ A}}\skt{ }\edtext{\skt{पशृङ्गिकषांश्च}}{\lemma{\skt{पशृङ्गिकषांश्}}\Afootnote{K H; \skt{पशृ\appbrk{}ङ्गि\appbrk{}स\appbrk{}न्नि\appbrk{}क\appbrk{}र्षांश्}\ A}}\skt{~।}\pend

\pstart \ledsidenote{\passref{4.24.92}}\edtext{\edtext{\skt{नाग्नि\-व्रा\-ह्म\-ण\-प्रे\-क्षा\-द\-म्प\-त्यो\-र\-न्त\-रे\-णा\-ती\-यात्}}{\lemma{\skt{नाग्नि\appbrk{}व्रा\appbrk{}ह्म\appbrk{}ण\appbrk{}प्रे\appbrk{}क्षा\appbrk{}द\appbrk{}म्प\appbrk{}त्योर्}}\Afootnote{K H; \skt{नाग्नि\appbrk{}गो\appbrk{}गु\appbrk{}रु\appbrk{}ब्रा\appbrk{}ह्म\appbrk{}ण\appbrk{}प्रे\appbrk{}ङ्खा\appbrk{}द\appbrk{}म्प\appbrk{}त्यन्}\ A}}\skt{~। न}}{\lemma{\skt{अन्त\appbrk{}रे\appbrk{}णा\appbrk{}ती\appbrk{}यात्~। न}}\Afootnote{K H; \skt{तरेण यायात्~। न}\ A}}\skt{ }\edtext{\edtext{\skt{शव\-म\-नु\-या\-यात्}}{\lemma{\skt{शवम्}}\Afootnote{A K; \skt{शव}\ H}}\skt{~। }\edtext{\skt{देव\-गो\-व्रा\-ह्म\-ण\-चै\-त्य\-ध्व\-ज\-रो\-गि\-प\-ति\-त\-पा\-प\-क\-र्त्तृ\-णाञ्च}}{\lemma{\skt{अध्व\appbrk{}ज\appbrk{}रो\appbrk{}गि\appbrk{}प\appbrk{}ति\appbrk{}त\appbrk{}पा\appbrk{}प\appbrk{}क\appbrk{}र्त्तृ\appbrk{}णाञ्च}}\Afootnote{\skt{अध्व\appbrk{}ज\appbrk{}रो\appbrk{}गि\appbrk{}प\appbrk{}ति\appbrk{}त\appbrk{}पा\appbrk{}प\appbrk{}का\appbrk{}रिणां च}\ A; \skt{त्यध्व\appbrk{}ज\appbrk{}रो\appbrk{}गि\appbrk{}प\appbrk{}ति\appbrk{}त\appbrk{}पा\appbrk{}प\appbrk{}क\appbrk{}र्त्तृ\appbrk{}णाञ्च}\ K H}}}{\lemma{\skt{अनुयायात्~। देव\appbrk{}गो\appbrk{}व्रा\appbrk{}ह्म\appbrk{}ण\appbrk{}चैत्य्}}\Afootnote{\skt{अनुयायात्~। देव\appbrk{}गो\appbrk{}ब्रा\appbrk{}ह्म\appbrk{}ण\appbrk{}चैत्य्}\ A; \skt{अनुयायात्~। देव\appbrk{}गो\appbrk{}व्रा\appbrk{}ह्म\appbrk{}णचै}\ K; \skt{मनुयायात्~। देवगोव्राह्मणं चै}\ H}}\skt{ }\edtext{\skt{छायान्नाक्रामेस्तं}}{\lemma{\skt{छायान्नाक्रामेस्तं}}\Afootnote{\skt{छायां नाक्रमेत~। नास्तं}\ A; \skt{छाया\appbrk{}न्ना\appbrk{}क्रा\appbrk{}मे\appbrk{}त्नास्तं}\ K H}}\skt{ }\edtext{\edtext{\skt{गच्छ\-न्त\-मु\-द्य\-न्त\-मा\-दित्यं}}{\lemma{\skt{उद्यन्तम्}}\Afootnote{K H; \skt{उद्यन्तं वा}\ A}}}{\lemma{\skt{आदित्यं}}\Afootnote{K; \skt{ऽऽआदित्यं}\ A; \skt{आ}\ H}}\skt{ }\edtext{\skt{पश्येत्~। गां}}{\lemma{\skt{पश्येत्~। गां}}\Afootnote{K; \skt{वीक्षेत~। गां धापयन्तीं}\ A; \skt{दित्यम्पश्येत्~। गान्}\ H}}\skt{ धयन्तीं }\edtext{\skt{परस्यं वा}}{\lemma{\skt{परस्यं वा}}\Afootnote{K H; \skt{परशस्यं वा}\ A}}\skt{ }\edtext{\edtext{\skt{चरन्तीम्परस्यै}}{\lemma{\skt{चरन्तीम्}}\Afootnote{A K; \skt{चयन्तीम्}\ H}}}{\lemma{\skt{परस्यै}}\Afootnote{K H; \skt{न कस्मै}\ A}}\skt{ }\edtext{\skt{नाचक्षीत}}{\lemma{\skt{नाचक्षीत}}\Afootnote{K H; \skt{चिदाचक्षीत, न}\ A}}\skt{~। }\edtext{\skt{नोल्का\-पा\-ता\-त्पा\-ते\-न्द्र\-ध\-नू\-षिवा}}{\lemma{\skt{नोल्का\appbrk{}पा\appbrk{}ता\appbrk{}त्पा\appbrk{}ते\appbrk{}न्द्र\appbrk{}ध\appbrk{}नू\appbrk{}षिवा}}\Afootnote{K; \skt{चोल्का\appbrk{}पा\appbrk{}तो\appbrk{}त्पा\appbrk{}ते\appbrk{}न्द्र\appbrk{}ध\appbrk{}नूंषि~।}\ A; \skt{नोल्का\appbrk{}पा\appbrk{}ता\appbrk{}त्पा\appbrk{}ते\appbrk{}न्द्र\appbrk{}ध\appbrk{}नूंषि वा~।}\ H}}\skt{~। नाग्निं }\edtext{\skt{मुखे\-नो\-प\-ध\-मेत्~। नायो}}{\lemma{\skt{मुखे\appbrk{}नो\appbrk{}प\appbrk{}ध\appbrk{}मेत्~। नायो}}\Afootnote{K; \skt{मुखे\appbrk{}नो\appbrk{}प\appbrk{}ध\appbrk{}मेत्~। नापो}\ A H}}\skt{ भूमिं वा पाणि\-पा\-दे\-ना\-भि\-ह\-न्यात्~।}\pend

\pstart \ledsidenote{\passref{4.24.93}}\edtext{\skt{न वेगान्धारयेत्~। न}}{\lemma{\skt{न वेगान्धारयेत्~। न}}\Afootnote{K H; \skt{न वेगा\appbrk{}न्धा\appbrk{}र\appbrk{}ये\appbrk{}द्वा\appbrk{}त\appbrk{}मू\appbrk{}त्र\appbrk{}पु\appbrk{}री\appbrk{}षा\appbrk{}दी\appbrk{}नाम्~। न}\ A}}\skt{ }\edtext{\skt{वहि\-र्भा\-गां\-ग्रा\-म\-न\-ग\-र\-दे\-व\-ता\-य\-त\-न\-श्म\-शा\-न\-प\-थ\-स\-लि\-ला\-श\-य\-स\-न्नि\-कृ\-ष्ठा\-नु\-त्सृ\-जेन्न}}{\lemma{\skt{वहि\appbrk{}र्भा\appbrk{}गां\appbrk{}ग्रा\appbrk{}म\appbrk{}न\appbrk{}ग\appbrk{}र\appbrk{}दे\appbrk{}व\appbrk{}ता\appbrk{}य\appbrk{}त\appbrk{}न\appbrk{}श्म\appbrk{}शा\appbrk{}न\appbrk{}प\appbrk{}थ\appbrk{}स\appbrk{}लि\appbrk{}ला\appbrk{}श\appbrk{}य\appbrk{}स\appbrk{}न्नि\appbrk{}कृ\appbrk{}ष्ठा\appbrk{}नु\appbrk{}त्सृ\appbrk{}जेन्न}}\Afootnote{K; \skt{बहि\appbrk{}र्वे\appbrk{}गा\appbrk{}न्ग्रा\appbrk{}म\appbrk{}न\appbrk{}ग\appbrk{}र\appbrk{}दे\appbrk{}व\appbrk{}ता\appbrk{}य\appbrk{}त\appbrk{}न\appbrk{}श्म\appbrk{}शा\appbrk{}न\appbrk{}च\appbrk{}तु\appbrk{}ष्प\appbrk{}थ\appbrk{}स\appbrk{}लि\appbrk{}ला\appbrk{}श\appbrk{}य\appbrk{}प\appbrk{}थि\appbrk{}स\appbrk{}न्नि\appbrk{}कृ\appbrk{}ष्टा\appbrk{}वु\appbrk{}त्सृ\appbrk{}जेन्न}\ A; \skt{वहि\appbrk{}र्व्भा\appbrk{}गा\appbrk{}ग्रा\appbrk{}म\appbrk{}न\appbrk{}ग\appbrk{}र\appbrk{}दे\appbrk{}व\appbrk{}ता\appbrk{}य\appbrk{}त\appbrk{}न\appbrk{}श्म\appbrk{}शा\appbrk{}न\appbrk{}प\appbrk{}थ\appbrk{}स\appbrk{}लि\appbrk{}ला\appbrk{}श\appbrk{}य\appbrk{}स\appbrk{}न्नि\appbrk{}कृ\appbrk{}ष्ठा\appbrk{}नु\appbrk{}त्सृजे}\ H}}\skt{ }\edtext{\skt{प्रकाशं~। न}}{\lemma{\skt{प्रकाशं~। न}}\Afootnote{A K; \skt{त्~। न}\ H}}\skt{ }\edtext{\skt{वाय्व\-ग्नि\-स\-लि\-ल\-सो\-मा\-र्क\-गो\-गुरूं}}{\lemma{\skt{वाय्व\appbrk{}ग्नि\appbrk{}स\appbrk{}लि\appbrk{}ल\appbrk{}सो\appbrk{}मा\appbrk{}र्क\appbrk{}गो\appbrk{}गुरूं}}\Afootnote{K H; \skt{वाय्व\appbrk{}ग्नि\appbrk{}स\appbrk{}लि\appbrk{}ल\appbrk{}सो\appbrk{}मा\appbrk{}र्क\appbrk{}गो\appbrk{}गुरु}\ A}}\skt{ }\edtext{\skt{प्रतिमुखम्वा}}{\lemma{\skt{वा}}\Afootnote{K H; \emph{om.}\ A}}\skt{~।}\pend

\pstart \ledsidenote{\passref{4.24.94}}\skt{न भूमिं विलिखेत्~। नासंवृतमुखः सदसि }\edtext{\skt{श्वास\-का\-स\-जृ\-म्भो\-द्भा\-र\-क्ष\-व\-थून्न}}{\lemma{\skt{श्वास\appbrk{}का\appbrk{}स\appbrk{}जृ\appbrk{}म्भो\appbrk{}द्भा\appbrk{}र\appbrk{}क्ष\appbrk{}व\appbrk{}थून्न}}\Afootnote{K; \skt{जृम्भो\appbrk{}द\appbrk{}गा\appbrk{}र\appbrk{}का\appbrk{}स\appbrk{}श्वा\appbrk{}स\appbrk{}क्ष\appbrk{}व\appbrk{}थूनु}\ A; \skt{स्वास\appbrk{}का\appbrk{}स\appbrk{}जृ\appbrk{}म्भो\appbrk{}द्भा\appbrk{}र\appbrk{}क्ष\appbrk{}व\appbrk{}थून्न}\ H}}\skt{ }\edtext{\skt{सृजेत्~। न}}{\lemma{\skt{सृजेत्~। न}}\Afootnote{K H; \skt{त्सृजेत्, न}\ A}}\skt{ }\edtext{\skt{पर्य\-व\-ष्टि\-का\-व\-ष्ट\-म्भ\-पा\-द\-प्र\-सा\-र\-णानि}}{\lemma{\skt{पर्य\appbrk{}व\appbrk{}ष्टि\appbrk{}का\appbrk{}व\appbrk{}ष्ट\appbrk{}म्भ\appbrk{}पा\appbrk{}द\appbrk{}प्र\appbrk{}सा\appbrk{}र\appbrk{}णानि}}\Afootnote{K H; \skt{पर्य\appbrk{}ङ्कि\appbrk{}का\appbrk{}व\appbrk{}ष्ट\appbrk{}म्भ\appbrk{}पा\appbrk{}द\appbrk{}प्र\appbrk{}सा\appbrk{}र\appbrk{}णानि}\ A}}\skt{ गुरुसन्निधौ कुर्यात्~।}\pend

\pstart \ledsidenote{\passref{4.24.95}}\edtext{\skt{न वा}}{\lemma{\skt{न वा}}\Afootnote{K H; \skt{न बा}\ A}}\skt{ }\edtext{\skt{कर्ण्ण\-न\-सा\-स्रो\-तो\-द\-श\-न\-वि\-व\-रा\-ण्य\-भि\-नि\-ष्कु\-ष्णी\-यात्~। न}}{\lemma{\skt{कर्ण्ण\appbrk{}न\appbrk{}सा\appbrk{}स्रो\appbrk{}तो\appbrk{}द\appbrk{}श\appbrk{}न\appbrk{}वि\appbrk{}व\appbrk{}रा\appbrk{}ण्य\appbrk{}भि\appbrk{}नि\appbrk{}ष्कु\appbrk{}ष्णी\appbrk{}यात्~। न}}\Afootnote{K H; \skt{लक\appbrk{}र्ण\appbrk{}ना\appbrk{}सा\appbrk{}स्रो\appbrk{}तो\appbrk{}द\appbrk{}श\appbrk{}ना\appbrk{}क्षि\appbrk{}वि\appbrk{}व\appbrk{}रा\appbrk{}ण्य\appbrk{}भि\appbrk{}कु\appbrk{}ष्णी\appbrk{}यात्, न}\ A}}\skt{ }\edtext{\skt{वेज\-ये\-त्के\-श\-न\-ख\-व\-स्त्र\-गा\-त्राणि~। न}}{\lemma{\skt{वेज\appbrk{}ये\appbrk{}त्के\appbrk{}श\appbrk{}न\appbrk{}ख\appbrk{}व\appbrk{}स्त्र\appbrk{}गा\appbrk{}त्राणि~। न}}\Afootnote{K H; \skt{वीज\appbrk{}ये\appbrk{}त्के\appbrk{}श\appbrk{}मु\appbrk{}ख\appbrk{}न\appbrk{}ख\appbrk{}व\appbrk{}स्त्र\appbrk{}गा\appbrk{}त्राणि, न}\ A}}\skt{ }\edtext{\skt{गात्र\-व\-क्त्र\-न\-ख\-वा\-दित्रः}}{\lemma{\skt{गात्र\appbrk{}व\appbrk{}क्त्र\appbrk{}न\appbrk{}ख\appbrk{}वा\appbrk{}दित्रः}}\Afootnote{K H; \skt{गात्र\appbrk{}न\appbrk{}ख\appbrk{}व\appbrk{}क्त्र\appbrk{}वा\appbrk{}दित्रं}\ A}}\skt{ कुर्यात्~। न }\edtext{\skt{काष्ठ\-लो\-ष्ठ\-तृ\-णा\-दी\-न\-भि\-ह\-न्यात्~। च्छिन्द्याद्भिन्द्याद्वा}}{\lemma{\skt{काष्ठ\appbrk{}लो\appbrk{}ष्ठ\appbrk{}तृ\appbrk{}णा\appbrk{}दी\appbrk{}न\appbrk{}भि\appbrk{}ह\appbrk{}न्यात्~। च्छिन्द्याद्भिन्द्याद्वा}}\Afootnote{K; \skt{काष्ठ\appbrk{}लो\appbrk{}ष्ट\appbrk{}तृ\appbrk{}णा\appbrk{}दी\appbrk{}न\appbrk{}भि\appbrk{}ह\appbrk{}न्या\appbrk{}च्छि\appbrk{}न्द्या\appbrk{}द्भि\appbrk{}न्द्याद्वा~॥}\ A; \skt{काष्ठ\appbrk{}लो\appbrk{}ष्ठ\appbrk{}तृ\appbrk{}णा\appbrk{}दी\appbrk{}ना\appbrk{}भि\appbrk{}ह\appbrk{}न्यात्~। छिन्द्याद्भिन्द्याद्वा~।}\ H}}\skt{~।}\pend

\pstart \ledsidenote{\passref{4.24.96}}\edtext{\edtext{\skt{नाति\-प्र\-वा\-ता\-त\-प\-मु\-प\-से\-वेत}}{\lemma{\skt{नाति\appbrk{}प्र\appbrk{}वा\appbrk{}ता\appbrk{}त\appbrk{}पम्}}\Afootnote{K H\textsuperscript{pc}; \skt{न प्रतिवातातपं}\ A; \skt{नाति\appbrk{}प्र\appbrk{}वा\appbrk{}वा\appbrk{}ना\appbrk{}ति\appbrk{}प्र\appbrk{}वा\appbrk{}ता\appbrk{}त\appbrk{}पम्}\ H\textsuperscript{ac}}}\skt{~। न}}{\lemma{\skt{उपसेवेत~। न}}\Afootnote{K H; \skt{सेवेत, न}\ A}}\skt{ }\edtext{\skt{मुक्त\-मा\-त्रा\-ग्नि\-मु\-पा\-सीत}}{\lemma{\skt{मुक्त\appbrk{}मा\appbrk{}त्रा\appbrk{}ग्नि\appbrk{}मु\appbrk{}पा\appbrk{}सीत}}\Afootnote{K; \skt{भुक्त\appbrk{}मा\appbrk{}त्रोऽ\appbrk{}ग्नि\appbrk{}मु\appbrk{}पा\appbrk{}सीत,}\ A; \skt{मुक्त\appbrk{}मा\appbrk{}त्राऽ\appbrk{}ग्नि\appbrk{}मु\appbrk{}पा\appbrk{}सीत~।}\ H}}\skt{~। }\edtext{\skt{नोत्कू\-ट्र\-का\-स\-ना\-ल्पा\-स\-ना\-न्द\-ध्या\-सीत~। न}}{\lemma{\skt{नोत्कू\appbrk{}ट्र\appbrk{}का\appbrk{}स\appbrk{}ना\appbrk{}ल्पा\appbrk{}स\appbrk{}ना\appbrk{}न्द\appbrk{}ध्या\appbrk{}सीत~। न}}\Afootnote{\skt{नोत्क\appbrk{}ट\appbrk{}का\appbrk{}ल्प\appbrk{}का\appbrk{}ष्ठा\appbrk{}स\appbrk{}न\appbrk{}म\appbrk{}ध्या\appbrk{}सीत, न}\ A; \skt{नोत्कू\appbrk{}ट्र\appbrk{}का\appbrk{}स\appbrk{}ना\appbrk{}ल्पा\appbrk{}स\appbrk{}ना\appbrk{}न्द\appbrk{}ध्या\appbrk{}सीत~। न}\ K; \skt{नोत्क\appbrk{}टु\appbrk{}का\appbrk{}स\appbrk{}ना\appbrk{}ल्पा\appbrk{}स\appbrk{}ना\appbrk{}न्य\appbrk{}ध्या\appbrk{}सीत्~। न}\ H}}\skt{ }\edtext{\skt{ग्रीवं}}{\lemma{\skt{ग्रीवं}}\Afootnote{K; \skt{ग्रीवां}\ A H}}\skt{ विषमं }\edtext{\skt{हरेत्~। न}}{\lemma{\skt{हरेत्~। न}}\Afootnote{K H; \skt{धारयेत्, न}\ A}}\skt{ }\edtext{\skt{विष\-म\-प\-ति\-त\-कायः}}{\lemma{\skt{विष\appbrk{}म\appbrk{}प\appbrk{}ति\appbrk{}त\appbrk{}कायः}}\Afootnote{K H; \skt{विषमकायः}\ A}}\skt{ क्रियां }\edtext{\skt{सेवे\-द्व\-से\-द्भु\-ञ्जीत वा~। न}}{\lemma{\skt{सेवे\appbrk{}द्व\appbrk{}से\appbrk{}द्भु\appbrk{}ञ्जीत वा~। न}}\Afootnote{K H; \skt{भजेत भुञ्जीत वा, न}\ A}}\skt{ }\edtext{\skt{प्रत\-त\-मु\-दी\-क्षेत}}{\lemma{\skt{प्रत\appbrk{}त\appbrk{}मु\appbrk{}दी\appbrk{}क्षेत}}\Afootnote{K; \skt{प्रततमीक्षेत}\ A; \skt{प्रत\appbrk{}त\appbrk{}मु\appbrk{}दी\appbrk{}क्षेत~।}\ H}}\skt{~। }\edtext{\skt{विशे\-षा\-र्ज्जो\-ति\-र्भा\-स्क\-र\-सू\-ष्म\-ब\-ल\-भ्रा\-न्तानि~। न}}{\lemma{\skt{विशे\appbrk{}षा\appbrk{}र्ज्जो\appbrk{}ति\appbrk{}र्भा\appbrk{}स्क\appbrk{}र\appbrk{}सू\appbrk{}ष्म\appbrk{}ब\appbrk{}ल\appbrk{}भ्रा\appbrk{}न्तानि~। न}}\Afootnote{\skt{विशे\appbrk{}षा\appbrk{}ज्ज्यो\appbrk{}ति\appbrk{}र्भा\appbrk{}स्क\appbrk{}र\appbrk{}सू\appbrk{}क्ष्म\appbrk{}च\appbrk{}ल\appbrk{}भ्रा\appbrk{}न्तानि, न}\ A; \skt{विशे\appbrk{}षा\appbrk{}र्ज्जो\appbrk{}ति\appbrk{}र्भा\appbrk{}स्क\appbrk{}र\appbrk{}सू\appbrk{}ष्म\appbrk{}व\appbrk{}ल\appbrk{}भ्रा\appbrk{}न्तानि~। न}\ K; \skt{विशे\appbrk{}षा\appbrk{}ज्जो\appbrk{}ति\appbrk{}र्भा\appbrk{}स्क\appbrk{}र\appbrk{}सू\appbrk{}ष्म\appbrk{}व\appbrk{}ल\appbrk{}भ्रा\appbrk{}न्तानि~। न}\ H}}\skt{ भारं शिरसा }\edtext{\skt{वहेत्~। न च}}{\lemma{\skt{वहेत्~। न च}}\Afootnote{K; \skt{वहेत्, न}\ A; \skt{भवेत्~। न च}\ H\textsuperscript{ac}; \skt{वहेभवेत्~। न च}\ H\textsuperscript{pc}}}\skt{ }\edtext{\edtext{\skt{स्वप्न\-जा\-ग\-र\-ण\-श\-य\-ना\-स\-न\-स्था\-न\-च\-ङ्क्र\-म\-न\-या\-न\-वा\-ह\-न\-प्र\-धा\-व\-न\-ल\-घ\-न\-प्र\-त\-र\-ण\-हा\-स्य\-भा\-ष्य\-व्य\-या\-य\-व्या\-या\-मा\-दी\-नु\-चि\-ता\-न\-प्य\-ति\-मात्रं}}{\lemma{\skt{स्वप्न\appbrk{}जा\appbrk{}ग\appbrk{}र\appbrk{}ण\appbrk{}श\appbrk{}य\appbrk{}ना\appbrk{}स\appbrk{}न\appbrk{}स्था\appbrk{}न\appbrk{}च\appbrk{}ङ्क्र\appbrk{}म\appbrk{}न\appbrk{}या\appbrk{}न\appbrk{}वा\appbrk{}ह\appbrk{}न\appbrk{}प्र\appbrk{}धा\appbrk{}व\appbrk{}न\appbrk{}ल\appbrk{}घ\appbrk{}न\appbrk{}प्र\appbrk{}त\appbrk{}र\appbrk{}ण\appbrk{}हा\appbrk{}स्य\appbrk{}भा\appbrk{}ष्य\appbrk{}व्य\appbrk{}या\appbrk{}य\appbrk{}व्या\appbrk{}या\appbrk{}मा\appbrk{}दी\appbrk{}नु\appbrk{}चि\appbrk{}ता\appbrk{}नप्य्}}\Afootnote{\skt{स्वप्न\appbrk{}जा\appbrk{}ग\appbrk{}र\appbrk{}ण\appbrk{}श\appbrk{}य\appbrk{}ना\appbrk{}स\appbrk{}न\appbrk{}स्था\appbrk{}न\appbrk{}च\appbrk{}ङ्क्र\appbrk{}म\appbrk{}ण\appbrk{}या\appbrk{}न\appbrk{}वा\appbrk{}ह\appbrk{}न\appbrk{}प्र\appbrk{}धा\appbrk{}व\appbrk{}न\appbrk{}ल\appbrk{}ङ्घ\appbrk{}न\appbrk{}प्ल\appbrk{}व\appbrk{}न\appbrk{}प्र\appbrk{}त\appbrk{}र\appbrk{}ण\appbrk{}हा\appbrk{}स्य\appbrk{}भा\appbrk{}ष्य\appbrk{}व्य\appbrk{}वा\appbrk{}य\appbrk{}व्या\appbrk{}या\appbrk{}मा\appbrk{}दी\appbrk{}नु\appbrk{}चि\appbrk{}ता\appbrk{}नप्य्}\ A; \skt{स्वप्न\appbrk{}जा\appbrk{}ग\appbrk{}र\appbrk{}ण\appbrk{}श\appbrk{}य\appbrk{}ना\appbrk{}स\appbrk{}न\appbrk{}स्था\appbrk{}न\appbrk{}च\appbrk{}ङ्क्र\appbrk{}म\appbrk{}न\appbrk{}या\appbrk{}न\appbrk{}वा\appbrk{}ह\appbrk{}न\appbrk{}प्र\appbrk{}धा\appbrk{}व\appbrk{}न\appbrk{}ल\appbrk{}घ\appbrk{}न\appbrk{}प्र\appbrk{}त\appbrk{}र\appbrk{}ण\appbrk{}हा\appbrk{}स्य\appbrk{}भा\appbrk{}ष्य\appbrk{}व्य\appbrk{}या\appbrk{}य\appbrk{}व्या\appbrk{}या\appbrk{}मा\appbrk{}दी\appbrk{}नु\appbrk{}चि\appbrk{}ता\appbrk{}नप्य्}\ K; \skt{स्वप्न\appbrk{}प्र\appbrk{}जा\appbrk{}ग\appbrk{}र\appbrk{}ण\appbrk{}श\appbrk{}य\appbrk{}ना\appbrk{}स\appbrk{}न\appbrk{}स्था\appbrk{}न\appbrk{}च\appbrk{}ङ्क्र\appbrk{}म\appbrk{}न\appbrk{}या\appbrk{}न\appbrk{}वा\appbrk{}ह\appbrk{}न\appbrk{}प्र\appbrk{}धा\appbrk{}व\appbrk{}न\appbrk{}लं\appbrk{}घ\appbrk{}न\appbrk{}प्र\appbrk{}त\appbrk{}र\appbrk{}ण\appbrk{}हा\appbrk{}स्य\appbrk{}भा\appbrk{}ष्य\appbrk{}व्य\appbrk{}याय व्याया\appbrk{}मा\appbrk{}दी\appbrk{}नु\appbrk{}चि\appbrk{}ता\appbrk{}नप्य्}\ H}}}{\lemma{\skt{अतिमात्रं}}\Afootnote{K H; \skt{अति}\ A}}\skt{ सेवेत~।}\pend

\pstart \ledsidenote{\passref{4.24.97}}\edtext{\edtext{\skt{उचि\-ता\-द\-प्य\-हि\-ता\-त्क्र\-मशो}}{\lemma{\skt{उचितादप्य्}}\Afootnote{A H; \skt{उचिताद}\ K}}}{\lemma{\skt{अहितात्क्रमशो}}\Afootnote{A H; \skt{प्यहितात्क्रमशो}\ K}}\skt{ }\edtext{\skt{विरमेत}}{\lemma{\skt{विरमेत}}\Afootnote{K H; \skt{विरमे}\ A}}\skt{~। }\edtext{\edtext{\edtext{\edtext{\skt{हित\-म\-नु\-चि\-त\-म\-प्या\-से\-वेत}}{\lemma{\skt{हितम्}}\Afootnote{H; \skt{त्, हितम्}\ A; \skt{हित}\ K}}}{\lemma{\skt{अनुचितम्}}\Afootnote{A H; \skt{मनुचित}\ K}}}{\lemma{\skt{अप्य्}}\Afootnote{A H; \skt{म}\ K}}}{\lemma{\skt{आसेवेत}}\Afootnote{H; \skt{असेवेत}\ A; \skt{प्यासेवेत}\ K}}\skt{ }\edtext{\skt{क्रमशः}}{\lemma{\skt{क्रमशः}}\Afootnote{K H; \skt{क्रमशः, न चैकान्ततः पादहीनात्~॥}\ A}}\skt{~।}\pend

\pstart \ledsidenote{\passref{4.24.98}}\edtext{\skt{नावा\-क्छि\-रा\-श\-यीत~। न}}{\lemma{\skt{नावा\appbrk{}क्छि\appbrk{}रा\appbrk{}श\appbrk{}यीत~। न}}\Afootnote{K H; \skt{नावाक्शिराः शयीत, न}\ A}}\skt{ भिन्नपात्रे }\edtext{\skt{भुञ्जीत}}{\lemma{\skt{भुञ्जीत}}\Afootnote{K H; \skt{भुञ्जीत, न}\ A}}\skt{~। }\edtext{\skt{नाञ्ज\-लि\-पु\-टे\-नापः}}{\lemma{\skt{नाञ्ज\appbrk{}लि\appbrk{}पु\appbrk{}टे\appbrk{}नापः}}\Afootnote{K H; \skt{विना पात्रेण, नाञ्ज\appbrk{}लि\appbrk{}पु\appbrk{}टे\appbrk{}नापः}\ A}}\skt{ }\edtext{\skt{पिवेत्~। काल\-हि\-त\-मि\-त\-स्नि\-ग्ध\-म\-धु\-र\-प्रा\-य\-मा\-हारं}}{\lemma{\skt{पिवेत्~। काल\appbrk{}हि\appbrk{}त\appbrk{}मि\appbrk{}त\appbrk{}स्नि\appbrk{}ग्ध\appbrk{}म\appbrk{}धु\appbrk{}र\appbrk{}प्रा\appbrk{}य\appbrk{}मा\appbrk{}हारं}}\Afootnote{\skt{पिबेत्, काले हित\appbrk{}मि\appbrk{}त\appbrk{}स्नि\appbrk{}ग्ध\appbrk{}म\appbrk{}धु\appbrk{}र\appbrk{}प्रा\appbrk{}य\appbrk{}मा\appbrk{}हारं}\ A; \skt{पिवेत्~। काल\appbrk{}हि\appbrk{}त\appbrk{}मि\appbrk{}त\appbrk{}स्नि\appbrk{}ग्ध\appbrk{}म\appbrk{}धु\appbrk{}र\appbrk{}प्रा\appbrk{}य\appbrk{}मारं}\ K\textsuperscript{ac}; \skt{पिवेत्~। काल\appbrk{}हि\appbrk{}त\appbrk{}मि\appbrk{}त\appbrk{}स्नि\appbrk{}ग्ध\appbrk{}म\appbrk{}धु\appbrk{}र\appbrk{}प्रा\appbrk{}य\appbrk{}मा\appbrk{}हारं}\ K\textsuperscript{pc}; \skt{पिवेत्~। काल\appbrk{}हि\appbrk{}त\appbrk{}मि\appbrk{}त\appbrk{}स्नि\appbrk{}ग्ध\appbrk{}म\appbrk{}धु\appbrk{}र\appbrk{}प्रा\appbrk{}य\appbrk{}मा\appbrk{}हार}\ H}}\skt{ }\edtext{\edtext{\skt{वैद्य\-प्र\-त्य\-वे\-क्षि\-त\-म\-श्नी\-यात्}}{\lemma{\skt{वैद्य\appbrk{}प्र\appbrk{}त्य\appbrk{}वे\appbrk{}क्षि\appbrk{}तम्}}\Afootnote{A K\textsuperscript{pc} H; \skt{वैद्य\appbrk{}प्र\appbrk{}त्य\appbrk{}वे\appbrk{}क्षिम्}\ K\textsuperscript{ac}}}\skt{~। ग्राम\-ग\-ण\-ग\-णि\-का\-प\-ति\-त\-श\-त्रु\-भो\-ज\-नानि}}{\lemma{\skt{अश्नीयात्~। ग्राम\appbrk{}ग\appbrk{}ण\appbrk{}ग\appbrk{}णि\appbrk{}का\appbrk{}प\appbrk{}ति\appbrk{}त\appbrk{}श\appbrk{}त्रु\appbrk{}भो\appbrk{}ज\appbrk{}नानि}}\Afootnote{K H; \skt{अश्नीयात्, ग्राम\appbrk{}ग\appbrk{}ण\appbrk{}ग\appbrk{}णि\appbrk{}का\appbrk{}प\appbrk{}णि\appbrk{}क\appbrk{}श\appbrk{}त्रु\appbrk{}स\appbrk{}त्र\appbrk{}श\appbrk{}ठ\appbrk{}प\appbrk{}ति\appbrk{}त\appbrk{}भो\appbrk{}ज\appbrk{}नानि}\ A}}\skt{ }\edtext{\skt{परिहरेत्~। शेषाण्यपि}}{\lemma{\skt{परिहरेत्~। शेषाण्य्}}\Afootnote{K H; \skt{परिकरेत्, शेषाण्य्}\ A}}\skt{ }\edtext{\skt{चानि\-ष्ट\-रू\-प\-र\-स\-ग\-न्ध\-श\-ब्द\-स्प\-र्श\-मा\-न\-सानि}}{\lemma{\skt{चानि\appbrk{}ष्ट\appbrk{}रू\appbrk{}प\appbrk{}र\appbrk{}स\appbrk{}ग\appbrk{}न्ध\appbrk{}श\appbrk{}ब्द\appbrk{}स्प\appbrk{}र्श\appbrk{}मा\appbrk{}न\appbrk{}सानि}}\Afootnote{K; \skt{चानि\appbrk{}ष्ट\appbrk{}रू\appbrk{}प\appbrk{}र\appbrk{}स\appbrk{}ग\appbrk{}न्ध\appbrk{}स्प\appbrk{}र्श\appbrk{}श\appbrk{}ब्द\appbrk{}मा\appbrk{}न\appbrk{}सानि,}\ A; \skt{चानि\appbrk{}ष्ट\appbrk{}रू\appbrk{}प\appbrk{}र\appbrk{}स\appbrk{}ग\appbrk{}न्ध\appbrk{}श\appbrk{}व्द\appbrk{}स्प\appbrk{}र्श\appbrk{}मा\appbrk{}न\appbrk{}सानि}\ H}}\skt{ }\edtext{\skt{अन्यानेवं}}{\lemma{\skt{अन्यानेवं}}\Afootnote{K H; \skt{अन्यान्येवं}\ A}}\skt{ }\edtext{\skt{गुणानपि}}{\lemma{\skt{गुणानपि}}\Afootnote{K H; \skt{गुणान्य}\ A}}\skt{ }\edtext{\skt{चास\-म्भ्र\-म\-द\-त्तानि}}{\lemma{\skt{चास\appbrk{}म्भ्र\appbrk{}म\appbrk{}द\appbrk{}त्तानि}}\Afootnote{K\textsuperscript{pc} H; \skt{पि संभ्रमदत्तानि,}\ A; \skt{चासम्भ्रदत्तानि}\ K\textsuperscript{ac}}}\skt{ }\edtext{\edtext{\skt{तान्यपि}}{\lemma{\skt{तान्य्}}\Afootnote{K H; \skt{(ता}\ A}}\skt{ न}}{\lemma{\skt{अपि न}}\Afootnote{K H; \skt{न्यपि)}\ A}}\skt{ }\edtext{\skt{मक्षि\-का\-वा\-लो\-प\-हि\-तानि}}{\lemma{\skt{मक्षि\appbrk{}का\appbrk{}वा\appbrk{}लो\appbrk{}प\appbrk{}हि\appbrk{}तानि}}\Afootnote{K H; \skt{मक्षि\appbrk{}का\appbrk{}बा\appbrk{}लो\appbrk{}प\appbrk{}ह\appbrk{}तानि,}\ A}}\skt{~। }\edtext{\skt{नाप्र\-क्षा\-लि\-त\-पा\-णि\-पादौ}}{\lemma{\skt{नाप्र\appbrk{}क्षा\appbrk{}लि\appbrk{}त\appbrk{}पा\appbrk{}णि\appbrk{}पादौ}}\Afootnote{K H; \skt{नाप्र\appbrk{}क्षा\appbrk{}लि\appbrk{}त\appbrk{}पादो}\ A}}\skt{ }\edtext{\skt{भुञ्जीत}}{\lemma{\skt{भुञ्जीत}}\Afootnote{K H; \skt{भुञ्जीत न}\ A}}\skt{~। }\edtext{\skt{मूत्रो\-च्चा\-र\-पी\-डितो वा~। न}}{\lemma{\skt{मूत्रो\appbrk{}च्चा\appbrk{}र\appbrk{}पी\appbrk{}डितो वा~। न}}\Afootnote{K H; \skt{मूत्रो\appbrk{}च्चा\appbrk{}र\appbrk{}पी\appbrk{}डितो न}\ A}}\skt{ }\edtext{\skt{सन्ध्य\-यो\-र्ना\-पा\-श्रितो}}{\lemma{\skt{नापाश्रितो}}\Afootnote{K H; \skt{नानुपाश्रितो}\ A}}\skt{ }\edtext{\edtext{\skt{नाति\-का\-ल\-ही\-न\-म\-ति\-मात्रं}}{\lemma{\skt{नाति\appbrk{}का\appbrk{}ल\appbrk{}ही\appbrk{}नम्}}\Afootnote{K; \skt{नातीतकालं हीनम्}\ A; \skt{नातिकालहीन}\ H}}}{\lemma{\skt{अतिमात्रं}}\Afootnote{A K; \skt{मतिमात्रं}\ H}}\skt{ }\edtext{\skt{वेति}}{\lemma{\skt{वेति}}\Afootnote{K H; \skt{(नोद्धृतस्नेहं) चेति~॥}\ A}}\skt{~॥}\pend

\pstart\versepar \ledsidenote{\passref{4.24.102}}\edtext{\skt{यस्मिन्यस्मिनृतौ ये ये}}{\lemma{\skt{यस्मिन्यस्मिनृतौ ये ये}}\Afootnote{K H; \skt{भवन्ति चात्र यस्मिन्यस्मिन्नृतौ ये ये}\ A}}\skt{ दोषाः कुप्यन्ति देहिनां~। }\pend
\pstart\versepar \skt{तेषु तेषु प्रदातव्या }\edtext{\skt{रसास्ते}}{\lemma{\skt{रसास्ते}}\Afootnote{K; \skt{रसास्ते ते}\ A; \skt{रसास्ते ते}\ H}}\skt{ विजानता~।}\pend

\pstart \ledsidenote{\passref{4.24.103}}\edtext{\skt{वर्षासु}}{\lemma{\skt{वर्षासु}}\Afootnote{K H; \skt{वर्षासु न पिबेत्तोयं पिबेच्छरदि मात्रया~। वर्षसु}\ A}}\skt{ चतुरो }\edtext{\skt{मासां}}{\lemma{\skt{मासां}}\Afootnote{K; \skt{मासा}\ A H}}\skt{ }\edtext{\skt{मात्रावदुकं}}{\lemma{\skt{मात्रावदुकं}}\Afootnote{K; \skt{न्मात्रावदुदकं}\ A; \skt{न्मात्रा\appbrk{}व\appbrk{}दु\appbrk{}द\appbrk{}कम्}\ H}}\skt{ }\edtext{\skt{पिवेत्}}{\lemma{\skt{पिवेत्}}\Afootnote{K H; \skt{पिबेत्~॥}\ A}}\skt{~।}\pend

\pstart\versepar \ledsidenote{\passref{4.24.104}}\skt{उष्णं }\edtext{\skt{हिमवसन्ते}}{\lemma{\skt{हिमवसन्ते}}\Afootnote{K; \skt{हैमे वसन्ते}\ A}}\skt{ }\edtext{\skt{चकासं}}{\lemma{\skt{चकासं}}\Afootnote{K\textsuperscript{pc}; \skt{च कामं}\ A; \skt{च}\ K\textsuperscript{ac}}}\skt{ }\edtext{\skt{ग्रीष्मे च}}{\lemma{\skt{ग्रीष्मे च}}\Afootnote{K\textsuperscript{pc}; \skt{श्रीष्मे तु}\ A; \emph{om.}\ K\textsuperscript{ac}}}\skt{ }\edtext{\skt{शीतलं}}{\lemma{\skt{शीतलं}}\Afootnote{A K\textsuperscript{pc}; \emph{om.}\ K\textsuperscript{ac}}}\skt{~। }\pend
\pstart\versepar \skt{}\edtext{\skt{हेमन्ते च}}{\lemma{\skt{हेमन्ते च}}\Afootnote{A K\textsuperscript{pc}; \emph{om.}\ K\textsuperscript{ac}; \skt{हिमं}\ H}}\skt{ }\edtext{\skt{वसन्ते च}}{\lemma{\skt{वसन्ते च}}\Afootnote{A K\textsuperscript{pc} H; \emph{om.}\ K\textsuperscript{ac}}}\skt{ }\edtext{\skt{सीध्वरिष्टौ}}{\lemma{\skt{सीध्वरिष्टौ}}\Afootnote{A K\textsuperscript{pc} H; \emph{om.}\ K\textsuperscript{ac}}}\skt{ }\edtext{\edtext{\skt{पिवेन्नरः}}{\lemma{\skt{पिवेन्}}\Afootnote{K\textsuperscript{pc} H; \skt{पिबेन्}\ A; \emph{om.}\ K\textsuperscript{ac}}}}{\lemma{\skt{नरः}}\Afootnote{A K\textsuperscript{pc} H; \emph{om.}\ K\textsuperscript{ac}}}\skt{~।}\edtext{}{\lemma{\skt{हिमवसन्ते… हे}}\Afootnote{\emph{om.}\ H}}\pend

\pstart\versepar \ledsidenote{\passref{4.24.105}}\edtext{\skt{श्रितशीतं}}{\lemma{\skt{श्रितशीतं}}\Afootnote{K H; \skt{शृतशीतं}\ A}}\skt{ }\edtext{\skt{पयो}}{\lemma{\skt{पयो}}\Afootnote{A K; \skt{पयो}\ H}}\skt{ }\edtext{\skt{ग्रीष्मे}}{\lemma{\skt{ग्रीष्मे}}\Afootnote{A K; \skt{ग्रिष्मे}\ H}}\skt{ प्रावृट्काले रसं पिवेत्~। }\pend
\pstart\versepar \skt{यूषं वर्षति तस्यान्ते }\edtext{\skt{प्रपिवेच्छीतलं}}{\lemma{\skt{प्रपिवेच्छीतलं}}\Afootnote{K H; \skt{प्रपिच्छीतलं}\ A}}\skt{ जलं~।}\edtext{}{\lemma{\skt{पिवेत्~। यूषं}}\Afootnote{K H; \skt{पिबेत्~। यूषं}\ A}}\pend

\pstart\versepar \ledsidenote{\passref{4.24.106}}\skt{स्वस्थ }\edtext{\skt{एव}}{\lemma{\skt{एव}}\Afootnote{K H; \skt{एवम्}\ A}}\skt{ }\edtext{\skt{मतो}}{\lemma{\skt{मतो}}\Afootnote{K H; \skt{अतो}\ A}}\skt{ }\edtext{\skt{न्यस्तु}}{\lemma{\skt{न्यस्तु}}\Afootnote{K H; \skt{ऽन्यस्तु}\ A}}\skt{ }\edtext{\skt{दोषा\-हा\-ग\-दा\-तुराः}}{\lemma{\skt{दोषा\appbrk{}हा\appbrk{}ग\appbrk{}दा\appbrk{}तुराः}}\Afootnote{K; \skt{दोषा\appbrk{}हा\appbrk{}र\appbrk{}ग\appbrk{}ता\appbrk{}नुगः~।}\ A; \skt{दोषा\appbrk{}हा\appbrk{}र\appbrk{}ग\appbrk{}दा\appbrk{}तुराः~।}\ H\textsuperscript{ac}; \skt{दोषा\appbrk{}हा\appbrk{}र\appbrk{}ग\appbrk{}ता\appbrk{}नु\appbrk{}गः\appbrk{}दा\appbrk{}तुराः~।}\ H\textsuperscript{pc}}}\skt{~। }\pend
\pstart\versepar \skt{स्नेहं }\edtext{\skt{पिप्पलिचूर्ण्णेन}}{\lemma{\skt{पिप्पलिचूर्ण्णेन}}\Afootnote{K H; \skt{सैन्धवचूर्णेन}\ A}}\skt{ }\edtext{\skt{सैन्धवेन च}}{\lemma{\skt{सैन्धवेन च}}\Afootnote{K H; \skt{पिप्पलीभिश्च}\ A}}\skt{ संयुतं~।}\pend

\pstart\versepar \ledsidenote{\passref{4.24.107}}\edtext{\skt{पिवे\-द\-ग्नि\-वि\-वृ\-ध्यर्थं न च}}{\lemma{\skt{पिवे\appbrk{}द\appbrk{}ग्नि\appbrk{}वि\appbrk{}वृ\appbrk{}ध्यर्थं न च}}\Afootnote{K H; \skt{पिबे\appbrk{}द\appbrk{}ग्नि\appbrk{}वि\appbrk{}वृ\appbrk{}द्ध्यर्थं न च}\ A}}\skt{ }\edtext{\skt{वेगां}}{\lemma{\skt{वेगां}}\Afootnote{K; \skt{वेगा}\ A H}}\skt{ विधारयेत्~। }\pend
\pstart\versepar \skt{}\edtext{\skt{अग्निदीप्तिरलं}}{\lemma{\skt{अलं}}\Afootnote{K H; \skt{करं}\ A}}\skt{ }\edtext{\skt{नृणां}}{\lemma{\skt{नृणां}}\Afootnote{K H; \skt{नॄणां}\ A}}\skt{ }\edtext{\skt{रोगानां}}{\lemma{\skt{रोगानां}}\Afootnote{K; \skt{रोगाणां}\ A H}}\skt{ }\edtext{\skt{स्रमनं}}{\lemma{\skt{स्रमनं}}\Afootnote{K H; \skt{शमनं}\ A}}\skt{ प्रति~।}\edtext{}{\lemma{\skt{विधारयेत्~। अग्निदीप्तिर्}}\Afootnote{K; \skt{न्विधारयेत्~। अग्निदीप्ति}\ A; \skt{न्विधारयेत्~। अग्निदीप्तिर्}\ H}}\pend

\pstart\versepar \ledsidenote{\passref{4.24.108}}\edtext{\skt{प्रावृ\-ट्छ\-र\-द्व\-स\-न्तेषु}}{\lemma{\skt{प्रावृ\appbrk{}ट्छ\appbrk{}र\appbrk{}द्व\appbrk{}स\appbrk{}न्तेषु}}\Afootnote{K; \skt{प्रावृ\appbrk{}ट्श\appbrk{}र\appbrk{}द्व\appbrk{}स\appbrk{}न्तेषु}\ A; \skt{प्रावृ\appbrk{}च्छ\appbrk{}र\appbrk{}द्व\appbrk{}स\appbrk{}न्तेषु}\ H}}\skt{ }\edtext{\skt{सम्य\-क्स्ने\-हा\-दि\-मा\-च\-रेत्}}{\lemma{\skt{सम्यक्स्नेहादिम्}}\Afootnote{K H; \skt{सम्यक्स्नेहादि}\ A}}\skt{~। }\pend
\pstart\versepar \skt{कफे प्रच्छर्दनं पित्ते }\edtext{\skt{विरो}}{\lemma{\skt{विरो}}\Afootnote{K; \skt{विरेको}\ A H}}\skt{ }\edtext{\edtext{\skt{वस्तिरीरणे}}{\lemma{\skt{वस्तिर्}}\Afootnote{K H; \skt{बस्ति}\ A}}}{\lemma{\skt{ईरणे}}\Afootnote{K H; \skt{रीरने~॥}\ A}}\skt{~।}\edtext{}{\lemma{\skt{आचरेत्~। कफे}}\Afootnote{K H; \skt{माचरेत्~। कफे}\ A}}\pend

\pstart \ledsidenote{\passref{4.24.109}}\skt{शस्यते }\edtext{\skt{त्रिष्वपिहितो}}{\lemma{\skt{त्रिष्वपिहितो}}\Afootnote{K H; \skt{त्रिष्वपि सदा}\ A}}\skt{ व्यायामो }\edtext{\skt{दोषनाशनः}}{\lemma{\skt{दोषनाशनः}}\Afootnote{K H; \skt{दोषनाशनः~। भुक्तं विरुद्धमप्यन्नं व्यायामान्न प्रदुष्यति~॥}\ A}}\pend

\pstart\versepar \ledsidenote{\passref{4.24.110}}\edtext{\skt{उत्स\-र्ग\-मै\-थु\-ना\-हा\-र\-शो\-धने}}{\lemma{\skt{उत्स\appbrk{}र्ग\appbrk{}मै\appbrk{}थु\appbrk{}ना\appbrk{}हा\appbrk{}र\appbrk{}शो\appbrk{}धने}}\Afootnote{K H; \skt{उत्स\appbrk{}र्ग\appbrk{}मै\appbrk{}थु\appbrk{}न\appbrk{}हा\appbrk{}र\appbrk{}शो\appbrk{}धने}\ A}}\skt{ स्यात्तु तन्मना~। }\pend
\pstart\versepar \skt{}\edtext{\skt{नेच्छेत्प्रतिभयां}}{\lemma{\skt{नेच्छेत्प्रतिभयां}}\Afootnote{K H; \skt{नेच्छेद्दोषचया}\ A}}\skt{ }\edtext{\skt{प्राप्तुं}}{\lemma{\skt{प्राप्तुं}}\Afootnote{K H; \skt{त्प्राज्ञः}\ A}}\skt{ }\edtext{\skt{पीडा\-शा\-री\-र\-मा\-नसीं}}{\lemma{\skt{पीडा\appbrk{}शा\appbrk{}री\appbrk{}र\appbrk{}मा\appbrk{}नसीं}}\Afootnote{K H; \skt{पीडां वा कायमानसीम्~॥}\ A}}\skt{~।}\pend

\pstart\versepar \ledsidenote{\passref{4.24.111}}\skt{अति\-स्त्री\-सं\-प्र\-यो\-गाच्च }\edtext{\skt{रक्षे\-दा\-त्मा\-न\-मा\-त्म\-वान्}}{\lemma{\skt{रक्षेदात्मानम्}}\Afootnote{A H; \skt{रक्षेदात्मानम्}\ K}}\skt{~। }\pend
\pstart\versepar \skt{शोषकारार्शांसि }\edtext{\skt{श्वास\-का\-श्र्या\-ति\-पा\-ण्डुता}}{\lemma{\skt{श्वास\appbrk{}का\appbrk{}श्र्या\appbrk{}ति\appbrk{}पा\appbrk{}ण्डुता}}\Afootnote{K; \skt{श्वास\appbrk{}का\appbrk{}र्श्य\appbrk{}पा\appbrk{}ण्ड्वा\appbrk{}म\appbrk{}य\appbrk{}क्षयाः~॥}\ A; \skt{श्वासकास्याति पाण्डुता~।}\ H}}\skt{~।}\edtext{}{\lemma{\skt{आत्मवान्~। शोषकारार्शांसि}}\Afootnote{\skt{आत्मवान्~। शूलकासज्वर}\ A; \skt{आत्म\appbrk{}वा\appbrk{}न्शो\appbrk{}ष\appbrk{}का\appbrk{}स\appbrk{}ज्व\appbrk{}रा\appbrk{}र्शांसि}\ K H}}\pend

\pstart\versepar \ledsidenote{\passref{4.24.112}}\edtext{\skt{अति\-व्य\-वा\-या\-ज्जा\-यन्ते}}{\lemma{\skt{अति\appbrk{}व्य\appbrk{}वा\appbrk{}या\appbrk{}ज्जा\appbrk{}यन्ते}}\Afootnote{K; \skt{अति\appbrk{}व्य\appbrk{}वा\appbrk{}जा\appbrk{}ज्जा\appbrk{}यन्ते}\ A; \skt{अति\appbrk{}व्य\appbrk{}वा\appbrk{}या\appbrk{}ज्जा\appbrk{}यन्ते}\ H}}\skt{ }\edtext{\skt{रोगा\-श्चा\-क्षे\-प\-गा\-दयः}}{\lemma{\skt{चाक्षेपगादयः}}\Afootnote{K; \skt{चाक्षेपकादयः~।}\ A H}}\skt{~। }\pend
\pstart\versepar \skt{आयुष्मन्तो }\edtext{\skt{मन्द\-ज\-रा\-व\-पु\-र्व\-र्ण्ण\-व\-ला\-धिकाः}}{\lemma{\skt{मन्द\appbrk{}ज\appbrk{}रा\appbrk{}व\appbrk{}पु\appbrk{}र्व\appbrk{}र्ण्ण\appbrk{}व\appbrk{}ला\appbrk{}धिकाः}}\Afootnote{K; \skt{मन्दजरा वपु\appbrk{}र्व\appbrk{}र्ण\appbrk{}ब\appbrk{}ला\appbrk{}न्विताः~॥}\ A; \skt{मन्दजरा आयु\appbrk{}र्व्व\appbrk{}र्ण्ण\appbrk{}व\appbrk{}ला\appbrk{}धिकाः~।}\ H}}\skt{~।}\pend

\pstart\versepar \ledsidenote{\passref{4.24.113}}\skt{स्थिरो\-प\-चि\-त\-मां\-साश्च भवन्ति स्त्रीषु }\edtext{\skt{संयताः}}{\lemma{\skt{संयताः}}\Afootnote{A K; \skt{संस्थिताः~॥}\ H}}\skt{~। }\pend
\pstart\versepar \skt{}\edtext{\edtext{\edtext{\skt{त्रिभि\-स्त्रि\-भि\-र\-हो\-भिस्तु}}{\lemma{\skt{त्रिभिस्}}\Afootnote{A H; \skt{त्रिभि}\ K}}}{\lemma{\skt{त्रिभिर्}}\Afootnote{A H; \skt{स्त्रिभि}\ K}}}{\lemma{\skt{अहोभिस्तु}}\Afootnote{H; \skt{अहोभिर्वा}\ A; \skt{रहोभिस्तु}\ K}}\skt{ }\edtext{\skt{सेवेत}}{\lemma{\skt{सेवेत}}\Afootnote{K H; \skt{समीया}\ A}}\skt{ }\edtext{\skt{प्रममदान्नरः}}{\lemma{\skt{प्रममदान्नरः}}\Afootnote{K\textsuperscript{pc}; \skt{त्प्रमदां नरः~॥}\ A; \skt{प्रमदान्नरः~।}\ K\textsuperscript{ac} H}}\skt{~।}\pend

\pstart\versepar \ledsidenote{\passref{4.24.114}}\skt{सर्वेष्वृतुषु }\edtext{\skt{ग्रीष्मे तु}}{\lemma{\skt{ग्रीष्मे तु}}\Afootnote{K H; \skt{घर्मेषु}\ A}}\skt{ }\edtext{\skt{पक्षा\-त्प\-क्षा\-द्व्र\-जे\-द्बुधः}}{\lemma{\skt{पक्षा\appbrk{}त्प\appbrk{}क्षा\appbrk{}द्व्र\appbrk{}जे\appbrk{}द्बुधः}}\Afootnote{A H; \skt{पक्षा\appbrk{}त्प\appbrk{}क्षा\appbrk{}द्व्र\appbrk{}जे\appbrk{}द्वुधः~।}\ K}}\skt{~। }\pend
\pstart\versepar \skt{}\edtext{\edtext{\skt{रज\-स्व\-ला\-म\-का\-माञ्च}}{\lemma{\skt{रजस्वलाम्}}\Afootnote{A; \skt{रजस्वला}\ K; \skt{रजस्वरा}\ H}}}{\lemma{\skt{अकामाञ्च}}\Afootnote{A; \skt{मकामाञ्च}\ K H}}\skt{ }\edtext{\edtext{\skt{मलि\-ना\-म\-प्र\-जा\-मपि}}{\lemma{\skt{अप्रजाम्}}\Afootnote{K H; \skt{अप्रियां}\ A}}}{\lemma{\skt{अपि}}\Afootnote{K H; \skt{तथा~॥}\ A}}\skt{~।}\pend

\pstart\versepar \ledsidenote{\passref{4.24.115}}\edtext{\skt{वर्ण्ण\-वृ\-द्धा\-म्व\-यो\-वृद्धं}}{\lemma{\skt{वयोवृद्धं}}\Afootnote{K; \skt{वयोवृद्धां}\ A H}}\skt{ तथा व्याधिप्रपीडितां~। }\pend
\pstart\versepar \skt{हीनाङ्गीं }\edtext{\skt{गर्भिणीं}}{\lemma{\skt{गर्भिणीं}}\Afootnote{A K; \skt{गर्व्भिणीं}\ H}}\skt{ द्वेष्यां योनि\-दो\-ष\-स\-म\-न्वितां~।}\pend

\pstart\versepar \ledsidenote{\passref{4.24.116}}\skt{सगोत्रां }\edtext{\skt{गुरुपत्नींञ्च}}{\lemma{\skt{गुरुपत्नींञ्च}}\Afootnote{K; \skt{गुरुपत्नीं च}\ A H}}\skt{ तथा प्रव्रजितामपि~। }\pend
\pstart\versepar \skt{}\edtext{\skt{सन्ध्ययोः}}{\lemma{\skt{सन्ध्ययोः}}\Afootnote{K H; \skt{सन्ध्याप}\ A}}\skt{ }\edtext{\skt{सर्वकाले च}}{\lemma{\skt{सर्वकाले च}}\Afootnote{K H; \skt{र्वस्वगम्यां च}\ A}}\skt{ नोपे\-या\-त्प्र\-म\-दा\-न्नरः~।}\pend

\pstart\versepar \ledsidenote{\passref{4.24.117}}\skt{गोसर्गे चार्धरात्रे च तथा मध्यन्दिने }\edtext{\skt{पि च}}{\lemma{\skt{पि च}}\Afootnote{K; \skt{षु च~।}\ A; \skt{ऽपि च~॥}\ H}}\skt{~। }\pend
\pstart\versepar \skt{लज्जासमावहे देशे विवृते }\edtext{\skt{शुद्ध}}{\lemma{\skt{शुद्ध}}\Afootnote{K H; \skt{ऽशुद्ध}\ A}}\skt{ एव च~।}\pend

\pstart\versepar \ledsidenote{\passref{4.24.118}}\skt{क्षुधितो व्याधितश्चैव }\edtext{\skt{क्षुब्धचित्तश्च}}{\lemma{\skt{क्षुब्धचित्तश्च}}\Afootnote{A; \skt{क्षुव्धचित्तश्च}\ K H}}\skt{ मानवः~। }\pend
\pstart\versepar \skt{}\edtext{\skt{प्राप्त\-मू\-त्र\-पु\-री\-शश्च}}{\lemma{\skt{प्राप्त\appbrk{}मू\appbrk{}त्र\appbrk{}पु\appbrk{}री\appbrk{}शश्च}}\Afootnote{K; \skt{वात\appbrk{}वि\appbrk{}ण्मू\appbrk{}त्र\appbrk{}वेगी च}\ A; \skt{प्राप्त\appbrk{}मू\appbrk{}त्र\appbrk{}पु\appbrk{}री\appbrk{}षश्च}\ H}}\skt{ }\edtext{\skt{पिपासी }\edtext{\skt{दुर्बलस्तथा}}{\lemma{\skt{तथा}}\Afootnote{K H; \emph{om.}\ A}}}{\lemma{\skt{पिपासी दुर्बलस्}}\Afootnote{\skt{पिपा\appbrk{}स\appbrk{}र\appbrk{}ति\appbrk{}दु\appbrk{}र्बलः~॥}\ A; \skt{पिपा\appbrk{}सी\appbrk{}दु\appbrk{}र्व\appbrk{}लस्}\ K; \skt{पिपा\appbrk{}सी\appbrk{}दू\appbrk{}ऊ\appbrk{}र्व्व\appbrk{}लस्}\ H}}\skt{~।}\pend

\pstart\versepar \ledsidenote{\passref{4.24.119}}\edtext{\edtext{\skt{तिर्य\-ग्यो\-ना\-व\-योनौ}}{\lemma{\skt{तिर्यग्योनाव्}}\Afootnote{A; \skt{तिर्यग्योना}\ K H}}\skt{ च}}{\lemma{\skt{अयोनौ च}}\Afootnote{A; \skt{वयोनौ च}\ K H}}\skt{ प्राप्त\-शु\-क्र\-वि\-धा\-रणं~। }\pend
\pstart\versepar \skt{}\edtext{\skt{स्गश्च}}{\lemma{\skt{स्गश्च}}\Afootnote{\skt{दुष्टयोनौ विसर्गं तु}\ A; \skt{दुष्टयोनौ विसर्गश्च}\ K; \skt{दुष्टयोनौ विसर्ग्गञ्च}\ H}}\skt{ }\edtext{\skt{बलवानपि}}{\lemma{\skt{बलवानपि}}\Afootnote{A K; \skt{वलवानपि}\ H}}\skt{ वर्जयेत्~।}\pend

\pstart \ledsidenote{\passref{4.24.120}}\edtext{\skt{स्थिता\-नु\-त्ता\-न\-श\-यने}}{\lemma{\skt{स्थिता\appbrk{}नु\appbrk{}त्ता\appbrk{}न\appbrk{}श\appbrk{}यने}}\Afootnote{K H; \skt{रेत\appbrk{}स\appbrk{}श्चा\appbrk{}ति\appbrk{}मात्रं तु मूर्धा\appbrk{}व\appbrk{}र\appbrk{}ण\appbrk{}मेव च~। स्थिता\appbrk{}वु\appbrk{}त्ता\appbrk{}न\appbrk{}श\appbrk{}यने}\ A}}\skt{ }\edtext{\skt{विशेषेण तु}}{\lemma{\skt{विशेषेण तु}}\Afootnote{K H; \skt{विशेषेणैव}\ A}}\skt{ गर्हितं~।}\pend

\pstart\versepar \ledsidenote{\passref{4.24.121}}\edtext{\skt{क्रीडयान्यश्च}}{\lemma{\skt{क्रीडयान्यश्च}}\Afootnote{K H; \skt{क्रीडायामपि}\ A}}\skt{ मेधावी }\edtext{\skt{हितार्थं}}{\lemma{\skt{हितार्थं}}\Afootnote{K H; \skt{हितार्थी}\ A}}\skt{ परिवर्जयेत्~। }\pend
\pstart\versepar \skt{रजस्वलां }\edtext{\skt{प्रारस्या नियतात्मनः}}{\lemma{\skt{प्रारस्या नियतात्मनः}}\Afootnote{\skt{प्राप्तवतो नर\appbrk{}स्या\appbrk{}नि\appbrk{}य\appbrk{}ता\appbrk{}त्मनः}\ A H K}}\skt{~।}\pend

\pstart\versepar \ledsidenote{\passref{4.24.122}}\skt{दृष्ट्यायुस्तेजसां }\edtext{\skt{हानिरधर्मश्च}}{\lemma{\skt{हानिरधर्मश्च}}\Afootnote{\skt{हानिरधर्मश्च}\ A K H}}\skt{ ततो भवेत्~। }\pend
\pstart\versepar \skt{लिङ्गिनीं गुरुपत्नीञ्च सगोत्रामथ पर्वसु~।}\pend

\pstart\versepar \ledsidenote{\passref{4.24.123}}\edtext{\skt{वृद्धां}}{\lemma{\skt{वृद्धां}}\Afootnote{K; \skt{वृद्धां च}\ A H}}\skt{ सन्ध्ययोश्चापि }\edtext{\skt{गर्छतो}}{\lemma{\skt{गर्छतो}}\Afootnote{K; \skt{गच्छतो}\ A H}}\skt{ जीवितक्षयः~। }\pend
\pstart\versepar \skt{}\edtext{\skt{गर्भिण्यां}}{\lemma{\skt{गर्भिण्यां}}\Afootnote{K; \skt{गर्भिण्या}\ A; \skt{गर्व्भिण्यां}\ H}}\skt{ }\edtext{\skt{गर्भ\-पी\-डा\-स्या\-द्व्या\-धि\-ताञ्च}}{\lemma{\skt{गर्भ\appbrk{}पी\appbrk{}डा\appbrk{}स्या\appbrk{}द्व्या\appbrk{}धि\appbrk{}ताञ्च}}\Afootnote{K; \skt{गर्भपीडा स्याद्व्याधितायां}\ A; \skt{गर्व्भीडा स्याद्व्याधिताञ्च}\ H}}\skt{ }\edtext{\skt{बलक्षयः}}{\lemma{\skt{बलक्षयः}}\Afootnote{A; \skt{वलक्षयः}\ K H}}\skt{~।}\pend

\pstart\versepar \ledsidenote{\passref{4.24.124}}\skt{हीनांगीं मलिनां द्वेष्यां }\edtext{\skt{कामां}}{\lemma{\skt{कामां}}\Afootnote{K H; \skt{कामं}\ A}}\skt{ }\edtext{\edtext{\skt{वन्ध्यामसंवृते}}{\lemma{\skt{वन्ध्याम्}}\Afootnote{K H; \skt{वन्ध्या}\ A}}}{\lemma{\skt{असंवृते}}\Afootnote{K H; \skt{मसंवृते~।}\ A}}\skt{~। }\pend
\pstart\versepar \skt{देशे }\edtext{\skt{शुद्धे च}}{\lemma{\skt{शुद्धे च}}\Afootnote{K H; \skt{ऽशुद्धे च}\ A}}\skt{ शुक्रस्य मनसश्च क्षयो भवेत्~।}\pend

\pstart\versepar \ledsidenote{\passref{4.24.125}}\skt{क्षुधितः }\edtext{\skt{क्षुव्धचित्तश्च}}{\lemma{\skt{क्षुव्धचित्तश्च}}\Afootnote{K H; \skt{क्षुब्धचित्तश्च}\ A}}\skt{ मध्याह्ने }\edtext{\skt{तृषितोबलः}}{\lemma{\skt{तृषितोबलः}}\Afootnote{\skt{तृषितोऽबलः~।}\ A; \skt{तृषितोवलः~।}\ K H}}\skt{~। }\pend
\pstart\versepar \skt{}\edtext{\skt{स्थितायं}}{\lemma{\skt{स्थितायं}}\Afootnote{K; \skt{स्थितश्च}\ A; \skt{स्थितायां}\ H}}\skt{ }\edtext{\skt{हानिमाप्नोति}}{\lemma{\skt{आप्नोति}}\Afootnote{K H; \emph{om.}\ A}}\skt{ }\edtext{\skt{शुक्रं}}{\lemma{\skt{शुक्रं}}\Afootnote{K H; \skt{शुक्रस्य}\ A}}\skt{ }\edtext{\skt{वायुश्च}}{\lemma{\skt{वायुश्च}}\Afootnote{K H; \skt{वायोः कोपं च}\ A}}\skt{ }\edtext{\skt{कुप्यति}}{\lemma{\skt{कुप्यति}}\Afootnote{K H; \skt{विन्दति~॥}\ A}}\skt{~।}\pend

\pstart \ledsidenote{\passref{4.24.126}}\edtext{\skt{व्याधितायां}}{\lemma{\skt{व्याधितायां}}\Afootnote{K H; \skt{अति\appbrk{}प्र\appbrk{}स\appbrk{}ङ्गा\appbrk{}द्भ\appbrk{}वति शोषः शुक्रक्षयावहः~। व्याधितस्य रुजा}\ A}}\skt{ }\edtext{\skt{प्लीहामयो}}{\lemma{\skt{प्लीहामयो}}\Afootnote{K H; \skt{प्लीह्नि मृत्युर्}\ A}}\skt{ मूर्च्छा }\edtext{\skt{दोषश्च}}{\lemma{\skt{दोषश्च}}\Afootnote{K H; \skt{च}\ A}}\skt{ जायते~।}\pend

\pstart\versepar \ledsidenote{\passref{4.24.127}}\edtext{\skt{प्रत्यूषे}}{\lemma{\skt{प्रत्यूषे}}\Afootnote{K H; \skt{प्रयूष}\ A}}\skt{ }\edtext{\skt{चार्धरात्रे च}}{\lemma{\skt{चार्धरात्रे च}}\Afootnote{K H; \skt{स्यर्धरात्रे च}\ A}}\skt{ }\edtext{\skt{वात\-पि\-त्त\-प्र\-कु\-प्यतः}}{\lemma{\skt{वात\appbrk{}पि\appbrk{}त्त\appbrk{}प्र\appbrk{}कु\appbrk{}प्यतः}}\Afootnote{K H; \skt{वातपित्ते प्रकुप्यतः~।}\ A}}\skt{~। }\pend
\pstart\versepar \skt{}\edtext{\skt{तिर्य\-ग्यो\-ना\-व\-योनौ च}}{\lemma{\skt{अयोनौ च}}\Afootnote{K H; \skt{अयनौ च}\ A}}\skt{ दुष्टयोनौ तथैव च~।}\pend

\pstart\versepar \ledsidenote{\passref{4.24.128}}\edtext{\skt{उपदंशस्तथा}}{\lemma{\skt{उपदंशस्}}\Afootnote{A K\textsuperscript{pc} H; \skt{उपशस्}\ K\textsuperscript{ac}}}\skt{ वायोः कोपः }\edtext{\skt{स्याद्रेतसः}}{\lemma{\skt{स्याद्रेतसः}}\Afootnote{K H; \skt{शुक्रस्य च}\ A}}\skt{ क्षयः~। }\pend
\pstart\versepar \skt{उच्चारिते मूत्रिते च रेतसश्च विधारणे~।}\pend

\pstart\versepar \ledsidenote{\passref{4.24.129}}\skt{उत्ताने च भवेच्छीघ्रं }\edtext{\skt{शुक्रार्श्मर्यास्तु}}{\lemma{\skt{शुक्रार्श्मर्यास्तु}}\Afootnote{K; \skt{शुक्राश्मर्यास्तु}\ A; \skt{शुक्रास्मर्यास्तु}\ H}}\skt{ सम्भवः~। }\pend
\pstart\versepar \skt{}\edtext{\skt{तस्मात्सर्वं}}{\lemma{\skt{तस्मात्सर्वं}}\Afootnote{K H; \skt{सर्वं}\ A}}\skt{ }\edtext{\skt{परि\-ह\-रे\-दे\-त\-ल्लो\-क\-द्व\-या\-हितं}}{\lemma{\skt{परि\appbrk{}ह\appbrk{}रे\appbrk{}दे\appbrk{}त\appbrk{}ल्लो\appbrk{}क\appbrk{}द्व\appbrk{}या\appbrk{}हितं}}\Afootnote{K; \skt{परि\appbrk{}ह\appbrk{}रे\appbrk{}त्त\appbrk{}स्मा\appbrk{}दे\appbrk{}त\appbrk{}ल्लो\appbrk{}क\appbrk{}द्व\appbrk{}येऽ\appbrk{}हि\appbrk{}तम्~॥}\ A; \skt{परि\appbrk{}ह\appbrk{}रे\appbrk{}दे\appbrk{}त\appbrk{}ल्लो\appbrk{}क\appbrk{}द्व\appbrk{}या\appbrk{}हितं~।}\ H}}\skt{~।}\pend

\pstart \ledsidenote{\passref{4.24.130}}\edtext{\edtext{\skt{वयोरूपगुणो}}{\lemma{\skt{वयोर्}}\Afootnote{K H; \skt{शुक्रं चोपस्थितं मोहान्न सन्धार्यं कथंचन~। वयोर्}\ A}}\skt{ }\edtext{\skt{प्येतां}}{\lemma{\skt{एतां}}\Afootnote{K; \skt{पेतां}\ A H}}}{\lemma{\skt{ऊपगुणो प्य्}}\Afootnote{K; \skt{ऊपगुणो}\ A H}}\skt{ }\edtext{\skt{कुल्य कुलान्वितां}}{\lemma{\skt{कुल्य कुलान्वितां}}\Afootnote{\skt{तुल्यशीलां कुलान्विताम्~॥}\ A; \skt{कुल्यशीलां कुलान्वितां}\ K H}}\skt{~।}\pend

\pstart\versepar \ledsidenote{\passref{4.24.131}}\edtext{\skt{अभि\-का\-मो\-प्ति\-का\-मांश्च}}{\lemma{\skt{अभि\appbrk{}का\appbrk{}मो\appbrk{}प्ति\appbrk{}का\appbrk{}मांश्च}}\Afootnote{K H; \skt{अभि\appbrk{}का\appbrk{}मोऽ\appbrk{}भि\appbrk{}कामां तु}\ A}}\skt{ हृष्टो }\edtext{\edtext{\skt{हृष्टामलङ्कृतां}}{\lemma{\skt{हृष्टाम्}}\Afootnote{A K; \skt{हृष्टा}\ H}}}{\lemma{\skt{अलङ्कृतां}}\Afootnote{A K; \skt{मलङ्कृताः~।}\ H}}\skt{~। }\pend
\pstart\versepar \skt{सेवेत }\edtext{\skt{प्रमदान्नित्यं}}{\lemma{\skt{नित्यं}}\Afootnote{K H; \skt{युक्त्या}\ A}}\skt{ }\edtext{\skt{व्याजी\-क\-र\-ण\-पी\-त\-वान्}}{\lemma{\skt{व्याजी\appbrk{}क\appbrk{}र\appbrk{}ण\appbrk{}पी\appbrk{}त\appbrk{}वान्}}\Afootnote{K H; \skt{वाजी\appbrk{}क\appbrk{}र\appbrk{}ण\appbrk{}बृं\appbrk{}हितः~॥}\ A}}\skt{~।}\pend

\pstart\versepar \ledsidenote{\passref{4.24.132}}\edtext{\skt{शर्करोपहितां}}{\lemma{\skt{शर्करोपहितां}}\Afootnote{K H; \skt{भक्ष्याः सशर्क}\ A}}\skt{ }\edtext{\skt{भक्ष्यां}}{\lemma{\skt{भक्ष्यां}}\Afootnote{K H; \skt{राः}\ A}}\skt{ }\edtext{\skt{क्षीरञ्चापि}}{\lemma{\skt{चापि}}\Afootnote{K\textsuperscript{pc} H; \skt{ससि}\ A; \skt{चा}\ K\textsuperscript{ac}}}\skt{ सशर्करं~। }\pend
\pstart\versepar \skt{व्यजनं }\edtext{\skt{स्वप्नसेवा च}}{\lemma{\skt{स्वप्नसेवा च}}\Afootnote{K H; \skt{स्वप्नो}\ A}}\skt{ व्यवायान्ते हितानि तु~।}\edtext{}{\lemma{\skt{सशर्करं~। व्यजनं}}\Afootnote{K; \skt{तं रस एव च~। स्नानं सव्यजनं}\ A; \skt{सशर्क्करां~। व्यजनं}\ H}}\pend

\pstart\versepar \ledsidenote{\passref{4.24.133}}\edtext{\skt{सुखमात्रं}}{\lemma{\skt{सुखमात्रं}}\Afootnote{K H; \skt{मुखमात्रं}\ A}}\skt{ समासेन }\edtext{\skt{सद्वृ\-त्त\-स्ये\-त\-दी\-रितं}}{\lemma{\skt{सद्वृ\appbrk{}त्त\appbrk{}स्ये\appbrk{}त\appbrk{}दी\appbrk{}रितं}}\Afootnote{K; \skt{सद्वृ\appbrk{}त्त\appbrk{}स्यै\appbrk{}त\appbrk{}दी\appbrk{}रि\appbrk{}तम्~।}\ A; \skt{समुद्वृत्तस्य दीरितं~।}\ H\textsuperscript{ac}; \skt{सद्वृ\appbrk{}त\appbrk{}स्ये\appbrk{}त\appbrk{}मु\appbrk{}द्वृ\appbrk{}त्तस्य दीरितं~।}\ H\textsuperscript{pc}}}\skt{~। }\pend
\pstart\versepar \skt{}\edtext{\edtext{\skt{आयु\-रा\-रो\-ग्य\-मर्थो}}{\lemma{\skt{आयुरारोग्यम्}}\Afootnote{K; \skt{आरोग्यम्}\ A; \skt{आयुरारोग्यम्}\ H}}\skt{ वा}}{\lemma{\skt{अर्थो वा}}\Afootnote{K H; \skt{आयुरर्थो वा}\ A}}\skt{ }\edtext{\skt{नासद्भिः}}{\lemma{\skt{नासद्भिः}}\Afootnote{K H; \skt{गासद्भिः}\ A}}\skt{ प्राप्यते }\edtext{\edtext{\skt{त्रिभिरिति}}{\lemma{\skt{त्रिभिर्}}\Afootnote{K H; \skt{नृभिः~॥}\ A}}}{\lemma{\skt{इति}}\Afootnote{K H; \emph{om.}\ A}}\skt{~॥}\pend

\pstart \ledsidenote{\passref{4.24.trailer}}\edtext{\skt{{\symfont ❈}}}{\lemma{\skt{{\symfont ❈}}}\Afootnote{K; \skt{इति सुश्रु\appbrk{}त\appbrk{}सं\appbrk{}हि\appbrk{}तायां चिकि\appbrk{}त्सा\appbrk{}स्था\appbrk{}नेऽ\appbrk{}ना\appbrk{}ग\appbrk{}ता\appbrk{}बा\appbrk{}ध\appbrk{}चि\appbrk{}कि\appbrk{}त्सितं नाम चतुर्विंशोऽध्यायः~॥२४~॥}\ A; \skt{॥ चिकित्सा}\ H}}\pend

\endnumbering
\beginnumbering
\eledsection*{Adhyāya 25}
\pstart \ledsidenote{\passref{4.25.1}}\skt{अथातो }\edtext{\skt{मिश्रि\-क\-चि\-कि\-त्सितं}}{\lemma{\skt{मिश्रि\appbrk{}क\appbrk{}चि\appbrk{}कि\appbrk{}त्सितं}}\Afootnote{K; \skt{मिश्र\appbrk{}क\appbrk{}चि\appbrk{}कि\appbrk{}त्सितं}\ A; \skt{मिश्रि\appbrk{}क\appbrk{}म\appbrk{}ध्या\appbrk{}यम्}\ H}}\skt{ }\edtext{\skt{व्याख्यास्यामः}}{\lemma{\skt{व्याख्यास्यामः}}\Afootnote{A H; \emph{om.}\ K}}\skt{~॥}\pend

\pstart\versepar \ledsidenote{\passref{4.25.3}}\skt{पाल्यामयास्तु विस्राव्या }\edtext{\skt{यः}}{\lemma{\skt{यः}}\Afootnote{K H; \skt{इत्य्}\ A}}\skt{ }\edtext{\skt{उक्ताः}}{\lemma{\skt{उक्ताः}}\Afootnote{K; \skt{उक्तं}\ A; \skt{उक्तः}\ H}}\skt{ }\edtext{\skt{प्राङ्निबोध तां}}{\lemma{\skt{प्राङ्निबोध तां}}\Afootnote{K H; \skt{प्राङ्गिबोध ता}\ A}}\skt{~। }\pend
\pstart\versepar \skt{}\edtext{\edtext{\skt{परि\-पो\-ट\-स्त\-थो\-त्पा\-दो\-व\-मन्थो}}{\lemma{\skt{परिपोटस्}}\Afootnote{K H; \skt{न्~। परिपोटस्}\ A}}}{\lemma{\skt{तथो\appbrk{}त्पा\appbrk{}दो\appbrk{}व\appbrk{}मन्थो}}\Afootnote{K; \skt{तथोत्पात उन्मन्थो}\ A; \skt{तथो\appbrk{}त्पा\appbrk{}दोऽ\appbrk{}व\appbrk{}मन्थो}\ H}}\skt{ }\edtext{\skt{दुःखवर्धनंः}}{\lemma{\skt{दुःखवर्धनंः}}\Afootnote{K; \skt{दुःखवर्धनः~॥}\ A; \skt{दुःखवर्द्धनं~॥}\ H}}\pend

\pstart\versepar \ledsidenote{\passref{4.25.4}}\skt{पञ्चमः परिलेही च कर्ण्णपाल्यां }\edtext{\skt{गदाः}}{\lemma{\skt{गदाः}}\Afootnote{A K; \skt{गताः}\ H}}\skt{ }\edtext{\skt{स्मृताः}}{\lemma{\skt{स्मृताः}}\Afootnote{A K H\textsuperscript{ac}; \skt{सधा स्मृताः~॥}\ H\textsuperscript{pc}}}\skt{~। }\pend
\pstart\versepar \skt{}\edtext{\skt{सौकु\-मा\-र्या\-च्चि\-रो\-त्सृष्टे}}{\lemma{\skt{सौकु\appbrk{}मा\appbrk{}र्या\appbrk{}च्चि\appbrk{}रो\appbrk{}त्सृष्टे}}\Afootnote{A K; \skt{सौकु\appbrk{}मा\appbrk{}र्या\appbrk{}च्चि\appbrk{}रो\appbrk{}त्सृष्टा}\ H}}\skt{ }\edtext{\skt{सह\-सा\-भि\-प्र\-व\-र्धिते}}{\lemma{\skt{सह\appbrk{}सा\appbrk{}भि\appbrk{}प्र\appbrk{}व\appbrk{}र्धिते}}\Afootnote{K H; \skt{सह\appbrk{}साऽ\appbrk{}भि\appbrk{}प्र\appbrk{}व\appbrk{}र्धिते~॥}\ A}}\skt{~।}\pend

\pstart\versepar \ledsidenote{\passref{4.25.5}}\edtext{\skt{कर्ण्णशोफो}}{\lemma{\skt{कर्ण्णशोफो}}\Afootnote{A K; \skt{कर्ण्णे शोफो}\ H}}\skt{ भवेत्पाल्यां सरुजः परिपोटवान्~। }\pend
\pstart\versepar \skt{कृष्णा\-रु\-णि\-नि\-भ\-स्तव्धः सवा\-ता\-त्प\-रि\-पो\-टकः~।}\edtext{}{\lemma{\skt{परिपोटवान्~। कृष्णा\appbrk{}रु\appbrk{}णि\appbrk{}नि\appbrk{}भ\appbrk{}स्तव्धः}}\Afootnote{K; \skt{परिपोटवान्~। कृष्णारुणनिभः स्तब्धः}\ A; \skt{परिपोटवान्~॥ कृष्णोऽ\appbrk{}रु\appbrk{}ण\appbrk{}नि\appbrk{}भ\appbrk{}स्तब्धः}\ H}}\pend

\pstart\versepar \ledsidenote{\passref{4.25.6}}\edtext{\skt{गुर्वा\-भ\-र\-ण\-सं\-यो\-गा\-त्ता\-ड\-नो\-द्घ\-र्ष\-णा\-दिभिः}}{\lemma{\skt{गुर्वा\appbrk{}भ\appbrk{}र\appbrk{}ण\appbrk{}सं\appbrk{}यो\appbrk{}गा\appbrk{}त्ता\appbrk{}ड\appbrk{}नो\appbrk{}द्घ\appbrk{}र्ष\appbrk{}णा\appbrk{}दिभिः}}\Afootnote{K H; \skt{गुर्वा\appbrk{}भ\appbrk{}र\appbrk{}ण\appbrk{}सं\appbrk{}यो\appbrk{}गा\appbrk{}त्ता\appbrk{}ड\appbrk{}ना\appbrk{}द्घ\appbrk{}र्ष\appbrk{}णा\appbrk{}दपि~।}\ A}}\skt{~। }\pend
\pstart\versepar \skt{शोफः पाल्यां भवेच्छ्यावो }\edtext{\skt{दाह\-पा\-क\-रु\-ज\-न्वितः}}{\lemma{\skt{दाह\appbrk{}पा\appbrk{}क\appbrk{}रु\appbrk{}ज\appbrk{}न्वितः}}\Afootnote{K; \skt{दाह\appbrk{}पा\appbrk{}क\appbrk{}रु\appbrk{}ग\appbrk{}न्वितः~॥}\ A; \skt{दाह\appbrk{}पा\appbrk{}क\appbrk{}रु\appbrk{}जा\appbrk{}न्वितः~।}\ H}}\skt{~।}\pend

\pstart\versepar \ledsidenote{\passref{4.25.7}}\skt{रक्तो वा }\edtext{\skt{रक्त\-पि\-त्ता\-भ्या\-मु\-न्मादः स}}{\lemma{\skt{उन्मादः स}}\Afootnote{K H; \skt{उत्पातः स}\ A}}\skt{ गदो मतः~। }\pend
\pstart\versepar \skt{}\edtext{\skt{कर्ण्णं}}{\lemma{\skt{कर्ण्णं}}\Afootnote{K H; \emph{om.}\ A}}\skt{ }\edtext{\skt{वलाद्वर्धयतः}}{\lemma{\skt{वलाद्वर्धयतः}}\Afootnote{K H; \skt{बलाद्वर्धयतः}\ A}}\skt{ }\edtext{\skt{पाल्यां}}{\lemma{\skt{पाल्यां}}\Afootnote{K; \skt{कर्णं पाल्यां}\ A; \skt{पाल्या}\ H}}\skt{ वायुः प्रकुप्यति~।}\pend

\pstart\versepar \ledsidenote{\passref{4.25.8}}\edtext{\skt{सकफं}}{\lemma{\skt{सकफं}}\Afootnote{K H; \skt{गृहीत्वा सकफं}\ A}}\skt{ }\edtext{\skt{प्राप्य}}{\lemma{\skt{प्राप्य}}\Afootnote{K H; \emph{om.}\ A}}\skt{ }\edtext{\skt{कुरुते}}{\lemma{\skt{कुरुते}}\Afootnote{K H; \skt{कुर्या}\ A}}\skt{ }\edtext{\skt{शोफं}}{\lemma{\skt{शोफं}}\Afootnote{K H; \skt{च्छोफं}\ A}}\skt{ }\edtext{\skt{स्तव्धमवेदनं}}{\lemma{\skt{स्तव्धमवेदनं}}\Afootnote{K; \skt{तद्व\appbrk{}र्ण\appbrk{}वे\appbrk{}द\appbrk{}नम्~।}\ A; \skt{स्तब्धमवेदनं~॥}\ H}}\skt{~। }\pend
\pstart\versepar \skt{}\edtext{\skt{सोवमन्थः}}{\lemma{\skt{सोवमन्थः}}\Afootnote{K H; \skt{उन्मन्थकः}\ A}}\skt{ }\edtext{\skt{सकण्डूको}}{\lemma{\skt{सकण्डूको}}\Afootnote{A K\textsuperscript{pc} H; \skt{सण्डूको}\ K\textsuperscript{ac}}}\skt{ }\edtext{\skt{विदारुः}}{\lemma{\skt{विदारुः}}\Afootnote{K H; \skt{विकारः}\ A}}\skt{ कफवातजः~।}\pend

\pstart\versepar \ledsidenote{\passref{4.25.9}}\edtext{\skt{वर्धमाने तु}}{\lemma{\skt{वर्धमाने तु}}\Afootnote{K H; \skt{वर्धमाने य}\ A}}\skt{ }\edtext{\skt{दुर्विद्ध}}{\lemma{\skt{दुर्विद्ध}}\Afootnote{K; \skt{दा कर्णे}\ A; \skt{दुर्व्विद्धः}\ H}}\skt{ }\edtext{\skt{कण्डू\-दा\-ह\-रु\-जा\-न्वितः}}{\lemma{\skt{कण्डू\appbrk{}दा\appbrk{}ह\appbrk{}रु\appbrk{}जा\appbrk{}न्वितः}}\Afootnote{K\textsuperscript{ac} H; \skt{कण्डू\appbrk{}दा\appbrk{}ह\appbrk{}रु\appbrk{}ग\appbrk{}न्वितः~।}\ A; \skt{कण्डू\appbrk{}दा\appbrk{}ह\appbrk{}रु\appbrk{}जा\appbrk{}न्वितः~।पि}\ K\textsuperscript{pc}}}\skt{ }\pend
\pstart\versepar \skt{शोफो भवति पाकश्च }\edtext{\skt{त्रिदोषो}}{\lemma{\skt{त्रिदोषो}}\Afootnote{K H; \skt{त्वक्स्थोऽसौ}\ A}}\skt{ दुःखवर्द्धनः~।}\pend

\pstart\versepar \ledsidenote{\passref{4.25.10}}\edtext{\skt{कफासृक्क्रिमयः}}{\lemma{\skt{कफासृक्क्रिमयः}}\Afootnote{K H; \skt{कफासृक्कृमयः}\ A}}\skt{ }\edtext{\skt{क्रुद्धाः}}{\lemma{\skt{क्रुद्धाः}}\Afootnote{K H; \skt{कुर्युः}\ A}}\skt{ }\edtext{\skt{सर्ष\-पा\-भा\-वि\-चा\-रिणः}}{\lemma{\skt{सर्ष\appbrk{}पा\appbrk{}भा\appbrk{}वि\appbrk{}चा\appbrk{}रिणः}}\Afootnote{K H\textsuperscript{pc}; \skt{सर्षपाभा विकारिणीः~।}\ A; \skt{सर्ष\appbrk{}पा\appbrk{}वि\appbrk{}चा\appbrk{}रिणः~।}\ H\textsuperscript{ac}}}\skt{~। }\pend
\pstart\versepar \skt{}\edtext{\skt{कुर्वन्ति}}{\lemma{\skt{कुर्वन्ति}}\Afootnote{K H; \emph{om.}\ A}}\skt{ }\edtext{\skt{पाल्यां}}{\lemma{\skt{पाल्यां}}\Afootnote{K H; \emph{om.}\ A}}\skt{ }\edtext{\skt{स्राविण्यः}}{\lemma{\skt{स्राविण्यः}}\Afootnote{K H; \skt{स्राविणीः}\ A}}\skt{ }\edtext{\skt{कण्डू\-दा\-ह\-रु\-जा\-न्विताः}}{\lemma{\skt{कण्डू\appbrk{}दा\appbrk{}ह\appbrk{}रु\appbrk{}जा\appbrk{}न्विताः}}\Afootnote{K; \skt{पिडकाः पाल्यां कण्डू\appbrk{}दा\appbrk{}ह\appbrk{}रु\appbrk{}ग\appbrk{}न्विताः~॥}\ A; \skt{कण्डू\appbrk{}दा\appbrk{}ह\appbrk{}र\appbrk{}जा\appbrk{}न्विताः~।}\ H}}\skt{~।}\pend

\pstart\versepar \ledsidenote{\passref{4.25.11}}\edtext{\skt{कफा\-सृ\-क्क्रि\-मि\-स\-म्भूतः स}}{\lemma{\skt{कफा\appbrk{}सृ\appbrk{}क्क्रि\appbrk{}मि\appbrk{}स\appbrk{}म्भूतः स}}\Afootnote{K H; \skt{कफा\appbrk{}सृ\appbrk{}क्कृ\appbrk{}मि\appbrk{}सं\appbrk{}भूतः स}\ A}}\skt{ }\edtext{\edtext{\skt{विस\-र्पा\-न्वि\-न\-त\-स्तनुः}}{\lemma{\skt{विस\appbrk{}र्पा\appbrk{}न्वि\appbrk{}न\appbrk{}तस्}}\Afootnote{K; \skt{विसर्पन्नित}\ A; \skt{विसर्पान्वितस्}\ H}}}{\lemma{\skt{तनुः}}\Afootnote{K H; \skt{स्ततः~।}\ A}}\skt{~। }\pend
\pstart\versepar \skt{}\edtext{\skt{लिहे\-त्स\-श\-ष्कु\-लिं\-पालिं}}{\lemma{\skt{लिहे\appbrk{}त्स\appbrk{}श\appbrk{}ष्कु\appbrk{}लिं\appbrk{}पालिं}}\Afootnote{K; \skt{लिह्यात्सशष्कुलीं पालीं}\ A; \skt{लिहे\appbrk{}त्स\appbrk{}ष\appbrk{}ष्कु\appbrk{}लिं\appbrk{}पा\appbrk{}लिम्}\ H}}\skt{ परिलेहीति स स्मृतः~।}\pend

\pstart\versepar \ledsidenote{\passref{4.25.12}}\skt{पाल्यामया ह्यमी घोरा नरस्या }\pend
\pstart\versepar \skt{मिथ्या\-हा\-र\-वि\-हा\-रस्य पालिं }\edtext{\edtext{\skt{हिंस्युरुपेक्षिताः}}{\lemma{\skt{हिंस्युर्}}\Afootnote{A K; \skt{हिस्यु}\ H}}}{\lemma{\skt{उपेक्षिताः}}\Afootnote{A K; \skt{रुपेक्षिताः~।}\ H}}\skt{~।}\edtext{}{\lemma{\skt{नरस्या मिथ्या\appbrk{}हा\appbrk{}र\appbrk{}वि\appbrk{}हा\appbrk{}रस्य}}\Afootnote{\skt{नर\appbrk{}स्या\appbrk{}प्र\appbrk{}ति\appbrk{}का\appbrk{}रिणः मिथ्या\appbrk{}हा\appbrk{}र\appbrk{}वि\appbrk{}हा\appbrk{}रस्य}\ A K; \skt{नर\appbrk{}स्या\appbrk{}प्र\appbrk{}ति\appbrk{}का\appbrk{}रणः~॥ मिथ्या\appbrk{}हा\appbrk{}र\appbrk{}वि\appbrk{}हा\appbrk{}रस्य}\ H}}\pend

\pstart\versepar \ledsidenote{\passref{4.25.13}}\skt{तस्मादाशु भिषक्तेषु स्नेहा\-दि\-क्र\-म\-मा\-च\-रेत्~। }\pend
\pstart\versepar \skt{तथा\-भ्य\-ङ्ग\-प\-री\-षे\-क\-प्र\-दे\-हा\-सृ\-ग्वि\-मो\-क्षणं~।}\edtext{}{\lemma{\skt{स्नेहा\appbrk{}दि\appbrk{}क्र\appbrk{}म\appbrk{}मा\appbrk{}च\appbrk{}रेत्~। तथा\appbrk{}भ्य\appbrk{}ङ्ग\appbrk{}प\appbrk{}री\appbrk{}षे\appbrk{}क\appbrk{}प्र\appbrk{}दे\appbrk{}हा\appbrk{}सृ\appbrk{}ग्वि\appbrk{}मो\appbrk{}क्षणं}}\Afootnote{K H; \skt{स्नेहा\appbrk{}दि\appbrk{}क्र\appbrk{}म\appbrk{}मा\appbrk{}च\appbrk{}रेत्~। तथाऽ\appbrk{}भ्य\appbrk{}ङ्ग\appbrk{}प\appbrk{}री\appbrk{}षे\appbrk{}क\appbrk{}प्र\appbrk{}दे\appbrk{}हा\appbrk{}सृ\appbrk{}ग्वि\appbrk{}मो\appbrk{}क्ष\appbrk{}णम्~॥}\ A}}\pend

\pstart\versepar \ledsidenote{\passref{4.25.14}}\skt{सामान्यतो विशेषाच्च वक्ष्याम्यभ्यञ्जनं प्रति~। }\pend
\pstart\versepar \skt{}\edtext{\skt{खर\-मं\-ड\-रि\-य\-ष्ट्या\-ह्व\-सै\-न्ध\-वा\-म\-र\-दा\-रुभिः}}{\lemma{\skt{खर\appbrk{}मं\appbrk{}ड\appbrk{}रि\appbrk{}य\appbrk{}ष्ट्या\appbrk{}ह्व\appbrk{}सै\appbrk{}न्ध\appbrk{}वा\appbrk{}म\appbrk{}र\appbrk{}दा\appbrk{}रुभिः}}\Afootnote{K; \skt{खर\appbrk{}म\appbrk{}ञ्ज\appbrk{}रि\appbrk{}य\appbrk{}ष्ट्या\appbrk{}ह्व\appbrk{}सै\appbrk{}न्ध\appbrk{}वा\appbrk{}म\appbrk{}र\appbrk{}दा\appbrk{}रुभिः~॥}\ A; \skt{खर\appbrk{}म\appbrk{}ञ्ज\appbrk{}रि\appbrk{}य\appbrk{}ष्ट्याह्वा सैन्ध\appbrk{}वा\appbrk{}म\appbrk{}र\appbrk{}दा\appbrk{}रुभिः~।}\ H}}\skt{~।}\pend

\pstart\versepar \ledsidenote{\passref{4.25.15}}\skt{सुपिष्टैः }\edtext{\skt{साश्वगन्धैश्च}}{\lemma{\skt{साश्वगन्धैश्च}}\Afootnote{A K; \skt{साश्वगन्धैश्च}\ H}}\skt{ }\edtext{\skt{मूल\-का\-व\-ल्गु\-जा\-फलैः}}{\lemma{\skt{मूल\appbrk{}का\appbrk{}व\appbrk{}ल्गु\appbrk{}जा\appbrk{}फलैः}}\Afootnote{K H; \skt{मूलकावल्गुजैः फलैः~।}\ A}}\skt{~। }\pend
\pstart\versepar \skt{}\edtext{\skt{सर्पि\-स्तै\-ल\-व\-सा\-म\-ज्जा\-म\-धू\-च्छि\-ष्टानि}}{\lemma{\skt{सर्पि\appbrk{}स्तै\appbrk{}ल\appbrk{}व\appbrk{}सा\appbrk{}म\appbrk{}ज्जा\appbrk{}म\appbrk{}धू\appbrk{}च्छि\appbrk{}ष्टानि}}\Afootnote{K; \skt{सर्पि\appbrk{}स्तै\appbrk{}ल\appbrk{}व\appbrk{}सा\appbrk{}म\appbrk{}ज्ज\appbrk{}म\appbrk{}धू\appbrk{}च्छि\appbrk{}ष्टानि}\ A; \skt{सर्प्पिस्तैलंं वसा\appbrk{}म\appbrk{}ज्जा\appbrk{}म\appbrk{}धू\appbrk{}च्छि\appbrk{}ष्टानि}\ H}}\skt{ }\edtext{\skt{चाहरेत्}}{\lemma{\skt{चाहरेत्}}\Afootnote{K H; \skt{पाचयेत्~॥}\ A}}\skt{~।}\pend

\pstart\versepar \ledsidenote{\passref{4.25.16}}\edtext{\skt{सक्षीराण्यथ}}{\lemma{\skt{सक्षीराण्यथ}}\Afootnote{\skt{सक्षीराण्यथ}\ A K H}}\skt{ तैःपालिं प्रदि\-ह्या\-त्प\-रि\-पो\-टके~। }\pend
\pstart\versepar \skt{}\edtext{\skt{मञ्जि\-ष्ठा\-ति\-ल\-य\-ष्ट्या\-ह्व\-सा\-रि\-वो\-त्प\-ल\-प\-द्मकैः}}{\lemma{\skt{मञ्जि\appbrk{}ष्ठा\appbrk{}ति\appbrk{}ल\appbrk{}य\appbrk{}ष्ट्या\appbrk{}ह्व\appbrk{}सा\appbrk{}रि\appbrk{}वो\appbrk{}त्प\appbrk{}ल\appbrk{}प\appbrk{}द्मकैः}}\Afootnote{A K; \skt{मञ्जि\appbrk{}ष्ठा\appbrk{}ति\appbrk{}ल\appbrk{}य\appbrk{}ष्ट्याह्वा शारि\appbrk{}वो\appbrk{}त्प\appbrk{}ल\appbrk{}प\appbrk{}द्मकैः~।}\ H}}\skt{~।}\pend

\pstart\versepar \ledsidenote{\passref{4.25.17}}\edtext{\skt{सलोध्रैः}}{\lemma{\skt{सलोध्रैः}}\Afootnote{K H; \skt{सरोध्रैः}\ A}}\skt{ }\edtext{\skt{सकदम्वैश्च}}{\lemma{\skt{सकदम्वैश्च}}\Afootnote{K H; \skt{सकदम्बैश्च}\ A}}\skt{ }\edtext{\skt{वला\-ज\-म्व्वा\-म्र\-प\-ल्लवैः}}{\lemma{\skt{वला\appbrk{}ज\appbrk{}म्व्वा\appbrk{}म्र\appbrk{}प\appbrk{}ल्लवैः}}\Afootnote{K H; \skt{बला\appbrk{}ज\appbrk{}म्ब्वा\appbrk{}म्र\appbrk{}प\appbrk{}ल्लवैः~।}\ A}}\skt{~। }\pend
\pstart\versepar \skt{सिद्धं }\edtext{\skt{धान्याम्व्लसंयुक्तं}}{\lemma{\skt{धान्याम्व्लसंयुक्तं}}\Afootnote{K; \skt{धान्याम्लसंयुक्तं}\ A H}}\skt{ }\edtext{\skt{तैल\-मु\-त्पा\-द\-ना\-शनं}}{\lemma{\skt{तैल\appbrk{}मु\appbrk{}त्पा\appbrk{}द\appbrk{}ना\appbrk{}शनं}}\Afootnote{K H; \skt{तैल\appbrk{}मु\appbrk{}त्पा\appbrk{}त\appbrk{}ना\appbrk{}श\appbrk{}नम्~॥}\ A}}\skt{~।}\pend

\pstart\versepar \ledsidenote{\passref{4.25.18}}\edtext{\skt{ताल\-प\-त्र्य\-श्व\-ग\-न्धाभ्यां}}{\lemma{\skt{ताल\appbrk{}प\appbrk{}त्र्य\appbrk{}श्व\appbrk{}ग\appbrk{}न्धाभ्यां}}\Afootnote{K H; \skt{ताल\appbrk{}प\appbrk{}त्र्य\appbrk{}श्व\appbrk{}ग\appbrk{}न्धा\appbrk{}र्क\appbrk{}बा\appbrk{}कु\appbrk{}ची\appbrk{}फल}\ A}}\skt{ }\edtext{\skt{सैन्ध\-वा\-व\-ल्गु\-जा\-फलैः}}{\lemma{\skt{सैन्ध\appbrk{}वा\appbrk{}व\appbrk{}ल्गु\appbrk{}जा\appbrk{}फलैः}}\Afootnote{K H; \skt{सैन्धवैः~।}\ A}}\skt{~। }\pend
\pstart\versepar \skt{तैलं कुलीरगोधाभ्यां वसया सह }\edtext{\skt{साधितं}}{\lemma{\skt{साधितं}}\Afootnote{K H; \skt{पाचितम्~॥}\ A}}\skt{~।}\pend

\pstart\versepar \ledsidenote{\passref{4.25.19}}\skt{सर\-ला\-ला\-ङ्ग\-ली\-भ्याञ्च }\edtext{\edtext{\skt{हित\-मु\-न्म\-न्थ\-ना\-शनं}}{\lemma{\skt{हितम्}}\Afootnote{K; \skt{हित}\ A H}}}{\lemma{\skt{उन्मन्थनाशनं}}\Afootnote{K; \skt{मुन्म\appbrk{}न्थ\appbrk{}ना\appbrk{}श\appbrk{}नम्~।}\ A; \skt{मुत्मन्थनाशनः~॥}\ H}}\skt{~। }\pend
\pstart\versepar \skt{}\edtext{\skt{अथा\-श्म\-न्त\-क\-ज\-म्व्वा\-म्र\-प\-त्र\-क्वा\-थ\-नि\-षे\-चितां}}{\lemma{\skt{अथा\appbrk{}श्म\appbrk{}न्त\appbrk{}क\appbrk{}ज\appbrk{}म्व्वा\appbrk{}म्र\appbrk{}प\appbrk{}त्र\appbrk{}क्वा\appbrk{}थ\appbrk{}नि\appbrk{}षे\appbrk{}चितां}}\Afootnote{K; \skt{तथाऽ\appbrk{}श्म\appbrk{}न्त\appbrk{}क\appbrk{}ज\appbrk{}म्ब्वा\appbrk{}म्र\appbrk{}प\appbrk{}त्र\appbrk{}क्वा\appbrk{}थेन सेचनम्~॥}\ A; \skt{अथा\appbrk{}श्म\appbrk{}न्त\appbrk{}क\appbrk{}ज\appbrk{}म्व्वाम्र पत्र\appbrk{}क्वा\appbrk{}थ\appbrk{}नि\appbrk{}षे\appbrk{}वितां~।}\ H}}\skt{~।}\pend

\pstart\versepar \ledsidenote{\passref{4.25.20}}\edtext{\skt{प्रपौ\-ण्डी\-री\-क\-म\-धु\-क\-म\-ञ्जिष्ठं}}{\lemma{\skt{प्रपौ\appbrk{}ण्डी\appbrk{}री\appbrk{}क\appbrk{}म\appbrk{}धु\appbrk{}क\appbrk{}म\appbrk{}ञ्जिष्ठं}}\Afootnote{K; \skt{प्रपौ\appbrk{}ण्ड\appbrk{}री\appbrk{}क\appbrk{}म\appbrk{}धु\appbrk{}क\appbrk{}म\appbrk{}ञ्जिष्ठा}\ A; \skt{प्रपौ\appbrk{}ण्ड\appbrk{}री\appbrk{}क\appbrk{}म\appbrk{}धु\appbrk{}क\appbrk{}म\appbrk{}ञ्जिष्ठं}\ H}}\skt{ }\edtext{\skt{रजनीकृतैः}}{\lemma{\skt{रजनीकृतैः}}\Afootnote{K H; \skt{रजनीद्वयैः~।}\ A}}\skt{~। }\pend
\pstart\versepar \skt{}\edtext{\skt{चूर्ण्णैर्दुर्गन्धनै}}{\lemma{\skt{दुर्गन्धनै}}\Afootnote{K; \skt{उद्वर्तनैः}\ A; \skt{दुर्ग्गन्धने}\ H}}\skt{ }\edtext{\skt{पालि\-न्तै\-ला\-क्ता\-म\-व\-चू\-र्ण्ण\-येत्}}{\lemma{\skt{पालि\appbrk{}न्तै\appbrk{}ला\appbrk{}क्ता\appbrk{}म\appbrk{}व\appbrk{}चू\appbrk{}र्ण्ण\appbrk{}येत्}}\Afootnote{\skt{पालीं तैला\appbrk{}क्ता\appbrk{}म\appbrk{}व\appbrk{}चू\appbrk{}ण\appbrk{}येत्~॥}\ A; \skt{पालि\appbrk{}न्तै\appbrk{}ला\appbrk{}क्ता\appbrk{}म\appbrk{}व\appbrk{}चू\appbrk{}र्ण्ण\appbrk{}येत्}\ K H}}\skt{~।}\pend

\pstart\versepar \ledsidenote{\passref{4.25.21}}\edtext{\edtext{\skt{चूर्ण्णै\-र्वि\-ड\-ङ्ग\-ला\-क्षाभ्यां}}{\lemma{\skt{चूर्ण्णैर्}}\Afootnote{K; \skt{लाक्षा}\ A; \skt{चूर्ण्णै}\ H}}}{\lemma{\skt{विडङ्गलाक्षाभ्यां}}\Afootnote{K; \skt{विडङ्गकल्केन}\ A; \skt{र्व्विड\appbrk{}ङ्ग\appbrk{}ला\appbrk{}क्षा\appbrk{}भ्यान्}\ H}}\skt{ तैलं पक्त्वा\-व\-चा\-र\-येत्~। }\pend
\pstart\versepar \skt{स्विन्नां }\edtext{\skt{गोमयनिःपिण्डैः}}{\lemma{\skt{गोमयनिःपिण्डैः}}\Afootnote{K H; \skt{गोमयपिण्डेन}\ A}}\skt{ }\edtext{\skt{प्रदि\-ह्या\-त्प\-रि\-हिके}}{\lemma{\skt{प्रदि\appbrk{}ह्या\appbrk{}त्प\appbrk{}रि\appbrk{}हिके}}\Afootnote{\skt{प्रदि\appbrk{}ह्या\appbrk{}त्प\appbrk{}रि\appbrk{}ले\appbrk{}हिके}\ A K H}}\skt{~।}\edtext{}{\lemma{\skt{पक्त्वा\appbrk{}व\appbrk{}चा\appbrk{}र\appbrk{}येत्~। स्विन्नां}}\Afootnote{K; \skt{पक्त्वाऽ\appbrk{}व\appbrk{}चा\appbrk{}र\appbrk{}येत्~। स्विन्नां}\ A; \skt{पक्वा\appbrk{}व\appbrk{}चा\appbrk{}र\appbrk{}येत्~॥ स्विन्नां}\ H}}\pend

\pstart\versepar \ledsidenote{\passref{4.25.22}}\edtext{\edtext{\skt{कृष्णै\-र्वि\-ड\-ङ्गे\-र\-थवा}}{\lemma{\skt{कृष्णैर्}}\Afootnote{K; \skt{पिष्टैर्}\ A; \skt{कृष्णै}\ H}}}{\lemma{\skt{विडङ्गेर्}}\Afootnote{K; \skt{विडङ्गैर्}\ A; \skt{र्व्विडङ्गैर्}\ H}}\skt{ }\edtext{\skt{प्रर\-त्री\-मू\-त्र\-पे\-षिमैः}}{\lemma{\skt{प्रर\appbrk{}त्री\appbrk{}मू\appbrk{}त्र\appbrk{}पे\appbrk{}षिमैः}}\Afootnote{K; \skt{त्रिवृ\appbrk{}च्छ्या\appbrk{}मा\appbrk{}र्क\appbrk{}सं\appbrk{}युतैः~।}\ A; \skt{पर\appbrk{}त्त्री\appbrk{}मू\appbrk{}त्र\appbrk{}पे\appbrk{}षिमैः~॥}\ H}}\skt{~। }\pend
\pstart\versepar \skt{}\edtext{\skt{कर\-ण्जे\-ङ्गु\-दि\-बी\-जैर्वा}}{\lemma{\skt{कर\appbrk{}ण्जे\appbrk{}ङ्गु\appbrk{}दि\appbrk{}बी\appbrk{}जैर्वा}}\Afootnote{A K; \skt{कर\appbrk{}ञ्जै\appbrk{}ङ्गु\appbrk{}दि\appbrk{}वी\appbrk{}जैर्व्वा}\ H}}\skt{ कुट\-जा\-र\-ग्व\-धा\-युतैः~।}\pend

\pstart\versepar \ledsidenote{\passref{4.25.23}}\edtext{\skt{सर्वैर्वा}}{\lemma{\skt{सर्वैर्वा}}\Afootnote{A H; \skt{सर्वैर्वा}\ K}}\skt{ सार्षपं तैलं सिद्धं }\edtext{\skt{मरिचसंयुतैः}}{\lemma{\skt{मरिचसंयुतैः}}\Afootnote{K H; \skt{मरि\appbrk{}च\appbrk{}सं\appbrk{}यु\appbrk{}तम्~।}\ A}}\skt{~। }\pend
\pstart\versepar \skt{}\edtext{\skt{सनि\-म्व\-प\-त्रै\-र\-भ्यङ्गे}}{\lemma{\skt{सनि\appbrk{}म्व\appbrk{}प\appbrk{}त्रै\appbrk{}र\appbrk{}भ्यङ्गे}}\Afootnote{H; \skt{सनि\appbrk{}म्ब\appbrk{}प\appbrk{}त्रै\appbrk{}र\appbrk{}भ्यङ्गे}\ A; \skt{सनि\appbrk{}म्व\appbrk{}प\appbrk{}त्रै\appbrk{}र\appbrk{}भ्यङ्गे}\ K}}\skt{ }\edtext{\skt{मधुच्छिष्टान्वितो}}{\lemma{\skt{मधुच्छिष्टान्वितो}}\Afootnote{K; \skt{मधूच्छिष्टान्वितं}\ A; \skt{मधूच्छिष्टान्वितो}\ H}}\skt{ }\edtext{\skt{हितः}}{\lemma{\skt{हितः}}\Afootnote{K H; \skt{हितम्~॥}\ A}}\skt{~।}\pend

\pstart\versepar \ledsidenote{\passref{4.25.24}}\edtext{\skt{व्याधिक्लिष्टासु}}{\lemma{\skt{व्याधिक्लिष्टासु}}\Afootnote{K H; \emph{om.}\ A}}\skt{ पालीषु }\edtext{\skt{तन्वीषु}}{\lemma{\skt{तन्वीषु}}\Afootnote{K H; \skt{व्याधियुक्तासु तन्वीषु}\ A}}\skt{ कठिनासु च~। }\pend
\pstart\versepar \skt{}\edtext{\skt{पुष्ट्य\-र्थ\-म्मा\-र्द\-वा\-र्थञ्च}}{\lemma{\skt{पुष्ट्यर्थम्}}\Afootnote{K H; \skt{पुष्ट्यार्थं}\ A}}\skt{ कुर्यादभ्यञ्जनं }\edtext{\skt{शुभं}}{\lemma{\skt{शुभं}}\Afootnote{K H; \skt{त्विदम्~॥}\ A}}\skt{~।}\pend

\pstart\versepar \ledsidenote{\passref{4.25.25}}\edtext{\skt{लोपा\-का\-नू\-प\-जा\-म\-ज्जा\-व\-सा\-तैलं}}{\lemma{\skt{लोपा\appbrk{}का\appbrk{}नू\appbrk{}प\appbrk{}जा\appbrk{}म\appbrk{}ज्जा\appbrk{}व\appbrk{}सा\appbrk{}तैलं}}\Afootnote{K H; \skt{लोपा\appbrk{}का\appbrk{}नू\appbrk{}प\appbrk{}म\appbrk{}ज्जानं वसां तैलं}\ A}}\skt{ नवं घृतं~। }\pend
\pstart\versepar \skt{}\edtext{\skt{दशक्षीरं}}{\lemma{\skt{दशक्षीरं}}\Afootnote{K H; \emph{om.}\ A}}\skt{ }\edtext{\skt{पचे\-त्स\-म्य\-गा\-वाप्य}}{\lemma{\skt{पचे\appbrk{}त्स\appbrk{}म्य\appbrk{}गा\appbrk{}वाप्य}}\Afootnote{K; \skt{पचेद्दशगुणं क्षीरमावाप्य}\ A; \skt{पचे\appbrk{}त्स\appbrk{}म्य\appbrk{}ग्भा\appbrk{}वाप्य}\ H}}\skt{ मधुरंगणं~।}\pend

\pstart\versepar \ledsidenote{\passref{4.25.26}}\edtext{\skt{अपा\-मा\-र्गा\-श्व\-गन्धा च}}{\lemma{\skt{अपा\appbrk{}मा\appbrk{}र्गा\appbrk{}श्व\appbrk{}गन्धा च}}\Afootnote{K H; \skt{अपा\appbrk{}मा\appbrk{}र्गा\appbrk{}श्व\appbrk{}गन्धे च}\ A}}\skt{ तथा लाक्षारसं शुभं~। }\pend
\pstart\versepar \skt{}\edtext{\skt{अथ}}{\lemma{\skt{अथ}}\Afootnote{K H; \skt{त}\ A}}\skt{ }\edtext{\skt{सिद्धञ्च}}{\lemma{\skt{सिद्धञ्च}}\Afootnote{K H; \skt{त्सिद्धं प}\ A}}\skt{ }\edtext{\skt{पूतञ्च}}{\lemma{\skt{पूतञ्च}}\Afootnote{K H; \skt{रिपूतं च}\ A}}\skt{ स्वनु\-गु\-प्त\-न्नि\-धा\-प\-येत्~।}\pend

\pstart\versepar \ledsidenote{\passref{4.25.27}}\edtext{\skt{तेनाभ्यञ्जात्सदा}}{\lemma{\skt{तेनाभ्यञ्जात्सदा}}\Afootnote{K H; \skt{तेनाभ्यञ्ज्यात्सदा}\ A}}\skt{ }\edtext{\skt{पालिं}}{\lemma{\skt{पालिं}}\Afootnote{K H; \skt{पालीं}\ A}}\skt{ }\edtext{\edtext{\skt{सुस्वि\-न्ना\-म\-ति\-म\-र्दितां}}{\lemma{\skt{सुस्विन्नाम्}}\Afootnote{A K; \skt{सुस्विन्ना}\ H}}}{\lemma{\skt{अतिमर्दितां}}\Afootnote{K; \skt{तिमर्दिताम्~।}\ A; \skt{मभिमर्द्दितां~॥}\ H}}\skt{~। }\pend
\pstart\versepar \skt{एतेन पाल्यो वर्धन्ते }\edtext{\skt{निरुजा}}{\lemma{\skt{निरुजा}}\Afootnote{K H; \skt{निरुजो}\ A}}\skt{ }\edtext{\skt{निरुपद्रवाः}}{\lemma{\skt{निरुपद्रवाः}}\Afootnote{A K; \skt{निरुपद्रवा~।}\ H}}\skt{~।}\pend

\pstart\versepar \ledsidenote{\passref{4.25.28}}\skt{मृद्व्यः पुष्टाः समाः स्निग्धा जायन्ते भूषणक्षमाः~॥ }\pend
\pstart\versepar \skt{}\edtext{\skt{नीली\-प\-त्रं\-भृ\-ङ्ग\-र\-जो\-र्जु\-न\-त्व\-क्तृ\-फ\-ला\-चूर्ण्णं}}{\lemma{\skt{नीली\appbrk{}प\appbrk{}त्रं\appbrk{}भृ\appbrk{}ङ्ग\appbrk{}र\appbrk{}जो\appbrk{}र्जु\appbrk{}न\appbrk{}त्व\appbrk{}क्तृ\appbrk{}फ\appbrk{}ला\appbrk{}चूर्ण्णं}}\Afootnote{K\textsuperscript{pc}; \skt{नीलीदलं भृङ्ग\appbrk{}र\appbrk{}जोऽ\appbrk{}र्जु\appbrk{}न\appbrk{}त्व\appbrk{}क्पि\appbrk{}ण्डी\appbrk{}तकं कृष्ण\appbrk{}म\appbrk{}यो\appbrk{}र\appbrk{}जश्च~। बीजोद्भवं}\ A; \skt{नीली\appbrk{}प\appbrk{}त्रं\appbrk{}भृ\appbrk{}ङ्ग\appbrk{}र\appbrk{}जो\appbrk{}र्जु\appbrk{}न\appbrk{}त्व\appbrk{}फ\appbrk{}ला\appbrk{}चूर्ण्णं}\ K\textsuperscript{ac}; \skt{नीलीपत्रं भृङ्ग\appbrk{}र\appbrk{}जोऽ\appbrk{}र्ज्जु\appbrk{}न\appbrk{}त्व\appbrk{}क्तृ\appbrk{}फ\appbrk{}ला\appbrk{}चूर्ण्णं}\ H}}\skt{ साहचरञ्च }\edtext{\skt{पुष्पं}}{\lemma{\skt{पुष्पं}}\Afootnote{K H; \emph{om.}\ A}}\skt{~। }\pend
\pstart\versepar \skt{}\edtext{\skt{श्यामापुष्पं}}{\lemma{\skt{श्यामापुष्पं}}\Afootnote{K H; \skt{पुष्पं}\ A}}\skt{ }\edtext{\skt{कृष्ण\-पि\-ण्डी\-त\-कञ्च }\edtext{\skt{सूक्ष्मञ्चूर्ण्णं}}{\lemma{\skt{सूक्ष्मञ्चूर्ण्णं}}\Afootnote{K; \skt{शूक्ष्मञ्चूर्ण्णं}\ H}}\skt{ सुकृतं }\edtext{\skt{यावदेतत्}}{\lemma{\skt{यावदेतत्}}\Afootnote{\skt{यावदेतत्~।}\ K; \skt{यावदेव~॥}\ H}}\skt{~।}}{\lemma{\skt{कृष्ण\appbrk{}पि\appbrk{}ण्डी\appbrk{}त\appbrk{}कञ्… एत}}\Afootnote{\emph{om.}\ A}}\pend

\pstart\versepar \ledsidenote{\passref{4.25.29}}\edtext{\edtext{\skt{तावान्देयः}}{\lemma{\skt{तावान्}}\Afootnote{K H; \skt{एकीकृतं सर्वमिदं प्रमाय पङ्केन तुल्यं नलिनीभवे}\ A}}}{\lemma{\skt{देयः}}\Afootnote{K H; \skt{न~। सं}\ A}}\skt{ }\edtext{\skt{कर्दमः}}{\lemma{\skt{कर्दमः}}\Afootnote{K H; \skt{योज्य प}\ A}}\skt{ }\edtext{\skt{पङ्कषण्डात्~। पक्त\-ञ्चै\-त\-त्स्था\-प\-ये\-ल्लो\-ह\-कुम्भ्यां}}{\lemma{\skt{पङ्कषण्डात्~। पक्त\appbrk{}ञ्चै\appbrk{}त\appbrk{}त्स्था\appbrk{}प\appbrk{}ये\appbrk{}ल्लो\appbrk{}ह\appbrk{}कुम्भ्यां}}\Afootnote{K; \skt{क्षं कलशे निधाय}\ A; \skt{पङ्क\appbrk{}ष\appbrk{}ण्डा\appbrk{}त्प\appbrk{}क्व\appbrk{}ञ्चै\appbrk{}त\appbrk{}त्स्था\appbrk{}प\appbrk{}ये\appbrk{}ल्लो\appbrk{}ह\appbrk{}कुम्भ्यां~।}\ H}}\skt{~। }\pend
\pstart\versepar \skt{}\edtext{\edtext{\skt{पचे\-दे\-त\-त्ति\-ल\-तै\-लेन}}{\lemma{\skt{पचे\appbrk{}दे\appbrk{}त\appbrk{}त्ति\appbrk{}ल\appbrk{}तै\appbrk{}लेन}}\Afootnote{H; \skt{लौहे विपचे}\ A; \skt{पचे\appbrk{}दे\appbrk{}त\appbrk{}त्ति\appbrk{}ल\appbrk{}तै\appbrk{}लेन}\ K}}\skt{ सार्धं दत्वा }\edtext{\skt{क्वाथं}}{\lemma{\skt{क्वाथं}}\Afootnote{K H; \skt{मिश्रं रसेन भृङ्}\ A}}\skt{ }\edtext{\skt{त्रिफ\-ला\-मा\-र्क\-वाभ्यां}}{\lemma{\skt{त्रिफ\appbrk{}ला\appbrk{}मा\appbrk{}र्क\appbrk{}वाभ्यां}}\Afootnote{K H; \skt{गत्रि\appbrk{}फ\appbrk{}ला\appbrk{}भ\appbrk{}वेन~। आसन्नपाके च परीक्षणार्थं पत्रं बला\appbrk{}का\appbrk{}भ\appbrk{}व\appbrk{}मा\appbrk{}क्षि\appbrk{}पेच्च~॥}\ A}}}{\lemma{\skt{द्… दत्वा}}\Afootnote{\emph{om.}\ A}}\skt{~।}\edtext{}{\lemma{\skt{त्… स्थाप\appbrk{}ये\appbrk{}ल्लो\appbrk{}ह\appbrk{}कुम्भ्यां}}\Afootnote{\emph{om.}\ A}}\pend

\pstart\versepar \ledsidenote{\passref{4.25.30}}\edtext{\edtext{\skt{मासादूर्ध्वं}}{\lemma{\skt{मासादूर्ध्वं}}\Afootnote{K; \skt{मासादूर्द्धम्}\ H}}\skt{ म्रक्षणेनैव तैलं }\edtext{\skt{कुर्या\-दे\-त\-त्प\-लितं}}{\lemma{\skt{कुर्या\appbrk{}दे\appbrk{}त\appbrk{}त्प\appbrk{}लितं}}\Afootnote{H; \skt{कुर्या\appbrk{}दे\appbrk{}त\appbrk{}त्प\appbrk{}लितं}\ K}}}{\lemma{\skt{… कुर्या}}\Afootnote{\emph{om.}\ A}}\skt{ }\edtext{\skt{स्निग्धनीलं}}{\lemma{\skt{स्निग्धनीलं}}\Afootnote{K H; \skt{लं}\ A}}\skt{~॥ }\pend
\pstart\versepar \skt{}\edtext{\skt{जम्वू\-पु\-ष्प\-का\-र्श्म\-री\-पा\-र्थि\-वञ्च}}{\lemma{\skt{जम्वू\appbrk{}पु\appbrk{}ष्प\appbrk{}का\appbrk{}र्श्म\appbrk{}री\appbrk{}पा\appbrk{}र्थि\appbrk{}वञ्च}}\Afootnote{K; \skt{विपचेद्विमिश्रं र}\ A; \skt{जम्वूपुष्पं काश्म\appbrk{}री\appbrk{}पा\appbrk{}र्थि\appbrk{}वञ्च}\ H}}\skt{ }\edtext{\skt{सैरीयकं}}{\lemma{\skt{सैरीयकं}}\Afootnote{K H; \skt{सेन भृङ्ग}\ A}}\skt{ }\edtext{\skt{त्रि}}{\lemma{\skt{त्रि}}\Afootnote{K; \skt{त्रिफलाभवेन। आसन्नपाके च परीक्षणार्थं पत्रं बला\appbrk{}का\appbrk{}भ\appbrk{}व\appbrk{}मा\appbrk{}क्षि\appbrk{}पेच्च॥}\ A; \skt{त्रिफला कर्द्दमञ्च~॥}\ H}}\pend

\pstart\versepar \ledsidenote{\passref{4.25.31}}\skt{द्वौव\-र्षा\-भू\-शा\-रि\-वा\-क\-ण्ट\-कार्यौ }\edtext{\skt{नीला\-स्ति\-ला\-स्रो\-त\-ज\-ञ्चो\-त्प\-लञ्च}}{\lemma{\skt{नीला\appbrk{}स्ति\appbrk{}ला\appbrk{}स्रो\appbrk{}त\appbrk{}ज\appbrk{}ञ्चो\appbrk{}त्प\appbrk{}लञ्च}}\Afootnote{K; \skt{नीला\appbrk{}स्ति\appbrk{}ला\appbrk{}स्रो\appbrk{}त्र\appbrk{}ज\appbrk{}ञ्चो\appbrk{}त्प\appbrk{}लञ्च~।}\ H}}\skt{~। }\pend
\pstart\versepar \skt{आय\-श्चू\-र्ण्ण\-ञ्चू\-त\-प\-क्वस्य मज्जा यष्ट्याह्वायं मार्कवो मोदयन्ती~।}\edtext{}{\lemma{\skt{… मोदयन्ती}}\Afootnote{\emph{om.}\ A}}\pend

\pstart\versepar \ledsidenote{\passref{4.25.32}}\edtext{\skt{कासीसमग्र्यं}}{\lemma{\skt{कासीसमग्र्यं}}\Afootnote{K\textsuperscript{pc} H; \skt{सैरी\appbrk{}य\appbrk{}ज\appbrk{}म्ब्व\appbrk{}र्जु\appbrk{}न\appbrk{}का\appbrk{}श्म\appbrk{}रीजं पुष्पं तिला\appbrk{}न्मा\appbrk{}र्क\appbrk{}व\appbrk{}चू\appbrk{}त\appbrk{}बीजे~। पुनर्नवे कर्द\appbrk{}म\appbrk{}क\appbrk{}ण्ट\appbrk{}कार्यौ कासी\appbrk{}स\appbrk{}पि\appbrk{}ण्डी\appbrk{}त\appbrk{}क\appbrk{}बीज}\ A; \skt{कासीसमग्र्यं स}\ K\textsuperscript{ac}}}\skt{ }\edtext{\edtext{\skt{सारश्चासनस्य}}{\lemma{\skt{सारश्}}\Afootnote{K H; \skt{सारम्~॥}\ A}}}{\lemma{\skt{चासनस्य}}\Afootnote{K H; \emph{om.}\ A}}\skt{ }\edtext{\skt{वराहसाह्वस्य}}{\lemma{\skt{वराहसाह्वस्य}}\Afootnote{K H; \emph{om.}\ A}}\skt{ }\edtext{\skt{फलञ्च}}{\lemma{\skt{फलञ्च}}\Afootnote{K H; \emph{om.}\ A}}\skt{ }\edtext{\skt{पाण्डूं}}{\lemma{\skt{पाण्डूं}}\Afootnote{K; \emph{om.}\ A; \skt{पाण्डुं~।}\ H}}\skt{~। }\pend
\pstart\versepar \skt{}\edtext{\skt{भागानेतां}}{\lemma{\skt{भागानेतां}}\Afootnote{K; \emph{om.}\ A; \skt{भागानेता}\ H}}\skt{ }\edtext{\skt{कार्षिकां}}{\lemma{\skt{कार्षिकां}}\Afootnote{K; \emph{om.}\ A; \skt{न्कार्षिका}\ H}}\skt{ }\edtext{\skt{पीषयित्वा}}{\lemma{\skt{पीषयित्वा}}\Afootnote{K; \emph{om.}\ A; \skt{न्पीषयित्वा}\ H}}\skt{ }\edtext{\skt{सारो\-द\-के\-ना\-स\-न\-स्यैव}}{\lemma{\skt{सारो\appbrk{}द\appbrk{}के\appbrk{}ना\appbrk{}स\appbrk{}न\appbrk{}स्यैव}}\Afootnote{K H; \emph{om.}\ A}}\skt{ }\edtext{\skt{विद्वान्}}{\lemma{\skt{विद्वान्}}\Afootnote{K H; \emph{om.}\ A}}\skt{~।}\pend

\pstart\versepar \ledsidenote{\passref{4.25.33}}\edtext{\skt{कल्कानेतां सप्त\-स\-प्ता\-ढ\-केषु}}{\lemma{\skt{कल्कानेतां सप्त\appbrk{}स\appbrk{}प्ता\appbrk{}ढ\appbrk{}केषु}}\Afootnote{K; \skt{फलत्रयं लोहरजोऽञ्जनं च यष्ट्याह्वयं नीरजसारिवे च~। पिष्ट्वाऽथ सर्वं सह मोदयन्त्या साराम्भ}\ A; \skt{कल्काने एता\appbrk{}न्स\appbrk{}प्त\appbrk{}स\appbrk{}प्ता\appbrk{}ढ\appbrk{}केषु}\ H\textsuperscript{ac}; \skt{कल्काने एता\appbrk{}न्स\appbrk{}प्त\appbrk{}सु\appbrk{}ह्या\appbrk{}स\appbrk{}प्ता\appbrk{}ढ\appbrk{}केषु}\ H\textsuperscript{pc}}}\skt{ }\edtext{\skt{क्षारो\-द\-क\-स्या\-क्ष\-तै\-लस्य}}{\lemma{\skt{क्षारो\appbrk{}द\appbrk{}क\appbrk{}स्या\appbrk{}क्ष\appbrk{}तै\appbrk{}लस्य}}\Afootnote{K; \skt{सा बीज\appbrk{}क\appbrk{}सं\appbrk{}भ\appbrk{}वेन~॥}\ A; \skt{सारो\appbrk{}द\appbrk{}क\appbrk{}स्या\appbrk{}क्ष\appbrk{}तै\appbrk{}लस्य}\ H}}\skt{ }\edtext{\skt{भागं}}{\lemma{\skt{भागं}}\Afootnote{K H; \emph{om.}\ A}}\skt{~। }\pend
\pstart\versepar \skt{}\edtext{\skt{दत्वा}}{\lemma{\skt{दत्वा}}\Afootnote{K H; \emph{om.}\ A}}\skt{ }\edtext{\skt{सम्य\-ग्चे\-ल्लो\-ह\-पात्रे}}{\lemma{\skt{सम्य\appbrk{}ग्चे\appbrk{}ल्लो\appbrk{}ह\appbrk{}पात्रे}}\Afootnote{\emph{om.}\ A; \skt{सम्य\appbrk{}ग्वि\appbrk{}प\appbrk{}चे\appbrk{}ल्लो\appbrk{}ह\appbrk{}पात्रे}\ K H}}\skt{ }\edtext{\skt{स्थितं पूर्वं}}{\lemma{\skt{स्थितं पूर्वं}}\Afootnote{K H; \emph{om.}\ A}}\skt{ }\edtext{\skt{लोहपात्रे}}{\lemma{\skt{लोहपात्रे}}\Afootnote{K H; \emph{om.}\ A}}\skt{ }\edtext{\skt{दशावहं}}{\lemma{\skt{दशावहं}}\Afootnote{K H; \emph{om.}\ A}}\skt{~।}\pend

\pstart\versepar \ledsidenote{\passref{4.25.34}}\skt{मृदुपाकं साधितं चावतार्य }\edtext{\skt{सेत्पात्रे}}{\lemma{\skt{सेत्पात्रे}}\Afootnote{\skt{न्यसेत्पात्रे}\ K H}}\skt{ }\edtext{\skt{बीजके}}{\lemma{\skt{बीजके}}\Afootnote{K; \skt{वीजके}\ H}}\skt{ वायसे वा~। }\pend
\pstart\versepar \skt{नरः शुद्धो }\edtext{\skt{नस्य\-मे\-त\-त्प्र\-यु\-ज्या\-त्यथा}}{\lemma{\skt{एत\appbrk{}त्प्र\appbrk{}यु\appbrk{}ज्या\appbrk{}त्यथा}}\Afootnote{K; \skt{एत\appbrk{}त्प्र\appbrk{}यु\appbrk{}ञ्जे\appbrk{}त्यथा}\ H}}\skt{ वत्स्यात्संयतो मासमेकं~।}\edtext{}{\lemma{\skt{… मासमेकं}}\Afootnote{\emph{om.}\ A}}\pend

\pstart \ledsidenote{\passref{4.25.35}}\edtext{\skt{भवेच्च}}{\lemma{\skt{भवेच्च}}\Afootnote{K H; \skt{पक्वं च लौहेऽभिनवे निधाय नस्यं विद\appbrk{}ध्या\appbrk{}त्प\appbrk{}रि\appbrk{}शु\appbrk{}द्ध\appbrk{}कायः~। अभ्यङ्गयोगैश्च}\ A}}\skt{ }\edtext{\skt{नित्यं}}{\lemma{\skt{नित्यं}}\Afootnote{K H; \skt{नि}\ A}}\skt{ }\edtext{\skt{कृसरोपसेवी}}{\lemma{\skt{कृसरोपसेवी}}\Afootnote{K H; \skt{युज्यमानं}\ A}}\skt{ भुञ्जीत }\edtext{\skt{चान्नं}}{\lemma{\skt{चान्नं}}\Afootnote{K H; \skt{माषा\appbrk{}न्कृ\appbrk{}श\appbrk{}रा\appbrk{}मथो}\ A}}\skt{ }\edtext{\skt{माषयूषेण}}{\lemma{\skt{माषयूषेण}}\Afootnote{K; \skt{वा~॥}\ A; \skt{मासयूषेण}\ H}}\skt{ }\edtext{\skt{चापि}}{\lemma{\skt{चापि}}\Afootnote{K H; \emph{om.}\ A}}\skt{~।}\pend

\pstart\versepar \ledsidenote{\passref{4.25.36}}\edtext{\skt{मासादूर्ध्वं}}{\lemma{\skt{मासादूर्ध्वं}}\Afootnote{K; \skt{मासोपरि}\ A; \emph{om.}\ H}}\skt{ }\edtext{\skt{तस्य}}{\lemma{\skt{तस्य}}\Afootnote{K; \skt{ष्टाद्घ\appbrk{}न\appbrk{}कु\appbrk{}ञ्जि\appbrk{}ताग्राः}\ A; \skt{मासादूर्द्ध्वन्तस्य}\ H}}\skt{ केशा भवन्ति }\edtext{\skt{यावज्जीवं}}{\lemma{\skt{यावज्जीवं}}\Afootnote{K H; \emph{om.}\ A}}\skt{ }\edtext{\skt{मृदु\-व\-त्स्व\-ञ्ज\-नाभं}}{\lemma{\skt{मृदु\appbrk{}व\appbrk{}त्स्व\appbrk{}ञ्ज\appbrk{}नाभं}}\Afootnote{K H; \skt{भ्रमराञ्जनाभाः~।}\ A}}\skt{~। }\pend
\pstart\versepar \skt{}\edtext{\skt{भवन्त्येते नेन्द्रलुप्ते च }\edtext{\skt{केशाः}}{\lemma{\skt{केशाः}}\Afootnote{K H; \skt{केशास्तथा}\ A}}\skt{~। }\edtext{\skt{जरा}}{\lemma{\skt{जरा}}\Afootnote{K H; \skt{ऽन्ये खलतौ भवेयुर्जरा न}\ A}}\skt{ चैनं}}{\lemma{\skt{भवन्त्य्… के}}\Afootnote{\emph{om.}\ A}}\skt{ सहसा }\edtext{\skt{नाभ्यपैति}}{\lemma{\skt{नाभ्यपैति}}\Afootnote{K; \skt{ऽभ्युपैति~॥}\ A; \skt{नाभ्युपैति~।}\ H}}\skt{~।}\pend

\pstart\versepar \ledsidenote{\passref{4.25.37}}\edtext{\skt{बलं}}{\lemma{\skt{बलं}}\Afootnote{A H; \skt{वलं}\ K}}\skt{ परं }\edtext{\skt{चेन्द्रियाणां}}{\lemma{\skt{चेन्द्रियाणां}}\Afootnote{K H; \skt{संभ\appbrk{}व\appbrk{}ती\appbrk{}न्दि\appbrk{}र्याणां}\ A}}\skt{ }\edtext{\skt{लभेत}}{\lemma{\skt{लभेत}}\Afootnote{K H; \skt{भवेच्च}\ A}}\skt{ सञ्जायन्ते निर्वलीकञ्च वक्त्रं~। }\pend
\pstart\versepar \skt{}\edtext{\skt{एतत्तैलं}}{\lemma{\skt{एतत्तैलं}}\Afootnote{K H; \skt{वलि\appbrk{}भि\appbrk{}र्वि\appbrk{}मु\appbrk{}क्तम्~। नाका\appbrk{}मि\appbrk{}नेऽ\appbrk{}ना\appbrk{}र्थिनि नाकृताय नैवारये तै}\ A}}\skt{ }\edtext{\skt{नाप्रियाय}}{\lemma{\skt{नाप्रियाय}}\Afootnote{K H; \skt{ल}\ A}}\skt{ }\edtext{\skt{प्रदेयं}}{\lemma{\skt{प्रदेयं}}\Afootnote{K H; \skt{मिदं प्रदेयम्~॥}\ A}}\skt{ }\edtext{\skt{राज्ञा}}{\lemma{\skt{राज्ञा}}\Afootnote{K; \emph{om.}\ A; \skt{राज्ञे}\ H}}\skt{ }\edtext{\skt{देयं}}{\lemma{\skt{देयं}}\Afootnote{K H; \emph{om.}\ A}}\skt{ }\edtext{\skt{कृष्णतैलं}}{\lemma{\skt{कृष्णतैलं}}\Afootnote{K H; \emph{om.}\ A}}\skt{ }\edtext{\skt{प्रधाणं}}{\lemma{\skt{प्रधाणं}}\Afootnote{K; \emph{om.}\ A; \skt{प्रधानं~॥}\ H}}\skt{~।}\edtext{}{\lemma{\skt{सञ्जायन्ते… व}}\Afootnote{\emph{om.}\ A}}\pend

\pstart\versepar \ledsidenote{\passref{4.25.38}}\edtext{\skt{मनिःशिला}}{\lemma{\skt{मनिःशिला}}\Afootnote{K; \skt{ला}\ A; \skt{मनःशिला}\ H}}\skt{ }\edtext{\skt{हरितालं}}{\lemma{\skt{हरितालं}}\Afootnote{K H; \skt{क्षा रोध्रं द्वे}\ A}}\skt{ हरिद्रे }\edtext{\skt{लाक्षारोध्रं}}{\lemma{\skt{लाक्षारोध्रं}}\Afootnote{K; \skt{शिलाले कुष्ठं नागं}\ A; \skt{लाक्षालोध्रं}\ H}}\skt{ }\edtext{\skt{गैरिको}}{\lemma{\skt{गैरिको}}\Afootnote{K H; \skt{गैरिका}\ A}}\skt{ }\edtext{\skt{वर्ण्णकश्च}}{\lemma{\skt{वर्ण्णकश्च}}\Afootnote{K H; \skt{वर्णकाश्च~।}\ A}}\skt{~। }\pend
\pstart\versepar \skt{}\edtext{\skt{सुरा\-ष्ट्र\-जा\-कु\-ष्ठ\-म\-ञ्जि\-ष्ठ\-कल्काः}}{\lemma{\skt{सुरा\appbrk{}ष्ट्र\appbrk{}जा\appbrk{}कु\appbrk{}ष्ठ\appbrk{}म\appbrk{}ञ्जि\appbrk{}ष्ठ\appbrk{}कल्काः}}\Afootnote{K H; \skt{मञ्जिष्ठोग्रा स्यात्सुराष्ट्रोद्भवा}\ A}}\skt{ }\edtext{\skt{सप\-त्त\-ङ्गा\-रो\-चना}}{\lemma{\skt{सप\appbrk{}त्त\appbrk{}ङ्गा\appbrk{}रो\appbrk{}चना}}\Afootnote{K H; \skt{च पत्तङ्गं वै रोचना}\ A}}\skt{ }\edtext{\skt{कुङ्कुमश्च}}{\lemma{\skt{कुङ्कुमश्}}\Afootnote{K; \skt{चाञ्जनं}\ A; \skt{कुङ्कुमञ्}\ H}}\skt{~।}\pend

\pstart\versepar \ledsidenote{\passref{4.25.39}}\edtext{\skt{सुव\-र्ण्णा\-ह्वा\-पा\-ण्डु\-पत्रं}}{\lemma{\skt{सुव\appbrk{}र्ण्णा\appbrk{}ह्वा\appbrk{}पा\appbrk{}ण्डु\appbrk{}पत्रं}}\Afootnote{K H; \skt{हेमा\appbrk{}ङ्ग\appbrk{}त्व\appbrk{}क्पा\appbrk{}ण्डु\appbrk{}पत्रं}\ A}}\skt{ वटस्य }\edtext{\skt{कालेयकं}}{\lemma{\skt{कालेयकं}}\Afootnote{K H; \skt{कालीयं स्या}\ A}}\skt{ }\edtext{\skt{पद्मकं}}{\lemma{\skt{पद्मकं}}\Afootnote{K H; \skt{त्पद्मकं}\ A}}\skt{ पद्ममध्यं~। }\pend
\pstart\versepar \skt{}\edtext{\skt{कूचन्दनं}}{\lemma{\skt{कूचन्दनं}}\Afootnote{K; \skt{रक्तं श्वेतं}\ A; \skt{कू}\ H}}\skt{ चन्दनं }\edtext{\skt{पारतञ्च}}{\lemma{\skt{पारतञ्च}}\Afootnote{K H; \skt{पारदं च}\ A}}\skt{ काको\-ल्या\-दिः\-क्षी\-र\-पि\-ष्टश्च }\edtext{\skt{सर्वः}}{\lemma{\skt{सर्वः}}\Afootnote{K; \skt{वर्गः~॥}\ A; \skt{सर्व्व}\ H}}\pend

\pstart\versepar \ledsidenote{\passref{4.25.40}}\edtext{\skt{वसामज्जा}}{\lemma{\skt{वसामज्जा}}\Afootnote{K\textsuperscript{pc} H; \skt{मेदो मज्जा}\ A; \skt{सवसामज्जा}\ K\textsuperscript{ac}}}\skt{ }\edtext{\skt{मधूंच्छिष्टं}}{\lemma{\skt{मधूंच्छिष्टं}}\Afootnote{K; \skt{सिक्थकं गो}\ A; \skt{मधूच्छिष्टं}\ H}}\skt{ घृतञ्च दुग्धं क्वाथः क्षीरिणाञ्च द्रुमाणां~। }\pend
\pstart\versepar \skt{}\edtext{\skt{एत\-त्स\-र्व\-म्प\-क्व\-मे\-क\-ध्य\-तस्तु}}{\lemma{\skt{एकध्यतस्तु}}\Afootnote{K H; \skt{ऐकध्यतस्तु}\ A}}\skt{ }\edtext{\skt{मुखाभ्यङ्गे}}{\lemma{\skt{मुखाभ्यङ्गे}}\Afootnote{K H; \skt{वक्त्राभ्यङ्गे}\ A}}\skt{ }\edtext{\skt{सर्पिरुक्तं}}{\lemma{\skt{उक्तं}}\Afootnote{A K; \skt{उक्त}\ H}}\skt{ प्रधानं~।}\pend

\pstart\versepar \ledsidenote{\passref{4.25.41}}\edtext{\skt{हन्याद्व्यङ्गं}}{\lemma{\skt{हन्याद्व्यङ्गं}}\Afootnote{A K; \skt{हन्याद्व्यङ्गा}\ H}}\skt{ }\edtext{\skt{नीलिकां च}}{\lemma{\skt{नीलिकां च}}\Afootnote{K H; \skt{नीलिकां चा}\ A}}\skt{ }\edtext{\skt{प्रवृद्धां}}{\lemma{\skt{प्रवृद्धां}}\Afootnote{K H; \skt{तिवृद्धां}\ A}}\skt{ वक्त्रेजाताः }\edtext{\skt{पिटकायाश्च}}{\lemma{\skt{पिटकायाश्च}}\Afootnote{K H; \skt{स्फोटिकाश्चापि}\ A}}\skt{ काश्चित्~। }\pend
\pstart\versepar \skt{पद्माकारं निर्वलीकञ्च वक्त्रं कुर्या\-दे\-त\-त्पी\-न\-गण्डं मनोज्ञं~।}\edtext{}{\lemma{\skt{काश्चित्~। पद्माकारं}}\Afootnote{K H; \skt{काश्चित्~। पद्माकारं}\ A}}\pend

\pstart \ledsidenote{\passref{4.25.42}}\edtext{\skt{राज्ञा\-मे\-त\-द्यो\-षि\-ता\-न्चापि}}{\lemma{\skt{राज्ञा\appbrk{}मे\appbrk{}त\appbrk{}द्यो\appbrk{}षि\appbrk{}ता\appbrk{}न्चापि}}\Afootnote{K; \skt{राज्ञा\appbrk{}मे\appbrk{}त\appbrk{}द्यो\appbrk{}षितां चापि}\ A; \skt{राज्ञा\appbrk{}मे\appbrk{}त\appbrk{}द्यो\appbrk{}षि\appbrk{}ता\appbrk{}म्वापि}\ H}}\skt{ }\edtext{\skt{तेषां}}{\lemma{\skt{तेषां}}\Afootnote{K H; \skt{नित्यं}\ A}}\skt{ }\edtext{\edtext{\edtext{\skt{कुर्या\-द्वै\-द्य\-स्त\-त्स\-मा\-नाञ्च}}{\lemma{\skt{कुर्याद्वैद्यस्}}\Afootnote{A K; \skt{कुर्या}\ H}}}{\lemma{\skt{तत्समानाञ्}}\Afootnote{A K; \emph{om.}\ H}}}{\lemma{\skt{च}}\Afootnote{K; \skt{नृणां च~। कुष्ठघ्नं वै सर्पि\appbrk{}रे\appbrk{}त\appbrk{}त्प्र\appbrk{}धानं येषां पादे सन्ति वैपादिकाश्च~॥}\ A; \skt{दे वैद्य\appbrk{}स्त\appbrk{}त्स\appbrk{}मा\appbrk{}नाञ्च}\ H}}\skt{ }\edtext{\skt{नित्यम्}}{\lemma{\skt{नित्यम्}}\Afootnote{K H; \emph{om.}\ A}}\pend

\pstart \ledsidenote{\passref{4.25.trailer}}\edtext{\skt{इति}}{\lemma{\skt{इति}}\Afootnote{H; \skt{इति सुश्रु\appbrk{}त\appbrk{}सं\appbrk{}हि\appbrk{}तायां}\ A; \skt{इति~॥ चि~॥ {\symfont ❈}~॥}\ K}}\skt{ }\edtext{\skt{चिकित्सास्थान~॥ {\symfont ❈}}}{\lemma{\skt{चिकित्सास्थान~॥ {\symfont ❈}}}\Afootnote{\skt{चिकित्सास्थाने मिश्र\appbrk{}क\appbrk{}चि\appbrk{}कि\appbrk{}त्सितं नाम पञ्चविंशोऽध्यायः~॥२५~॥}\ A; \emph{om.}\ K; \skt{चिकित्सा}\ H}}\skt{~॥}\pend

\endnumbering
\beginnumbering
\eledsection*{Adhyāya 26}
\pstart \ledsidenote{\passref{4.26.1}}\skt{अथातः }\edtext{\skt{क्षीणवलीयं}}{\lemma{\skt{क्षीणवलीयं}}\Afootnote{K H; \skt{क्षीणबलीयं}\ A}}\skt{ }\edtext{\skt{व्याजी\-क\-र\-ण\-चि\-कि\-त्सितं}}{\lemma{\skt{व्याजी\appbrk{}क\appbrk{}र\appbrk{}ण\appbrk{}चि\appbrk{}कि\appbrk{}त्सितं}}\Afootnote{K H; \skt{वाजी\appbrk{}क\appbrk{}र\appbrk{}ण\appbrk{}चि\appbrk{}कि\appbrk{}त्सितं}\ A}}\skt{ }\edtext{\skt{व्याख्यास्यामः}}{\lemma{\skt{व्याख्यास्यामः}}\Afootnote{A H; \skt{व्या व्याख्यास्यामः~॥}\ K}}\skt{~॥}\pend

\pstart\versepar \ledsidenote{\passref{4.26.3}}\edtext{\skt{कल\-स्यो\-द\-ग्र\-व\-यसो}}{\lemma{\skt{कल\appbrk{}स्यो\appbrk{}द\appbrk{}ग्र\appbrk{}व\appbrk{}यसो}}\Afootnote{K H; \skt{कल्य\appbrk{}स्यो\appbrk{}द\appbrk{}ग्र\appbrk{}व\appbrk{}यसो}\ A}}\skt{ }\edtext{\skt{व्याजी\-क\-र\-ण\-से\-विनः}}{\lemma{\skt{व्याजी\appbrk{}क\appbrk{}र\appbrk{}ण\appbrk{}से\appbrk{}विनः}}\Afootnote{K; \skt{वाजी\appbrk{}क\appbrk{}र\appbrk{}ण\appbrk{}से\appbrk{}विनः~।}\ A; \skt{व्याजी\appbrk{}क\appbrk{}र\appbrk{}ण\appbrk{}से\appbrk{}विनः ः।}\ H}}\skt{~। }\pend
\pstart\versepar \skt{}\edtext{\skt{सर्वे\-ष्वृ\-तु\-ष्व\-ह\-रहः}}{\lemma{\skt{अहरहः}}\Afootnote{K H; \skt{अहरहर्}\ A}}\skt{ व्यवायो न }\edtext{\skt{निवारितः}}{\lemma{\skt{निवारितः}}\Afootnote{A K\textsuperscript{pc} H; \skt{निवातरितः~।}\ K\textsuperscript{ac}}}\skt{~।}\pend

\pstart\versepar \ledsidenote{\passref{4.26.4}}\edtext{\edtext{\skt{स्त्रीष्वक्षर्य्यं}}{\lemma{\skt{अक्षर्य्यं}}\Afootnote{K; \skt{अक्षार्य्यम्}\ H}}\skt{ मृगयतां }\edtext{\skt{वृद्धानाञ्च}}{\lemma{\skt{वृद्धानाञ्च}}\Afootnote{K H\textsuperscript{pc}; \skt{स्थविराणां}\ A; \skt{वृ}\ H\textsuperscript{ac}}}}{\lemma{\skt{… वृ}}\Afootnote{\emph{om.}\ A}}\skt{ }\edtext{\skt{रिरंसतां}}{\lemma{\skt{रिरंसतां}}\Afootnote{\skt{रिरंसूनां स्त्रीणां वाल्ल\appbrk{}भ्य\appbrk{}मि\appbrk{}च्छ\appbrk{}ताम्~। योषि\appbrk{}त्प्र\appbrk{}स\appbrk{}ङ्गा\appbrk{}त्क्षी\appbrk{}णानां}\ A; \skt{रिरांसतां~।}\ K; \skt{द्धानाञ्रिरंशतां~।}\ H\textsuperscript{ac}; \skt{रितं सतां~। रिरंशतां~।}\ H\textsuperscript{pc}}}\skt{~। }\pend
\pstart\versepar \skt{}\edtext{\edtext{\skt{क्लीवा\-णा\-म\-ल्प\-शु\-क्राणां}}{\lemma{\skt{क्लीवाणाम्}}\Afootnote{K; \skt{क्लीबानाम्}\ A; \skt{क्लीवानाम्}\ H}}}{\lemma{\skt{अल्पशुक्राणां}}\Afootnote{K H; \skt{अल्परेतसाम्~॥}\ A}}\skt{ }\edtext{\skt{स्त्रीषु}}{\lemma{\skt{स्त्रीषु}}\Afootnote{K H; \emph{om.}\ A}}\skt{ }\edtext{\skt{क्षीणाश्च ये}}{\lemma{\skt{क्षीणाश्च ये}}\Afootnote{K H; \emph{om.}\ A}}\skt{ }\edtext{\skt{नराः}}{\lemma{\skt{नराः}}\Afootnote{K H; \emph{om.}\ A}}\skt{~॥}\pend

\pstart\versepar \ledsidenote{\passref{4.26.5}}\skt{विला\-सि\-ना\-म\-र्थ\-वतां रूप\-यौ\-व\-न\-शा\-लिनां~। }\pend
\pstart\versepar \skt{}\edtext{\skt{नृणा\-म्व\-ह्वी\-प\-ती\-नाञ्च}}{\lemma{\skt{वह्वीपतीनाञ्च}}\Afootnote{K H; \skt{च ब}\ A}}\skt{ }\edtext{\skt{योगा}}{\lemma{\skt{योगा}}\Afootnote{K H; \skt{हुभार्याणां योगा}\ A}}\skt{ }\edtext{\skt{व्यजीकरो}}{\lemma{\skt{व्यजीकरो}}\Afootnote{K; \skt{वाजीकरा}\ A; \skt{व्याजीकरो}\ H}}\skt{ हिताः~।}\pend

\pstart\versepar \ledsidenote{\passref{4.26.7}}\edtext{\skt{भोजनानि च}}{\lemma{\skt{भोजनानि च}}\Afootnote{K H; \skt{भोजनानि वि}\ A}}\skt{ चित्राणि पानानि विविधानि च~। }\pend
\pstart\versepar \skt{वाचः }\edtext{\skt{स्रोत्रो }\edtext{\skt{नुगा\-मि\-न्य\-स्त्व\-क्सु\-ख\-स्प\-र्श\-नानि}}{\lemma{\skt{त्वक्सु\appbrk{}ख\appbrk{}स्प\appbrk{}र्श\appbrk{}नानि}}\Afootnote{K H; \skt{त्वचः स्पर्शसुखास्}\ A}}}{\lemma{\skt{स्रोत्रो नुगामिन्यस्}}\Afootnote{\skt{श्रोत्रा\appbrk{}नु\appbrk{}गा\appbrk{}मि\appbrk{}न्यस्}\ A; \skt{स्रोत्रो\appbrk{}नु\appbrk{}गा\appbrk{}मि\appbrk{}न्यस्}\ K; \skt{स्रोतो\appbrk{}नु\appbrk{}गा\appbrk{}मि\appbrk{}न्यस्}\ H}}\skt{ }\edtext{\skt{च}}{\lemma{\skt{च}}\Afootnote{K H; \skt{तथा~॥}\ A}}\skt{~।}\pend

\pstart\versepar \ledsidenote{\passref{4.26.9}}\edtext{\skt{गन्धान्मनोज्ञं}}{\lemma{\skt{गन्धान्मनोज्ञं}}\Afootnote{K H; \skt{गन्धा मनोज्ञा}\ A}}\skt{ रूपाणि चित्रा\-ण्यु\-प\-व\-नानि च~। }\pend
\pstart\versepar \skt{}\edtext{\skt{मन\-स\-श्चा\-प्र\-ती\-वातो}}{\lemma{\skt{चाप्रतीवातो}}\Afootnote{K H; \skt{चाप्रतीघातो}\ A}}\skt{ }\edtext{\skt{व्याजीकुर्वन्ति}}{\lemma{\skt{व्याजीकुर्वन्ति}}\Afootnote{K H; \skt{वाजीकुर्वन्ति}\ A}}\skt{ मानवं~। }\pend
\pstart\versepar \skt{तैस्तै\-र्भा\-वै\-र\-हृ\-द्यैस्तु रिरंसोर्मनसि क्षते~।}\pend

\pstart\versepar \ledsidenote{\passref{4.26.10}}\edtext{\skt{ध्वजः}}{\lemma{\skt{ध्वजः}}\Afootnote{K H; \skt{द्वेष्यस्त्रीसं}\ A}}\skt{ }\edtext{\edtext{\skt{पतत्यधो}}{\lemma{\skt{पतत्य्}}\Afootnote{K H; \skt{प्रयोगा}\ A}}}{\lemma{\skt{अधो}}\Afootnote{K H; \skt{च्च}\ A}}\skt{ }\edtext{\skt{नृणां}}{\lemma{\skt{नृणां}}\Afootnote{K H; \emph{om.}\ A}}\skt{ }\edtext{\skt{क्लैव्यं}}{\lemma{\skt{क्लैव्यं}}\Afootnote{K H; \skt{क्लैब्यं}\ A}}\skt{ }\edtext{\skt{समुपजायते}}{\lemma{\skt{समुपजायते}}\Afootnote{K H; \skt{तन्मान}\ A}}\skt{~। }\pend
\pstart\versepar \skt{}\edtext{\edtext{\skt{अन्नै\-र\-म्ल्वौ\-ष्ण\-ल\-व\-णै\-र\-ति\-मा\-त्रो\-प\-से\-वितैः}}{\lemma{\skt{अन्नैर्}}\Afootnote{K H; \skt{सं स्मृ}\ A}}}{\lemma{\skt{अम्ल्वौ\appbrk{}ष्ण\appbrk{}ल\appbrk{}व\appbrk{}णैर्}}\Afootnote{K; \skt{तम्~। कटु\appbrk{}का\appbrk{}म्लो\appbrk{}ष्ण\appbrk{}ल\appbrk{}व\appbrk{}णैर्}\ A; \skt{अम्त्वो\appbrk{}ष्ण\appbrk{}ल\appbrk{}व\appbrk{}णैर्}\ H}}\skt{~।}\pend

\pstart\versepar \ledsidenote{\passref{4.26.11}}\skt{सौम्यधातुक्षयो दृष्टः }\edtext{\skt{क्लैव्यं}}{\lemma{\skt{क्लैव्यं}}\Afootnote{K H; \skt{क्लैब्यं}\ A}}\skt{ तदपरं स्मृतं~। }\pend
\pstart\versepar \skt{}\edtext{\skt{अति\-व्य\-वा\-य\-शीलो वा न च}}{\lemma{\skt{अति\appbrk{}व्य\appbrk{}वा\appbrk{}य\appbrk{}शीलो वा न च}}\Afootnote{K H; \skt{अति\appbrk{}व्य\appbrk{}वा\appbrk{}य\appbrk{}शीलो यो न च}\ A}}\skt{ }\edtext{\skt{व्याजीक्रियारतः}}{\lemma{\skt{व्याजीक्रियारतः}}\Afootnote{K H; \skt{वाजीक्रियारतः~॥}\ A}}\skt{~।}\pend

\pstart\versepar \ledsidenote{\passref{4.26.12}}\skt{ध्वज\-भ\-ङ्ग\-म\-वा\-प्नोति }\edtext{\skt{सशु\-क्र\-क्ष\-य\-हे\-तुकं}}{\lemma{\skt{सशु\appbrk{}क्र\appbrk{}क्ष\appbrk{}य\appbrk{}हे\appbrk{}तुकं}}\Afootnote{K H; \skt{तच्छु\appbrk{}क्र\appbrk{}क्ष\appbrk{}य\appbrk{}हे\appbrk{}तु\appbrk{}कम्~।}\ A}}\skt{~। }\pend
\pstart\versepar \skt{}\edtext{\skt{मेढ्ररोगेण}}{\lemma{\skt{मेढ्ररोगेण}}\Afootnote{K H; \skt{महता मेढ्ररोगेण}\ A}}\skt{ }\edtext{\skt{महता}}{\lemma{\skt{महता}}\Afootnote{K H; \emph{om.}\ A}}\skt{ }\edtext{\skt{मर्म्चछेदेन वा}}{\lemma{\skt{मर्म्चछेदेन वा}}\Afootnote{K; \skt{मर्मच्छेदेन वा}\ A H}}\skt{ पुनः~।}\pend

\pstart \ledsidenote{\passref{4.26.13}}\edtext{\skt{क्लिव्यं}}{\lemma{\skt{क्लिव्यं}}\Afootnote{K H; \skt{क्लैब्यमेत}\ A}}\skt{ }\edtext{\edtext{\skt{चतुर्थम्भवति}}{\lemma{\skt{चतुर्थम्}}\Afootnote{K H; \skt{च्चतुर्थं}\ A}}}{\lemma{\skt{भवति}}\Afootnote{K H; \skt{स्यान्}\ A}}\skt{ }\edtext{\skt{नृणां}}{\lemma{\skt{नृणां}}\Afootnote{K H; \skt{नॄणां}\ A}}\skt{ }\edtext{\skt{पुंस्त्वोपघातजं}}{\lemma{\skt{पुंस्त्वोपघातजं}}\Afootnote{K H; \skt{पुंस्त्वो\appbrk{}प\appbrk{}घा\appbrk{}त\appbrk{}जम्~। जन्मप्रभृति यः क्लीबः क्लैब्यं तत्सहजं स्मृतम्~॥}\ A}}\skt{~।}\pend

\pstart\versepar \ledsidenote{\passref{4.26.15}}\skt{असाध्यं सहजं }\edtext{\skt{क्लैव्यं}}{\lemma{\skt{क्लैव्यं}}\Afootnote{K H; \skt{क्लैब्यं}\ A}}\skt{ मर्मच्छेदाच्च यद्भवेत्~। }\pend
\pstart\versepar \skt{}\edtext{\skt{साध्या\-ना\-म\-व\-शि\-ष्टानां}}{\lemma{\skt{अवशिष्टानां}}\Afootnote{K H; \skt{इतरेषां तु}\ A}}\skt{ कार्यो }\edtext{\skt{व्याजीकरो}}{\lemma{\skt{व्याजीकरो}}\Afootnote{K H; \emph{om.}\ A}}\skt{ }\edtext{\skt{विधिः}}{\lemma{\skt{विधिः}}\Afootnote{K H; \skt{हेतुविपर्ययः~॥}\ A}}\skt{~।}\pend

\pstart \ledsidenote{\passref{4.26.16}}\edtext{\skt{तिल\-मा\-ष\-वि\-दा\-रीणां}}{\lemma{\skt{तिल\appbrk{}मा\appbrk{}ष\appbrk{}वि\appbrk{}दा\appbrk{}रीणां}}\Afootnote{K H; \skt{विधिर्वाजीकरो यस्तु तं प्रवक्ष्याम्यतः परम्~। तिल\appbrk{}मा\appbrk{}ष\appbrk{}वि\appbrk{}दा\appbrk{}रीणां}\ A}}\skt{ शालीनां चूर्ण्णमेव }\edtext{\skt{च}}{\lemma{\skt{च}}\Afootnote{K H; \skt{वा~॥}\ A}}\skt{~।}\pend

\pstart\versepar \ledsidenote{\passref{4.26.17}}\edtext{\edtext{\edtext{\skt{रसै\-रि\-क्षु\-र\-सै\-र्वापि}}{\lemma{\skt{रसैर्}}\Afootnote{K H; \skt{पौण्}\ A}}}{\lemma{\skt{इक्षुरसैर्}}\Afootnote{K H; \skt{ड्रकेक्षुरसै}\ A}}}{\lemma{\skt{वापि}}\Afootnote{K H; \skt{रार्द्रं}\ A}}\skt{ मर्दितं }\edtext{\skt{सैन्धवान्वित}}{\lemma{\skt{सैन्धवान्वित}}\Afootnote{K; \skt{सैन्धवान्वितम्}\ A H}}\skt{~। }\pend
\pstart\versepar \skt{वराहमेदसा }\edtext{\skt{युक्तं}}{\lemma{\skt{युक्तं}}\Afootnote{K H; \skt{युक्तां}\ A}}\skt{ }\edtext{\skt{घृतेनोत्कारिकां}}{\lemma{\skt{घृतेनोत्कारिकां}}\Afootnote{A K; \skt{घृते\appbrk{}नो\appbrk{}त्का\appbrk{}रि\appbrk{}कम्}\ H}}\skt{ पचेत्~।}\pend

\pstart\versepar \ledsidenote{\passref{4.26.18}}\skt{तां भक्षयित्वा }\edtext{\skt{गच्छेयुः}}{\lemma{\skt{गच्छेयुः}}\Afootnote{K H; \emph{om.}\ A}}\skt{ पुरुषः }\edtext{\edtext{\skt{षष्ठिमङ्गनां}}{\lemma{\skt{षष्ठिम्}}\Afootnote{K H; \skt{गच्छेत्तु प्रम}\ A}}}{\lemma{\skt{अङ्गनां}}\Afootnote{K H; \skt{दाशतम्~।}\ A}}\skt{~॥ }\pend
\pstart\versepar \skt{}\edtext{\skt{वस्त्रा\-ण्ड\-सि\-द्ध\-प\-यसि}}{\lemma{\skt{वस्त्रा\appbrk{}ण्ड\appbrk{}सि\appbrk{}द्ध\appbrk{}प\appbrk{}यसि}}\Afootnote{K; \skt{बस्ताण्डसिद्धे पयसि}\ A; \skt{वस्त्राण्डसिद्धं पयसि}\ H}}\skt{ भाविता न }\edtext{\skt{सकृत्तिलां}}{\lemma{\skt{सकृत्तिलां}}\Afootnote{\skt{सकृत्तिलान्~॥}\ A; \skt{सकृत्तिलां}\ K H}}\skt{~।}\pend

\pstart\versepar \ledsidenote{\passref{4.26.19}}\edtext{\skt{शिशु\-मा\-र\-व\-सा\-पक्वा }\edtext{\skt{तैस्तिलैः}}{\lemma{\skt{तिलैः}}\Afootnote{K H; \emph{om.}\ A}}}{\lemma{\skt{शिशु\appbrk{}मा\appbrk{}र\appbrk{}व\appbrk{}सा\appbrk{}पक्वा तैस्}}\Afootnote{K; \skt{शिशु\appbrk{}मा\appbrk{}र\appbrk{}व\appbrk{}सा\appbrk{}पक्वाः}\ A; \skt{शिशु\appbrk{}मा\appbrk{}र\appbrk{}व\appbrk{}सा\appbrk{}पक्व तैस्}\ H}}\skt{ शष्कुलीशुभां~। }\pend
\pstart\versepar \skt{यः }\edtext{\skt{खादेत्स}}{\lemma{\skt{खादेत्स}}\Afootnote{K H; \skt{स्वादेत्स}\ A}}\skt{ }\edtext{\skt{पुमां}}{\lemma{\skt{पुमां}}\Afootnote{K; \skt{पुमा}\ A H}}\skt{ }\edtext{\skt{गच्छेत्स्त्रीणां}}{\lemma{\skt{गच्छेत्स्त्रीणां}}\Afootnote{K; \skt{न्गच्छेत्स्त्रीणां}\ A H}}\skt{ शत\-म\-पू\-र्व\-वत्~॥}\edtext{}{\lemma{\skt{शष्कुलीशुभां~। यः}}\Afootnote{K; \skt{शष्कुल्यस्तिस्तैलैः कृताः~। यः}\ A; \skt{शष्कुलीशुभाः~॥ यः}\ H}}\pend

\pstart\versepar \ledsidenote{\passref{4.26.20}}\edtext{\skt{पिप्प\-ली\-ल\-व\-णौ\-पेतौ}}{\lemma{\skt{पिप्प\appbrk{}ली\appbrk{}ल\appbrk{}व\appbrk{}णौ\appbrk{}पेतौ}}\Afootnote{K; \skt{पिप्प\appbrk{}ली\appbrk{}ल\appbrk{}व\appbrk{}णो\appbrk{}पेते}\ A; \skt{पिप्प\appbrk{}ली\appbrk{}ल\appbrk{}व\appbrk{}णो\appbrk{}पेतौ}\ H}}\skt{ }\edtext{\skt{वस्ताण्डौ}}{\lemma{\skt{वस्ताण्डौ}}\Afootnote{K H; \skt{बस्ताण्डे}\ A}}\skt{ क्षीरसर्पिषि~। }\pend
\pstart\versepar \skt{}\edtext{\skt{साधितौ}}{\lemma{\skt{साधितौ}}\Afootnote{K H; \skt{साधिते}\ A}}\skt{ भक्षयेद्यस्तु स गच्छे\-त्प्र\-म\-दा\-शतं~॥}\pend

\pstart\versepar \ledsidenote{\passref{4.26.21}}\edtext{\skt{माष\-पि\-प्प\-लि\-शा\-लीनां}}{\lemma{\skt{माष\appbrk{}पि\appbrk{}प्प\appbrk{}लि\appbrk{}शा\appbrk{}लीनां}}\Afootnote{K H; \skt{पिप्प\appbrk{}ली\appbrk{}मा\appbrk{}ष\appbrk{}शा\appbrk{}लीनां}\ A}}\skt{ यव\-गो\-धू\-म\-यो\-स्तथा~। }\pend
\pstart\versepar \skt{चूर्ण्णभागैः }\edtext{\skt{समैस्तैस्तु}}{\lemma{\skt{समैस्तैस्तु}}\Afootnote{A K; \skt{समस्तैस्तु}\ H}}\skt{ घृते पूपलिकां पचेत्~।}\pend

\pstart\versepar \ledsidenote{\passref{4.26.22}}\skt{तां भक्षयित्वा }\edtext{\skt{पीत्वा च}}{\lemma{\skt{पीत्वा च}}\Afootnote{K H; \skt{पीत्वा तु}\ A}}\skt{ शर्क\-रा\-म\-धु\-र\-म्पयः~। }\pend
\pstart\versepar \skt{नर\-श्च\-ट\-क\-व\-द्ग\-च्छे\-द्द\-श\-वा\-रा\-न्नि\-र\-न्तरं~॥}\pend

\pstart\versepar \ledsidenote{\passref{4.26.23}}\edtext{\skt{चूर्ण्णं}}{\lemma{\skt{चूर्ण्णं}}\Afootnote{K H; \emph{om.}\ A}}\skt{ विदार्याः }\edtext{\skt{सकृतं}}{\lemma{\skt{सकृतं}}\Afootnote{K; \skt{सुकृतं}\ A H}}\skt{ }\edtext{\skt{स्वरसेनैव}}{\lemma{\skt{स्वरसेनैव}}\Afootnote{K H; \skt{चूर्णं स्वरसेनैव}\ A}}\skt{ भावितं~। }\pend
\pstart\versepar \skt{}\edtext{\skt{सर्पिःक्षौद्रयुतं}}{\lemma{\skt{सर्पिःक्षौद्रयुतं}}\Afootnote{K H; \skt{सर्पिर्मधुयुतं}\ A}}\skt{ लीढ्वा }\edtext{\skt{दश\-ना\-र्यो\-धि\-रो\-हित}}{\lemma{\skt{दश\appbrk{}ना\appbrk{}र्यो\appbrk{}धि\appbrk{}रो\appbrk{}हित}}\Afootnote{K; \skt{दश स्त्रीरधिगच्छति~॥}\ A; \skt{दश\appbrk{}ना\appbrk{}र्योऽ\appbrk{}धि\appbrk{}रो\appbrk{}हितः~॥}\ H}}\skt{~॥}\pend

\pstart\versepar \ledsidenote{\passref{4.26.24}}\skt{एवमामलकं चूर्ण्णं स्वरसेनैव भावितं~। }\pend
\pstart\versepar \skt{}\edtext{\skt{शर्क\-रा\-म\-धु\-स\-र्पि\-र्भि\-र्युक्तं}}{\lemma{\skt{शर्क\appbrk{}रा\appbrk{}म\appbrk{}धु\appbrk{}स\appbrk{}र्पि\appbrk{}र्भिर्}}\Afootnote{K; \skt{शर्क\appbrk{}रा\appbrk{}म\appbrk{}धु\appbrk{}स\appbrk{}र्पि\appbrk{}भिर्}\ A; \skt{शर्क्क\appbrk{}रा\appbrk{}म\appbrk{}धु\appbrk{}स\appbrk{}र्प्पि\appbrk{}र्व्भिर्}\ H}}\skt{ लीढ्वा पयः }\edtext{\skt{पिवेत्}}{\lemma{\skt{पिवेत्}}\Afootnote{K H; \skt{पिबेत्~॥}\ A}}\skt{~।}\pend

\pstart\versepar \ledsidenote{\passref{4.26.25}}\edtext{\skt{एते\-ना\-शी\-ति\-वर्षो}}{\lemma{\skt{एते\appbrk{}ना\appbrk{}शी\appbrk{}ति\appbrk{}वर्षो}}\Afootnote{A K H\textsuperscript{ac}; \skt{एते\appbrk{}ना\appbrk{}शी\appbrk{}ति\appbrk{}व\appbrk{}र्षा\appbrk{}यु\appbrk{}र्भ\appbrk{}वेच्च वर्षो}\ H\textsuperscript{pc}}}\skt{ पि }\edtext{\skt{युवेव}}{\lemma{\skt{युवेव}}\Afootnote{A K; \skt{युते च}\ H}}\skt{ परिहृष्यति~॥ }\pend
\pstart\versepar \skt{}\edtext{\skt{पिप्प\-ली\-ल\-व\-णौ\-पेतौ}}{\lemma{\skt{पिप्प\appbrk{}ली\appbrk{}ल\appbrk{}व\appbrk{}णौ\appbrk{}पेतौ}}\Afootnote{K; \skt{पिप्प\appbrk{}ली\appbrk{}ल\appbrk{}व\appbrk{}णो\appbrk{}पेते}\ A; \skt{पिप्प\appbrk{}ली\appbrk{}ल\appbrk{}व\appbrk{}णो\appbrk{}पेतौ}\ H}}\skt{ }\edtext{\skt{वस्ताण्डौ}}{\lemma{\skt{वस्ताण्डौ}}\Afootnote{K H; \skt{बस्ताण्डे}\ A}}\skt{ }\edtext{\skt{घृतसाधितौ}}{\lemma{\skt{घृतसाधितौ}}\Afootnote{K H; \skt{घृतसाधिते~॥}\ A}}\skt{~।}\pend

\pstart\versepar \ledsidenote{\passref{4.26.26}}\skt{शिशुमारस्य वा }\edtext{\skt{खादेत्तौ तु}}{\lemma{\skt{खादेत्तौ तु}}\Afootnote{K H; \skt{खादेत्ते तु}\ A}}\skt{ }\edtext{\skt{व्याजीकरौ}}{\lemma{\skt{व्याजीकरौ}}\Afootnote{K H; \skt{वाजीकरे}\ A}}\skt{ भृशं~॥ }\pend
\pstart\versepar \skt{कुली\-र\-कू\-र्म\-न\-क्रा\-णा\-म\-ण्डा\-न्येवं तु भक्षयेत्~।}\pend

\pstart\versepar \ledsidenote{\passref{4.26.27}}\edtext{\edtext{\skt{महि\-ष\-र्ष\-भ\-व\-स्ता\-ना\-म्पि\-वे\-च्छु\-क्राणि}}{\lemma{\skt{महि\appbrk{}ष\appbrk{}र्ष\appbrk{}भ\appbrk{}व\appbrk{}स्ता\appbrk{}नाम्}}\Afootnote{H; \skt{महि\appbrk{}ष\appbrk{}र्ष\appbrk{}भ\appbrk{}ब\appbrk{}स्तानां}\ A; \skt{महिषर्षभ}\ K}}}{\lemma{\skt{पिवेच्छुक्राणि}}\Afootnote{H; \skt{पिबेच्छुक्राणि}\ A; \emph{om.}\ K}}\skt{ }\edtext{\skt{वानरः}}{\lemma{\skt{वानरः}}\Afootnote{A H; \emph{om.}\ K}}\skt{~॥ }\pend
\pstart\versepar \skt{}\edtext{\skt{अश्व\-त्थ\-फ\-ल\-मू\-लत्वं}}{\lemma{\skt{अश्व\appbrk{}त्थ\appbrk{}फ\appbrk{}ल\appbrk{}मू\appbrk{}लत्वं}}\Afootnote{H; \skt{अश्व\appbrk{}त्थ\appbrk{}फ\appbrk{}ल\appbrk{}मू\appbrk{}लत्व}\ A; \emph{om.}\ K}}\skt{ }\edtext{\edtext{\edtext{\skt{छुङ्ग\-सि\-द्ध\-म्प\-योऽपि}}{\lemma{\skt{छुङ्गसिद्धम्}}\Afootnote{H; \skt{क्शुङ्गासिद्धं}\ A; \emph{om.}\ K}}}{\lemma{\skt{पयो}}\Afootnote{A H; \emph{om.}\ K}}}{\lemma{\skt{ऽपि}}\Afootnote{H; \skt{नरः~॥}\ A; \emph{om.}\ K}}\skt{ }\edtext{\skt{वा}}{\lemma{\skt{वा}}\Afootnote{H; \emph{om.}\ A K}}\skt{~।}\pend

\pstart\versepar \ledsidenote{\passref{4.26.28}}\skt{पीत्वा सशर्क्कराक्षौद्रं कुलिङ्ग इव हृष्यति~॥ }\pend
\pstart\versepar \skt{}\edtext{\skt{विदा\-री\-मू\-ल\-क\-ल्कन्तु}}{\lemma{\skt{विदा\appbrk{}री\appbrk{}मू\appbrk{}ल\appbrk{}क\appbrk{}ल्कन्तु}}\Afootnote{H; \skt{विदा\appbrk{}रि\appbrk{}मू\appbrk{}ल\appbrk{}कल्कं तु}\ A}}\skt{ }\edtext{\skt{स्रितेन}}{\lemma{\skt{स्रितेन}}\Afootnote{H; \skt{शृतेन}\ A}}\skt{ पयसा नरः~।}\pend

\pstart\versepar \ledsidenote{\passref{4.26.29}}\edtext{\skt{उदु\-म्व\-र\-र\-स\-म्पीत्वा}}{\lemma{\skt{उदु\appbrk{}म्व\appbrk{}र\appbrk{}र\appbrk{}सम्}}\Afootnote{H; \skt{उडुम्बरसमं}\ A}}\skt{ वृद्धोऽपि तरुणायते~॥ }\pend
\pstart\versepar \skt{माषा\-णा\-म्प\-ल\-मे\-कन्तु }\edtext{\skt{संयु\-क्त\-म्म\-धु\-स\-र्प्पिषा}}{\lemma{\skt{मधुसर्प्पिषा}}\Afootnote{H; \skt{क्षौद्रसर्पिषा~॥}\ A}}\skt{~।}\pend

\pstart \ledsidenote{\passref{4.26.30}}\edtext{\skt{तिं}}{\lemma{\skt{तिं}}\Afootnote{H; \skt{अवलिह्य पयः}\ A}}\skt{ }\edtext{\skt{लीढ्वा\-नु\-पि\-वे\-त्क्षीरं}}{\lemma{\skt{लीढ्वा\appbrk{}नु\appbrk{}पि\appbrk{}वे\appbrk{}त्क्षीरं}}\Afootnote{H; \skt{पीत्वा}\ A}}\skt{ तेन }\edtext{\skt{व्याजीभवेत्नरः}}{\lemma{\skt{व्याजीभवेत्नरः}}\Afootnote{H; \skt{वाजी भवेन्नरः~। क्षीरपक्वांस्तु गोधू\appbrk{}मा\appbrk{}ना\appbrk{}त्म\appbrk{}गु\appbrk{}प्ता\appbrk{}फलैः सह~॥}\ A}}\skt{~॥}\pend

\pstart \ledsidenote{\passref{4.26.31}}\edtext{\skt{नक्र\-मू\-षि\-क\-म\-ण्डू\-क\-च\-ट\-का\-ण्ड\-कृतं}}{\lemma{\skt{नक्र\appbrk{}मू\appbrk{}षि\appbrk{}क\appbrk{}म\appbrk{}ण्डू\appbrk{}क\appbrk{}च\appbrk{}ट\appbrk{}का\appbrk{}ण्ड\appbrk{}कृतं}}\Afootnote{H; \skt{शीता\appbrk{}न्घृ\appbrk{}त\appbrk{}यु\appbrk{}ता\appbrk{}न्खा\appbrk{}दे\appbrk{}त्ततः पश्चात्पयः पिबेत्~। नक्र\appbrk{}मू\appbrk{}षि\appbrk{}क\appbrk{}म\appbrk{}ण्डू\appbrk{}क\appbrk{}च\appbrk{}ट\appbrk{}का\appbrk{}ण्ड\appbrk{}कृतं}\ A}}\skt{ घृतं~।}\pend

\pstart\versepar \ledsidenote{\passref{4.26.32}}\skt{पादाभ्यङ्गेन कुरुते }\edtext{\skt{व्याजीभूमिन्तु न}}{\lemma{\skt{व्याजीभूमिन्तु न}}\Afootnote{H; \skt{बलं भूमिं तु न}\ A}}\skt{ स्पृशेत्~॥ }\pend
\pstart\versepar \skt{यावन्न }\edtext{\skt{स्पृशते}}{\lemma{\skt{स्पृशते}}\Afootnote{H; \skt{त्स्पृशति नो}\ A}}\skt{ भूमि\-न्ता\-व\-द्ग\-च्छे\-न्नि\-र\-न्तरं~।}\edtext{}{\lemma{\skt{स्पृशेत्~॥ यावन्न}}\Afootnote{H; \skt{स्पृशेत्~। याव}\ A}}\pend

\pstart\versepar \ledsidenote{\passref{4.26.33}}\skt{स्वयं }\edtext{\edtext{\skt{गुप्ते\-क्षु\-र\-क\-यो\-र्वी\-ज\-पूर्ण्णं}}{\lemma{\skt{गुप्ते\appbrk{}क्षु\appbrk{}र\appbrk{}क\appbrk{}योर्}}\Afootnote{A; \skt{गुप्तेक्षुरकयो}\ H}}}{\lemma{\skt{वीजपूर्ण्णं}}\Afootnote{\skt{फलचूर्णं}\ A; \skt{र्व्वीजपूर्ण्णं}\ H}}\skt{ सशर्क्करं~॥ }\pend
\pstart\versepar \skt{}\edtext{\skt{धारोष्णेन}}{\lemma{\skt{धारोष्णेन}}\Afootnote{A H\textsuperscript{pc}; \skt{धारोष्णे}\ H\textsuperscript{ac}}}\skt{ }\edtext{\skt{नर}}{\lemma{\skt{नर}}\Afootnote{H\textsuperscript{pc}; \skt{न}\ A H\textsuperscript{ac}}}\skt{ }\edtext{\skt{पयः}}{\lemma{\skt{पयः}}\Afootnote{H; \skt{रः}\ A}}\skt{ पीत्वा पयसा न क्षयं व्रजेत्~।}\pend

\pstart\versepar \ledsidenote{\passref{4.26.34}}\skt{उच्च\-टा\-चू\-र्ण्ण\-म\-प्येव }\edtext{\skt{क्षीरे\-णो\-त्त\-म\-मु\-च्यते}}{\lemma{\skt{उच्यते}}\Afootnote{H\textsuperscript{pc}; \skt{उच्य}\ H\textsuperscript{ac}}}\skt{~॥ }\pend
\pstart\versepar \skt{}\edtext{\edtext{\edtext{\skt{शता\-व\-र्यु\-च्च\-टा\-चू\-र्ण्ण\-म्पयः}}{\lemma{\skt{शता\appbrk{}व\appbrk{}र्यु\appbrk{}च्च\appbrk{}टा\appbrk{}चू\appbrk{}र्ण्णम्}}\Afootnote{\skt{शता\appbrk{}व\appbrk{}र्यु\appbrk{}च्च\appbrk{}टा\appbrk{}चू\appbrk{}र्ण्णम्}\ H\textsuperscript{pc}}}}{\lemma{\skt{पयः}}\Afootnote{\skt{पेयम्}\ A; \skt{ते~॥}\ H\textsuperscript{ac}; \skt{पेयपयः}\ H\textsuperscript{pc}}}\skt{ }\edtext{\skt{मेदो}}{\lemma{\skt{मेदो}}\Afootnote{H\textsuperscript{pc}; \skt{एवं}\ A}}}{\lemma{\skt{ते… उच्चटाचूर्ण्णम्}}\Afootnote{\emph{om.}\ H\textsuperscript{ac}}}\skt{ }\edtext{\skt{सुखार्थिना}}{\lemma{\skt{सुखार्थिना}}\Afootnote{H\textsuperscript{pc}; \skt{बलार्थिना~।}\ A}}\skt{~। }\pend
\pstart\versepar \skt{}\edtext{\skt{स्वयं\-गु\-प्ता\-फ\-ल\-युतं}}{\lemma{\skt{स्वयं\appbrk{}गु\appbrk{}प्ता\appbrk{}फ\appbrk{}ल\appbrk{}युतं}}\Afootnote{H; \skt{स्वयं\appbrk{}गु\appbrk{}प्ता\appbrk{}फ\appbrk{}लै\appbrk{}र्युक्तं}\ A}}\skt{ }\edtext{\edtext{\skt{माष\-यू\-ष\-म्पि\-वे\-त्नरः}}{\lemma{\skt{माषयूषम्}}\Afootnote{H; \skt{माषसूपं}\ A}}}{\lemma{\skt{पिवेत्नरः}}\Afootnote{H; \skt{पिबेन्नरः~॥}\ A}}\skt{~॥}\edtext{}{\lemma{\skt{… उ}}\Afootnote{\emph{om.}\ A}}\edtext{}{\lemma{\skt{पयः… स्व}}\Afootnote{\emph{om.}\ H\textsuperscript{ac}}}\pend

\pstart\versepar \ledsidenote{\passref{4.26.37}}\skt{गृष्टीनां वृद्धवत्सानां }\edtext{\skt{माष\-चू\-र्ण्ण\-भृ\-ता\-ङ्गवां}}{\lemma{\skt{माष\appbrk{}चू\appbrk{}र्ण्ण\appbrk{}भृ\appbrk{}ता\appbrk{}ङ्गवां}}\Afootnote{H; \skt{माषपर्णभृतां गवाम्~।}\ A}}\skt{~। }\pend
\pstart\versepar \skt{यत्क्षी\-र\-न्त\-त्प्र\-श\-न्सन्ति }\edtext{\skt{बलकामेषु}}{\lemma{\skt{बलकामेषु}}\Afootnote{A; \skt{वलकामेषु}\ H}}\skt{ जन्तुषु~॥}\pend

\pstart\versepar \ledsidenote{\passref{4.26.39}}\skt{एते }\edtext{\skt{व्याजीकरा}}{\lemma{\skt{व्याजीकरा}}\Afootnote{H; \skt{वाजीकरा}\ A}}\skt{ योगाः }\edtext{\skt{प्रीत्य\-प\-त्य\-ब\-ल\-प्रदाः}}{\lemma{\skt{प्रीत्य\appbrk{}प\appbrk{}त्य\appbrk{}ब\appbrk{}ल\appbrk{}प्रदाः}}\Afootnote{A; \skt{प्रीत्य\appbrk{}प\appbrk{}त्य\appbrk{}व\appbrk{}ल\appbrk{}प्रदाः~।}\ H}}\skt{~। }\pend
\pstart\versepar \skt{सेव्या }\edtext{\skt{विशुद्धापचित}}{\lemma{\skt{विशुद्धापचित}}\Afootnote{H\textsuperscript{pc}; \skt{विशुद्धो}\ A; \emph{om.}\ H\textsuperscript{ac}}}\skt{ }\edtext{\skt{विशुद्धा}}{\lemma{\skt{विशुद्धा}}\Afootnote{H; \emph{om.}\ A}}\skt{ }\edtext{\skt{परिचित}}{\lemma{\skt{परिचित}}\Afootnote{H; \skt{पचित}\ A}}\skt{ देहैः }\edtext{\edtext{\skt{काला\-द्य\-पे\-क्ष\-येति}}{\lemma{\skt{कालाद्य्}}\Afootnote{H; \skt{काला}\ A}}}{\lemma{\skt{अपेक्षयेति}}\Afootnote{H; \skt{द्यपेक्षया~॥}\ A}}\skt{~॥}\pend

\endnumbering
\beginnumbering
\eledsection*{Adhyāya 27}
\pstart \ledsidenote{\passref{4.27.1}}\skt{अथातः }\edtext{\skt{सर्व्वा\-वा\-ध\-सं\-श\-म\-नीयं}}{\lemma{\skt{सर्व्वा\appbrk{}वा\appbrk{}ध\appbrk{}सं\appbrk{}श\appbrk{}म\appbrk{}नीयं}}\Afootnote{H; \skt{सर्वो\appbrk{}प\appbrk{}घा\appbrk{}त\appbrk{}श\appbrk{}म\appbrk{}नीयं}\ A}}\skt{ रसा\-य\-न\-म्व्या\-ख्या\-स्यामः~॥}\pend

\pstart\versepar \ledsidenote{\passref{4.27.3}}\skt{पूर्व्वे वयसि मध्ये वा मनुष्यस्य रसायनं~। }\pend
\pstart\versepar \skt{प्रयुञ्जीत भिषक्प्राज्ञः स्निग्धशुद्धतनोः सदा~॥}\pend

\pstart\versepar \ledsidenote{\passref{4.27.4}}\edtext{\skt{अवि\-शु\-द्ध\-श\-री\-रस्य}}{\lemma{\skt{अवि\appbrk{}शु\appbrk{}द्ध\appbrk{}श\appbrk{}री\appbrk{}रस्य}}\Afootnote{H; \skt{नावि\appbrk{}शु\appbrk{}द्ध\appbrk{}श\appbrk{}री\appbrk{}रस्य}\ A}}\skt{ युक्तो रासायनो विधिः~। }\pend
\pstart\versepar \skt{न भाति वाससि क्लिष्टे रङ्गयोग इवाहितः~॥}\pend

\pstart\versepar \ledsidenote{\passref{4.27.5}}\edtext{\skt{शरी\-र\-स्यो\-प\-घा\-ताय}}{\lemma{\skt{शरी\appbrk{}र\appbrk{}स्यो\appbrk{}प\appbrk{}घा\appbrk{}ताय}}\Afootnote{H; \skt{शरी\appbrk{}र\appbrk{}स्यो\appbrk{}प\appbrk{}घाता ये}\ A}}\skt{ दोषजा मानसास्तथा~। }\pend
\pstart\versepar \skt{}\edtext{\skt{उप\-दि\-ष्टो\-प\-दे\-शेन}}{\lemma{\skt{उप\appbrk{}दि\appbrk{}ष्टो\appbrk{}प\appbrk{}दे\appbrk{}शेन}}\Afootnote{H; \skt{उपदिष्टाः प्रदेशेषु}\ A}}\skt{ तेषां वक्ष्यामि वारणं~॥}\pend

\pstart\versepar \ledsidenote{\passref{4.27.6}}\skt{शीतोदकम्पयः क्षौद्रं }\edtext{\skt{घृतमेकैकशो}}{\lemma{\skt{घृतमेकैकशो}}\Afootnote{H; \skt{सर्पिरित्येकशो}\ A}}\skt{ द्विशः~। }\pend
\pstart\versepar \skt{त्रिशः }\edtext{\edtext{\skt{समस्तमथवा}}{\lemma{\skt{समस्तम्}}\Afootnote{H; \skt{समस्त}\ A}}}{\lemma{\skt{अथवा}}\Afootnote{H; \skt{मथवा}\ A}}\skt{ प्राक्पीतं स्थापयेद्वयः~॥}\pend

\pstart \ledsidenote{\passref{4.27.7}}\skt{तत्र }\edtext{\edtext{\skt{विड\-ङ्ग\-त\-ण्डु\-ल\-मा\-हृत्य}}{\lemma{\skt{विड\appbrk{}ङ्ग\appbrk{}त\appbrk{}ण्डु\appbrk{}लम्}}\Afootnote{H; \skt{विडङ्गतण्डुल}\ A}}}{\lemma{\skt{आहृत्य}}\Afootnote{H; \skt{चूर्णमाहृत्य}\ A}}\skt{ यष्टी\-म\-धु\-क\-म\-धु\-युक्तं }\edtext{\skt{यथाबलं}}{\lemma{\skt{यथाबलं}}\Afootnote{A; \skt{यथावलं}\ H}}\skt{ }\edtext{\skt{शीततोयेन}}{\lemma{\skt{शीततोयेन}}\Afootnote{H; \skt{शीततोयेनो}\ A}}\skt{ }\edtext{\skt{उपयुञ्जीत}}{\lemma{\skt{उपयुञ्जीत}}\Afootnote{H; \skt{पयुञ्जीत}\ A}}\skt{ शीततोयं }\edtext{\skt{चानुपिवेत्~। अह\-र\-ह\-र्म्मा\-स\-न्त\-देव}}{\lemma{\skt{चानुपिवेत्~। अहरहर्}}\Afootnote{H; \skt{चानु\appbrk{}पि\appbrk{}बे\appbrk{}दे\appbrk{}व\appbrk{}म\appbrk{}ह\appbrk{}र\appbrk{}हर्}\ A}}\skt{ मधुयुक्तं }\edtext{\skt{भल्ला\-त\-क\-क्वा\-थेन}}{\lemma{\skt{भल्ला\appbrk{}त\appbrk{}क\appbrk{}क्वा\appbrk{}थेन}}\Afootnote{H; \skt{भल्ला\appbrk{}त\appbrk{}क\appbrk{}क्वा\appbrk{}थेन वा,}\ A}}\skt{~। }\edtext{\skt{द्राक्षाक्वाथेन वा}}{\lemma{\skt{द्राक्षाक्वाथेन वा}}\Afootnote{H; \skt{मधुद्राक्षाक्वा}\ A}}\skt{ }\edtext{\skt{मधुसंयुक्तं}}{\lemma{\skt{मधुसंयुक्तं}}\Afootnote{H; \skt{थयुक्तं वा,}\ A}}\skt{~। मध्वा\-म\-ल\-क\-र\-सा\-भ्याम्वा~। गुडूचीक्वाथेन वा~। एवमेते पञ्चप्रयोगा भवन्ति~। जीर्ण्णे }\edtext{\skt{मुद्ग\-यू\-षे\-णा\-ल\-व\-णेन}}{\lemma{\skt{मुद्ग\appbrk{}यू\appbrk{}षे\appbrk{}णा\appbrk{}ल\appbrk{}व\appbrk{}णेन}}\Afootnote{H; \skt{मुद्गा\appbrk{}म\appbrk{}ल\appbrk{}क\appbrk{}यू\appbrk{}षे\appbrk{}णा\appbrk{}ल\appbrk{}व\appbrk{}णेना}\ A}}\skt{ }\edtext{\edtext{\edtext{\skt{घृतं\-न्त\-मो\-द\-न\-म\-श्नी\-या\-दिति}}{\lemma{\skt{घृतंन्तम्}}\Afootnote{H\textsuperscript{pc}; \emph{om.}\ A; \skt{घृ}\ H\textsuperscript{ac}}}}{\lemma{\skt{ओदनम्}}\Afootnote{H\textsuperscript{pc}; \emph{om.}\ A; \skt{तं}\ H\textsuperscript{ac}}}}{\lemma{\skt{अश्नीयादिति}}\Afootnote{\skt{ल्पस्नेहेन}\ A; \emph{om.}\ H\textsuperscript{ac}; \skt{अश्नीयादिति}\ H\textsuperscript{pc}}}\skt{ }\edtext{\skt{घृत\-व\-न्त\-मो\-द\-न\-म\-श्नी\-यात्~। एतैः}}{\lemma{\skt{अश्नीयात्~। एतैः}}\Afootnote{H; \skt{अश्नीयात्; एते}\ A}}\skt{ }\edtext{\skt{खलु}}{\lemma{\skt{खलु}}\Afootnote{H; \skt{स्वल्व}\ A}}\skt{ }\edtext{\skt{दुर्न्नामानः}}{\lemma{\skt{दुर्न्नामानः}}\Afootnote{H; \skt{र्शांसि}\ A}}\skt{ }\edtext{\skt{क्षयं}}{\lemma{\skt{क्षयं}}\Afootnote{H; \skt{क्षप}\ A}}\skt{ }\edtext{\skt{यान्ति\-क्रि\-म\-य\-श्चो\-प\-शा\-म्यन्ति}}{\lemma{\skt{यान्ति\appbrk{}क्रि\appbrk{}म\appbrk{}य\appbrk{}श्चो\appbrk{}प\appbrk{}शा\appbrk{}म्यन्ति}}\Afootnote{\skt{यन्ति, कृमीनुपघ्नन्ति,}\ A; \skt{यान्ति\appbrk{}क्रि\appbrk{}म\appbrk{}य\appbrk{}श्चो\appbrk{}प\appbrk{}शा\appbrk{}म्यन्ति~।}\ H}}\skt{~। }\edtext{\skt{ग्रह\-ण\-धा\-र\-ण\-श\-क्तिश्च}}{\lemma{\skt{ग्रह\appbrk{}ण\appbrk{}धा\appbrk{}र\appbrk{}ण\appbrk{}श\appbrk{}क्तिश्च}}\Afootnote{H; \skt{ग्रह\appbrk{}ण\appbrk{}धा\appbrk{}र\appbrk{}ण\appbrk{}शक्तिं ज}\ A}}\skt{ }\edtext{\skt{भवति}}{\lemma{\skt{भवति}}\Afootnote{H; \skt{नयन्ति,}\ A}}\skt{~। }\edtext{\skt{मासि}}{\lemma{\skt{मासि}}\Afootnote{H; \skt{मासे}\ A}}\skt{ }\edtext{\skt{मासि च}}{\lemma{\skt{मासि च}}\Afootnote{H; \skt{मासे च}\ A}}\skt{ प्रयोगे वर्षशतं }\edtext{\skt{वर्षशतं}}{\lemma{\skt{वर्षशतं}}\Afootnote{H; \skt{वर्षशत}\ A}}\skt{ }\edtext{\edtext{\edtext{\skt{आयु\-षोऽ\-भि\-वृ\-द्धि\-र्व्भ\-वति}}{\lemma{\skt{आयुषो}}\Afootnote{H; \skt{मायुषो}\ A}}}{\lemma{\skt{ऽभिवृद्धिर्}}\Afootnote{A; \skt{ऽभिवृद्धि}\ H}}}{\lemma{\skt{व्भवति}}\Afootnote{\skt{भवति~॥}\ A; \skt{र्व्भवति~॥}\ H}}\skt{~॥}\pend

\pstart \ledsidenote{\passref{4.27.8}}\edtext{\skt{विड\-ङ्ग\-त\-ण्डु\-ल\-द्रोणं}}{\lemma{\skt{विड\appbrk{}ङ्ग\appbrk{}त\appbrk{}ण्डु\appbrk{}ल\appbrk{}द्रोणं}}\Afootnote{H; \skt{विड\appbrk{}ङ्ग\appbrk{}त\appbrk{}ण्डु\appbrk{}लानां द्रोणं}\ A}}\skt{ }\edtext{\skt{पिष्ट\-प\-व\-न\-पि\-ष्ट\-व\-दु\-स्वेद्य}}{\lemma{\skt{पिष्ट\appbrk{}प\appbrk{}व\appbrk{}न\appbrk{}पि\appbrk{}ष्ट\appbrk{}व\appbrk{}दु\appbrk{}स्वेद्य}}\Afootnote{H; \skt{पिष्टपचने पिष्ट\appbrk{}व\appbrk{}दु\appbrk{}प\appbrk{}स्वेद्य}\ A}}\skt{ विगतकषायं }\edtext{\edtext{\skt{सुस्वि\-न्न\-म\-व\-तार्य}}{\lemma{\skt{सुस्विन्नम्}}\Afootnote{H; \skt{स्विन्न}\ A}}}{\lemma{\skt{अवतार्य}}\Afootnote{H; \skt{मवतार्य}\ A}}\skt{ }\edtext{\skt{दृशदि}}{\lemma{\skt{दृशदि}}\Afootnote{H; \skt{दृषदि}\ A}}\skt{ }\edtext{\skt{प्रपिष्टमायस्यां}}{\lemma{\skt{प्रपिष्टमायस्यां}}\Afootnote{H; \skt{पिष्टमायसे}\ A}}\skt{ }\edtext{\skt{कुम्भ्यां}}{\lemma{\skt{कुम्भ्यां}}\Afootnote{H; \skt{दृढे कुम्भे}\ A}}\skt{ }\edtext{\edtext{\skt{मधू\-त्त\-रो\-द\-क\-ङ्कृत्वा}}{\lemma{\skt{मधू\appbrk{}त्त\appbrk{}रो\appbrk{}द\appbrk{}कङ्}}\Afootnote{H; \skt{मधूदकोत्तरं}\ A}}}{\lemma{\skt{कृत्वा}}\Afootnote{H; \skt{प्रावृषि}\ A}}\skt{~। }\edtext{\skt{भस्मराशौ}}{\lemma{\skt{भस्मराशौ}}\Afootnote{H; \skt{भस्मराशा}\ A}}\skt{ }\edtext{\skt{प्रावृ\-ष्य\-भ्य\-त\-र\-तश्च}}{\lemma{\skt{प्रावृ\appbrk{}ष्य\appbrk{}भ्य\appbrk{}त\appbrk{}र\appbrk{}तश्च}}\Afootnote{H; \skt{वन्तर्गृहे}\ A}}\skt{ चतु रो }\edtext{\skt{मासान्निदध्यात्~। वर्षाव्यवाये}}{\lemma{\skt{मासान्निदध्यात्~। वर्षाव्यवाये}}\Afootnote{H; \skt{मासान्निदध्यात्, वर्षाविगमे}\ A}}\skt{ चोद्धृ\-त्यो\-प\-सं\-स्कृ\-त\-श\-रीरः सह\-स्र\-स\-म्पा\-ता\-भि\-हुतं कृत्वा प्रातः }\edtext{\skt{प्रात\-र्य\-था\-ब\-ल\-मु\-प\-यु\-ञ्जीत}}{\lemma{\skt{प्रात\appbrk{}र्य\appbrk{}था\appbrk{}ब\appbrk{}लम्}}\Afootnote{\skt{प्रात\appbrk{}र्य\appbrk{}था\appbrk{}ब\appbrk{}लम्}\ A; \skt{प्रात\appbrk{}र्य\appbrk{}था\appbrk{}व\appbrk{}लम्}\ H}}\skt{~। जीर्ण्णे }\edtext{\skt{मुद्गयूषेण}}{\lemma{\skt{मुद्गयूषेण}}\Afootnote{H; \skt{मुद्गा\appbrk{}म\appbrk{}ल\appbrk{}क\appbrk{}यू\appbrk{}षे\appbrk{}णा\appbrk{}ल\appbrk{}व\appbrk{}णेन}\ A}}\skt{ }\edtext{\skt{घृत\-मो\-द\-न\-म\-श्नी\-यात्}}{\lemma{\skt{घृतम्}}\Afootnote{H; \skt{घृतवन्तम्}\ A}}\skt{~। पांशुशय्यायां }\edtext{\skt{शाययीत}}{\lemma{\skt{शाययीत}}\Afootnote{H; \skt{शयीत,}\ A}}\skt{~। तस्य }\edtext{\skt{मासादूर्द्धं}}{\lemma{\skt{मासादूर्द्धं}}\Afootnote{H; \skt{मासादूर्ध्वं}\ A}}\skt{ सर्व्वाङ्गेभ्यः }\edtext{\skt{क्रिमयो}}{\lemma{\skt{क्रिमयो}}\Afootnote{H; \skt{कृमयो}\ A}}\skt{ }\edtext{\skt{निःक्रामन्ति}}{\lemma{\skt{निःक्रामन्ति}}\Afootnote{H; \skt{निष्क्रामन्ति,}\ A}}\skt{~। }\edtext{\skt{तान\-नु\-तै\-ले\-ना\-भ्य\-क्तस्य}}{\lemma{\skt{तान\appbrk{}नु\appbrk{}तै\appbrk{}ले\appbrk{}ना\appbrk{}भ्य\appbrk{}क्तस्य}}\Afootnote{H; \skt{तान\appbrk{}णु\appbrk{}तै\appbrk{}ले\appbrk{}ना\appbrk{}भ्य\appbrk{}क्तस्य}\ A}}\skt{ }\edtext{\skt{विद\-ले\-ना\-प\-ह\-रेत्~। द्वितीये}}{\lemma{\skt{विद\appbrk{}ले\appbrk{}ना\appbrk{}प\appbrk{}ह\appbrk{}रेत्~। द्वितीये}}\Afootnote{H; \skt{वंश\appbrk{}वि\appbrk{}द\appbrk{}ले\appbrk{}ना\appbrk{}प\appbrk{}ह\appbrk{}रेत्, द्वितीये}\ A}}\skt{ }\edtext{\edtext{\skt{पिपी\-लि\-का\-स्तृ\-तीये}}{\lemma{\skt{पिपीलिकास्}}\Afootnote{H; \skt{पिपीलिका}\ A}}}{\lemma{\skt{तृतीये}}\Afootnote{H; \skt{स्तृतीये}\ A}}\skt{ }\edtext{\edtext{\skt{यूका\-ना\-न्त\-थै\-वा\-प\-ह\-रेत्}}{\lemma{\skt{यूकानान्}}\Afootnote{H; \skt{यूका}\ A}}\skt{~। चतुर्थे}}{\lemma{\skt{तथै\appbrk{}वा\appbrk{}प\appbrk{}ह\appbrk{}रेत्~। चतुर्थे}}\Afootnote{H; \skt{स्तथै\appbrk{}वा\appbrk{}प\appbrk{}ह\appbrk{}रेत्, चतुर्थे}\ A}}\skt{ }\edtext{\skt{दन्त\-न\-ख\-रो\-मा\-ण्य\-व\-सी\-र्यन्ते}}{\lemma{\skt{दन्त\appbrk{}न\appbrk{}ख\appbrk{}रो\appbrk{}मा\appbrk{}ण्य\appbrk{}व\appbrk{}सी\appbrk{}र्यन्ते}}\Afootnote{H; \skt{दन्त\appbrk{}न\appbrk{}ख\appbrk{}रो\appbrk{}मा\appbrk{}ण्य\appbrk{}व\appbrk{}शी\appbrk{}र्यन्ते;}\ A}}\skt{~। पञ्चमे प्रशस्त गुणलक्षणानि जायन्ते~। }\edtext{\skt{अमा\-नु\-ष\-ञ्चा\-दि\-त्य\-प्र\-का\-श\-म्व\-पु\-र\-धि\-ग\-च्छति}}{\lemma{\skt{अमा\appbrk{}नु\appbrk{}ष\appbrk{}ञ्चा\appbrk{}दि\appbrk{}त्य\appbrk{}प्र\appbrk{}का\appbrk{}शम्}}\Afootnote{H; \skt{अमानुषं चादित्यप्रकाशं}\ A}}\skt{~। }\edtext{\skt{दूराच्छ्रवण}}{\lemma{\skt{दूराच्छ्रवण}}\Afootnote{H; \skt{दूराच्छ्रवणानि}\ A}}\skt{ दर्शनानि चास्य }\edtext{\skt{भवति}}{\lemma{\skt{भवति}}\Afootnote{H; \skt{भवन्ति,}\ A}}\skt{~। रजस्तमसी चापोह्य }\edtext{\skt{सत्व\-म\-धि\-ति\-ष्ठते}}{\lemma{\skt{सत्व\appbrk{}म\appbrk{}धि\appbrk{}ति\appbrk{}ष्ठते}}\Afootnote{H; \skt{सत्त्व\appbrk{}म\appbrk{}धि\appbrk{}ति\appbrk{}ष्ठति,}\ A}}\skt{~। }\edtext{\skt{श्रुतनिगादी}}{\lemma{\skt{श्रुतनिगादी}}\Afootnote{H; \skt{श्रुतनिगा}\ A}}\skt{~। }\edtext{\skt{अपूर्व्वोत्पादी}}{\lemma{\skt{अपूर्व्वोत्पादी}}\Afootnote{H; \skt{द्यपूर्वोत्पादी}\ A}}\skt{ }\edtext{\skt{पुष्टो}}{\lemma{\skt{पुष्टो}}\Afootnote{H; \emph{om.}\ A}}\skt{ }\edtext{\skt{वृष\-भ\-व\-लोऽ\-श्व\-जवः}}{\lemma{\skt{वृषभवलो}}\Afootnote{H; \skt{गजबलो}\ A}}\skt{ }\edtext{\skt{पुनर्युवाष्टौ}}{\lemma{\skt{पुनर्युवाष्टौ}}\Afootnote{H; \skt{पुनर्युवाऽष्टौ}\ A}}\skt{ वर्ष\-श\-ता\-न्या\-यु\-र\-वा\-प्नोति~। }\edtext{\skt{तस्या\-नु\-तै\-ल\-म\-भ्य\-ङ्गार्थे}}{\lemma{\skt{तस्या\appbrk{}नु\appbrk{}तै\appbrk{}ल\appbrk{}म\appbrk{}भ्य\appbrk{}ङ्गार्थे}}\Afootnote{H; \skt{तस्या\appbrk{}णु\appbrk{}तै\appbrk{}ल\appbrk{}म\appbrk{}भ्यङ्गा}\ A}}\skt{~। कषा\-य\-मा\-ज\-क\-र्ण्ण\-मु\-त्सा\-द\-नार्थे~। सौशीरं }\edtext{\skt{कूनूपोदकं}}{\lemma{\skt{कूनूपोदकं}}\Afootnote{H\textsuperscript{pc}; \skt{कूपोदकं}\ H\textsuperscript{ac}}}\skt{ स्नानार्थे~। }\edtext{\edtext{\skt{चन्द\-न\-म\-नु\-ले\-प\-नार्थे}}{\lemma{\skt{चन्द\appbrk{}न\appbrk{}म\appbrk{}नु\appbrk{}ले\appbrk{}प\appbrk{}नार्थे}}\Afootnote{H; \skt{चन्द\appbrk{}न\appbrk{}मु\appbrk{}प\appbrk{}ले\appbrk{}प\appbrk{}नार्थे,}\ A}}\skt{~। }\edtext{\skt{भल्लाततक}}{\lemma{\skt{भल्लाततक}}\Afootnote{H; \skt{भल्लातक}\ A}}\skt{ }\edtext{\skt{विधा\-न\-व\-दु\-प\-योगः}}{\lemma{\skt{विधा\appbrk{}न\appbrk{}व\appbrk{}दु\appbrk{}प\appbrk{}योगः}}\Afootnote{H; \skt{विधा\appbrk{}न\appbrk{}व\appbrk{}दा\appbrk{}हारः}\ A}}\skt{ परिहारश्च~।}}{\lemma{\skt{र्थे… स्नानार्थे}}\Afootnote{\emph{om.}\ A}}\pend

\pstart \ledsidenote{\passref{4.27.9}}\edtext{\skt{काश्म\-री\-णा\-न्नि\-ष्कु\-ली\-कृ\-ता\-ना\-मेष}}{\lemma{\skt{काश्म\appbrk{}री\appbrk{}णा\appbrk{}न्नि\appbrk{}ष्कु\appbrk{}ली\appbrk{}कृ\appbrk{}ता\appbrk{}ना\appbrk{}मेष}}\Afootnote{H; \skt{काश्मर्याणां निष्कु\appbrk{}ली\appbrk{}कृ\appbrk{}ता\appbrk{}ना\appbrk{}मेष}\ A}}\skt{ एव कल्पः~। }\edtext{\skt{पांशु\-श\-य्या\-भो\-ज\-न\-वर्ज्यं}}{\lemma{\skt{पांशु\appbrk{}श\appbrk{}य्या\appbrk{}भो\appbrk{}ज\appbrk{}न\appbrk{}वर्ज्यं}}\Afootnote{H; \skt{पांशु\appbrk{}श\appbrk{}य्या\appbrk{}भो\appbrk{}ज\appbrk{}न\appbrk{}वर्जं;}\ A}}\skt{~। }\edtext{\skt{तत्रापि}}{\lemma{\skt{तत्रापि}}\Afootnote{H; \skt{अत्र हि}\ A}}\skt{ पयसा }\edtext{\skt{श्रितेन}}{\lemma{\skt{श्रितेन}}\Afootnote{H; \skt{शृतेन}\ A}}\skt{ भोक्तव्यं~। }\edtext{\skt{आशिषश्च}}{\lemma{\skt{आशिषश्च}}\Afootnote{H; \skt{समा\appbrk{}न\appbrk{}म\appbrk{}न्य\appbrk{}त्पू\appbrk{}र्वे\appbrk{}णा\appbrk{}शि\appbrk{}षश्च~।}\ A}}\skt{ }\edtext{\skt{समानाः}}{\lemma{\skt{समानाः}}\Afootnote{H; \emph{om.}\ A}}\skt{ }\edtext{\skt{पूर्व्वेण}}{\lemma{\skt{पूर्व्वेण}}\Afootnote{H; \emph{om.}\ A}}\skt{~। शोणि\-त\-पि\-त्त\-नि\-मि\-त्तेषु विका\-रे\-ष्वे\-ते\-षा\-मु\-प\-योगः~।}\pend

\pstart \ledsidenote{\passref{4.27.10}}\skt{यथो\-क्त\-मा\-गा\-र\-म्प्र\-विश्य बला\-मू\-ला\-र्द्ध\-प\-ल\-म्प\-लम्वा }\edtext{\skt{पयस्यालोड्य}}{\lemma{\skt{पयस्यालोड्य}}\Afootnote{H; \skt{पयसाऽऽआलोड्य}\ A}}\skt{ पिबेत्~। जीर्ण्णे पयः सर्प्पिरोदन इत्याहारः~। एवं }\edtext{\skt{द्वाद\-श\-रा\-त्र\-मु\-प\-यु\-ञ्जीत ततो}}{\lemma{\skt{उपयुञ्जीत ततो}}\Afootnote{H; \skt{उपयुज्य द्वादश वर्षाणि वयस्तिष्ठति; एवं दिव\appbrk{}स\appbrk{}श\appbrk{}त\appbrk{}मु\appbrk{}प\appbrk{}युज्य}\ A}}\skt{ वर्ष\-श\-त\-म्व\-य\-स्ति\-ष्ठति~। }\edtext{\skt{एष}}{\lemma{\skt{एष}}\Afootnote{H; \skt{एवम्}\ A}}\skt{ }\edtext{\skt{एव}}{\lemma{\skt{एव}}\Afootnote{H; \skt{ए}\ A}}\skt{ }\edtext{\skt{बला\-ति\-ब\-ला\-ना\-ग\-ब\-ला\-वि\-दा\-रीश}}{\lemma{\skt{बला\appbrk{}ति\appbrk{}ब\appbrk{}ला\appbrk{}ना\appbrk{}ग\appbrk{}ब\appbrk{}ला\appbrk{}वि\appbrk{}दा\appbrk{}रीश}}\Afootnote{H; \skt{वाति\appbrk{}ब\appbrk{}ला\appbrk{}ना\appbrk{}ग\appbrk{}ब\appbrk{}ला\appbrk{}वि\appbrk{}दा\appbrk{}रीश}\ A}}\skt{ ताव\-री\-णा\-मु\-प\-योगः~। }\edtext{\skt{विशेषतस्तु}}{\lemma{\skt{विशेषतस्तु}}\Afootnote{H; \skt{विशेषत}\ A}}\skt{ }\edtext{\edtext{\skt{अति\-ब\-ला\-मु\-द\-केन}}{\lemma{\skt{अतिबलाम्}}\Afootnote{H; \skt{स्त्वतिबला}\ A}}}{\lemma{\skt{उदकेन}}\Afootnote{H; \skt{मुदकेन,}\ A}}\skt{~। नागबलाचूर्ण्णं }\edtext{\skt{मधुक्षीरेण}}{\lemma{\skt{मधुक्षीरेण}}\Afootnote{H; \skt{मधुना, विदारीचूर्णं क्षीरेण,}\ A}}\skt{~। शता\-व\-री\-म\-प्येवं~। }\edtext{\skt{पूर्व्वेणाशिषश्च}}{\lemma{\skt{पूर्व्वेणाशिषश्च}}\Afootnote{H; \skt{पूर्वे\appbrk{}णा\appbrk{}न्य\appbrk{}त्स\appbrk{}मा\appbrk{}न\appbrk{}मा\appbrk{}शि\appbrk{}षश्च}\ A}}\skt{ }\edtext{\skt{समानाः}}{\lemma{\skt{समानाः}}\Afootnote{H; \skt{समाः~।}\ A}}\skt{~। }\edtext{\skt{एते }\edtext{\skt{ह्यौषधप्रयोगाः}}{\lemma{\skt{औषधप्रयोगाः}}\Afootnote{H; \skt{स्त्वौषधयो}\ A}}}{\lemma{\skt{एते ह्य्}}\Afootnote{H; \skt{एता}\ A}}\skt{~। बलकामानां शोषिणां रक्त\-पि\-त्तो\-प\-सृ\-ष्टानां शोणि\-त\-ञ्च्छ\-र्द्द\-य\-ता\-म्वि\-रि\-च्य\-मा\-ना\-ना\-ञ्चो\-प\-दि\-श्यन्ते~॥}\pend

\pstart \ledsidenote{\passref{4.27.11}}\edtext{\skt{वारा\-ही\-मू\-ल\-न्तु\-ला\-चू\-र्ण्ण\-ङ्कृत्वा}}{\lemma{\skt{वाराहीमूलन्}}\Afootnote{H; \skt{वाराहीमूल}\ A}}\skt{ ततो }\edtext{\edtext{\skt{मात्रा\-म्म\-धु\-यु\-क्त\-म्प\-य\-सा\-लोड्य}}{\lemma{\skt{मधुयुक्तम्}}\Afootnote{H; \skt{मधुयुक्तां}\ A}}}{\lemma{\skt{पयसालोड्य}}\Afootnote{H; \skt{पयसाऽऽआलोड्य}\ A}}\skt{ पिबेत्~। जीर्ण्णे पयः सर्प्पिरोदन इत्याहारः~। }\edtext{\edtext{\edtext{\skt{प्रयो\-ग\-मि\-म\-मु\-प\-ये\-व\-मानो}}{\lemma{\skt{प्रयोगम्}}\Afootnote{H; \skt{प्रतिबेधोऽत्र पूर्ववत्; प्रयोग}\ A}}}{\lemma{\skt{इमम्}}\Afootnote{H; \skt{मिमम्}\ A}}}{\lemma{\skt{उपयेवमानो}}\Afootnote{H; \skt{उपसेवमानो}\ A}}\skt{ }\edtext{\edtext{\edtext{\skt{वर्ष\-श\-त\-मा\-यु\-र\-वा\-प्नोति}}{\lemma{\skt{वर्षशतम्}}\Afootnote{H; \skt{वर्षशत}\ A}}}{\lemma{\skt{आयुर्}}\Afootnote{H; \skt{मायु}\ A}}}{\lemma{\skt{अवाप्नोति}}\Afootnote{H; \skt{रवाप्नोति}\ A}}\skt{~॥ स्त्रीषु }\edtext{\skt{चाक्षयता}}{\lemma{\skt{चाक्षयता}}\Afootnote{H; \skt{चाक्षयताम्,}\ A}}\skt{~। एतेनैव }\edtext{\skt{चूर्ण्णएन}}{\lemma{\skt{चूर्ण्णएन}}\Afootnote{H; \skt{चूर्णेन}\ A}}\skt{ }\edtext{\edtext{\edtext{\skt{पयो\-व\-चू\-र्ण्णा\-श्रि\-त\-श्रि\-त\-शी\-त\-म\-भि\-म\-थ्या\-ज्य\-मु\-त्पाद्य}}{\lemma{\skt{पयो\appbrk{}व\appbrk{}चू\appbrk{}र्ण्णा\appbrk{}श्रि\appbrk{}त\appbrk{}श्रि\appbrk{}त\appbrk{}शी\appbrk{}तम्}}\Afootnote{H; \skt{पयोऽवचूर्ण्य शृतशीतम्}\ A}}}{\lemma{\skt{अभिमथ्याज्यम्}}\Afootnote{H; \skt{अभिमथ्याज्य}\ A}}}{\lemma{\skt{उत्पाद्य}}\Afootnote{H; \skt{मुत्पाद्य}\ A}}\skt{~। }\edtext{\skt{मधु\-यु\-क्त\-मु\-प\-यु\-ञ्जीत}}{\lemma{\skt{मधुयुक्तम्}}\Afootnote{H; \skt{मधुयुतम्}\ A}}\skt{~। }\edtext{\skt{जीर्ण्णे}}{\lemma{\skt{जीर्ण्णे}}\Afootnote{H; \skt{सायं\appbrk{}प्रा\appbrk{}त\appbrk{}रे\appbrk{}क\appbrk{}कालं वा, जीर्णे}\ A}}\skt{ पयः सर्प्पिरोदन }\edtext{\skt{इत्याहारः}}{\lemma{\skt{आहारः}}\Afootnote{H; \skt{आहारः, एवं मासमुपयुज्य सौपर्णं भव\appbrk{}त्य\appbrk{}न\appbrk{}ल्प\appbrk{}बलः स्त्रीषु चाक्षयो वर्ष\appbrk{}श\appbrk{}ता\appbrk{}यु\appbrk{}र्भ\appbrk{}व\appbrk{}तीति~॥}\ A}}\skt{~॥}\pend

\pstart \ledsidenote{\passref{4.27.12}}\skt{चक्षुः कामः प्राणकामो वा }\edtext{\skt{जीव\-वी\-ज\-क\-सा\-रा\-ग्नि\-म\-न्थ\-मू\-लानि}}{\lemma{\skt{जीव\appbrk{}वी\appbrk{}ज\appbrk{}क\appbrk{}सा\appbrk{}रा\appbrk{}ग्नि\appbrk{}म\appbrk{}न्थ\appbrk{}मू\appbrk{}लानि}}\Afootnote{H\textsuperscript{pc}; \skt{बीज\appbrk{}क\appbrk{}सा\appbrk{}रा\appbrk{}ग्नि\appbrk{}म\appbrk{}न्थ\appbrk{}मूलं नि}\ A; \skt{जीव\appbrk{}क\appbrk{}सा\appbrk{}रा\appbrk{}ग्नि\appbrk{}म\appbrk{}न्थ\appbrk{}मू\appbrk{}लानि}\ H\textsuperscript{ac}}}\skt{ }\edtext{\skt{क्वाथैर्म्माषप्रस्थं}}{\lemma{\skt{क्वाथैर्म्माषप्रस्थं}}\Afootnote{H; \skt{ष्क्वाथ्य माषप्रस्थं}\ A}}\skt{ साधयेत्~। तस्मिन्सिध्यति }\edtext{\skt{चित्र\-क\-मू\-ला\-ना\-म\-क्ष\-मात्रं}}{\lemma{\skt{चित्र\appbrk{}क\appbrk{}मू\appbrk{}ला\appbrk{}ना\appbrk{}म\appbrk{}क्ष\appbrk{}मात्रं}}\Afootnote{H; \skt{चित्र\appbrk{}क\appbrk{}मू\appbrk{}ला\appbrk{}ना\appbrk{}म\appbrk{}क्ष\appbrk{}मात्रं}\ A}}\skt{ }\edtext{\skt{कल्कन्दद्यात्~। आम\-ल\-क\-र\-स\-च\-तु\-र्व्भा\-ग\-न्ततः}}{\lemma{\skt{दद्यात्~। आम\appbrk{}ल\appbrk{}क\appbrk{}र\appbrk{}स\appbrk{}च\appbrk{}तु\appbrk{}र्व्भा\appbrk{}गन्}}\Afootnote{H; \skt{दद्या\appbrk{}दा\appbrk{}म\appbrk{}ल\appbrk{}क\appbrk{}र\appbrk{}स\appbrk{}च\appbrk{}तु\appbrk{}र्थ\appbrk{}भागं,}\ A}}\skt{ }\edtext{\skt{सुस्वि\-न्न\-म\-व\-तार्य}}{\lemma{\skt{सुस्विन्नम्}}\Afootnote{H; \skt{स्विन्नम्}\ A}}\skt{ }\edtext{\skt{शीती\-भू\-त\-म्म\-धु\-स\-र्प्पिर्ब्भ्यां}}{\lemma{\skt{शीतीभूतम्}}\Afootnote{H; \skt{सह\appbrk{}स्र\appbrk{}स्ं\appbrk{}पा\appbrk{}ता\appbrk{}भि\appbrk{}हुतं कृत्वा शीतीभूतं}\ A}}\skt{ संसृ\-ज्यो\-प\-यु\-ञ्जीत~। }\edtext{\skt{जीर्ण्णे}}{\lemma{\skt{जीर्ण्णे}}\Afootnote{H; \skt{यथाबलं, यथासात्म्यं च लवणं परि\appbrk{}ह\appbrk{}र\appbrk{}न्भ\appbrk{}क्ष\appbrk{}येत्। जीर्णे}\ A}}\skt{ }\edtext{\skt{मुद्गयूषेणा}}{\lemma{\skt{मुद्गयूषेणा}}\Afootnote{H; \skt{मुद्गा\appbrk{}म\appbrk{}ल\appbrk{}क\appbrk{}यू\appbrk{}षेणा}\ A}}\skt{ लवणेन }\edtext{\edtext{\skt{घृत\-व\-न्त\-मो\-द\-न\-म\-श्नी\-यात्}}{\lemma{\skt{घृत\appbrk{}व\appbrk{}न्त\appbrk{}मो\appbrk{}द\appbrk{}नम्}}\Afootnote{H; \skt{घृत\appbrk{}व\appbrk{}न्त\appbrk{}मो\appbrk{}दन}\ A}}\skt{~। }\edtext{\skt{मास\-मा\-त्र\-मे\-व\-म\-नेन}}{\lemma{\skt{अनेन}}\Afootnote{H; \skt{वमाभ्यां}\ A}}}{\lemma{\skt{अश्नीयात्~। मास\appbrk{}मा\appbrk{}त्र\appbrk{}मे\appbrk{}वम्}}\Afootnote{H; \skt{मश्नीयात्पयसा वा मासत्रयम्; ए}\ A}}\skt{ }\edtext{\skt{प्रयोगेन च क्षुः}}{\lemma{\skt{प्रयोगेन च क्षुः}}\Afootnote{H; \skt{प्रयोगाभ्यां चक्षुः}\ A}}\skt{ सौपर्ण्णम्भवति~। }\edtext{\skt{अल्पबलाबलां}}{\lemma{\skt{अल्पबलाबलां}}\Afootnote{H; \skt{अनल्पबलः}\ A}}\skt{ स्त्रीषु }\edtext{\skt{चाक्ष्य्यो}}{\lemma{\skt{चाक्ष्य्यो}}\Afootnote{H; \skt{चाक्षयो}\ A}}\skt{ }\edtext{\skt{वर्ष\-श\-ता\-यु\-र्ब्भ\-वति}}{\lemma{\skt{वर्ष\appbrk{}श\appbrk{}ता\appbrk{}यु\appbrk{}र्ब्भ\appbrk{}वति}}\Afootnote{H; \skt{वर्ष\appbrk{}श\appbrk{}ता\appbrk{}यु\appbrk{}र्भ\appbrk{}व\appbrk{}तीति॥}\ A}}\skt{~॥~॥}\pend

\pstart \ledsidenote{\passref{4.27.trailer}}\edtext{\skt{चिकित्सा}}{\lemma{\skt{चिकित्सा}}\Afootnote{H; \skt{इति सुश्रु\appbrk{}त\appbrk{}सं\appbrk{}हि\appbrk{}तायां चिकित्सा}\ A}}\skt{ }\edtext{\skt{थ}}{\lemma{\skt{थ}}\Afootnote{H; \skt{स्थाने सर्वो\appbrk{}प\appbrk{}घा\appbrk{}त\appbrk{}श\appbrk{}म\appbrk{}नीयं रसा\appbrk{}य\appbrk{}न\appbrk{}चि\appbrk{}कि\appbrk{}त्सितं नाम सप्तविंशोऽध्यायः~॥२७~॥}\ A}}\pend

\endnumbering
\beginnumbering
\eledsection*{Adhyāya 28}
\pstart \ledsidenote{\passref{4.28.1}}\skt{अथातो }\edtext{\skt{मेधायुःकामीयं}}{\lemma{\skt{मेधायुःकामीयं}}\Afootnote{H; \skt{मेधायुष्कामीयं}\ A}}\skt{ }\edtext{\skt{रसा\-य\-न\-म्व्या\-ख्या\-स्यामः}}{\lemma{\skt{रसायनम्}}\Afootnote{H; \skt{रसा\appbrk{}य\appbrk{}न\appbrk{}चि\appbrk{}कि\appbrk{}त्सितं}\ A}}\skt{~॥}\pend

\pstart \ledsidenote{\passref{4.28.3}}\edtext{\skt{मेधायुःकामः}}{\lemma{\skt{मेधायुःकामः}}\Afootnote{A; \skt{मेहायुः कासः}\ H}}\skt{ }\edtext{\skt{श्वेता\-व\-ल्गु\-ज\-फ\-ला\-द्या\-त\-प\-शु\-ष्का\-ण्या\-दाय}}{\lemma{\skt{श्वेता\appbrk{}व\appbrk{}ल्गु\appbrk{}ज\appbrk{}फ\appbrk{}ला\appbrk{}द्या\appbrk{}त\appbrk{}प\appbrk{}शु\appbrk{}ष्का\appbrk{}ण्या\appbrk{}दाय}}\Afootnote{\skt{श्वेता\appbrk{}व\appbrk{}ल्गु\appbrk{}ज\appbrk{}फ\appbrk{}ला\appbrk{}न्या\appbrk{}त\appbrk{}प\appbrk{}प\appbrk{}रि\appbrk{}शु\appbrk{}ष्का\appbrk{}ण्या\appbrk{}दाय}\ A; \skt{स्वेता\appbrk{}व\appbrk{}ल्गु\appbrk{}ज\appbrk{}फ\appbrk{}ला\appbrk{}त\appbrk{}प\appbrk{}शु\appbrk{}ष्का\appbrk{}ण्या\appbrk{}दाय}\ H}}\skt{ }\edtext{\skt{चूर्णानि}}{\lemma{\skt{चूर्णानि}}\Afootnote{H; \skt{सूक्ष्मचूर्णानि}\ A}}\skt{ कृत्वा गुडेन }\edtext{\skt{सहालोड्य}}{\lemma{\skt{सहालोड्य}}\Afootnote{A; \skt{साहालोड्य}\ H}}\skt{ }\edtext{\skt{स्नेहकुम्भे}}{\lemma{\skt{स्नेहकुम्भे}}\Afootnote{A; \skt{सेन्हकुम्भे}\ H}}\skt{ }\edtext{\skt{निदध्यात्~। ततः}}{\lemma{\skt{निदध्यात्~। ततः}}\Afootnote{H; \skt{सप्तरात्रं धान्यराशौ निदध्यात्, स}\ A}}\skt{ }\edtext{\skt{सप्त\-रा\-त्रा\-दु\-द्धृत्य}}{\lemma{\skt{सप्त\appbrk{}रा\appbrk{}त्रा\appbrk{}दु\appbrk{}द्धृत्य}}\Afootnote{H; \skt{प्तरात्रादुद्धृत्य}\ A}}\skt{ हृतदोषस्य }\edtext{\skt{यथाबलं}}{\lemma{\skt{यथाबलं}}\Afootnote{A; \skt{यथावलम्}\ H}}\skt{ पिण्डं }\edtext{\skt{प्रयच्छेत्~। अनुदित्सूर्ये,}}{\lemma{\skt{प्रयच्छेत्~। अनुदित्सूर्ये}}\Afootnote{\skt{प्रय\appbrk{}च्छे\appbrk{}द\appbrk{}नु\appbrk{}दिते सूर्ये,}\ A; \skt{प्रयच्छेत्~। अनुदितसूर्ये}\ H}}\skt{ }\edtext{\skt{उष्णो\-द\-क\-ञ्चा\-नु\-पि\-बेत्~। भल्ला\-त\-क\-वि\-धा\-न\-व\-च्चा\-गा\-र\-प्र\-वेशः}}{\lemma{\skt{चानुपिबेत्~। भल्ला\appbrk{}त\appbrk{}क\appbrk{}वि\appbrk{}धा\appbrk{}न\appbrk{}व\appbrk{}च्चा\appbrk{}गा\appbrk{}र\appbrk{}प्र\appbrk{}वेशः}}\Afootnote{A; \skt{चानुपिवेत्~। भल्ला\appbrk{}त\appbrk{}क\appbrk{}वि\appbrk{}धा\appbrk{}न\appbrk{}व\appbrk{}च्चा\appbrk{}गा\appbrk{}र\appbrk{}प्र\appbrk{}वेशः~।}\ H}}\skt{~। }\edtext{\skt{जीर्णौ\-ष\-ध\-श्चा\-प\-राह्णे}}{\lemma{\skt{चापराह्णे}}\Afootnote{A; \skt{चापराण्ह्वे}\ H}}\skt{ }\edtext{\skt{विहिमाभिरद्भिः}}{\lemma{\skt{विहिमाभिर्}}\Afootnote{H; \skt{हिमाभिर्}\ A}}\skt{~। परिषिक्तगात्रः शालीनां }\edtext{\skt{षष्टिकानाम्वा}}{\lemma{\skt{षष्टिकानाम्वा}}\Afootnote{H; \skt{षष्टिकानां च}\ A}}\skt{ पयसा शर्क\-रा\-म\-धु\-रे\-णौ\-द\-न\-म\-श्नी\-यात्~। एवं }\edtext{\skt{षण्मा\-सा\-नु\-प\-युञ्ज्य}}{\lemma{\skt{षण्मा\appbrk{}सा\appbrk{}नु\appbrk{}प\appbrk{}युञ्ज्य}}\Afootnote{\skt{षण्मा\appbrk{}सा\appbrk{}नु\appbrk{}प\appbrk{}युज्य}\ A H}}\skt{ विगतपाप्मा }\edtext{\skt{वलोपेतः}}{\lemma{\skt{वलोपेतः}}\Afootnote{\skt{बलवर्णोपेतः}\ A; \skt{वलायेतः}\ H}}\skt{ श्रुतनिगादी }\edtext{\skt{भवति}}{\lemma{\skt{भवति}}\Afootnote{H; \skt{स्मृतिमानरोगो वर्ष\appbrk{}श\appbrk{}ता\appbrk{}यु\appbrk{}र्भ\appbrk{}वति~।}\ A}}\skt{~। }\edtext{\edtext{\skt{कुष्ठ\-रो\-गि\-णा\-मु\-द\-रि\-णाम्वा}}{\lemma{\skt{कुष्ठरोगिणाम्}}\Afootnote{H; \skt{कुष्ठिनं पाण्डुरोगिणम्}\ A}}}{\lemma{\skt{उदरिणाम्वा}}\Afootnote{H; \skt{उदरिणं वा}\ A}}\skt{ }\edtext{\skt{कृष्णानि}}{\lemma{\skt{कृष्णानि}}\Afootnote{H; \skt{कृष्णाया}\ A}}\skt{ }\edtext{\skt{गोमू\-त्रा\-लो\-ड्या\-र्द्ध\-प\-लिकं}}{\lemma{\skt{गोमू\appbrk{}त्रा\appbrk{}लो\appbrk{}ड्या\appbrk{}र्द्ध\appbrk{}प\appbrk{}लिकं}}\Afootnote{\skt{गोमू\appbrk{}त्रे\appbrk{}णा\appbrk{}लो\appbrk{}ड्या\appbrk{}र्ध\appbrk{}प\appbrk{}लिकं}\ A H}}\skt{ }\edtext{\skt{पिण्डमुदिते}}{\lemma{\skt{उदिते}}\Afootnote{\skt{विगतलौहित्ये}\ A; \skt{उदित}\ H}}\skt{ }\edtext{\skt{सूर्ये}}{\lemma{\skt{सूर्ये}}\Afootnote{\skt{सवितरि}\ A; \skt{सूर्य}\ H}}\skt{ }\edtext{\skt{पाययेत्~। सायाह्ने}}{\lemma{\skt{पाययेत्~। सायाह्ने}}\Afootnote{\skt{पाययेत्, पराह्णे}\ A; \skt{पाययेत्~। पा}\ H}}\skt{ चाल\-व\-णे\-ना\-म\-ल\-क\-यू\-षेण }\edtext{\edtext{\skt{सर्पि\-ष्म\-न्त\-मो\-द\-न\-म\-श्नी\-यात्}}{\lemma{\skt{सर्पिष्मन्तम्}}\Afootnote{A; \skt{सर्प्पिष्मन्त}\ H}}}{\lemma{\skt{ओदनम्}}\Afootnote{A; \skt{मोदनम्}\ H}}\skt{~। एवं मासमुपयुज्य स्मृतिमानरोगो वर्ष\-श\-ता\-यु\-र्भ\-वति~। एष }\edtext{\skt{एवो\-प\-यो\-ग\-श्चि\-त्र\-क\-मू\-लानां}}{\lemma{\skt{चित्रकमूलानां}}\Afootnote{A; \skt{श्चित्रकमूलानां}\ H}}\skt{ }\edtext{\skt{रजन्याश्च}}{\lemma{\skt{रजन्याश्च}}\Afootnote{H; \skt{रजन्याश्च; चित्रकमूले विशेषो द्विपलिकं पिण्डं परं प्रमाणं, शेषं पूर्ववत्~॥}\ A}}\skt{~।}\pend

\pstart \ledsidenote{\passref{4.28.4}}\skt{हृतदोष }\edtext{\skt{एवा\-गा\-र\-म्प्र\-विश्य}}{\lemma{\skt{एवागारम्}}\Afootnote{H; \skt{एव प्रति\appbrk{}सं\appbrk{}सृ\appbrk{}ष्ट\appbrk{}भक्तो यथा\appbrk{}क्र\appbrk{}म\appbrk{}मा\appbrk{}गारं}\ A}}\skt{ }\edtext{\edtext{\skt{मण्डू\-क\-प\-र्णी\-स्व\-र\-स\-मा\-दाय}}{\lemma{\skt{मण्डू\appbrk{}क\appbrk{}प\appbrk{}र्णी\appbrk{}स्व\appbrk{}र\appbrk{}सम्}}\Afootnote{A; \skt{मण्डू\appbrk{}क\appbrk{}प\appbrk{}र्ण्णी\appbrk{}स्व\appbrk{}रस}\ H}}}{\lemma{\skt{आदाय}}\Afootnote{A; \skt{मादाय}\ H}}\skt{ सह\-स्र\-स\-म्पा\-ता\-भि\-हुतं कृत्वा }\edtext{\edtext{\skt{यथा\-ब\-ल\-म्प\-य\-सा\-लोड्य}}{\lemma{\skt{यथाबलम्}}\Afootnote{A; \skt{यथावलम्}\ H}}}{\lemma{\skt{पयसालोड्य}}\Afootnote{\skt{पयसाऽऽआलोड्य}\ A; \skt{पयसा}\ H}}\skt{ }\edtext{\skt{पिबेत्~। तिलैर्वा}}{\lemma{\skt{पिबेत्~। तिलैर्वा}}\Afootnote{\skt{पिबे\appbrk{}त्प\appbrk{}योऽ\appbrk{}नु\appbrk{}पानं वा, तस्यां जीर्णायां यवान्नं पय\appbrk{}सो\appbrk{}प\appbrk{}यु\appbrk{}ञ्जीत; तिलैर्वा}\ A; \skt{लो पिवेत्~। तिलैर्व्वा}\ H}}\skt{ सह }\edtext{\skt{भक्षयेत्~। पयोऽनुपानं}}{\lemma{\skt{भक्षयेत्~। पयो}}\Afootnote{H; \skt{भक्ष\appbrk{}ये\appbrk{}त्त्री\appbrk{}न्मा\appbrk{}सा\appbrk{}न्पयो}\ A}}\skt{~। जीर्णे पयः सर्पिरोदन }\edtext{\skt{इत्याहारः}}{\lemma{\skt{आहारः}}\Afootnote{H; \skt{आहारः; एव\appbrk{}मु\appbrk{}प\appbrk{}यु\appbrk{}ञ्जानो ब्रह्मवर्चसी श्रुतनिगादी भवति वर्ष\appbrk{}श\appbrk{}त\appbrk{}मा\appbrk{}यु\appbrk{}र\appbrk{}वा\appbrk{}प्नोति~। त्रिरा\appbrk{}त्रो\appbrk{}पो\appbrk{}षि\appbrk{}तश्च त्रिरात्रमेनां भक्षयेत्, त्रिरात्रादूर्ध्वं पयः}\ A}}\skt{~। }\edtext{\skt{बिल्वमात्रम्वा}}{\lemma{\skt{बिल्वमात्रम्वा}}\Afootnote{\skt{सर्पिरिति चोपयुञ्जीत~। बिल्वमात्रं}\ A; \skt{विल्वमात्रम्वा}\ H}}\skt{ }\edtext{\skt{पि }\edtext{\skt{पिण्ड\-म्प\-य\-सा\-लोड्य}}{\lemma{\skt{पयसालोड्य}}\Afootnote{\skt{पयसाऽऽआलोड्य}\ A; \skt{पयसालो}\ H}}}{\lemma{\skt{पि पिण्डम्}}\Afootnote{\skt{पिण्डं वा}\ A; \skt{पिण्डम्}\ H}}\skt{ }\edtext{\skt{पिबेत्~। }\edtext{\skt{एव\-न्द\-श\-रा\-त्र\-मु\-प\-युज्य}}{\lemma{\skt{दशरात्रम्}}\Afootnote{H; \skt{द्वादशरात्रम्}\ A}}}{\lemma{\skt{पिबेत्~। एवन्}}\Afootnote{A; \skt{पिवेत्~। एवन्}\ H}}\skt{ मेधावी वर्ष\-श\-ता\-यु\-र्भ\-वति~॥}\pend

\pstart \ledsidenote{\passref{4.28.5}}\skt{हृतदोष एवा\-गा\-र\-म्प्र\-विश्य }\edtext{\skt{प्रति\-सं\-सृ\-ष्ट\-भक्तो}}{\lemma{\skt{प्रति\appbrk{}सं\appbrk{}सृ\appbrk{}ष्ट\appbrk{}भक्तो}}\Afootnote{A; \skt{प्रति\appbrk{}सं\appbrk{}सृ\appbrk{}ष्ट\appbrk{}भक्तौ}\ H}}\skt{ }\edtext{\skt{ब्राह्मी\-स्व\-र\-स\-मा\-दाय}}{\lemma{\skt{ब्राह्मी\appbrk{}स्व\appbrk{}र\appbrk{}स\appbrk{}मा\appbrk{}दाय}}\Afootnote{A; \skt{व्रह्मी\appbrk{}स्व\appbrk{}र\appbrk{}स\appbrk{}मा\appbrk{}दाय}\ H}}\skt{ सह\-स्र\-स\-म्पा\-ता\-भि\-हुतं कृत्वा }\edtext{\edtext{\skt{यथा\-ब\-ल\-म्प्र\-यु\-ञ्जीतौ}}{\lemma{\skt{यथाबलम्}}\Afootnote{A; \skt{यथावलम्}\ H}}}{\lemma{\skt{प्रयुञ्जीतौ}}\Afootnote{\skt{पयुञ्जीत,}\ A; \skt{उपञ्जीत~।}\ H\textsuperscript{ac}; \skt{उपयुञ्जीत~।}\ H\textsuperscript{pc}}}\skt{~। }\edtext{\skt{जीर्णौ\-ष\-ध\-श्चा\-प\-राह्णे}}{\lemma{\skt{चापराह्णे}}\Afootnote{A; \skt{चापराण्हे}\ H}}\skt{ }\edtext{\edtext{\skt{यवा\-गू\-म\-ल\-व\-ण\-म्पि\-बेत्}}{\lemma{\skt{यवा\appbrk{}गू\appbrk{}म\appbrk{}ल\appbrk{}व\appbrk{}णम्}}\Afootnote{A; \skt{यवा\appbrk{}गू\appbrk{}म\appbrk{}ल\appbrk{}व\appbrk{}णाम्}\ H}}\skt{~॥ }\edtext{\skt{क्षीरसात्म्योऽथवा}}{\lemma{\skt{ऽथवा}}\Afootnote{\skt{वा}\ A H}}}{\lemma{\skt{पिबेत्~॥ क्षीरसात्म्यो}}\Afootnote{A; \skt{पिवेत्~॥ क्षीरसात्म्यो}\ H}}\skt{ पयसा भुञ्जीत~॥ एवं }\edtext{\skt{सप्त\-रा\-त्र\-मु\-प\-युज्य}}{\lemma{\skt{उपयुज्य}}\Afootnote{A; \skt{उपयुंज्य}\ H}}\skt{ }\edtext{\skt{व्रह्मवर्च्चसी}}{\lemma{\skt{व्रह्मवर्च्चसी}}\Afootnote{\skt{ब्रह्मवर्चसी}\ A; \skt{व्रह्मवर्च्चस्वी}\ H}}\skt{ मेधावी भवति~। द्वितीयं सप्त\-रा\-त्र\-मु\-प\-युज्य }\edtext{\skt{ग्रन्थ\-मी\-प्सि\-त\-मु\-त्पा\-द\-यति}}{\lemma{\skt{उत्पादयति}}\Afootnote{A; \skt{उत्पादय्ति}\ H}}\skt{ }\edtext{\skt{नष्टञ्चास्य}}{\lemma{\skt{नष्टञ्}}\Afootnote{A; \skt{नष्टश्}\ H}}\skt{ प्रादुर्भवति~। }\edtext{\skt{तृती\-य\-मु\-च्चा\-रितं}}{\lemma{\skt{उच्चारितं}}\Afootnote{H; \skt{सप्त\appbrk{}रा\appbrk{}त्र\appbrk{}मु\appbrk{}प\appbrk{}युज्य द्विरुच्चारितं}\ A}}\skt{ }\edtext{\skt{शत\-म\-प्य\-व\-धा\-र\-यति}}{\lemma{\skt{अवधारयति}}\Afootnote{A; \skt{उपधारयति~।}\ H}}\skt{~। }\edtext{\skt{एव\-मे\-क\-विं\-श\-ति\-रा\-त्र\-मु\-प\-युज्य}}{\lemma{\skt{उपयुज्य}}\Afootnote{H; \skt{उपयुज्या}\ A}}\skt{ }\edtext{\edtext{\skt{लक्ष्मी\-र\-प\-क्रा\-मति}}{\lemma{\skt{लक्ष्मीर्}}\Afootnote{A; \skt{लक्षी}\ H\textsuperscript{ac}; \skt{लक्ष्मी}\ H\textsuperscript{pc}}}}{\lemma{\skt{अपक्रामति}}\Afootnote{A; \skt{रपक्रामति~।}\ H}}\skt{~। }\edtext{\skt{शरी\-रा\-त्मू\-र्त्ति\-मती}}{\lemma{\skt{शरी\appbrk{}रा\appbrk{}त्मू\appbrk{}र्त्ति\appbrk{}मती}}\Afootnote{H; \skt{मूर्तिमती}\ A}}\skt{ }\edtext{\skt{चैन\-म्वा\-ग्दे\-वी\-म\-नु\-प्र\-वि\-शति}}{\lemma{\skt{वाग्देवीम्}}\Afootnote{H; \skt{वाग्देव्य्}\ A}}\skt{~। सर्वाश्चैनं }\edtext{\edtext{\skt{श्रुत\-योऽ\-व\-ति\-ष्ठते}}{\lemma{\skt{श्रुतयो}}\Afootnote{H; \skt{श्रुतय}\ A}}}{\lemma{\skt{ऽवतिष्ठते}}\Afootnote{\skt{उपतिष्ठन्ति,}\ A; \skt{ऽवतिष्ठेते~।}\ H}}\skt{~। }\edtext{\skt{श्रुतिधरः}}{\lemma{\skt{श्रुतिधरः}}\Afootnote{H; \skt{श्रुतधरः}\ A}}\skt{ पञ्च\-व\-र्ष\-श\-ता\-यु\-र्भ\-वति~॥}\pend

\pstart \ledsidenote{\passref{4.28.6}}\edtext{\skt{व्रह्मी\-स्व\-र\-स\-प्र\-स्थ\-द्वये}}{\lemma{\skt{व्रह्मी\appbrk{}स्व\appbrk{}र\appbrk{}स\appbrk{}प्र\appbrk{}स्थ\appbrk{}द्वये}}\Afootnote{H; \skt{ब्राह्मी\appbrk{}स्व\appbrk{}र\appbrk{}स\appbrk{}प्र\appbrk{}स्थ\appbrk{}द्वये}\ A}}\skt{ }\edtext{\edtext{\skt{घृत\-प्र\-स्थ\-म्वि\-ड\-ङ्ग\-त\-ण्डु\-ला\-ना\-ङ्कु\-ड\-व\-म्व\-चाभृ}}{\lemma{\skt{कुडवम्}}\Afootnote{H; \skt{कुडवं द्वे द्वे}\ A}}}{\lemma{\skt{वचाभृ}}\Afootnote{H; \skt{पले वचा}\ A}}\skt{ }\edtext{\skt{वृतयोः}}{\lemma{\skt{वृतयोः}}\Afootnote{H; \skt{मृतयोर्}\ A}}\skt{~। }\edtext{\skt{द्वादश}}{\lemma{\skt{द्वादश}}\Afootnote{H; \emph{om.}\ A}}\skt{ }\edtext{\skt{द्वादश च}}{\lemma{\skt{द्वादश च}}\Afootnote{H; \skt{द्वादश}\ A}}\skt{ }\edtext{\skt{हरी\-त\-क्या\-म\-ल\-क\-वि\-भी\-त\-कानि}}{\lemma{\skt{हरी\appbrk{}त\appbrk{}क्या\appbrk{}म\appbrk{}ल\appbrk{}क\appbrk{}वि\appbrk{}भी\appbrk{}त\appbrk{}कानि}}\Afootnote{H; \skt{हरी\appbrk{}त\appbrk{}क्या\appbrk{}म\appbrk{}ल\appbrk{}क\appbrk{}बि\appbrk{}भी\appbrk{}त\appbrk{}कानि}\ A}}\skt{ }\edtext{\skt{कल्क\-पि\-ष्टा\-न्य\-वा\-प्यै\-कध्यं}}{\lemma{\skt{कल्क\appbrk{}पि\appbrk{}ष्टा\appbrk{}न्य\appbrk{}वा\appbrk{}प्यै\appbrk{}कध्यं}}\Afootnote{H; \skt{श्लक्ष्ण\appbrk{}पि\appbrk{}ष्टा\appbrk{}न्या\appbrk{}वा\appbrk{}प्यै\appbrk{}कध्यं}\ A}}\skt{ साधयित्वा स्वनु\-गु\-प्त\-न्नि\-द\-ध्यात्~। ततः }\edtext{\skt{पूर्व्व\-व\-द्य\-था\-ब\-ल\-मु\-प\-यु\-ञ्जीत}}{\lemma{\skt{पूर्व्व\appbrk{}व\appbrk{}द्य\appbrk{}था\appbrk{}ब\appbrk{}लम्}}\Afootnote{\skt{पूर्वविधानेन मात्रां यथाबलम्}\ A; \skt{पूर्व्व\appbrk{}व\appbrk{}द्य\appbrk{}था\appbrk{}व\appbrk{}लम्}\ H}}\skt{~। }\edtext{\skt{जीर्ण्ण\-पू\-र्व्व\-व\-दा\-हारः}}{\lemma{\skt{जीर्ण्ण\appbrk{}पू\appbrk{}र्व्व\appbrk{}व\appbrk{}दा\appbrk{}हारः}}\Afootnote{H; \skt{जीर्णे पयः सर्पिरोदन इत्याहारः, पूर्ववच्चात्र}\ A}}\skt{ }\edtext{\skt{परिहारश्च}}{\lemma{\skt{परिहारश्च}}\Afootnote{H; \skt{परीहारः,}\ A}}\skt{~। }\edtext{\skt{एते\-नो\-र्ध्व\-म\-धश्च}}{\lemma{\skt{च}}\Afootnote{H; \skt{तिर्य}\ A}}\skt{ }\edtext{\skt{क्रिमयोगच्छन्ति}}{\lemma{\skt{क्रिमयोगच्छन्ति}}\Afootnote{H; \skt{क्कृमयो निष्क्रामन्ति,}\ A}}\skt{~। }\edtext{\skt{अलक्ष्मीर}}{\lemma{\skt{अलक्ष्मीर}}\Afootnote{H; \skt{अलक्ष्मीर्}\ A}}\skt{ अपक्रामति~। पुष्करवर्ण्णः }\edtext{\skt{श्रुतनिगादी}}{\lemma{\skt{श्रुतनिगादी}}\Afootnote{H; \skt{स्थिरवयाः श्रुतनिगादी}\ A}}\skt{ त्रिव\-र्ष\-श\-ता\-यु\-र्भ\-वति~। एतदेव }\edtext{\edtext{\skt{विष\-कु\-ष्ठ\-वि\-ष\-म\-ज्व\-र\-म\-प\-स्मा\-रो\-न्मा\-द\-भू\-त\-ग्र\-हेषु}}{\lemma{\skt{विष\appbrk{}कु\appbrk{}ष्ठ\appbrk{}वि\appbrk{}ष\appbrk{}म\appbrk{}ज्व\appbrk{}रम्}}\Afootnote{H; \skt{कुष्ठ\appbrk{}वि\appbrk{}ष\appbrk{}म\appbrk{}ज्वरा}\ A}}}{\lemma{\skt{अप\appbrk{}स्मा\appbrk{}रो\appbrk{}न्मा\appbrk{}द\appbrk{}भू\appbrk{}त\appbrk{}ग्र\appbrk{}हेषु}}\Afootnote{H; \skt{पस्मा\appbrk{}रो\appbrk{}न्मा\appbrk{}द\appbrk{}वि\appbrk{}ष\appbrk{}भू\appbrk{}त\appbrk{}ग्र\appbrk{}हे\appbrk{}ष्व\appbrk{}न्येषु च महाव्याधिषु}\ A}}\skt{ }\edtext{\skt{संशो\-ध\-न\-मु\-प\-दि\-शन्ति}}{\lemma{\skt{उपदिशन्ति}}\Afootnote{H; \skt{आदिशन्ति~॥}\ A}}\skt{~।}\pend

\pstart \ledsidenote{\passref{4.28.7}}\skt{हृतदोष }\edtext{\skt{एवा\-गा\-र\-म्प्र\-विश्य}}{\lemma{\skt{प्रविश्य}}\Afootnote{A; \skt{प्रनिश्य}\ H}}\skt{ हैमवत्या वचायाः }\edtext{\edtext{\edtext{\skt{पिण्ड\-मा\-म\-ल\-क\-मा\-त्र\-म\-भि\-हु\-त\-न्कृत्वा}}{\lemma{\skt{पिण्डम्}}\Afootnote{H; \skt{पिण्ड}\ A}}}{\lemma{\skt{आम\appbrk{}ल\appbrk{}क\appbrk{}मा\appbrk{}त्रम्}}\Afootnote{H; \skt{माम\appbrk{}ल\appbrk{}क\appbrk{}मा\appbrk{}त्रम्}\ A}}}{\lemma{\skt{अभिहुतन्कृत्वा}}\Afootnote{H; \skt{अभिहुतं}\ A}}\skt{ }\edtext{\skt{पयसालो पिवेत्~। जीर्ण्णे}}{\lemma{\skt{पयसालो पिवेत्~। जीर्ण्णे}}\Afootnote{\skt{पयसाऽऽआलोड्य पिबेत्, जीर्णे}\ A; \skt{पय\appbrk{}सा\appbrk{}लो\appbrk{}पि\appbrk{}वेत्~। जीर्ण्णे}\ H}}\skt{ पयः सर्प्पिरोदन इत्याहारः~। एवं द्वाद\-श\-रा\-त्र\-मु\-प\-युं\-जीत~। ततोऽस्य श्रोत्रं }\edtext{\skt{विधियते}}{\lemma{\skt{विधियते}}\Afootnote{H; \skt{विव्रियते,}\ A}}\skt{~। द्विर\-भ्या\-सा\-त्स्मृ\-ति\-मा\-न्भ\-वति~। }\edtext{\skt{त्रिरभ्यस्य}}{\lemma{\skt{त्रिरभ्यस्य}}\Afootnote{H; \skt{त्रिरभ्यासा}\ A}}\skt{ }\edtext{\edtext{\skt{श्रुतमाधत्ते}}{\lemma{\skt{श्रुतम्}}\Afootnote{H; \skt{च्छ्रुत}\ A}}}{\lemma{\skt{आधत्ते}}\Afootnote{H; \skt{मादत्ते,}\ A}}\skt{~। }\edtext{\skt{चतु\-र्द्वा\-द\-श\-रा\-त्र\-म\-भ्यस्य}}{\lemma{\skt{अभ्यस्य}}\Afootnote{H; \skt{उपयुज्य सर्वं तर}\ A}}\skt{ }\edtext{\edtext{\edtext{\edtext{\skt{तत\-स्ता\-र्क्ष्य\-न्द\-र्श\-न\-मु\-त्प\-द्यते}}{\lemma{\skt{ततस्}}\Afootnote{H; \skt{ति किल्बिषं,}\ A}}}{\lemma{\skt{तार्क्ष्यन्}}\Afootnote{H; \skt{तार्क्ष्य}\ A}}}{\lemma{\skt{दर्शनम्}}\Afootnote{H; \skt{दर्शन}\ A}}}{\lemma{\skt{उत्पद्यते}}\Afootnote{H; \skt{मुत्पद्यते,}\ A}}\skt{~। शतायुश्च }\edtext{\skt{भवति~। द्वे}}{\lemma{\skt{भवति~। द्वे}}\Afootnote{H; \skt{भवति~। द्वे द्वे}\ A}}\skt{ पले }\edtext{\skt{इतस्या}}{\lemma{\skt{इतस्या}}\Afootnote{H; \skt{इतरस्या}\ A}}\skt{ }\edtext{\skt{वचायाः}}{\lemma{\skt{वचायाः}}\Afootnote{H; \skt{वचाया}\ A}}\skt{ }\edtext{\edtext{\skt{क्वाथ्य\-म्पि\-वे\-त्प\-यसा}}{\lemma{\skt{क्वाथ्यम्}}\Afootnote{H; \skt{विष्क्वाथ्य}\ A}}}{\lemma{\skt{पिवेत्पयसा}}\Afootnote{H; \skt{पिबेत्पयसा,}\ A}}\skt{ }\edtext{\skt{पूर्व्वशाशिषश्च}}{\lemma{\skt{पूर्व्वशाशिषश्च}}\Afootnote{H; \skt{समानं भोजनं, समाः पूर्वेणाशिषश्च~॥}\ A}}\skt{ }\edtext{\skt{समानाः}}{\lemma{\skt{समानाः}}\Afootnote{H; \emph{om.}\ A}}\skt{~।}\pend

\pstart \ledsidenote{\passref{4.28.8}}\skt{वचा\-श\-त\-पा\-कम्वा }\edtext{\edtext{\edtext{\skt{सर्प्पि\-र्द्रो\-ण\-मु\-पय्ञ्ज्य}}{\lemma{\skt{सर्प्पिर्}}\Afootnote{H; \skt{सर्पि}\ A}}}{\lemma{\skt{द्रोणम्}}\Afootnote{H; \skt{र्द्रोणम्}\ A}}}{\lemma{\skt{उपय्ञ्ज्य}}\Afootnote{H; \skt{उपयुज्य}\ A}}\skt{ }\edtext{\edtext{\skt{पञ्च\-व\-र्ष\-श\-ता\-यु\-र्भ\-वति}}{\lemma{\skt{पञ्च\appbrk{}व\appbrk{}र्ष\appbrk{}श\appbrk{}ता\appbrk{}युर्}}\Afootnote{H; \skt{पञ्च\appbrk{}व\appbrk{}र्ष\appbrk{}श\appbrk{}तायु}\ A}}}{\lemma{\skt{भवति}}\Afootnote{H; \skt{र्भवति,}\ A}}\skt{~। }\edtext{\edtext{\skt{गल\-ग\-ण्डा\-प\-ची\-श्ली\-प\-द\-स्व\-र\-र\-भे\-दां\-श्चा\-प\-हन्ति}}{\lemma{\skt{गल\appbrk{}ग\appbrk{}ण्डा\appbrk{}प\appbrk{}ची\appbrk{}श्ली\appbrk{}प\appbrk{}द\appbrk{}स्व\appbrk{}र\appbrk{}र\appbrk{}भे\appbrk{}दांश्}}\Afootnote{H; \skt{गल\appbrk{}ग\appbrk{}ण्डा\appbrk{}प\appbrk{}ची\appbrk{}श्ली\appbrk{}प\appbrk{}द\appbrk{}स्व\appbrk{}र\appbrk{}मे\appbrk{}दांश्}\ A}}}{\lemma{\skt{चापहन्ति}}\Afootnote{H; \skt{चापहन्ती}\ A}}\skt{~॥ }\edtext{\skt{भवन्ति}}{\lemma{\skt{भवन्ति}}\Afootnote{H; \skt{ति~॥}\ A}}\skt{~॥}\pend

\pstart\versepar \ledsidenote{\passref{4.28.26}}\edtext{\edtext{\skt{पाप्मा\-न\-मु\-प\-ह\-न्त्येताः}}{\lemma{\skt{उपहन्त्य्}}\Afootnote{\skt{नाशयन्त्य्}\ A; \skt{उपहन्}\ H}}}{\lemma{\skt{एताः}}\Afootnote{\skt{एता}\ A; \emph{om.}\ H}}\skt{ }\edtext{\skt{श्रियं}}{\lemma{\skt{श्रियं}}\Afootnote{\skt{दद्युश्चौषधयः श्रियम्~।}\ A; \skt{त्येताः श्रियं}\ H}}\skt{ }\edtext{\skt{दद्यु\-त्त\-था\-उ\-षधीः}}{\lemma{\skt{दद्यु\appbrk{}त्त\appbrk{}था\appbrk{}उ\appbrk{}षधीः}}\Afootnote{\emph{om.}\ A; \skt{दद्यु\appbrk{}त्त\appbrk{}थ\appbrk{}औ\appbrk{}षधीः~।}\ H}}\skt{~। }\pend
\pstart\versepar \skt{}\edtext{\edtext{\edtext{\skt{कुर्यु\-न्ना\-ग\-ब\-ल\-म्वापि}}{\lemma{\skt{कुर्युन्}}\Afootnote{H; \skt{कुर्युर्}\ A}}}{\lemma{\skt{नागबलम्}}\Afootnote{A; \skt{नागवलम्}\ H}}}{\lemma{\skt{वापि}}\Afootnote{H; \skt{चापि}\ A}}\skt{ मनु\-ष्य\-म\-म\-रो\-पमं~॥}\pend

\pstart\versepar \ledsidenote{\passref{4.28.27}}\skt{सत\-ता\-ध्य\-य\-न\-म्वादः पर\-त\-न्त्रा\-व\-लो\-कनं~। }\pend
\pstart\versepar \skt{}\edtext{\skt{तद्वि\-द्या\-चा\-र्य\-सेवा च}}{\lemma{\skt{तद्वि\appbrk{}द्या\appbrk{}चा\appbrk{}र्य\appbrk{}सेवा च}}\Afootnote{H; \skt{तद्विद्या चार्यसेवा च}\ A}}\skt{ }\edtext{\skt{बुद्धिमेधाकराः}}{\lemma{\skt{बुद्धिमेधाकराः}}\Afootnote{\skt{बुद्धिमेधाकरो}\ A; \skt{वुद्धिमेधाकराः}\ H}}\skt{ }\edtext{\skt{गणः}}{\lemma{\skt{गणः}}\Afootnote{H; \skt{गुणः(गणः)~॥}\ A}}\skt{~।}\pend

\pstart\versepar \ledsidenote{\passref{4.28.28}}\skt{आयुष्यं भोजनञ्जीर्ण्णे }\edtext{\skt{वेगानाञ्च}}{\lemma{\skt{वेगानाञ्च}}\Afootnote{H; \skt{वेगानां चा}\ A}}\skt{ विधारणं~। }\pend
\pstart\versepar \skt{}\edtext{\skt{ब्रह्म\-च\-र्य\-म\-हिंसा}}{\lemma{\skt{ब्रह्मचर्यम्}}\Afootnote{A; \skt{व्रह्मचर्यम्}\ H}}\skt{ च साहसानाञ्च }\edtext{\skt{वर्ज्जनमिति}}{\lemma{\skt{इति}}\Afootnote{H; \emph{om.}\ A}}\skt{~॥~॥}\pend

\pstart \ledsidenote{\passref{4.28.trailer}}\edtext{\skt{॥ २८~॥}}{\lemma{\skt{॥ २८~॥}}\Afootnote{\skt{इति सुश्रु\appbrk{}त\appbrk{}सं\appbrk{}हि\appbrk{}तायां चिकित्सास्थाने मेधायुष्कामीयं रसायनं नामा\appbrk{}ष्टा\appbrk{}विं\appbrk{}शोऽ\appbrk{}ध्यायः~॥२८~॥}\ A; \skt{सा थ~॥~॥}\ H}}\pend

\endnumbering
\beginnumbering
\eledsection*{Adhyāya 29}
\pstart \ledsidenote{\passref{4.29.1}}\skt{अथातः स्वभा\-व\-व्या\-धि\-प्र\-ति\-षे\-ध\-नीयं रसा\-य\-न\-म्व्या\-ख्या\-स्यामः~॥}\pend

\pstart\versepar \ledsidenote{\passref{4.29.3}}\edtext{\skt{व्रह्मा}}{\lemma{\skt{व्रह्मा}}\Afootnote{H; \skt{ब्रह्मा}\ A}}\skt{ }\edtext{\skt{यद\-जृ\-त्पू\-र्व्व\-म\-मृतं}}{\lemma{\skt{यद\appbrk{}जृ\appbrk{}त्पू\appbrk{}र्व्व\appbrk{}म\appbrk{}मृतं}}\Afootnote{H; \skt{दयोऽ\appbrk{}सृ\appbrk{}ज\appbrk{}न्पू\appbrk{}र्व\appbrk{}म\appbrk{}मृतं}\ A}}\skt{ सोमसंज्ञितं~। }\pend
\pstart\versepar \skt{जरा\-मृ\-त्यु\-वि\-ना\-शाय विधानन्तस्य वक्ष्यते॥}\pend

\pstart \ledsidenote{\passref{4.29.4}}\skt{एक एव खलु भगवान्सोमः }\edtext{\edtext{\edtext{\edtext{\skt{स्थाना\-ना\-मा\-कृ\-ति\-र्व्वि\-शे\-षै\-श्च\-तु\-र्व्विं\-ङ्श\-तिधा}}{\lemma{\skt{स्थानानाम्}}\Afootnote{H; \skt{स्थानना}\ A}}}{\lemma{\skt{आकृतिर्}}\Afootnote{H; \skt{माकृति}\ A}}}{\lemma{\skt{व्विशेषैश्}}\Afootnote{H; \skt{वीर्यविशेषैश्}\ A}}}{\lemma{\skt{चतुर्व्विंङ्शतिधा}}\Afootnote{H; \skt{चतुर्विंशतिधा}\ A}}\skt{ भिद्यते~॥}\pend

\pstart \ledsidenote{\passref{4.29.5}}\edtext{\skt{तद्यथा}}{\lemma{\skt{तद्यथा}}\Afootnote{H; \emph{om.}\ A}}\skt{~॥ }\edtext{\skt{अंशु\-मा\-न्मू\-ञ्ज\-वा\-श्चैव}}{\lemma{\skt{अंशु\appbrk{}मा\appbrk{}न्मू\appbrk{}ञ्ज\appbrk{}वाश्}}\Afootnote{H; \skt{अंशु\appbrk{}मा\appbrk{}न्मु\appbrk{}ञ्ज\appbrk{}वांश्}\ A}}\skt{ }\edtext{\skt{चन्द्र\-मा\-र\-ज\-त\-प्रभः}}{\lemma{\skt{चन्द्र\appbrk{}मा\appbrk{}र\appbrk{}ज\appbrk{}त\appbrk{}प्रभः}}\Afootnote{H; \skt{चन्द्रमा रजतप्रभः~। दूर्वासोमः कनीयांश्च श्वेताक्षः कनकप्रभः~॥}\ A}}\skt{~।}\pend

\pstart\versepar \ledsidenote{\passref{4.29.6}}\edtext{\skt{प्रता\-न\-वा\-स्ता\-ल\-वृन्तः}}{\lemma{\skt{तालवृन्तः}}\Afootnote{H; \skt{स्तालवृन्तः}\ A}}\skt{ }\edtext{\skt{कर\-वी\-रा\-ङ्शु\-स\-न्निभाः}}{\lemma{\skt{कर\appbrk{}वी\appbrk{}रा\appbrk{}ङ्शु\appbrk{}स\appbrk{}न्निभाः}}\Afootnote{H; \skt{कर\appbrk{}वी\appbrk{}रोऽं\appbrk{}श\appbrk{}वा\appbrk{}नपि~।}\ A}}\skt{~॥ }\pend
\pstart\versepar \skt{स्वयम्प्रभो महासोमो यश्चापि गरुडाहृतः~।}\pend

\pstart\versepar \ledsidenote{\passref{4.29.7}}\edtext{\skt{गायत्र्यस्त्रैष्टुभः}}{\lemma{\skt{गायत्र्यस्}}\Afootnote{H; \skt{गायत्रस्}\ A}}\skt{ }\edtext{\skt{पांक्त्यो}}{\lemma{\skt{पांक्त्यो}}\Afootnote{H; \skt{पाङ्क्तो}\ A}}\skt{ }\edtext{\skt{जाग्रतः}}{\lemma{\skt{जाग्रतः}}\Afootnote{H; \skt{जागतः}\ A}}\skt{ }\edtext{\skt{सात्करस्तथा}}{\lemma{\skt{सात्करस्}}\Afootnote{H; \skt{शाक्वरस्}\ A}}\skt{~॥ }\pend
\pstart\versepar \skt{अग्निष्टोमो }\edtext{\skt{रेवतश्च }\edtext{\skt{यश्चोक्थ}}{\lemma{\skt{चोक्थ}}\Afootnote{H; \skt{थोक्त}\ A}}}{\lemma{\skt{रेवतश्च यश्}}\Afootnote{H; \skt{रैवतश्च य}\ A}}\skt{ इति संज्ञितः~।}\pend

\pstart\versepar \ledsidenote{\passref{4.29.8}}\skt{गायत्र्या त्रिपदा युक्तो यश्चो\-डु\-प\-ति\-रु\-च्यते~॥ }\pend
\pstart\versepar \skt{एते सोमाः समाख्याता }\edtext{\skt{वेदोक्तैर्न्नामभिः}}{\lemma{\skt{न्नामभिः}}\Afootnote{H; \skt{नामभिः}\ A}}\skt{ शुभैः~।}\pend

\pstart\versepar \ledsidenote{\passref{4.29.9}}\skt{सर्व्वेषामेव }\edtext{\skt{चैतेषामेको}}{\lemma{\skt{चैतेषामेको}}\Afootnote{H; \skt{चैतेषामेको}\ A}}\skt{ विधिरुपासने~॥ }\pend
\pstart\versepar \skt{सर्व्वे तुल्यगुणाश्चैव विधानं }\edtext{\skt{सम्प्रवक्ष्यते}}{\lemma{\skt{सम्प्रवक्ष्यते}}\Afootnote{H; \skt{तेषु वक्ष्यते~॥}\ A}}\skt{~॥}\pend

\pstart \ledsidenote{\passref{4.29.10}}\edtext{\skt{अथो}}{\lemma{\skt{अथो}}\Afootnote{H; \skt{अतो}\ A}}\skt{ }\edtext{\skt{न्यतम}}{\lemma{\skt{न्यतम}}\Afootnote{H; \skt{ऽन्यतमं}\ A}}\skt{ }\edtext{\skt{सोम\-मु\-प\-यु\-युंक्षुः}}{\lemma{\skt{उपयुयुंक्षुः}}\Afootnote{H; \skt{उपयुयुक्षुः}\ A}}\skt{ सर्व्वो\-प\-क\-र\-ण\-प\-रि\-चा\-र\-को\-पेतः~। प्रशस्ते देशे }\edtext{\skt{त्रिवृ\-त\-म\-गा\-र\-ङ्का\-र\-यित्वा}}{\lemma{\skt{त्रिवृ\appbrk{}त\appbrk{}म\appbrk{}गा\appbrk{}र\appbrk{}ङ्का\appbrk{}र\appbrk{}यित्वा}}\Afootnote{H; \skt{त्रिवृतमागारं कारयित्वा}\ A}}\skt{ हृतदोषः प्रति\-सं\-सृ\-ष्ट\-भक्तः }\edtext{\skt{प्रशस्ते}}{\lemma{\skt{प्रशस्ते}}\Afootnote{H; \skt{प्रशस्तेषु}\ A}}\skt{ }\edtext{\edtext{\skt{तिथि\-क\-र\-ण\-मु\-हू\-र्त्त\-ङ\-अ\-शु\-म\-न्त\-मा\-दा\-या\-ध्व\-र\-क\-ल्पे\-ना\-भि\-ष्टुत्य}}{\lemma{\skt{तिथि\appbrk{}क\appbrk{}र\appbrk{}ण\appbrk{}मु\appbrk{}हू\appbrk{}र्त्त\appbrk{}ङ\appbrk{}अ\appbrk{}शु\appbrk{}म\appbrk{}न्तम्}}\Afootnote{H; \skt{तिथि\appbrk{}क\appbrk{}र\appbrk{}ण\appbrk{}मु\appbrk{}हू\appbrk{}र्त\appbrk{}न\appbrk{}क्ष\appbrk{}त्रेषु अंशुमन्तम्}\ A}}}{\lemma{\skt{आदा\appbrk{}या\appbrk{}ध्व\appbrk{}र\appbrk{}क\appbrk{}ल्पे\appbrk{}ना\appbrk{}भि\appbrk{}ष्टुत्य}}\Afootnote{H; \skt{आदा\appbrk{}या\appbrk{}ध्व\appbrk{}र\appbrk{}क\appbrk{}ल्पे\appbrk{}ना\appbrk{}हृ\appbrk{}त\appbrk{}म\appbrk{}भि\appbrk{}षु\appbrk{}त\appbrk{}म\appbrk{}भि\appbrk{}हुतं}\ A}}\skt{ }\edtext{\skt{यथो\-क्ता\-गा\-र\-स्कृ\-त\-म\-ङ्गलः}}{\lemma{\skt{यथो\appbrk{}क्ता\appbrk{}गा\appbrk{}र\appbrk{}स्कृ\appbrk{}त\appbrk{}म\appbrk{}ङ्गलः}}\Afootnote{H; \skt{चान्तरागारे कृत\appbrk{}म\appbrk{}ङ्ग\appbrk{}ल\appbrk{}स्व\appbrk{}स्ति\appbrk{}वा\appbrk{}चनः}\ A}}\skt{~। }\edtext{\skt{सोम\-च\-न्द्रं\-क\-न्दं\-सु\-व\-र्ण्ण\-सूच्या}}{\lemma{\skt{सोम\appbrk{}च\appbrk{}न्द्रं\appbrk{}क\appbrk{}न्दं\appbrk{}सु\appbrk{}व\appbrk{}र्ण्ण\appbrk{}सूच्या}}\Afootnote{H\textsuperscript{pc}; \skt{सोमकन्दं सुवर्णसूच्या}\ A; \skt{सोम\appbrk{}च\appbrk{}न्द्रं\appbrk{}सु\appbrk{}व\appbrk{}र्ण्ण\appbrk{}सूच्या}\ H\textsuperscript{ac}}}\skt{ विदार्य पयो गृह्णीयात्~। सौवर्ण्णे }\edtext{\skt{राजते}}{\lemma{\skt{राजते}}\Afootnote{H; \skt{(राजते वा)}\ A}}\skt{ }\edtext{\skt{पात्रेञ्जलिमात्रं}}{\lemma{\skt{पात्रेञ्जलिमात्रं}}\Afootnote{H; \skt{पात्रेऽञ्जलिमात्रं,}\ A}}\skt{ }\edtext{\skt{ततो}}{\lemma{\skt{ततो}}\Afootnote{H; \skt{ततः}\ A}}\skt{ }\edtext{\skt{सकृदेवोप}}{\lemma{\skt{सकृदेवोप}}\Afootnote{H; \skt{सकृदेवोप}\ A}}\skt{ }\edtext{\skt{युञ्जीतो}}{\lemma{\skt{युञ्जीतो}}\Afootnote{H; \skt{युञ्जीत}\ A}}\skt{ }\edtext{\skt{नास्वा\-द\-य\-मा\-नस्त}}{\lemma{\skt{नास्वा\appbrk{}द\appbrk{}य\appbrk{}मा\appbrk{}नस्त}}\Afootnote{H; \skt{नास्वादयन्, तत}\ A}}\skt{ उपस्पृश्य }\edtext{\skt{शेष\-म\-प्स्व\-स्त्व\-व\-साद्य }\edtext{\skt{यम\-नि\-य\-मा\-भ्या\-मा\-त्मानं}}{\lemma{\skt{अनि\appbrk{}य\appbrk{}मा\appbrk{}भ्या\appbrk{}मा\appbrk{}त्मानं}}\Afootnote{H; \skt{मनि\appbrk{}य\appbrk{}मा\appbrk{}भ्या\appbrk{}मा\appbrk{}त्मानं}\ A}}}{\lemma{\skt{अस्त्ववसाद्य यम्}}\Afootnote{H\textsuperscript{pc}; \skt{अवसाद्य य}\ A; \skt{अवसाद्य यम्}\ H\textsuperscript{ac}}}\skt{ संयोज्य }\edtext{\skt{वाग्यतोऽभ्यन्तरः}}{\lemma{\skt{ऽभ्यन्तरः}}\Afootnote{H; \skt{ऽभ्यन्तरतः}\ A}}\skt{~। }\edtext{\skt{सुहृ\-द्भि\-र\-न्वा\-स्य\-मानो}}{\lemma{\skt{अन्वास्यमानो}}\Afootnote{H; \skt{उपास्यमानो}\ A}}\skt{ विहरेत्~।}\pend

\pstart\versepar \ledsidenote{\passref{4.29.11}}\skt{रसा\-य\-न\-म्पी\-त\-वांस्तु }\edtext{\skt{निर्व्वाते}}{\lemma{\skt{निर्व्वाते}}\Afootnote{H; \skt{निवाते}\ A}}\skt{ }\edtext{\skt{नियतात्मना}}{\lemma{\skt{नियतात्मना}}\Afootnote{H; \skt{तन्मनाः}\ A}}\skt{ }\pend
\pstart\versepar \skt{}\edtext{\skt{शुचि\-रा\-सी\-त\-स्ति\-ष्ठे\-च्चं\-क्र\-म्पान्न}}{\lemma{\skt{शुचि\appbrk{}रा\appbrk{}सी\appbrk{}त\appbrk{}स्ति\appbrk{}ष्ठे\appbrk{}च्चं\appbrk{}क्र\appbrk{}म्पान्न}}\Afootnote{H; \skt{शुचिः~। आसीत तिष्ठेत्क्रामेच्च न}\ A}}\skt{ कथञ्चन सम्विशेत्~।}\pend

\pstart \ledsidenote{\passref{4.29.12}}\edtext{\skt{न}}{\lemma{\skt{न}}\Afootnote{K; \skt{सायं वा भुक्त\appbrk{}वा\appbrk{}नु\appbrk{}प\appbrk{}श्रु\appbrk{}त\appbrk{}शान्तिः कुशशय्यायां कृष्णा\appbrk{}जि\appbrk{}नो\appbrk{}त्त\appbrk{}रायां सुहृ\appbrk{}द्भि\appbrk{}रु\appbrk{}पा\appbrk{}स्य\appbrk{}मानः शयीत, तृषितो वा शीतोदकमात्रां पिबेत् (अशनायितो वा क्षीरं); ततः प्रात\appbrk{}रु\appbrk{}त्था\appbrk{}यो\appbrk{}प\appbrk{}श्रु\appbrk{}त\appbrk{}शान्तिः कृतमङ्गलो गां स्पृष्ट्वा ततहैवासीत, तस्य जीर्णे सोमे छर्दिरुत्पद्यते, ततः शोणिताक्तं कृमिव्यामिश्रं छर्दितवते सायं शृतशीतं क्षीरं वितरेत्; तत\appbrk{}स्तृ\appbrk{}ती\appbrk{}येऽ\appbrk{}हनि कृमि\appbrk{}व्या\appbrk{}मि\appbrk{}श्र\appbrk{}म\appbrk{}ति\appbrk{}सा\appbrk{}र्यते, स तेना\appbrk{}नि\appbrk{}ष्ट\appbrk{}प्र\appbrk{}ति\appbrk{}ग्र\appbrk{}ह\appbrk{}भु\appbrk{}क्त\appbrk{}प्र\appbrk{}भृ\appbrk{}ति\appbrk{}भि\appbrk{}र्वि\appbrk{}शे\appbrk{}षै\appbrk{}र्वि\appbrk{}नि\appbrk{}र्मुक्तः शुद्ध\appbrk{}त\appbrk{}नु\appbrk{}र्भ\appbrk{}वति, ततः सायं स्नाताय पूर्ववदेव क्षीरं वितरेत्, क्षौम\appbrk{}व\appbrk{}स्त्रा\appbrk{}स्तृ\appbrk{}तायां चैनं शय्यायां शाययेत्; तत\appbrk{}श्च\appbrk{}तु\appbrk{}र्थेऽ\appbrk{}हनि तस्य श्वय\appbrk{}थु\appbrk{}रु\appbrk{}त्प\appbrk{}द्यते, ततः सर्वाङ्गेभ्यः कृमयो निष्क्रामन्ति, तदहश्च शय्यायां पांशु\appbrk{}भि\appbrk{}र\appbrk{}व\appbrk{}की\appbrk{}र्य\appbrk{}माणः शयीत, ततः सायं पूर्ववदेव क्षीरं वितरेत्; एवं पञ्च\appbrk{}म\appbrk{}ष\appbrk{}ष्ठ\appbrk{}यो\appbrk{}र्दि\appbrk{}व\appbrk{}स\appbrk{}यो\appbrk{}र्व\appbrk{}र्तेत, केव\appbrk{}ल\appbrk{}मु\appbrk{}भ\appbrk{}य\appbrk{}का\appbrk{}ल\appbrk{}मस्मै क्षीरं वितरेत्; ततः सप्तमेऽहनि निर्मां\appbrk{}स\appbrk{}स्त्व\appbrk{}ग\appbrk{}स्थि\appbrk{}भूतः केवलं सोम\appbrk{}प\appbrk{}रि\appbrk{}ग्र\appbrk{}हा\appbrk{}दे\appbrk{}वो\appbrk{}च्छ्व\appbrk{}सिति, तदहश्च क्षीरेण सुखोष्णेन परिषिच्य तिल\appbrk{}म\appbrk{}धु\appbrk{}क\appbrk{}च\appbrk{}न्द\appbrk{}ना\appbrk{}नु\appbrk{}लि\appbrk{}प्त\appbrk{}देहं पयः पाययेत्; ततोऽष्टमेऽहनि प्रातरेव क्षीरपरिषिक्तं चन्द\appbrk{}न\appbrk{}प्र\appbrk{}दि\appbrk{}ग्धा\appbrk{}गात्रं पयः पाययित्वा पांशुशय्यां समुत्सृज्य क्षौम\appbrk{}व\appbrk{}स्त्रा\appbrk{}स्तृ\appbrk{}तायां शय्यायां शाययेत्, ततोऽस्य मांसमाप्याय्यते, त्वक्चावदलति, दन्त\appbrk{}न\appbrk{}ख\appbrk{}रो\appbrk{}माणि चास्य पतन्ति; तस्य नव\appbrk{}म\appbrk{}दि\appbrk{}व\appbrk{}सा\appbrk{}त्प्र\appbrk{}भृ\appbrk{}त्य\appbrk{}णु\appbrk{}तै\appbrk{}ला\appbrk{}भ्यङ्गः सोम\appbrk{}व\appbrk{}ल्क\appbrk{}क\appbrk{}षा\appbrk{}य\appbrk{}प\appbrk{}रि\appbrk{}षेकः; ततो दश\appbrk{}मेऽ\appbrk{}ह\appbrk{}न्ये\appbrk{}त\appbrk{}देव वितरेत्, ततोऽस्य त्वक्स्थि\appbrk{}र\appbrk{}ता\appbrk{}मु\appbrk{}पैति; एव\appbrk{}मे\appbrk{}का\appbrk{}द\appbrk{}श\appbrk{}द्वा\appbrk{}द\appbrk{}श\appbrk{}यो\appbrk{}र्व\appbrk{}र्तेत; तत\appbrk{}स्त्र\appbrk{}यो\appbrk{}द\appbrk{}शा\appbrk{}त्प्र\appbrk{}भृति सोम\appbrk{}व\appbrk{}ल्क\appbrk{}क\appbrk{}षा\appbrk{}य\appbrk{}प\appbrk{}रि\appbrk{}षेकः, एव\appbrk{}मा\appbrk{}षो\appbrk{}ड\appbrk{}शा\appbrk{}द्व\appbrk{}र्तेत; ततः सप्त\appbrk{}द\appbrk{}शा\appbrk{}ष्टा\appbrk{}द\appbrk{}श\appbrk{}यो\appbrk{}र्दि\appbrk{}व\appbrk{}स\appbrk{}यो\appbrk{}र्द\appbrk{}शना जायन्ते शिखरिणः स्निग्ध\appbrk{}व\appbrk{}ज्र\appbrk{}वै\appbrk{}दू\appbrk{}र्य\appbrk{}स्फ\appbrk{}टि\appbrk{}क\appbrk{}प्र\appbrk{}काशाः समाः स्थिराः सहिष्णवः, तदा प्रभृति चानवैः शालितण्डुलैः क्षीर\appbrk{}य\appbrk{}वा\appbrk{}गू\appbrk{}मु\appbrk{}प\appbrk{}से\appbrk{}वेत याव\appbrk{}त्प\appbrk{}ञ्च\appbrk{}विं\appbrk{}श\appbrk{}ति\appbrk{}रिति; ततोऽस्मै दद्याच्छाल्योदनं मृदूभयकालं पयसा, ततोऽस्य नखा जायन्ते विद्रु\appbrk{}मे\appbrk{}न्द्र\appbrk{}गो\appbrk{}प\appbrk{}क\appbrk{}त\appbrk{}रु\appbrk{}णा\appbrk{}दि\appbrk{}त्य\appbrk{}प्र\appbrk{}काशाः, स्थिराः स्निग्धा लक्षणसंपन्नाः केशाश्च सूक्ष्मा जायन्ते, त्वक्च नीलो\appbrk{}त्प\appbrk{}ला\appbrk{}त\appbrk{}सी\appbrk{}पु\appbrk{}ष्प\appbrk{}वै\appbrk{}दू\appbrk{}र्य\appbrk{}प्र\appbrk{}काशा; ऊर्ध्वं च मासा\appbrk{}त्के\appbrk{}शा\appbrk{}न्वा\appbrk{}प\appbrk{}येत्, वापयित्वा चोशी\appbrk{}र\appbrk{}च\appbrk{}न्द\appbrk{}न\appbrk{}कृ\appbrk{}ष्ण\appbrk{}ति\appbrk{}ल\appbrk{}कल्कैः शिरः प्रदिह्यात्पयसा वा स्नापयेत्; ततोऽस्यानन्तरं सप्तरात्रात्केशा जायन्ते भ्रम\appbrk{}रा\appbrk{}ञ्ज\appbrk{}न\appbrk{}निभाः कुञ्चिताः स्थिराः स्निग्धाः; तत\appbrk{}स्त्रि\appbrk{}रा\appbrk{}त्रा\appbrk{}त्प्र\appbrk{}थ\appbrk{}मा\appbrk{}व\appbrk{}स\appbrk{}थ\appbrk{}प\appbrk{}रि\appbrk{}स\appbrk{}रा\appbrk{}न्नि\appbrk{}ष्क्रम्य मुहूर्तं स्थित्वा पुनरेवान्तः प्रविशेत्, ततोऽस्य बला\appbrk{}तै\appbrk{}ल\appbrk{}म\appbrk{}भ्य\appbrk{}ङ्गा\appbrk{}र्थेऽ\appbrk{}व\appbrk{}चार्यं, यव\appbrk{}पि\appbrk{}ष्ट\appbrk{}मु\appbrk{}द्व\appbrk{}र्त\appbrk{}नार्थे, सुखोष्णं च पयः परिषेकार्थे, अज\appbrk{}क\appbrk{}र्ण\appbrk{}क\appbrk{}षा\appbrk{}य\appbrk{}मु\appbrk{}त्सा\appbrk{}द\appbrk{}नार्थे, सोशीरं कूपोदकं स्नानार्थे, चन्द\appbrk{}न\appbrk{}म\appbrk{}नु\appbrk{}ले\appbrk{}पार्थे, आम\appbrk{}ल\appbrk{}क\appbrk{}र\appbrk{}स\appbrk{}वि\appbrk{}मि\appbrk{}श्रा\appbrk{}श्चास्य यूष\appbrk{}सू\appbrk{}प\appbrk{}वि\appbrk{}कल्पाः, क्षीर\appbrk{}म\appbrk{}धु\appbrk{}क\appbrk{}सिद्धं च कृष्ण\appbrk{}ति\appbrk{}ल\appbrk{}म\appbrk{}व\appbrk{}चा\appbrk{}र\appbrk{}णार्थे, एवं दशरात्रं; ततोऽ\appbrk{}न्य\appbrk{}द्द\appbrk{}श\appbrk{}रात्रं द्वितीये परिसरे वर्तेत; ततस्तृतीये परिसरे स्थिरी\appbrk{}कु\appbrk{}र्व\appbrk{}न्ना\appbrk{}त्मा\appbrk{}न\appbrk{}म\appbrk{}न्य\appbrk{}द्द\appbrk{}श\appbrk{}रा\appbrk{}त्र\appbrk{}मा\appbrk{}सीत, किञ्चि\appbrk{}दा\appbrk{}त\appbrk{}प\appbrk{}प\appbrk{}व\appbrk{}नान्वा सेवेत, पुनश्चान्तः प्रविशेत्, न}\ A; \skt{सायं वा भुक्त\appbrk{}वा\appbrk{}नु\appbrk{}प\appbrk{}श्रु\appbrk{}त\appbrk{}शा\appbrk{}न्तिः\appbrk{}कु\appbrk{}श\appbrk{}श\appbrk{}य्यायां कृष्णा\appbrk{}जि\appbrk{}नो\appbrk{}त्त\appbrk{}रायां सुहृ\appbrk{}द्भि\appbrk{}र\appbrk{}न्वा\appbrk{}स्य\appbrk{}मानः शयीत~। तृषितो वा शीतोदकं मात्राम्पिवेत्~। ततः प्रात\appbrk{}रु\appbrk{}त्था\appbrk{}यो\appbrk{}प\appbrk{}श्रु\appbrk{}त\appbrk{}शान्तिः कृतमङ्गलो गां स्पृष्ट्वा तथैवासीत्~। तस्य जीर्ण्णे सोमे च्छर्द्दिरुत्पद्यते~। ततः शोणि\appbrk{}ता\appbrk{}क्त\appbrk{}ङ्क्रि\appbrk{}मि\appbrk{}व्या\appbrk{}मि\appbrk{}श्र\appbrk{}ञ्छ\appbrk{}र्द्दि\appbrk{}त\appbrk{}वतः~। सायं स्रितशीतं क्षीर\appbrk{}म्वि\appbrk{}त\appbrk{}रे\appbrk{}त्त\appbrk{}त\appbrk{}स्तृ\appbrk{}ती\appbrk{}येऽ\appbrk{}हनि क्रिमि\appbrk{}व्या\appbrk{}मि\appbrk{}श्र\appbrk{}म\appbrk{}ति\appbrk{}सा\appbrk{}र्यते~। स तेना\appbrk{}ति\appbrk{}सृ\appbrk{}ष्ट\appbrk{}प्र\appbrk{}ति\appbrk{}ग्र\appbrk{}ह\appbrk{}भु\appbrk{}क्त\appbrk{}प्र\appbrk{}भृ\appbrk{}भि\appbrk{}र्वि\appbrk{}शैषैः प्रतिमुक्तः शुद्ध\appbrk{}त\appbrk{}नु\appbrk{}र्व्भ\appbrk{}वति~। ततः सायं स्नातस्य पूर्व्ववदेव क्षीरम्वितरेत्~। क्षौम\appbrk{}व\appbrk{}स्त्रा\appbrk{}स्तृ\appbrk{}ता\appbrk{}या\appbrk{}ञ्चैनं शय्यायां शाययीत~। तत\appbrk{}श्च\appbrk{}तु\appbrk{}र्थेऽ\appbrk{}हनि तस्य श्वय\appbrk{}थु\appbrk{}रु\appbrk{}प\appbrk{}जा\appbrk{}यते~। ततः सर्व्वाङ्गेभ्यः क्रिमयो निःक्रामन्ति~। तदहश्च पांशु शय्यायां पांशु\appbrk{}भि\appbrk{}र\appbrk{}व\appbrk{}की\appbrk{}र्य\appbrk{}माणः शयीत~। ततः साय\appbrk{}म्पू\appbrk{}व्व\appbrk{}व\appbrk{}देव क्षीर\appbrk{}म्वि\appbrk{}त\appbrk{}रे\appbrk{}देवं पञ्च\appbrk{}म\appbrk{}ष\appbrk{}ष्ठ\appbrk{}यो\appbrk{}र्द्दि\appbrk{}व\appbrk{}स\appbrk{}यो\appbrk{}र्व्व\appbrk{}र्तेत~। केव\appbrk{}ल\appbrk{}न्नू\appbrk{}भ\appbrk{}यतः कालमस्मै क्षीर\appbrk{}म्वि\appbrk{}त\appbrk{}रे\appbrk{}त्ततः~। सप्तमेहनि निर्म्मां\appbrk{}स\appbrk{}त्व\appbrk{}ग\appbrk{}स्ति\appbrk{}भूतः~। सोम\appbrk{}प्र\appbrk{}ति\appbrk{}ग्र\appbrk{}हा\appbrk{}दे\appbrk{}वो\appbrk{}च्छ\appbrk{}सिति~। तदहः क्षीक्षीरेण सुखोष्णेन परिषिच्य तिल\appbrk{}म\appbrk{}धु\appbrk{}क\appbrk{}च\appbrk{}न्द\appbrk{}ना\appbrk{}नु\appbrk{}लि\appbrk{}प्त\appbrk{}देहं पयः पाययेत्~। ततोऽष्टमेहनि प्रातरेव क्षीर\appbrk{}षि\appbrk{}क्त\appbrk{}श्च\appbrk{}न्द\appbrk{}न\appbrk{}प\appbrk{}रि\appbrk{}दि\appbrk{}ग्ध\appbrk{}गात्रः पयः पीत्वा पांशु\appbrk{}श\appbrk{}य्या\appbrk{}मु\appbrk{}त्सृज्य क्षौमास्तृतायां शयीत~। तस्य मांसान्याप्याय्यन्ते त्वक्चावदलति~। दन्त\appbrk{}न\appbrk{}ख\appbrk{}रो\appbrk{}माणि चास्यशीर्यन्ते~। तस्य नवम दिव सात्प्र\appbrk{}भृ\appbrk{}त्य\appbrk{}नु\appbrk{}तै\appbrk{}ल\appbrk{}म\appbrk{}भ्यङ्गः~। सोम\appbrk{}व\appbrk{}ल्क\appbrk{}क\appbrk{}षायः परिषेकः~। ततोऽस्य दश\appbrk{}मे\appbrk{}ह\appbrk{}न्ये\appbrk{}त\appbrk{}देव वितरेत्ततोस्य त्वक्स्थि\appbrk{}र\appbrk{}ता\appbrk{}मु\appbrk{}पैति~। एव मेकाद शद्वा\appbrk{}द\appbrk{}श\appbrk{}यो\appbrk{}र्व्व\appbrk{}र्त्तेत~। तत\appbrk{}स्त्र\appbrk{}यो\appbrk{}द\appbrk{}शा\appbrk{}त्प्र\appbrk{}भृ\appbrk{}ति\appbrk{}नैव सोम\appbrk{}व\appbrk{}ल्क\appbrk{}क\appbrk{}षायो दातव्यः~। एव\appbrk{}मा\appbrk{}षो\appbrk{}ड\appbrk{}शा\appbrk{}द्व\appbrk{}र्त्तेत~। तत\appbrk{}स्स\appbrk{}प्त\appbrk{}द\appbrk{}शा\appbrk{}ष्टा\appbrk{}द\appbrk{}श\appbrk{}यो\appbrk{}र्द्द\appbrk{}श\appbrk{}यो\appbrk{}र्द्द\appbrk{}न्ताजा येत शिखरिणः स्निग्धा\appbrk{}व\appbrk{}ज्र\appbrk{}वै\appbrk{}दू\appbrk{}र्य\appbrk{}स्फ\appbrk{}टि\appbrk{}क\appbrk{}स\appbrk{}न्नि\appbrk{}काशाः स्थिराः समाः सहिष्णवः~। तदा प्रभृति चानवैः शालितण्डुलैः क्षीरयवागूमु पसेवेत~। याव\appbrk{}दा\appbrk{}प\appbrk{}ञ्च\appbrk{}विं\appbrk{}श\appbrk{}ति\appbrk{}रिति~। ततोऽस्मै दद्याच्छाल्योदनं मृदुमुभयतः कालं~। ततोस्य नखा जायन्ते विद्रु\appbrk{}मे\appbrk{}न्द्र\appbrk{}गो\appbrk{}प\appbrk{}क\appbrk{}त\appbrk{}रु\appbrk{}णा\appbrk{}दित्य प्रकाशाः स्निग्धाः स्थिराः~। लक्षणसम्पन्नाः केशाश्च शूक्ष्मा जायन्ते~। त्वक्च नीलो\appbrk{}त्प\appbrk{}ला\appbrk{}त\appbrk{}सी\appbrk{}पु\appbrk{}ष्प\appbrk{}वै\appbrk{}दू\appbrk{}र्य\appbrk{}प्र\appbrk{}काशाः~। उर्द्ध्वञ्च मासा\appbrk{}त्के\appbrk{}शा\appbrk{}न्वा\appbrk{}प\appbrk{}येत~। वापयित्वा चोशी\appbrk{}र\appbrk{}च\appbrk{}न्द\appbrk{}न\appbrk{}ति\appbrk{}ल\appbrk{}कल्कैः शिरः प्रदिदेत~। पयसा च स्नेहयीत~। ततोस्यानन्तरं सप्तरात्रात्केशा जायन्ते~। अञ्ज\appbrk{}न\appbrk{}वै\appbrk{}दू\appbrk{}र्य\appbrk{}भ्र\appbrk{}म\appbrk{}र\appbrk{}निभाः कुञ्चिताग्राः~। स्निग्धा\appbrk{}शि\appbrk{}र\appbrk{}सो\appbrk{}प\appbrk{}भो\appbrk{}गा\appbrk{}ना\appbrk{}मा\appbrk{}ह\appbrk{}र्त्तारैः~। तत\appbrk{}स्त्रि\appbrk{}रा\appbrk{}त्रा\appbrk{}त्प्र\appbrk{}थ\appbrk{}म\appbrk{}प\appbrk{}रि\appbrk{}स\appbrk{}रा\appbrk{}न्नि\appbrk{}ष्क्रम्य मुहूर्त्तं स्थित्वा पुनरेवान्तः प्रविशेत्~। ततोऽस्य वला\appbrk{}तै\appbrk{}ल\appbrk{}म\appbrk{}भ्य\appbrk{}ङ्गोऽ\appbrk{}व\appbrk{}चार्य यव\appbrk{}पि\appbrk{}ष्ट\appbrk{}मु\appbrk{}द्व\appbrk{}र्त्त\appbrk{}नार्थे~। सुखोष्णम्पयश्च परिषेकार्थे~। कषा\appbrk{}य\appbrk{}मा\appbrk{}ज\appbrk{}क\appbrk{}र्ण्ण\appbrk{}मु\appbrk{}त्सा\appbrk{}द\appbrk{}नार्थे~। आम\appbrk{}ल\appbrk{}क\appbrk{}मि\appbrk{}श्रा\appbrk{}श्चास्य यूषविकल्पाः~। क्षीर\appbrk{}म\appbrk{}धु\appbrk{}सि\appbrk{}द्धञ्च कृष्णतिल मभिहाराणार्थे~। एवं दशरात्रं द्वितीये परि\appbrk{}स\appbrk{}रे\appbrk{}व\appbrk{}र्त्तेत~। ततस्तृतीये परिसरे स्थिरी\appbrk{}कु\appbrk{}र्व्व\appbrk{}न्ना\appbrk{}त्मा\appbrk{}न\appbrk{}म\appbrk{}न्य\appbrk{}द्दे\appbrk{}श\appbrk{}रा\appbrk{}त्र\appbrk{}मा\appbrk{}सीत~। किञ्चि\appbrk{}द्दा\appbrk{}ता\appbrk{}त\appbrk{}प\appbrk{}प\appbrk{}व\appbrk{}न\appbrk{}मा\appbrk{}से\appbrk{}वेत~। पुनश्चान्तः प्रविशेत्~। न}\ H\textsuperscript{ac}; \skt{सायं वा भुक्त\appbrk{}वा\appbrk{}नु\appbrk{}प\appbrk{}श्रु\appbrk{}त\appbrk{}शा\appbrk{}न्तिः\appbrk{}कु\appbrk{}श\appbrk{}श\appbrk{}य्यायां कृष्णा\appbrk{}जि\appbrk{}नो\appbrk{}त्त\appbrk{}रायां सुहृ\appbrk{}द्भि\appbrk{}र\appbrk{}न्वा\appbrk{}स्य\appbrk{}मानः शयीत~। तृषितो वा शीतोदकं मात्राम्पिवेत्~। ततः प्रात\appbrk{}रु\appbrk{}त्था\appbrk{}यो\appbrk{}प\appbrk{}श्रु\appbrk{}त\appbrk{}शान्तिः कृतमङ्गलो गां स्पृष्ट्वा तथैवासीत्~। तस्य जीर्ण्णे सोमे च्छर्द्दिरुत्पद्यते~। ततः शोणि\appbrk{}ता\appbrk{}क्त\appbrk{}ङ्क्रि\appbrk{}मि\appbrk{}व्या\appbrk{}मि\appbrk{}श्र\appbrk{}ञ्छ\appbrk{}र्द्दि\appbrk{}त\appbrk{}वतः~। सायं स्रितशीतं क्षीर\appbrk{}म्वि\appbrk{}त\appbrk{}रे\appbrk{}त्त\appbrk{}त\appbrk{}स्तृ\appbrk{}ती\appbrk{}येऽ\appbrk{}हनि क्रिमि\appbrk{}व्या\appbrk{}मि\appbrk{}श्र\appbrk{}म\appbrk{}ति\appbrk{}सा\appbrk{}र्यते~। स तेना\appbrk{}ति\appbrk{}सृ\appbrk{}ष्ट\appbrk{}प्र\appbrk{}ति\appbrk{}ग्र\appbrk{}ह\appbrk{}भु\appbrk{}क्त\appbrk{}प्र\appbrk{}भृ\appbrk{}भि\appbrk{}र्वि\appbrk{}शैषैः प्रतिमुक्तः शुद्ध\appbrk{}त\appbrk{}नु\appbrk{}र्व्भ\appbrk{}वति~। ततः सायं स्नातस्य पूर्व्ववदेव क्षीरम्वितरेत्~। क्षौम\appbrk{}व\appbrk{}स्त्रा\appbrk{}स्तृ\appbrk{}ता\appbrk{}या\appbrk{}ञ्चैनं शय्यायां शाययीत~। तत\appbrk{}श्च\appbrk{}तु\appbrk{}र्थेऽ\appbrk{}हनि तस्य श्वय\appbrk{}थु\appbrk{}रु\appbrk{}प\appbrk{}जा\appbrk{}यते~। ततः सर्व्वाङ्गेभ्यः क्रिमयो निःक्रामन्ति~। तदहश्च पांशु शय्यायां पांशु\appbrk{}भि\appbrk{}र\appbrk{}व\appbrk{}की\appbrk{}र्य\appbrk{}माणः शयीत~। ततः साय\appbrk{}म्पू\appbrk{}व्व\appbrk{}व\appbrk{}देव क्षीर\appbrk{}म्वि\appbrk{}त\appbrk{}रे\appbrk{}देवं पञ्च\appbrk{}म\appbrk{}ष\appbrk{}ष्ठ\appbrk{}यो\appbrk{}र्द्दि\appbrk{}व\appbrk{}स\appbrk{}यो\appbrk{}र्व्व\appbrk{}र्तेत~। केव\appbrk{}ल\appbrk{}न्नू\appbrk{}भ\appbrk{}यतः कालमस्मै क्षीर\appbrk{}म्वि\appbrk{}त\appbrk{}रे\appbrk{}त्ततः~। सप्तमेहनि निर्म्मां\appbrk{}स\appbrk{}त्व\appbrk{}ग\appbrk{}स्ति\appbrk{}भूतः~। सोम\appbrk{}प्र\appbrk{}ति\appbrk{}ग्र\appbrk{}हा\appbrk{}दे\appbrk{}वो\appbrk{}च्छ\appbrk{}सिति~। तदहः क्षीहक्षीरेण सुखोष्णेन परिषिच्य तिल\appbrk{}म\appbrk{}धु\appbrk{}क\appbrk{}च\appbrk{}न्द\appbrk{}ना\appbrk{}नु\appbrk{}लि\appbrk{}प्त\appbrk{}देहं पयः पाययेत्~। ततोऽष्टमेहनि प्रातरेव क्षीर\appbrk{}षि\appbrk{}क्त\appbrk{}श्च\appbrk{}न्द\appbrk{}न\appbrk{}प\appbrk{}रि\appbrk{}दि\appbrk{}ग्ध\appbrk{}गात्रः पयः पीत्वा पांशु\appbrk{}श\appbrk{}य्या\appbrk{}मु\appbrk{}त्सृज्य क्षौमास्तृतायां शयीत~। तस्य मांसान्याप्याय्यन्ते त्वक्चावदलति~। दन्त\appbrk{}न\appbrk{}ख\appbrk{}रो\appbrk{}माणि चास्यशीर्यन्ते~। तस्य नवम दिव सात्प्र\appbrk{}भृ\appbrk{}त्य\appbrk{}नु\appbrk{}तै\appbrk{}ल\appbrk{}म\appbrk{}भ्यङ्गः~। सोम\appbrk{}व\appbrk{}ल्क\appbrk{}क\appbrk{}षायः परिषेकः~। ततोऽस्य दश\appbrk{}मे\appbrk{}ह\appbrk{}न्ये\appbrk{}त\appbrk{}देव वितरेत्ततोस्य त्वक्स्थि\appbrk{}र\appbrk{}ता\appbrk{}मु\appbrk{}पैति~। एव मेकाद शद्वा\appbrk{}द\appbrk{}श\appbrk{}यो\appbrk{}र्व्व\appbrk{}र्त्तेत~। तत\appbrk{}स्त्र\appbrk{}यो\appbrk{}द\appbrk{}शा\appbrk{}त्प्र\appbrk{}भृ\appbrk{}ति\appbrk{}नैव सोम\appbrk{}व\appbrk{}ल्क\appbrk{}क\appbrk{}षायो दातव्यः~। एव\appbrk{}मा\appbrk{}षो\appbrk{}ड\appbrk{}शा\appbrk{}द्व\appbrk{}र्त्तेत~। तत\appbrk{}स्स\appbrk{}प्त\appbrk{}द\appbrk{}शा\appbrk{}ष्टा\appbrk{}द\appbrk{}श\appbrk{}यो\appbrk{}र्द्द\appbrk{}श\appbrk{}यो\appbrk{}र्द्द\appbrk{}न्ताजा येत शिखरिणः स्निग्धा\appbrk{}व\appbrk{}ज्र\appbrk{}वै\appbrk{}दू\appbrk{}र्य\appbrk{}स्फ\appbrk{}टि\appbrk{}क\appbrk{}स\appbrk{}न्नि\appbrk{}काशाः स्थिराः समाः सहिष्णवः~। तदा प्रभृति चानवैः शालितण्डुलैः क्षीरयवागूमु पसेवेत~। याव\appbrk{}दा\appbrk{}प\appbrk{}ञ्च\appbrk{}विं\appbrk{}श\appbrk{}ति\appbrk{}रिति~। ततोऽस्मै दद्याच्छाल्योदनं मृदुमुभयतः कालं~। ततोस्य नखा जायन्ते विद्रु\appbrk{}मे\appbrk{}न्द्र\appbrk{}गो\appbrk{}प\appbrk{}क\appbrk{}त\appbrk{}रु\appbrk{}णा\appbrk{}दित्य प्रकाशाः स्निग्धाः स्थिराः~। लक्षणसम्पन्नाः केशाश्च शूक्ष्मा जायन्ते~। त्वक्च नीलो\appbrk{}त्प\appbrk{}ला\appbrk{}त\appbrk{}सी\appbrk{}पु\appbrk{}ष्प\appbrk{}वै\appbrk{}दू\appbrk{}र्य\appbrk{}प्र\appbrk{}काशाः~। उर्द्ध्वञ्च मासा\appbrk{}त्के\appbrk{}शा\appbrk{}न्वा\appbrk{}प\appbrk{}येत~। वापयित्वा चोशी\appbrk{}र\appbrk{}च\appbrk{}न्द\appbrk{}न\appbrk{}ति\appbrk{}ल\appbrk{}कल्कैः शिरः प्रदिदेत~। पयसा च स्नेहयीत~। ततोस्यानन्तरं सप्तरात्रात्केशा जायन्ते~। अञ्ज\appbrk{}न\appbrk{}वै\appbrk{}दू\appbrk{}र्य\appbrk{}भ्र\appbrk{}म\appbrk{}र\appbrk{}निभाः कुञ्चिताग्राः~। स्निग्धा\appbrk{}शि\appbrk{}र\appbrk{}सो\appbrk{}प\appbrk{}भो\appbrk{}गा\appbrk{}ना\appbrk{}मा\appbrk{}ह\appbrk{}र्त्तारैः~। तत\appbrk{}स्त्रि\appbrk{}रा\appbrk{}त्रा\appbrk{}त्तत्र शिरा प्रथ\appbrk{}म\appbrk{}प\appbrk{}रि\appbrk{}स\appbrk{}रा\appbrk{}न्नि\appbrk{}ष्क्रम्य मुहूर्त्तं स्थित्वा पुनरेवान्तः प्रविशेत्~। ततोऽस्य वला\appbrk{}तै\appbrk{}ल\appbrk{}म\appbrk{}भ्य\appbrk{}ङ्गोऽ\appbrk{}व\appbrk{}चार्य यव\appbrk{}पि\appbrk{}ष्ट\appbrk{}मु\appbrk{}द्व\appbrk{}र्त्त\appbrk{}नार्थे~। सुखोष्णम्पयश्च परिषेकार्थे~। कषा\appbrk{}य\appbrk{}मा\appbrk{}ज\appbrk{}क\appbrk{}र्ण्ण\appbrk{}मु\appbrk{}त्सा\appbrk{}द\appbrk{}नार्थे~। आम\appbrk{}ल\appbrk{}क\appbrk{}मि\appbrk{}श्रा\appbrk{}श्चास्य यूषविकल्पाः~। क्षीर\appbrk{}म\appbrk{}धु\appbrk{}सि\appbrk{}द्धञ्च कृष्णतिलतैल मभिहाराणार्थे~। एवं दशरात्रं द्वितीये परि\appbrk{}स\appbrk{}रे\appbrk{}व\appbrk{}र्त्तेत~। ततस्तृतीये परिसरे स्थिरी\appbrk{}कु\appbrk{}र्व्व\appbrk{}न्ना\appbrk{}त्मा\appbrk{}न\appbrk{}म\appbrk{}न्य\appbrk{}द्दे\appbrk{}श\appbrk{}रा\appbrk{}त्र\appbrk{}मा\appbrk{}सीत~। किञ्चि\appbrk{}द्दा\appbrk{}ता\appbrk{}त\appbrk{}प\appbrk{}प\appbrk{}व\appbrk{}न\appbrk{}मा\appbrk{}से\appbrk{}वेत~। पुनश्चान्तः प्रविशेत्~। न}\ H\textsuperscript{pc}}}\skt{ चात्मानमादर्शे }\edtext{\skt{स्यु वा}}{\lemma{\skt{स्यु वा}}\Afootnote{K H; \skt{ऽप्सु वा}\ A}}\skt{ निरीक्षेत~। रूप\-शा\-लि\-त्वा\-त्ततो }\edtext{\skt{न्यद्दशरात्रं}}{\lemma{\skt{न्यद्दशरात्रं}}\Afootnote{\skt{ऽन्यद्दशरात्रं}\ A; \skt{न्यद्दशरात्रं}\ K; \skt{ऽन्य\appbrk{}द\appbrk{}द्द\appbrk{}श\appbrk{}रात्रं}\ H}}\skt{ }\edtext{\skt{क्रोधा\-दी\-न्प\-रि\-ह\-रे\-देष}}{\lemma{\skt{क्रोधा\appbrk{}दी\appbrk{}न्प\appbrk{}रि\appbrk{}ह\appbrk{}रे\appbrk{}देष}}\Afootnote{K H; \skt{क्रोधा\appbrk{}दी\appbrk{}न्प\appbrk{}रि\appbrk{}ह\appbrk{}रेत्, एवं}\ A}}\skt{ }\edtext{\skt{सर्वे\-षा\-मु\-प\-योगः}}{\lemma{\skt{उपयोगः}}\Afootnote{K H; \skt{उपयोग}\ A}}\skt{~। }\edtext{\skt{विशेषतस्तु}}{\lemma{\skt{विशेषतस्तु}}\Afootnote{K H; \skt{विकल्पः~। विशेषतस्तु}\ A}}\skt{ }\edtext{\skt{वल्ली\-प्र\-ता\-न\-क्षु\-पा\-दयः}}{\lemma{\skt{वल्ली\appbrk{}प्र\appbrk{}ता\appbrk{}न\appbrk{}क्षु\appbrk{}पा\appbrk{}दयः}}\Afootnote{K H; \skt{वल्ली\appbrk{}प्र\appbrk{}ता\appbrk{}न\appbrk{}क्षु\appbrk{}प\appbrk{}का\appbrk{}दयः}\ A}}\skt{ सोमा }\edtext{\skt{भक्षयितव्याः}}{\lemma{\skt{भक्षयितव्याः}}\Afootnote{K H; \skt{ब्राह्म\appbrk{}ण\appbrk{}क्ष\appbrk{}त्रि\appbrk{}य\appbrk{}वै\appbrk{}श्यै\appbrk{}र्भ\appbrk{}क्ष\appbrk{}यि\appbrk{}तव्याः~।}\ A}}\skt{~। }\edtext{\skt{भवन्ति}}{\lemma{\skt{भवन्ति}}\Afootnote{K H; \skt{तेषां तु}\ A}}\skt{ }\edtext{\skt{तेषां\-म\-व\-च\-तुर्थो मुष्टयः}}{\lemma{\skt{तेषां\appbrk{}म\appbrk{}व\appbrk{}च\appbrk{}तुर्थो मुष्टयः}}\Afootnote{\skt{प्रमा\appbrk{}ण\appbrk{}म\appbrk{}र्ध\appbrk{}च\appbrk{}तु\appbrk{}ष्क\appbrk{}मु\appbrk{}ष्टयः~॥}\ A; \skt{तेषां\appbrk{}म\appbrk{}व\appbrk{}च\appbrk{}तुर्थो मुष्टयः~।}\ K; \skt{तेषा\appbrk{}म\appbrk{}र्थ\appbrk{}च\appbrk{}तु\appbrk{}र्था\appbrk{}मु\appbrk{}ष्टयः~।}\ H}}\skt{~।}\pend

\pstart \ledsidenote{\passref{4.29.13}}\skt{अङ्शुमन्तं }\edtext{\skt{सौवर्ण्णो}}{\lemma{\skt{सौवर्ण्णो}}\Afootnote{K; \skt{सौवर्णे}\ A; \skt{सौवर्ण्ण}\ H}}\skt{ }\edtext{\skt{पात्रे\-भि\-षु\-श\-या\-च्च\-न्द्र\-मसं}}{\lemma{\skt{पात्रे\appbrk{}भि\appbrk{}षु\appbrk{}श\appbrk{}या\appbrk{}च्च\appbrk{}न्द्र\appbrk{}मसं}}\Afootnote{K; \skt{पात्रेऽ\appbrk{}भि\appbrk{}षु\appbrk{}णु\appbrk{}यात्, चन्द्रमसं}\ A; \skt{पात्रेऽ\appbrk{}भि\appbrk{}षु\appbrk{}या\appbrk{}या\appbrk{}च्च\appbrk{}न्द्र\appbrk{}मसं}\ H}}\skt{ राजते~। }\edtext{\edtext{\skt{तावु\-प\-यो\-ज्या\-ष्ट\-गु\-ण\-मै\-श्व\-र्य\-म\-वा\-प्यै\-शानं}}{\lemma{\skt{तावु\appbrk{}प\appbrk{}यो\appbrk{}ज्या\appbrk{}ष्ट\appbrk{}गु\appbrk{}णम्}}\Afootnote{K; \skt{तावु\appbrk{}प\appbrk{}यु\appbrk{}ज्या\appbrk{}ष्ट\appbrk{}गु\appbrk{}णम्}\ A; \skt{तावु\appbrk{}प\appbrk{}यो\appbrk{}ज्या\appbrk{}ष्ट\appbrk{}गु\appbrk{}णम्}\ H}}}{\lemma{\skt{अवाप्यैशानं}}\Afootnote{K; \skt{अवाप्येशानं}\ A; \skt{अवाप्यैशानन्}\ H}}\skt{ देवं }\edtext{\skt{पश्यति}}{\lemma{\skt{पश्यति}}\Afootnote{K H; \skt{अनुप्रविशति,}\ A}}\skt{~। }\edtext{\skt{शेषांस्तु}}{\lemma{\skt{शेषांस्तु}}\Afootnote{A K; \skt{शेषास्तु}\ H}}\skt{ ताम्रमये }\edtext{\skt{मृन्मये}}{\lemma{\skt{मृन्मये}}\Afootnote{K H; \skt{मृन्मये वा}\ A}}\skt{ }\edtext{\skt{वा}}{\lemma{\skt{वा}}\Afootnote{K H; \skt{रोहिते वा}\ A}}\skt{ }\edtext{\skt{शूद्रवर्ज्यञ्च}}{\lemma{\skt{शूद्रवर्ज्यञ्च}}\Afootnote{K H; \skt{चर्मणि वितते; शूद्रवर्जं}\ A}}\skt{ }\edtext{\skt{त्रिभिर्वर्ण्णैः}}{\lemma{\skt{त्रिभिर्}}\Afootnote{A K; \skt{भभिर्}\ H}}\skt{ }\edtext{\skt{सोम}}{\lemma{\skt{सोम}}\Afootnote{K H; \skt{सोमा}\ A}}\skt{ }\edtext{\skt{उपयोक्तव्यः}}{\lemma{\skt{उपयोक्तव्यः}}\Afootnote{K H; \skt{उपयोक्तव्याः~।}\ A}}\skt{~। ततश्चतुर्थे }\edtext{\skt{मासि}}{\lemma{\skt{मासि}}\Afootnote{K H; \skt{मासे}\ A}}\skt{ पौर्ण्णमास्यां }\edtext{\skt{व्राह्म\-ण\-म\-र्च\-यित्वा}}{\lemma{\skt{व्राह्म\appbrk{}ण\appbrk{}म\appbrk{}र्च\appbrk{}यित्वा}}\Afootnote{K; \skt{शुचौ देशे ब्राह्म\appbrk{}णा\appbrk{}न\appbrk{}र्च\appbrk{}यित्वा}\ A; \skt{व्राह्म\appbrk{}णा\appbrk{}न\appbrk{}र्च्च\appbrk{}यित्वा}\ H}}\skt{ }\edtext{\skt{कृतमङ्गला}}{\lemma{\skt{कृतमङ्गला}}\Afootnote{K H; \skt{कृतमङ्गलो}\ A}}\skt{ }\edtext{\skt{निःक्रम्य}}{\lemma{\skt{निःक्रम्य}}\Afootnote{K H; \skt{निष्क्रम्य}\ A}}\skt{ }\edtext{\skt{यथेष्टं}}{\lemma{\skt{यथेष्टं}}\Afootnote{K H; \skt{यथोक्तं}\ A}}\skt{ }\edtext{\skt{व्रजेदिति}}{\lemma{\skt{व्रजेदिति}}\Afootnote{K H; \skt{व्रजेदिति~॥}\ A}}\skt{~॥}\pend

\pstart\versepar \ledsidenote{\passref{4.29.14}}\edtext{\skt{ओषधीनां}}{\lemma{\skt{ओषधीनां}}\Afootnote{A K; \skt{औषधीनाम्}\ H}}\skt{ पतिं }\edtext{\edtext{\skt{सोम\-मु\-प\-यु\-ज्यै\-व\-म\-क्षत}}{\lemma{\skt{सोमम्}}\Afootnote{A K; \skt{सोममम्}\ H}}}{\lemma{\skt{उप\appbrk{}यु\appbrk{}ज्यै\appbrk{}व\appbrk{}म\appbrk{}क्षत}}\Afootnote{\skt{उपयुज्य विचक्षणः~।}\ A; \skt{उप\appbrk{}यु\appbrk{}ज्यै\appbrk{}व\appbrk{}म\appbrk{}क्षत}\ K H}}\skt{~। }\pend
\pstart\versepar \skt{दश\-व\-र्ष\-स\-ह\-स्राणि नवान्धारयते तनुं~।}\pend

\pstart\versepar \ledsidenote{\passref{4.29.15}}\edtext{\skt{नाग्निर्न}}{\lemma{\skt{नाग्निर्न}}\Afootnote{A K; \skt{नाग्निर्न्न}\ H}}\skt{ }\edtext{\skt{तेजं न}}{\lemma{\skt{तेजं न}}\Afootnote{K H; \skt{तोयं न}\ A}}\skt{ विषं न शस्त्रं }\edtext{\skt{नाशस्त्रमेव}}{\lemma{\skt{नाशस्त्रम्}}\Afootnote{K H; \skt{नास्त्रम्}\ A}}\skt{ च~। }\pend
\pstart\versepar \skt{}\edtext{\skt{तस्यालमायुः}}{\lemma{\skt{तस्यालमायुः}}\Afootnote{K; \skt{तस्यालमायुः}\ A H}}\skt{ क्षपणे समर्थानि भवन्ति }\edtext{\skt{हि}}{\lemma{\skt{हि}}\Afootnote{K H; \emph{om.}\ A}}\skt{~।}\pend

\pstart\versepar \ledsidenote{\passref{4.29.16}}\skt{भद्राणां षष्टिवर्षाणां }\edtext{\skt{प्रसू\-ता\-ना\-म\-ने\-कधा}}{\lemma{\skt{प्रसू\appbrk{}ता\appbrk{}ना\appbrk{}म\appbrk{}ने\appbrk{}कधा}}\Afootnote{K; \skt{प्रस्रु\appbrk{}ता\appbrk{}ना\appbrk{}म\appbrk{}ने\appbrk{}कधा~।}\ A; \skt{प्रसू\appbrk{}ता\appbrk{}ना\appbrk{}म\appbrk{}ने\appbrk{}कधा~।}\ H}}\skt{~। }\pend
\pstart\versepar \skt{कुञ्जराणां सहस्रस्य }\edtext{\skt{बलं}}{\lemma{\skt{बलं}}\Afootnote{A; \skt{वलं}\ K H}}\skt{ }\edtext{\skt{समधिगच्छति}}{\lemma{\skt{समधिगच्छति}}\Afootnote{A K\textsuperscript{pc} H; \skt{सम\appbrk{}धि\appbrk{}ग\appbrk{}ग\appbrk{}च्छति~॥}\ K\textsuperscript{ac}}}\skt{~॥}\pend

\pstart\versepar \ledsidenote{\passref{4.29.17}}\skt{क्षीरोदं }\edtext{\skt{शक्रसदनं}}{\lemma{\skt{शक्रसदनं}}\Afootnote{A K; \skt{शक्रसदन}\ H}}\skt{ }\edtext{\skt{उत्तरांश्च}}{\lemma{\skt{उत्तरांश्च}}\Afootnote{A K; \skt{मुत्तराश्च}\ H}}\skt{ कुरूनपि }\pend
\pstart\versepar \skt{यत्रेच्छति स }\edtext{\skt{गन्तुं वै}}{\lemma{\skt{गन्तुं वै}}\Afootnote{K H; \skt{गन्तुं वा}\ A}}\skt{ }\edtext{\skt{तत्रा\-स्या\-प्र\-ति\-घा\-सतिः}}{\lemma{\skt{तत्रा\appbrk{}स्या\appbrk{}प्र\appbrk{}ति\appbrk{}घा\appbrk{}सतिः}}\Afootnote{K; \skt{तत्राप्रतिहता गतिः~॥}\ A; \skt{तत्रा\appbrk{}स्या\appbrk{}प्र\appbrk{}ति\appbrk{}घा\appbrk{}सति~॥}\ H}}\skt{~।}\pend

\pstart\versepar \ledsidenote{\passref{4.29.18}}\skt{कंदर्प इव रूपेण }\edtext{\skt{कन्त्या}}{\lemma{\skt{कन्त्या}}\Afootnote{K; \skt{कान्त्या}\ A H}}\skt{ चन्द्र इवापरः~। }\pend
\pstart\versepar \skt{प्रह्लादयति भूतानां मनांसि स महाद्युतिः~।}\pend

\pstart\versepar \ledsidenote{\passref{4.29.19}}\edtext{\skt{सांगोपांगांश्च}}{\lemma{\skt{सांगोपांगांश्च}}\Afootnote{A K; \skt{साङ्गोपाङ्गाश्च}\ H}}\skt{ }\edtext{\skt{निखिलां}}{\lemma{\skt{निखिलां}}\Afootnote{K H; \skt{निखिला}\ A}}\skt{ }\edtext{\skt{वेदां}}{\lemma{\skt{वेदां}}\Afootnote{K; \skt{न्वेदा}\ A; \skt{वेदा}\ H}}\skt{ }\edtext{\skt{विंदति}}{\lemma{\skt{विंदति}}\Afootnote{K; \skt{न्विन्दति}\ A H}}\skt{ तत्वतः~। }\pend
\pstart\versepar \skt{चर\-त्य\-मो\-घ\-स\-ङ्कल्पो }\edtext{\skt{देववच्चाखिलं}}{\lemma{\skt{देववच्चाखिलं}}\Afootnote{K\textsuperscript{pc}; \skt{देववच्चाखिलं}\ A H; \skt{देवच्चाखिलं}\ K\textsuperscript{ac}}}\skt{ जगत्~॥}\pend

\pstart\versepar \ledsidenote{\passref{4.29.20}}\edtext{\edtext{\skt{सर्वेषामेव}}{\lemma{\skt{सर्वेषाम्}}\Afootnote{A H; \skt{सर्वेषा}\ K}}}{\lemma{\skt{एव}}\Afootnote{A H; \skt{मेव}\ K}}\skt{ सोमानां पत्राणि दश पञ्च च~। }\pend
\pstart\versepar \skt{तानि शुक्ले च कृष्णे च }\edtext{\skt{जायन्ते च}}{\lemma{\skt{जायन्ते च}}\Afootnote{K H; \skt{जायन्ते नि}\ A}}\skt{ पतन्ति च~।}\pend

\pstart\versepar \ledsidenote{\passref{4.29.21}}\skt{एकैकं जायते पत्रं }\edtext{\edtext{\skt{सोम\-स्या\-ह\-र\-ह\-स्तथा}}{\lemma{\skt{सोम\appbrk{}स्या\appbrk{}ह\appbrk{}र\appbrk{}हस्}}\Afootnote{A K; \skt{सोमस्याहर}\ H\textsuperscript{ac}; \skt{सोमस्याहरह}\ H\textsuperscript{pc}}}}{\lemma{\skt{तथा}}\Afootnote{K; \skt{तदा~।}\ A; \skt{स्तथा~।}\ H}}\skt{~। }\pend
\pstart\versepar \skt{}\edtext{\skt{शुक्लस्य}}{\lemma{\skt{शुक्लस्य}}\Afootnote{A K; \skt{शुक्लस्या}\ H}}\skt{ }\edtext{\skt{पूर्ण्णमास्यान्तु}}{\lemma{\skt{पूर्ण्णमास्यान्तु}}\Afootnote{K H; \skt{पौर्णमास्यां तु}\ A}}\skt{ भवे\-त्प\-ञ्च\-द\-श\-च्छदः~।}\pend

\pstart\versepar \ledsidenote{\passref{4.29.22}}\edtext{\skt{दिवसे}}{\lemma{\skt{दिवसे}}\Afootnote{K H; \skt{शीर्यते पत्रमेकैकं दिवसे}\ A}}\skt{ दिवसे }\edtext{\skt{पत्रमेकैकं}}{\lemma{\skt{पत्रमेकैकं}}\Afootnote{K H; \emph{om.}\ A}}\skt{ }\edtext{\skt{शीर्यते}}{\lemma{\skt{शीर्यते}}\Afootnote{K H; \emph{om.}\ A}}\skt{ पुनः~। }\pend
\pstart\versepar \skt{}\edtext{\skt{कृष्णे}}{\lemma{\skt{कृष्णे}}\Afootnote{K H; \skt{कृष्ण}\ A}}\skt{ }\edtext{\skt{कृष्णद्वयं}}{\lemma{\skt{कृष्णद्वयं}}\Afootnote{K H; \skt{पक्षक्षये}\ A}}\skt{ चापि लता भवति केवला~॥}\pend

\pstart\versepar \ledsidenote{\passref{4.29.23}}\skt{अंशु\-मा\-ना\-ज्य\-ग\-न्धस्तु कन्द\-वा\-न्र\-ज\-त\-प्रभः~। }\pend
\pstart\versepar \skt{}\edtext{\skt{कद\-ल्या\-का\-र\-क\-न्दस्तु}}{\lemma{\skt{कद\appbrk{}ल्या\appbrk{}का\appbrk{}र\appbrk{}क\appbrk{}न्दस्तु}}\Afootnote{A K; \skt{कद\appbrk{}ल्या\appbrk{}का\appbrk{}र\appbrk{}क\appbrk{}न्दस्तु}\ H}}\skt{ }\edtext{\skt{मुञ्ज\-वा\-न्ल\-शु\-न\-च्छदः}}{\lemma{\skt{मुञ्ज\appbrk{}वा\appbrk{}न्ल\appbrk{}शु\appbrk{}न\appbrk{}च्छदः}}\Afootnote{K H; \skt{मुञ्जवां ल्लशुनच्छदः~॥}\ A}}\skt{~।}\pend

\pstart\versepar \ledsidenote{\passref{4.29.24}}\edtext{\skt{चन्द्रमा}}{\lemma{\skt{चन्द्रमा}}\Afootnote{K H; \skt{चन्द्रमाः}\ A}}\skt{ }\edtext{\skt{रजताभासो}}{\lemma{\skt{रजताभासो}}\Afootnote{K\textsuperscript{pc} H; \skt{कनकाभासो}\ A; \skt{रजभासो}\ K\textsuperscript{ac}}}\skt{ जले चरति सर्वदा~। }\pend
\pstart\versepar \skt{गरु\-डा\-हृ\-त\-नामा च }\edtext{\skt{श्येनाख्यौ}}{\lemma{\skt{श्येनाख्यौ}}\Afootnote{K H; \skt{श्वेताक्षश्}\ A}}\skt{ चापि पाण्डुरौ~।}\pend

\pstart\versepar \ledsidenote{\passref{4.29.25}}\skt{सर्प\-नि\-र्मो\-क\-स\-दृशौ तौ वृक्षा }\edtext{\skt{ग्रावलम्विनौ}}{\lemma{\skt{ग्रावलम्विनौ}}\Afootnote{K H; \skt{ग्रावलम्बिनौ~।}\ A}}\skt{~। }\pend
\pstart\versepar \skt{}\edtext{\skt{अथान्ये}}{\lemma{\skt{अथान्ये}}\Afootnote{K H; \skt{तथाऽन्ये}\ A}}\skt{ मण्ड\-लै\-श्चि\-त्रै\-श्चि\-त्रिता इव भान्ति ते~।}\pend

\pstart \ledsidenote{\passref{4.29.26}}\edtext{\skt{क्षीर\-क\-न्द\-ल\-ता\-वन्तः}}{\lemma{\skt{क्षीर\appbrk{}क\appbrk{}न्द\appbrk{}ल\appbrk{}ता\appbrk{}वन्तः}}\Afootnote{K H; \skt{सर्व एव तु विज्ञेयाः सोमाः पञ्चदशच्छदाः~। क्षीर\appbrk{}क\appbrk{}न्द\appbrk{}ल\appbrk{}ता\appbrk{}वन्तः}\ A}}\skt{ }\edtext{\skt{पत्रै\-र्ना\-ना\-वि\-धै\-र्वृता}}{\lemma{\skt{पत्रै\appbrk{}र्ना\appbrk{}ना\appbrk{}वि\appbrk{}धै\appbrk{}र्वृता}}\Afootnote{\skt{पत्रैर्नानाविधैः स्मृताः~॥}\ A; \skt{पत्रै\appbrk{}र्ना\appbrk{}ना\appbrk{}वि\appbrk{}धै\appbrk{}र्वृता~।}\ K; \skt{मन्त्रै\appbrk{}र्न्ना\appbrk{}ना\appbrk{}वि\appbrk{}धै\appbrk{}र्व्वृता~।}\ H}}\skt{~।}\pend

\pstart \ledsidenote{\passref{4.29.27}}\edtext{\skt{हिमवत्यर्वुदे}}{\lemma{\skt{अर्वुदे}}\Afootnote{K H; \skt{अर्बुदे}\ A}}\skt{ सह्ये महेन्द्रे मलये }\edtext{\skt{तथा}}{\lemma{\skt{तथा}}\Afootnote{K H; \skt{तथा~। श्रीपर्वते देवगिरौ गिरौ देवसहे तथा~॥}\ A}}\skt{~।}\pend

\pstart\versepar \ledsidenote{\passref{4.29.28}}\skt{पारियात्रे च विन्ध्ये च }\edtext{\skt{देवसूते}}{\lemma{\skt{देवसूते}}\Afootnote{K; \skt{देवसुन्दे}\ A; \skt{देवसूत्रे}\ H}}\skt{ }\edtext{\skt{हृदे}}{\lemma{\skt{हृदे}}\Afootnote{K H; \skt{ह्रदे}\ A}}\skt{ तथा~। }\pend
\pstart\versepar \skt{उत्तरेण वितस्तायाः }\edtext{\skt{प्रभवस्य}}{\lemma{\skt{प्रभवस्य}}\Afootnote{K; \skt{प्रवृद्धा ये}\ A; \skt{प्रभवत्य}\ H}}\skt{ महीधराः~।}\pend

\pstart \ledsidenote{\passref{4.29.28.add}}\edtext{\skt{पञ्च}}{\lemma{\skt{पञ्च}}\Afootnote{K H\textsuperscript{ac}; \skt{पञ्चपत्रे}\ H\textsuperscript{pc}}}\skt{ तेषामतो मध्ये सिन्धुर्नाम महाह्रदः~।}\pend

\pstart\versepar \ledsidenote{\passref{4.29.29}}\edtext{\skt{हठवत्प्लवते}}{\lemma{\skt{हठवत्प्लवते}}\Afootnote{K H; \skt{हष्ठवत्प्लवते}\ A}}\skt{ तत्र }\edtext{\skt{चन्द्रमा}}{\lemma{\skt{चन्द्रमा}}\Afootnote{K; \skt{चन्द्रमाः}\ A H}}\skt{ सोमसत्तमः~। }\pend
\pstart\versepar \skt{तस्योद्देशेषु }\edtext{\skt{चाप्यस्ति}}{\lemma{\skt{अस्ति}}\Afootnote{K H; \skt{अन्ति}\ A}}\skt{ मुञ्ज\-वा\-नं\-शु\-मा\-नपि~।}\pend

\pstart \ledsidenote{\passref{4.29.30}}\edtext{\skt{कश्मीरे च}}{\lemma{\skt{कश्मीरे च}}\Afootnote{K H; \skt{तस्योद्देशेषु चाप्यस्ति मुञ्ज\appbrk{}वा\appbrk{}नं\appbrk{}शु\appbrk{}मा\appbrk{}नपि~। काश्मीरेषु}\ A}}\skt{ सरो }\edtext{\skt{यन्तु}}{\lemma{\skt{यन्तु}}\Afootnote{K H; \skt{दिव्यं}\ A}}\skt{ नाम्ना क्षुद्रकमानसं~।}\pend

\pstart\versepar \ledsidenote{\passref{4.29.31}}\edtext{\edtext{\skt{गाय\-त्र्य\-श्त्रै\-ष्टु\-भ\-श्चापि}}{\lemma{\skt{गाय\appbrk{}त्र्य\appbrk{}श्त्रै\appbrk{}ष्टु\appbrk{}भश्}}\Afootnote{K; \skt{गायत्रस्त्रैष्टुभः}\ A; \skt{गाय\appbrk{}त्र्य\appbrk{}स्त्रै\appbrk{}ष्टु\appbrk{}भश्}\ H}}}{\lemma{\skt{चापि}}\Afootnote{K H; \skt{पाङ्क्तो}\ A}}\skt{ }\edtext{\skt{जाग्रतः}}{\lemma{\skt{जाग्रतः}}\Afootnote{K H; \skt{जागतः}\ A}}\skt{ }\edtext{\skt{सत्करस्तथा}}{\lemma{\skt{सत्करस्}}\Afootnote{K H; \skt{शाक्वरस्}\ A}}\skt{~। }\pend
\pstart\versepar \skt{}\edtext{\skt{तत्र}}{\lemma{\skt{तत्र}}\Afootnote{K H; \skt{अत्र}\ A}}\skt{ }\edtext{\skt{सन्त्यपरे}}{\lemma{\skt{सन्त्यपरे}}\Afootnote{K; \skt{सन्त्यपरे}\ A H}}\skt{ चापि सोमाः सोमसमप्रभाः~॥}\pend

\pstart\versepar \ledsidenote{\passref{4.29.32}}\edtext{\skt{न तां}}{\lemma{\skt{न तां}}\Afootnote{K H; \skt{यैश्चात्र मन्दभाग्यैस्ते भिष\appbrk{}ज\appbrk{}श्चा\appbrk{}प\appbrk{}मा\appbrk{}निताः~। न ता}\ A}}\skt{ }\edtext{\skt{पश्यन्त्यधर्मिष्ठाः}}{\lemma{\skt{पश्यन्त्य्}}\Afootnote{K H; \skt{न्पश्यन्त्य्}\ A}}\skt{ }\edtext{\skt{कृतघ्नाश्चापि ये}}{\lemma{\skt{चापि ये}}\Afootnote{K H; \skt{चापि मा}\ A}}\skt{ }\edtext{\skt{नराः}}{\lemma{\skt{नराः}}\Afootnote{K H; \skt{नवाः~।}\ A}}\skt{~। }\pend
\pstart\versepar \skt{भेष\-ज\-द्वे\-षि\-ण\-श्चापि }\edtext{\edtext{\skt{व्राह्म\-ण\-द्वे\-षि\-ण\-स्त\-थेति}}{\lemma{\skt{व्राह्म\appbrk{}ण\appbrk{}द्वे\appbrk{}षि\appbrk{}णस्}}\Afootnote{K H; \skt{ब्राह्म\appbrk{}ण\appbrk{}द्वे\appbrk{}षि\appbrk{}णस्}\ A}}}{\lemma{\skt{तथेति}}\Afootnote{K; \skt{तथा~॥}\ A H}}\skt{~॥}\pend

\endnumbering
\beginnumbering
\eledsection*{Adhyāya 30}
\pstart \ledsidenote{\passref{4.30.1}}\skt{अथातो निवृ\-त्त\-स\-न्ता\-पीयं रसायनं }\edtext{\skt{व्याख्यास्यामः}}{\lemma{\skt{व्याख्यास्यामः}}\Afootnote{A H; \emph{om.}\ K}}\skt{~॥}\pend

\pstart\versepar \ledsidenote{\passref{4.30.3}}\skt{यथा निवृत्तसन्तापा मोदन्तेदिवि देवताः~। }\pend
\pstart\versepar \skt{}\edtext{\skt{तथौषध्यस्त्विमाः}}{\lemma{\skt{तथौषध्यस्त्विमाः}}\Afootnote{K H; \skt{तथौषधीरिमाः}\ A}}\skt{ प्राप्य मोदन्ते भुवि मानवाः~॥}\pend

\pstart \ledsidenote{\passref{4.30.4}}\skt{अथ खलु }\edtext{\skt{सप्तविधाः}}{\lemma{\skt{सप्तविधाः}}\Afootnote{K H; \skt{सप्त}\ A}}\skt{ पुरुषा रसायनं नोपयुञ्जीरन्~। तद्यथा~। अना\-त्म\-वा\-न\-लसो दरिद्रः प्रमादी व्यसनी }\edtext{\skt{पापिष्ठो}}{\lemma{\skt{पापिष्ठो}}\Afootnote{K H; \skt{पापकृ}\ A}}\skt{ }\edtext{\skt{भेषजापमानी}}{\lemma{\skt{भेषजापमानी}}\Afootnote{K H; \skt{द्भेषजापमानी}\ A}}\skt{ चेति~॥ सप्तभिः }\edtext{\skt{कारणैर्न}}{\lemma{\skt{कारणैर्न}}\Afootnote{K; \skt{एव कारणैर्न}\ A; \skt{कारणैर्न्न}\ H}}\skt{ }\edtext{\skt{संपद्यन्ते}}{\lemma{\skt{संपद्यन्ते}}\Afootnote{K H; \skt{संपद्यते;}\ A}}\skt{~। तद्यथा~। }\edtext{\skt{अज्ञा\-ना\-द\-र\-म्भा\-द\-स्थि\-र\-चि\-त्त\-त्वा\-द्दा\-रि\-द्रा\-द\-ना\-य\-त\-ना\-दा\-या\-सा\-दौ\-ष\-धा\-ला\-भा\-च्चेति}}{\lemma{\skt{अज्ञा\appbrk{}ना\appbrk{}द\appbrk{}र\appbrk{}म्भा\appbrk{}द\appbrk{}स्थि\appbrk{}र\appbrk{}चि\appbrk{}त्त\appbrk{}त्वा\appbrk{}द्दा\appbrk{}रि\appbrk{}द्रा\appbrk{}द\appbrk{}ना\appbrk{}य\appbrk{}त\appbrk{}ना\appbrk{}दा\appbrk{}या\appbrk{}सा\appbrk{}दौ\appbrk{}ष\appbrk{}धा\appbrk{}ला\appbrk{}भा\appbrk{}च्चेति}}\Afootnote{K; \skt{अज्ञा\appbrk{}ना\appbrk{}द\appbrk{}ना\appbrk{}र\appbrk{}म्भा\appbrk{}द\appbrk{}स्थि\appbrk{}र\appbrk{}चि\appbrk{}त्त\appbrk{}त्वा\appbrk{}द्दा\appbrk{}रि\appbrk{}द्व्या\appbrk{}द\appbrk{}ना\appbrk{}य\appbrk{}त्त\appbrk{}त्वा\appbrk{}द\appbrk{}ध\appbrk{}र्मा\appbrk{}दौ\appbrk{}ष\appbrk{}धा\appbrk{}ला\appbrk{}भा\appbrk{}च्चेति॥}\ A; \skt{अज्ञा\appbrk{}ना\appbrk{}द\appbrk{}ना\appbrk{}रं\appbrk{}भा\appbrk{}द\appbrk{}स्थि\appbrk{}र\appbrk{}चि\appbrk{}त्त\appbrk{}त्वा\appbrk{}द्दा\appbrk{}रि\appbrk{}द्र्या\appbrk{}द\appbrk{}ना\appbrk{}य\appbrk{}त\appbrk{}ना\appbrk{}दा\appbrk{}या\appbrk{}सा\appbrk{}दौ\appbrk{}ष\appbrk{}धा\appbrk{}ला\appbrk{}भा\appbrk{}च्चेति~॥}\ H}}\skt{~॥}\pend

\pstart \ledsidenote{\passref{4.30.5}}\edtext{\skt{अथौ\-ष\-धी\-न्व\-क्ष्यामः}}{\lemma{\skt{अथौ\appbrk{}ष\appbrk{}धी\appbrk{}न्व\appbrk{}क्ष्यामः}}\Afootnote{K H; \skt{अथौ\appbrk{}ष\appbrk{}धी\appbrk{}व्या\appbrk{}ख्या\appbrk{}स्यामः}\ A}}\skt{~। तत्राजगरी }\edtext{\skt{स्वेतकापोती}}{\lemma{\skt{स्वेतकापोती}}\Afootnote{K H; \skt{श्वेतकापोती,}\ A}}\skt{ }\edtext{\skt{गोनसी}}{\lemma{\skt{गोनसी}}\Afootnote{K H; \emph{om.}\ A}}\skt{ }\edtext{\skt{कृष्णकापोती}}{\lemma{\skt{कृष्णकापोती}}\Afootnote{A K; \skt{कृष्णकपोती}\ H}}\skt{ }\edtext{\skt{वाराही}}{\lemma{\skt{वाराही}}\Afootnote{K H; \skt{गोनसी, वाराही,}\ A}}\skt{ }\edtext{\skt{च्छत्रातिच्छत्रा}}{\lemma{\skt{च्छत्रातिच्छत्रा}}\Afootnote{K\textsuperscript{pc} H; \skt{कन्या, छत्रा, अतिच्छत्रा,}\ A; \skt{च्छत्रा}\ K\textsuperscript{ac}}}\skt{ }\edtext{\skt{कन्या}}{\lemma{\skt{कन्या}}\Afootnote{K H; \emph{om.}\ A}}\skt{ }\edtext{\skt{करेणु}}{\lemma{\skt{करेणु}}\Afootnote{K H; \skt{करेणुः,}\ A}}\skt{ अजा }\edtext{\skt{तक्र\-का\-दि\-त्य\-प\-र्ण्णिनी}}{\lemma{\skt{तक्र\appbrk{}का\appbrk{}दि\appbrk{}त्य\appbrk{}प\appbrk{}र्ण्णिनी}}\Afootnote{\skt{चक्रका, आदित्यपर्णी,}\ A; \skt{तक्र\appbrk{}का\appbrk{}दि\appbrk{}त्य\appbrk{}प\appbrk{}र्ण्णिनी}\ K; \skt{तक्रक आदित्यपर्ण्णिनी}\ H}}\skt{ }\edtext{\skt{व्रह्मसुवर्चला}}{\lemma{\skt{व्रह्मसुवर्चला}}\Afootnote{K H; \skt{ब्रह्मसुवर्चला,}\ A}}\skt{ }\edtext{\skt{महास्रावणी}}{\lemma{\skt{महास्रावणी}}\Afootnote{K H; \skt{श्रावणी, महाश्रावणी,}\ A}}\skt{ }\edtext{\skt{गोलामी}}{\lemma{\skt{गोलामी}}\Afootnote{K H; \skt{गोलोमी, अजलोमी,}\ A}}\skt{ महावेगवती चेत्यष्टादश }\edtext{\skt{सोम\-स\-म\-वी\-र्या\-म\-हौ\-ष\-धयो}}{\lemma{\skt{सोम\appbrk{}स\appbrk{}म\appbrk{}वी\appbrk{}र्या\appbrk{}म\appbrk{}हौ\appbrk{}ष\appbrk{}धयो}}\Afootnote{A K; \skt{सोमवीर्या महौषधयो}\ H}}\skt{ }\edtext{\skt{व्याख्याः}}{\lemma{\skt{व्याख्याः}}\Afootnote{K H; \skt{व्याख्या}\ A}}\skt{~॥ }\edtext{\skt{तत्र}}{\lemma{\skt{तत्र}}\Afootnote{K H; \skt{ताः।}\ A}}\skt{ तासां }\edtext{\skt{सोम\-व\-त्क्रि\-या\-सी\-स्तु\-त\-तयः}}{\lemma{\skt{सोम\appbrk{}व\appbrk{}त्क्रि\appbrk{}या\appbrk{}सी\appbrk{}स्तु\appbrk{}त\appbrk{}तयः}}\Afootnote{K; \skt{सोम\appbrk{}व\appbrk{}त्क्रि\appbrk{}या\appbrk{}शीः\appbrk{}स्तु\appbrk{}तयः}\ A; \skt{सोम\appbrk{}व\appbrk{}त्क्रि\appbrk{}या\appbrk{}सी\appbrk{}स्तु\appbrk{}तयः}\ H}}\skt{ }\edtext{\skt{शास्त्रेभिहिताः}}{\lemma{\skt{शास्त्रेभिहिताः}}\Afootnote{K; \skt{शास्त्रेअभिहिताः}\ A H}}\skt{~। }\edtext{\skt{तासा\-म\-गा\-रे\-भि\-हि\-तानां याः}}{\lemma{\skt{तासा\appbrk{}म\appbrk{}गा\appbrk{}रे\appbrk{}भि\appbrk{}हि\appbrk{}तानां याः}}\Afootnote{K; \skt{तासा\appbrk{}मा\appbrk{}गा\appbrk{}रे+\appbrk{}अ\appbrk{}भि\appbrk{}हु\appbrk{}तानां याः}\ A; \skt{तासा\appbrk{}म\appbrk{}गा\appbrk{}रेऽ\appbrk{}भि\appbrk{}हि\appbrk{}तानां याः}\ H}}\skt{ }\edtext{\skt{क्षीरवत्यस्तासां}}{\lemma{\skt{क्षीरवत्यस्तासां}}\Afootnote{A K\textsuperscript{pc}; \skt{क्षीरत्यस्तासां}\ K\textsuperscript{ac}; \skt{क्षीरवत्यस्तासां}\ H}}\skt{ क्षीरकुडवं }\edtext{\skt{सकृ\-दे\-वो\-पु\-य\-ञ्जीत}}{\lemma{\skt{सकृ\appbrk{}दे\appbrk{}वो\appbrk{}पु\appbrk{}य\appbrk{}ञ्जीत}}\Afootnote{K; \skt{सकृ\appbrk{}दे\appbrk{}वो\appbrk{}प\appbrk{}यु\appbrk{}ञ्जीत}\ A H}}\skt{~। यास्त्व\-क्षी\-रा\-मू\-ल\-व\-त्य\-स्तासां प्रदेशिनी }\edtext{\skt{प्रमाणानी}}{\lemma{\skt{प्रमाणानी}}\Afootnote{K; \skt{प्रमाणानि}\ A; \skt{प्रमाणी}\ H}}\skt{ }\edtext{\skt{त्रीणि}}{\lemma{\skt{त्रीणि}}\Afootnote{A K; \skt{त्रिणि}\ H}}\skt{ काण्डानि प्रमा\-ण\-मु\-प\-योगे~। }\edtext{\skt{स्वेतकापोती}}{\lemma{\skt{स्वेतकापोती}}\Afootnote{K H; \skt{श्वेतकापोती}\ A}}\skt{ }\edtext{\skt{समूलवृत्ता}}{\lemma{\skt{समूलवृत्ता}}\Afootnote{K H; \skt{समूलपत्रा}\ A}}\skt{ भक्षयितव्या~। }\edtext{\skt{गोनस्याजगरी}}{\lemma{\skt{गोनस्याजगरी}}\Afootnote{K H; \skt{गोनस्यजगरी}\ A}}\skt{ कृष्णकापोतीनां }\edtext{\skt{सनखं}}{\lemma{\skt{सनखं}}\Afootnote{K H; \skt{सनख}\ A}}\skt{ मुष्टिं }\edtext{\skt{काण्डसः}}{\lemma{\skt{काण्डसः}}\Afootnote{K H; \skt{खण्डशः}\ A}}\skt{ कल्पयित्वा क्षीरेण }\edtext{\skt{विस्राव्य}}{\lemma{\skt{विस्राव्य}}\Afootnote{K H; \skt{विपाच्य}\ A}}\skt{ }\edtext{\skt{प्रति\-सा\-रि\-त\-म\-भि\-घा\-रि\-त\-म\-भि\-हु\-तञ्च}}{\lemma{\skt{प्रति\appbrk{}सा\appbrk{}रि\appbrk{}त\appbrk{}म\appbrk{}भि\appbrk{}घा\appbrk{}रि\appbrk{}त\appbrk{}म\appbrk{}भि\appbrk{}हु\appbrk{}तञ्च}}\Afootnote{K H; \skt{परि\appbrk{}स्रा\appbrk{}व्या\appbrk{}भि\appbrk{}घा\appbrk{}रि\appbrk{}त\appbrk{}म\appbrk{}भि\appbrk{}हुतं च}\ A}}\skt{ सकृ\-दे\-वो\-प\-यु\-ञ्जीत~। }\edtext{\skt{ततस्तासां}}{\lemma{\skt{ततस्तासां}}\Afootnote{K H; \skt{चक्र}\ A}}\skt{ }\edtext{\skt{यान्यापयसा}}{\lemma{\skt{यान्यापयसा}}\Afootnote{K H\textsuperscript{pc}; \skt{कायाः पयः}\ A}}\skt{ सकृदेव }\edtext{\skt{व्रह्मसुवर्चला}}{\lemma{\skt{व्रह्मसुवर्चला}}\Afootnote{K H\textsuperscript{pc}; \skt{ब्रह्मसुवर्चला}\ A}}\skt{ सप्त\-रा\-त्र\-मु\-प\-यो\-क्तव्या~। }\edtext{\edtext{\skt{भक्षकल्पेन}}{\lemma{\skt{भक्षकल्पेन}}\Afootnote{K H\textsuperscript{pc}; \skt{भक्ष्यकल्पेन,}\ A}}\skt{~। शेषाणां पञ्च पञ्च पलानि }\edtext{\skt{क्षीरा\-ढ\-क\-क\-थि\-तानि}}{\lemma{\skt{क्षीरा\appbrk{}ढ\appbrk{}क\appbrk{}क\appbrk{}थि\appbrk{}तानि}}\Afootnote{K H\textsuperscript{pc}; \skt{क्षीरा\appbrk{}ढ\appbrk{}क\appbrk{}क्व\appbrk{}थि\appbrk{}तानि}\ A}}\skt{ }\edtext{\skt{प्रस्थेवशिष्टे परिस्राव्य}}{\lemma{\skt{प्रस्थेवशिष्टे परिस्राव्य}}\Afootnote{K; \skt{प्रस्थे+\appbrk{}अ\appbrk{}व\appbrk{}शि\appbrk{}ष्टे+\appbrk{}अ\appbrk{}व\appbrk{}तार्य परिस्राव्य}\ A; \skt{प्रस्थेवसिषे परिस्राव्य}\ H\textsuperscript{pc}}}\skt{ सकृ\-दे\-वो\-प\-यु\-ञ्जीत~। }\edtext{\skt{सोमेन}}{\lemma{\skt{सोमेन}}\Afootnote{K H; \skt{सोमव}\ A}}\skt{ }\edtext{\skt{चाहारविहाराः}}{\lemma{\skt{चाहारविहाराः}}\Afootnote{K; \skt{दाहारविहारौ}\ A; \skt{चारविहारा}\ H}}\skt{ }\edtext{\skt{व्याख्याताः}}{\lemma{\skt{व्याख्याताः}}\Afootnote{K; \skt{व्याख्यातौ,}\ A; \skt{व्याख्याता~॥}\ H}}\skt{~। }\edtext{\skt{केवलं तु}}{\lemma{\skt{केवलं तु}}\Afootnote{K H; \skt{केवलं}\ A}}\skt{ नवनीतमभ्यं गार्थे~। शेषं सोम\-व\-दा\-नि\-र्ग\-मा\-दिति}}{\lemma{\skt{यान्यापयसा… सो}}\Afootnote{\emph{om.}\ H\textsuperscript{ac}}}\skt{~॥}\pend

\pstart\versepar \ledsidenote{\passref{4.30.6}}\edtext{\skt{भ}}{\lemma{\skt{भ}}\Afootnote{K H; \skt{भवन्ति चात्र}\ A}}\skt{~॥ युवानं }\edtext{\skt{सिहंविक्रान्तं}}{\lemma{\skt{सिहंविक्रान्तं}}\Afootnote{K; \skt{सिंहविक्रान्तं}\ A H}}\skt{ }\edtext{\skt{सकृ\-च्छ्रु\-भ\-नि\-गा\-दितं}}{\lemma{\skt{सकृ\appbrk{}च्छ्रु\appbrk{}भ\appbrk{}नि\appbrk{}गा\appbrk{}दितं}}\Afootnote{K H; \skt{कान्तं श्रुत\appbrk{}नि\appbrk{}गा\appbrk{}दि\appbrk{}नम्।}\ A}}\skt{~। }\pend
\pstart\versepar \skt{कुर्युरेताः }\edtext{\skt{क्रमेणैवं}}{\lemma{\skt{क्रमेणैवं}}\Afootnote{K; \skt{क्रमेणैव}\ A H}}\skt{ द्विसहस्रायुषं नरं~॥}\pend

\pstart\versepar \ledsidenote{\passref{4.30.7}}\skt{अङ्गदी कुण्डली मौली }\edtext{\skt{दिव्य\-स्र\-क्च\-न्द\-ना\-म्वरः}}{\lemma{\skt{दिव्य\appbrk{}स्र\appbrk{}क्च\appbrk{}न्द\appbrk{}ना\appbrk{}म्वरः}}\Afootnote{K\textsuperscript{pc} H; \skt{दिव्य\appbrk{}स्र\appbrk{}क्च\appbrk{}न्द\appbrk{}ना\appbrk{}म्बरः।}\ A; \skt{दिव्य\appbrk{}स्र\appbrk{}क्च\appbrk{}न्द\appbrk{}ना\appbrk{}म्व\appbrk{}रःरः}\ K\textsuperscript{ac}}}\skt{ }\pend
\pstart\versepar \skt{चर\-त्य\-मो\-घ\-स\-ङ्कल्पो }\edtext{\skt{नभ\-स्यं\-र्वु\-द\-नि\-र्गमे}}{\lemma{\skt{नभ\appbrk{}स्यं\appbrk{}र्वु\appbrk{}द\appbrk{}नि\appbrk{}र्गमे}}\Afootnote{K; \skt{नभ\appbrk{}स्य\appbrk{}म्बु\appbrk{}द\appbrk{}दु\appbrk{}र्गमे॥}\ A; \skt{नभ\appbrk{}स्य\appbrk{}म्वु\appbrk{}द\appbrk{}नि\appbrk{}र्ग्गमे~॥}\ H}}\skt{~।}\pend

\pstart\versepar \ledsidenote{\passref{4.30.8}}\skt{व्रजन्ति पक्षिणो येन }\edtext{\skt{जललम्वाश्च}}{\lemma{\skt{जललम्वाश्च}}\Afootnote{K H; \skt{जललम्बाश्च}\ A}}\skt{ तोयदाः~। }\pend
\pstart\versepar \skt{}\edtext{\skt{गति\-स्त्वौ\-ष\-ध\-सि\-द्धस्य}}{\lemma{\skt{गति\appbrk{}स्त्वौ\appbrk{}ष\appbrk{}ध\appbrk{}सि\appbrk{}द्धस्य}}\Afootnote{K H; \skt{गतिः सौषधिसिद्धस्य}\ A}}\skt{ }\edtext{\skt{सोमसिद्धिगतिः}}{\lemma{\skt{सोमसिद्धिगतिः}}\Afootnote{K H; \skt{सोमसिद्धे गतिः}\ A}}\skt{ }\edtext{\skt{पराः~॥ >}}{\lemma{\skt{पराः~॥ >}}\Afootnote{\skt{परा॥}\ A; \skt{पराः~॥}\ K H}}\pend

\pstart\versepar \ledsidenote{\passref{4.30.9}}\skt{अथ वक्ष्यामि }\edtext{\skt{विज्ञा\-न\-मौ\-ष\-धीनां}}{\lemma{\skt{विज्ञा\appbrk{}न\appbrk{}मौ\appbrk{}ष\appbrk{}धीनां}}\Afootnote{A K; \skt{विज्ञा\appbrk{}न\appbrk{}मो\appbrk{}ष\appbrk{}धी\appbrk{}नाम्}\ H}}\skt{ पृथक्पृथक्~। }\pend
\pstart\versepar \skt{मण्डलैः }\edtext{\skt{कपिलैश्चित्रा}}{\lemma{\skt{कपिलैश्चित्रा}}\Afootnote{K H; \skt{कपिलैश्चित्रैः}\ A}}\skt{ सर्पाभा पञ्चपर्ण्णिनी~।}\pend

\pstart\versepar \ledsidenote{\passref{4.30.10}}\edtext{\skt{पञ्चा\-र\-न्नि\-प्र\-माणा च}}{\lemma{\skt{पञ्चा\appbrk{}र\appbrk{}न्नि\appbrk{}प्र\appbrk{}माणा च}}\Afootnote{K H; \skt{पञ्चा\appbrk{}र\appbrk{}त्नि\appbrk{}प्र\appbrk{}माणा च}\ A}}\skt{ }\edtext{\skt{विज्ञे\-या\-ज\-ग\-री\-र्वुधैः}}{\lemma{\skt{विज्ञे\appbrk{}या\appbrk{}ज\appbrk{}ग\appbrk{}री\appbrk{}र्वुधैः}}\Afootnote{K H; \skt{विज्ञे\appbrk{}या+\appbrk{}अ\appbrk{}ज\appbrk{}गरी बुधैः।}\ A}}\skt{~॥ }\pend
\pstart\versepar \skt{}\edtext{\skt{निःपत्राः}}{\lemma{\skt{निःपत्राः}}\Afootnote{K H; \skt{निष्पत्रा}\ A}}\skt{ }\edtext{\skt{कनकाभाशा}}{\lemma{\skt{कनकाभाशा}}\Afootnote{K; \skt{कनकाभासा}\ A H}}\skt{ }\edtext{\skt{मूलि\-न्य\-ङ्गु\-ष्ठ\-स\-म्मिता}}{\lemma{\skt{मूलि\appbrk{}न्य\appbrk{}ङ्गु\appbrk{}ष्ठ\appbrk{}स\appbrk{}म्मिता}}\Afootnote{K H; \skt{मूले द्व्यङ्गुलसंमिता॥}\ A}}\skt{~।}\pend

\pstart\versepar \ledsidenote{\passref{4.30.11}}\skt{सर्पाकारा }\edtext{\skt{लोहिताक्षी}}{\lemma{\skt{लोहिताक्षी}}\Afootnote{K H; \skt{लोहितान्ता}\ A}}\skt{ }\edtext{\skt{स्वेत\-का\-पो\-ति\-रु\-च्यते}}{\lemma{\skt{स्वेत\appbrk{}का\appbrk{}पो\appbrk{}ति\appbrk{}रु\appbrk{}च्यते}}\Afootnote{\skt{श्वेत\appbrk{}का\appbrk{}पो\appbrk{}ति\appbrk{}रु\appbrk{}च्यते।}\ A; \skt{स्वेत\appbrk{}का\appbrk{}पो\appbrk{}ति\appbrk{}रु\appbrk{}च्यते}\ K H}}\skt{~॥ }\pend
\pstart\versepar \skt{}\edtext{\skt{द्विपर्ण्णिनी}}{\lemma{\skt{द्विपर्ण्णिनी}}\Afootnote{K H; \skt{द्विपर्णिनीं}\ A}}\skt{ मूल\-भ\-वा\-म\-रुणां }\edtext{\skt{कृष्णमण्डला}}{\lemma{\skt{कृष्णमण्डला}}\Afootnote{K H; \skt{कृष्णमण्डलाम्॥}\ A}}\skt{~।}\pend

\pstart\versepar \ledsidenote{\passref{4.30.12}}\edtext{\skt{द्विर\-न्नि\-प्र\-मा\-णा\-ञ्जा\-नी\-या\-द्गो\-नसीं}}{\lemma{\skt{द्विर\appbrk{}न्नि\appbrk{}प्र\appbrk{}मा\appbrk{}णा\appbrk{}ञ्जा\appbrk{}नी\appbrk{}या\appbrk{}द्गो\appbrk{}नसीं}}\Afootnote{K; \skt{द्व्यरत्निमात्रां जानीयाद्गोनसीं}\ A; \skt{द्विर\appbrk{}न्नि\appbrk{}प्र\appbrk{}मा\appbrk{}णा\appbrk{}ञ्जा\appbrk{}नी\appbrk{}या\appbrk{}ङ्गो\appbrk{}न\appbrk{}सीङ्}\ H}}\skt{ गोनसाकृतिं~॥ }\pend
\pstart\versepar \skt{सक्षीरां }\edtext{\skt{रोमसां}}{\lemma{\skt{रोमसां}}\Afootnote{K H; \skt{रोमशां}\ A}}\skt{ }\edtext{\skt{मृद्वी}}{\lemma{\skt{मृद्वी}}\Afootnote{K; \skt{मृद्वीं}\ A H}}\skt{ }\edtext{\skt{रसे\-चे\-क्षु\-र\-सो\-पमां}}{\lemma{\skt{रसे\appbrk{}चे\appbrk{}क्षु\appbrk{}र\appbrk{}सो\appbrk{}पमां}}\Afootnote{K H; \skt{रसे\appbrk{}ने\appbrk{}क्षु\appbrk{}र\appbrk{}सो\appbrk{}प\appbrk{}माम्॥}\ A}}\skt{~।}\pend

\pstart\versepar \ledsidenote{\passref{4.30.13}}\skt{एवं }\edtext{\skt{रूपरसां}}{\lemma{\skt{रूपरसां}}\Afootnote{A K; \skt{रूपरसं}\ H}}\skt{ चापि कृष्ण\-का\-पो\-ति\-मा\-दि\-शेत्~॥ }\pend
\pstart\versepar \skt{कृष्ण\-स\-र्प\-स्य\-रू\-पेण वाराही कन्दसम्भवा~।}\edtext{}{\lemma{\skt{कृष्ण\appbrk{}का\appbrk{}पो\appbrk{}ति\appbrk{}मा\appbrk{}दि\appbrk{}शेत्~॥ कृष्ण\appbrk{}स\appbrk{}र्प\appbrk{}स्य\appbrk{}रू\appbrk{}पेण}}\Afootnote{\skt{कृष्ण\appbrk{}का\appbrk{}पो\appbrk{}ति\appbrk{}मा\appbrk{}दि\appbrk{}शेत्। कृष्ण\appbrk{}स\appbrk{}र्प\appbrk{}स्व\appbrk{}रू\appbrk{}पेण}\ A; \skt{कृष्ण\appbrk{}का\appbrk{}पो\appbrk{}ति\appbrk{}मा\appbrk{}दि\appbrk{}शे\appbrk{}त्कृ\appbrk{}ष्ण\appbrk{}स\appbrk{}र्प\appbrk{}स्य\appbrk{}रू\appbrk{}पेण}\ K H}}\pend

\pstart\versepar \ledsidenote{\passref{4.30.14}}\edtext{\skt{एकपत्र्या}}{\lemma{\skt{एकपत्र्या}}\Afootnote{K H; \skt{एकपत्रा}\ A}}\skt{ महावीर्या भिन्ना\-ञ्ज\-न\-स\-म\-प्रभा~। }\pend
\pstart\versepar \skt{}\edtext{\skt{च्छत्रातिच्छत्रिके}}{\lemma{\skt{च्छत्रातिच्छत्रिके}}\Afootnote{K H; \skt{छत्रातिच्छत्रके}\ A}}\skt{ }\edtext{\skt{विद्याद्रक्षोघ्ना}}{\lemma{\skt{विद्याद्रक्षोघ्ना}}\Afootnote{K H; \skt{विद्याद्रक्षोघ्ने}\ A}}\skt{ }\edtext{\skt{कन्दसम्भवे}}{\lemma{\skt{कन्दसम्भवे}}\Afootnote{A K; \skt{कन्दसम्भवा~।}\ H}}\skt{~॥}\pend

\pstart\versepar \ledsidenote{\passref{4.30.15}}\skt{जरा\-मृ\-त्यु\-नि\-वा\-रिण्यौ }\edtext{\skt{स्वेत\-का\-पो\-ति\-सं\-स्थिते}}{\lemma{\skt{स्वेत\appbrk{}का\appbrk{}पो\appbrk{}ति\appbrk{}सं\appbrk{}स्थिते}}\Afootnote{K H; \skt{श्वेत\appbrk{}का\appbrk{}पो\appbrk{}ति\appbrk{}सं\appbrk{}स्थिते।}\ A}}\skt{~। }\pend
\pstart\versepar \skt{}\edtext{\edtext{\skt{कान्तैर्द्वादशभिः}}{\lemma{\skt{कान्तैर्}}\Afootnote{K H; \skt{कान्तै}\ A}}}{\lemma{\skt{द्वादशभिः}}\Afootnote{K H; \skt{र्द्वादशभिः}\ A}}\skt{ }\edtext{\skt{पत्रै}}{\lemma{\skt{पत्रै}}\Afootnote{A K; \skt{पत्त्रैर्}\ H}}\skt{ }\edtext{\skt{मयू\-र\-व\-र\-हो\-पमैः}}{\lemma{\skt{मयू\appbrk{}र\appbrk{}व\appbrk{}र\appbrk{}हो\appbrk{}पमैः}}\Afootnote{K H; \skt{र्मयू\appbrk{}रा\appbrk{}ङ्ग\appbrk{}रु\appbrk{}हो\appbrk{}पमैः॥}\ A}}\skt{~।}\pend

\pstart\versepar \ledsidenote{\passref{4.30.16}}\skt{कन्दजा काञ्चनक्षीरी कन्या नाम }\edtext{\skt{महौषधी}}{\lemma{\skt{महौषधी}}\Afootnote{A K; \skt{महौषधीः~।}\ H}}\skt{~। }\pend
\pstart\versepar \skt{करेणुः }\edtext{\skt{सुवहुक्षीरा}}{\lemma{\skt{सुवहुक्षीरा}}\Afootnote{K H; \skt{सुबहुक्षीरा}\ A}}\skt{ कन्देन }\edtext{\skt{गजरूपिणां}}{\lemma{\skt{गजरूपिणां}}\Afootnote{K H; \skt{गजरूपिणी॥}\ A}}\skt{~।}\pend

\pstart\versepar \ledsidenote{\passref{4.30.17}}\edtext{\skt{हस्ति\-क\-र्ण्ण\-प\-ला\-श\-स्या\-त्तू\-ल्य\-प\-त्रा\-द्वि\-प\-र्ण्णिनी}}{\lemma{\skt{हस्ति\appbrk{}क\appbrk{}र्ण्ण\appbrk{}प\appbrk{}ला\appbrk{}श\appbrk{}स्या\appbrk{}त्तू\appbrk{}ल्य\appbrk{}प\appbrk{}त्रा\appbrk{}द्वि\appbrk{}प\appbrk{}र्ण्णिनी}}\Afootnote{K; \skt{हस्ति\appbrk{}क\appbrk{}र्ण\appbrk{}प\appbrk{}ला\appbrk{}शस्य तुल्यपर्णा द्विपर्णिनी।}\ A; \skt{हस्ति\appbrk{}क\appbrk{}र्ण्ण\appbrk{}प\appbrk{}लासः स्यात्तू\appbrk{}ल्य\appbrk{}प\appbrk{}त्रा\appbrk{}द्वि\appbrk{}प\appbrk{}र्ण्णिनी~।}\ H}}\skt{~। }\pend
\pstart\versepar \skt{}\edtext{\skt{अजा\-न\-ना\-भ\-कन्दा तु}}{\lemma{\skt{अजा\appbrk{}न\appbrk{}ना\appbrk{}भ\appbrk{}कन्दा तु}}\Afootnote{K H; \skt{अजा\appbrk{}स्त\appbrk{}ना\appbrk{}भ\appbrk{}कन्दा तु}\ A}}\skt{ }\edtext{\skt{सक्षी\-रा\-क्षु\-प\-रू\-पिणी}}{\lemma{\skt{सक्षी\appbrk{}रा\appbrk{}क्षु\appbrk{}प\appbrk{}रू\appbrk{}पिणी}}\Afootnote{K H; \skt{सक्षीरा श्रुपरूपिणी॥}\ A}}\skt{~।}\pend

\pstart\versepar \ledsidenote{\passref{4.30.18}}\skt{अजा महौषधी ज्ञेया शङ्खकुन्देन्दु पाण्डुरा~। }\pend
\pstart\versepar \skt{}\edtext{\skt{स्वेता}}{\lemma{\skt{स्वेता}}\Afootnote{K H; \skt{श्वेतां}\ A}}\skt{ }\edtext{\skt{विचित्रकुसुमा}}{\lemma{\skt{विचित्रकुसुमा}}\Afootnote{K; \skt{विचित्रकुसुमां}\ A; \skt{विचित्रा कुसुमा}\ H}}\skt{ }\edtext{\skt{काकादनि}}{\lemma{\skt{काकादनि}}\Afootnote{K H; \skt{काकादन्या}\ A}}\skt{ }\edtext{\skt{समच्छदाः}}{\lemma{\skt{समच्छदाः}}\Afootnote{K H; \skt{समां क्षुपाम्॥}\ A}}\pend

\pstart\versepar \ledsidenote{\passref{4.30.19}}\edtext{\skt{तक्रकामौषधीं}}{\lemma{\skt{तक्रकामौषधीं}}\Afootnote{K H; \skt{चक्रकामोषधीं}\ A}}\skt{ }\edtext{\skt{विद्या\-ज्ज\-रा\-मृ\-त्यु\-वि\-ना\-शनी}}{\lemma{\skt{विद्या\appbrk{}ज्ज\appbrk{}रा\appbrk{}मृ\appbrk{}त्यु\appbrk{}वि\appbrk{}ना\appbrk{}शनी}}\Afootnote{K; \skt{विद्या\appbrk{}ज्ज\appbrk{}रा\appbrk{}मृ\appbrk{}त्यु\appbrk{}नि\appbrk{}वा\appbrk{}रि\appbrk{}णीम्।}\ A; \skt{विद्या\appbrk{}ज्ज\appbrk{}रा\appbrk{}मृ\appbrk{}त्यु\appbrk{}वि\appbrk{}ना\appbrk{}शनीं~।}\ H}}\skt{~। }\pend
\pstart\versepar \skt{}\edtext{\skt{मूलिनीं}}{\lemma{\skt{मूलिनीं}}\Afootnote{K H; \skt{मूलिनी}\ A}}\skt{ पञ्चभिः पत्रैः सुर\-क्तां\-शु\-क\-को\-मलैः~।}\pend

\pstart\versepar \ledsidenote{\passref{4.30.20}}\edtext{\edtext{\skt{आदित्यपर्ण्णिनी}}{\lemma{\skt{आदित्य्}}\Afootnote{A; \skt{आदि}\ K H}}}{\lemma{\skt{अपर्ण्णिनी}}\Afootnote{A; \skt{त्यपर्ण्णिनी}\ K H}}\skt{ ज्ञेया }\edtext{\edtext{\skt{सदा\-दि\-त्या\-नु\-व\-र्त्तिनी}}{\lemma{\skt{सदादित्य्}}\Afootnote{H; \skt{सदा++आदित्य्}\ A; \skt{सदादि}\ K}}}{\lemma{\skt{आनुवर्त्तिनी}}\Afootnote{A H; \skt{त्यानुवर्त्तिनी~।}\ K}}\skt{~। }\pend
\pstart\versepar \skt{कनकाभा जलान्तेषु सर्वतः परिसर्पति~।}\pend

\pstart\versepar \ledsidenote{\passref{4.30.21}}\skt{सक्षीरा पद्मिनीप्रख्या }\edtext{\skt{देवी\-व्र\-ह्म\-सु\-व\-र्चला}}{\lemma{\skt{देवी\appbrk{}व्र\appbrk{}ह्म\appbrk{}सु\appbrk{}व\appbrk{}र्चला}}\Afootnote{K H; \skt{देवी ब्रह्मसुवर्चला।}\ A}}\skt{~। }\pend
\pstart\versepar \skt{}\edtext{\edtext{\skt{रन्नि\-प्र\-मा\-णा\-क्षु\-प\-व\-त्प\-त्रै\-र्द्व्यं\-गु\-ल\-स\-म्मितैः}}{\lemma{\skt{रन्नि\appbrk{}प्र\appbrk{}मा\appbrk{}णा\appbrk{}क्षु\appbrk{}प\appbrk{}व\appbrk{}त्प\appbrk{}त्रैर्}}\Afootnote{K; \skt{अर\appbrk{}त्नि\appbrk{}मा\appbrk{}त्र\appbrk{}क्षु\appbrk{}पका पत्रै}\ A; \skt{रत्नि\appbrk{}प्र\appbrk{}मा\appbrk{}णा\appbrk{}क्षु\appbrk{}प\appbrk{}व\appbrk{}त्प\appbrk{}त्त्रैर्}\ H}}}{\lemma{\skt{द्व्यंगुलसम्मितैः}}\Afootnote{K H; \skt{र्द्व्यङ्गुलसंमितैः॥}\ A}}\skt{~।}\pend

\pstart\versepar \ledsidenote{\passref{4.30.22}}\edtext{\skt{पुष्पै}}{\lemma{\skt{पुष्पै}}\Afootnote{A K; \skt{पुष्पैर्}\ H}}\skt{ }\edtext{\skt{नीलोत्पलाकारैः}}{\lemma{\skt{नीलोत्पलाकारैः}}\Afootnote{K H; \skt{र्नीलोत्पलाकारैः}\ A}}\skt{ }\edtext{\skt{फलै\-श्चा\-ञ्ज\-न\-स\-प्रभै}}{\lemma{\skt{चाञ्जनसप्रभै}}\Afootnote{K; \skt{चाञ्जनसन्निभैः।}\ A; \skt{चाञ्जनसप्रभैः~॥}\ H}}\skt{~॥ }\pend
\pstart\versepar \skt{}\edtext{\skt{स्रावणी}}{\lemma{\skt{स्रावणी}}\Afootnote{K H; \skt{श्रावणी}\ A}}\skt{ महती ज्ञेया कनकाभा पयस्विनी~।}\pend

\pstart\versepar \ledsidenote{\passref{4.30.23}}\edtext{\skt{स्रावणीपाण्डुरा तु}}{\lemma{\skt{स्रावणीपाण्डुरा तु}}\Afootnote{K H; \skt{श्रावणी पाण्डुराभा}\ A}}\skt{ }\edtext{\skt{स्यात्म\-हा\-स्रा\-व\-णि\-ल\-क्षणा}}{\lemma{\skt{स्यात्म\appbrk{}हा\appbrk{}स्रा\appbrk{}व\appbrk{}णि\appbrk{}ल\appbrk{}क्षणा}}\Afootnote{K; \skt{सा महा\appbrk{}श्रा\appbrk{}व\appbrk{}णि\appbrk{}ल\appbrk{}क्षणा।}\ A; \skt{स्यात्म\appbrk{}हा\appbrk{}स्रा\appbrk{}व\appbrk{}णी\appbrk{}ल\appbrk{}क्षणा~॥}\ H}}\skt{~। }\pend
\pstart\versepar \skt{गोलोमी चाजलोमी च रोमशे कन्दसम्भवे~।}\pend

\pstart\versepar \ledsidenote{\passref{4.30.24}}\edtext{\skt{हंसपादी च}}{\lemma{\skt{हंसपादी च}}\Afootnote{K H; \skt{हंसपादीव}\ A}}\skt{ विच्छिन्नैः }\edtext{\edtext{\skt{पत्रै\-र्मू\-ल\-स\-मु\-द्भवैः}}{\lemma{\skt{पत्रैर्}}\Afootnote{K H; \skt{पत्रै}\ A}}}{\lemma{\skt{मूलसमुद्भवैः}}\Afootnote{K H; \skt{र्मूलसमुद्भवैः।}\ A}}\skt{~। }\pend
\pstart\versepar \skt{अथवा }\edtext{\skt{शंखपुष्पायाः}}{\lemma{\skt{शंखपुष्पायाः}}\Afootnote{K H; \skt{शङ्खपुष्प्या च}\ A}}\skt{ }\edtext{\skt{समानाः}}{\lemma{\skt{समानाः}}\Afootnote{K; \skt{समाना}\ A H}}\skt{ सर्वरूपतः~।}\pend

\pstart\versepar \ledsidenote{\passref{4.30.25}}\edtext{\skt{वेगेन}}{\lemma{\skt{वेगेन}}\Afootnote{K H; \skt{वेतेन}\ A}}\skt{ }\edtext{\skt{महताविष्टा}}{\lemma{\skt{महताविष्टा}}\Afootnote{K H; \skt{महता++आविष्टा}\ A}}\skt{ सर्प\-नि\-र्मो\-क\-स\-न्निभा~। }\pend
\pstart\versepar \skt{एषा वेगवती नाम जायते }\edtext{\skt{ह्यम्वुदक्षये}}{\lemma{\skt{अम्वुदक्षये}}\Afootnote{K H; \skt{अम्बुदक्षये॥}\ A}}\skt{~।}\pend

\pstart\versepar \ledsidenote{\passref{4.30.26}}\skt{सप्तादौ }\edtext{\skt{सर्परूपिण्यौ}}{\lemma{\skt{सर्परूपिण्यौ}}\Afootnote{K; \skt{सर्परूपिण्यो}\ A H}}\skt{ }\edtext{\skt{यस्त्वौषध्यः}}{\lemma{\skt{यस्त्वौषध्यः}}\Afootnote{K H; \skt{ह्यौषध्यो याः}\ A}}\skt{ }\edtext{\skt{प्रकीर्त्तिता}}{\lemma{\skt{प्रकीर्त्तिता}}\Afootnote{K; \skt{प्रकीर्तिताः}\ A H}}\skt{~। }\pend
\pstart\versepar \skt{तासामुद्धरणं कार्यं मन्त्रेणानेन सर्वदा~॥}\pend

\pstart\versepar \ledsidenote{\passref{4.30.27}}\edtext{\skt{महे\-न्द\-र\-रा\-म\-कृ\-ष्णानां}}{\lemma{\skt{महे\appbrk{}न्द\appbrk{}र\appbrk{}रा\appbrk{}म\appbrk{}कृ\appbrk{}ष्णानां}}\Afootnote{K; \skt{महे\appbrk{}न्द्र\appbrk{}रा\appbrk{}म\appbrk{}कृ\appbrk{}ष्णानां}\ A H}}\skt{ }\edtext{\edtext{\skt{व्रावामपि}}{\lemma{\skt{व्रावाम्}}\Afootnote{\skt{ब्राह्मणानां गवा}\ A; \skt{व्राह्मणानां गवाम्}\ K H}}}{\lemma{\skt{अपि}}\Afootnote{K H; \skt{मपि।}\ A}}\skt{~। }\pend
\pstart\versepar \skt{तपसा तेजसा }\edtext{\skt{चैव}}{\lemma{\skt{चैव}}\Afootnote{K H; \skt{वा+अपि}\ A}}\skt{ प्रशाम्यध्वं }\edtext{\skt{शिवाय}}{\lemma{\skt{शिवाय}}\Afootnote{A K; \skt{शिताय}\ H\textsuperscript{ac}; \skt{शिवाताय}\ H\textsuperscript{pc}}}\skt{ वै~॥}\pend

\pstart\versepar \ledsidenote{\passref{4.30.28}}\edtext{\skt{अनेन}}{\lemma{\skt{अनेन}}\Afootnote{K H; \skt{नत्रेणानेन}\ A}}\skt{ }\edtext{\skt{मन्त्रेण}}{\lemma{\skt{मन्त्रेण}}\Afootnote{K H; \skt{मतिमा}\ A}}\skt{ }\edtext{\skt{सर्व}}{\lemma{\skt{सर्व}}\Afootnote{K H; \skt{न्सर्वा}\ A}}\skt{ }\edtext{\skt{भेषजानि}}{\lemma{\skt{भेषजानि}}\Afootnote{K H; \skt{एवाभि}\ A}}\skt{ }\edtext{\skt{परिजपेदिति}}{\lemma{\skt{परिजपेदिति}}\Afootnote{K H; \skt{मन्त्रयेत्। अ}\ A}}\skt{~। }\pend
\pstart\versepar \skt{}\edtext{\skt{अश्र\-द्द\-धा\-नै\-र\-लसैः}}{\lemma{\skt{अश्र\appbrk{}द्द\appbrk{}धा\appbrk{}नै\appbrk{}र\appbrk{}लसैः}}\Afootnote{K H; \skt{श्रद्द\appbrk{}धा\appbrk{}नै\appbrk{}र\appbrk{}लसैः}\ A}}\skt{ कृतघ्नैः }\edtext{\skt{पापकर्तृभिः}}{\lemma{\skt{पापकर्तृभिः}}\Afootnote{K H; \skt{पापकर्मभिः॥}\ A}}\skt{~।}\pend

\pstart\versepar \ledsidenote{\passref{4.30.29}}\edtext{\skt{न शक्य\-मा\-सा\-द\-यितुं}}{\lemma{\skt{न शक्य\appbrk{}मा\appbrk{}सा\appbrk{}द\appbrk{}यितुं}}\Afootnote{K H; \skt{नैवासादयितुं}\ A}}\skt{ }\edtext{\skt{सोमाः}}{\lemma{\skt{सोमाः}}\Afootnote{K H; \skt{शक्याः सोमाः}\ A}}\skt{ }\edtext{\edtext{\skt{सोमासमास्तथा}}{\lemma{\skt{सोमासमास्}}\Afootnote{K; \skt{सोमसमा}\ A H}}}{\lemma{\skt{तथा}}\Afootnote{K; \skt{स्तथा}\ A H}}\skt{~। }\pend
\pstart\versepar \skt{पीता\-व\-शे\-ष\-म\-मृतं }\edtext{\edtext{\skt{देवै\-र्व्र\-ह्म\-पु\-रो\-गमैः}}{\lemma{\skt{देवैर्}}\Afootnote{K H; \skt{देवै}\ A}}}{\lemma{\skt{व्रह्मपुरोगमैः}}\Afootnote{K H; \skt{र्ब्रह्मपुरोगमैः॥}\ A}}\skt{~।}\pend

\pstart\versepar \ledsidenote{\passref{4.30.30}}\skt{निहितं सोमवीर्यासु सोमे }\edtext{\skt{चाती}}{\lemma{\skt{चाती}}\Afootnote{\skt{चाप्योषधीपतौ।}\ A; \skt{चाप्यौषधीपती~॥}\ K; \skt{चाप्यौषधीपतौ~॥}\ H}}\skt{~॥ }\pend
\pstart\versepar \skt{}\edtext{\skt{देवसूत्रे}}{\lemma{\skt{देवसूत्रे}}\Afootnote{K H; \skt{देवसुन्दे}\ A}}\skt{ ह्रदवरे तथा }\edtext{\skt{सिन्धु}}{\lemma{\skt{सिन्धु}}\Afootnote{K H; \skt{सिन्धौ}\ A}}\skt{ महानदे~।}\pend

\pstart\versepar \ledsidenote{\passref{4.30.31}}\skt{दृश्यते च जलान्तेषु }\edtext{\skt{देवी}}{\lemma{\skt{देवी}}\Afootnote{K H; \skt{मेध्या}\ A}}\skt{ }\edtext{\skt{व्रह्मसुर्वर्चला}}{\lemma{\skt{व्रह्मसुर्वर्चला}}\Afootnote{K; \skt{ब्रह्मसुवर्चला।}\ A; \skt{व्रह्मसुवर्च्चला~॥}\ H}}\skt{~। }\pend
\pstart\versepar \skt{}\edtext{\skt{आदित्यपर्ण्णिनी}}{\lemma{\skt{आदित्यपर्ण्णिनी}}\Afootnote{K H; \skt{आदित्यपर्णिनी}\ A}}\skt{ }\edtext{\skt{चैव}}{\lemma{\skt{चैव}}\Afootnote{K H; \skt{ज्ञेया}\ A}}\skt{ तथैव हिमसं क्षये~।}\pend

\pstart\versepar \ledsidenote{\passref{4.30.32}}\edtext{\skt{दृश्यतेजगरी}}{\lemma{\skt{दृश्यतेजगरी}}\Afootnote{K; \skt{दृश्य\appbrk{}ते\appbrk{}अ\appbrk{}ज\appbrk{}गरी}\ A H}}\skt{ नित्यं गोनसी }\edtext{\skt{चाम्वुदागमेः}}{\lemma{\skt{चाम्वुदागमेः}}\Afootnote{K; \skt{चाम्बुदागमे।}\ A; \skt{चाम्वुदागमे~॥}\ H}}\skt{~। }\pend
\pstart\versepar \skt{}\edtext{\skt{कश्मीरे तु}}{\lemma{\skt{कश्मीरे तु}}\Afootnote{K H; \skt{काश्मीरेषु}\ A}}\skt{ }\edtext{\skt{सदा}}{\lemma{\skt{सदा}}\Afootnote{K H; \skt{सरो दिव्यं नाम्ना क्षु}\ A}}\skt{ }\edtext{\skt{यन्त्रमानसं}}{\lemma{\skt{यन्त्रमानसं}}\Afootnote{K H; \skt{द्रकमानसम्॥}\ A}}\skt{ }\edtext{\skt{नामविश्रुतं}}{\lemma{\skt{नामविश्रुतं}}\Afootnote{K H; \emph{om.}\ A}}\skt{~।}\pend

\pstart\versepar \ledsidenote{\passref{4.30.33}}\skt{करेणुस्तत्र कन्या च }\edtext{\skt{च्छत्रातिच्छत्रिके}}{\lemma{\skt{च्छत्रातिच्छत्रिके}}\Afootnote{K H; \skt{छत्रातिच्छत्रके}\ A}}\skt{ तथा~। }\pend
\pstart\versepar \skt{गोलोमी चाजलोमी च महती श्रावणी तथा~।}\pend

\pstart\versepar \ledsidenote{\passref{4.30.34}}\edtext{\skt{हेमन्ते}}{\lemma{\skt{हेमन्ते}}\Afootnote{K H; \skt{वसन्ते}\ A}}\skt{ }\edtext{\skt{कृष्णसर्पाभा}}{\lemma{\skt{कृष्णसर्पाभा}}\Afootnote{K H; \skt{कृष्णसर्पाख्या}\ A}}\skt{ }\edtext{\skt{वसन्ते}}{\lemma{\skt{वसन्ते}}\Afootnote{K H; \skt{गोनसी}\ A}}\skt{ }\edtext{\skt{चात्र}}{\lemma{\skt{चात्र}}\Afootnote{K H; \skt{च प्र}\ A}}\skt{ दृश्यते~। }\pend
\pstart\versepar \skt{}\edtext{\skt{नदीं}}{\lemma{\skt{नदीं}}\Afootnote{K H; \emph{om.}\ A}}\skt{ }\edtext{\skt{कौशिकीमुत्तीर्य}}{\lemma{\skt{उत्तीर्य}}\Afootnote{K H; \skt{सरितं तीर्त्वा}\ A}}\skt{ }\edtext{\skt{संजयन्त्यास्तु}}{\lemma{\skt{संजयन्त्यास्तु}}\Afootnote{K; \skt{सञ्जयन्त्यास्तु}\ A; \skt{संजयन्यास्तु}\ H\textsuperscript{ac}; \skt{संज\appbrk{}य\appbrk{}न्या\appbrk{}वै\appbrk{}ज\appbrk{}य\appbrk{}न्त्यास्तु}\ H\textsuperscript{pc}}}\skt{ पूर्वतः~।}\pend

\pstart\versepar \ledsidenote{\passref{4.30.35}}\skt{क्षितिप्रदेशो }\edtext{\edtext{\skt{वल्मीकैरावृतो}}{\lemma{\skt{वल्मीकैर्}}\Afootnote{K H; \skt{वल्मीकै}\ A}}}{\lemma{\skt{आवृतो}}\Afootnote{K H; \skt{राचितो}\ A}}\skt{ योजनत्रयं~॥ }\pend
\pstart\versepar \skt{}\edtext{\skt{विधेया}}{\lemma{\skt{विधेया}}\Afootnote{K H; \skt{विज्ञेया}\ A}}\skt{ तत्र कापोती }\edtext{\skt{स्वेता}}{\lemma{\skt{स्वेता}}\Afootnote{K; \skt{श्वेता}\ A; \skt{स्वेत}\ H}}\skt{ }\edtext{\skt{वल्मीकमूधसु}}{\lemma{\skt{वल्मीकमूधसु}}\Afootnote{K; \skt{वल्मीकमूर्धसु}\ A H}}\skt{~।}\pend

\pstart\versepar \ledsidenote{\passref{4.30.36}}\skt{मलये नलसेतौ च }\edtext{\skt{वेगवत्यौषधी}}{\lemma{\skt{वेगवत्यौषधी}}\Afootnote{A K; \skt{वेगवत्यौषधी}\ H}}\skt{ ध्रुवा~। }\pend
\pstart\versepar \skt{कार्त्तिक्यां }\edtext{\skt{पौर्ण्णमास्यान्तु}}{\lemma{\skt{पौर्ण्णमास्यान्तु}}\Afootnote{K H; \skt{पौर्णमास्यां}\ A}}\skt{ }\edtext{\skt{दृश्यते}}{\lemma{\skt{दृश्यते}}\Afootnote{K H; \emph{om.}\ A}}\skt{ }\edtext{\skt{नात्र}}{\lemma{\skt{नात्र}}\Afootnote{K H; \emph{om.}\ A}}\skt{ }\edtext{\skt{संशयः}}{\lemma{\skt{संशयः}}\Afootnote{K H; \skt{च भक्ष\appbrk{}ये\appbrk{}त्ता\appbrk{}मु\appbrk{}पो\appbrk{}षितः॥}\ A}}\skt{~।}\pend

\pstart\versepar \ledsidenote{\passref{4.30.37}}\skt{सोमवच्चात्र वर्त्तेत }\edtext{\skt{विधि\-रा\-नि\-र्ग\-मा\-दिति}}{\lemma{\skt{विधि\appbrk{}रा\appbrk{}नि\appbrk{}र्ग\appbrk{}मा\appbrk{}दिति}}\Afootnote{K; \skt{फलं तावच्च कीर्तितम्।}\ A; \skt{विधिना निर्ग्गमादिति~॥}\ H}}\skt{~॥ }\pend
\pstart\versepar \skt{सर्वा }\edtext{\skt{विधेयास्त्वौषध्यः}}{\lemma{\skt{विधेयास्त्वौषध्यः}}\Afootnote{\skt{विचेयास्त्वोपध्यः}\ A; \skt{विधेयास्त्वौषध्यः}\ K H}}\skt{ }\edtext{\edtext{\edtext{\skt{सोमश्चाप्यर्वुदे}}{\lemma{\skt{सोमश्}}\Afootnote{K H; \skt{सोमाश्}\ A}}}{\lemma{\skt{चाप्य्}}\Afootnote{A; \skt{चा}\ K H}}}{\lemma{\skt{अर्वुदे}}\Afootnote{\skt{अर्बुदे}\ A; \skt{प्यर्वुदे}\ K H}}\skt{ गिरौ~।}\pend

\pstart\versepar \ledsidenote{\passref{4.30.38}}\edtext{\skt{स }\edtext{\edtext{\skt{शृङ्गै\-र्दे\-व\-च\-रि\-तै\-र\-म्वु\-दा\-नी\-क\-भे\-दिभिः}}{\lemma{\skt{देवचरितैर्}}\Afootnote{K H; \skt{र्देवचरितै}\ A}}}{\lemma{\skt{अम्वु\appbrk{}दा\appbrk{}नी\appbrk{}क\appbrk{}भे\appbrk{}दिभिः}}\Afootnote{K H; \skt{रम्बु\appbrk{}दा\appbrk{}नी\appbrk{}क\appbrk{}भे\appbrk{}दिभिः।}\ A}}}{\lemma{\skt{स शृङ्गैर्}}\Afootnote{K H; \skt{स शृङ्गै}\ A}}\skt{~। }\pend
\pstart\versepar \skt{}\edtext{\skt{व्याप्तस्तीथैश्च}}{\lemma{\skt{व्याप्तस्तीथैश्च}}\Afootnote{K H; \skt{व्याप्तस्तीर्थैश्च}\ A}}\skt{ विख्यातैः सिद्ध\-र्षि\-सु\-र\-से\-वितैः~।}\pend

\pstart\versepar \ledsidenote{\passref{4.30.39}}\edtext{\edtext{\skt{गुहा\-भि\-र्भी\-म\-रू\-पाभिः}}{\lemma{\skt{गुहाभिर्}}\Afootnote{A K; \skt{गुहाभि}\ H}}}{\lemma{\skt{भीमरूपाभिः}}\Afootnote{A K; \skt{र्व्भीमरूपाभिः}\ H}}\skt{ }\edtext{\skt{सिंहो\-न्ना\-दि\-त\-ना\-दिभिः}}{\lemma{\skt{सिंहो\appbrk{}न्ना\appbrk{}दि\appbrk{}त\appbrk{}ना\appbrk{}दिभिः}}\Afootnote{K\textsuperscript{pc} H\textsuperscript{pc}; \skt{सिंहोन्नादित}\ A; \skt{सिंहो\appbrk{}न्ना\appbrk{}दि\appbrk{}त\appbrk{}नादि}\ K\textsuperscript{ac} H\textsuperscript{ac}}}\skt{~। }\pend
\pstart\versepar \skt{गजा\-लो\-डि\-त\-तो\-या\-भि\-रा\-प\-गाभिः समन्ततः~। }\pend
\pstart\versepar \skt{}\edtext{\edtext{\edtext{\skt{विवि\-धै\-र्धा\-तु\-भि\-श्चि\-त्रै\-स्स\-र्व\-त्रै\-वो\-प\-शो\-भितः}}{\lemma{\skt{विविधैर्}}\Afootnote{K H; \skt{विविधै}\ A}}}{\lemma{\skt{धातुभिश्}}\Afootnote{K H; \skt{र्धातुभि}\ A}}}{\lemma{\skt{चित्रैस्}}\Afootnote{K H; \skt{श्चित्रैः}\ A}}\skt{~।}\edtext{}{\lemma{\skt{नादिभिः… समन्}}\Afootnote{\emph{om.}\ A}}\edtext{}{\lemma{\skt{भिः… समन्ततः}}\Afootnote{\emph{om.}\ K\textsuperscript{ac}}}\edtext{}{\lemma{\skt{भिः… समन्ततः}}\Afootnote{\emph{om.}\ H\textsuperscript{ac}}}\pend

\pstart\versepar \ledsidenote{\passref{4.30.40}}\skt{नदीषु शैलेषु सरस्सु चापि पुण्येष्वरण्येषु }\edtext{\skt{तथाश्रमेषु}}{\lemma{\skt{तथाश्रमेषु}}\Afootnote{K H; \skt{तथा++आश्रमेषु।}\ A}}\skt{~। }\pend
\pstart\versepar \skt{सर्वत्र सर्वाः परिमार्गितव्याः सर्वत्र }\edtext{\skt{भूमिर्हि}}{\lemma{\skt{भूमिर्हि}}\Afootnote{K H; \skt{भूमिर्हि}\ A}}\skt{ }\edtext{\edtext{\skt{वसून्निधत्त}}{\lemma{\skt{वसून्}}\Afootnote{K H; \skt{वसू}\ A}}}{\lemma{\skt{निधत्त}}\Afootnote{K; \skt{नि धत्ते॥}\ A; \skt{निधत्ता}\ H}}\skt{ }\edtext{\skt{इति}}{\lemma{\skt{इति}}\Afootnote{K H; \emph{om.}\ A}}\skt{~॥ }\edtext{\skt{चिल ०~॥ {\symfont ❈}}}{\lemma{\skt{चिल ०~॥ {\symfont ❈}}}\Afootnote{K; \emph{om.}\ A H}}\skt{~॥}\pend

\pstart\versepar \ledsidenote{\passref{4.30.40.add}}\edtext{\skt{शुक्रं}}{\lemma{\skt{शुक्रं}}\Afootnote{K; \skt{शूकम्}\ H}}\skt{ मुखं शोफ\-म\-ना\-ग\-तञ्च मिश्र\-न्त\-था\-क्षी\-ण\-व\-ले\-न्द्रि\-यञ्च~। }\pend
\pstart\versepar \skt{आवाधसर्व प्रतिषेधमेधा }\edtext{\skt{स्वभा\-व\-न्ता\-प\-नि\-वृ\-त्त\-कञ्च}}{\lemma{\skt{स्वभा\appbrk{}व\appbrk{}न्ता\appbrk{}प\appbrk{}नि\appbrk{}वृ\appbrk{}त्त\appbrk{}कञ्}}\Afootnote{K; \skt{स्वभा\appbrk{}व\appbrk{}स\appbrk{}न्ता\appbrk{}प\appbrk{}नि\appbrk{}वृ\appbrk{}त्त\appbrk{}कञ्}\ H}}\skt{~॥}\pend

\pstart \ledsidenote{\passref{4.30.trailer}}\edtext{\skt{चिकित्सितेषु}}{\lemma{\skt{चिकित्सितेषु}}\Afootnote{K; \skt{इति सुश्रु\appbrk{}त\appbrk{}सं\appbrk{}हि\appbrk{}तायां चिकित्सास्थाने निवृत्तसं}\ A; \skt{चिकित्सितेषु तृतीयोदश~॥~॥चिकित्साल~॥~॥}\ H}}\skt{ }\edtext{\skt{तृतीयो}}{\lemma{\skt{तृतीयो}}\Afootnote{K; \skt{तापीयं}\ A; \emph{om.}\ H}}\skt{ }\edtext{\skt{दशः}}{\lemma{\skt{दशः}}\Afootnote{K; \skt{रसायनं नाम त्रिंशोऽध्यायः~॥३०॥}\ A; \emph{om.}\ H}}\skt{~॥ ०~॥}\pend

\endnumbering
\beginnumbering
\eledsection*{Adhyāya 31}
\pstart \ledsidenote{\passref{4.31.1}}\skt{अथातः }\edtext{\skt{स्नेहो\-प\-यो\-ग\-चि\-कि\-त्सितं}}{\lemma{\skt{स्नेहो\appbrk{}प\appbrk{}यो\appbrk{}ग\appbrk{}चि\appbrk{}कि\appbrk{}त्सितं}}\Afootnote{K H; \skt{स्नेहो\appbrk{}प\appbrk{}यौ\appbrk{}गि\appbrk{}क\appbrk{}चि\appbrk{}कि\appbrk{}त्सितं}\ A}}\skt{ }\edtext{\skt{व्याख्यास्यामः}}{\lemma{\skt{व्याख्यास्यामः}}\Afootnote{A H; \emph{om.}\ K}}\pend

\pstart \ledsidenote{\passref{4.31.4}}\edtext{\skt{द्वियो\-नि\-श्च\-तु\-र्वि\-क\-ल्पो\-भि\-हितः}}{\lemma{\skt{द्वियो\appbrk{}नि\appbrk{}श्च\appbrk{}तु\appbrk{}र्वि\appbrk{}क\appbrk{}ल्पो\appbrk{}भि\appbrk{}हितः}}\Afootnote{K; \skt{तत्र द्वियो\appbrk{}नि\appbrk{}श्च\appbrk{}तु\appbrk{}र्वि\appbrk{}क\appbrk{}ल्पोऽ\appbrk{}भि\appbrk{}हितः}\ A; \skt{द्वियो\appbrk{}नि\appbrk{}श्च\appbrk{}तु\appbrk{}र्व्वि\appbrk{}क\appbrk{}ल्पोऽ\appbrk{}भि\appbrk{}हितः}\ H}}\skt{ }\edtext{\skt{स्नेहः}}{\lemma{\skt{स्नेहः}}\Afootnote{A K; \emph{om.}\ H}}\skt{ स्नेहगुणाश्च~। तत्र जङ्गमेभ्यो गव्यं }\edtext{\skt{स्नेहं}}{\lemma{\skt{स्नेहं}}\Afootnote{K H; \skt{घृतं}\ A}}\skt{ }\edtext{\skt{प्रधाणतमं}}{\lemma{\skt{प्रधाणतमं}}\Afootnote{K; \skt{प्रधानं,}\ A; \skt{प्रधानतमं~।}\ H}}\skt{ }\edtext{\skt{स्थाव\-रे\-भ्य\-स्ति\-ल\-स्नेह}}{\lemma{\skt{तिलस्नेह}}\Afootnote{K H; \skt{तिलतैलं प्रधानम्}\ A}}\skt{ }\edtext{\skt{इति}}{\lemma{\skt{इति}}\Afootnote{A K; \skt{इति~॥ ओ~॥}\ H}}\skt{~॥}\pend

\pstart \ledsidenote{\passref{4.31.5}}\skt{अत ऊर्ध्वं }\edtext{\skt{यथा\-प्र\-यो\-ज\-नतो}}{\lemma{\skt{यथा\appbrk{}प्र\appbrk{}यो\appbrk{}ज\appbrk{}नतो}}\Afootnote{K H; \skt{यथाप्रयोजनं}\ A}}\skt{ }\edtext{\skt{यथा\-प्र\-धा\-न\-तश्च}}{\lemma{\skt{यथा\appbrk{}प्र\appbrk{}धा\appbrk{}न\appbrk{}तश्च}}\Afootnote{K H; \skt{यथाप्रधानं च}\ A}}\skt{~। }\edtext{\skt{स्थावरस्नेहां}}{\lemma{\skt{स्थावरस्नेहां}}\Afootnote{K; \skt{स्थावरस्नेहा}\ A; \skt{स्थावरस्ने}\ H}}\skt{ }\edtext{\skt{वक्ष्यामः}}{\lemma{\skt{वक्ष्यामः}}\Afootnote{K; \skt{नुपदेक्ष्यामः}\ A; \skt{हं वक्ष्यामः~॥}\ H}}\skt{~। तत्र }\edtext{\skt{तिल्व\-कै\-र\-ण्ड\-कु\-शा\-म्र\-द\-न्ती\-प\-ला\-श\-श\-ङ्खि\-नी\-वि\-षा\-णि\-ग\-वा\-क्षी\-क\-म्पि\-ल्य\-क\-श\-म्पा\-क\-नी\-लि\-नी\-स्नेहा}}{\lemma{\skt{तिल्व\appbrk{}कै\appbrk{}र\appbrk{}ण्ड\appbrk{}कु\appbrk{}शा\appbrk{}म्र\appbrk{}द\appbrk{}न्ती\appbrk{}प\appbrk{}ला\appbrk{}श\appbrk{}श\appbrk{}ङ्खि\appbrk{}नी\appbrk{}वि\appbrk{}षा\appbrk{}णि\appbrk{}ग\appbrk{}वा\appbrk{}क्षी\appbrk{}क\appbrk{}म्पि\appbrk{}ल्य\appbrk{}क\appbrk{}श\appbrk{}म्पा\appbrk{}क\appbrk{}नी\appbrk{}लि\appbrk{}नी\appbrk{}स्नेहा}}\Afootnote{K; \skt{तिल्व\appbrk{}कै\appbrk{}र\appbrk{}ण्ड\appbrk{}को\appbrk{}शा\appbrk{}म्र\appbrk{}द\appbrk{}न्ती\appbrk{}द्र\appbrk{}व\appbrk{}न्ती\appbrk{}स\appbrk{}प्त\appbrk{}ला\appbrk{}श\appbrk{}ङ्खि\appbrk{}नी\appbrk{}प\appbrk{}ला\appbrk{}श\appbrk{}वि\appbrk{}षा\appbrk{}णि\appbrk{}का\appbrk{}ग\appbrk{}वा\appbrk{}क्षी\appbrk{}क\appbrk{}म्पि\appbrk{}ल्ल\appbrk{}क\appbrk{}श\appbrk{}म्पा\appbrk{}क\appbrk{}नी\appbrk{}लि\appbrk{}नी\appbrk{}स्नेहा}\ A; \skt{तिल्व\appbrk{}कै\appbrk{}र\appbrk{}र\appbrk{}ण्ड\appbrk{}कु\appbrk{}शा\appbrk{}म्र\appbrk{}द\appbrk{}न्ती\appbrk{}प\appbrk{}ला\appbrk{}श\appbrk{}शं\appbrk{}खि\appbrk{}नी\appbrk{}वि\appbrk{}षा\appbrk{}णि\appbrk{}का\appbrk{}ग\appbrk{}वा\appbrk{}क्षी\appbrk{}कं\appbrk{}पि\appbrk{}ल्ल्य\appbrk{}क\appbrk{}श\appbrk{}म्पा\appbrk{}क\appbrk{}नी\appbrk{}लि\appbrk{}नी\appbrk{}स्नेहा}\ H}}\skt{ }\edtext{\skt{विरेचयन्ति}}{\lemma{\skt{विरेचयन्ति}}\Afootnote{A K; \skt{विरचयन्ति~।}\ H}}\skt{~।. }\edtext{\skt{जीमू\-त\-क\-कू\-ट\-ज\-कृ\-त\-वे\-ध\-ने\-क्ष्वा\-कु\-म\-द\-न\-स्नेहा}}{\lemma{\skt{जीमू\appbrk{}त\appbrk{}क\appbrk{}कू\appbrk{}ट\appbrk{}ज\appbrk{}कृ\appbrk{}त\appbrk{}वे\appbrk{}ध\appbrk{}ने\appbrk{}क्ष्वा\appbrk{}कु\appbrk{}म\appbrk{}द\appbrk{}न\appbrk{}स्नेहा}}\Afootnote{K; \skt{जीमू\appbrk{}त\appbrk{}क\appbrk{}कु\appbrk{}ट\appbrk{}ज\appbrk{}कृ\appbrk{}त\appbrk{}वे\appbrk{}ध\appbrk{}ने\appbrk{}क्ष्वा\appbrk{}कु\appbrk{}धा\appbrk{}मा\appbrk{}र्ग\appbrk{}व\appbrk{}म\appbrk{}द\appbrk{}न\appbrk{}स्नेहा}\ A; \skt{जीमू\appbrk{}त\appbrk{}क\appbrk{}कु\appbrk{}ट\appbrk{}ज\appbrk{}कृ\appbrk{}त\appbrk{}वे\appbrk{}ध\appbrk{}ने\appbrk{}क्ष्वा\appbrk{}कु\appbrk{}म\appbrk{}द\appbrk{}न\appbrk{}स्नेहा}\ H}}\skt{ वामयन्ति~।. }\edtext{\skt{विडङ्गस्थू}}{\lemma{\skt{विडङ्गस्थू}}\Afootnote{K H; \skt{विडङ्गख}\ A}}\skt{ }\edtext{\skt{लम\-ञ्ज\-री\-म\-धु\-शि\-ग्रु\-सू\-र्य\-व\-ल्ली\-पी\-लु\-सि\-द्धा\-र्थ\-क\-ज्यो\-ति\-ष्म\-ती\-स्नेहाः}}{\lemma{\skt{लम\appbrk{}ञ्ज\appbrk{}री\appbrk{}म\appbrk{}धु\appbrk{}शि\appbrk{}ग्रु\appbrk{}सू\appbrk{}र्य\appbrk{}व\appbrk{}ल्ली\appbrk{}पी\appbrk{}लु\appbrk{}सि\appbrk{}द्धा\appbrk{}र्थ\appbrk{}क\appbrk{}ज्यो\appbrk{}ति\appbrk{}ष्म\appbrk{}ती\appbrk{}स्नेहाः}}\Afootnote{K H; \skt{रम\appbrk{}ञ्ज\appbrk{}री\appbrk{}म\appbrk{}धु\appbrk{}शि\appbrk{}ग्रु\appbrk{}सू\appbrk{}र्य\appbrk{}व\appbrk{}ल्ली\appbrk{}पी\appbrk{}लु\appbrk{}सि\appbrk{}द्धा\appbrk{}र्थ\appbrk{}क\appbrk{}ज्यो\appbrk{}ति\appbrk{}ष्म\appbrk{}ती\appbrk{}स्नेहाः}\ A}}\skt{ शिरो विरेचयन्ति~। }\edtext{\skt{कर\-ञ्ज\-पू\-ती\-क\-ल\-व\-ण\-मा\-तु\-लु\-ङ्गे\-ङ्गु\-दी\-कि\-रा\-त\-ति\-क्त\-क\-स्नेहाः}}{\lemma{\skt{कर\appbrk{}ञ्ज\appbrk{}पू\appbrk{}ती\appbrk{}क\appbrk{}ल\appbrk{}व\appbrk{}ण\appbrk{}मा\appbrk{}तु\appbrk{}लु\appbrk{}ङ्गे\appbrk{}ङ्गु\appbrk{}दी\appbrk{}कि\appbrk{}रा\appbrk{}त\appbrk{}ति\appbrk{}क्त\appbrk{}क\appbrk{}स्नेहाः}}\Afootnote{K H; \skt{कर\appbrk{}ञ्ज\appbrk{}पू\appbrk{}ती\appbrk{}क\appbrk{}कृ\appbrk{}त\appbrk{}मा\appbrk{}ल\appbrk{}मा\appbrk{}तु\appbrk{}लु\appbrk{}ङ्गे\appbrk{}ङ्गु\appbrk{}दी\appbrk{}कि\appbrk{}रा\appbrk{}त\appbrk{}ति\appbrk{}क्त\appbrk{}क\appbrk{}स्नेहा}\ A}}\skt{ }\edtext{\skt{दुष्ट\-व्र\-णे\-षू\-प\-योज्याः}}{\lemma{\skt{दुष्ट\appbrk{}व्र\appbrk{}णे\appbrk{}षू\appbrk{}प\appbrk{}योज्याः}}\Afootnote{K H; \skt{दुष्ट\appbrk{}व्र\appbrk{}णे\appbrk{}षू\appbrk{}प\appbrk{}यु\appbrk{}ज्यन्ते,}\ A}}\skt{~।. }\edtext{\skt{तुरु\-व\-क\-क\-पि\-त्थ\-भ\-ल्ला\-त\-क\-प\-टो\-ल\-स्नेहाः}}{\lemma{\skt{तुरु\appbrk{}व\appbrk{}क\appbrk{}क\appbrk{}पि\appbrk{}त्थ\appbrk{}भ\appbrk{}ल्ला\appbrk{}त\appbrk{}क\appbrk{}प\appbrk{}टो\appbrk{}ल\appbrk{}स्नेहाः}}\Afootnote{K; \skt{तुव\appbrk{}र\appbrk{}क\appbrk{}क\appbrk{}पि\appbrk{}त्थ\appbrk{}क\appbrk{}म्पि\appbrk{}ल्ल\appbrk{}क\appbrk{}भ\appbrk{}ल्ला\appbrk{}त\appbrk{}क\appbrk{}प\appbrk{}टो\appbrk{}ल\appbrk{}स्नेहा}\ A; \skt{तुर\appbrk{}व\appbrk{}क\appbrk{}क\appbrk{}पि\appbrk{}त्थक भल्ला\appbrk{}त\appbrk{}क\appbrk{}प\appbrk{}टो\appbrk{}ल\appbrk{}स्नेहा}\ H}}\skt{ महाव्याधिषु~।. }\edtext{\skt{त्रपु\-स\-त्वा\-रु\-के\-र्का\-र\-तु\-म्वु\-रु\-क\-ष्मा\-ण्ड\-स्नेहाः}}{\lemma{\skt{त्रपु\appbrk{}स\appbrk{}त्वा\appbrk{}रु\appbrk{}के\appbrk{}र्का\appbrk{}र\appbrk{}तु\appbrk{}म्वु\appbrk{}रु\appbrk{}क\appbrk{}ष्मा\appbrk{}ण्ड\appbrk{}स्नेहाः}}\Afootnote{K; \skt{त्रपु\appbrk{}सै\appbrk{}र्वा\appbrk{}रु\appbrk{}क\appbrk{}क\appbrk{}र्का\appbrk{}रु\appbrk{}क\appbrk{}तु\appbrk{}म्बी\appbrk{}कू\appbrk{}ष्मा\appbrk{}ण्ड\appbrk{}स्नेहा}\ A; \skt{त्रपु\appbrk{}षे\appbrk{}र्व्वा\appbrk{}रु\appbrk{}के\appbrk{}र्क्का\appbrk{}रु\appbrk{}तु\appbrk{}म्वु\appbrk{}रु\appbrk{}कू\appbrk{}ष्मा\appbrk{}ण्ड\appbrk{}स्नेहा}\ H}}\skt{ }\edtext{\skt{मूत्रसङ्गेषु~॥क\-पो\-त\-व\-ङ्क\-ह\-री\-त\-की\-स्नेहाः}}{\lemma{\skt{मूत्रसङ्गेषु~॥क\appbrk{}पो\appbrk{}त\appbrk{}व\appbrk{}ङ्क\appbrk{}ह\appbrk{}री\appbrk{}त\appbrk{}की\appbrk{}स्नेहाः}}\Afootnote{K H; \skt{मूत्रसङ्गेषु, कपो\appbrk{}त\appbrk{}व\appbrk{}ङ्का\appbrk{}व\appbrk{}ल्गु\appbrk{}ज\appbrk{}ह\appbrk{}री\appbrk{}त\appbrk{}की\appbrk{}स्नेहाः}\ A}}\skt{ }\edtext{\skt{शर्करार्श्मरिषु}}{\lemma{\skt{शर्करार्श्मरिषु}}\Afootnote{K; \skt{शर्कराश्मरीषु,}\ A; \skt{शर्क्कराश्मरिषु~।}\ H}}\skt{~।. }\edtext{\skt{कुसु\-म्भ\-स\-र्ष\-पा\-त\-सी\-पि\-चु\-म\-र्दा\-ति\-मु\-क्त\-क\-का\-ण्डी\-र\-स्नेहाः}}{\lemma{\skt{कुसु\appbrk{}म्भ\appbrk{}स\appbrk{}र्ष\appbrk{}पा\appbrk{}त\appbrk{}सी\appbrk{}पि\appbrk{}चु\appbrk{}म\appbrk{}र्दा\appbrk{}ति\appbrk{}मु\appbrk{}क्त\appbrk{}क\appbrk{}का\appbrk{}ण्डी\appbrk{}र\appbrk{}स्नेहाः}}\Afootnote{K H; \skt{कुसु\appbrk{}म्भ\appbrk{}स\appbrk{}र्ष\appbrk{}पा\appbrk{}त\appbrk{}सी\appbrk{}पि\appbrk{}चु\appbrk{}म\appbrk{}र्दा\appbrk{}ति\appbrk{}मु\appbrk{}क्त\appbrk{}क\appbrk{}भा\appbrk{}ण्डी\appbrk{}क\appbrk{}टु\appbrk{}तु\appbrk{}म्बी\appbrk{}क\appbrk{}ट\appbrk{}भी\appbrk{}स्नेहाः}\ A}}\skt{ प्रमेहेषु~।. }\edtext{\skt{नाडी\-के\-र\-प\-न\-स\-पि\-या\-ल\-वि\-ल्व\-म\-धूक}}{\lemma{\skt{नाडी\appbrk{}के\appbrk{}र\appbrk{}प\appbrk{}न\appbrk{}स\appbrk{}पि\appbrk{}या\appbrk{}ल\appbrk{}वि\appbrk{}ल्व\appbrk{}म\appbrk{}धूक}}\Afootnote{K; \skt{ताल\appbrk{}ना\appbrk{}रि\appbrk{}के\appbrk{}ल\appbrk{}प\appbrk{}न\appbrk{}स\appbrk{}मो\appbrk{}च\appbrk{}प्रि\appbrk{}या\appbrk{}ल\appbrk{}बि\appbrk{}ल्व\appbrk{}म\appbrk{}धूक}\ A; \skt{नालि\appbrk{}के\appbrk{}र\appbrk{}प\appbrk{}न\appbrk{}स\appbrk{}पि\appbrk{}या\appbrk{}ल\appbrk{}बि\appbrk{}ल्व\appbrk{}म\appbrk{}धूक}\ H}}\skt{ }\edtext{\skt{श्लेए\-ष्मा\-त\-क\-स्नेहाः}}{\lemma{\skt{श्लेए\appbrk{}ष्मा\appbrk{}त\appbrk{}क\appbrk{}स्नेहाः}}\Afootnote{K; \skt{श्लेष्मा\appbrk{}त\appbrk{}का\appbrk{}म्रा\appbrk{}त\appbrk{}क\appbrk{}फ\appbrk{}ल\appbrk{}स्नेहाः पित्तसंसृष्टे वायौ, बिभीतकभल्ला}\ A; \skt{श्लेष्मान्तकस्नेहाः पित्तसंसृष्टे वायौ~। विभी\appbrk{}त\appbrk{}क\appbrk{}पि\appbrk{}ण्डी\appbrk{}त\appbrk{}क\appbrk{}भ\appbrk{}ल्ला\appbrk{}त\appbrk{}क\appbrk{}स्नेहाः}\ H}}\skt{ पित्तसंसृष्टे वायौ~॥ }\edtext{\skt{विभी\-त\-क\-पि\-ण्डी\-त\-क\-भ\-ल्ला\-त\-क\-स्नेहाः }\edtext{\skt{कृष्णीकरणेषु}}{\lemma{\skt{कृष्णीकरणेषु}}\Afootnote{K H; \skt{कृष्णीकरणे,}\ A}}\skt{~।.}}{\lemma{\skt{तकस्नेहाः… विभीत}}\Afootnote{\emph{om.}\ A}}\skt{ }\edtext{\skt{श्रव\-ण\-कं\-गु\-क\-स्नेहाः}}{\lemma{\skt{श्रव\appbrk{}ण\appbrk{}कं\appbrk{}गु\appbrk{}क\appbrk{}स्नेहाः}}\Afootnote{K; \skt{श्रव\appbrk{}ण\appbrk{}क\appbrk{}ङ्गु\appbrk{}क\appbrk{}टु\appbrk{}ण्टु\appbrk{}क\appbrk{}स्नेहाः}\ A; \skt{स्रवणकंगुका स्नेहाः}\ H}}\skt{ }\edtext{\skt{पाण्डूकरनेषु}}{\lemma{\skt{पाण्डूकरनेषु}}\Afootnote{K; \skt{पाण्डूकरणे,}\ A; \skt{पाण्डूकरणेषु~।}\ H}}\skt{~। }\edtext{\skt{सर\-ला\-पी\-त\-दा\-रु\-शिं\-श\-पा\-स\-न\-सा\-र\-स्नेहाः}}{\lemma{\skt{सर\appbrk{}ला\appbrk{}पी\appbrk{}त\appbrk{}दा\appbrk{}रु\appbrk{}शिं\appbrk{}श\appbrk{}पा\appbrk{}स\appbrk{}न\appbrk{}सा\appbrk{}र\appbrk{}स्नेहाः}}\Afootnote{K; \skt{सर\appbrk{}ल\appbrk{}पी\appbrk{}त\appbrk{}दा\appbrk{}रु\appbrk{}शिं\appbrk{}श\appbrk{}पा\appbrk{}गु\appbrk{}रु\appbrk{}सा\appbrk{}र\appbrk{}स्नेहा}\ A; \skt{सर\appbrk{}ला\appbrk{}पी\appbrk{}त\appbrk{}दा\appbrk{}रु\appbrk{}शिं\appbrk{}स\appbrk{}पा\appbrk{}स\appbrk{}न\appbrk{}शा\appbrk{}न\appbrk{}सा\appbrk{}र\appbrk{}स्नेहाः~।}\ H}}\skt{ }\edtext{\skt{दर्द्रुकिटिभेषु}}{\lemma{\skt{दर्द्रुकिटिभेषु}}\Afootnote{K; \skt{दद्रु\appbrk{}कु\appbrk{}ष्ठ\appbrk{}कि\appbrk{}टि\appbrk{}मेषु,}\ A; \skt{दर्द्द्रूकिटिभेषु~।}\ H}}\skt{~। सर्व एव स्नेहा }\edtext{\skt{वातमुपहन्ति}}{\lemma{\skt{उपहन्ति}}\Afootnote{K H; \skt{उपघ्नन्ति,}\ A}}\skt{~। }\edtext{\skt{तैलगुणाः}}{\lemma{\skt{तैलगुणाः}}\Afootnote{K H; \skt{तैलगुणाश्च}\ A}}\skt{ समासेन व्याख्याताः~॥}\pend

\pstart \ledsidenote{\passref{4.31.6}}\skt{अत ऊर्ध्वं }\edtext{\skt{कषा\-य\-पा\-क\-क्र\-म\-मु\-प\-दे\-क्ष्यामः}}{\lemma{\skt{कषा\appbrk{}य\appbrk{}पा\appbrk{}क\appbrk{}क्र\appbrk{}मम्}}\Afootnote{K H; \skt{कषा\appbrk{}य\appbrk{}स्ने\appbrk{}ह\appbrk{}पा\appbrk{}क\appbrk{}क्र\appbrk{}मम्}\ A}}\skt{~॥ तत्र }\edtext{\skt{केचि\-दा\-हु\-स्त्व\-क्प\-त्र\-पु\-ष्प\-फ\-ला\-दीनां}}{\lemma{\skt{त्वक्प\appbrk{}त्र\appbrk{}पु\appbrk{}ष्प\appbrk{}फ\appbrk{}ला\appbrk{}दीनां}}\Afootnote{K H; \skt{त्वक्प\appbrk{}त्र\appbrk{}फ\appbrk{}ल\appbrk{}मू\appbrk{}ला\appbrk{}दीनां}\ A}}\skt{ }\edtext{\skt{भाग\-स्त\-च्च\-तु\-र्गु\-ण\-मु\-दकं}}{\lemma{\skt{भाग\appbrk{}स्त\appbrk{}च्च\appbrk{}तु\appbrk{}र्गु\appbrk{}ण\appbrk{}मु\appbrk{}दकं}}\Afootnote{K; \skt{भाग\appbrk{}स्त\appbrk{}च्च\appbrk{}तु\appbrk{}र्गुणं जलं}\ A; \skt{भाग\appbrk{}स्त\appbrk{}त्च\appbrk{}तु\appbrk{}र्ग्गु\appbrk{}ण\appbrk{}मु\appbrk{}दकं}\ H}}\skt{ }\edtext{\skt{चतु\-र्थां\-शा\-व\-शे\-षितं}}{\lemma{\skt{चतु\appbrk{}र्थां\appbrk{}शा\appbrk{}व\appbrk{}शे\appbrk{}षितं}}\Afootnote{K H; \skt{चतु\appbrk{}र्भा\appbrk{}गा\appbrk{}व\appbrk{}शेषं नि}\ A}}\skt{ }\edtext{\edtext{\skt{क्वाथ्या\-व\-ता\-र\-ये\-दि\-त्येष}}{\lemma{\skt{क्वाथ्या\appbrk{}व\appbrk{}ता\appbrk{}र\appbrk{}ये\appbrk{}दित्य्}}\Afootnote{\skt{ष्क्वाथ्या\appbrk{}प\appbrk{}ह\appbrk{}रे\appbrk{}दित्य्}\ A; \skt{क्वाथ्या\appbrk{}व\appbrk{}ता\appbrk{}र\appbrk{}येदि}\ K H}}}{\lemma{\skt{एष}}\Afootnote{A; \skt{त्येष}\ K H}}\skt{ }\edtext{\skt{कषायकल्पः}}{\lemma{\skt{कषायकल्पः}}\Afootnote{K H; \skt{कषा\appbrk{}य\appbrk{}पा\appbrk{}क\appbrk{}कल्पः; स्नेहप्रसृतेषु षट्सु चतुर्गुणं द्रवमावप्य चतु\appbrk{}र\appbrk{}श्चा\appbrk{}क्ष\appbrk{}स\appbrk{}मा\appbrk{}न्भे\appbrk{}ष\appbrk{}ज\appbrk{}पि\appbrk{}ण्डा\appbrk{}नि\appbrk{}त्येष स्नेहपाककल्पः। ए}\ A}}\skt{~। तत्तु न }\edtext{\skt{सम्य\-क्क\-स्मा\-दा\-ग\-म\-सि\-द्ध\-त्वात्}}{\lemma{\skt{सम्य\appbrk{}क्क\appbrk{}स्मा\appbrk{}दा\appbrk{}ग\appbrk{}म\appbrk{}सि\appbrk{}द्ध\appbrk{}त्वात्}}\Afootnote{K H; \skt{सम्यक्; कस्मा\appbrk{}ता\appbrk{}ग\appbrk{}मा\appbrk{}सि\appbrk{}द्ध\appbrk{}त्वात्~॥}\ A}}\pend

\pstart \ledsidenote{\passref{4.31.7}}\edtext{\edtext{\edtext{\skt{पल\-कु\-ड\-वा\-दी\-ना\-मि\-त्ये\-त\-च्चो\-प\-दे\-क्ष्यामः}}{\lemma{\skt{पल\appbrk{}कु\appbrk{}ड\appbrk{}वा\appbrk{}दी\appbrk{}नाम्}}\Afootnote{\skt{पल\appbrk{}कु\appbrk{}ड\appbrk{}वा\appbrk{}दीना}\ A K; \skt{पल\appbrk{}कु\appbrk{}ड\appbrk{}वा\appbrk{}डी\appbrk{}नाम्}\ H}}}{\lemma{\skt{इत्य्}}\Afootnote{\skt{मतो}\ A; \skt{मित्ये}\ K; \skt{इत्ये}\ H}}}{\lemma{\skt{एत\appbrk{}च्चो\appbrk{}प\appbrk{}दे\appbrk{}क्ष्यामः}}\Afootnote{\skt{मानं तु व्याख्यास्यामः}\ A; \skt{तच्चोपदेक्ष्यामः~।}\ K; \skt{तत्च उपदेक्ष्यामः~।}\ H}}\skt{~। तत्र द्वाद\-श\-धा\-न्य\-माषा }\edtext{\skt{मध्यमा}}{\lemma{\skt{मध्यमा}}\Afootnote{K H; \skt{मध्यमाः}\ A}}\skt{ }\edtext{\skt{सुव\-र्ण्ण\-मा\-ष\-कस्ते}}{\lemma{\skt{सुव\appbrk{}र्ण्ण\appbrk{}मा\appbrk{}ष\appbrk{}कस्ते}}\Afootnote{K; \skt{सुवर्णमाषकः ते}\ A H}}\skt{ }\edtext{\skt{शोडषसुवर्ण्णाः}}{\lemma{\skt{शोडषसुवर्ण्णाः}}\Afootnote{K; \skt{षोडश सुवर्णम्;}\ A; \skt{षोडशसुवर्ण्णः~।}\ H}}\skt{~। }\edtext{\skt{अतश्चोर्ध्वं}}{\lemma{\skt{अतश्चोर्ध्वं}}\Afootnote{K; \skt{अथवा मध्यमनिष्पावा एको\appbrk{}न\appbrk{}विं\appbrk{}श\appbrk{}ति\appbrk{}र्ध\appbrk{}रणं, ताव्य\appbrk{}र्ध\appbrk{}तृ\appbrk{}ती\appbrk{}यानि कर्षः; ततश्चोर्ध्वं}\ A; \skt{अतश्च ऊर्ध्वञ्}\ H}}\skt{ }\edtext{\skt{चतु\-र्गु\-ण\-म\-भि\-व\-र्ध\-यतः}}{\lemma{\skt{चतु\appbrk{}र्गु\appbrk{}ण\appbrk{}म\appbrk{}भि\appbrk{}व\appbrk{}र्ध\appbrk{}यतः}}\Afootnote{K; \skt{चतु\appbrk{}र्गु\appbrk{}ण\appbrk{}म\appbrk{}भि\appbrk{}व\appbrk{}र्ध\appbrk{}यन्तः}\ A; \skt{चतु\appbrk{}र्ग्गु\appbrk{}ण\appbrk{}म\appbrk{}भि\appbrk{}व\appbrk{}र्द्ध\appbrk{}यतः~।}\ H}}\skt{ पल\-कु\-ड\-व\-प्र\-स्था\-ढ\-क\-द्रोणा }\edtext{\skt{अभिनिष्पद्यन्ते}}{\lemma{\skt{अभिनिष्पद्यन्ते}}\Afootnote{K H; \skt{इत्य\appbrk{}भि\appbrk{}नि\appbrk{}ष्प\appbrk{}द्यन्ते,}\ A}}\skt{~। तुला पुनः पलशतं }\edtext{\edtext{\edtext{\skt{ताभि\-र्विं\-श\-ति\-र्भारः}}{\lemma{\skt{ताभिर्}}\Afootnote{K H; \skt{ताः पुन}\ A}}}{\lemma{\skt{विंशतिर्}}\Afootnote{K; \skt{र्विंशति}\ A; \skt{व्विंशति}\ H}}}{\lemma{\skt{भारः}}\Afootnote{K; \skt{र्भारः}\ A H}}\skt{ }\edtext{\edtext{\skt{शुष्काणामार्द्राणां}}{\lemma{\skt{शुष्काणाम्}}\Afootnote{K H; \skt{शुष्काणा}\ A}}}{\lemma{\skt{आर्द्राणां}}\Afootnote{K; \skt{मिदं मानम्, आर्द्र}\ A; \skt{आद्राणाम्}\ H}}\skt{ द्रवाणां च }\edtext{\edtext{\skt{द्विगुणमिति}}{\lemma{\skt{द्विगुणम्}}\Afootnote{K H; \skt{द्विगुण}\ A}}}{\lemma{\skt{इति}}\Afootnote{K H; \skt{मिति॥}\ A}}\skt{~।}\pend

\pstart \ledsidenote{\passref{4.31.8}}\edtext{\skt{स्नेह\-कु\-कु\-ड\-वे\-सा\-ण्य\-क्वा\-थ्य\-द्र\-व्य\-दु\-र्व्य\-प्रस्थो}}{\lemma{\skt{स्नेह\appbrk{}कु\appbrk{}कु\appbrk{}ड\appbrk{}वे\appbrk{}सा\appbrk{}ण्य\appbrk{}क्वा\appbrk{}थ्य\appbrk{}द्र\appbrk{}व्य\appbrk{}दु\appbrk{}र्व्य\appbrk{}प्रस्थो}}\Afootnote{K; \skt{तत्रा\appbrk{}न्य\appbrk{}त\appbrk{}म\appbrk{}प\appbrk{}रि\appbrk{}मा\appbrk{}ण\appbrk{}सं\appbrk{}मि\appbrk{}तानां यथायोगं त्वक्प\appbrk{}त्र\appbrk{}फ\appbrk{}ल\appbrk{}मू\appbrk{}ला\appbrk{}दी\appbrk{}ना\appbrk{}मा\appbrk{}त\appbrk{}प\appbrk{}प\appbrk{}रि\appbrk{}शो\appbrk{}षि\appbrk{}तानां छेद्यानि खण्ड\appbrk{}श\appbrk{}श्छे\appbrk{}द\appbrk{}यित्वा भेद्यान्यणुशो भेद\appbrk{}यि\appbrk{}त्वा+\appbrk{}अ\appbrk{}व\appbrk{}कु\appbrk{}ट्या\appbrk{}ष्ट\appbrk{}गु\appbrk{}णेन षोडशगुणेन वा+अ\appbrk{}म्भ\appbrk{}सा+\appbrk{}अ\appbrk{}भि\appbrk{}षि\appbrk{}च्य\appbrk{}स्थाल्यां चतु\appbrk{}र्भा\appbrk{}गा\appbrk{}व\appbrk{}शिष्टं क्वाथ\appbrk{}यि\appbrk{}त्वा+\appbrk{}अ\appbrk{}प\appbrk{}ह\appbrk{}रे\appbrk{}दि\appbrk{}त्येष कषा\appbrk{}य\appbrk{}या\appbrk{}क\appbrk{}कल्पः। स्नेहा\appbrk{}च्च\appbrk{}तु\appbrk{}र्भा\appbrk{}गा\appbrk{}व\appbrk{}शिष्टं क्वाथ\appbrk{}यि\appbrk{}त्वा+\appbrk{}अ\appbrk{}प\appbrk{}ह\appbrk{}रे\appbrk{}दि\appbrk{}त्येष कषा\appbrk{}य\appbrk{}पा\appbrk{}क\appbrk{}कल्पः। स्नेहाच्चतुर्गुणो द्रवः, स्नेह}\ A; \skt{स्नेहकुडवे साध्ये क्वाथ्यद्रव्यप्रस्थो}\ H}}\skt{ }\edtext{\edtext{\skt{विधे\-य\-स्त\-च्च\-तु\-र्गु\-ण\-मु\-दकं}}{\lemma{\skt{विधेयस्}}\Afootnote{K H; \emph{om.}\ A}}}{\lemma{\skt{तच्च\appbrk{}तु\appbrk{}र्गु\appbrk{}ण\appbrk{}मु\appbrk{}दकं}}\Afootnote{K; \skt{चतुर्थंशो भेषजकल्कः, तदैकध्यं संऋज्य विपचेदित्येष स्नेहपाककल्पः। अथवा तत्रोदकद्रोणे त्वक्प\appbrk{}त्र\appbrk{}फ\appbrk{}ल\appbrk{}मू\appbrk{}ला\appbrk{}दीनां}\ A; \skt{तत्च\appbrk{}तु\appbrk{}र्गु\appbrk{}ण\appbrk{}मु\appbrk{}दकं}\ H}}\skt{ }\edtext{\skt{चतु\-र्था\-ङ्गा\-व\-शे\-षितं}}{\lemma{\skt{चतु\appbrk{}र्था\appbrk{}ङ्गा\appbrk{}व\appbrk{}शे\appbrk{}षितं}}\Afootnote{K; \skt{तुलामावाप्य चतु\appbrk{}र्भा\appbrk{}गा\appbrk{}व\appbrk{}शिष्टं नि}\ A; \skt{चतु\appbrk{}र्थां\appbrk{}शा\appbrk{}व\appbrk{}शे\appbrk{}षितं}\ H}}\skt{ }\edtext{\edtext{\skt{क्वाथ्या\-व\-ता\-र\-ये\-दि\-त्येष}}{\lemma{\skt{क्वाथ्या\appbrk{}व\appbrk{}ता\appbrk{}र\appbrk{}ये\appbrk{}दित्य्}}\Afootnote{\skt{ष्क्वाथ्या\appbrk{}प\appbrk{}ह\appbrk{}रे\appbrk{}दित्य्}\ A; \skt{क्वाथ्या\appbrk{}व\appbrk{}ता\appbrk{}र\appbrk{}येदि}\ K H}}}{\lemma{\skt{एष}}\Afootnote{A; \skt{त्येष}\ K H}}\skt{ }\edtext{\skt{कषायकल्पः}}{\lemma{\skt{कषायकल्पः}}\Afootnote{K H; \skt{कषा\appbrk{}य\appbrk{}पा\appbrk{}क\appbrk{}कल्पः;}\ A}}\skt{~। }\edtext{\skt{स्नेह\-कु\-ड\-वे\-साण्य}}{\lemma{\skt{स्नेह\appbrk{}कु\appbrk{}ड\appbrk{}वे\appbrk{}साण्य}}\Afootnote{K; \skt{स्नेहकुडवे}\ A; \skt{स्नेहकुडवे साध्ये}\ H}}\skt{ भेषजपलं }\edtext{\edtext{\skt{कल्कमिष्टं}}{\lemma{\skt{कल्कम्}}\Afootnote{\skt{पिष्टं}\ A; \skt{कल्क}\ K H}}}{\lemma{\skt{इष्टं}}\Afootnote{\skt{कल्कं}\ A; \skt{मिष्टं}\ K; \skt{पिष्टञ्}\ H}}\skt{ }\edtext{\skt{चतुर्गुणञ्च}}{\lemma{\skt{चतुर्गुणञ्च}}\Afootnote{K H; \skt{चतुर्गुणं द्र}\ A}}\skt{ }\edtext{\edtext{\skt{कषायमावाप्य}}{\lemma{\skt{कषायम्}}\Afootnote{H; \skt{वमा}\ A; \skt{कषाय}\ K}}}{\lemma{\skt{आवाप्य}}\Afootnote{H; \skt{वाप्य}\ A; \skt{मावाप्य}\ K}}\skt{ }\edtext{\edtext{\skt{विपचेदित्येष}}{\lemma{\skt{विपचेदित्य्}}\Afootnote{\skt{विपचेदित्य्}\ A; \skt{विपचेदि}\ K; \skt{विपचेति}\ H}}}{\lemma{\skt{एष}}\Afootnote{A; \skt{त्येष}\ K H}}\skt{ }\edtext{\skt{स्नेहपाककल्प}}{\lemma{\skt{स्नेहपाककल्प}}\Afootnote{K H; \skt{स्नेहपाककल्पः॥}\ A}}\skt{ }\edtext{\skt{इति}}{\lemma{\skt{इति}}\Afootnote{K H; \emph{om.}\ A}}\skt{~॥}\pend

\pstart\versepar \ledsidenote{\passref{4.31.9}}\edtext{\skt{भ}}{\lemma{\skt{भ}}\Afootnote{K; \skt{भवतश्चात्र}\ A; \skt{भव~॥}\ H}}\skt{~।. स्नेह\-भे\-ष\-ज\-तो\-यानां }\edtext{\skt{मानं}}{\lemma{\skt{मानं}}\Afootnote{K; \skt{प्रमाणं}\ A; \skt{मान}\ H}}\skt{ }\edtext{\skt{यत्र तु}}{\lemma{\skt{यत्र तु}}\Afootnote{K; \skt{यत्र}\ A; \skt{यन्त्रं तु}\ H}}\skt{ नेरितं~। }\pend
\pstart\versepar \skt{}\edtext{\skt{तत्रायं}}{\lemma{\skt{तत्रायं}}\Afootnote{A K; \skt{तत्रायं वियं}\ H}}\skt{ }\edtext{\skt{विधिरास्थेयो}}{\lemma{\skt{विधिरास्थेयो}}\Afootnote{\skt{विधिरास्थेयो}\ A K H}}\skt{ निर्दिष्टे }\edtext{\skt{तद्वदेव}}{\lemma{\skt{तद्वदेव}}\Afootnote{H; \skt{तद्वदेव}\ A K}}\skt{ तु~॥}\pend

\pstart\versepar \ledsidenote{\passref{4.31.10}}\skt{अनुक्ते द्रवकार्ये तु सर्वत्र सलिलं }\edtext{\skt{स्मृतं}}{\lemma{\skt{स्मृतं}}\Afootnote{K H; \skt{मतम्।}\ A}}\skt{~। }\pend
\pstart\versepar \skt{}\edtext{\skt{कल्क\-क्वा\-था\-व\-नि\-र्देशे}}{\lemma{\skt{कल्क\appbrk{}क्वा\appbrk{}था\appbrk{}व\appbrk{}नि\appbrk{}र्देशे}}\Afootnote{A K; \skt{कल्क\appbrk{}क्वा\appbrk{}थ्या\appbrk{}व\appbrk{}नि\appbrk{}र्द्देशे}\ H}}\skt{ }\edtext{\skt{गणा\-त्त\-स्मा\-त्स\-मा\-ध\-येत्}}{\lemma{\skt{गणा\appbrk{}त्त\appbrk{}स्मा\appbrk{}त्स\appbrk{}मा\appbrk{}ध\appbrk{}येत्}}\Afootnote{K; \skt{गणा\appbrk{}त्त\appbrk{}स्मा\appbrk{}त्प्र\appbrk{}यो\appbrk{}ज\appbrk{}येत्॥}\ A; \skt{गणा\appbrk{}त्त\appbrk{}स्मा\appbrk{}त्स\appbrk{}मा\appbrk{}व\appbrk{}येत्~॥}\ H}}\skt{~॥}\pend

\pstart \ledsidenote{\passref{4.31.11}}\skt{अत ऊर्ध्वं स्नेह\-पा\-क\-क्र\-म\-मु\-प\-दे\-क्ष्यामः~। स तु }\edtext{\edtext{\skt{पाकस्तृविधो}}{\lemma{\skt{पाकस्}}\Afootnote{K H; \emph{om.}\ A}}}{\lemma{\skt{तृविधो}}\Afootnote{K; \skt{त्रिविधः}\ A H}}\skt{ }\edtext{\skt{मृदुश्चिक्कनः}}{\lemma{\skt{मृदुश्चिक्कनः}}\Afootnote{K; \skt{तद्यथा मृदुः, मध्यमः,}\ A; \skt{मृदुश्चिक्कनः}\ H}}\skt{ }\edtext{\skt{खरचिक्कण}}{\lemma{\skt{खरचिक्कण}}\Afootnote{K; \skt{खर}\ A; \skt{खरचिक्कन}\ H}}\skt{ इति~। }\edtext{\skt{धविवेकमात्रं}}{\lemma{\skt{धविवेकमात्रं}}\Afootnote{\skt{तत्र स्नेहौ\appbrk{}ष\appbrk{}धि\appbrk{}वि\appbrk{}वे\appbrk{}क\appbrk{}मात्रं}\ A; \skt{स्नेहौ\appbrk{}ष\appbrk{}ध\appbrk{}वि\appbrk{}वे\appbrk{}क\appbrk{}मात्रं}\ K H}}\skt{ यत्र भेषजं स }\edtext{\skt{मृदुः}}{\lemma{\skt{मृदुः}}\Afootnote{K H; \skt{मृदु}\ A}}\skt{~। }\edtext{\skt{मधू\-च्छि\-ष्ट\-मि\-व\-ले\-प\-यति}}{\lemma{\skt{मधू\appbrk{}च्छि\appbrk{}ष्ट\appbrk{}मि\appbrk{}व\appbrk{}ले\appbrk{}प\appbrk{}यति}}\Afootnote{K H; \skt{रिति, मधूच्छिष्टमिव विश\appbrk{}द\appbrk{}म\appbrk{}वि\appbrk{}लेपि}\ A}}\skt{ यत्र भेषजं स }\edtext{\skt{चिक्कणः}}{\lemma{\skt{चिक्कणः}}\Afootnote{K; \skt{मध्यमः,}\ A; \skt{चिक्कनः~।}\ H}}\skt{~। }\edtext{\skt{फेनानि मात्र}}{\lemma{\skt{फेनानि मात्र}}\Afootnote{K; \skt{फेनातिमात्र}\ H}}\skt{ }\edtext{\skt{वैलस्य}}{\lemma{\skt{वैलस्य}}\Afootnote{K; \skt{तैलस्य}\ H}}\skt{ शेषं }\edtext{\edtext{\skt{घृत\-व\-दा\-दि\-शे\-त्कृ\-ष्ण\-म\-व\-स\-न्न\-मी\-ष\-द्वि\-शदं}}{\lemma{\skt{घृत\appbrk{}व\appbrk{}दा\appbrk{}दि\appbrk{}शे\appbrk{}त्कृ\appbrk{}ष्ण\appbrk{}म\appbrk{}व\appbrk{}स\appbrk{}न्न\appbrk{}मी\appbrk{}ष\appbrk{}द्वि\appbrk{}शदं}}\Afootnote{K; \skt{घृतवदा दिशेत्~। कृष्ण\appbrk{}म\appbrk{}व\appbrk{}स\appbrk{}न्न\appbrk{}मी\appbrk{}ष\appbrk{}द्वि\appbrk{}षदं}\ H}}\skt{ चिक्कणं}}{\lemma{\skt{फेनानि… कृ}}\Afootnote{\emph{om.}\ A}}\skt{ च यत्र }\edtext{\skt{स}}{\lemma{\skt{स}}\Afootnote{K H; \skt{भेषजं स}\ A}}\skt{ }\edtext{\skt{खरचिक्कण}}{\lemma{\skt{खरचिक्कण}}\Afootnote{K H; \skt{खर}\ A}}\skt{ इति~।. अत ऊर्ध्वं }\edtext{\skt{प्रदग्धस्नेह}}{\lemma{\skt{प्रदग्धस्नेह}}\Afootnote{K; \skt{दग्धस्नेहो}\ A; \skt{प्रदग्धस्नेहे}\ H}}\skt{ भवति~। तं पुनः }\edtext{\skt{साधयेत्~। तत्र}}{\lemma{\skt{साधयेत्~। तत्र}}\Afootnote{K H; \skt{साधु साधयेत्। तत्र}\ A}}\skt{ }\edtext{\skt{पानाभ्य}}{\lemma{\skt{पानाभ्य}}\Afootnote{A K; \skt{यानाभ्य}\ H\textsuperscript{ac}; \skt{पायानाभ्य}\ H\textsuperscript{pc}}}\skt{ }\edtext{\skt{व नस्या\-भ्य\-ञ्ज\-नी\-य\-श्चि\-क्कणः}}{\lemma{\skt{व नस्या\appbrk{}भ्य\appbrk{}ञ्ज\appbrk{}नी\appbrk{}य\appbrk{}श्चि\appbrk{}क्कणः}}\Afootnote{\skt{वहारयोर्मृदुः, नस्या\appbrk{}भ्य\appbrk{}ङ्ग\appbrk{}यो\appbrk{}र्म\appbrk{}ध्यमः,}\ A; \skt{वहारयोर्म्मृदुः नस्या\appbrk{}भ्य\appbrk{}ञ्ज\appbrk{}नी\appbrk{}य\appbrk{}श्चि\appbrk{}क्कणः~।}\ K; \skt{वहारयोर्म्मृदुः~॥ नस्या\appbrk{}ञ्ज\appbrk{}नी\appbrk{}य\appbrk{}श्चि\appbrk{}क्कणः~।}\ H}}\skt{~। }\edtext{\skt{वस्तिस्नेहाः}}{\lemma{\skt{वस्तिस्नेहाः}}\Afootnote{K H; \skt{बस्ति}\ A}}\skt{ }\edtext{\skt{खरचिक्कण}}{\lemma{\skt{खरचिक्कण}}\Afootnote{K H; \skt{कर्णपूरण}\ A}}\skt{ }\edtext{\skt{इति}}{\lemma{\skt{इति}}\Afootnote{K H; \skt{योस्तु खर इति॥}\ A}}\skt{~॥}\edtext{}{\lemma{\skt{चिक्कणं च}}\Afootnote{A K; \skt{चिक्कनञ्च}\ H}}\pend

\pstart\versepar \ledsidenote{\passref{4.31.12}}\edtext{\skt{शब्दस्योपरमे}}{\lemma{\skt{शब्दस्योपरमे}}\Afootnote{K; \skt{भवतश्चात्र शब्दस्योपरमे}\ A; \skt{शब्दव्यूपरमे}\ H}}\skt{ प्राप्ते }\edtext{\skt{फेनस्योपरमे}}{\lemma{\skt{फेनस्योपरमे}}\Afootnote{K; \skt{फेनस्योपशमे}\ A; \skt{फेनस्योपदमे}\ H}}\skt{ तथा~।. }\pend
\pstart\versepar \skt{}\edtext{\skt{गन्ध\-व\-र्ण्ण\-र\-सा\-नाञ्च}}{\lemma{\skt{गन्ध\appbrk{}व\appbrk{}र्ण्ण\appbrk{}र\appbrk{}सा\appbrk{}नाञ्च}}\Afootnote{K H; \skt{गन्ध\appbrk{}व\appbrk{}र्ण\appbrk{}र\appbrk{}सा\appbrk{}दीनां}\ A}}\skt{ सम्पत्तौ }\edtext{\skt{सिद्धमादिशेत्}}{\lemma{\skt{सिद्धमादिशेत्}}\Afootnote{\skt{सिद्धिमादिशेत्॥}\ A; \skt{सिद्धमादिशेत्}\ K H}}\skt{~॥}\pend

\pstart \ledsidenote{\passref{4.31.14}}\edtext{\skt{अतः}}{\lemma{\skt{अतः}}\Afootnote{K H; \skt{अत ऊर्ध्वं}\ A}}\skt{ स्नेह\-पा\-न\-क्र\-म\-मु\-प\-दे\-क्ष्यामः~। अथ }\edtext{\skt{लघु\-को\-ष्ठा\-या\-तु\-राय}}{\lemma{\skt{लघु\appbrk{}को\appbrk{}ष्ठा\appbrk{}या\appbrk{}तु\appbrk{}राय}}\Afootnote{K H; \skt{खलु लघु\appbrk{}को\appbrk{}ष्ठा\appbrk{}या\appbrk{}तु\appbrk{}राय}\ A}}\skt{ }\edtext{\skt{कृत\-मं\-ग\-ला\-यो\-द\-य\-गि\-रि\-शि\-ख\-र\-प्र\-स्थिते}}{\lemma{\skt{कृत\appbrk{}मं\appbrk{}ग\appbrk{}ला\appbrk{}यो\appbrk{}द\appbrk{}य\appbrk{}गि\appbrk{}रि\appbrk{}शि\appbrk{}ख\appbrk{}र\appbrk{}प्र\appbrk{}स्थिते}}\Afootnote{K; \skt{कृत\appbrk{}म\appbrk{}ङ्ग\appbrk{}ल\appbrk{}स्व\appbrk{}स्ति\appbrk{}वा\appbrk{}च\appbrk{}ना\appbrk{}यो\appbrk{}द\appbrk{}य\appbrk{}गि\appbrk{}रि\appbrk{}शि\appbrk{}ख\appbrk{}र\appbrk{}सं\appbrk{}स्थिते}\ A; \skt{कृत\appbrk{}मं\appbrk{}ग\appbrk{}ला\appbrk{}यो\appbrk{}द\appbrk{}य\appbrk{}गि\appbrk{}रि\appbrk{}शि\appbrk{}क\appbrk{}र\appbrk{}स\appbrk{}म्प्र\appbrk{}स्थिते}\ H}}\skt{ }\edtext{\skt{प्रत\-प्त\-क\-न\-क\-पी\-त\-लो\-हिते}}{\lemma{\skt{प्रत\appbrk{}प्त\appbrk{}क\appbrk{}न\appbrk{}क\appbrk{}पी\appbrk{}त\appbrk{}लो\appbrk{}हिते}}\Afootnote{K H; \skt{प्रत\appbrk{}प्त\appbrk{}क\appbrk{}न\appbrk{}क\appbrk{}नि\appbrk{}क\appbrk{}र\appbrk{}पी\appbrk{}त\appbrk{}लो\appbrk{}हिते}\ A}}\skt{ सवितरि }\edtext{\skt{यथाबलं}}{\lemma{\skt{यथाबलं}}\Afootnote{A; \skt{यथावलं}\ K H}}\skt{ तैलस्य घृतस्य वा }\edtext{\skt{मात्रं}}{\lemma{\skt{मात्रं}}\Afootnote{K H; \skt{मात्रां}\ A}}\skt{ }\edtext{\skt{प्रयच्छेत्}}{\lemma{\skt{प्रयच्छेत्}}\Afootnote{K H; \skt{पातुं प्रयच्छेत्। पीतमात्रे चोष्णो\appbrk{}द\appbrk{}के\appbrk{}नो\appbrk{}प\appbrk{}स्पृश्य सोपानत्को यथासुखं विहरेत्॥}\ A}}\skt{~।}\pend

\pstart\versepar \ledsidenote{\passref{4.31.19}}\skt{केवलं पैत्तिके }\edtext{\edtext{\skt{सर्पिर्वातिके}}{\lemma{\skt{सर्पिर्}}\Afootnote{K H; \skt{सर्पि}\ A}}}{\lemma{\skt{वातिके}}\Afootnote{K H\textsuperscript{pc}; \skt{र्वातिके}\ A; \skt{व्वाति}\ H\textsuperscript{ac}}}\skt{ लवणान्वितं~। }\pend
\pstart\versepar \skt{देयं }\edtext{\skt{वहुकफे}}{\lemma{\skt{वहुकफे}}\Afootnote{K; \skt{बहुकफे}\ A H}}\skt{ }\edtext{\skt{तैलं}}{\lemma{\skt{तैलं}}\Afootnote{K H; \skt{चा}\ A}}\skt{ }\edtext{\skt{सव्यो\-ष\-क्षा\-र\-सं\-युतं}}{\lemma{\skt{सव्यो\appbrk{}ष\appbrk{}क्षा\appbrk{}र\appbrk{}सं\appbrk{}युतं}}\Afootnote{K H; \skt{पि व्योष\appbrk{}क्षा\appbrk{}र\appbrk{}स\appbrk{}मा\appbrk{}यु\appbrk{}तम्॥}\ A}}\skt{~।}\pend

\pstart\versepar \ledsidenote{\passref{4.31.21}}\edtext{\skt{स्नेहसात्म्यः}}{\lemma{\skt{स्नेहसात्म्यः}}\Afootnote{A K; \skt{स्नेहः सात्म्यः}\ H}}\skt{ क्लेशसहो }\edtext{\skt{दृढः}}{\lemma{\skt{दृढः}}\Afootnote{K H; \emph{om.}\ A}}\skt{ काले }\edtext{\skt{च}}{\lemma{\skt{च}}\Afootnote{K H; \skt{नात्युष्ण}\ A}}\skt{ शीतले }\pend
\pstart\versepar \skt{}\edtext{\edtext{\skt{अच्छ\-मे\-व\-पि\-वे\-त्स्ने\-ह\-म\-च्छ\-पानं}}{\lemma{\skt{अच्छ\appbrk{}मे\appbrk{}व\appbrk{}पि\appbrk{}वे\appbrk{}त्स्ने\appbrk{}हम्}}\Afootnote{K H; \skt{अच्छमेव पिबेत्स्नेह}\ A}}\skt{ हि}}{\lemma{\skt{अच्छपानं हि}}\Afootnote{K H; \skt{मच्छपानं हि}\ A}}\skt{ }\edtext{\skt{शोभणं}}{\lemma{\skt{शोभणं}}\Afootnote{K; \skt{पूजितम्॥}\ A; \skt{शोभनं~॥}\ H}}\skt{~।}\pend

\pstart\versepar \ledsidenote{\passref{4.31.22}}\edtext{\skt{शीते}}{\lemma{\skt{शीते}}\Afootnote{K H; \skt{शीत}\ A}}\skt{ काले }\edtext{\skt{दिवास्नेहमुष्णे}}{\lemma{\skt{दिवास्नेहमुष्णे}}\Afootnote{K; \skt{दिवा स्नेहमुष्ण}\ A; \skt{दिवा स्नेहमुष्ण}\ H}}\skt{ काले }\edtext{\skt{पिवेन्निशि}}{\lemma{\skt{पिवेन्निशि}}\Afootnote{K H; \skt{पिबेन्निशि।}\ A}}\skt{~। }\pend
\pstart\versepar \skt{}\edtext{\skt{वातपित्ताधिको}}{\lemma{\skt{वातपित्ताधिको}}\Afootnote{A K; \skt{वातपित्ताधिके}\ H}}\skt{ रात्रौ~। }\edtext{\skt{वातश्लेष्माधिको}}{\lemma{\skt{वातश्लेष्माधिको}}\Afootnote{A K; \skt{वातश्लेष्माधिके}\ H}}\skt{ दिवा~।}\pend

\pstart\versepar \ledsidenote{\passref{4.31.23}}\edtext{\skt{वात\-पि\-त्था\-धि\-क\-स्योष्णे}}{\lemma{\skt{वात\appbrk{}पि\appbrk{}त्था\appbrk{}धि\appbrk{}क\appbrk{}स्योष्णे}}\Afootnote{K; \skt{वात\appbrk{}पि\appbrk{}त्ता\appbrk{}धि\appbrk{}क\appbrk{}स्योष्णे}\ A H}}\skt{ }\edtext{\skt{मूर्च्छोन्माद}}{\lemma{\skt{मूर्च्छोन्माद}}\Afootnote{K; \skt{तृणमूर्च्छोन्माद}\ A; \skt{मूर्च्छोत्साद}\ H}}\skt{ }\edtext{\skt{तृषावहः}}{\lemma{\skt{तृषावहः}}\Afootnote{K H; \skt{कारकः।}\ A}}\skt{~। }\pend
\pstart\versepar \skt{}\edtext{\skt{शीते}}{\lemma{\skt{शीते}}\Afootnote{A K; \skt{शेत}\ H}}\skt{ वातकफार्त्तस्य }\edtext{\skt{गौरवारुच}}{\lemma{\skt{गौरवारुच}}\Afootnote{K H; \skt{गौरवारुचि}\ A}}\skt{ शूलकृत्~।}\pend

\pstart\versepar \ledsidenote{\passref{4.31.24}}\edtext{\skt{स्नेहपीतस्तु}}{\lemma{\skt{स्नेहपीतस्तु}}\Afootnote{K H; \skt{स्नेहपीतस्य}\ A}}\skt{ }\edtext{\skt{तृष्णायां}}{\lemma{\skt{तृष्णायां}}\Afootnote{K H; \skt{चेत्तृष्णा}\ A}}\skt{ }\edtext{\skt{पिवे\-दु\-ष्णो\-द\-क\-न्नरं}}{\lemma{\skt{पिवे\appbrk{}दु\appbrk{}ष्णो\appbrk{}द\appbrk{}क\appbrk{}न्नरं}}\Afootnote{K; \skt{पिबेदुष्णोदकं नरः।}\ A; \skt{पिवे\appbrk{}दु\appbrk{}ष्णो\appbrk{}द\appbrk{}क\appbrk{}न्नरं~॥}\ H}}\skt{~। }\pend
\pstart\versepar \skt{}\edtext{\skt{एव\-ञ्चा\-नु\-प\-शा\-म्यन्ती}}{\lemma{\skt{चानुपशाम्यन्ती}}\Afootnote{K H\textsuperscript{pc}; \skt{चानुपशाम्यन्त्यां}\ A; \skt{चानुशाम्यन्ती}\ H\textsuperscript{ac}}}\skt{ }\edtext{\skt{स्नेह\-मु\-ष्णा\-म्बु\-नो\-द्ध\-रेत्}}{\lemma{\skt{स्नेह\appbrk{}मु\appbrk{}ष्णा\appbrk{}म्बु\appbrk{}नो\appbrk{}द्ध\appbrk{}रेत्}}\Afootnote{\skt{स्नेहमुष्णाम्बुना वमेत्॥}\ A; \skt{स्नेह\appbrk{}मु\appbrk{}ष्णा\appbrk{}म्वु\appbrk{}नो\appbrk{}द्ध\appbrk{}रेत्~।}\ K; \skt{स्नेह\appbrk{}मु\appbrk{}ष्णा\appbrk{}म्वु\appbrk{}नो\appbrk{}द्ध\appbrk{}रेत्~।}\ H}}\skt{~।}\pend

\pstart \ledsidenote{\passref{4.31.25cd}}\skt{या मात्रा परिजीर्येत }\edtext{\skt{चतुर्भागगते न}}{\lemma{\skt{चतुर्भागगते न}}\Afootnote{K; \skt{चतु\appbrk{}र्भा\appbrk{}ग\appbrk{}ग\appbrk{}ते+अ}\ A; \skt{चतुर्व्भागगते न}\ H}}\skt{ }\edtext{\skt{हि}}{\lemma{\skt{हि}}\Afootnote{K H; \skt{हनि॥}\ A}}\pend

\pstart\versepar \ledsidenote{\passref{4.31.26}}\skt{सा मात्रा }\edtext{\edtext{\skt{दीप\-य\-त्य\-ग्नि\-म\-ल्प\-दोषे}}{\lemma{\skt{दीपयत्यग्निम्}}\Afootnote{K H; \skt{दीपयत्यग्नि}\ A}}\skt{ च}}{\lemma{\skt{अल्पदोषे च}}\Afootnote{K H; \skt{मल्पदोषे च}\ A}}\skt{ पूजिता~। }\pend
\pstart\versepar \skt{या मात्रा परिजीर्येत }\edtext{\skt{तत्रार्धदिवसे}}{\lemma{\skt{तत्रार्धदिवसे}}\Afootnote{K H; \skt{तथा+\appbrk{}अ\appbrk{}र्ध\appbrk{}दि\appbrk{}वसे}\ A}}\skt{ गते}\pend

\pstart\versepar \ledsidenote{\passref{4.31.27}}\edtext{\skt{सा वृष्यावृंहणीया च}}{\lemma{\skt{सा वृष्यावृंहणीया च}}\Afootnote{K; \skt{सा वृष्या बृंहणी या च}\ A; \skt{सा वृष्णावृंहनीया च}\ H}}\skt{ मध्यदोषे च पूजिता~। }\pend
\pstart\versepar \skt{या मात्रा परिजीर्येत }\edtext{\skt{चतु\-र्भा\-गा\-व\-शे\-षिते}}{\lemma{\skt{चतु\appbrk{}र्भा\appbrk{}गा\appbrk{}व\appbrk{}शे\appbrk{}षिते}}\Afootnote{A K; \skt{चतु\appbrk{}र्व्भा\appbrk{}गा\appbrk{}व\appbrk{}शे\appbrk{}षिते~॥}\ H}}\pend

\pstart\versepar \ledsidenote{\passref{4.31.28}}\skt{स्नेहनीया च सा मात्रा }\edtext{\skt{वहुदोषे च}}{\lemma{\skt{वहुदोषे च}}\Afootnote{K H; \skt{बहुदोषे च}\ A}}\skt{ शस्यते~। }\pend
\pstart\versepar \skt{या }\edtext{\skt{स्नेह\-मा\-त्रा\-जी\-र्येत}}{\lemma{\skt{स्नेह\appbrk{}मा\appbrk{}त्रा\appbrk{}जी\appbrk{}र्येत}}\Afootnote{K H; \skt{मात्रा परिजीर्येत्तु}\ A}}\skt{ }\edtext{\skt{जन्तोः}}{\lemma{\skt{जन्तोः}}\Afootnote{K H; \skt{तथा}\ A}}\skt{ }\edtext{\skt{परिणतेहनि}}{\lemma{\skt{परिणतेहनि}}\Afootnote{K H; \skt{परि\appbrk{}ण\appbrk{}ते+\appbrk{}अ\appbrk{}हनि॥}\ A}}\edtext{}{\lemma{\skt{शस्यते~। या}}\Afootnote{K H; \skt{पूजिता। या}\ A}}\pend

\pstart\versepar \ledsidenote{\passref{4.31.29}}\edtext{\skt{ग्लानिमूर्च्छामदां}}{\lemma{\skt{ग्लानिमूर्च्छामदां}}\Afootnote{K; \skt{ग्लानिमूर्च्छामदा}\ A H}}\skt{ }\edtext{\skt{हित्वा सा}}{\lemma{\skt{हित्वा सा}}\Afootnote{K H; \skt{न्हित्वा सा}\ A}}\skt{ मात्रा }\edtext{\skt{स्नेहनी}}{\lemma{\skt{स्नेहनी}}\Afootnote{K H; \skt{पूजिता}\ A}}\skt{ }\edtext{\skt{मता}}{\lemma{\skt{मता}}\Afootnote{K H; \skt{भवे}\ A}}\skt{~। }\pend
\pstart\versepar \skt{}\edtext{\skt{अहो\-रा\-त्रा\-द\-सं\-दुष्टा या}}{\lemma{\skt{अहो\appbrk{}रा\appbrk{}त्रा\appbrk{}द\appbrk{}सं\appbrk{}दुष्टा या}}\Afootnote{K; \skt{त्। अहो\appbrk{}रा\appbrk{}त्रा\appbrk{}द\appbrk{}सं\appbrk{}दुष्टा या}\ A; \skt{अहो\appbrk{}रा\appbrk{}त्रा\appbrk{}द\appbrk{}सं\appbrk{}तुष्टा या}\ H}}\skt{ मात्रा परिजीर्यति}\pend

\pstart\versepar \ledsidenote{\passref{4.31.30}}\skt{सा तु }\edtext{\skt{कुष्ठ\-वि\-षो\-न्मा\-द\-ग्र\-हा\-प\-स्मा\-र\-ना\-शनी}}{\lemma{\skt{कुष्ठ\appbrk{}वि\appbrk{}षो\appbrk{}न्मा\appbrk{}द\appbrk{}ग्र\appbrk{}हा\appbrk{}प\appbrk{}स्मा\appbrk{}र\appbrk{}ना\appbrk{}शनी}}\Afootnote{K H; \skt{कुष्ठ\appbrk{}वो\appbrk{}षो\appbrk{}न्मा\appbrk{}द\appbrk{}ग्र\appbrk{}हा\appbrk{}प\appbrk{}स्मा\appbrk{}र\appbrk{}ना\appbrk{}शिनी।}\ A}}\skt{~। }\pend
\pstart\versepar \skt{}\edtext{\skt{यथाग्निं}}{\lemma{\skt{यथाग्निं}}\Afootnote{K H; \skt{यथाग्नि}\ A}}\skt{ }\edtext{\skt{प्रतिमां}}{\lemma{\skt{प्रतिमां}}\Afootnote{K H; \skt{प्रथमां}\ A}}\skt{ मात्रां }\edtext{\skt{पाययेत}}{\lemma{\skt{पाययेत}}\Afootnote{A K; \skt{यापयेत}\ H}}\skt{ विचक्षणः~।}\pend

\pstart\versepar \ledsidenote{\passref{4.31.31}}\skt{पीतो }\edtext{\skt{ह्यतिवहुस्नेहो}}{\lemma{\skt{अतिवहुस्नेहो}}\Afootnote{K H; \skt{अतिबहुः स्नेहो}\ A}}\skt{ }\edtext{\skt{जन\-ये\-त्प्रा\-ण\-सं\-शयं}}{\lemma{\skt{जन\appbrk{}ये\appbrk{}त्प्रा\appbrk{}ण\appbrk{}सं\appbrk{}शयं}}\Afootnote{A K; \skt{जन\appbrk{}ये\appbrk{}त्प्रा\appbrk{}ण\appbrk{}सं\appbrk{}शयः~।}\ H}}\skt{~। }\pend
\pstart\versepar \skt{}\edtext{\skt{मिथ्या\-चा\-रा\-द्व\-हु\-त्वाद्वा}}{\lemma{\skt{मिथ्या\appbrk{}चा\appbrk{}रा\appbrk{}द्व\appbrk{}हु\appbrk{}त्वाद्वा}}\Afootnote{K H; \skt{मिथ्या\appbrk{}चा\appbrk{}रा\appbrk{}द्ब\appbrk{}हु\appbrk{}त्वाद्वा}\ A}}\skt{ यस्य स्नेहो न जीर्यति~।}\pend

\pstart\versepar \ledsidenote{\passref{4.31.32}}\skt{विष्टभ्य }\edtext{\skt{वापि}}{\lemma{\skt{वापि}}\Afootnote{K H; \skt{चापि}\ A}}\skt{ }\edtext{\skt{जीर्येत}}{\lemma{\skt{जीर्येत}}\Afootnote{K H; \skt{जीर्येत्तं}\ A}}\skt{ }\edtext{\skt{वारिणौष्णेन}}{\lemma{\skt{वारिणौष्णेन}}\Afootnote{K; \skt{वारिणोष्णेन}\ A H}}\skt{ वामयेत्~। }\pend
\pstart\versepar \skt{ततः स्नेहं }\edtext{\skt{पुन\-र्द\-द्या\-ल्ल\-घु\-को\-ष्ठाय}}{\lemma{\skt{पुन\appbrk{}र्द\appbrk{}द्या\appbrk{}ल्ल\appbrk{}घु\appbrk{}को\appbrk{}ष्ठाय}}\Afootnote{A K; \skt{पुन\appbrk{}र्द्द\appbrk{}द्या\appbrk{}ल्ल\appbrk{}घु\appbrk{}को\appbrk{}ष्ठाय}\ H}}\skt{ देहिने~। }\pend
\pstart\versepar \skt{}\edtext{\skt{जीर्ण्णा\-जी\-र्ण्ण\-वि\-शं\-कायां}}{\lemma{\skt{जीर्ण्णा\appbrk{}जी\appbrk{}र्ण्ण\appbrk{}वि\appbrk{}शं\appbrk{}कायां}}\Afootnote{A K; \skt{जीर्ण्णे\appbrk{}जी\appbrk{}र्ण्ण\appbrk{}वि\appbrk{}श\appbrk{}ङ्कायां}\ H}}\skt{ }\edtext{\skt{पिवेदुष्णोदकं}}{\lemma{\skt{पिवेदुष्णोदकं}}\Afootnote{K H; \skt{स्नेहस्योष्णोदकं}\ A}}\skt{ }\edtext{\skt{नरः}}{\lemma{\skt{नरः}}\Afootnote{K H; \skt{पिबेत्॥}\ A}}\skt{~।}\pend

\pstart \ledsidenote{\passref{4.31.33ab}}\skt{तेनोद्गारो }\edtext{\skt{भवेत्छुद्धो}}{\lemma{\skt{भवेत्छुद्धो}}\Afootnote{K; \skt{भवेच्छुद्धो}\ A H}}\skt{ }\edtext{\edtext{\skt{रुचिश्चान्नं}}{\lemma{\skt{रुचिश्}}\Afootnote{K H; \emph{om.}\ A}}}{\lemma{\skt{चान्नं}}\Afootnote{K H; \emph{om.}\ A}}\skt{ }\edtext{\skt{भवेत्प्रति}}{\lemma{\skt{भवेत्प्रति}}\Afootnote{K H; \skt{भक्तं रुचिस्तथा।}\ A}}\skt{~।}\pend

\pstart\versepar \ledsidenote{\passref{4.31.34}}\skt{परि\-षि\-च्या\-द्भि\-रु\-ष्णाभिः }\edtext{\skt{जीर्ण्णस्नेहे}}{\lemma{\skt{जीर्ण्णस्नेहे}}\Afootnote{K; \skt{जीर्णस्नेहं}\ A; \skt{जीर्ण्णे स्नेहो}\ H}}\skt{ }\edtext{\skt{ततोनरं}}{\lemma{\skt{ततोनरं}}\Afootnote{K H; \skt{ततो नरम्।}\ A}}\skt{~। }\pend
\pstart\versepar \skt{}\edtext{\skt{यवागू}}{\lemma{\skt{यवागू}}\Afootnote{K; \skt{यवागूं}\ A H}}\skt{ }\edtext{\skt{पाययेदुष्णां}}{\lemma{\skt{पाययेदुष्णां}}\Afootnote{K; \skt{पाययेच्चोष्णां}\ A; \skt{पाययेदुष्णां}\ H}}\skt{ }\edtext{\skt{सुक्लि\-न्ना\-म\-ल्प\-त\-ण्डुलां}}{\lemma{\skt{सुक्लि\appbrk{}न्ना\appbrk{}म\appbrk{}ल्प\appbrk{}त\appbrk{}ण्डुलां}}\Afootnote{K; \skt{कामं क्लिन्ना\appbrk{}ल्प\appbrk{}त\appbrk{}ण्डु\appbrk{}लाम्॥}\ A; \skt{सुक्लि\appbrk{}न्ना\appbrk{}म\appbrk{}ल्प\appbrk{}त\appbrk{}ण्डुलां~॥}\ H}}\skt{~।}\pend

\pstart\versepar \ledsidenote{\passref{4.31.35}}\skt{देयौ }\edtext{\skt{युषरसो}}{\lemma{\skt{युषरसो}}\Afootnote{K; \skt{यूषर}\ A; \skt{यूषरौ}\ H}}\skt{ }\edtext{\edtext{\skt{स्यातामकृतौ}}{\lemma{\skt{स्याताम्}}\Afootnote{K H; \skt{सौ वा+}\ A}}}{\lemma{\skt{अकृतौ}}\Afootnote{K H; \skt{अपि सुगन्धी}\ A}}\skt{ }\edtext{\skt{सैन्धवान्वितौ}}{\lemma{\skt{सैन्धवान्वितौ}}\Afootnote{K H; \skt{स्नेहवर्जितौ।}\ A}}\skt{~। }\pend
\pstart\versepar \skt{कृतौ }\edtext{\skt{चाप्यल्पसर्पिष्कौ}}{\lemma{\skt{चाप्यल्पसर्पिष्कौ}}\Afootnote{\skt{वा+अ\appbrk{}त्य\appbrk{}ल्प\appbrk{}स\appbrk{}र्पिष्कौ}\ A; \skt{चाप्यल्पसर्पिष्कौ}\ K H}}\skt{ }\edtext{\skt{विलेपीता}}{\lemma{\skt{विलेपीता}}\Afootnote{K H; \skt{विलेपी वा}\ A}}\skt{ विधीयते~।}\pend

\pstart\versepar \ledsidenote{\passref{4.31.36}}\edtext{\skt{पिवेद्द्व्यहं}}{\lemma{\skt{पिवेद्द्व्यहं}}\Afootnote{K H; \skt{पिबेत्त्र्यहं}\ A}}\skt{ }\edtext{\skt{त्र्यहं}}{\lemma{\skt{त्र्यहं}}\Afootnote{K H; \skt{चतु}\ A}}\skt{ }\edtext{\skt{वापि}}{\lemma{\skt{वापि}}\Afootnote{K H; \skt{रहं}\ A}}\skt{ पञ्चाहं }\edtext{\skt{षडहानि वा}}{\lemma{\skt{षडहानि वा}}\Afootnote{K H; \skt{षडहं तथा।}\ A}}\skt{~। }\pend
\pstart\versepar \skt{सप्तरात्रात्परं स्नेहः }\edtext{\skt{सात्मीभावाय}}{\lemma{\skt{सात्मीभावाय}}\Afootnote{K; \skt{सात्म्यीभ}\ A; \skt{सात्मीभावाय}\ H}}\skt{ }\edtext{\skt{कल्पते}}{\lemma{\skt{कल्पते}}\Afootnote{K H; \skt{वति सेवितः॥}\ A}}\skt{~।}\pend

\pstart\versepar \ledsidenote{\passref{4.31.37}}\edtext{\skt{स्नेहद्विषः}}{\lemma{\skt{स्नेहद्विषः}}\Afootnote{K H; \skt{सुकुमारं कृशं वृद्धं शिशुं स्नेहद्वि}\ A}}\skt{ }\edtext{\skt{कृषां}}{\lemma{\skt{कृषां}}\Afootnote{K H; \skt{षं}\ A}}\skt{ }\edtext{\skt{वृद्धां सुकुमारं शिशूनपि}}{\lemma{\skt{वृद्धां सुकुमारं शिशूनपि}}\Afootnote{K; \skt{वृद्धां\appbrk{}सु\appbrk{}कु\appbrk{}मा\appbrk{}रा\appbrk{}न्शि\appbrk{}शू\appbrk{}नपि~।}\ H}}\skt{~। }\pend
\pstart\versepar \skt{}\edtext{\skt{तृष्णा\-लू\-ञ्चो\-ष्ण\-काले च}}{\lemma{\skt{तृष्णा\appbrk{}लू\appbrk{}ञ्चो\appbrk{}ष्ण\appbrk{}काले च}}\Afootnote{K; \skt{था। तृष्णा\appbrk{}र्त\appbrk{}मु\appbrk{}ष्ण\appbrk{}काले च}\ A; \skt{कृष्णा\appbrk{}लू\appbrk{}ञ्चो\appbrk{}ष्ण\appbrk{}काले च}\ H\textsuperscript{ac}; \skt{कृतृ\appbrk{}ष्णा\appbrk{}लू\appbrk{}ञ्चो\appbrk{}ष्ण\appbrk{}काले च}\ H\textsuperscript{pc}}}\skt{ सह भक्तेन पाययेत्~॥}\edtext{}{\lemma{\skt{वृद्धां… शिशूनपि}}\Afootnote{\emph{om.}\ A}}\pend

\pstart\versepar \ledsidenote{\passref{4.31.38}}\skt{पिप्पल्यो लवणं }\edtext{\edtext{\skt{स्नेहाश्चत्वारो}}{\lemma{\skt{स्नेहाश्}}\Afootnote{K H; \skt{स्नेहा}\ A}}}{\lemma{\skt{चत्वारो}}\Afootnote{K H; \skt{श्चत्वारो}\ A}}\skt{ }\edtext{\skt{दधिमस्तु च}}{\lemma{\skt{दधिमस्तु च}}\Afootnote{K H; \skt{दधिमस्तुकः।}\ A}}\skt{~। }\pend
\pstart\versepar \skt{}\edtext{\skt{पीत\-मे\-क\-ध्य\-मे\-तत्तु}}{\lemma{\skt{पीत\appbrk{}मे\appbrk{}क\appbrk{}ध्य\appbrk{}मे\appbrk{}तत्तु}}\Afootnote{K H; \skt{पीत\appbrk{}मै\appbrk{}क\appbrk{}ध्य\appbrk{}मे\appbrk{}तद्धि}\ A}}\skt{ }\edtext{\skt{सद्य}}{\lemma{\skt{सद्य}}\Afootnote{K; \skt{सद्यः}\ A H}}\skt{ }\edtext{\skt{स्नेहनमुच्यते}}{\lemma{\skt{स्नेहनमुच्यते}}\Afootnote{A K; \skt{स्नेहनमुच्यते~॥}\ H}}\skt{~।}\pend

\pstart\versepar \ledsidenote{\passref{4.31.39}}\edtext{\skt{भृष्टमांसरसे}}{\lemma{\skt{भृष्टमांसरसे}}\Afootnote{K; \skt{भृष्टा मांसरसे}\ A; \skt{भृष्टमांसरस}\ H}}\skt{ स्निग्धा }\edtext{\skt{यवागू}}{\lemma{\skt{यवागू}}\Afootnote{K H; \skt{यवागूः}\ A}}\skt{ }\edtext{\skt{सूपल्पिता}}{\lemma{\skt{सूपल्पिता}}\Afootnote{K; \skt{सूपकल्पिता}\ A H}}\skt{~। }\pend
\pstart\versepar \skt{}\edtext{\skt{सक्षौद्राः}}{\lemma{\skt{सक्षौद्राः}}\Afootnote{K H; \skt{प्रक्षुद्रा}\ A}}\skt{ पीयमाना तु }\edtext{\skt{सद्य\-स्ने\-ह\-न\-मु\-च्यते}}{\lemma{\skt{सद्य\appbrk{}स्ने\appbrk{}ह\appbrk{}न\appbrk{}मु\appbrk{}च्यते}}\Afootnote{K; \skt{सद्यः\appbrk{}स्ने\appbrk{}ह\appbrk{}न\appbrk{}मु\appbrk{}च्यते॥}\ A; \skt{सद्यः स्नेहनमुच्यते~॥}\ H}}\skt{~॥}\pend

\pstart\versepar \ledsidenote{\passref{4.31.42}}\edtext{\skt{शर्क\-रा\-घृ\-त\-सं\-सृष्टे}}{\lemma{\skt{शर्क\appbrk{}रा\appbrk{}घृ\appbrk{}त\appbrk{}सं\appbrk{}सृष्टे}}\Afootnote{K H; \skt{शर्क\appbrk{}रा\appbrk{}चू\appbrk{}र्ण\appbrk{}सं\appbrk{}सृष्टे}\ A}}\skt{ }\edtext{\skt{दुह्या\-द्गां\-क\-ल\-शेथ}}{\lemma{\skt{दुह्या\appbrk{}द्गां\appbrk{}क\appbrk{}ल\appbrk{}शेथ}}\Afootnote{K H; \skt{दोहनस्थे घृते}\ A}}\skt{ तत्~। }\pend
\pstart\versepar \skt{पाय\-ये\-दृ\-क्ष\-मे\-तधि }\edtext{\skt{सद्य}}{\lemma{\skt{सद्य}}\Afootnote{K; \skt{सद्यः}\ A H}}\skt{ }\edtext{\skt{स्नेहनमुच्यते}}{\lemma{\skt{स्नेहनमुच्यते}}\Afootnote{A K; \skt{स्नेहनमुच्यते~॥}\ H}}\skt{~॥}\edtext{}{\lemma{\skt{तत्~। पाय\appbrk{}ये\appbrk{}दृ\appbrk{}क्ष\appbrk{}मे\appbrk{}तधि}}\Afootnote{K; \skt{तु गाम्। दुग्ध्वा क्षीरं पिबेद्रूक्षः}\ A; \skt{तत्~। पाय\appbrk{}ये\appbrk{}दृ\appbrk{}क्ष\appbrk{}मे\appbrk{}तद्धि}\ H}}\pend

\pstart\versepar \ledsidenote{\passref{4.31.43}}\skt{यव\-को\-ल\-कु\-ल\-त्थानां क्वाथो }\edtext{\skt{मागधिका}}{\lemma{\skt{मागधिका}}\Afootnote{K H; \skt{मागधिकान्}\ A}}\skt{ }\edtext{\skt{युताः}}{\lemma{\skt{युताः}}\Afootnote{K H; \skt{वितः।}\ A}}\skt{~। }\pend
\pstart\versepar \skt{पयो दधि सुरा चेति }\edtext{\skt{घृतमप्यष्टमं}}{\lemma{\skt{घृतमप्यष्टमं}}\Afootnote{\skt{घृतमप्यष्टमं}\ A K H}}\skt{ भवेत्~।}\pend

\pstart\versepar \ledsidenote{\passref{4.31.44}}\edtext{\skt{सिद्धमेतद्घृतं}}{\lemma{\skt{सिद्धमेतद्घृतं}}\Afootnote{K H; \skt{सिद्धमेतैर्घृतं}\ A}}\skt{ पीतं }\edtext{\skt{सद्य\-स्ने\-ह\-न\-मु\-च्यते}}{\lemma{\skt{सद्य\appbrk{}स्ने\appbrk{}ह\appbrk{}न\appbrk{}मु\appbrk{}च्यते}}\Afootnote{K; \skt{सद्यः\appbrk{}स्ने\appbrk{}ह\appbrk{}न\appbrk{}मु\appbrk{}त्त\appbrk{}मम्।}\ A; \skt{सद्यः स्नेहनमुच्यते~।}\ H}}\skt{~।. }\pend
\pstart\versepar \skt{}\edtext{\skt{नृपाणां}}{\lemma{\skt{नृपाणां}}\Afootnote{K; \skt{राज्ञे रा}\ A; \skt{नृपानान्}\ H}}\skt{ }\edtext{\skt{तत्समानानां}}{\lemma{\skt{तत्समानानां}}\Afootnote{K H; \skt{जसमेभ्यो वा}\ A}}\skt{ }\edtext{\skt{देय\-मे\-त\-द्घृ\-तो\-त्तमं}}{\lemma{\skt{देय\appbrk{}मे\appbrk{}त\appbrk{}द्घृ\appbrk{}तो\appbrk{}त्तमं}}\Afootnote{K; \skt{देय\appbrk{}मे\appbrk{}त\appbrk{}द्घृ\appbrk{}तो\appbrk{}त्त\appbrk{}मम्॥}\ A; \skt{देय\appbrk{}मे\appbrk{}त\appbrk{}द्घृ\appbrk{}तो\appbrk{}त्तमं~।}\ H}}\skt{~।}\pend

\pstart\versepar \ledsidenote{\passref{4.31.45}}\edtext{\skt{बलहीनेषु}}{\lemma{\skt{बलहीनेषु}}\Afootnote{A K; \skt{वलहीनेषु}\ H}}\skt{ वृद्धेषु मृद्वग्निस्त्री }\edtext{\skt{महात्मनि}}{\lemma{\skt{महात्मनि}}\Afootnote{K H; \skt{हतात्मसु।}\ A}}\skt{~। }\pend
\pstart\versepar \skt{अल्पदोषेषु }\edtext{\skt{योज्या स्यु ये}}{\lemma{\skt{योज्या स्यु ये}}\Afootnote{K; \skt{योज्याः स्युर्ये}\ A; \skt{योज्याश्च ये}\ H}}\skt{ योगाः सम्यगीरिताः~॥}\pend

\pstart\versepar \ledsidenote{\passref{4.31.46}}\edtext{\skt{विव\-र्ज\-ये\-त्स्ने\-ह\-पा\-न\-म\-जीर्ण्णी}}{\lemma{\skt{विव\appbrk{}र्ज\appbrk{}ये\appbrk{}त्स्ने\appbrk{}ह\appbrk{}पा\appbrk{}न\appbrk{}म\appbrk{}जीर्ण्णी}}\Afootnote{A K; \skt{विव\appbrk{}र्ज्ज\appbrk{}ये\appbrk{}त्स्ने\appbrk{}ह\appbrk{}पा\appbrk{}न\appbrk{}म\appbrk{}जीर्ण्णी}\ H}}\skt{ तरुणज्वरी~। }\pend
\pstart\versepar \skt{}\edtext{\skt{दुर्वलारोचकी}}{\lemma{\skt{दुर्वलारोचकी}}\Afootnote{K H; \skt{दुर्बलोऽरोचकी}\ A}}\skt{ स्थूलो मूर्च्छार्त्तो मदपीडितः~।}\pend

\pstart\versepar \ledsidenote{\passref{4.31.47}}\edtext{\skt{छर्द्या\-भि\-भू\-त\-स्तृ\-षितः}}{\lemma{\skt{छर्द्या\appbrk{}भि\appbrk{}भू\appbrk{}त\appbrk{}स्तृ\appbrk{}षितः}}\Afootnote{K; \skt{छर्द्यर्दितः पिपासार्तः}\ A; \skt{छर्द्द्या\appbrk{}भि\appbrk{}भू\appbrk{}त\appbrk{}स्तृ\appbrk{}षितः}\ H}}\skt{ श्रान्तः पानक्लमान्वितः~। }\pend
\pstart\versepar \skt{}\edtext{\skt{वस्तिदत्तो}}{\lemma{\skt{वस्तिदत्तो}}\Afootnote{K H; \skt{दत्तबस्तिर्}\ A}}\skt{ }\edtext{\skt{विरिक्तश्च}}{\lemma{\skt{विरिक्तश्च}}\Afootnote{A K; \skt{विरेकाश्च}\ H}}\skt{ }\edtext{\skt{वान्तो यश्चापि}}{\lemma{\skt{वान्तो यश्}}\Afootnote{A K; \skt{वान्यो यश्}\ H}}\skt{ }\edtext{\skt{मानवः}}{\lemma{\skt{मानवः}}\Afootnote{A K; \skt{मानवाः~॥}\ H}}\skt{~।}\pend

\pstart \ledsidenote{\passref{4.31.48cd}}\skt{अकाले च प्रसूता स्त्री स्नेहपानं विवर्जयेत्~।}\pend

\pstart\versepar \ledsidenote{\passref{4.31.49}}\edtext{\skt{स्नेह\-पा\-ना\-द्भ\-वे\-द्नॄणां}}{\lemma{\skt{स्नेह\appbrk{}पा\appbrk{}ना\appbrk{}द्भ\appbrk{}वे\appbrk{}द्नॄणां}}\Afootnote{K; \skt{स्नेह\appbrk{}पा\appbrk{}ना\appbrk{}द्भ\appbrk{}व\appbrk{}न्त्येषां नॄणां}\ A H}}\skt{ नानाविधा गदाः~। }\pend
\pstart\versepar \skt{गदा वा कृच्छ्रतां यान्ति न }\edtext{\skt{सिध्यन्त्यत्र वा}}{\lemma{\skt{सिध्यन्त्यत्र वा}}\Afootnote{K H; \skt{सिध्यन्त्यथवा}\ A}}\skt{ }\edtext{\skt{यतः}}{\lemma{\skt{यतः}}\Afootnote{K H; \skt{पुनः॥}\ A}}\skt{~॥}\pend

\pstart \ledsidenote{\passref{4.31.51cd}}\edtext{\skt{रूक्षं}}{\lemma{\skt{रूक्षं}}\Afootnote{K H; \emph{om.}\ A}}\skt{ }\edtext{\skt{पुरीशं}}{\lemma{\skt{पुरीशं}}\Afootnote{K; \skt{पुरीषं}\ A H}}\skt{ ग्रथितं }\edtext{\skt{भुक्तं}}{\lemma{\skt{भुक्तं}}\Afootnote{K H; \skt{रूक्षं}\ A}}\skt{ }\edtext{\skt{कृच्छ्रेण}}{\lemma{\skt{कृच्छ्रेण}}\Afootnote{K H; \skt{कृच्छ्रादन्नं}\ A}}\skt{ }\edtext{\skt{पच्यते}}{\lemma{\skt{पच्यते}}\Afootnote{K H; \skt{विपच्यते॥}\ A}}\skt{~।}\pend

\pstart \ledsidenote{\passref{4.31.52cd}}\skt{दुर्वर्ण्णो }\edtext{\skt{दुर्बलश्चैव}}{\lemma{\skt{दुर्बलश्}}\Afootnote{A; \skt{दुर्वलश्}\ K H}}\skt{ रूक्षो भवति मानवः~॥}\pend

\pstart\versepar \ledsidenote{\passref{4.31.53}}\edtext{\skt{ग्लानिः}}{\lemma{\skt{ग्लानिः}}\Afootnote{K H; \skt{सुस्निग्धा त्वग्विट्शैथिल्यं दीप्तोऽ\appbrk{}ग्नि\appbrk{}र्मृ\appbrk{}दु\appbrk{}गा\appbrk{}त्रता। ग्लानि}\ A}}\skt{ }\edtext{\edtext{\skt{सद\-न\-म\-ङ्गा\-ना\-म\-ध\-स्ता\-त्स्ने\-ह\-द\-र्शनं}}{\lemma{\skt{सद\appbrk{}न\appbrk{}म\appbrk{}ङ्गा\appbrk{}नाम्}}\Afootnote{K; \skt{र्लाघवमङ्गाना}\ A; \skt{सदनमङ्गाना}\ H}}}{\lemma{\skt{अध\appbrk{}स्ता\appbrk{}त्स्ने\appbrk{}ह\appbrk{}द\appbrk{}र्शनं}}\Afootnote{K; \skt{मध\appbrk{}स्ता\appbrk{}त्स्ने\appbrk{}ह\appbrk{}द\appbrk{}र्श\appbrk{}नम्}\ A H}}\skt{~। }\pend
\pstart\versepar \skt{सम्यक्स्निग्धस्य लिङ्गानि स्नेहो\-द्वे\-ग\-स्त\-थैव च~॥}\pend

\pstart\versepar \ledsidenote{\passref{4.31.54}}\skt{भक्तद्वेषो }\edtext{\skt{मुखास्रावो}}{\lemma{\skt{मुखास्रावो}}\Afootnote{K H; \skt{मुखस्रावो}\ A}}\skt{ }\edtext{\skt{गुददाह}}{\lemma{\skt{गुददाह}}\Afootnote{K H; \skt{गुददाहः}\ A}}\skt{ }\edtext{\skt{प्रवाहिका}}{\lemma{\skt{प्रवाहिका}}\Afootnote{A K; \skt{प्रवाहिकाः~।}\ H}}\skt{~। }\pend
\pstart\versepar \skt{}\edtext{\skt{पुरी\-शा\-ति\-प्र\-वृ\-त्तिश्च}}{\lemma{\skt{पुरी\appbrk{}शा\appbrk{}ति\appbrk{}प्र\appbrk{}वृ\appbrk{}त्तिश्च}}\Afootnote{K; \skt{पुरी\appbrk{}षा\appbrk{}ति\appbrk{}प्र\appbrk{}वृ\appbrk{}त्तिश्च}\ A H}}\skt{ }\edtext{\skt{तद\-ति\-स्नि\-ग्ध\-ल\-क्षणं}}{\lemma{\skt{तद\appbrk{}ति\appbrk{}स्नि\appbrk{}ग्ध\appbrk{}ल\appbrk{}क्षणं}}\Afootnote{K H; \skt{भृशस्निग्धस्य लक्षणम्॥}\ A}}\skt{~॥}\pend

\pstart\versepar \ledsidenote{\passref{4.31.55}}\skt{रूक्षस्य स्नेहनं स्नेहै\-र\-ति\-स्नि\-ग्धस्य रूक्षणं~। }\pend
\pstart\versepar \skt{श्यामा\-क\-को\-र\-दू\-षा\-न्न\-त\-क्र\-पि\-ण्या\-क\-श\-क्तुभिः~॥}\pend

\pstart\versepar \ledsidenote{\passref{4.31.56}}\skt{दीप्तान्तरग्निः परिशुद्धकोष्ठः }\edtext{\skt{प्रत्य\-ग्र\-धा\-तु\-र्ब\-ल\-व\-र्ण्ण\-युक्तः}}{\lemma{\skt{प्रत्य\appbrk{}ग्र\appbrk{}धा\appbrk{}तु\appbrk{}र्ब\appbrk{}ल\appbrk{}व\appbrk{}र्ण्ण\appbrk{}युक्तः}}\Afootnote{A K; \skt{प्रत्य\appbrk{}ग्र\appbrk{}धा\appbrk{}तु\appbrk{}र्व्व\appbrk{}ल\appbrk{}व\appbrk{}र्ण्ण\appbrk{}युक्तः~।}\ H}}\skt{~। }\pend
\pstart\versepar \skt{दृढेन्द्रियो }\edtext{\skt{मन्दरजाः}}{\lemma{\skt{मन्दरजाः}}\Afootnote{K; \skt{मन्दजरः}\ A; \skt{मन्दरुजः}\ H}}\skt{ शतायुः स्नेहोपसेवी पुरुषो }\edtext{\skt{भवेधि}}{\lemma{\skt{भवेधि}}\Afootnote{K H; \skt{भवेत्तु॥}\ A}}\skt{~।}\pend

\pstart\versepar \ledsidenote{\passref{4.31.57}}\skt{स्नेहो हितो }\edtext{\skt{दुर्ब\-ल\-व\-ह्नि\-दे\-हे\-स\-न्धु\-क्षणे}}{\lemma{\skt{दुर्ब\appbrk{}ल\appbrk{}व\appbrk{}ह्नि\appbrk{}दे\appbrk{}हे\appbrk{}स\appbrk{}न्धु\appbrk{}क्षणे}}\Afootnote{\skt{दुर्ब\appbrk{}ल\appbrk{}ह्नि\appbrk{}दे\appbrk{}ह\appbrk{}स\appbrk{}न्धु\appbrk{}क्षणे}\ A; \skt{दुर्व\appbrk{}ल\appbrk{}व\appbrk{}ह्नि\appbrk{}दे\appbrk{}हे\appbrk{}स\appbrk{}न्धु\appbrk{}क्षणे}\ K H}}\skt{ व्याधि\-नि\-पी\-डि\-तस्य~। }\pend
\pstart\versepar \skt{}\edtext{\skt{वलान्वितो}}{\lemma{\skt{वलान्वितो}}\Afootnote{K; \skt{बलान्वितौ}\ A; \skt{वलान्विन्वितो}\ H}}\skt{ भोज\-न\-दो\-ष\-जातैः प्रमर्दितुं तौ सहसा न }\edtext{\skt{शक्याविति}}{\lemma{\skt{शक्याविति}}\Afootnote{K; \skt{साध्यौ॥}\ A; \skt{शक्यौ}\ H}}\skt{~॥}\pend

\pstart \ledsidenote{\passref{4.31.trailer}}\edtext{\skt{ल}}{\lemma{\skt{ल}}\Afootnote{K; \skt{इति सुश्रुत}\ A; \skt{इति~॥~॥}\ H}}\skt{ }\edtext{\skt{ल..{\symfont ❈}॥}}{\lemma{\skt{ल..{\symfont ❈}}}\Afootnote{K; \skt{संहितायां चिकित्सास्थाने स्नेहो\appbrk{}प\appbrk{}यौ\appbrk{}गि\appbrk{}क\appbrk{}चि\appbrk{}कि\appbrk{}त्सितं नामै\appbrk{}क\appbrk{}त्रिं\appbrk{}श\appbrk{}त्त\appbrk{}मोऽ\appbrk{}ध्यायः~॥३१॥}\ A; \skt{चिकित्सा एकविंशत्~॥॥}\ H}}\pend

\endnumbering
\beginnumbering
\eledsection*{Adhyāya 32}
\pstart \ledsidenote{\passref{4.32.1}}\skt{अथातः स्वेदा\-व\-चा\-र\-णीयं }\edtext{\skt{चिकित्सितं}}{\lemma{\skt{चिकित्सितं}}\Afootnote{A K; \emph{om.}\ H}}\skt{ }\edtext{\skt{व्याख्यास्यामः॥}}{\lemma{\skt{व्याख्यास्यामः}}\Afootnote{A H; \emph{om.}\ K}}\pend

\pstart \ledsidenote{\passref{4.32.3}}\edtext{\skt{तत्र}}{\lemma{\skt{तत्र}}\Afootnote{K H; \emph{om.}\ A}}\skt{ }\edtext{\skt{चतुर्विधः}}{\lemma{\skt{चतुर्विधः}}\Afootnote{A K; \skt{चतुर्व्विधाः}\ H}}\skt{ }\edtext{\skt{स्वेदा}}{\lemma{\skt{स्वेदा}}\Afootnote{K H; \skt{स्वे}\ A}}\skt{ }\edtext{\skt{भवति}}{\lemma{\skt{भवति}}\Afootnote{K; \skt{दः}\ A; \skt{भवन्ति~॥}\ H}}\skt{~।. तद्यथा }\edtext{\skt{तापस्वेदो}}{\lemma{\skt{तापस्वेदो}}\Afootnote{K H; \skt{तापस्वेद,}\ A}}\skt{ }\edtext{\skt{बाष्पस्वेदः~।उ\-प\-ना\-ह\-स्वेदो}}{\lemma{\skt{बाष्पस्वेदः~।उ\appbrk{}प\appbrk{}ना\appbrk{}ह\appbrk{}स्वेदो}}\Afootnote{K H; \skt{ऊष्मस्वेद, उपनाहस्वेदो,}\ A}}\skt{ }\edtext{\edtext{\skt{द्रवस्वेदश्चेति}}{\lemma{\skt{द्रवस्वेदश्}}\Afootnote{K H; \skt{द्रवस्वेद}\ A}}}{\lemma{\skt{चेति}}\Afootnote{K H; \skt{इति}\ A}}\skt{~। }\edtext{\skt{तत्र}}{\lemma{\skt{तत्र}}\Afootnote{K H; \skt{अत्र}\ A}}\skt{ }\edtext{\skt{स्वेद\-वि\-क\-ल्पा\-व\-रोधः}}{\lemma{\skt{स्वेद\appbrk{}वि\appbrk{}क\appbrk{}ल्पा\appbrk{}व\appbrk{}रोधः}}\Afootnote{K H; \skt{सर्व\appbrk{}स्वे\appbrk{}द\appbrk{}वि\appbrk{}क\appbrk{}ल्पा\appbrk{}व\appbrk{}रोधः~॥}\ A}}\skt{~।}\pend

\pstart \ledsidenote{\passref{4.32.4}}\skt{तत्र तापस्वेदो }\edtext{\skt{नामपालिकां}}{\lemma{\skt{नामपालिकां}}\Afootnote{K H; \skt{पाणिकां}\ A}}\skt{ }\edtext{\skt{सक\-न्द\-फ\-ल\-वां\-लि\-का\-वस्त्रैः}}{\lemma{\skt{सक\appbrk{}न्द\appbrk{}फ\appbrk{}ल\appbrk{}वां\appbrk{}लि\appbrk{}का\appbrk{}वस्त्रैः}}\Afootnote{K; \skt{स्यक\appbrk{}न्दु\appbrk{}क\appbrk{}क\appbrk{}पा\appbrk{}ल\appbrk{}वा\appbrk{}लु\appbrk{}का\appbrk{}वस्त्रैः}\ A; \skt{सक\appbrk{}न्द\appbrk{}फ\appbrk{}ल\appbrk{}वा\appbrk{}लि\appbrk{}का\appbrk{}वस्त्रैः}\ H}}\skt{ प्रयुज्यते शयानस्य }\edtext{\skt{बाष्पतापोङ्गारैः}}{\lemma{\skt{बाष्पतापोङ्गारैः}}\Afootnote{K H; \skt{चाङ्गतापो बहुशः खादि\appbrk{}रा\appbrk{}ङ्गा\appbrk{}रै\appbrk{}रिति~॥}\ A}}\skt{~॥}\pend

\pstart \ledsidenote{\passref{4.32.5}}\edtext{\skt{बाष्पस्वेदस्तु}}{\lemma{\skt{बाष्पस्वेदस्तु}}\Afootnote{K H; \skt{ऊष्मस्वेदस्तु}\ A}}\skt{ }\edtext{\edtext{\skt{कपा\-ल\-पा\-षा\-णो\-ष्ट\-क\-लो\-ह\-पि\-ण्डा\-न\-ग्नि\-व\-र्णा\-न\-द्भि\-रा\-सि\-च्य\-तै\-रा\-र्द्र\-व\-स्त्र\-प\-रि\-वे\-ष्टितैः}}{\lemma{\skt{कपा\appbrk{}ल\appbrk{}पा\appbrk{}षा\appbrk{}णो\appbrk{}ष्ट\appbrk{}क\appbrk{}लो\appbrk{}ह\appbrk{}पि\appbrk{}ण्डा\appbrk{}न\appbrk{}ग्नि\appbrk{}व\appbrk{}र्णा\appbrk{}न\appbrk{}द्भिर्}}\Afootnote{K; \skt{कपा\appbrk{}ल\appbrk{}पा\appbrk{}षा\appbrk{}शे\appbrk{}ष्ट\appbrk{}का\appbrk{}लो\appbrk{}ह\appbrk{}पि\appbrk{}ण्डा\appbrk{}न\appbrk{}ग्नि\appbrk{}व\appbrk{}र्णा\appbrk{}न\appbrk{}द्भिर्}\ A; \skt{कपा\appbrk{}ल\appbrk{}पा\appbrk{}षा\appbrk{}णे\appbrk{}ष्ट\appbrk{}क\appbrk{}लो\appbrk{}ह\appbrk{}पि\appbrk{}ण्डा\appbrk{}न\appbrk{}ग्नि\appbrk{}व\appbrk{}र्ण्णा\appbrk{}न\appbrk{}द्भिर्}\ H}}}{\lemma{\skt{आसि\appbrk{}च्य\appbrk{}तै\appbrk{}रा\appbrk{}र्द्र\appbrk{}व\appbrk{}स्त्र\appbrk{}प\appbrk{}रि\appbrk{}वे\appbrk{}ष्टितैः}}\Afootnote{K; \skt{आसि\appbrk{}ञ्चे\appbrk{}द\appbrk{}म्ल\appbrk{}द्र\appbrk{}व्यैर्वा, तैरा\appbrk{}र्द्रा\appbrk{}ल\appbrk{}क्त\appbrk{}क\appbrk{}प\appbrk{}रि\appbrk{}वे\appbrk{}ष्टि\appbrk{}तैर्}\ A; \skt{आसिच्यते राद्र\appbrk{}व\appbrk{}स्त्र\appbrk{}प\appbrk{}रि\appbrk{}वे\appbrk{}ष्टितैः~।}\ H}}\skt{~। }\edtext{\skt{स्वेद\-ये\-द\-ङ्ग\-प्र\-देशा}}{\lemma{\skt{स्वेद\appbrk{}ये\appbrk{}द\appbrk{}ङ्ग\appbrk{}प्र\appbrk{}देशा}}\Afootnote{K H; \skt{अङ्गप्रदेशं}\ A}}\skt{ }\edtext{\skt{एवं}}{\lemma{\skt{एवं}}\Afootnote{K; \skt{स्वेदये}\ A; \skt{न्~। एवम्}\ H}}\skt{ }\edtext{\skt{मांस\-र\-स\-प\-यो\-द\-धि\-स्ने\-ह\-धा\-न्या\-म्व्ल\-वा\-त\-ह\-र\-प\-त्र\-भं\-ग\-क्वा\-थ\-पू\-र्ण्णम्वा}}{\lemma{\skt{मांस\appbrk{}र\appbrk{}स\appbrk{}प\appbrk{}यो\appbrk{}द\appbrk{}धि\appbrk{}स्ने\appbrk{}ह\appbrk{}धा\appbrk{}न्या\appbrk{}म्व्ल\appbrk{}वा\appbrk{}त\appbrk{}ह\appbrk{}र\appbrk{}प\appbrk{}त्र\appbrk{}भं\appbrk{}ग\appbrk{}क्वा\appbrk{}थ\appbrk{}पू\appbrk{}र्ण्णम्वा}}\Afootnote{K; \skt{त्~। मांस\appbrk{}र\appbrk{}स\appbrk{}प\appbrk{}यो\appbrk{}द\appbrk{}धि\appbrk{}स्ने\appbrk{}ह\appbrk{}धा\appbrk{}न्या\appbrk{}म्ल\appbrk{}वा\appbrk{}त\appbrk{}ह\appbrk{}र\appbrk{}प\appbrk{}त्र\appbrk{}भ\appbrk{}ङ्ग\appbrk{}क्वा\appbrk{}थ\appbrk{}पूर्णां वा}\ A; \skt{मांस\appbrk{}र\appbrk{}स\appbrk{}प\appbrk{}यो\appbrk{}द\appbrk{}धि\appbrk{}स्ने\appbrk{}ह\appbrk{}धा\appbrk{}न्या\appbrk{}म्ल\appbrk{}वा\appbrk{}त\appbrk{}ह\appbrk{}र\appbrk{}प\appbrk{}त्त्र\appbrk{}भ\appbrk{}ङ्ग\appbrk{}क्वा\appbrk{}पू\appbrk{}र्ण्णम्वा}\ H}}\skt{ }\edtext{\edtext{\skt{कुम्भी\-म\-भि\-प्र\-तप्तां}}{\lemma{\skt{कुम्भीम्}}\Afootnote{K H; \skt{कुम्भी}\ A}}}{\lemma{\skt{अभिप्रतप्तां}}\Afootnote{K H; \skt{मनुतप्तां}\ A}}\skt{ प्रावृत्योष्माणं गृह्णीयात्~। पार्श्वच्छिद्रेण वा }\edtext{\skt{कुम्भेन}}{\lemma{\skt{कुम्भेन}}\Afootnote{K H; \skt{कुम्भेना}\ A}}\skt{ }\edtext{\skt{तस्याः}}{\lemma{\skt{तस्याः}}\Afootnote{K H; \skt{धो}\ A}}\skt{ }\edtext{\skt{कुम्भ्यामुखेन}}{\lemma{\skt{कुम्भ्यामुखेन}}\Afootnote{K H; \skt{मुखेन}\ A}}\skt{ }\edtext{\edtext{\skt{मुख\-म\-भि\-स\-न्धाय}}{\lemma{\skt{मुखमभिस्}}\Afootnote{K; \skt{तस्या मुखमभिस्}\ A; \skt{मुखमभि}\ H}}}{\lemma{\skt{अन्धाय}}\Afootnote{A K; \skt{सन्धाय}\ H}}\skt{ }\edtext{\skt{तस्मि\-न\-ञ्छि\-द्रे\-ह\-स्ति\-शु\-ण्डा\-कारं}}{\lemma{\skt{तस्मि\appbrk{}न\appbrk{}ञ्छि\appbrk{}द्रे\appbrk{}ह\appbrk{}स्ति\appbrk{}शु\appbrk{}ण्डा\appbrk{}कारं}}\Afootnote{K; \skt{तस्मि\appbrk{}ञ्छि\appbrk{}द्रे\appbrk{}ह\appbrk{}स्ति\appbrk{}शु\appbrk{}ण्डा\appbrk{}कारां}\ A; \skt{तस्मिञ्छिद्रे हस्तिशुण्डाकारं}\ H}}\skt{ नाडीं प्रणिधाय }\edtext{\skt{शयानं}}{\lemma{\skt{शयानं}}\Afootnote{K H; \skt{तं}\ A}}\skt{ स्वेदयेत्}\pend

\pstart\versepar \ledsidenote{\passref{4.32.6}}\skt{सुखोपविष्टं स्वभ्यक्तं }\edtext{\skt{प्रावृतम्वा~।. भ}}{\lemma{\skt{प्रावृतम्वा~।. भ}}\Afootnote{K H; \skt{गुरु\appbrk{}प्रा\appbrk{}व\appbrk{}र\appbrk{}णा\appbrk{}वृ\appbrk{}तम्~।}\ A}}\skt{~॥ }\pend
\pstart\versepar \skt{हस्तिशुण्डिकया नाड्या }\edtext{\skt{स्वेद\-ये\-द्वा\-त\-रो\-गिणं}}{\lemma{\skt{स्वेद\appbrk{}ये\appbrk{}द्वा\appbrk{}त\appbrk{}रो\appbrk{}गिणं}}\Afootnote{A K; \skt{स्वेद\appbrk{}ये\appbrk{}द्वा\appbrk{}त\appbrk{}रो\appbrk{}गिनं~।}\ H}}\skt{~। }\pend
\pstart\versepar \skt{}\edtext{\skt{सुखं}}{\lemma{\skt{सुखं}}\Afootnote{K H; \skt{सुखा}\ A}}\skt{ }\edtext{\skt{सर्वानुगा}}{\lemma{\skt{सर्वानुगा}}\Afootnote{K H; \skt{सर्वाङ्गगा}\ A}}\skt{ ह्येषा न च }\edtext{\edtext{\skt{संक्लेषयेन्नरं}}{\lemma{\skt{संक्लेषयेन्}}\Afootnote{K; \skt{क्लिश्नाति मा}\ A; \skt{संङ्क्लेशयेन्}\ H}}}{\lemma{\skt{नरं}}\Afootnote{K H\textsuperscript{pc}; \skt{नवम्~॥}\ A; \skt{नं~॥}\ H\textsuperscript{ac}}}\skt{~॥}\pend

\pstart\versepar \ledsidenote{\passref{4.32.7}}\edtext{\skt{व्यामा\-र्ध\-मा\-त्र\-त्रि\-व\-क्र\-ह\-स्ति\-ह\-स्त\-स\-मा\-कृतिः}}{\lemma{\skt{व्यामा\appbrk{}र्ध\appbrk{}मा\appbrk{}त्र\appbrk{}त्रि\appbrk{}व\appbrk{}क्र\appbrk{}ह\appbrk{}स्ति\appbrk{}ह\appbrk{}स्त\appbrk{}स\appbrk{}मा\appbrk{}कृतिः}}\Afootnote{K H; \skt{व्यामार्धमात्रा त्रिर्वक्रा हस्ति\appbrk{}ह\appbrk{}स्त\appbrk{}स\appbrk{}मा\appbrk{}कृतिः~।}\ A}}\skt{~। }\pend
\pstart\versepar \skt{स्वेदनार्थे हिता नाडी कैलिञ्जी हस्तिशुण्डिका~॥}\pend

\pstart \ledsidenote{\passref{4.32.8ab}}\edtext{\skt{पुरु\-षा\-या\-म\-मा\-त्रम्वा}}{\lemma{\skt{पुरु\appbrk{}षा\appbrk{}या\appbrk{}म\appbrk{}मा\appbrk{}त्रम्वा}}\Afootnote{K; \skt{पुरु\appbrk{}षा\appbrk{}या\appbrk{}म\appbrk{}मात्रां च}\ A; \skt{पुरु\appbrk{}षा\appbrk{}या\appbrk{}म\appbrk{}मा\appbrk{}त्राम्वा}\ H}}\skt{ }\edtext{\edtext{\skt{भूमिमुत्कीर्य}}{\lemma{\skt{भूमिम्}}\Afootnote{K H; \skt{भूमि}\ A}}}{\lemma{\skt{उत्कीर्य}}\Afootnote{K H; \skt{मुत्कीर्य}\ A}}\skt{ खादिरैः}\pend

\pstart \ledsidenote{\passref{4.32.8cd}}\edtext{\edtext{\skt{काष्टै\-र\-व\-दा\-ह्या\-भ्युक्ष्य}}{\lemma{\skt{काष्टैर्}}\Afootnote{K H; \skt{काष्ठैर्}\ A}}}{\lemma{\skt{अवदाह्याभ्युक्ष्य}}\Afootnote{K; \skt{दग्ध्वा तथाऽभ्युक्ष्य}\ A; \skt{अवदाह्यभ्युक्ष्य}\ H}}\skt{ }\edtext{\skt{वातहर}}{\lemma{\skt{वातहर}}\Afootnote{K H; \skt{क्षीर\appbrk{}धा\appbrk{}न्या\appbrk{}म्ल\appbrk{}वा\appbrk{}रिभिः~॥}\ A}}\pend

\pstart\versepar \ledsidenote{\passref{4.32.9}}\edtext{\skt{पत्र\-भ\-ङ्गै\-र\-व\-च्छाद्य}}{\lemma{\skt{पत्र\appbrk{}भ\appbrk{}ङ्गै\appbrk{}र\appbrk{}व\appbrk{}च्छाद्य}}\Afootnote{K H; \skt{पत्र\appbrk{}भ\appbrk{}ङ्गै\appbrk{}र\appbrk{}व\appbrk{}च्छाद्य}\ A}}\skt{ शयानं स्वेदयेत्~। }\pend
\pstart\versepar \skt{}\edtext{\skt{शिला\-त\-ल\-म\-व\-दा\-ह्या\-पोह्य भस्मं च }\edtext{\skt{पूर्ववत्}}{\lemma{\skt{पूर्ववत्}}\Afootnote{K H; \skt{पूर्व\appbrk{}व\appbrk{}त्स्वे\appbrk{}द\appbrk{}ये\appbrk{}द्दग्ध्वा भस्मापोह्यापि वा शिलाम्~॥}\ A}}}{\lemma{\skt{त्… भस्मं}}\Afootnote{\emph{om.}\ A}}\skt{~।}\edtext{}{\lemma{\skt{स्वेदयेत्~। शिला\appbrk{}त\appbrk{}ल\appbrk{}म\appbrk{}व\appbrk{}दा\appbrk{}ह्या\appbrk{}पोह्य}}\Afootnote{K H; \skt{स्वेदये}\ A}}\pend

\pstart \ledsidenote{\passref{4.32.10}}\edtext{\skt{कुटीम्वा}}{\lemma{\skt{कुटीम्वा}}\Afootnote{K H; \skt{पूर्ववत्कुटीं वा}\ A}}\skt{ चतुर्द्वारां कृत्वा }\edtext{\skt{तस्या\-मु\-प\-वि\-ष्ट\-स्या\-न्त\-श्च\-तु\-र्वां\-गा\-रा\-न\-नु\-प\-रि\-धाय}}{\lemma{\skt{चतु\appbrk{}र्वां\appbrk{}गा\appbrk{}रा\appbrk{}न\appbrk{}नु\appbrk{}प\appbrk{}रि\appbrk{}धाय}}\Afootnote{K; \skt{चतु\appbrk{}र्द्वा\appbrk{}रेऽ\appbrk{}ङ्गा\appbrk{}रा\appbrk{}नु\appbrk{}प\appbrk{}स\appbrk{}न्धाय तं}\ A; \skt{चतु\appbrk{}र्द्वा\appbrk{}रा\appbrk{}न\appbrk{}नु\appbrk{}प\appbrk{}रि\appbrk{}धाय}\ H}}\skt{ स्वेदयेत्~।}\pend

\pstart \ledsidenote{\passref{4.32.11}}\edtext{\skt{काशधान्यानि वा}}{\lemma{\skt{काशधान्यानि वा}}\Afootnote{K H; \skt{कोशधान्यानि वा}\ A}}\skt{ }\edtext{\skt{सम्य\-गु\-स्वे\-द्या\-स्तीर्य}}{\lemma{\skt{सम्य\appbrk{}गु\appbrk{}स्वे\appbrk{}द्या\appbrk{}स्तीर्य}}\Afootnote{K H; \skt{सम्य\appbrk{}गु\appbrk{}प\appbrk{}स्वे\appbrk{}द्या\appbrk{}स्तीर्य}\ A}}\skt{ }\edtext{\skt{किलिञ्जेन्यस्मिं वा}}{\lemma{\skt{किलिञ्जेन्यस्मिं वा}}\Afootnote{K; \skt{किलिञ्जेऽन्यस्मिन्वा}\ A; \skt{किलिञ्जेन्यस्मिन्वा}\ H}}\skt{ तत्प्रतिरूपके शयानं }\edtext{\skt{स्वेदयेदेवं}}{\lemma{\skt{स्वेदयेदेवं}}\Afootnote{K H; \skt{प्रावृत्य स्वेदयेत्; एवं}\ A}}\skt{ पांशु }\edtext{\skt{गोशकृत्तुषप}}{\lemma{\skt{गोशकृत्तुषप}}\Afootnote{K; \skt{गोश\appbrk{}कृ\appbrk{}त्तु\appbrk{}ष\appbrk{}बु\appbrk{}सप}\ A; \skt{गोश\appbrk{}कृ\appbrk{}त्तु\appbrk{}ष\appbrk{}वु\appbrk{}षप}\ H}}\skt{ }\edtext{\skt{लालोष्मभिः}}{\lemma{\skt{लालोष्मभिः}}\Afootnote{K; \skt{लालोष्मभिः}\ A; \skt{लालोस्मभिः}\ H}}\skt{ }\edtext{\skt{स्वेदयेदिति}}{\lemma{\skt{स्वेदयेदिति}}\Afootnote{K H; \skt{स्वेदयेत्~॥}\ A}}\skt{~॥}\pend

\pstart \ledsidenote{\passref{4.32.12}}\skt{उप\-ना\-ह\-स्वे\-दस्तु }\edtext{\edtext{\skt{वात\-ह\-र\-द्र\-व्य\-क\-ल्कै\-र\-म्व्ल\-पि\-ष्टै\-र्ल\-व\-ण\-प्र\-गाढैः}}{\lemma{\skt{वात\appbrk{}ह\appbrk{}र\appbrk{}द्र\appbrk{}व्य\appbrk{}क\appbrk{}ल्कैर्}}\Afootnote{K; \skt{वात\appbrk{}ह\appbrk{}र\appbrk{}मू\appbrk{}ल\appbrk{}क\appbrk{}ल्कैर्}\ A; \skt{वात\appbrk{}ह\appbrk{}र\appbrk{}द्र\appbrk{}व्य\appbrk{}कल्कै}\ H}}}{\lemma{\skt{अम्व्लपिष्टैर्}}\Afootnote{K; \skt{अम्लपिष्टैर्}\ A; \skt{रम्लपिष्टैर्}\ H}}\skt{ सुस्निग्धैः }\edtext{\skt{सुखोष्णैश्च}}{\lemma{\skt{सुखोष्णैश्च}}\Afootnote{K H; \skt{सुखोष्णैः}\ A}}\skt{ प्रदिह्य स्वेदयेदेवं }\edtext{\skt{काको\-ल्या\-दि\-रे\-लासु}}{\lemma{\skt{काको\appbrk{}ल्या\appbrk{}दि\appbrk{}रे\appbrk{}लासु}}\Afootnote{K H; \skt{काको\appbrk{}ल्या\appbrk{}दि\appbrk{}भि\appbrk{}रे\appbrk{}ला\appbrk{}दिभिः सु}\ A}}\skt{ }\edtext{\edtext{\skt{रसादिमाहृत्य}}{\lemma{\skt{रसादिम्}}\Afootnote{K H; \skt{रसादि}\ A}}}{\lemma{\skt{आहृत्य}}\Afootnote{K H; \skt{भिस्}\ A}}\skt{ }\edtext{\skt{तिला\-त\-सी\-स\-र्ष\-प\-क\-ल्कैर्वा}}{\lemma{\skt{तिला\appbrk{}त\appbrk{}सी\appbrk{}स\appbrk{}र्ष\appbrk{}प\appbrk{}क\appbrk{}ल्कैर्वा}}\Afootnote{K H; \skt{तिला\appbrk{}त\appbrk{}सी\appbrk{}स\appbrk{}र्ष\appbrk{}प\appbrk{}कल्कैः}\ A}}\skt{ }\edtext{\skt{कृस\-र\-वे\-श\-वा\-र\-पा\-य\-सो\-त्का\-रि\-का\-भिर्वा}}{\lemma{\skt{कृस\appbrk{}र\appbrk{}वे\appbrk{}श\appbrk{}वा\appbrk{}र\appbrk{}पा\appbrk{}य\appbrk{}सो\appbrk{}त्का\appbrk{}रि\appbrk{}का\appbrk{}भिर्वा}}\Afootnote{K H; \skt{कृश\appbrk{}रा\appbrk{}पा\appbrk{}य\appbrk{}सो\appbrk{}त्का\appbrk{}रि\appbrk{}का\appbrk{}भिर्वे}\ A}}\skt{ }\edtext{\edtext{\skt{एताभिरेव}}{\lemma{\skt{एताभिर्}}\Afootnote{K H; \emph{om.}\ A}}}{\lemma{\skt{एव}}\Afootnote{K H; \skt{श}\ A}}\skt{ }\edtext{\skt{वायो}}{\lemma{\skt{वायो}}\Afootnote{K H; \skt{वारैः}\ A}}\skt{ }\edtext{\skt{दूलिकां}}{\lemma{\skt{दूलिकां}}\Afootnote{K H; \skt{साल्वणैर्वा}\ A}}\skt{ }\edtext{\skt{वध्वा}}{\lemma{\skt{वध्वा}}\Afootnote{K H; \skt{तनु\appbrk{}व\appbrk{}स्त्रा\appbrk{}व\appbrk{}नद्धैः}\ A}}\skt{ स्वेदयेत्~॥}\pend

\pstart \ledsidenote{\passref{4.32.13}}\skt{द्रवस्वेदस्तु }\edtext{\skt{वात\-ह\-र\-द्र\-व्य\-क्वा\-थ\-पू\-र्ण्ण\-को\-ष्ट\-क\-टाहे}}{\lemma{\skt{वात\appbrk{}ह\appbrk{}र\appbrk{}द्र\appbrk{}व्य\appbrk{}क्वा\appbrk{}थ\appbrk{}पू\appbrk{}र्ण्ण\appbrk{}को\appbrk{}ष्ट\appbrk{}क\appbrk{}टाहे}}\Afootnote{K H; \skt{वात\appbrk{}ह\appbrk{}र\appbrk{}द्र\appbrk{}व्य\appbrk{}क्वा\appbrk{}थ\appbrk{}पूर्णे कोष्णकटाहे}\ A}}\skt{ }\edtext{\skt{द्रोण्याम्वावगाह्य}}{\lemma{\skt{वावगाह्य}}\Afootnote{K H; \skt{वाऽवगाह्य}\ A}}\skt{ स्वेदयेत्~। एवं }\edtext{\skt{पयो\-मां\-स\-र\-स\-यू\-ष\-तै\-ल\-घृ\-त\-धा\-न्या\-म्व्ल\-मू\-त्रे\-ष्वे\-तै\-रेव}}{\lemma{\skt{पयो\appbrk{}मां\appbrk{}स\appbrk{}र\appbrk{}स\appbrk{}यू\appbrk{}ष\appbrk{}तै\appbrk{}ल\appbrk{}घृ\appbrk{}त\appbrk{}धा\appbrk{}न्या\appbrk{}म्व्ल\appbrk{}मू\appbrk{}त्रे\appbrk{}ष्वे\appbrk{}तैर्}}\Afootnote{K; \skt{पयो\appbrk{}मां\appbrk{}स\appbrk{}र\appbrk{}स\appbrk{}यू\appbrk{}ष\appbrk{}तै\appbrk{}ल\appbrk{}धा\appbrk{}न्या\appbrk{}म्ल\appbrk{}घृ\appbrk{}त\appbrk{}व\appbrk{}सा\appbrk{}मू\appbrk{}त्रे\appbrk{}ष्व\appbrk{}व\appbrk{}गा\appbrk{}हेत; एतैर्}\ A; \skt{पयो\appbrk{}मां\appbrk{}स\appbrk{}र\appbrk{}स\appbrk{}यू\appbrk{}ष\appbrk{}तै\appbrk{}ल\appbrk{}घृ\appbrk{}त\appbrk{}धा\appbrk{}न्या\appbrk{}म्ल\appbrk{}मू\appbrk{}त्रे\appbrk{}ष्वे\appbrk{}तैर्}\ H}}\skt{ सुखोष्णैः }\edtext{\skt{परिषिञ्चेदिति}}{\lemma{\skt{परिषिञ्चेदिति}}\Afootnote{K H; \skt{कषायैश्च परिषिञ्चेदिति~॥}\ A}}\skt{~॥}\pend

\pstart \ledsidenote{\passref{4.32.14}}\skt{तत्र }\edtext{\skt{तापोष्मास्वेदे}}{\lemma{\skt{तापोष्मास्वेदे}}\Afootnote{K; \skt{तापोष्मस्वेदौ}\ A H}}\skt{ विशेषतः श्लेष्मघ्नौ~। उपनाहस्वेदो वातघ्नः~। }\edtext{\skt{अन्यतरस्मिं}}{\lemma{\skt{अन्यतरस्मिं}}\Afootnote{K; \skt{अन्यतरस्मि}\ A H}}\skt{ }\edtext{\skt{पित्त\-सं\-सृ\-ष्ट\-द्र\-व\-स्वेद}}{\lemma{\skt{पित्त\appbrk{}सं\appbrk{}सृ\appbrk{}ष्ट\appbrk{}द्र\appbrk{}व\appbrk{}स्वेद}}\Afootnote{K; \skt{न्पित्तसंसृष्टे द्रवस्वेद}\ A H}}\skt{ इति~॥}\pend

\pstart\versepar \ledsidenote{\passref{4.32.16}}\edtext{\skt{चतुर्विधो }\edtext{\skt{योऽभिहितो}}{\lemma{\skt{ऽभिहितो}}\Afootnote{A; \skt{भिहितो}\ K H}}}{\lemma{\skt{चतुर्विधो यो}}\Afootnote{\skt{भवन्ति चात्र चतुर्विधो यो}\ A; \skt{चतुर्विधो यो}\ K H}}\skt{ }\edtext{\skt{द्वाधा}}{\lemma{\skt{द्वाधा}}\Afootnote{K; \skt{द्विधा}\ A H}}\skt{ स्वेदः प्रयुज्यते~। }\pend
\pstart\versepar \skt{}\edtext{\skt{स्वस्मिन्नेव}}{\lemma{\skt{स्वस्मिन्नेव}}\Afootnote{\skt{सर्वस्मिन्नेव}\ A H; \skt{स्वस्मिन्नेव}\ K}}\skt{ देहे तु देहस्यावयवे तथा~।}\pend

\pstart\versepar \ledsidenote{\passref{4.32.17}}\skt{येषान्नस्यं विधातव्यं }\edtext{\skt{बस्तिश्चापि हि}}{\lemma{\skt{चापि हि}}\Afootnote{K H\textsuperscript{pc}; \skt{चैव हि}\ A; \skt{चापि}\ H\textsuperscript{ac}}}\skt{ देहिनां~। }\pend
\pstart\versepar \skt{शोधनीयाश्च ये }\edtext{\skt{केचित्पूर्वं}}{\lemma{\skt{केचित्पूर्वं}}\Afootnote{A K; \skt{केचित्पूर्व्व}\ H}}\skt{ }\edtext{\skt{स्वेद्यास्तु मे}}{\lemma{\skt{स्वेद्यास्तु मे}}\Afootnote{K; \skt{स्वेद्यास्तु ते}\ A; \skt{स्वेदास्तु ते}\ H}}\skt{ }\edtext{\skt{मताः}}{\lemma{\skt{मताः}}\Afootnote{A K; \skt{मता~।}\ H}}\skt{~॥}\pend

\pstart\versepar \ledsidenote{\passref{4.32.18}}\edtext{\skt{पश्चात्स्वेद्याहृते}}{\lemma{\skt{पश्चात्स्वेद्याहृते}}\Afootnote{A K; \skt{पश्चात्स्वेदाहृते}\ H}}\skt{ शल्ये }\edtext{\skt{मूढ\-ग\-र्भा\-नु\-प\-द्रवा}}{\lemma{\skt{मूढ\appbrk{}ग\appbrk{}र्भा\appbrk{}नु\appbrk{}प\appbrk{}द्रवा}}\Afootnote{K; \skt{मूढ\appbrk{}ग\appbrk{}र्भाऽ\appbrk{}नु\appbrk{}प\appbrk{}द्रवा~।}\ A; \skt{मूढ\appbrk{}ग\appbrk{}र्व्भा\appbrk{}नु\appbrk{}प\appbrk{}द्रवा~॥}\ H}}\skt{~। }\pend
\pstart\versepar \skt{सम्यक्प्रजाता }\edtext{\skt{काले च }\edtext{\skt{पश्चात्स्वेद्यैव}}{\lemma{\skt{चात्स्वेद्यैव}}\Afootnote{K H; \skt{चात्स्वेद्या वि}\ A}}}{\lemma{\skt{काले च पश्}}\Afootnote{K H; \skt{काले या पश्}\ A}}\skt{ }\edtext{\skt{जानतां}}{\lemma{\skt{जानतां}}\Afootnote{K H; \skt{जानता~॥}\ A}}\skt{~॥}\pend

\pstart\versepar \ledsidenote{\passref{4.32.19}}\edtext{\skt{पूर्वां}}{\lemma{\skt{पूर्वां}}\Afootnote{K; \emph{om.}\ A; \skt{पूर्व्वा}\ H}}\skt{ }\edtext{\skt{स्वेद्या}}{\lemma{\skt{स्वेद्या}}\Afootnote{K; \skt{स्वेद्यः}\ A; \skt{स्वेद्याश्}\ H}}\skt{ }\edtext{\edtext{\skt{पश्चाच्च}}{\lemma{\skt{पश्}}\Afootnote{K; \skt{पूर्वं च पश्}\ A; \skt{च पश्}\ H}}}{\lemma{\skt{चाच्च}}\Afootnote{K H; \skt{चाच्च}\ A}}\skt{ }\edtext{\skt{भग\-न्द\-र्य\-र्श\-श\-स्तथा}}{\lemma{\skt{भग\appbrk{}न्द\appbrk{}र्य\appbrk{}र्श\appbrk{}शस्}}\Afootnote{K; \skt{भग\appbrk{}न्द\appbrk{}र्य\appbrk{}र्श\appbrk{}सस्}\ A H}}\skt{~। }\pend
\pstart\versepar \skt{}\edtext{\skt{अश्मर्यां}}{\lemma{\skt{अश्मर्यां}}\Afootnote{K H; \skt{अश्मर्या}\ A}}\skt{ चातुरो }\edtext{\skt{जन्तु}}{\lemma{\skt{जन्तु}}\Afootnote{K H; \skt{जन्तुः}\ A}}\skt{ }\edtext{\skt{शेषाञ्चास्त्रे}}{\lemma{\skt{शेषाञ्चास्त्रे}}\Afootnote{K H; \skt{शेषाञ्छास्त्रे}\ A}}\skt{ }\edtext{\skt{प्रवक्ष्यते}}{\lemma{\skt{प्रवक्ष्यते}}\Afootnote{K H; \skt{प्रचक्ष्महे~॥}\ A}}\skt{~॥}\pend

\pstart\versepar \ledsidenote{\passref{4.32.20}}\skt{नानभ्यक्ते नापि चा स्निग्धदेहे }\edtext{\skt{योज्यः}}{\lemma{\skt{योज्यः}}\Afootnote{K H; \emph{om.}\ A}}\skt{ }\edtext{\skt{स्वेद्यः}}{\lemma{\skt{स्वेद्यः}}\Afootnote{K H; \skt{स्वेदो}\ A}}\skt{ }\edtext{\skt{स्वेदविद्भिः}}{\lemma{\skt{स्वेदविद्भिः}}\Afootnote{K H\textsuperscript{ac}; \skt{योज्यः स्वेदविद्भिः}\ A; \skt{द स्वेदविद्भिः}\ H\textsuperscript{pc}}}\skt{ कथञ्चित्~। }\pend
\pstart\versepar \skt{दृष्टं लोके }\edtext{\edtext{\edtext{\skt{काष्ठ\-म\-स्नि\-ग्ध\-मा\-शु\-र्या\-या\-त्भङ्गं}}{\lemma{\skt{काष्ठम्}}\Afootnote{K H; \skt{काष्ठ}\ A}}}{\lemma{\skt{अस्निग्धमाशुर्}}\Afootnote{K; \skt{मस्निग्धमाशु}\ A; \skt{अस्निग्धमाशु}\ H}}}{\lemma{\skt{यायात्भङ्गं}}\Afootnote{K H; \skt{गच्छेद्भङ्गं}\ A}}\skt{ स्वेद\-यो\-गै\-र्गृ\-हीतं~॥}\pend

\pstart\versepar \ledsidenote{\passref{4.32.21}}\edtext{\skt{स्वेदक्लिन्ना}}{\lemma{\skt{स्वेदक्लिन्ना}}\Afootnote{K H\textsuperscript{pc}; \skt{स्नेहक्लिन्ना}\ A; \skt{स्वेदःक्लिन्ना}\ H\textsuperscript{ac}}}\skt{ }\edtext{\skt{धातुसंस्था तु}}{\lemma{\skt{धातुसंस्था तु}}\Afootnote{K H; \skt{धातुसंस्थाश्च}\ A}}\skt{ दोषाः स्वस्थानस्था ये च मार्गेषु लीनाः~। }\pend
\pstart\versepar \skt{सम्य\-क्स्वे\-दै\-र्यो\-जि\-तैस्ते द्रवत्वं प्राप्ताः कोष्ठं }\edtext{\skt{शोध\-नै\-र्या\-त्य\-शेषाः}}{\lemma{\skt{यात्यशेषाः}}\Afootnote{K H; \skt{यान्त्यशेषम्~॥}\ A}}\skt{~॥}\pend

\pstart\versepar \ledsidenote{\passref{4.32.22}}\edtext{\skt{अग्ने}}{\lemma{\skt{अग्ने}}\Afootnote{A K; \skt{अग्नेर्}\ H}}\skt{ }\edtext{\skt{दीप्तिन्मार्दवं}}{\lemma{\skt{दीप्तिन्मार्दवं}}\Afootnote{K; \skt{र्दीप्तिं मार्दवं}\ A; \skt{द्दीपिम्मार्द्दवं}\ H}}\skt{ त्वक्प्रसादं भक्तश्रद्धां स्रोतसां निर्मलत्वं~। }\pend
\pstart\versepar \skt{}\edtext{\edtext{\skt{कुर्या\-त्स्वे\-द\-स्त\-न्द्री\-निद्रे}}{\lemma{\skt{कुर्यात्स्वेदस्}}\Afootnote{A K; \skt{कुर्यात्स्वेद}\ H}}\skt{ च}}{\lemma{\skt{तन्द्रीनिद्रे च}}\Afootnote{K; \skt{हन्ति निद्रां स}\ A; \skt{स्तन्द्रिनिद्रे च}\ H}}\skt{ }\edtext{\skt{हन्यात्~। }\edtext{\skt{सन्धी\-म्क\-त्वां\-श्चे\-ष्ट\-ये\-दा\-शु\-युक्तः}}{\lemma{\skt{कत्वां\appbrk{}श्चे\appbrk{}ष्ट\appbrk{}ये\appbrk{}दा\appbrk{}शु\appbrk{}युक्तः}}\Afootnote{K H; \skt{न्स्तब्धां\appbrk{}श्चे\appbrk{}ष्ट\appbrk{}ये\appbrk{}दाशु युक्तः~॥}\ A}}}{\lemma{\skt{हन्यात्~। सन्धीम्}}\Afootnote{K; \skt{तन्द्रां सन्धी}\ A; \skt{हन्यात्~। सन्धी}\ H}}\skt{~।}\pend

\pstart\versepar \ledsidenote{\passref{4.32.23}}\skt{स्वेदास्रावो व्याधि\-हा\-नि\-र्ल\-घुत्वं शीतार्थित्वं मार्दवं चातुरस्य~। }\pend
\pstart\versepar \skt{सम्यक्स्विन्ने लक्षणं }\edtext{\skt{प्राहुरेतत्~। मिथ्यास्विन्ने}}{\lemma{\skt{प्राहुरेतत्~। मिथ्यास्विन्ने}}\Afootnote{K; \skt{प्राहु\appbrk{}रे\appbrk{}त\appbrk{}न्मि\appbrk{}थ्या\appbrk{}स्विन्ने}\ A H}}\skt{ }\edtext{\skt{विपरीतं}}{\lemma{\skt{विपरीतं}}\Afootnote{K H; \skt{व्यत्ययेनै}\ A}}\skt{ तदेव~॥}\pend

\pstart\versepar \ledsidenote{\passref{4.32.24}}\edtext{\skt{स्फोटोत्पत्तिः}}{\lemma{\skt{स्फोटोत्पत्तिः}}\Afootnote{K H; \skt{स्विन्नेऽत्यर्थं सन्धिपीडा विदाहः स्फोटोत्पत्तिः}\ A}}\skt{ पित्त\-र\-क्त\-प्र\-कोपो }\edtext{\skt{मदा\-मू\-र्च्छा\-भ्र\-म\-दाहो}}{\lemma{\skt{मदा\appbrk{}मू\appbrk{}र्च्छा\appbrk{}भ्र\appbrk{}म\appbrk{}दाहो}}\Afootnote{K H; \skt{मूर्च्छा भ्रान्}\ A}}\skt{ }\edtext{\skt{क्लमश्च}}{\lemma{\skt{क्लमश्च}}\Afootnote{K H; \emph{om.}\ A}}\skt{~। }\pend
\pstart\versepar \skt{}\edtext{\skt{अतिस्विन्ने}}{\lemma{\skt{अतिस्विन्ने}}\Afootnote{K H; \skt{तिर्दा}\ A}}\skt{ }\edtext{\skt{सन्धिपीडा च}}{\lemma{\skt{सन्धिपीडा च}}\Afootnote{K H; \skt{ह}\ A}}\skt{ }\edtext{\skt{तृट्च}}{\lemma{\skt{तृट्च}}\Afootnote{K H; \skt{तृष्णे}\ A}}\skt{ }\edtext{\skt{शीता}}{\lemma{\skt{शीता}}\Afootnote{K H\textsuperscript{ac}; \skt{क्लमश्च}\ A; \skt{तृष्णा शीता}\ H\textsuperscript{pc}}}\skt{ }\edtext{\edtext{\skt{क्रियास्तत्र}}{\lemma{\skt{क्रियास्}}\Afootnote{K H; \emph{om.}\ A}}}{\lemma{\skt{तत्र}}\Afootnote{K H; \emph{om.}\ A}}\skt{ }\edtext{\skt{कुर्याद्विधिज्ञः}}{\lemma{\skt{कुर्याद्विधिज्ञः}}\Afootnote{K; \skt{कुर्यात्तूर्णं तत्र शीतं विधानम्~॥}\ A; \skt{कुर्याद्विधिज्ञः~॥}\ H}}\skt{~॥}\pend

\pstart\versepar \ledsidenote{\passref{4.32.25}}\edtext{\skt{पाण्डुमेही}}{\lemma{\skt{पाण्डुमेही}}\Afootnote{K H; \skt{पाण्डुर्मेही}\ A}}\skt{ }\edtext{\skt{रक्तपित्ती}}{\lemma{\skt{रक्तपित्ती}}\Afootnote{K; \skt{पित्तरक्ती}\ A; \skt{रक्तमेही}\ H}}\skt{ }\edtext{\skt{तृषार्त्तः}}{\lemma{\skt{तृषार्त्तः}}\Afootnote{K H; \skt{क्षयार्तः}\ A}}\skt{ }\edtext{\skt{क्षुतक्षीणो}}{\lemma{\skt{क्षुतक्षीणो}}\Afootnote{K H\textsuperscript{ac}; \skt{क्षामो}\ A; \skt{क्षुक्षतक्षीणो}\ H\textsuperscript{pc}}}\skt{ }\edtext{\skt{दुर्वलो}}{\lemma{\skt{दुर्वलो}}\Afootnote{K H; \skt{ऽ}\ A}}\skt{ }\edtext{\skt{जीर्ण्णभक्तः}}{\lemma{\skt{जीर्ण्णभक्तः}}\Afootnote{K; \skt{जीर्णी}\ A; \skt{जीर्ण्णभोक्ताः~।}\ H}}\skt{~॥ }\pend
\pstart\versepar \skt{}\edtext{\skt{दको\-द\-री\-ग\-र्भिणी}}{\lemma{\skt{दको\appbrk{}द\appbrk{}री\appbrk{}ग\appbrk{}र्भिणी}}\Afootnote{K; \skt{चोदरार्तो विषार्तः~। तृड्च्छर्द्यार्तो गर्भिणी}\ A; \skt{दको\appbrk{}द\appbrk{}री\appbrk{}ग\appbrk{}र्व्भिणी}\ H}}\skt{ }\edtext{\skt{पानपञ्च}}{\lemma{\skt{पानपञ्च}}\Afootnote{K H; \skt{पीतमद्यो}\ A}}\skt{ नैते स्वेद्या }\edtext{\skt{यश्च\-म\-र्त्यो\-ति\-सारी}}{\lemma{\skt{यश्च\appbrk{}म\appbrk{}र्त्यो\appbrk{}ति\appbrk{}सारी}}\Afootnote{K; \skt{यश्च मर्त्योअतिसारी~।}\ A; \skt{यश्च\appbrk{}म\appbrk{}र्त्योऽ\appbrk{}ति\appbrk{}सारी~॥}\ H}}\skt{~। }\pend
\pstart\versepar \skt{}\edtext{\skt{स्वेदादेषा}}{\lemma{\skt{स्वेदादेषा}}\Afootnote{K; \skt{स्वेदादेषां}\ A; \skt{स्वेदा तेषां}\ H}}\skt{ यांति देहाविनाशं }\edtext{\skt{असाध्यत्वं}}{\lemma{\skt{असाध्यत्वं}}\Afootnote{K; \skt{नोसाध्यत्वं}\ A; \skt{मसाध्यत्वं}\ H}}\skt{ यान्ति चैषां विकाराः~। }\pend
\pstart\versepar \skt{स्वेदैः }\edtext{\skt{साध्यो}}{\lemma{\skt{साध्यो}}\Afootnote{A K; \skt{साद्यो}\ H}}\skt{ }\edtext{\skt{दुर्व\-लो\-जी\-र्ण्ण\-भक्ता}}{\lemma{\skt{दुर्व\appbrk{}लो\appbrk{}जी\appbrk{}र्ण्ण\appbrk{}भक्ता}}\Afootnote{K H; \skt{दुर्ब\appbrk{}लोऽ\appbrk{}जी\appbrk{}र्ण\appbrk{}भक्तः}\ A}}\skt{ }\edtext{\skt{यदि}}{\lemma{\skt{यदि}}\Afootnote{K H; \emph{om.}\ A}}\skt{ स्यातां }\edtext{\skt{स्वेदनीयौ}}{\lemma{\skt{स्वेदनीयौ}}\Afootnote{K H; \skt{चेद्द्वौ स्वेदनीयौ}\ A}}\skt{ ततस्तौ~।}\pend

\pstart\versepar \ledsidenote{\passref{4.32.27}}\edtext{\skt{सर्वस्वेदां}}{\lemma{\skt{सर्वस्वेदां}}\Afootnote{K H; \skt{सर्वान्स्वेदान्}\ A}}\skt{ }\edtext{\skt{निवाते तु}}{\lemma{\skt{निवाते तु}}\Afootnote{K H; \skt{निवाते}\ A}}\skt{ }\edtext{\skt{जीर्ण्णान्ने}}{\lemma{\skt{जीर्ण्णान्ने}}\Afootnote{K; \skt{च जीर्णान्}\ A; \skt{जीर्ण्णान्चा}\ H\textsuperscript{ac}; \skt{जीर्ण्णान्जीर्ण्णे न}\ H\textsuperscript{pc}}}\skt{ चावचारयेत्~। }\pend
\pstart\versepar \skt{स्नेहा\-भ्य\-क्त\-श\-री\-रस्य }\edtext{\skt{शीतैराच्छाद्य}}{\lemma{\skt{शीतैराच्छाद्य}}\Afootnote{\skt{शीतैराच्छाद्य}\ A K H}}\skt{ चक्षुषी~।}\edtext{}{\lemma{\skt{चावचारयेत्~। स्नेहा\appbrk{}भ्य\appbrk{}क्त\appbrk{}श\appbrk{}री\appbrk{}रस्य}}\Afootnote{K\textsuperscript{pc}; \skt{नस्या\appbrk{}व\appbrk{}चा\appbrk{}र\appbrk{}येत्~। स्नेहा\appbrk{}भ्य\appbrk{}क्त\appbrk{}श\appbrk{}री\appbrk{}रस्य}\ A; \skt{चावचारयेत्~। स्नेहा\appbrk{}भ्य\appbrk{}श\appbrk{}री\appbrk{}रस्य}\ K\textsuperscript{ac}; \skt{वचारयेत्~। स्नेहा\appbrk{}भ्य\appbrk{}क्त\appbrk{}श\appbrk{}री\appbrk{}रस्य}\ H\textsuperscript{ac}; \skt{वाचा\appbrk{}व\appbrk{}चा\appbrk{}र\appbrk{}येत्~। स्नेहा\appbrk{}भ्य\appbrk{}क्त\appbrk{}श\appbrk{}री\appbrk{}रस्य}\ H\textsuperscript{pc}}}\pend

\pstart\versepar \ledsidenote{\passref{4.32.28}}\edtext{\skt{अति\-स्वि\-न्न\-म\-था\-स्वाभ्य}}{\lemma{\skt{अति\appbrk{}स्वि\appbrk{}न्न\appbrk{}म\appbrk{}था\appbrk{}स्वाभ्य}}\Afootnote{\skt{स्विद्यमानस्य च मुहुर्हृदयं}\ A; \skt{अति\appbrk{}स्वि\appbrk{}न्न\appbrk{}था\appbrk{}स्वाभ्य}\ K\textsuperscript{ac}; \skt{अति\appbrk{}स्वि\appbrk{}न्न\appbrk{}म\appbrk{}था\appbrk{}स्वाभ्य}\ K\textsuperscript{pc} H}}\skt{ }\edtext{\skt{शीतांबुत्रासनं}}{\lemma{\skt{शीतांबुत्रासनं}}\Afootnote{K; \skt{शीतलैः स्पृशेत्~। सम्यक्स्विन्नं विमृ}\ A; \skt{शीताम्वुत्रासनं}\ H}}\skt{ }\edtext{\skt{हितं}}{\lemma{\skt{हितं}}\Afootnote{K H; \skt{दितं}\ A}}\skt{~। }\pend
\pstart\versepar \skt{}\edtext{\edtext{\skt{स्नानमुष्णाम्बुना}}{\lemma{\skt{स्नानम्}}\Afootnote{K H; \skt{स्नात}\ A}}}{\lemma{\skt{उष्णाम्बुना}}\Afootnote{K H; \skt{मुष्णाम्बु}\ A}}\skt{ }\edtext{\skt{चैव}}{\lemma{\skt{चैव}}\Afootnote{K H; \emph{om.}\ A}}\skt{ }\edtext{\skt{निव\-त\-ञ्चा\-ल\-य\-येत्}}{\lemma{\skt{निव\appbrk{}त\appbrk{}ञ्चा\appbrk{}ल\appbrk{}य\appbrk{}येत्}}\Afootnote{K; \skt{भिः शनैः~॥}\ A; \skt{विवातं चालयन्नयेत्~॥}\ H}}\skt{~।}\pend

\pstart \ledsidenote{\passref{4.32.29cd}}\edtext{\skt{भोज\-ये\-च्चा\-न\-भि\-ष्यन्दि}}{\lemma{\skt{भोज\appbrk{}ये\appbrk{}च्चा\appbrk{}न\appbrk{}भि\appbrk{}ष्यन्दि}}\Afootnote{K H; \skt{भोज\appbrk{}ये\appbrk{}द\appbrk{}न\appbrk{}भि\appbrk{}ष्यन्दि}\ A}}\skt{ सर्वं }\edtext{\edtext{\skt{वाचा\-र\-मा\-दि\-शे\-दिति}}{\lemma{\skt{वाचारम्}}\Afootnote{\skt{चाचारम्}\ A; \skt{वाचार}\ K H}}}{\lemma{\skt{आदिशेदिति}}\Afootnote{\skt{आदिशेत्~॥}\ A; \skt{मादिशेदिति~॥}\ K H}}\skt{~॥}\pend

\pstart \ledsidenote{\passref{4.32.trailer}}\edtext{\skt{चि॥ {\symfont ❈}॥}}{\lemma{\skt{चि॥ {\symfont ❈}}}\Afootnote{\skt{इति सुश्रु\appbrk{}त\appbrk{}सं\appbrk{}हि\appbrk{}तायां चिकित्सास्थाने स्वेदा\appbrk{}व\appbrk{}चा\appbrk{}र\appbrk{}णीयं चिकित्सितं नाम द्वात्रिंशोऽध्यायः~॥३२~॥}\ A; \skt{चि>॥{\symfont ❈}॥}\ K; \skt{॥ चिकित्सालद्वि॥}\ H}}\pend

\endnumbering
\beginnumbering
\eledsection*{Adhyāya 33}
\pstart \ledsidenote{\passref{4.33.1}}\skt{अथातो }\edtext{\skt{वम\-न\-वि\-रे\-च\-न\-सा\-ध्यो\-प\-क्र\-म\-चि\-कि\-त्सितं}}{\lemma{\skt{वम\appbrk{}न\appbrk{}वि\appbrk{}रे\appbrk{}च\appbrk{}न\appbrk{}सा\appbrk{}ध्यो\appbrk{}प\appbrk{}क्र\appbrk{}म\appbrk{}चि\appbrk{}कि\appbrk{}त्सितं}}\Afootnote{K H; \skt{वम\appbrk{}न\appbrk{}वि\appbrk{}रे\appbrk{}च\appbrk{}न\appbrk{}सा\appbrk{}ध्यो\appbrk{}प\appbrk{}द्र\appbrk{}व\appbrk{}चि\appbrk{}कि\appbrk{}त्सितं}\ A}}\skt{ }\edtext{\skt{व्याख्यास्यामः}}{\lemma{\skt{व्याख्यास्यामः}}\Afootnote{A H; \emph{om.}\ K}}\pend

\pstart \ledsidenote{\passref{4.33.3}}\skt{दोषाः क्षीणा बृंहयितव्याः }\edtext{\skt{चलिताः}}{\lemma{\skt{चलिताः}}\Afootnote{K H; \skt{कुपिताः}\ A}}\skt{ }\edtext{\skt{प्रशमयितव्याः}}{\lemma{\skt{प्रशमयितव्याः}}\Afootnote{K H; \skt{प्रशमयितव्याः,}\ A}}\skt{ }\edtext{\skt{प्रवृद्धा}}{\lemma{\skt{प्रवृद्धा}}\Afootnote{K H; \skt{वृद्धा}\ A}}\skt{ }\edtext{\skt{निर्हत्तव्याः}}{\lemma{\skt{निर्हत्तव्याः}}\Afootnote{K; \skt{निर्हर्तव्याः}\ A H}}\skt{ }\edtext{\skt{स्वस्थानस्थारक्ष्या}}{\lemma{\skt{स्वस्थानस्थारक्ष्या}}\Afootnote{K; \skt{समाः परिपाल्या}\ A; \skt{स्वस्था\appbrk{}न\appbrk{}स्था\appbrk{}न\appbrk{}रक्ष्या}\ H}}\skt{ इति सिद्धान्तः~॥}\pend

\pstart \ledsidenote{\passref{4.33.4}}\skt{प्राधान्येन }\edtext{\skt{वम\-न\-वि\-रे\-च\-ने\-नि\-र्ह\-रणे}}{\lemma{\skt{वम\appbrk{}न\appbrk{}वि\appbrk{}रे\appbrk{}च\appbrk{}ने\appbrk{}नि\appbrk{}र्ह\appbrk{}रणे}}\Afootnote{K H; \skt{वम\appbrk{}न\appbrk{}वि\appbrk{}रे\appbrk{}चने वर्तते निर्हरणे}\ A}}\skt{ }\edtext{\skt{वर्त्तते}}{\lemma{\skt{वर्त्तते}}\Afootnote{K H; \skt{दोषाणाम्~।}\ A}}\skt{~। }\edtext{\skt{तस्मा\-त्त\-यो\-र्वि\-धा\-न\-मु\-च्य\-मा\-न\-मु\-प\-धा\-र\-यस्व}}{\lemma{\skt{उपधारयस्व}}\Afootnote{K H; \skt{उपधारय~॥}\ A}}\skt{~।}\pend

\pstart \ledsidenote{\passref{4.33.5}}\edtext{\skt{तद्यथातुरं}}{\lemma{\skt{तद्यथातुरं}}\Afootnote{K H; \skt{अथातुरं}\ A}}\skt{ }\edtext{\edtext{\skt{स्निग्ध\-स्वि\-न्न\-म\-न\-भि\-ष्य\-न्दि\-भि\-रा\-हा\-रै\-र\-न\-व\-ब\-द्ध\-दो\-ष\-मा\-लोच्य}}{\lemma{\skt{स्निग्ध\appbrk{}स्वि\appbrk{}न्न\appbrk{}म\appbrk{}न\appbrk{}भि\appbrk{}ष्य\appbrk{}न्दि\appbrk{}भिर्}}\Afootnote{\skt{स्निग्धं स्विन्न\appbrk{}म\appbrk{}भि\appbrk{}ष्य\appbrk{}न्दि\appbrk{}भिर्}\ A; \skt{स्निग्ध\appbrk{}स्वि\appbrk{}न्न\appbrk{}म\appbrk{}न\appbrk{}भि\appbrk{}ष्य\appbrk{}न्दि\appbrk{}भिर्}\ K H}}}{\lemma{\skt{आहा\appbrk{}रै\appbrk{}र\appbrk{}न\appbrk{}व\appbrk{}ब\appbrk{}द्ध\appbrk{}दो\appbrk{}ष\appbrk{}मा\appbrk{}लोच्य}}\Afootnote{\skt{आहा\appbrk{}रै\appbrk{}र\appbrk{}न\appbrk{}व\appbrk{}ब\appbrk{}द्ध\appbrk{}दो\appbrk{}ष\appbrk{}म\appbrk{}व\appbrk{}लोक्य}\ A; \skt{आहा\appbrk{}रै\appbrk{}र\appbrk{}न\appbrk{}व\appbrk{}ब\appbrk{}द्ध\appbrk{}दो\appbrk{}ष\appbrk{}मा\appbrk{}लोच्य}\ K H}}\skt{ }\edtext{\skt{स्वो}}{\lemma{\skt{स्वो}}\Afootnote{K H\textsuperscript{ac}; \skt{श्वो}\ A; \skt{स्वोस्तो}\ H\textsuperscript{pc}}}\skt{ वमनं }\edtext{\skt{पाययितास्मीति}}{\lemma{\skt{पाययितास्मीति}}\Afootnote{K H; \skt{पाययिताऽस्मीति}\ A}}\skt{ }\edtext{\skt{संभोजयेत्~। संभोज्यस्तु}}{\lemma{\skt{संभोजयेत्~। संभोज्यस्तु}}\Afootnote{K H; \skt{संभोजये}\ A}}\skt{ }\edtext{\edtext{\skt{तीक्ष्णा\-ग्नि\-र्ब\-ल\-वा\-न्ब\-हु\-दोषो}}{\lemma{\skt{तीक्ष्णाग्निर्}}\Afootnote{K H; \skt{त्तीक्ष्णाग्निं}\ A}}}{\lemma{\skt{बल\appbrk{}वा\appbrk{}न्ब\appbrk{}हु\appbrk{}दोषो}}\Afootnote{K H; \skt{बलवन्तं बहुदोषं}\ A}}\skt{ }\edtext{\skt{महा\-व्या\-धि\-प\-रि\-गृ\-हीतो}}{\lemma{\skt{महा\appbrk{}व्या\appbrk{}धि\appbrk{}प\appbrk{}रि\appbrk{}गृ\appbrk{}हीतो}}\Afootnote{K H; \skt{महा\appbrk{}व्या\appbrk{}धि\appbrk{}प\appbrk{}रीतं}\ A}}\skt{ }\edtext{\skt{वमनसात्म्यश्च}}{\lemma{\skt{वमनसात्म्यश्}}\Afootnote{K H; \skt{वमनसात्म्यं}\ A}}\skt{~॥}\pend

\pstart\versepar \ledsidenote{\passref{4.33.6}}\edtext{\edtext{\edtext{\skt{पेश\-लै\-र्वि\-वि\-धै\-र\-न्नै\-र्दो\-ष\-मु\-त्क्लिश्य}}{\lemma{\skt{पेशलैर्}}\Afootnote{K H; \skt{भवति चात्र पेशलैर्}\ A}}}{\lemma{\skt{विविधैरन्नैर्}}\Afootnote{K H; \skt{विविधैरन्नैर्}\ A}}}{\lemma{\skt{दोषमुत्क्लिश्य}}\Afootnote{K H; \skt{दोषानुत्क्लेश्य}\ A}}\skt{ देहिनः~। }\pend
\pstart\versepar \skt{स्निग्धस्विन्नाय वमनं दत्तं सम्यक्प्रवर्त्तते~।}\pend

\pstart \ledsidenote{\passref{4.33.7}}\skt{अथापरेद्युः पूर्वाह्णे साधारणे काले }\edtext{\skt{वम\-न\-द्र\-व्य\-क\-षा\-य\-क\-ल्क\-चू\-र्ण्ण\-स्ने\-ह\-ना\-न्य\-त\-मस्य}}{\lemma{\skt{वम\appbrk{}न\appbrk{}द्र\appbrk{}व्य\appbrk{}क\appbrk{}षा\appbrk{}य\appbrk{}क\appbrk{}ल्क\appbrk{}चू\appbrk{}र्ण्ण\appbrk{}स्ने\appbrk{}ह\appbrk{}ना\appbrk{}न्य\appbrk{}त\appbrk{}मस्य}}\Afootnote{K; \skt{वम\appbrk{}न\appbrk{}द्र\appbrk{}व्य\appbrk{}क\appbrk{}षा\appbrk{}य\appbrk{}क\appbrk{}ल्क\appbrk{}चू\appbrk{}र्ण\appbrk{}स्ने\appbrk{}हा\appbrk{}ना\appbrk{}म\appbrk{}न्य\appbrk{}त\appbrk{}मस्य}\ A; \skt{वम\appbrk{}न\appbrk{}द्र\appbrk{}व्य\appbrk{}क\appbrk{}षा\appbrk{}य\appbrk{}क\appbrk{}ल्क\appbrk{}चू\appbrk{}र्ण्ण\appbrk{}स्ने\appbrk{}हे\appbrk{}ना\appbrk{}न्य\appbrk{}त\appbrk{}मस्य}\ H}}\skt{ मात्रां पाययित्वा }\edtext{\skt{वामयेत्~। असा\-त्म्य\-बी\-भ\-त्स\-दु\-द\-र्श\-न\-दु\-र्ग\-न्धानि}}{\lemma{\skt{वामयेत्~। असा\appbrk{}त्म्य\appbrk{}बी\appbrk{}भ\appbrk{}त्स\appbrk{}दु\appbrk{}द\appbrk{}र्श\appbrk{}न\appbrk{}दु\appbrk{}र्ग\appbrk{}न्धानि}}\Afootnote{K; \skt{वाम\appbrk{}ये\appbrk{}द्य\appbrk{}था\appbrk{}योगं कोष्ठ\appbrk{}वि\appbrk{}शे\appbrk{}ष\appbrk{}म\appbrk{}वेक्ष्य; असा\appbrk{}त्म्य\appbrk{}बी\appbrk{}भ\appbrk{}त्स\appbrk{}दु\appbrk{}र्ग\appbrk{}न्ध\appbrk{}दु\appbrk{}र्द\appbrk{}र्श\appbrk{}नानि}\ A; \skt{वामयेत्~। असा\appbrk{}त्म्य\appbrk{}बी\appbrk{}भ\appbrk{}त्स\appbrk{}दु\appbrk{}र्द्द\appbrk{}र्श\appbrk{}न\appbrk{}दु\appbrk{}र्ग्ग\appbrk{}न्धानि}\ H}}\skt{ }\edtext{\skt{वामनीयानि}}{\lemma{\skt{वामनीयानि}}\Afootnote{K H; \skt{च वमनानि}\ A}}\skt{ विदध्यादतो }\edtext{\skt{विपरीतानि}}{\lemma{\skt{विपरीतानि}}\Afootnote{K; \skt{विपरीतानि विरेचनानि}\ A H}}\skt{~।. तत्र }\edtext{\skt{सुकुमारां}}{\lemma{\skt{सुकुमारां}}\Afootnote{K H; \skt{सुकुमारं}\ A}}\skt{ }\edtext{\skt{कृशं}}{\lemma{\skt{कृशं}}\Afootnote{A K; \skt{कृशां}\ H}}\skt{ बालं }\edtext{\skt{वृद्धां}}{\lemma{\skt{वृद्धां}}\Afootnote{K H; \skt{वृद्धं}\ A}}\skt{ भीरुम्वा }\edtext{\skt{वम\-न\-सा\-ध्य\-वि\-का\-रेषु}}{\lemma{\skt{वम\appbrk{}न\appbrk{}सा\appbrk{}ध्य\appbrk{}वि\appbrk{}का\appbrk{}रेषु}}\Afootnote{K H; \skt{वमनसाध्येषु विकारेषु}\ A}}\skt{ क्षीर\-द\-धि\-त\-क्र\-य\-वा\-गू\-ना\-म\-न्य\-त\-म\-मा\-कण्ठं }\edtext{\skt{पाययित्वा}}{\lemma{\skt{पाययित्वा}}\Afootnote{K H; \skt{पाय}\ A}}\skt{ }\edtext{\skt{वामयेत्~। }\edtext{\skt{पीतौ\-ष\-ध\-ञ्च\-पा\-णि\-भि\-र\-ग्नि\-त\-प्तै\-रु\-प\-स्पृ\-श्य\-मानं}}{\lemma{\skt{उपस्पृश्यमानं}}\Afootnote{K H; \skt{प्रताप्यमानं}\ A}}}{\lemma{\skt{वामयेत्~। पीतौ\appbrk{}ष\appbrk{}ध\appbrk{}ञ्च\appbrk{}पा\appbrk{}णि\appbrk{}भिर्}}\Afootnote{K H; \skt{येत्, पीतौषधं च पाणिभिर्}\ A}}\skt{ }\edtext{\edtext{\skt{मुहू\-र्त्त\-मु\-पे\-क्षेत}}{\lemma{\skt{मुहूर्त्तम्}}\Afootnote{K; \skt{मुहूर्त}\ A H}}}{\lemma{\skt{उपेक्षेत}}\Afootnote{K; \skt{मुपेक्षेत}\ A H}}\skt{~। }\edtext{\skt{तत्र}}{\lemma{\skt{तत्र}}\Afootnote{K H; \skt{तस्य च स्वेद\appbrk{}प्रा\appbrk{}दु\appbrk{}र्भा\appbrk{}वेण शिथि\appbrk{}ल\appbrk{}ता\appbrk{}मा\appbrk{}पन्नं स्वेभ्यः स्थानेभ्यः प्रचलितं कुक्षिमनुसृतं जानीयात्, ततः}\ A}}\skt{ }\edtext{\skt{प्रवृत्तहृल्लासं}}{\lemma{\skt{प्रवृत्तहृल्लासं}}\Afootnote{A K; \skt{प्रवृत्तं हृल्लासं}\ H}}\skt{ ज्ञात्वा }\edtext{\edtext{\skt{जानु\-मा\-त्रा\-स\-नो\-प\-वि\-ष्ट\-म\-न्येन}}{\lemma{\skt{जानु\appbrk{}मा\appbrk{}त्रा\appbrk{}स\appbrk{}नो\appbrk{}प\appbrk{}वि\appbrk{}ष्टम्}}\Afootnote{K H; \skt{जानु\appbrk{}मा\appbrk{}त्रा\appbrk{}स\appbrk{}नो\appbrk{}प\appbrk{}विष्ट}\ A}}}{\lemma{\skt{अन्येन}}\Afootnote{K H; \skt{माप्तैर्ल}\ A}}\skt{ }\edtext{\skt{हृत्पृष्ठयोः}}{\lemma{\skt{हृत्पृष्ठयोः}}\Afootnote{K H; \skt{लाटे पृष्ठे पार्श्वयोः कण्ठे च}\ A}}\skt{ }\edtext{\skt{पाणिभ्यां}}{\lemma{\skt{पाणिभ्यां}}\Afootnote{K H; \skt{पाणिभिः}\ A}}\skt{ }\edtext{\edtext{\skt{उप\-रि\-गृ\-ही\-त\-म\-ङ्गु\-लि\-ग\-न्ध\-र्व\-ह\-स्त\-प\-त्र\-वृ\-त्तो\-त्प\-ल\-ना\-ला\-ना\-म\-न्य\-त\-मेन}}{\lemma{\skt{उप\appbrk{}रि\appbrk{}गृ\appbrk{}ही\appbrk{}तम्}}\Afootnote{K; \skt{सुपरिगृहीत}\ A; \skt{उपरि गृहीत}\ H}}}{\lemma{\skt{अङ्गु\appbrk{}लि\appbrk{}ग\appbrk{}न्ध\appbrk{}र्व\appbrk{}ह\appbrk{}स्त\appbrk{}प\appbrk{}त्र\appbrk{}वृ\appbrk{}त्तो\appbrk{}त्प\appbrk{}ल\appbrk{}ना\appbrk{}ला\appbrk{}नाम्}}\Afootnote{K; \skt{मङ्गु\appbrk{}ली\appbrk{}ग\appbrk{}न्ध\appbrk{}र्व\appbrk{}ह\appbrk{}स्तो\appbrk{}त्प\appbrk{}ल\appbrk{}ना\appbrk{}ला\appbrk{}नाम्}\ A; \skt{मङ्गु\appbrk{}लि\appbrk{}ग\appbrk{}न्ध\appbrk{}र्व्व\appbrk{}ह\appbrk{}स्त\appbrk{}प\appbrk{}त्त्र\appbrk{}वृ\appbrk{}त्येत्ना लानाम्}\ H}}\skt{ }\edtext{\edtext{\skt{कण्ठ\-म\-न\-भि\-स्पृ\-श\-न्वा\-म\-ये\-त्या\-व\-त्स\-म्य\-ग्वा\-न्त\-ल\-क्ष\-णानि}}{\lemma{\skt{कण्ठमनभिस्}}\Afootnote{K H; \skt{कण्ठमभिस्}\ A}}}{\lemma{\skt{पृश\appbrk{}न्वा\appbrk{}म\appbrk{}ये\appbrk{}त्या\appbrk{}व\appbrk{}त्स\appbrk{}म्य\appbrk{}ग्वा\appbrk{}न्त\appbrk{}ल\appbrk{}क्ष\appbrk{}णानि}}\Afootnote{K H; \skt{पृशन्तं वाम\appbrk{}ये\appbrk{}त्ता\appbrk{}व\appbrk{}द्या\appbrk{}व\appbrk{}त्स\appbrk{}म्य\appbrk{}ग्वा\appbrk{}न्त\appbrk{}लि\appbrk{}ङ्गा\appbrk{}नीति~॥}\ A}}\skt{ }\edtext{\skt{भवन्ति}}{\lemma{\skt{भवन्ति}}\Afootnote{K H; \emph{om.}\ A}}\skt{~॥}\pend

\pstart\versepar \ledsidenote{\passref{4.33.8}}\edtext{\skt{कफप्रसेकं}}{\lemma{\skt{कफप्रसेकं}}\Afootnote{K H; \skt{भवतश्चात्र कफप्रसेकं}\ A}}\skt{ }\edtext{\skt{हृदयाविशुद्धिं}}{\lemma{\skt{हृदयाविशुद्धिं}}\Afootnote{K H; \skt{हृद्याविशुद्धिं}\ A}}\skt{ कण्डूञ्च दुश्छ\-र्दि\-त\-लि\-ङ्ग\-माहुः~। }\pend
\pstart\versepar \skt{पित्तातियोगञ्च }\edtext{\skt{विसङ्ज्ञताञ्च}}{\lemma{\skt{विसङ्ज्ञताञ्च}}\Afootnote{K H; \skt{चिसंज्ञतां च}\ A}}\skt{ हृत्क\-ण्ठ\-पी\-डा\-मपि चातिवान्ते~।}\pend

\pstart\versepar \ledsidenote{\passref{4.33.9}}\skt{पित्ते }\edtext{\skt{कफस्यान्द्रमुखं}}{\lemma{\skt{कफस्यान्द्रमुखं}}\Afootnote{K H; \skt{कफस्यानु सुखं}\ A}}\skt{ प्रवृत्ते शुद्धेषु हृत्कण्ठशिरस्सु चापि~। }\pend
\pstart\versepar \skt{लघौ च देहे कफसंस्रवे च स्थिते }\edtext{\skt{सुवान्तं}}{\lemma{\skt{सुवान्तं}}\Afootnote{A K\textsuperscript{pc} H; \skt{सुवान्तंसुवातं}\ K\textsuperscript{ac}}}\skt{ पुरुषं }\edtext{\skt{व्यववस्येत्}}{\lemma{\skt{व्यववस्येत्}}\Afootnote{K; \skt{व्यवस्येत्}\ A H}}\skt{~।}\pend

\pstart \ledsidenote{\passref{4.33.10}}\skt{सम्यग्वान्तं }\edtext{\edtext{\skt{चैन\-म\-भि\-स\-मीक्ष्य}}{\lemma{\skt{चैनम्}}\Afootnote{A K; \skt{चैन}\ H}}}{\lemma{\skt{अभिस्}}\Afootnote{A K; \skt{मभिस्}\ H}}\skt{ }\edtext{\skt{स्नैहि\-क\-वै\-रे\-च\-नि\-को\-प\-श\-म\-नी\-यानां}}{\lemma{\skt{स्नैहि\appbrk{}क\appbrk{}वै\appbrk{}रे\appbrk{}च\appbrk{}नि\appbrk{}को\appbrk{}प\appbrk{}श\appbrk{}म\appbrk{}नी\appbrk{}यानां}}\Afootnote{K H; \skt{स्नेह\appbrk{}न\appbrk{}वि\appbrk{}रे\appbrk{}च\appbrk{}न\appbrk{}श\appbrk{}म\appbrk{}नानां}\ A}}\skt{ }\edtext{\skt{धूमानामन्यतं}}{\lemma{\skt{अन्यतं}}\Afootnote{K; \skt{अन्यतमं}\ A H}}\skt{ सामर्थ्यतः पाययित्वा }\edtext{\edtext{\skt{चारि\-क\-मु\-प\-दि\-शेत्}}{\lemma{\skt{चारिकम्}}\Afootnote{K H; \skt{ऽऽआचारिकम्}\ A}}}{\lemma{\skt{उपदिशेत्}}\Afootnote{K H; \skt{आदिशेत्~॥}\ A}}\skt{~॥}\pend

\pstart\versepar \ledsidenote{\passref{4.33.11}}\edtext{\skt{ततो}}{\lemma{\skt{ततो}}\Afootnote{K H; \skt{भवन्ति चात्र ततो}\ A}}\skt{ पराह्णे }\edtext{\edtext{\edtext{\skt{सुवि\-भ\-क्त\-दे\-ह\-मु\-ष्णा\-भि\-रह्निः}}{\lemma{\skt{सुवि\appbrk{}भ\appbrk{}क्त\appbrk{}दे\appbrk{}हम्}}\Afootnote{K H; \skt{शुचिशुद्धदेह}\ A}}}{\lemma{\skt{उष्णाभिर्}}\Afootnote{K; \skt{मुष्णाभिर्}\ A; \skt{उष्णाभि}\ H}}}{\lemma{\skt{अह्निः}}\Afootnote{K; \skt{अद्द्भिः}\ A; \skt{वह्निः}\ H}}\skt{ परिषिक्तगात्रं~। }\pend
\pstart\versepar \skt{}\edtext{\edtext{\skt{कुल\-त्थ\-मु\-द्ग\-ढ\-कि\-जा\-ङ्ग\-लानां}}{\lemma{\skt{कुलत्थम्}}\Afootnote{K; \skt{कुलत्थ}\ A H}}}{\lemma{\skt{उद्ग\appbrk{}ढ\appbrk{}कि\appbrk{}जा\appbrk{}ङ्ग\appbrk{}लानां}}\Afootnote{K; \skt{मुद्गा\appbrk{}ढ\appbrk{}कि\appbrk{}जा\appbrk{}ङ्ग\appbrk{}लानां}\ A H}}\skt{ }\edtext{\skt{युषै}}{\lemma{\skt{युषै}}\Afootnote{K H; \skt{यूषै}\ A}}\skt{ }\edtext{\edtext{\edtext{\skt{रसै\-र्वा\-प्यु\-प\-भो\-ज\-यीत}}{\lemma{\skt{रसैर्}}\Afootnote{A K; \skt{रसै}\ H}}}{\lemma{\skt{वाप्य्}}\Afootnote{K; \skt{वाऽ}\ A; \skt{र्व्वा}\ H}}}{\lemma{\skt{उपभोजयीत}}\Afootnote{K; \skt{प्युपभोजयेत्तु~॥}\ A; \skt{प्युपभोजयीत~॥}\ H}}\skt{~।}\pend

\pstart\versepar \ledsidenote{\passref{4.33.12}}\edtext{\skt{कफप्रसेकः}}{\lemma{\skt{कफप्रसेकः}}\Afootnote{K H; \skt{कासोपलेप}\ A}}\skt{ }\edtext{\skt{स्वरभेदतन्द्रा}}{\lemma{\skt{स्वरभेदतन्द्रा}}\Afootnote{K H; \skt{स्वरभेदनिद्रा}\ A}}\skt{ }\edtext{\skt{निद्रा\-स्य\-दौ\-र्ग\-न्ध्य\-वि\-षो\-प\-सर्गाः}}{\lemma{\skt{निद्रा\appbrk{}स्य\appbrk{}दौ\appbrk{}र्ग\appbrk{}न्ध्य\appbrk{}वि\appbrk{}षो\appbrk{}प\appbrk{}सर्गाः}}\Afootnote{K H; \skt{तन्द्रा\appbrk{}स्य\appbrk{}दौ\appbrk{}र्ग\appbrk{}न्ध्य\appbrk{}वि\appbrk{}षो\appbrk{}प\appbrk{}सर्गाः~।}\ A}}\skt{~। }\pend
\pstart\versepar \skt{}\edtext{\skt{गुरु\-त्व\-क\-ण्डू\-ग्र\-ह\-णी\-प्र\-दोषा न}}{\lemma{\skt{गुरु\appbrk{}त्व\appbrk{}क\appbrk{}ण्डू\appbrk{}ग्र\appbrk{}ह\appbrk{}णी\appbrk{}प्र\appbrk{}दोषा न}}\Afootnote{K H; \skt{कफ\appbrk{}प्र\appbrk{}से\appbrk{}क\appbrk{}ग्र\appbrk{}ह\appbrk{}णी\appbrk{}प्र\appbrk{}दोषा न}\ A}}\skt{ सन्ति\-ज\-न्तो\-र्व\-मतः कदाचित्~।}\pend

\pstart\versepar \ledsidenote{\passref{4.33.13}}\skt{छिन्ने तरौ पुष्प\-फ\-ल\-प्र\-रोहा यथाविनाशं सहसा व्रजन्ति~। }\pend
\pstart\versepar \skt{तथाहृते श्लेष्मणि }\edtext{\skt{च्छर्द्दनेन}}{\lemma{\skt{च्छर्द्दनेन}}\Afootnote{K; \skt{शोधनेन}\ A; \skt{च्छर्द्दनेन}\ H}}\skt{ तज्जा }\edtext{\skt{विकारा}}{\lemma{\skt{विकारा}}\Afootnote{K H; \skt{विकाराः}\ A}}\skt{ }\edtext{\skt{विलयं}}{\lemma{\skt{विलयं}}\Afootnote{K H; \skt{प्रशमं}\ A}}\skt{ }\edtext{\skt{व्रजन्ति}}{\lemma{\skt{व्रजन्ति}}\Afootnote{K H; \skt{प्रयान्ति~॥}\ A}}\skt{~॥}\pend

\pstart\versepar \ledsidenote{\passref{4.33.14}}\edtext{\skt{न वाम\-ये\-त्तै\-मि\-रि\-क\-न्न\-गु\-ल्मिनं न}}{\lemma{\skt{न वाम\appbrk{}ये\appbrk{}त्तै\appbrk{}मि\appbrk{}रि\appbrk{}क\appbrk{}न्न\appbrk{}गु\appbrk{}ल्मिनं न}}\Afootnote{K H; \skt{न वाम\appbrk{}ये\appbrk{}त्तै\appbrk{}मि\appbrk{}रि\appbrk{}को\appbrk{}र्ध्व\appbrk{}वा\appbrk{}त\appbrk{}गुल्मो}\ A}}\skt{ }\edtext{\skt{चैव\-पा\-ण्डू\-द\-र\-रो\-ग\-पी\-डितं}}{\lemma{\skt{चैव\appbrk{}पा\appbrk{}ण्डू\appbrk{}द\appbrk{}र\appbrk{}रो\appbrk{}ग\appbrk{}पी\appbrk{}डितं}}\Afootnote{K H; \skt{दर\appbrk{}प्ली\appbrk{}ह\appbrk{}कृ\appbrk{}मि\appbrk{}श्र\appbrk{}मार्ता}\ A}}\skt{~। }\pend
\pstart\versepar \skt{}\edtext{\skt{स्थूल\-क्ष\-त\-क्षी\-ण\-कृ\-शा\-ति\-वृद्धा}}{\lemma{\skt{स्थूल\appbrk{}क्ष\appbrk{}त\appbrk{}क्षी\appbrk{}ण\appbrk{}कृ\appbrk{}शा\appbrk{}ति\appbrk{}वृद्धा}}\Afootnote{K; \skt{न्~। स्थूल\appbrk{}क्ष\appbrk{}त\appbrk{}क्षी\appbrk{}ण\appbrk{}कृ\appbrk{}शा\appbrk{}ति\appbrk{}वृद्ध}\ A; \skt{स्थूल\appbrk{}क्ष\appbrk{}त\appbrk{}क्षी\appbrk{}ण\appbrk{}कृ\appbrk{}षा\appbrk{}ति\appbrk{}वृद्धा}\ H}}\skt{ }\edtext{\skt{नर्शोर्दिता}}{\lemma{\skt{नर्शोर्दिता}}\Afootnote{K H; \skt{मूत्रातुरा}\ A}}\skt{ }\edtext{\edtext{\skt{क्षेप\-क\-पी\-डि\-तांश्च}}{\lemma{\skt{क्षेप\appbrk{}क\appbrk{}पी\appbrk{}डि\appbrk{}तांश्}}\Afootnote{K; \skt{न्केवलवा}\ A; \skt{क्षेप\appbrk{}क\appbrk{}पी\appbrk{}डि\appbrk{}ताश्}\ H}}}{\lemma{\skt{च}}\Afootnote{K H; \skt{तरोगान्~॥}\ A}}\skt{~।}\pend

\pstart\versepar \ledsidenote{\passref{4.33.15}}\edtext{\skt{रूक्षं}}{\lemma{\skt{रूक्षं}}\Afootnote{K H; \skt{स्वरो\appbrk{}प\appbrk{}घा\appbrk{}ता\appbrk{}ध्य\appbrk{}य\appbrk{}न\appbrk{}प्र\appbrk{}स\appbrk{}क्त\appbrk{}दु\appbrk{}श्छ\appbrk{}र्दि\appbrk{}दुः\appbrk{}को\appbrk{}ष्ठ\appbrk{}तृ\appbrk{}डा\appbrk{}र्त\appbrk{}बा\appbrk{}लाम्~। ऊर्ध्वा\appbrk{}स्र\appbrk{}पि\appbrk{}त्ति\appbrk{}क्षु\appbrk{}धि\appbrk{}ता\appbrk{}ति\appbrk{}रूक्ष}\ A}}\skt{ }\edtext{\skt{प्रमेहे}}{\lemma{\skt{प्रमेहे}}\Afootnote{K H; \emph{om.}\ A}}\skt{ }\edtext{\skt{तरणे च}}{\lemma{\skt{तरणे च}}\Afootnote{K; \emph{om.}\ A; \skt{तरुणे च}\ H}}\skt{ }\edtext{\skt{गर्भे}}{\lemma{\skt{गर्भे}}\Afootnote{K H; \skt{गर्भि}\ A}}\skt{ }\edtext{\skt{गच्छत्यधोर्ध्वं}}{\lemma{\skt{गच्छत्यधोर्ध्वं}}\Afootnote{K; \skt{ण्युदा\appbrk{}व\appbrk{}र्ति\appbrk{}रू\appbrk{}हि\appbrk{}तांश्च~॥}\ A; \skt{गच्छत्यधोर्द्धं}\ H}}\skt{ }\edtext{\skt{रुधिरे च}}{\lemma{\skt{रुधिरे च}}\Afootnote{K H; \emph{om.}\ A}}\skt{ }\edtext{\skt{तीव्रे}}{\lemma{\skt{तीव्रे}}\Afootnote{K H; \emph{om.}\ A}}\skt{~। }\pend
\pstart\versepar \skt{}\edtext{\skt{तुष्टे च}}{\lemma{\skt{तुष्टे च}}\Afootnote{K; \emph{om.}\ A; \skt{दुष्टे च}\ H}}\skt{ }\edtext{\skt{कोष्ठे}}{\lemma{\skt{कोष्ठे}}\Afootnote{K H; \emph{om.}\ A}}\skt{ }\edtext{\edtext{\skt{क्रिमिभिर्मनुष्यं}}{\lemma{\skt{क्रिमिभिर्}}\Afootnote{K H; \emph{om.}\ A}}\skt{ न}}{\lemma{\skt{मनुष्यं न}}\Afootnote{K H; \emph{om.}\ A}}\skt{ }\edtext{\skt{वामयेद्वर्चसि}}{\lemma{\skt{वामयेद्वर्चसि}}\Afootnote{K H; \emph{om.}\ A}}\skt{ }\edtext{\skt{चातिवृद्धे}}{\lemma{\skt{चातिवृद्धे}}\Afootnote{K H; \emph{om.}\ A}}\skt{~॥}\pend

\pstart\versepar \ledsidenote{\passref{4.33.17}}\edtext{\skt{एते }\edtext{\skt{प्यजीर्ण्णव्यथिता}}{\lemma{\skt{अजीर्ण्णव्यथिता}}\Afootnote{K; \skt{प्यजीर्णव्यथिता}\ A H}}}{\lemma{\skt{एते प्य्}}\Afootnote{K; \skt{एतेऽ}\ A; \skt{एते}\ H}}\skt{ वाम्या ये च विषातुराः~। }\pend
\pstart\versepar \skt{}\edtext{\skt{अती\-व\-चो\-ल्ब\-ण\-क\-फास्ते च }\edtext{\skt{स्युर्मधुकाम्बुना}}{\lemma{\skt{मधुकाम्बुना}}\Afootnote{A H; \skt{र्मधुकाम्बुना~॥}\ K}}}{\lemma{\skt{अती\appbrk{}व\appbrk{}चो\appbrk{}ल्ब\appbrk{}ण\appbrk{}क\appbrk{}फास्ते च स्युर्}}\Afootnote{\skt{अतीव चोल्बणकफास्ते च स्युर्}\ A H; \skt{अती\appbrk{}व\appbrk{}चो\appbrk{}ल्ब\appbrk{}ण\appbrk{}क\appbrk{}फास्ते च स्यु}\ K}}\skt{~॥}\pend

\pstart \ledsidenote{\passref{4.33.19}}\skt{विरेचनमपि }\edtext{\skt{स्निग्धस्विन्नाय च}}{\lemma{\skt{स्निग्धस्विन्नाय च}}\Afootnote{K H; \skt{स्निग्धस्विन्नाय वान्ताय च}\ A}}\skt{ देयं~। अवान्तस्य हि सम्य\-ग्वि\-रि\-क्त\-स्यापि }\edtext{\skt{सतोवस्रस्तः}}{\lemma{\skt{सतोवस्रस्तः}}\Afootnote{K H; \skt{सतोऽधः स्रस्तः}\ A}}\skt{ }\edtext{\skt{श्लेष्माग्र}}{\lemma{\skt{श्लेष्माग्र}}\Afootnote{A K; \skt{श्लेष्मग्र}\ H}}\skt{ }\edtext{\skt{हणीञ्छादयति}}{\lemma{\skt{छादयति}}\Afootnote{A K; \skt{चादयति~।}\ H}}\skt{~। }\edtext{\edtext{\skt{गौर\-व\-मा\-पा\-द\-यति}}{\lemma{\skt{गौरवम्}}\Afootnote{A K; \skt{गौरव}\ H}}}{\lemma{\skt{आपादयति}}\Afootnote{A K; \skt{मापादयति~।}\ H}}\skt{ प्रवाहिकाम्वा }\edtext{\skt{कुरुते}}{\lemma{\skt{कुरुते}}\Afootnote{K H; \skt{जनय}\ A}}\skt{~। }\edtext{\skt{तस्मा\-त्स्नि\-ग्ध\-स्वि\-न्नाय}}{\lemma{\skt{तस्मा\appbrk{}त्स्नि\appbrk{}ग्ध\appbrk{}स्वि\appbrk{}न्नाय}}\Afootnote{K H; \skt{ति~॥}\ A}}\skt{ }\edtext{\skt{वान्ताय च}}{\lemma{\skt{वान्ताय च}}\Afootnote{K H; \emph{om.}\ A}}\skt{ }\edtext{\skt{देयम्}}{\lemma{\skt{देयम्}}\Afootnote{K H; \emph{om.}\ A}}\skt{~।}\pend

\pstart \ledsidenote{\passref{4.33.20}}\edtext{\skt{अथातुरं स्वो}}{\lemma{\skt{अथातुरं स्वो}}\Afootnote{K H; \skt{अथातुरं श्वो}\ A}}\skt{ विरेचनं }\edtext{\skt{पाययितास्मीति}}{\lemma{\skt{पाययितास्मीति}}\Afootnote{K H; \skt{पाययिताऽस्मीति}\ A}}\skt{ पूर्वाह्णे }\edtext{\skt{लघुभोजयेत्~। फला\-म्व्ल\-मु\-ष्णो\-द\-क\-श्चै\-न\-म\-नु\-पा\-य\-येत्~। अथापरेहनि}}{\lemma{\skt{लघुभोजयेत्~। फला\appbrk{}म्व्ल\appbrk{}मु\appbrk{}ष्णो\appbrk{}द\appbrk{}क\appbrk{}श्चै\appbrk{}न\appbrk{}म\appbrk{}नु\appbrk{}पा\appbrk{}य\appbrk{}येत्~। अथापरेहनि}}\Afootnote{K; \skt{लघु भोजयेत् , फला\appbrk{}म्ल\appbrk{}मु\appbrk{}ष्णो\appbrk{}दकं चैन\appbrk{}म\appbrk{}नु\appbrk{}पा\appbrk{}य\appbrk{}येत्~। अथापरेऽहनि}\ A; \skt{लघु भोजयेत्~। फला\appbrk{}म्ल\appbrk{}मु\appbrk{}ष्णो\appbrk{}द\appbrk{}क\appbrk{}ञ्चै\appbrk{}न\appbrk{}म\appbrk{}नु\appbrk{}पा\appbrk{}य\appbrk{}येत्~। अथापरेऽहनि}\ H}}\skt{ }\edtext{\edtext{\skt{विग\-त\-श्ले\-ष्म\-धा\-ता\-वा\-तु\-रो\-प\-क्र\-म\-णी\-या\-द\-पे\-क्ष्या\-तु\-र\-म\-थास्मै}}{\lemma{\skt{विग\appbrk{}त\appbrk{}श्ले\appbrk{}ष्म\appbrk{}धा\appbrk{}ता\appbrk{}वा\appbrk{}तु\appbrk{}रो\appbrk{}प\appbrk{}क्र\appbrk{}म\appbrk{}णी\appbrk{}या\appbrk{}द\appbrk{}पे\appbrk{}क्ष्या\appbrk{}तु\appbrk{}रम्}}\Afootnote{K; \skt{विग\appbrk{}त\appbrk{}श्ले\appbrk{}ष्म\appbrk{}धा\appbrk{}तु\appbrk{}मा\appbrk{}तु\appbrk{}रो\appbrk{}प\appbrk{}क्र\appbrk{}म\appbrk{}णी\appbrk{}या\appbrk{}द\appbrk{}वे\appbrk{}क्ष्या\appbrk{}तुर}\ A; \skt{विग\appbrk{}त\appbrk{}श्ले\appbrk{}ष्म\appbrk{}धा\appbrk{}ता\appbrk{}वा\appbrk{}तु\appbrk{}रो\appbrk{}प\appbrk{}क्र\appbrk{}म\appbrk{}णी\appbrk{}या\appbrk{}द\appbrk{}पे\appbrk{}क्ष्या\appbrk{}तुर}\ H}}}{\lemma{\skt{अथास्मै}}\Afootnote{K; \skt{मथास्मै}\ A H}}\skt{ }\edtext{\skt{विरेचनं}}{\lemma{\skt{विरेचनं}}\Afootnote{K H; \skt{औषध}\ A}}\skt{ }\edtext{\skt{मात्रां}}{\lemma{\skt{मात्रां}}\Afootnote{A K; \skt{मात्रम्}\ H}}\skt{ }\edtext{\skt{प्रयच्छेत्}}{\lemma{\skt{प्रयच्छेत्}}\Afootnote{K H; \skt{पातुं प्रयच्छेत्~॥}\ A}}\skt{~।}\pend

\pstart \ledsidenote{\passref{4.33.21}}\skt{तत्र मृदुः क्रूरो }\edtext{\skt{मध्य}}{\lemma{\skt{मध्य}}\Afootnote{K H; \skt{मध्यम}\ A}}\skt{ इति त्रिविधः कोष्ठो भवति~। तत्र बहुपित्तो मृदुः स }\edtext{\skt{दुग्धेना}}{\lemma{\skt{दुग्धेना}}\Afootnote{K H; \skt{दुग्धेनापि}\ A}}\skt{ विरिच्यते~। बहुवातश्लेष्मा }\edtext{\skt{क्रूरो}}{\lemma{\skt{क्रूरो}}\Afootnote{K H; \skt{क्रूरः, स}\ A}}\skt{ }\edtext{\skt{दुर्विरेच्यः}}{\lemma{\skt{दुर्विरेच्यः}}\Afootnote{A K; \skt{दुर्व्विरिच्यः~।}\ H}}\skt{~। }\edtext{\skt{सगदोषो}}{\lemma{\skt{सगदोषो}}\Afootnote{K; \skt{समदोषो}\ A H}}\skt{ }\edtext{\skt{मध्यः}}{\lemma{\skt{मध्यः}}\Afootnote{K H; \skt{मध्यमः, स}\ A}}\skt{ साधारण इति~। तत्र मृदौ }\edtext{\skt{मृद्वी}}{\lemma{\skt{मृद्वी}}\Afootnote{K H; \skt{मात्रा}\ A}}\skt{ }\edtext{\skt{मात्रा}}{\lemma{\skt{मात्रा}}\Afootnote{K H; \skt{मृद्वी,}\ A}}\skt{~। }\edtext{\skt{तीक्ष्णक्रूरे}}{\lemma{\skt{तीक्ष्णक्रूरे}}\Afootnote{K H; \skt{तीक्ष्णा क्रूरे,}\ A}}\skt{~। मध्ये मध्या कर्त्तव्येति~। }\edtext{\skt{पीतौषधश्च}}{\lemma{\skt{पीतौषधश्च}}\Afootnote{A K; \skt{पीतौषधश्च}\ H}}\skt{ तन्मनाः }\edtext{\skt{शय्याभ्यासे}}{\lemma{\skt{शय्याभ्यासे}}\Afootnote{K H; \skt{शय्याभ्याशे}\ A}}\skt{ }\edtext{\skt{विरिच्येत~। न}}{\lemma{\skt{विरिच्येत~। न}}\Afootnote{K H\textsuperscript{pc}; \skt{विरेच्यते~॥}\ A; \skt{विरिच्येत}\ H\textsuperscript{ac}}}\skt{ }\edtext{\skt{प्राप्तवेगं}}{\lemma{\skt{प्राप्तवेगं}}\Afootnote{K H\textsuperscript{pc}; \emph{om.}\ A H\textsuperscript{ac}}}\skt{ }\edtext{\skt{विधारयेत्~। न}}{\lemma{\skt{विधारयेत्~। न}}\Afootnote{K H\textsuperscript{pc}; \emph{om.}\ A H\textsuperscript{ac}}}\skt{ }\edtext{\skt{चाप्राप्तं}}{\lemma{\skt{चाप्राप्तं}}\Afootnote{K H\textsuperscript{pc}; \emph{om.}\ A H\textsuperscript{ac}}}\skt{ }\edtext{\skt{प्राणेनाकांक्षेत}}{\lemma{\skt{प्राणेनाकांक्षेत}}\Afootnote{K H\textsuperscript{pc}; \emph{om.}\ A H\textsuperscript{ac}}}\skt{~।}\pend

\pstart \ledsidenote{\passref{4.33.23}}\skt{यथा च वमने }\edtext{\skt{क्रमेण}}{\lemma{\skt{क्रमेण}}\Afootnote{K H; \emph{om.}\ A}}\skt{ }\edtext{\skt{प्रसेकौषधः}}{\lemma{\skt{प्रसेकौषधः}}\Afootnote{K H; \skt{प्रसेकौषध}\ A}}\skt{ }\edtext{\edtext{\skt{पित्ता\-नि\-ला\-ग\-च्छ\-न्त्येवं}}{\lemma{\skt{पित्ता\appbrk{}नि\appbrk{}ला\appbrk{}ग\appbrk{}च्छन्त्य्}}\Afootnote{K; \skt{कफपित्तानिलाः क्रमेण गच्छन्ति,}\ A; \skt{पित्तानिला गच्छन्}\ H}}}{\lemma{\skt{एवं}}\Afootnote{A K; \skt{त्येवम्}\ H}}\skt{ विरेचने }\edtext{\skt{वात\-मू\-त्र\-पु\-री\-ष\-पि\-त्त\-कफा}}{\lemma{\skt{वात\appbrk{}मू\appbrk{}त्र\appbrk{}पु\appbrk{}री\appbrk{}ष\appbrk{}पि\appbrk{}त्त\appbrk{}कफा}}\Afootnote{K H; \skt{मूत्र\appbrk{}पु\appbrk{}री\appbrk{}ष\appbrk{}पि\appbrk{}त्तौ\appbrk{}ष\appbrk{}ध\appbrk{}कफा}\ A}}\skt{ इति~॥}\pend

\pstart\versepar \ledsidenote{\passref{4.33.24}}\edtext{\skt{भ~॥हृत्कुक्ष्यशुद्धिः}}{\lemma{\skt{भ~॥हृत्कुक्ष्यशुद्धिः}}\Afootnote{K; \skt{भवन्ति चात्र हृत्कुक्ष्यशुद्धिः}\ A; \skt{भ~॥ हृत्कुक्ष्याशुद्धिः}\ H\textsuperscript{ac}; \skt{भ~॥ हृत्कुक्ष्याक्षिशुद्धिः}\ H\textsuperscript{pc}}}\skt{ }\edtext{\skt{परि\-दा\-ह\-क\-ण्डू\-वि\-ण्मू\-त्र\-सं\-गश्च न}}{\lemma{\skt{परि\appbrk{}दा\appbrk{}ह\appbrk{}क\appbrk{}ण्डू\appbrk{}वि\appbrk{}ण्मू\appbrk{}त्र\appbrk{}सं\appbrk{}गश्च न}}\Afootnote{K H; \skt{परि\appbrk{}दा\appbrk{}ह\appbrk{}क\appbrk{}ण्डू\appbrk{}वि\appbrk{}ण्मू\appbrk{}त्र\appbrk{}स\appbrk{}ङ्गाश्च न}\ A}}\skt{ }\edtext{\skt{सद्विरिक्ते}}{\lemma{\skt{सद्विरिक्ते}}\Afootnote{A K; \skt{सिद्विरिक्ते~।}\ H\textsuperscript{ac}; \skt{सिद्विसद्विरिक्ते~।}\ H\textsuperscript{pc}}}\skt{~। }\pend
\pstart\versepar \skt{}\edtext{\skt{मूर्च्छा\-गु\-द\-भ्रं\-श\-क\-फा\-ति\-योगाः}}{\lemma{\skt{मूर्च्छा\appbrk{}गु\appbrk{}द\appbrk{}भ्रं\appbrk{}श\appbrk{}क\appbrk{}फा\appbrk{}ति\appbrk{}योगाः}}\Afootnote{A K; \skt{मूर्च्छा\appbrk{}गु\appbrk{}द\appbrk{}भ्रं\appbrk{}स\appbrk{}क\appbrk{}फा\appbrk{}ति\appbrk{}योगाः}\ H}}\skt{ }\edtext{\skt{शूलो\-द्ग\-मा\-श्चा\-ति\-वि\-रि\-क्त\-लिङ्गं}}{\lemma{\skt{शूलोद्गमाश्}}\Afootnote{K H; \skt{शूलोद्गमश्}\ A}}\skt{~॥}\pend

\pstart\versepar \ledsidenote{\passref{4.33.25}}\skt{गतेषु दोषेषु }\edtext{\skt{कफान्तकेषु}}{\lemma{\skt{कफान्तकेषु}}\Afootnote{K; \skt{कफान्वितेषु}\ A; \skt{कफान्त्त्रकेषु}\ H}}\skt{ }\edtext{\skt{ग्लान्येलघुत्वे}}{\lemma{\skt{ग्लान्येलघुत्वे}}\Afootnote{K H; \skt{नाभ्या लघुत्वे}\ A}}\skt{ मनसश्च तुष्टौ~। }\pend
\pstart\versepar \skt{}\edtext{\skt{गतेनिले}}{\lemma{\skt{गतेनिले}}\Afootnote{K; \skt{गतेऽनिले}\ A H}}\skt{ चाप्य\-नु\-लो\-म\-भावं सम्यग्विरिक्तं }\edtext{\skt{पुरुषं}}{\lemma{\skt{पुरुषं}}\Afootnote{K H; \skt{मनुजं}\ A}}\skt{ व्यवस्येत्~।}\pend

\pstart\versepar \ledsidenote{\passref{4.33.26}}\edtext{\skt{मन्दा\-ग्नि\-म\-क्षी\-ण\-म\-स\-द्वि\-रिक्तं न}}{\lemma{\skt{मन्दा\appbrk{}ग्नि\appbrk{}म\appbrk{}क्षी\appbrk{}ण\appbrk{}म\appbrk{}स\appbrk{}द्वि\appbrk{}रिक्तं न}}\Afootnote{A K; \skt{मन्दा\appbrk{}ग्नि\appbrk{}म\appbrk{}क्षी\appbrk{}ण\appbrk{}म\appbrk{}स\appbrk{}द्वि\appbrk{}रि\appbrk{}क्तेन्न}\ H}}\skt{ }\edtext{\skt{पाय\-ये\-त्त\-द्दि\-वसं}}{\lemma{\skt{पाय\appbrk{}ये\appbrk{}त्त\appbrk{}द्दि\appbrk{}वसं}}\Afootnote{K H; \skt{पाययेताहनि त}\ A}}\skt{ }\edtext{\skt{यवागूं}}{\lemma{\skt{यवागूं}}\Afootnote{K H; \skt{त्र पेयाम्~।}\ A}}\skt{~। }\pend
\pstart\versepar \skt{क्षीणं तृषार्त्तं }\edtext{\skt{सुविरेचितश्च}}{\lemma{\skt{सुविरेचितश्च}}\Afootnote{K; \skt{सुविरेचितं च}\ A; \skt{सुचिरेचितञ्च}\ H}}\skt{ }\edtext{\skt{तन्वी}}{\lemma{\skt{तन्वी}}\Afootnote{K; \skt{तन्वीं}\ A H}}\skt{ }\edtext{\skt{सुखोष्णं}}{\lemma{\skt{सुखोष्णं}}\Afootnote{K H; \skt{सुखोष्णां}\ A}}\skt{ लघु }\edtext{\skt{पाययीत}}{\lemma{\skt{पाययीत}}\Afootnote{K H; \skt{पाययेच्च~॥}\ A}}\skt{~।}\pend

\pstart\versepar \ledsidenote{\passref{4.33.27}}\skt{बुद्धेः प्रसादं }\edtext{\skt{बलमिन्द्रियाणां}}{\lemma{\skt{बलमिन्द्रियाणां}}\Afootnote{K H; \skt{बलमिन्दिर्याणां}\ A}}\skt{ धातुस्थिरत्वं }\edtext{\skt{ज्वलनाभिवृद्धिं}}{\lemma{\skt{ज्वलनाभिवृद्धिं}}\Afootnote{K\textsuperscript{pc} H; \skt{बल\appbrk{}म\appbrk{}ग्नि\appbrk{}दी\appbrk{}प्तिम्~।}\ A; \skt{ज्वलभिवृद्धिं}\ K\textsuperscript{ac}}}\skt{ }\pend
\pstart\versepar \skt{चिराच्च पाकं वयसः करोति विरेचनं सम्य\-गु\-पा\-स्य\-मानं~।}\pend

\pstart\versepar \ledsidenote{\passref{4.33.28}}\edtext{\skt{यथोदकानां}}{\lemma{\skt{यथोदकानां}}\Afootnote{K H; \skt{यथौदकाना}\ A}}\skt{ }\edtext{\skt{तु विजङ्गमानां}}{\lemma{\skt{तु विजङ्गमानां}}\Afootnote{K; \skt{मुदकेऽ}\ A; \skt{सुविजङ्गमानां}\ H}}\skt{ }\edtext{\skt{रोधेवदीर्ण्णे}}{\lemma{\skt{रोधेवदीर्ण्णे}}\Afootnote{K H; \skt{पनीते}\ A}}\skt{ }\edtext{\skt{ध्रुव}}{\lemma{\skt{ध्रुव}}\Afootnote{K H\textsuperscript{ac}; \skt{चरस्थिराणां भव}\ A; \skt{ध्रुव एव}\ H\textsuperscript{pc}}}\skt{ }\edtext{\skt{एव}}{\lemma{\skt{एव}}\Afootnote{K H\textsuperscript{ac}; \skt{ति प्र}\ A; \emph{om.}\ H\textsuperscript{pc}}}\skt{ }\edtext{\skt{नाशः}}{\lemma{\skt{नाशः}}\Afootnote{K H; \skt{णाशः~।}\ A}}\skt{~। }\pend
\pstart\versepar \skt{पित्ते }\edtext{\skt{हृते\-न्धे\-व\-मु\-प\-द्र\-वाणां}}{\lemma{\skt{हृतेन्धेवम्}}\Afootnote{K H; \skt{हृते त्वेवम्}\ A}}\skt{ }\edtext{\skt{पित्तात्मिकाकानां}}{\lemma{\skt{पित्तात्मिकाकानां}}\Afootnote{K; \skt{पित्तात्मकानां भ}\ A; \skt{पित्तात्मिनाम्}\ H\textsuperscript{ac}; \skt{पित्तात्मिकानाम्}\ H\textsuperscript{pc}}}\skt{ }\edtext{\skt{विहितो}}{\lemma{\skt{विहितो}}\Afootnote{K H; \skt{वति प्र}\ A}}\skt{ }\edtext{\skt{विनाशः}}{\lemma{\skt{विनाशः}}\Afootnote{K H; \skt{णाशः~॥}\ A}}\skt{~॥}\pend

\pstart\versepar \ledsidenote{\passref{4.33.29}}\edtext{\skt{क्षीणाः\-कृ\-शा\-रू\-क्षि\-त\-बा\-ल\-वृ\-द्धा\-दी\-नोथ}}{\lemma{\skt{क्षीणाः\appbrk{}कृ\appbrk{}शा\appbrk{}रू\appbrk{}क्षि\appbrk{}त\appbrk{}बा\appbrk{}ल\appbrk{}वृ\appbrk{}द्धा\appbrk{}दी\appbrk{}नोथ}}\Afootnote{K; \skt{मन्दा\appbrk{}ग्न्य\appbrk{}ति\appbrk{}स्ने\appbrk{}हि\appbrk{}त\appbrk{}बा\appbrk{}ल\appbrk{}वृ\appbrk{}द्ध\appbrk{}स्थूलाः क्ष}\ A; \skt{क्षीणा\appbrk{}कृ\appbrk{}शा\appbrk{}रू\appbrk{}क्षि\appbrk{}त\appbrk{}बा\appbrk{}ल\appbrk{}वृ\appbrk{}द्धा\appbrk{}दी\appbrk{}नोथ}\ H}}\skt{ }\edtext{\skt{शोषी}}{\lemma{\skt{शोषी}}\Afootnote{K H; \skt{तक्षी}\ A}}\skt{ }\edtext{\skt{भयशोकतप्तः}}{\lemma{\skt{भयशोकतप्तः}}\Afootnote{K H; \skt{णभयोपतप्ताः~।}\ A}}\skt{~। }\pend
\pstart\versepar \skt{}\edtext{\edtext{\skt{शान्तस्तृषार्त्तो}}{\lemma{\skt{शान्तस्}}\Afootnote{K; \skt{श्रान्त}\ A; \skt{शान्त}\ H}}}{\lemma{\skt{तृषार्त्तो}}\Afootnote{K; \skt{स्तृषार्तो}\ A H}}\skt{ परिजीर्ण्णभक्तो गर्भिण्यधो गच्छति यस्य चासृक्~।}\pend

\pstart\versepar \ledsidenote{\passref{4.33.30}}\edtext{\skt{नव\-प्र\-ति\-श्या\-य\-प\-री\-त\-देहो}}{\lemma{\skt{नव\appbrk{}प्र\appbrk{}ति\appbrk{}श्या\appbrk{}य\appbrk{}प\appbrk{}री\appbrk{}त\appbrk{}देहो}}\Afootnote{K; \skt{नव\appbrk{}प्र\appbrk{}वि\appbrk{}श्या\appbrk{}य\appbrk{}म\appbrk{}दा\appbrk{}त्ययी च}\ A; \skt{नवप्रविश्याय परीतदेहो}\ H}}\skt{ }\edtext{\skt{नवज्वरीया च}}{\lemma{\skt{नवज्वरीया च}}\Afootnote{K H; \skt{नवज्रवी या च}\ A}}\skt{ नवप्रसूता~। }\pend
\pstart\versepar \skt{}\edtext{\edtext{\skt{कषायनित्यान}}{\lemma{\skt{कषायनित्य्}}\Afootnote{\skt{शल्यार्दिताश्चा}\ A; \skt{कषायनित्या}\ K H}}}{\lemma{\skt{आन}}\Afootnote{\skt{प्य}\ A; \skt{न}\ K H}}\skt{ }\edtext{\skt{विरेचनीया}}{\lemma{\skt{विरेचनीया}}\Afootnote{K H\textsuperscript{pc}; \skt{विरेचनीयाः}\ A; \skt{विरेनीया}\ H\textsuperscript{ac}}}\skt{ स्नेहादिभिर्ये त्वनुपस्कृताश्च~।}\pend

\pstart\versepar \ledsidenote{\passref{4.33.31}}\edtext{\skt{अत्य\-र्थ\-पि\-त्ता\-भि\-प\-री\-त\-देहां}}{\lemma{\skt{अत्य\appbrk{}र्थ\appbrk{}पि\appbrk{}त्ता\appbrk{}भि\appbrk{}प\appbrk{}री\appbrk{}त\appbrk{}देहां}}\Afootnote{K; \skt{अत्य\appbrk{}र्थ\appbrk{}पि\appbrk{}त्ता\appbrk{}भि\appbrk{}प\appbrk{}री\appbrk{}त\appbrk{}देहा}\ H}}\skt{ }\edtext{\skt{विरे\-च\-ये\-त्ता\-नपि}}{\lemma{\skt{विरे\appbrk{}च\appbrk{}ये\appbrk{}त्ता\appbrk{}नपि}}\Afootnote{K; \skt{न्विरे\appbrk{}च\appbrk{}ये\appbrk{}त्ता\appbrk{}नपि}\ H}}\skt{ मन्दमन्दं~। }\pend
\pstart\versepar \skt{विरेचनैर्यांति नरा }\edtext{\edtext{\edtext{\skt{विना\-श\-म\-ज्ञ\-प्र\-यु\-क्तै\-र\-वि\-रे\-च\-नीया}}{\lemma{\skt{विनाशम्}}\Afootnote{K H; \skt{विनाश}\ A}}}{\lemma{\skt{अज्ञप्रयुक्तैर्}}\Afootnote{K H; \skt{मज्ञप्रयुक्तैर्}\ A}}}{\lemma{\skt{अविरेचनीया}}\Afootnote{K; \skt{अविरेचनीयाः~॥}\ A H}}\skt{~॥}\edtext{}{\lemma{\skt{… वि}}\Afootnote{\emph{om.}\ A}}\pend

\pstart\versepar \ledsidenote{\passref{4.33.33}}\edtext{\skt{सर\-त्व\-सौ\-क्ष्म्य\-तै\-क्ष्ण्योष्णे}}{\lemma{\skt{सर\appbrk{}त्व\appbrk{}सौ\appbrk{}क्ष्म्य\appbrk{}तै\appbrk{}क्ष्ण्योष्णे}}\Afootnote{K; \skt{सर\appbrk{}त्व\appbrk{}सौ\appbrk{}क्ष्म्य\appbrk{}तै\appbrk{}क्ष्ण्यौष्ण्य}\ A; \skt{सर\appbrk{}त्व\appbrk{}सौ\appbrk{}क्ष्म्य\appbrk{}तै\appbrk{}क्ष्ण्यौष्णे}\ H}}\skt{ }\edtext{\skt{विका\-शि\-त्वा\-द्वि\-रे\-चनं}}{\lemma{\skt{विका\appbrk{}शि\appbrk{}त्वा\appbrk{}द्वि\appbrk{}रे\appbrk{}चनं}}\Afootnote{K H; \skt{विका\appbrk{}शि\appbrk{}त्वै\appbrk{}र्वि\appbrk{}रे\appbrk{}च\appbrk{}नम्~।}\ A}}\skt{~। }\pend
\pstart\versepar \skt{वमनं तु }\edtext{\skt{हरेद्दोषां}}{\lemma{\skt{हरेद्दोषां}}\Afootnote{K; \skt{हरेद्दोषं}\ A; \skt{हरेद्दोषा}\ H}}\skt{ }\edtext{\skt{सम्यगुक्तं}}{\lemma{\skt{सम्यगुक्तं}}\Afootnote{K; \skt{प्रकृत्या ग}\ A; \skt{न्सम्यगुक्तम्}\ H\textsuperscript{ac}; \skt{न्सम्यगुग्युक्तम्}\ H\textsuperscript{pc}}}\skt{ }\edtext{\skt{वृथान्यथा}}{\lemma{\skt{वृथान्यथा}}\Afootnote{K H; \skt{तमन्यथा~॥}\ A}}\skt{~।}\pend

\pstart\versepar \ledsidenote{\passref{4.33.34}}\edtext{\skt{मादाय}}{\lemma{\skt{मादाय}}\Afootnote{\skt{यात्यधो दोषमादाय}\ A; \skt{यात्य\appbrk{}ध\appbrk{}र्द्दो\appbrk{}ष\appbrk{}मा\appbrk{}दाय}\ K H}}\skt{ }\edtext{\skt{प्यमानं}}{\lemma{\skt{प्यमानं}}\Afootnote{K; \skt{पच्यमानं}\ A H}}\skt{ विरेचनं~। }\pend
\pstart\versepar \skt{}\edtext{\edtext{\skt{गुणो\-द्रे\-का\-द्व्र\-जे\-दू\-र्ध्व\-म\-पक्वं}}{\lemma{\skt{गुणो\appbrk{}द्रे\appbrk{}का\appbrk{}द्व्र\appbrk{}जे\appbrk{}दू\appbrk{}र्ध्वम्}}\Afootnote{K H; \skt{गुणो\appbrk{}त्क\appbrk{}र्षा\appbrk{}ह्व्र\appbrk{}ज\appbrk{}त्यूर्ध्व}\ A}}}{\lemma{\skt{अपक्वं}}\Afootnote{K H; \skt{मपक्वं}\ A}}\skt{ वमनं पुनः~।}\pend

\pstart\versepar \ledsidenote{\passref{4.33.35}}\edtext{\skt{मृदुकोष्ठस्य}}{\lemma{\skt{मृदुकोष्ठस्य}}\Afootnote{A K; \skt{मृदुकोष्णस्य}\ H}}\skt{ }\edtext{\edtext{\skt{दीप्ताग्नेर्दत्तं}}{\lemma{\skt{दीप्ताग्नेर्}}\Afootnote{K H; \skt{दीप्ताग्ने}\ A}}}{\lemma{\skt{दत्तं}}\Afootnote{K H; \skt{रति}\ A}}\skt{ तीक्ष्णं विरेचनं~। }\pend
\pstart\versepar \skt{न }\edtext{\skt{सम्य\-ग्नि\-र्ह\-रे\-द्दो\-षा\-न\-ति\-वे\-ग\-प्र\-धा\-वितं}}{\lemma{\skt{सम्य\appbrk{}ग्नि\appbrk{}र्ह\appbrk{}रे\appbrk{}द्दो\appbrk{}षा\appbrk{}न\appbrk{}ति\appbrk{}वे\appbrk{}ग\appbrk{}प्र\appbrk{}धा\appbrk{}वितं}}\Afootnote{K; \skt{सम्य\appbrk{}ङ्नि\appbrk{}र्ह\appbrk{}रे\appbrk{}द्दो\appbrk{}षा\appbrk{}न\appbrk{}ति\appbrk{}वे\appbrk{}ग\appbrk{}प्र\appbrk{}धा\appbrk{}वि\appbrk{}तम्~॥}\ A; \skt{सम्य\appbrk{}ङ्नि\appbrk{}र्ह\appbrk{}रे\appbrk{}द्दो\appbrk{}षा\appbrk{}न\appbrk{}ति\appbrk{}वे\appbrk{}ग\appbrk{}प्र\appbrk{}धा\appbrk{}वितं~॥}\ H\textsuperscript{ac}; \skt{सम्य\appbrk{}ङ्नि\appbrk{}र्ह\appbrk{}रे\appbrk{}द्दो\appbrk{}षा\appbrk{}न्स\appbrk{}न्त्य\appbrk{}ग्नि\appbrk{}र्द्दहे इति पाठः~। अति\appbrk{}वे\appbrk{}ग\appbrk{}प्र\appbrk{}धा\appbrk{}वितं~॥}\ H\textsuperscript{pc}}}\skt{~।}\pend

\pstart\versepar \ledsidenote{\passref{4.33.36}}\edtext{\edtext{\skt{प्रातर्यदौषधं}}{\lemma{\skt{प्रातर्}}\Afootnote{K H; \skt{पीतं}\ A}}}{\lemma{\skt{यदौषधं}}\Afootnote{K H; \skt{यदौषधं}\ A}}\skt{ }\edtext{\skt{पीतं}}{\lemma{\skt{पीतं}}\Afootnote{K H; \skt{प्रातर्}\ A}}\skt{ }\edtext{\skt{भक्तपाकसमे}}{\lemma{\skt{भक्तपाकसमे}}\Afootnote{K H; \skt{भुक्तपाकसमे}\ A}}\skt{ क्षणे~। }\pend
\pstart\versepar \skt{}\edtext{\skt{पक्तं}}{\lemma{\skt{पक्तं}}\Afootnote{K; \skt{पक्तिं}\ A H}}\skt{ गच्छति }\edtext{\skt{दोषांस्तु}}{\lemma{\skt{दोषांस्तु}}\Afootnote{K H; \skt{दोषांश्च}\ A}}\skt{ }\edtext{\edtext{\skt{निर्ह\-रे\-न्नि\-र\-व\-शे\-षतः}}{\lemma{\skt{निर्हरेन्}}\Afootnote{K; \skt{निर्हरे}\ A H}}}{\lemma{\skt{निरवशेषतः}}\Afootnote{K; \skt{त्तत्प्रशस्यते~॥}\ A; \skt{दवशेषतः~॥}\ H}}\skt{~॥}\pend

\pstart\versepar \ledsidenote{\passref{4.33.37}}\edtext{\skt{दुर्बलस्य}}{\lemma{\skt{दुर्बलस्य}}\Afootnote{A H; \skt{दुर्वलस्य}\ K}}\skt{ }\edtext{\skt{नः}}{\lemma{\skt{नः}}\Afootnote{\skt{चला\appbrk{}न्दो\appbrk{}षा\appbrk{}न\appbrk{}ल्पा\appbrk{}न\appbrk{}ल्पा\appbrk{}न्पुनः}\ A K; \skt{चला\appbrk{}न्दो\appbrk{}षा\appbrk{}म\appbrk{}ल्पा\appbrk{}न\appbrk{}ल्पा\appbrk{}न्पुनः}\ H}}\skt{ पुनः~। }\pend
\pstart\versepar \skt{}\edtext{\edtext{\skt{हरे\-त्प्र\-भू\-ता\-न\-ल्पांस्तु}}{\lemma{\skt{हरे\appbrk{}त्प्र\appbrk{}भू\appbrk{}ता\appbrk{}न\appbrk{}ल्पांस्}}\Afootnote{K; \skt{हरेत्प्रभूतान}\ A; \skt{प्रभूतानल्पतां}\ H}}}{\lemma{\skt{तु}}\Afootnote{K; \skt{क्पांस्तु}\ A; \skt{यान्तु}\ H}}\skt{ }\edtext{\skt{सम\-ये\-त्प्र\-च्यु\-ता\-नथ}}{\lemma{\skt{सम\appbrk{}ये\appbrk{}त्प्र\appbrk{}च्यु\appbrk{}ता\appbrk{}नथ}}\Afootnote{K H; \skt{शम\appbrk{}ये\appbrk{}त्प्र\appbrk{}च्यु\appbrk{}ता\appbrk{}नपि~॥}\ A}}\skt{~।}\pend

\pstart\versepar \ledsidenote{\passref{4.33.38}}\edtext{\skt{हरेद्दोषां}}{\lemma{\skt{हरेद्दोषां}}\Afootnote{K H; \skt{हरेद्दोषांश्}\ A}}\skt{ }\edtext{\skt{चलान्पक्वं}}{\lemma{\skt{चलान्पक्वं}}\Afootnote{K H; \skt{चलान्पक्वा}\ A}}\skt{ }\edtext{\skt{बलिनो}}{\lemma{\skt{बलिनो}}\Afootnote{K H; \skt{न्बलिनो}\ A}}\skt{ दुर्बलस्य वा~। }\pend
\pstart\versepar \skt{}\edtext{\skt{चला }\edtext{\skt{ह्युपेक्षिता}}{\lemma{\skt{उपेक्षिता}}\Afootnote{A K; \skt{ह्युपेक्षिता}\ H}}}{\lemma{\skt{चला ह्य्}}\Afootnote{A K; \skt{चला}\ H}}\skt{ दोषाः }\edtext{\edtext{\skt{क्लेश\-ये\-यु\-श्चि\-र\-न्नरं}}{\lemma{\skt{क्लेशयेयुश्}}\Afootnote{K H; \skt{क्लेशयेयु}\ A}}}{\lemma{\skt{चिरन्}}\Afootnote{K H; \skt{श्चिरं}\ A}}\skt{~।}\pend

\pstart\versepar \ledsidenote{\passref{4.33.39}}\skt{मन्दाग्निं क्रूरकोष्ठञ्च }\edtext{\skt{सक्षा\-र\-ल\-व\-णै\-र्घृतैः}}{\lemma{\skt{सक्षा\appbrk{}र\appbrk{}ल\appbrk{}व\appbrk{}णै\appbrk{}र्घृतैः}}\Afootnote{K; \skt{सक्षा\appbrk{}र\appbrk{}ल\appbrk{}व\appbrk{}णै\appbrk{}र्घृतैः~।}\ A; \skt{सक्षा\appbrk{}र\appbrk{}ल\appbrk{}व\appbrk{}णै\appbrk{}र्घृतं~।}\ H}}\skt{~। }\pend
\pstart\versepar \skt{}\edtext{\skt{सन्धुक्षिताग्नि}}{\lemma{\skt{सन्धुक्षिताग्नि}}\Afootnote{K; \skt{सन्धुक्षिताग्निं}\ A H}}\skt{ स्निग्धञ्च स्विन्नञ्चैव }\edtext{\skt{विशोधयेत्}}{\lemma{\skt{विशोधयेत्}}\Afootnote{K H; \skt{विरेचयेत्~॥}\ A}}\skt{~।}\pend

\pstart\versepar \ledsidenote{\passref{4.33.40}}\skt{स्निग्धस्विन्नस्य }\edtext{\skt{भैषज्यैर्दोषस्तु}}{\lemma{\skt{दोषस्तु}}\Afootnote{K H; \skt{दोषस्तू}\ A}}\skt{ }\edtext{\skt{क्लेशितो}}{\lemma{\skt{क्लेशितो}}\Afootnote{K H; \skt{त्क्लेशितो}\ A}}\skt{ भवेत्~। }\pend
\pstart\versepar \skt{न }\edtext{\skt{चालीयेत}}{\lemma{\skt{चालीयेत}}\Afootnote{K H; \skt{लीयते न}\ A}}\skt{ मार्गेषु स्निग्धे भाण्ड इवोदकं~।}\edtext{}{\lemma{\skt{भवेत्~। न}}\Afootnote{K H; \skt{बलात्~। नि}\ A}}\pend

\pstart\versepar \ledsidenote{\passref{4.33.41}}\edtext{\skt{न चाति\-स्नि\-ग्ध\-का\-याय}}{\lemma{\skt{न चाति\appbrk{}स्नि\appbrk{}ग्ध\appbrk{}का\appbrk{}याय}}\Afootnote{K H; \skt{न चाति\appbrk{}स्ने\appbrk{}ह\appbrk{}पी\appbrk{}तस्तु}\ A}}\skt{ }\edtext{\skt{दद्या\-त्स्ने\-र्ह\-वि\-रे\-चनं}}{\lemma{\skt{दद्या\appbrk{}त्स्ने\appbrk{}र्ह\appbrk{}वि\appbrk{}रे\appbrk{}चनं}}\Afootnote{K; \skt{पिबे\appbrk{}त्स्ने\appbrk{}ह\appbrk{}वि\appbrk{}रे\appbrk{}च\appbrk{}नम्~।}\ A; \skt{दद्या\appbrk{}त्स्ने\appbrk{}ह\appbrk{}वि\appbrk{}रे\appbrk{}चनं~।}\ H}}\skt{~। }\pend
\pstart\versepar \skt{दोषाः }\edtext{\skt{प्रच्याविता}}{\lemma{\skt{प्रच्याविता}}\Afootnote{K H; \skt{प्रचलिताः स्थाना}\ A}}\skt{ }\edtext{\skt{भूयो}}{\lemma{\skt{भूयो}}\Afootnote{K H; \skt{द्भूयः}\ A}}\skt{ }\edtext{\skt{लीयन्ते}}{\lemma{\skt{लीयन्ते}}\Afootnote{K H; \skt{श्लिष्यन्}\ A}}\skt{ }\edtext{\skt{तेन}}{\lemma{\skt{तेन}}\Afootnote{K H; \skt{ति}\ A}}\skt{ }\edtext{\skt{वर्त्मसु}}{\lemma{\skt{वर्त्मसु}}\Afootnote{A K; \skt{वर्त्म्मसु~॥}\ H}}\skt{~।}\pend

\pstart\versepar \ledsidenote{\passref{4.33.43}}\skt{विरूक्ष्य स्नेहसात्म्यन्तु भूयः संस्नेह्य रेचयेत्~। }\pend
\pstart\versepar \skt{तेन दोषा }\edtext{\edtext{\skt{हृतास्तस्य}}{\lemma{\skt{हृतास्}}\Afootnote{A; \skt{हृता}\ K H}}}{\lemma{\skt{तस्य}}\Afootnote{A; \skt{स्तस्य}\ K H}}\skt{ भवन्ति }\edtext{\skt{स्नेहबन्धनाः}}{\lemma{\skt{स्नेहबन्धनाः}}\Afootnote{K H; \skt{बलवर्धनाः~॥}\ A}}\skt{~।}\edtext{}{\lemma{\skt{रेचयेत्~। तेन}}\Afootnote{K H; \skt{शोधयेत्~। तेन}\ A}}\pend

\pstart\versepar \ledsidenote{\passref{4.33.44}}\edtext{\skt{प्रागधीतौषधं}}{\lemma{\skt{प्रागधीतौषधं}}\Afootnote{K; \skt{प्रागपीतं नरं}\ A; \skt{प्रागपीतौषधं}\ H}}\skt{ }\edtext{\skt{शोध्य}}{\lemma{\skt{शोध्य}}\Afootnote{K H; \skt{शोध्यं}\ A}}\skt{ }\edtext{\skt{पाय\-ये\-त्मृ\-दु\-शो\-धणं}}{\lemma{\skt{पाय\appbrk{}ये\appbrk{}त्मृ\appbrk{}दु\appbrk{}शो\appbrk{}धणं}}\Afootnote{K; \skt{पाययेतौषधं मृदु~।}\ A; \skt{पाययेत्मृद शोधनं~।}\ H}}\skt{~। }\pend
\pstart\versepar \skt{ततो विज्ञातकोष्ठस्य कार्यं }\edtext{\skt{संशोधणं}}{\lemma{\skt{संशोधणं}}\Afootnote{K; \skt{संशोधनं}\ A H}}\skt{ पुनः~।}\pend

\pstart\versepar \ledsidenote{\passref{4.33.45}}\skt{सुखं }\edtext{\skt{दृढफलं}}{\lemma{\skt{दृढफलं}}\Afootnote{K H; \skt{दृष्टफलं}\ A}}\skt{ }\edtext{\edtext{\skt{हृद्यमल्पमात्रं}}{\lemma{\skt{हृद्यम्}}\Afootnote{K H; \skt{हृद्य}\ A}}}{\lemma{\skt{अल्पमात्रं}}\Afootnote{K H; \skt{मल्पमात्रं}\ A}}\skt{ महागुणं~। }\pend
\pstart\versepar \skt{व्याप\-त्स्व\-ल्पा\-त्य\-य\-ञ्चापि }\edtext{\skt{पिबे\-न्नृ\-प\-ति\-रौ\-षधं}}{\lemma{\skt{नृपतिरौषधं}}\Afootnote{\skt{नृप\appbrk{}ति\appbrk{}रौ\appbrk{}ष\appbrk{}धम्}\ A K H}}\skt{~॥}\pend

\pstart\versepar \ledsidenote{\passref{4.33.46}}\edtext{\skt{स्नेह\-स्वे\-दा\-व\-कृ\-त्वाग्रे यस्तु}}{\lemma{\skt{स्नेह\appbrk{}स्वे\appbrk{}दा\appbrk{}व\appbrk{}कृ\appbrk{}त्वाग्रे यस्तु}}\Afootnote{K H; \skt{स्नेह\appbrk{}स्वे\appbrk{}दा\appbrk{}व\appbrk{}न\appbrk{}भ्यस्य यस्तु}\ A}}\skt{ संशोधनं पिबेत्~। }\pend
\pstart\versepar \skt{दारु शुष्कमिवानामे }\edtext{\edtext{\skt{देहस्तस्य}}{\lemma{\skt{देहस्}}\Afootnote{A; \skt{देह}\ K H}}}{\lemma{\skt{तस्य}}\Afootnote{A; \skt{स्तस्य}\ K H}}\skt{ विशीर्यते~॥}\pend

\pstart\versepar \ledsidenote{\passref{4.33.47}}\edtext{\skt{स्नेह\-स्वे\-द\-प्र\-च\-लितो}}{\lemma{\skt{स्नेह\appbrk{}स्वे\appbrk{}द\appbrk{}प्र\appbrk{}च\appbrk{}लितो}}\Afootnote{K; \skt{स्नेह\appbrk{}स्वे\appbrk{}द\appbrk{}प्र\appbrk{}च\appbrk{}लिता}\ A; \skt{स्नेह\appbrk{}स्वे\appbrk{}द\appbrk{}प्र\appbrk{}व\appbrk{}लितो}\ H\textsuperscript{ac}; \skt{स्नेह\appbrk{}स्वे\appbrk{}द\appbrk{}प्र\appbrk{}च\appbrk{}व\appbrk{}लितो}\ H\textsuperscript{pc}}}\skt{ रसैः स्निग्धैरुदीरिता~। }\pend
\pstart\versepar \skt{दोषाः कोष्ठगता जन्तोः }\edtext{\skt{सुखं}}{\lemma{\skt{सुखं}}\Afootnote{K H; \skt{सुखा}\ A}}\skt{ हर्तुं }\edtext{\skt{विरेचनैरिति}}{\lemma{\skt{विरेचनैरिति}}\Afootnote{K; \skt{विशोधनैः~॥}\ A; \skt{विरचनैरिति~॥}\ H}}\skt{~॥}\pend

\endnumbering
\beginnumbering
\eledsection*{Adhyāya 34}
\pstart \ledsidenote{\passref{4.34.1}}\skt{अथातो }\edtext{\edtext{\skt{वम\-न\-वि\-रे\-क\-व्या\-प\-च्चि\-कि\-त्सि\-त\-म्व्या\-ख्या\-स्यामः}}{\lemma{\skt{वम\appbrk{}न\appbrk{}वि\appbrk{}रे\appbrk{}क\appbrk{}व्या\appbrk{}प\appbrk{}च्चि\appbrk{}कि\appbrk{}त्सि\appbrk{}तम्}}\Afootnote{K H; \skt{वम\appbrk{}न\appbrk{}वि\appbrk{}रे\appbrk{}च\appbrk{}न\appbrk{}व्या\appbrk{}प\appbrk{}च्चि\appbrk{}कि\appbrk{}त्सितं}\ A}}}{\lemma{\skt{व्याख्यास्यामः}}\Afootnote{A H; \emph{om.}\ K}}\skt{~।}\pend

\pstart \ledsidenote{\passref{4.34.3}}\edtext{\skt{वैद्या\-तु\-र\-नि\-मि\-त्त\-म्व\-म\-न\-म्वि\-रे\-च\-नञ्च}}{\lemma{\skt{विरेचनञ्च}}\Afootnote{A K H\textsuperscript{pc}; \skt{विरेचनञ्}\ H\textsuperscript{ac}}}\skt{~॥ }\edtext{\skt{पञ्चदशधा}}{\lemma{\skt{पञ्चदशधा}}\Afootnote{A K H\textsuperscript{pc}; \skt{च~॥}\ H\textsuperscript{ac}}}\skt{ }\edtext{\skt{व्यापद्यते}}{\lemma{\skt{व्यापद्यते}}\Afootnote{A K H\textsuperscript{pc}; \emph{om.}\ H\textsuperscript{ac}}}\skt{ तत्र वमनस्याधो }\edtext{\edtext{\edtext{\skt{गच्छ\-त्यू\-र्द्ध्व\-म्वि\-रे\-च\-नस्य}}{\lemma{\skt{गच्छत्य्}}\Afootnote{K H; \skt{गति}\ A}}}{\lemma{\skt{ऊर्द्ध्वम्}}\Afootnote{K H; \skt{रूर्ध्वं}\ A}}}{\lemma{\skt{विरेचनस्य}}\Afootnote{K H; \skt{विरेचनस्येति}\ A}}\skt{~। }\edtext{\edtext{\skt{सामान्यमुभयोः}}{\lemma{\skt{सामान्यम्}}\Afootnote{K H; \skt{पृथक्; सामान्य}\ A}}}{\lemma{\skt{उभयोः}}\Afootnote{K H; \skt{मुभयोः}\ A}}\skt{ }\edtext{\edtext{\skt{साव\-शे\-षौ\-ष\-ध\-त्व\-म\-ल्प\-दो\-ष\-हृ\-तत्वं}}{\lemma{\skt{साव\appbrk{}शे\appbrk{}षौ\appbrk{}ष\appbrk{}ध\appbrk{}त्वम्}}\Afootnote{K H; \skt{सावशेषौधत्वं, जीर्णौषधत्वं,}\ A}}}{\lemma{\skt{अल्प\appbrk{}दो\appbrk{}ष\appbrk{}हृ\appbrk{}तत्वं}}\Afootnote{K H; \skt{हीन\appbrk{}दो\appbrk{}षा\appbrk{}प\appbrk{}हृ\appbrk{}तत्वं,}\ A}}\skt{ }\edtext{\edtext{\skt{जीर्ण्णौ\-ष\-ध\-म्वा\-त\-शू\-ल\-म\-यो\-गा\-ति\-योगौ}}{\lemma{\skt{जीर्ण्णौषधम्}}\Afootnote{K H; \emph{om.}\ A}}}{\lemma{\skt{अयोगातियोगौ}}\Afootnote{K H; \skt{अयोगो, अतियोगो,}\ A}}\skt{ }\edtext{\edtext{\skt{जीवा\-दा\-न\-मा\-ध्मा\-न\-म्प\-रि\-कर्त्तः}}{\lemma{\skt{जीवा\appbrk{}दा\appbrk{}न\appbrk{}मा\appbrk{}ध्मा\appbrk{}नम्}}\Afootnote{H; \skt{जीवा\appbrk{}दा\appbrk{}न\appbrk{}मा\appbrk{}ध्मानं}\ A K}}}{\lemma{\skt{परिकर्त्तः}}\Afootnote{K H; \skt{परिकर्तिका,}\ A}}\skt{ परिस्रावः प्रवाहिका }\edtext{\skt{हृदयसरण}}{\lemma{\skt{हृदयसरण}}\Afootnote{H; \skt{हृद\appbrk{}यो\appbrk{}प\appbrk{}स\appbrk{}रणं,}\ A; \skt{हृदयसरणं}\ K}}\skt{ विबन्ध }\edtext{\skt{इति}}{\lemma{\skt{इति}}\Afootnote{K H; \skt{अङ्गप्रग्रह इति~॥}\ A}}\skt{~।}\pend

\pstart \ledsidenote{\passref{4.34.4}}\edtext{\skt{तत्र}}{\lemma{\skt{तत्र}}\Afootnote{A H; \skt{तत्र च}\ K}}\skt{ }\edtext{\skt{बुभु\-क्षा\-पी\-डि\-तस्य}}{\lemma{\skt{बुभु\appbrk{}क्षा\appbrk{}पी\appbrk{}डि\appbrk{}तस्य}}\Afootnote{K H; \skt{बुभु\appbrk{}क्षा\appbrk{}पी\appbrk{}डि\appbrk{}तस्या}\ A}}\skt{ }\edtext{\edtext{\skt{तीक्ष्णा\-ग्ने\-र्म्मृ\-दु\-को\-ष्ठस्य}}{\lemma{\skt{तीक्ष्णाग्नेर्}}\Afootnote{K H; \skt{तितीक्ष्णाग्ने}\ A}}}{\lemma{\skt{म्मृदुकोष्ठस्य}}\Afootnote{K H; \skt{र्मृदुकोष्ठस्य}\ A}}\skt{ }\edtext{\skt{वावतिष्ठमानं}}{\lemma{\skt{वावतिष्ठमानं}}\Afootnote{K H; \skt{चावतिष्ठमानं}\ A}}\skt{ }\edtext{\skt{दुर्व्वान्तस्य वा}}{\lemma{\skt{दुर्व्वान्तस्य वा}}\Afootnote{K H; \skt{दुर्वमस्य वा}\ A}}\skt{ }\edtext{\edtext{\skt{गुण\-सा\-मा\-न्या\-द्व\-म\-न\-मधो}}{\lemma{\skt{गुण\appbrk{}सा\appbrk{}मा\appbrk{}न्या\appbrk{}द्व\appbrk{}म\appbrk{}नम्}}\Afootnote{K H; \skt{गुण\appbrk{}सा\appbrk{}मा\appbrk{}न्य\appbrk{}भा\appbrk{}वा\appbrk{}द्व\appbrk{}मन}\ A}}}{\lemma{\skt{अधो}}\Afootnote{K H; \skt{मधो}\ A}}\skt{ गच्छति~। }\edtext{\edtext{\skt{तत्रे\-प्सि\-ता\-न\-वा\-प्ति\-र्द्दोषो}}{\lemma{\skt{तत्रे\appbrk{}प्सि\appbrk{}ता\appbrk{}न\appbrk{}वा\appbrk{}प्तिर्}}\Afootnote{H; \skt{तत्रे\appbrk{}प्सि\appbrk{}ता\appbrk{}न\appbrk{}वाप्ति}\ A; \skt{तत्रे\appbrk{}प्ति\appbrk{}सा\appbrk{}न\appbrk{}वा\appbrk{}प्तिर्}\ K}}}{\lemma{\skt{द्दोषो}}\Afootnote{K H; \skt{र्दोषो}\ A}}\skt{ }\edtext{\skt{क्लेशश्च }\edtext{\skt{तमा\-शु\-स्ने\-ह\-यित्वा}}{\lemma{\skt{आशुस्नेहयित्वा}}\Afootnote{K H; \skt{माशु स्नेहयित्वा}\ A}}}{\lemma{\skt{क्लेशश्च तम्}}\Afootnote{K H; \skt{त्क्लेशश्च; त}\ A}}\skt{ }\edtext{\edtext{\skt{भूय\-स्ती\-क्ष्ण\-त\-रै\-र्व्वा\-म\-ये\-तिति}}{\lemma{\skt{भूय\appbrk{}स्ती\appbrk{}क्ष्ण\appbrk{}त\appbrk{}रैर्}}\Afootnote{A H; \skt{भूय\appbrk{}स्ती\appbrk{}क्ष\appbrk{}त\appbrk{}रैर्}\ K}}}{\lemma{\skt{व्वामयेतिति}}\Afootnote{K H; \skt{वामयेत्~॥}\ A}}\skt{~॥}\pend

\pstart \ledsidenote{\passref{4.34.5}}\skt{अप\-रि\-शु\-द्धा\-मा\-श\-य\-स्यो\-त्क्लि\-ष्ट\-श्ले\-ष्मणः }\edtext{\skt{सावशेषान्नस्य}}{\lemma{\skt{सावशेषान्}}\Afootnote{K H; \skt{सशेषान्}\ A}}\skt{ }\edtext{\edtext{\skt{वाऽहृ\-द्य\-म्ब\-हु\-धा\-व\-चा\-रि\-तम्वा}}{\lemma{\skt{वाऽहृद्यम्}}\Afootnote{K H; \skt{वाऽहृद्य}\ A}}}{\lemma{\skt{बहु\appbrk{}धा\appbrk{}व\appbrk{}चा\appbrk{}रि\appbrk{}तम्वा}}\Afootnote{K H; \skt{मतिप्रभूतं वा}\ A}}\skt{ }\edtext{\edtext{\skt{विरे\-च\-न\-मू\-र्द्ध\-म्ति\-ष्ठति}}{\lemma{\skt{ऊर्द्धम्}}\Afootnote{H; \skt{पीतमूर्ध्वं}\ A; \skt{ऊर्ध्वम्}\ K}}}{\lemma{\skt{तिष्ठति}}\Afootnote{K H; \skt{गच्छति,}\ A}}\skt{ }\edtext{\skt{तत्रे\-प्सि\-ता\-न\-वा\-प्ति\-ब\-स्ति\-दो\-षा\-न्क्ले\-श\-त\-मु\-ल्ब\-ण\-श्ले\-ष्मा\-ण\-माशु}}{\lemma{\skt{तत्रे\appbrk{}प्सि\appbrk{}ता\appbrk{}न\appbrk{}वा\appbrk{}प्ति\appbrk{}ब\appbrk{}स्ति\appbrk{}दो\appbrk{}षा\appbrk{}न्क्ले\appbrk{}श\appbrk{}त\appbrk{}मु\appbrk{}ल्ब\appbrk{}ण\appbrk{}श्ले\appbrk{}ष्मा\appbrk{}ण\appbrk{}माशु}}\Afootnote{H\textsuperscript{pc}; \skt{तत्रे\appbrk{}प्सि\appbrk{}ता\appbrk{}न\appbrk{}वा\appbrk{}प्ति\appbrk{}र्दो\appbrk{}षो\appbrk{}क्ले\appbrk{}शश्च; तत्रा\appbrk{}शु\appbrk{}द्धा\appbrk{}मा\appbrk{}श\appbrk{}य\appbrk{}मु\appbrk{}ल्ब\appbrk{}ण\appbrk{}श्ले\appbrk{}ष्मा\appbrk{}ण\appbrk{}माशु}\ A; \skt{तत्रे\appbrk{}प्सि\appbrk{}ता\appbrk{}न\appbrk{}वा\appbrk{}प्ति\appbrk{}बा\appbrk{}स्ति\appbrk{}र्दो\appbrk{}षो\appbrk{}त्क्ले\appbrk{}श\appbrk{}त\appbrk{}मु\appbrk{}ल्ब\appbrk{}ण\appbrk{}श्ले\appbrk{}ष्मा\appbrk{}ण\appbrk{}माशु}\ K; \skt{तत्रे\appbrk{}प्सि\appbrk{}ता\appbrk{}न\appbrk{}ब\appbrk{}स्ति\appbrk{}दो\appbrk{}षा\appbrk{}न्क्ले\appbrk{}श\appbrk{}त\appbrk{}मु\appbrk{}ल्ब\appbrk{}ण\appbrk{}श्ले\appbrk{}ष्मा\appbrk{}ण\appbrk{}माशु~॥ १॥}\ H\textsuperscript{ac}}}\skt{~॥ १॥ वामयित्वा }\edtext{\edtext{\edtext{\skt{भूय\-स्ती\-क्ष्ण\-त\-रै\-र्व्वि\-रे\-च\-ये\-तिति}}{\lemma{\skt{भूयस्}}\Afootnote{K H; \skt{भूय}\ A}}}{\lemma{\skt{तीक्ष्णतरैर्}}\Afootnote{K H; \skt{स्तीक्ष्णतरैर्}\ A}}}{\lemma{\skt{व्विरेचयेतिति}}\Afootnote{K H; \skt{विरेचयेत्, आ}\ A}}\skt{~। }\edtext{\skt{आमाशये}}{\lemma{\skt{आमाशये}}\Afootnote{K H; \skt{मान्वये}\ A}}\skt{ }\edtext{\skt{त्वाम\-व\-त्स\-म्वि\-धानं}}{\lemma{\skt{त्वाम\appbrk{}व\appbrk{}त्स\appbrk{}म्वि\appbrk{}धानं}}\Afootnote{A H; \skt{त्वाम\appbrk{}व\appbrk{}त्स\appbrk{}म्वि\appbrk{}धाणं~।}\ K}}\skt{~। }\edtext{\skt{अहृ\-द्ये\-ति\-प्र\-भूते च}}{\lemma{\skt{अहृ\appbrk{}द्ये\appbrk{}ति\appbrk{}प्र\appbrk{}भूते च}}\Afootnote{K H; \skt{अहृ\appbrk{}द्येऽ\appbrk{}ति\appbrk{}प्र\appbrk{}भूते च}\ A}}\skt{ हृद्य\-म्प्र\-मा\-ण\-यु\-क्तञ्च~। }\edtext{\skt{तथा}}{\lemma{\skt{तथा}}\Afootnote{K H; \skt{अ}\ A}}\skt{ }\edtext{\edtext{\skt{तथाप्युत्तिष्ठति}}{\lemma{\skt{तथाप्य्}}\Afootnote{H; \skt{त ऊर्ध्व}\ A; \skt{तथा}\ K}}}{\lemma{\skt{उत्तिष्ठति}}\Afootnote{H; \skt{मुत्तिष्ठत्यौ}\ A; \skt{प्युत्तिष्ठति}\ K}}\skt{ }\edtext{\skt{तृतीये न्न}}{\lemma{\skt{तृतीये न्न}}\Afootnote{H; \skt{षधे न तृतीयं}\ A; \skt{तृतीयन्न}\ K}}\skt{ }\edtext{\edtext{\edtext{\skt{पाय\-ये\-त\-त\-श्चै\-न\-म्म\-धु\-घृ\-त\-फा\-णि\-त\-यु\-क्तै\-र्ल्ले\-ह\-यित्वा}}{\lemma{\skt{पाययेततश्}}\Afootnote{H; \skt{पाययेत्, तत}\ A; \skt{पाययेदत}\ K}}}{\lemma{\skt{चैनम्}}\Afootnote{H; \skt{स्त्वेनं}\ A; \skt{श्चैनं}\ K}}}{\lemma{\skt{मधु\appbrk{}घृ\appbrk{}त\appbrk{}फा\appbrk{}णि\appbrk{}त\appbrk{}यु\appbrk{}क्तै\appbrk{}र्ल्ले\appbrk{}ह\appbrk{}यित्वा}}\Afootnote{H; \skt{मधु\appbrk{}घृ\appbrk{}त\appbrk{}फा\appbrk{}णि\appbrk{}त\appbrk{}यु\appbrk{}क्तै\appbrk{}र्ले\appbrk{}हैर्}\ A; \skt{मधु\appbrk{}घृ\appbrk{}त\appbrk{}फा\appbrk{}णि\appbrk{}त\appbrk{}यु\appbrk{}क्तै\appbrk{}र्ले\appbrk{}ह\appbrk{}यित्वा}\ K}}\skt{ }\edtext{\skt{विरेचयेतिति}}{\lemma{\skt{विरेचयेतिति}}\Afootnote{K H; \skt{विरेचयेत्~॥}\ A}}\skt{~।}\pend

\pstart \ledsidenote{\passref{4.34.6}}\edtext{\edtext{\edtext{\skt{दोष\-वि\-ग्र\-थि\-त\-म\-ल्प\-मौ\-ष\-ध\-म\-व\-स्थि\-त\-मू\-र्द्ध\-भा\-गि\-क\-म\-धो\-भा\-गि\-कम्वा}}{\lemma{\skt{दोष\appbrk{}वि\appbrk{}ग्र\appbrk{}थि\appbrk{}त\appbrk{}म\appbrk{}ल्प\appbrk{}मौ\appbrk{}ष\appbrk{}धम्}}\Afootnote{K H; \skt{दोष\appbrk{}वि\appbrk{}ग्र\appbrk{}थि\appbrk{}त\appbrk{}म\appbrk{}ल्प\appbrk{}मौ\appbrk{}षध}\ A}}}{\lemma{\skt{अव\appbrk{}स्थि\appbrk{}त\appbrk{}मू\appbrk{}र्द्ध\appbrk{}भा\appbrk{}गि\appbrk{}कम्}}\Afootnote{H; \skt{मव\appbrk{}स्थि\appbrk{}त\appbrk{}मू\appbrk{}र्ध्व\appbrk{}भा\appbrk{}गिक}\ A; \skt{अव\appbrk{}स्थि\appbrk{}त\appbrk{}मू\appbrk{}र्ध्व\appbrk{}भा\appbrk{}गि\appbrk{}कम्}\ K}}\skt{ न}}{\lemma{\skt{अधोभागिकम्वा न}}\Afootnote{K H; \skt{मधोभागिकं वा न}\ A}}\skt{ स्रंस\-य\-ति\-दो\-षान्~। तत्र }\edtext{\skt{तृष्णा\-पा\-र्श्व\-शू\-ल\-ञ्च्छ\-र्द्दि\-मू\-र्च्छा\-प\-र्व्व\-भे\-दो\-हृ\-ल्ला\-सो\-रु\-चि\-रु\-द्ग\-रा\-वि\-शु\-द्धिश्च}}{\lemma{\skt{तृष्णा\appbrk{}पा\appbrk{}र्श्व\appbrk{}शू\appbrk{}ल\appbrk{}ञ्च्छ\appbrk{}र्द्दि\appbrk{}मू\appbrk{}र्च्छा\appbrk{}प\appbrk{}र्व्व\appbrk{}भे\appbrk{}दो\appbrk{}हृ\appbrk{}ल्ला\appbrk{}सो\appbrk{}रु\appbrk{}चि\appbrk{}रु\appbrk{}द्ग\appbrk{}रा\appbrk{}वि\appbrk{}शु\appbrk{}द्धिश्च}}\Afootnote{H; \skt{तृष्णा पार्श्वशूलं छर्दिर्मूर्च्छा पर्वभेदो हृल्ला\appbrk{}सोऽ\appbrk{}र\appbrk{}ति\appbrk{}रु\appbrk{}द्गा\appbrk{}रा\appbrk{}वि\appbrk{}शु\appbrk{}द्धिश्च}\ A; \skt{तुष्णा\appbrk{}पा\appbrk{}र्श्व\appbrk{}शू\appbrk{}लं\appbrk{}श्छ\appbrk{}र्दि\appbrk{}र्मू\appbrk{}र्च्छा\appbrk{}प\appbrk{}र्व\appbrk{}भेदो हृल्लासो रुचि\appbrk{}रु\appbrk{}द्ग\appbrk{}रा\appbrk{}वि\appbrk{}शु\appbrk{}द्धिश्च}\ K}}\skt{ }\edtext{\skt{भवति~। }\edtext{\skt{तमु\-ष्णा\-भि\-र\-द्भि\-राशु}}{\lemma{\skt{उष्णाभिर्}}\Afootnote{K H; \skt{मुष्णाभिर्}\ A}}}{\lemma{\skt{भवति~। तम्}}\Afootnote{K H; \skt{भवति; त}\ A}}\skt{ }\edtext{\skt{वामयेत्~। }\edtext{\edtext{\skt{ऊर्द्ध्व\-भा\-गि\-क\-म\-धो\-भा\-गि\-कम्वा}}{\lemma{\skt{अधोभागिकम्}}\Afootnote{K H; \skt{अधोभागिके}\ A}}}{\lemma{\skt{वा}}\Afootnote{K H; \skt{ऽपि च}\ A}}}{\lemma{\skt{वामयेत्~। ऊर्द्ध्वभागिकम्}}\Afootnote{K H; \skt{वाम\appbrk{}ये\appbrk{}दू\appbrk{}र्ध्व\appbrk{}भा\appbrk{}गिके,}\ A}}\skt{ }\edtext{\edtext{\edtext{\edtext{\edtext{\edtext{\edtext{\edtext{\skt{साव\-शे\-षौ\-ष\-ध\-म\-प्र\-धा\-वि\-त\-दो\-ष\-म\-ति\-ब\-ल\-म\-स\-म्य\-ग्वि\-रि\-क्त\-ल\-क्ष\-ण\-मु\-प\-यु\-क्त\-म\-ल्प\-म\-प्ये\-व\-म्वा\-म\-ये\-तिति}}{\lemma{\skt{साव\appbrk{}शे\appbrk{}षौ\appbrk{}ष\appbrk{}धम्}}\Afootnote{K H; \skt{सावशेषौषध}\ A}}}{\lemma{\skt{अप्र\appbrk{}धा\appbrk{}वि\appbrk{}त\appbrk{}दो\appbrk{}षम्}}\Afootnote{K H; \skt{मति\appbrk{}प्र\appbrk{}धा\appbrk{}वि\appbrk{}त\appbrk{}दोष}\ A}}}{\lemma{\skt{अति\appbrk{}ब\appbrk{}ल\appbrk{}म\appbrk{}स\appbrk{}म्य\appbrk{}ग्वि\appbrk{}रि\appbrk{}क्त\appbrk{}ल\appbrk{}क्ष\appbrk{}णम्}}\Afootnote{H; \skt{मति\appbrk{}ब\appbrk{}ल\appbrk{}म\appbrk{}स\appbrk{}म्य\appbrk{}ग्वि\appbrk{}रि\appbrk{}क्त\appbrk{}लक्ष}\ A; \skt{अति\appbrk{}ब\appbrk{}ल\appbrk{}म\appbrk{}स\appbrk{}म्य\appbrk{}ग्वि\appbrk{}रि\appbrk{}क्त\appbrk{}ल\appbrk{}क्ष\appbrk{}णम्}\ K}}}{\lemma{\skt{उपयुक्तम्}}\Afootnote{K H; \skt{ण}\ A}}}{\lemma{\skt{अल्पम्}}\Afootnote{K H; \emph{om.}\ A}}}{\lemma{\skt{अप्य्}}\Afootnote{K H; \skt{म}\ A}}}{\lemma{\skt{एवम्}}\Afootnote{K H; \skt{प्येवं}\ A}}}{\lemma{\skt{वामयेतिति}}\Afootnote{K H; \skt{वामयेत्~॥}\ A}}\skt{~॥}\pend

\pstart \ledsidenote{\passref{4.34.7}}\edtext{\edtext{\edtext{\skt{क्रूर\-को\-ष्ठ\-स्या\-ति\-ती\-क्ष्णा\-ग्ने\-र\-ल्प\-म\-प्यौ\-ष\-ध\-म\-ल्प\-गु\-णम्वा}}{\lemma{\skt{क्रूर\appbrk{}को\appbrk{}ष्ठ\appbrk{}स्या\appbrk{}ति\appbrk{}ती\appbrk{}क्ष्णा\appbrk{}ग्नेर्}}\Afootnote{H; \skt{क्रूर\appbrk{}को\appbrk{}ष्ठ\appbrk{}स्या\appbrk{}ति\appbrk{}ती\appbrk{}क्ष्णाग्ने}\ A K}}}{\lemma{\skt{अल्प\appbrk{}म\appbrk{}प्यौ\appbrk{}ष\appbrk{}धम्}}\Afootnote{H; \skt{रल्पमौषध}\ A; \skt{रल्प\appbrk{}म\appbrk{}प्यौ\appbrk{}ष\appbrk{}धम्}\ K}}}{\lemma{\skt{अल्पगुणम्वा}}\Afootnote{K H; \skt{मल्पगुणं वा}\ A}}\skt{ भक्त\-व\-त्पा\-क\-मु\-पैति~। तत्र }\edtext{\skt{समुदीर्ण्ण}}{\lemma{\skt{समुदीर्ण्ण}}\Afootnote{K H; \skt{समुदीर्णा}\ A}}\skt{ दोषाः~। }\edtext{\edtext{\skt{यथा\-का\-ल\-म\-नि\-र्ह्रि\-य\-माणा}}{\lemma{\skt{यथाकालम्}}\Afootnote{K H; \skt{यथाकाल}\ A}}}{\lemma{\skt{अनिर्ह्रियमाणा}}\Afootnote{K H; \skt{मनिर्ह्रीयमाणा}\ A}}\skt{ }\edtext{\skt{विभ्रमं}}{\lemma{\skt{विभ्रमं}}\Afootnote{K H; \skt{व्याधिविभ्रमं बलविभ्रंशं}\ A}}\skt{ }\edtext{\skt{कुर्व्वन्ति}}{\lemma{\skt{कुर्व्वन्ति}}\Afootnote{K H; \skt{चापादयन्ति,}\ A}}\skt{~। }\edtext{\skt{तम\-न\-ल्प\-म\-म\-न्द\-ञ्चौ\-ष\-ध\-म्पा\-य\-येत्}}{\lemma{\skt{तम\appbrk{}न\appbrk{}ल्प\appbrk{}म\appbrk{}म\appbrk{}न्द\appbrk{}ञ्चौ\appbrk{}ष\appbrk{}धम्}}\Afootnote{H; \skt{तम\appbrk{}न\appbrk{}ल्प\appbrk{}म\appbrk{}म\appbrk{}न्द\appbrk{}मौ\appbrk{}षधं च}\ A; \skt{तम\appbrk{}न\appbrk{}ल्प\appbrk{}म\appbrk{}म\appbrk{}न्द\appbrk{}ञ्चौ\appbrk{}षधं}\ K}}\skt{~।}\pend

\pstart \ledsidenote{\passref{4.34.8}}\edtext{\edtext{\skt{स्निग्ध\-स्वि\-न्ने\-ना\-ल्प\-म\-ल्प\-गु\-णम्वा}}{\lemma{\skt{स्निग्ध\appbrk{}स्वि\appbrk{}न्ने\appbrk{}ना\appbrk{}ल्पम्}}\Afootnote{K H; \skt{अस्निग्धस्विन्नेना}\ A}}}{\lemma{\skt{अल्पगुणम्वा}}\Afootnote{K H; \skt{ल्पगुणं वा}\ A}}\skt{ }\edtext{\edtext{\edtext{\skt{भेष\-ज\-मु\-प\-यु\-क्त\-म\-ल्पा\-न्दो\-षा\-न्ह\-रति}}{\lemma{\skt{उपयुक्तम्}}\Afootnote{K H; \skt{उपयुक्त}\ A}}}{\lemma{\skt{अल्पान्}}\Afootnote{K H; \skt{मल्पान्}\ A}}}{\lemma{\skt{दोषान्हरति}}\Afootnote{K H; \skt{दोषान्हन्ति; त}\ A}}\skt{~। }\edtext{\edtext{\edtext{\skt{तद्दो\-ष\-शे\-ष\-म्व\-म\-न\-गौ\-र\-व\-मु\-प\-क्लेश}}{\lemma{\skt{तद्दोषशेषम्}}\Afootnote{K H; \skt{त्र वमने दोषशेषो}\ A}}}{\lemma{\skt{वम\appbrk{}न\appbrk{}गौ\appbrk{}र\appbrk{}वम्}}\Afootnote{K H; \skt{गौरव}\ A}}}{\lemma{\skt{उपक्लेश}}\Afootnote{H; \skt{मुत्क्लेशं}\ A; \skt{उपक्लेशं}\ K}}\skt{ }\edtext{\skt{हृदयां}}{\lemma{\skt{हृदयां}}\Afootnote{H; \skt{हृदया}\ A K}}\skt{ विशुद्धिं }\edtext{\skt{व्याधिविवृद्धिञ्च}}{\lemma{\skt{व्याधिविवृद्धिञ्च}}\Afootnote{K H; \skt{व्याधिवृद्धिं च}\ A}}\skt{ करोति }\edtext{\skt{रोग)।}}{\lemma{\skt{रोग)}}\Afootnote{H; \skt{त}\ A K}}\skt{ }\edtext{\skt{तत्र}}{\lemma{\skt{तत्र}}\Afootnote{H; \skt{त्र तं}\ A; \skt{त्र}\ K}}\skt{ यथा\-यो\-ग\-म्पा\-य\-यित्वा वाम\-ये\-द्दृ\-ढ\-तरं~। }\edtext{\skt{विरे\-च\-न\-दो\-ष\-शेषं}}{\lemma{\skt{विरे\appbrk{}च\appbrk{}न\appbrk{}दो\appbrk{}ष\appbrk{}शेषं}}\Afootnote{K H; \skt{विरेचने तु}\ A}}\skt{ }\edtext{\edtext{\skt{गुद\-प\-रि\-क\-र्त्त\-मा\-ध्मानं}}{\lemma{\skt{गुद\appbrk{}प\appbrk{}रि\appbrk{}क\appbrk{}र्त्तम्}}\Afootnote{K H; \skt{गुदपरिकर्त}\ A}}}{\lemma{\skt{आध्मानं}}\Afootnote{K H; \skt{नमाध्मानं}\ A}}\skt{ }\edtext{\edtext{\skt{गौर\-व\-म\-नि\-स\-र\-ण\-म्वा\-यो\-र्व्व्या\-धि\-वि\-वृ\-द्धिञ्च}}{\lemma{\skt{गौर\appbrk{}व\appbrk{}म\appbrk{}नि\appbrk{}स\appbrk{}र\appbrk{}णम्}}\Afootnote{H; \skt{शिरो\appbrk{}गौ\appbrk{}र\appbrk{}व\appbrk{}म\appbrk{}निः\appbrk{}स\appbrk{}रणं वा}\ A; \skt{गौर\appbrk{}व\appbrk{}म\appbrk{}नि\appbrk{}स\appbrk{}र\appbrk{}णम्}\ K}}}{\lemma{\skt{व्व्याधिविवृद्धिञ्च}}\Afootnote{K H; \skt{व्याधिवृद्धिं च}\ A}}\skt{ करोति~। तमुपपाद्य }\edtext{\skt{स्नेहस्वेदाभ्यां}}{\lemma{\skt{स्नेहस्वेदाभ्यां}}\Afootnote{K H; \skt{भूयः स्नेहस्वेदाभ्यां}\ A}}\skt{ }\edtext{\skt{भूयो}}{\lemma{\skt{भूयो}}\Afootnote{K H; \emph{om.}\ A}}\skt{ विरे\-च\-ये\-द्दृ\-ढ\-तरं~। दृढं }\edtext{\skt{प्रच\-लि\-त\-दो\-षम्वा}}{\lemma{\skt{प्रच\appbrk{}लि\appbrk{}त\appbrk{}दो\appbrk{}षम्वा}}\Afootnote{K H; \skt{बहु\appbrk{}प्र\appbrk{}च\appbrk{}लि\appbrk{}त\appbrk{}दोषं वा}\ A}}\skt{ तृतीये }\edtext{\skt{दिवस}}{\lemma{\skt{दिवस}}\Afootnote{K H; \skt{दिवसे}\ A}}\skt{ }\edtext{\skt{इति}}{\lemma{\skt{इति}}\Afootnote{K H; \skt{ऽल्पगुणं चेति~॥}\ A}}\skt{~।}\pend

\pstart \ledsidenote{\passref{4.34.9}}\edtext{\skt{अस्निग्धस्विन्ने}}{\lemma{\skt{अस्निग्धस्विन्ने}}\Afootnote{H; \skt{अस्निग्धस्विन्नेन}\ A K}}\skt{ }\edtext{\skt{रूक्ष\-मौ\-ष\-ध\-मु\-प\-यु\-क्त\-म\-ब्र\-ह्म\-चा\-रिणा}}{\lemma{\skt{रूक्षमौषधम्}}\Afootnote{K H; \skt{रूक्षौषधम्}\ A}}\skt{ वा वायुं कोपयति~। तत्र वायुः प्रकुपितः }\edtext{\skt{पार्श्व\-पृ\-ष्ठ\-शो\-णि\-त\-म\-न्या\-मर्म्म}}{\lemma{\skt{पार्श्व\appbrk{}पृ\appbrk{}ष्ठ\appbrk{}शो\appbrk{}णि\appbrk{}त\appbrk{}म\appbrk{}न्या\appbrk{}मर्म्म}}\Afootnote{K H; \skt{पार्श्व\appbrk{}पृ\appbrk{}ष्ठ\appbrk{}श्रो\appbrk{}णि\appbrk{}म\appbrk{}न्या\appbrk{}मर्म}\ A}}\skt{ }\edtext{\skt{मुशूलं}}{\lemma{\skt{मुशूलं}}\Afootnote{H; \skt{शूलं}\ A; \skt{मुशुलं}\ K}}\skt{ }\edtext{\skt{मूर्च्छा}}{\lemma{\skt{मूर्च्छा}}\Afootnote{H; \skt{मूर्च्छां}\ A K}}\skt{ }\edtext{\skt{भ्रम}}{\lemma{\skt{भ्रम}}\Afootnote{H; \skt{भ्रमं मदं}\ A; \skt{भ्रमं}\ K}}\skt{ संज्ञानाशञ्च करोति~। }\edtext{\edtext{\skt{तद्वा\-त\-शू\-ल\-मि\-त्या\-च\-क्षते}}{\lemma{\skt{तद्वातशूलम्}}\Afootnote{K H; \skt{तं वातशूल}\ A}}}{\lemma{\skt{इत्य्}}\Afootnote{K H; \skt{मित्य्}\ A}}\skt{~। तमभ्यज्य }\edtext{\skt{धान्यस्वेदेन}}{\lemma{\skt{धान्यस्वेदेन}}\Afootnote{K H; \skt{धाव्यस्वेदेन}\ A}}\skt{ स्वेदयित्वा }\edtext{\skt{यष्टी\-म\-धु\-क\-सि\-द्धेन}}{\lemma{\skt{यष्टी\appbrk{}म\appbrk{}धु\appbrk{}क\appbrk{}सि\appbrk{}द्धेन}}\Afootnote{K H; \skt{यष्टी\appbrk{}म\appbrk{}धु\appbrk{}क\appbrk{}वि\appbrk{}प\appbrk{}क्वेन}\ A}}\skt{ }\edtext{\skt{तैलेनानु वा}}{\lemma{\skt{तैलेनानु वा}}\Afootnote{H\textsuperscript{pc}; \skt{तैलेनानु}\ A K H\textsuperscript{ac}}}\skt{ }\edtext{\skt{वामयेदिति पाठः~।सयति}}{\lemma{\skt{वामयेदिति पाठः~।सयति}}\Afootnote{H\textsuperscript{pc}; \skt{वासयेत्~॥}\ A; \skt{वामयेदिति~॥}\ K; \skt{वा सयति~।}\ H\textsuperscript{ac}}}\skt{~।}\pend

\pstart \ledsidenote{\passref{4.34.10}}\edtext{\skt{स्नेह\-स्वे\-दा\-भ्या\-म\-वि\-भा\-वित}}{\lemma{\skt{स्नेहस्वेदाभ्याम्}}\Afootnote{A H; \skt{स्नेहस्वादाभ्याम्}\ K}}\skt{ }\edtext{\skt{शरी\-रे\-णौ\-ष\-ध\-म\-ल्प\-गु\-णम्वा}}{\lemma{\skt{शरी\appbrk{}रे\appbrk{}णौ\appbrk{}ष\appbrk{}ध\appbrk{}म\appbrk{}ल्प\appbrk{}गु\appbrk{}णम्वा}}\Afootnote{H; \skt{शरी\appbrk{}रे\appbrk{}णा\appbrk{}ल्प\appbrk{}मौ\appbrk{}ष\appbrk{}ध\appbrk{}म\appbrk{}ल्प\appbrk{}गुणं वा}\ A; \skt{शरी\appbrk{}रे\appbrk{}णौ\appbrk{}ष\appbrk{}ध\appbrk{}म\appbrk{}ल्प\appbrk{}गु\appbrk{}णंवा}\ K}}\skt{ }\edtext{\skt{पीतमूर्द्धमधो वा}}{\lemma{\skt{पीतमूर्द्धमधो वा}}\Afootnote{H; \skt{पीतमूर्ध्वमधो वा}\ A; \skt{पीतं ऊर्ध्वमधो वा}\ K}}\skt{ नाभ्येति }\edtext{\skt{दोषाश्चोपक्लेश्य तैः}}{\lemma{\skt{दोषाश्चोपक्लेश्य तैः}}\Afootnote{H; \skt{दोषांश्चोत्क्लेश्य तैः}\ A; \skt{दोषांश्चोत्क्लेश्य तैः}\ K}}\skt{ सह }\edtext{\skt{बल\-क्ष\-य\-मा\-पा\-द\-यति}}{\lemma{\skt{बल\appbrk{}क्ष\appbrk{}य\appbrk{}मा\appbrk{}पा\appbrk{}द\appbrk{}यति}}\Afootnote{H; \skt{बल\appbrk{}क्ष\appbrk{}य\appbrk{}मा\appbrk{}पा\appbrk{}द\appbrk{}यति}\ A K}}\skt{~। तत्राध्मानं }\edtext{\skt{हृद\-य\-ग्र\-ह\-स्तृ\-ष्णा\-मू\-र्च्छा\-दा\-हश्च}}{\lemma{\skt{हृद\appbrk{}य\appbrk{}ग्र\appbrk{}ह\appbrk{}स्तृ\appbrk{}ष्णा\appbrk{}मू\appbrk{}र्च्छा\appbrk{}दा\appbrk{}हश्च}}\Afootnote{A H; \skt{हृद\appbrk{}य\appbrk{}ग्र\appbrk{}ह\appbrk{}स्तृ\appbrk{}ष्णा\appbrk{}मू\appbrk{}र्च्छा\appbrk{}दा\appbrk{}हश्च}\ K}}\skt{ भवति~। }\edtext{\edtext{\edtext{\skt{तम\-यो\-ग\-मि\-त्या\-च\-क्षते}}{\lemma{\skt{अयोगम्}}\Afootnote{H; \skt{अयोग}\ A; \skt{अयोगिम्}\ K}}}{\lemma{\skt{इत्य्}}\Afootnote{K H; \skt{मित्य्}\ A}}\skt{ }\edtext{\skt{।तमाशु}}{\lemma{\skt{आशु}}\Afootnote{K H; \skt{माशु}\ A}}}{\lemma{\skt{आचक्षते~।तम्}}\Afootnote{K H; \skt{आचक्षते; त}\ A}}\skt{ }\edtext{\skt{वाम\-ये\-द्वि\-रे\-च\-येद्वा}}{\lemma{\skt{वाम\appbrk{}ये\appbrk{}द्वि\appbrk{}रे\appbrk{}च\appbrk{}येद्वा}}\Afootnote{K H; \skt{वाम\appbrk{}ये\appbrk{}न्म\appbrk{}द\appbrk{}न\appbrk{}फ\appbrk{}ल\appbrk{}ल\appbrk{}व\appbrk{}णा\appbrk{}म्बु\appbrk{}भि\appbrk{}र्वि\appbrk{}रे\appbrk{}च\appbrk{}येत्ती}\ A}}\skt{ }\edtext{\skt{दुर्व्वान्तस्य तु}}{\lemma{\skt{दुर्व्वान्तस्य तु}}\Afootnote{K H; \skt{क्ष्णतरैः कषायैश्च~। दुर्वान्तस्य तु}\ A}}\skt{ }\edtext{\skt{समुत्क्लिष्टा}}{\lemma{\skt{समुत्क्लिष्टा}}\Afootnote{A H; \skt{समुत्क्लिश्टा}\ K}}\skt{ दोषाः~। व्याप्यशरीरं }\edtext{\skt{कण्डू\-श्व\-य\-थु\-को\-ठ\-पि\-ट\-क\-ज्व\-रा\-ङ्ग\-सा\-द\-नानि}}{\lemma{\skt{कण्डू\appbrk{}श्व\appbrk{}य\appbrk{}थु\appbrk{}को\appbrk{}ठ\appbrk{}पि\appbrk{}ट\appbrk{}क\appbrk{}ज्व\appbrk{}रा\appbrk{}ङ्ग\appbrk{}सा\appbrk{}द\appbrk{}नानि}}\Afootnote{K H; \skt{कण्डू\appbrk{}श्व\appbrk{}य\appbrk{}थु\appbrk{}कु\appbrk{}ष्ठ\appbrk{}पि\appbrk{}ड\appbrk{}का\appbrk{}ज्व\appbrk{}रा\appbrk{}ङ्ग\appbrk{}म\appbrk{}र्द\appbrk{}नि\appbrk{}स्तो\appbrk{}द\appbrk{}नानि}\ A}}\skt{ कुर्व्वन्ति }\edtext{\edtext{\skt{तत\-स्ता\-न\-शे\-षा\-न्स\-हौ\-ष\-धे\-ना\-प\-ह\-रेत्}}{\lemma{\skt{ततस्}}\Afootnote{K H; \skt{तत}\ A}}\skt{~। दुर्व्विरिक्तस्य}}{\lemma{\skt{तान\appbrk{}शे\appbrk{}षा\appbrk{}न्स\appbrk{}हौ\appbrk{}ष\appbrk{}धे\appbrk{}ना\appbrk{}प\appbrk{}ह\appbrk{}रेत्~। दुर्व्विरिक्तस्य}}\Afootnote{H; \skt{स्तान\appbrk{}शे\appbrk{}षा\appbrk{}न्म\appbrk{}दौ\appbrk{}ष\appbrk{}धे\appbrk{}ना\appbrk{}प\appbrk{}ह\appbrk{}रेत्~। अस्निग्धस्विन्नस्य दुर्विरिक्तस्या}\ A; \skt{तानशेषां सहौ\appbrk{}ष\appbrk{}धे\appbrk{}ना\appbrk{}प\appbrk{}ह\appbrk{}रेत्~। दुर्विरिक्तस्य}\ K}}\skt{ }\edtext{\skt{तु}}{\lemma{\skt{तु}}\Afootnote{K H; \skt{धोनाभेः}\ A}}\skt{ स्तब्ध\-पू\-र्ण्णो\-द\-रता }\edtext{\skt{बस्ति\-शू\-लो\-मे\-ढ्र\-तोदो}}{\lemma{\skt{बस्ति\appbrk{}शू\appbrk{}लो\appbrk{}मे\appbrk{}ढ्र\appbrk{}तोदो}}\Afootnote{H; \skt{शूलं}\ A; \skt{बस्तिशूलो मेढ्रतोदो}\ K}}\skt{ }\edtext{\skt{वात\-मू\-त्र\-पु\-री\-ष\-स\-ङ्गशः}}{\lemma{\skt{वात\appbrk{}मू\appbrk{}त्र\appbrk{}पु\appbrk{}री\appbrk{}ष\appbrk{}स\appbrk{}ङ्गशः}}\Afootnote{H; \skt{वात\appbrk{}पु\appbrk{}री\appbrk{}ष\appbrk{}सङ्गः}\ A; \skt{वात\appbrk{}मू\appbrk{}त्र\appbrk{}पु\appbrk{}री\appbrk{}श\appbrk{}सं\appbrk{}गशः}\ K}}\skt{ }\edtext{\skt{कण्डू\-र्म\-ण्ड\-ल\-प्रा\-दु\-र्ब्भा\-वश्च}}{\lemma{\skt{कण्डू\appbrk{}र्म\appbrk{}ण्ड\appbrk{}ल\appbrk{}प्रा\appbrk{}दु\appbrk{}र्ब्भा\appbrk{}वश्च}}\Afootnote{K H; \skt{कण्डू\appbrk{}म\appbrk{}ण्ड\appbrk{}ल\appbrk{}प्रा\appbrk{}दु\appbrk{}र्भावो वा}\ A}}\skt{ भवति~। }\edtext{\skt{तमास्थप्य}}{\lemma{\skt{तमास्थप्य}}\Afootnote{H; \skt{तमास्थाप्य}\ A K}}\skt{ पुनः }\edtext{\skt{स्नेहयित्वा}}{\lemma{\skt{स्नेहयित्वा}}\Afootnote{K H; \skt{संस्रेह्य}\ A}}\skt{ }\edtext{\skt{विरे\-च\-ये\-त्ती\-क्ष्णे\-नेति}}{\lemma{\skt{विरे\appbrk{}च\appbrk{}ये\appbrk{}त्ती\appbrk{}क्ष्णे\appbrk{}नेति}}\Afootnote{K H; \skt{विरे\appbrk{}च\appbrk{}ये\appbrk{}त्ती\appbrk{}क्ष्णेन~।}\ A}}\skt{~। नाति\-प्र\-व\-र्त्त\-माने तिष्ठति वा दुष्टसंशोधने }\edtext{\edtext{\skt{तत्स\-न्ते\-ज\-ना\-र्थ\-मु\-ष्णो\-द\-क\-म्पा\-य\-येत्}}{\lemma{\skt{तत्स\appbrk{}न्ते\appbrk{}ज\appbrk{}ना\appbrk{}र्थम्}}\Afootnote{K H; \skt{तत्सन्तेजनार्थ}\ A}}}{\lemma{\skt{उष्णोदकम्}}\Afootnote{K H; \skt{मुष्णोदकं}\ A}}\skt{~। पाणितापैश्च पार्श्वो\-द\-र\-मु\-प\-स्वे\-द\-ये\-त्ततः }\edtext{\skt{प्रवर्त्तते}}{\lemma{\skt{प्रवर्त्तते}}\Afootnote{H; \skt{प्रवर्तन्ते}\ A K}}\skt{ दोषाः~। }\edtext{\edtext{\skt{अनुप्रवृत्तेऽल्पे}}{\lemma{\skt{अनुप्रवृत्ते}}\Afootnote{A K\textsuperscript{pc} H; \skt{अनुवृप्रवृत्ते}\ K\textsuperscript{ac}}}}{\lemma{\skt{ऽल्पे}}\Afootnote{H; \skt{चाल्प}\ A; \skt{ल्पे}\ K}}\skt{ दोषे }\edtext{\skt{जीर्ण्णेचौषधे च}}{\lemma{\skt{जीर्ण्णेचौषधे च}}\Afootnote{H; \skt{जीर्णौषधं}\ A; \skt{जीर्ण्णे चौषधै}\ K}}\skt{ }\edtext{\edtext{\skt{बहुदोषमन्तः}}{\lemma{\skt{बहुदोषम्}}\Afootnote{H; \skt{बहुदोष}\ A; \skt{बहुदेष}\ K}}}{\lemma{\skt{अन्तः}}\Afootnote{H; \skt{महः}\ A; \skt{मन्तः}\ K}}\skt{ }\edtext{\skt{शैल\-म्ब\-ल\-ञ्चा\-वेक्ष्य}}{\lemma{\skt{शैलम्}}\Afootnote{H; \skt{शेषं}\ A K}}\skt{ भूयो }\edtext{\skt{मात्राञ्च}}{\lemma{\skt{मात्राञ्च}}\Afootnote{\skt{मात्रां}\ K H\textsuperscript{pc}; \skt{मात्रांञ्च}\ H\textsuperscript{ac}}}\skt{ }\edtext{\skt{विद\-ध्या\-द\-नु\-प्र\-वृत्त}}{\lemma{\skt{विद\appbrk{}ध्या\appbrk{}द\appbrk{}नु\appbrk{}प्र\appbrk{}वृत्त}}\Afootnote{H; \skt{विद\appbrk{}ध्या\appbrk{}द\appbrk{}प्र\appbrk{}वृत्त}\ K}}\skt{ }\edtext{\skt{भेन्यूभं वा}}{\lemma{\skt{भेन्यूभं वा}}\Afootnote{H; \skt{भेषजम्वा}\ K}}\skt{ }\edtext{\edtext{\edtext{\skt{दश\-जा\-त्रा\-दू\-र्द्ध\-मु\-प\-स्थि\-त\-देहं}}{\lemma{\skt{दश\appbrk{}जा\appbrk{}त्रा\appbrk{}दू\appbrk{}र्द्धम्}}\Afootnote{H; \skt{दश\appbrk{}रा\appbrk{}त्रा\appbrk{}दू\appbrk{}र्ध्वम्}\ A; \skt{दश\appbrk{}रा\appbrk{}त्रा\appbrk{}दू\appbrk{}र्ध्वम्}\ K}}}{\lemma{\skt{उपस्थितदेहं}}\Afootnote{H; \skt{उप\appbrk{}सं\appbrk{}स्कृ\appbrk{}त\appbrk{}देहं}\ A; \skt{उप}\ K}}\skt{ }\edtext{\skt{भूयः}}{\lemma{\skt{भूयः}}\Afootnote{H; \skt{स्नेहस्वेदाभ्यां भूयः}\ A; \skt{स्थितदेहं भूयः}\ K}}\skt{ }\edtext{\skt{शोध\-ये\-द्दु\-र्व्वि\-रि\-क्त\-मा\-स्थाप्य}}{\lemma{\skt{शोध\appbrk{}ये\appbrk{}द्दु\appbrk{}र्व्वि\appbrk{}रि\appbrk{}क्त\appbrk{}मा\appbrk{}स्थाप्य}}\Afootnote{K H; \skt{शोधयेत्~। दुर्वि\appbrk{}रे\appbrk{}च्य\appbrk{}मा\appbrk{}स्थाप्य}\ A}}\skt{ पुनः }\edtext{\skt{स्नेहयित्वा}}{\lemma{\skt{स्नेहयित्वा}}\Afootnote{K H; \skt{संस्नेह्य}\ A}}\skt{ विरेचयेत्~।}}{\lemma{\skt{बलञ्… द}}\Afootnote{\emph{om.}\ A}}\skt{ ह्रीभयलोभैश्च वेगाघातशीलाः }\edtext{\skt{प्रायशःस्त्रियो}}{\lemma{\skt{प्रायशःस्त्रियो}}\Afootnote{A H; \skt{प्रायतः स्त्रियो}\ K}}\skt{ राजसमीपस्था वणिजः श्रोत्रियाश्च भवन्ति~। }\edtext{\skt{तस्मा\-दे\-तै\-र्द्दु\-र्व्वि\-रेच्याः}}{\lemma{\skt{तस्मादेतैर्}}\Afootnote{H; \skt{तस्मादेते}\ A K}}\skt{~। }\edtext{\skt{बहुवातत्वात्~। अतश्चैनां}}{\lemma{\skt{बहुवातत्वात्~। अतश्चैनां}}\Afootnote{K H; \skt{बहुवातत्वात् ; अत}\ A}}\skt{ }\edtext{\skt{सुस्वि\-न्न\-स्नि\-ग्धा\-न्वि\-रे\-च\-ये\-दिति}}{\lemma{\skt{सुस्वि\appbrk{}न्न\appbrk{}स्नि\appbrk{}ग्धा\appbrk{}न्वि\appbrk{}रे\appbrk{}च\appbrk{}ये\appbrk{}दिति}}\Afootnote{H; \skt{सुस्विन्नस्निग्धां विरेचयेदिति~।}\ K}}\skt{~॥}\edtext{}{\lemma{\skt{विरेचयेत्~। ह्रीभयलोभैश्च}}\Afootnote{K H; \skt{विरेचयेत्~। ह्रीभयलोभैर्}\ A}}\pend

\pstart \ledsidenote{\passref{4.34.11}}\edtext{\edtext{\skt{स्निग्ध\-स्वि\-न्न\-स्या\-ति\-मा\-त्र\-म\-ति\-मृ\-दु\-को\-ष्ठ\-स्या\-ति\-ती\-क्ष्ण\-म\-धि\-क\-म्वा\-भे\-ष\-ज\-म\-ति\-यो\-ग\-ङ्कु\-र्यात्}}{\lemma{\skt{स्निग्ध\appbrk{}स्वि\appbrk{}न्न\appbrk{}स्या\appbrk{}ति\appbrk{}मा\appbrk{}त्र\appbrk{}म\appbrk{}ति\appbrk{}मृ\appbrk{}दु\appbrk{}को\appbrk{}ष्ठ\appbrk{}स्या\appbrk{}ति\appbrk{}ती\appbrk{}क्ष्ण\appbrk{}म\appbrk{}धि\appbrk{}कम्}}\Afootnote{H; \skt{स्न्गिध\appbrk{}स्वि\appbrk{}न्न\appbrk{}स्या\appbrk{}ति\appbrk{}मा\appbrk{}त्र\appbrk{}म\appbrk{}ति\appbrk{}मृ\appbrk{}दु\appbrk{}को\appbrk{}ष्ठस्य वाऽति\appbrk{}ती\appbrk{}क्ष्ण\appbrk{}म\appbrk{}धिकं}\ A; \skt{स्निग्ध\appbrk{}स्वि\appbrk{}न्न\appbrk{}स्या\appbrk{}ति\appbrk{}मा\appbrk{}त्र\appbrk{}म\appbrk{}ति\appbrk{}मृ\appbrk{}दु\appbrk{}को\appbrk{}ष्ठ\appbrk{}स्या\appbrk{}ति\appbrk{}ती\appbrk{}क्ष्ण\appbrk{}म\appbrk{}धिकं}\ K}}}{\lemma{\skt{वाभे\appbrk{}ष\appbrk{}ज\appbrk{}म\appbrk{}ति\appbrk{}यो\appbrk{}गङ्}}\Afootnote{H; \skt{वा दत्त\appbrk{}मौ\appbrk{}ष\appbrk{}ध\appbrk{}म\appbrk{}ति\appbrk{}योगं}\ A; \skt{वा भेष\appbrk{}झ\appbrk{}म\appbrk{}ति\appbrk{}योगं}\ K}}\skt{~। तत्र वमनातियोगे }\edtext{\edtext{\skt{पित्ता\-ति\-प्र\-वृ\-त्ति\-र्ब्ब\-ल\-वि\-स्रंसो}}{\lemma{\skt{पित्ता\appbrk{}ति\appbrk{}प्र\appbrk{}वृ\appbrk{}त्तिर्}}\Afootnote{K H; \skt{पित्तातिप्रवृत्ति}\ A}}}{\lemma{\skt{ब्बलविस्रंसो}}\Afootnote{K H; \skt{र्बलविस्रंसो}\ A}}\skt{ }\edtext{\skt{वातकोपो}}{\lemma{\skt{वातकोपो}}\Afootnote{A H; \skt{वातकोपो}\ K}}\skt{ }\edtext{\skt{बलवाश्च}}{\lemma{\skt{बलवाश्च}}\Afootnote{K H; \skt{च बलवा}\ A}}\skt{ }\edtext{\skt{भवति}}{\lemma{\skt{भवति}}\Afootnote{K H; \skt{न्भवति, तं}\ A}}\skt{~। }\edtext{\skt{तम\-भ्य\-ज्या\-व\-गाह्य च}}{\lemma{\skt{तम\appbrk{}भ्य\appbrk{}ज्या\appbrk{}व\appbrk{}गाह्य च}}\Afootnote{K H; \skt{घृते\appbrk{}ना\appbrk{}भ्य\appbrk{}ज्या\appbrk{}व\appbrk{}गाह्य}\ A}}\skt{ शीता स्वप्सु }\edtext{\skt{शीतैः}}{\lemma{\skt{शीतैः}}\Afootnote{K H; \emph{om.}\ A}}\skt{ }\edtext{\edtext{\skt{शर्क्क\-रा\-मि\-श्रै\-र्ल्ले\-है\-रु\-प\-च\-रे\-दिति}}{\lemma{\skt{शर्क्करामिश्रैर्}}\Afootnote{K H; \skt{शर्क\appbrk{}रा\appbrk{}म\appbrk{}धु\appbrk{}मि\appbrk{}श्रैर्}\ A}}}{\lemma{\skt{उपचरेदिति}}\Afootnote{H; \skt{उप\appbrk{}च\appbrk{}रे\appbrk{}द्य\appbrk{}थास्वं;}\ A; \skt{उपचरेदिति~।}\ K}}\skt{~। विरे\-च\-ना\-ति\-योगे }\edtext{\skt{कफा\-ति\-प्र\-वृ\-त्ति\-रु\-त्त\-र\-का\-लञ्च}}{\lemma{\skt{कफा\appbrk{}ति\appbrk{}प्र\appbrk{}वृ\appbrk{}त्ति\appbrk{}रु\appbrk{}त्त\appbrk{}र\appbrk{}का\appbrk{}लञ्च}}\Afootnote{K H; \skt{कफ\appbrk{}स्या\appbrk{}ति\appbrk{}प्र\appbrk{}वृ\appbrk{}त्ति\appbrk{}रु\appbrk{}त्त\appbrk{}र\appbrk{}कालं च}\ A}}\skt{ सरक्तस्य तत्रापि }\edtext{\skt{बलविस्रंसा}}{\lemma{\skt{बलविस्रंसा}}\Afootnote{K H; \skt{बलविस्रंसो}\ A}}\skt{ }\edtext{\skt{वातकोपो}}{\lemma{\skt{वातकोपो}}\Afootnote{A K\textsuperscript{pc} H; \skt{तकोपो}\ K\textsuperscript{ac}}}\skt{ }\edtext{\skt{बलवाश्च}}{\lemma{\skt{बलवाश्च}}\Afootnote{K H; \skt{च बलवा}\ A}}\skt{ }\edtext{\skt{भवति}}{\lemma{\skt{भवति}}\Afootnote{K H; \skt{न्भवति,}\ A}}\skt{~। }\edtext{\skt{तमभ्यज्य}}{\lemma{\skt{तमभ्यज्य}}\Afootnote{K H; \skt{तमति}\ A}}\skt{ }\edtext{\edtext{\skt{शीताभिरद्भिः}}{\lemma{\skt{शीताभिर्}}\Afootnote{K H; \skt{शीताम्बु}\ A}}}{\lemma{\skt{अद्भिः}}\Afootnote{K H; \skt{भिः}\ A}}\skt{ }\edtext{\skt{परि\-षि\-च्या\-व\-गाह्य}}{\lemma{\skt{परि\appbrk{}षि\appbrk{}च्या\appbrk{}व\appbrk{}गाह्य}}\Afootnote{H; \skt{परि\appbrk{}षि\appbrk{}च्या\appbrk{}व\appbrk{}गाह्य वा}\ A; \skt{परि\appbrk{}शि\appbrk{}च्या\appbrk{}व\appbrk{}गाह्य}\ K}}\skt{ }\edtext{\edtext{\skt{शीतै\-स्त\-ण्डु\-ला\-म्बु\-भि\-र्म्म\-धु\-मि\-श्रै\-र्व्वा\-म\-येत्}}{\lemma{\skt{म्मधुमिश्रैर्}}\Afootnote{K H; \skt{मधुमिश्रै}\ A}}\skt{~। पिच्छाबस्तिञ्चास्मै}}{\lemma{\skt{व्वामयेत्~। पिच्छाबस्तिञ्चास्मै}}\Afootnote{K H; \skt{श्छर्दयेत्, पिच्छाबस्तिं चास्मै}\ A}}\skt{ }\edtext{\skt{दद्यात्~। क्षीर\-र\-स\-यो\-श्चै\-न\-म\-न्य\-त\-मेन}}{\lemma{\skt{दद्यात्~। क्षीर\appbrk{}र\appbrk{}स\appbrk{}यो\appbrk{}श्चै\appbrk{}नम्}}\Afootnote{K H; \skt{दद्या}\ A}}\skt{ }\edtext{\skt{भोज\-ये\-त्क्षी\-र\-स\-र्प्पिषा }\edtext{\edtext{\skt{चैन\-म\-नु\-वा\-स\-येत्}}{\lemma{\skt{चैनम्}}\Afootnote{K H; \skt{चैन}\ A}}\skt{~। प्रियङ्ग्वादिञ्चास्मै}}{\lemma{\skt{अनुवासयेत्~। प्रियङ्ग्वादिञ्चास्मै}}\Afootnote{K H; \skt{मनुवासयेत्, प्रियङ्ग्वादिं चास्मै}\ A}}\skt{ }\edtext{\skt{तण्डुलाम्बुना}}{\lemma{\skt{तण्डुलाम्बुना}}\Afootnote{A H; \skt{तण्डुलाम्बुनां}\ K}}}{\lemma{\skt{त्… क्षी}}\Afootnote{\emph{om.}\ A}}\skt{ }\edtext{\skt{पातुम्प्रयच्छेत्}}{\lemma{\skt{प्रयच्छेत्}}\Afootnote{K H; \skt{प्रयच्छेत्, क्षीर\appbrk{}र\appbrk{}स\appbrk{}यो\appbrk{}श्चा\appbrk{}न्य\appbrk{}त\appbrk{}रेण भोजयेत्~॥}\ A}}\skt{~।}\pend

\pstart \ledsidenote{\passref{4.34.12}}\edtext{\skt{तस्मिनेव च}}{\lemma{\skt{तस्मिनेव च}}\Afootnote{K H; \skt{तसिमिन्नेव}\ A}}\skt{ }\edtext{\skt{वमनोतियोगे}}{\lemma{\skt{वमनोतियोगे}}\Afootnote{H; \skt{वमनातियोगे}\ A K}}\skt{ }\edtext{\skt{तिप्रवृत्ते}}{\lemma{\skt{तिप्रवृत्ते}}\Afootnote{K H; \skt{प्रवृद्धे}\ A}}\skt{ शोणितं }\edtext{\skt{ष्ठीवति}}{\lemma{\skt{ष्ठीवति}}\Afootnote{A H; \skt{क्षीवति}\ K}}\skt{ च्छर्द्दयति वा }\edtext{\skt{रक्तं}}{\lemma{\skt{रक्तं}}\Afootnote{K H; \emph{om.}\ A}}\skt{ तत्र }\edtext{\edtext{\skt{जिह्वा\-निः\-स\-र\-ण\-म\-व\-स\-र्प्प\-ण\-ञ्चा\-क्ष्णो\-र्व्व्या\-वृ\-त्ति\-ह\-नु\-सं\-ह\-ननं}}{\lemma{\skt{जिह्वा\appbrk{}निः\appbrk{}स\appbrk{}र\appbrk{}ण\appbrk{}म\appbrk{}व\appbrk{}स\appbrk{}र्प्प\appbrk{}ण\appbrk{}ञ्चा\appbrk{}क्ष्णोर्}}\Afootnote{H; \skt{जिह्वा\appbrk{}निः\appbrk{}स\appbrk{}र\appbrk{}ण\appbrk{}म\appbrk{}प\appbrk{}स\appbrk{}र\appbrk{}ण\appbrk{}मक्ष्णो}\ A; \skt{जिह्वा\appbrk{}निः\appbrk{}स\appbrk{}र\appbrk{}ण\appbrk{}म\appbrk{}व\appbrk{}स\appbrk{}र्प\appbrk{}ण\appbrk{}ञ्चा\appbrk{}क्ष्णोर्}\ K}}}{\lemma{\skt{व्व्यावृ\appbrk{}त्ति\appbrk{}ह\appbrk{}नु\appbrk{}सं\appbrk{}ह\appbrk{}ननं}}\Afootnote{H; \skt{र्व्यावृ\appbrk{}त्ति\appbrk{}र्ह\appbrk{}नु\appbrk{}सं\appbrk{}ह\appbrk{}ननं}\ A; \skt{व्यावृ\appbrk{}र्त्ति\appbrk{}ह\appbrk{}नु\appbrk{}सं\appbrk{}ह\appbrk{}ननं}\ K}}\skt{ तृष्णा }\edtext{\skt{हिक्काङ्गरौ}}{\lemma{\skt{हिक्काङ्गरौ}}\Afootnote{K H; \skt{हिक्का ज्वरो}\ A}}\skt{ }\edtext{\skt{विसं }\edtext{\edtext{\skt{ज्ञत्व\-मि\-त्यु\-प\-द्रवा}}{\lemma{\skt{इत्य्}}\Afootnote{K H; \skt{मित्य्}\ A}}}{\lemma{\skt{उपद्रवा}}\Afootnote{K H; \skt{औपद्रवा}\ A}}}{\lemma{\skt{विसं ज्ञत्वम्}}\Afootnote{K H; \skt{वैसंज्ञ्य}\ A}}\skt{ भवन्ति~। }\edtext{\skt{तम\-ति\-वि\-स्रु\-त\-शो\-णि\-त\-वि\-धा\-नेन नोपचरेत्~। जिह्वा\-न्निः\-स\-र्प्पि\-ता\-न्त्रि\-क\-टु\-क\-ल\-व\-ण\-चू\-र्ण्ण\-प्र\-घृ\-ष्टा\-मन्तः}}{\lemma{\skt{तम\appbrk{}ति\appbrk{}वि\appbrk{}स्रु\appbrk{}त\appbrk{}शो\appbrk{}णि\appbrk{}त\appbrk{}वि\appbrk{}धा\appbrk{}नेन नोपचरेत्~। जिह्वा\appbrk{}न्निः\appbrk{}स\appbrk{}र्प्पि\appbrk{}ता\appbrk{}न्त्रि\appbrk{}क\appbrk{}टु\appbrk{}क\appbrk{}ल\appbrk{}व\appbrk{}ण\appbrk{}चू\appbrk{}र्ण्ण\appbrk{}प्र\appbrk{}घृ\appbrk{}ष्टा\appbrk{}मन्तः}}\Afootnote{H; \skt{तम\appbrk{}जा\appbrk{}सृ\appbrk{}क्च\appbrk{}न्द\appbrk{}नो\appbrk{}शी\appbrk{}रा\appbrk{}ञ्ज\appbrk{}न\appbrk{}ला\appbrk{}ज\appbrk{}चूर्णैः सश\appbrk{}र्क\appbrk{}रो\appbrk{}द\appbrk{}कै\appbrk{}र्मन्थं पाययेत्, फलरसैर्वा सघृ\appbrk{}त\appbrk{}क्षौ\appbrk{}द्र\appbrk{}श\appbrk{}र्करैः शुङ्गाभिर्वा वटादीनां पेयां सिद्धां सक्षौद्रां वर्चोग्राहिभिर्वा, पयसा जाङ्गलरसेन वा भोजयेत् , अति\appbrk{}स्रु\appbrk{}त\appbrk{}शो\appbrk{}णि\appbrk{}त\appbrk{}वि\appbrk{}धा\appbrk{}ने\appbrk{}नो\appbrk{}प\appbrk{}च\appbrk{}रेत्; जिह्वा\appbrk{}म\appbrk{}ति\appbrk{}स\appbrk{}र्पितां कटु\appbrk{}क\appbrk{}ल\appbrk{}व\appbrk{}ण\appbrk{}चू\appbrk{}र्ण\appbrk{}प्र\appbrk{}घृष्टां तिल\appbrk{}द्रा\appbrk{}क्षा\appbrk{}प्र\appbrk{}लिप्तां वाऽन्तः}\ A; \skt{तम\appbrk{}ति\appbrk{}वि\appbrk{}स्रु\appbrk{}त\appbrk{}शो\appbrk{}णि\appbrk{}त\appbrk{}वि\appbrk{}धा\appbrk{}णे\appbrk{}नो\appbrk{}प\appbrk{}च\appbrk{}रेत्~। जिह्वां निःसर्पितां त्रिक\appbrk{}टु\appbrk{}क\appbrk{}व\appbrk{}ल\appbrk{}व\appbrk{}ण\appbrk{}चूर्ण्ण प्रघृष्टामन्तः}\ K\textsuperscript{ac}; \skt{तम\appbrk{}ति\appbrk{}वि\appbrk{}स्रु\appbrk{}त\appbrk{}शो\appbrk{}णि\appbrk{}त\appbrk{}वि\appbrk{}धा\appbrk{}णे\appbrk{}नो\appbrk{}प\appbrk{}च\appbrk{}रेत्~। जिह्वां निःसर्पितां त्रिक\appbrk{}टु\appbrk{}क\appbrk{}ल\appbrk{}व\appbrk{}ण\appbrk{}चूर्ण्ण प्रघृष्टामन्तः}\ K\textsuperscript{pc}}}\skt{ }\edtext{\edtext{\skt{पीड\-ये\-त्प्र\-वि\-ष्टा\-या\-म\-न्ये\-म्ल\-मस्य}}{\lemma{\skt{पीड\appbrk{}ये\appbrk{}त्प्र\appbrk{}वि\appbrk{}ष्टा\appbrk{}याम्}}\Afootnote{K H; \skt{पीडयेत्, अन्तः प्रविष्टाया}\ A}}}{\lemma{\skt{अन्येम्लमस्य}}\Afootnote{H; \skt{मम्लमन्ये तस्य}\ A; \skt{अन्येम्व्लस्य}\ K}}\skt{ }\edtext{\skt{पुरस्तात्खादेयुः}}{\lemma{\skt{पुरस्तात्खादेयुः}}\Afootnote{K H; \skt{पुर\appbrk{}स्ता\appbrk{}त्खा\appbrk{}द\appbrk{}येयुः;}\ A}}\skt{~। व्यावृत्ते चाक्षिणी }\edtext{\skt{घृताभ्यक्तैः}}{\lemma{\skt{घृताभ्यक्तैः}}\Afootnote{H; \skt{घृताभ्यक्ते}\ A; \skt{घृताभ्यक्तै}\ K}}\skt{ }\edtext{\skt{पीडयेत्~। हनुसंहनने}}{\lemma{\skt{पीडयेत्~। हनुसंहनने}}\Afootnote{K H; \skt{पीडये}\ A}}\skt{ वात\-श्ले\-ष्म\-ह\-र\-न्नस्यं }\edtext{\skt{स्वेदञ्च विदध्यात्~। तृष्णादिषु च यथास्वं प्रतिकुर्व्वीत~। }\edtext{\skt{विसंज्ञम्वा}}{\lemma{\skt{विसंज्ञम्वा}}\Afootnote{K H; \skt{विसंज्ञे}\ A}}\skt{ }\edtext{\skt{वेणु\-वी\-णा\-गी\-त\-स्वेनं}}{\lemma{\skt{वेणु\appbrk{}वी\appbrk{}णा\appbrk{}गी\appbrk{}त\appbrk{}स्वेनं}}\Afootnote{\skt{वेणु\appbrk{}वी\appbrk{}णा\appbrk{}गी\appbrk{}त\appbrk{}स्वनं}\ A; \skt{वेणु\appbrk{}वी\appbrk{}णा\appbrk{}गी\appbrk{}त\appbrk{}स्वेनं}\ K H}}}{\lemma{\skt{त्… तृ}}\Afootnote{\emph{om.}\ A}}\skt{ }\edtext{\skt{स्रावयेदिति}}{\lemma{\skt{स्रावयेदिति}}\Afootnote{K H; \skt{श्रावयेत्~॥}\ A}}\skt{~।}\pend

\pstart \ledsidenote{\passref{4.34.13}}\edtext{\skt{विरे\-च\-ना\-ति\-योगे}}{\lemma{\skt{विरे\appbrk{}च\appbrk{}ना\appbrk{}ति\appbrk{}योगे}}\Afootnote{K H; \skt{विरे\appbrk{}च\appbrk{}ना\appbrk{}ति\appbrk{}योगे च}\ A}}\skt{ }\edtext{\edtext{\skt{सच\-न्द्र\-कं\-स\-लि\-ल\-मधः}}{\lemma{\skt{सच\appbrk{}न्द्र\appbrk{}कं\appbrk{}स\appbrk{}लि\appbrk{}लम्}}\Afootnote{K H; \skt{सचन्द्रकं सलिल}\ A}}}{\lemma{\skt{अधः}}\Afootnote{K H; \skt{मधः}\ A}}\skt{ स्रवति ततो मांस\-धा\-व\-न\-प्र\-का\-श\-मु\-त्त\-र\-कालं जीवशोणितञ्च~। }\edtext{\skt{तत्र}}{\lemma{\skt{तत्र}}\Afootnote{K H; \skt{ततो}\ A}}\skt{ गुदनिःसरणं }\edtext{\edtext{\skt{वेप\-थु\-र्व्व\-म\-ना\-ति\-यो\-गो\-प\-द्र\-वाश्च}}{\lemma{\skt{वेपथुर्}}\Afootnote{K H; \skt{वेपथु}\ A}}}{\lemma{\skt{व्वम\appbrk{}ना\appbrk{}ति\appbrk{}यो\appbrk{}गो\appbrk{}प\appbrk{}द्र\appbrk{}वाश्च}}\Afootnote{K H; \skt{र्वम\appbrk{}ना\appbrk{}ति\appbrk{}यो\appbrk{}गो\appbrk{}प\appbrk{}द्र\appbrk{}वा\appbrk{}श्चास्य}\ A}}\skt{ }\edtext{\skt{भवन्ति~। तम\-ति\-स्रु\-त\-वि\-धा\-ने\-नै\-वो\-प\-च\-रेत्}}{\lemma{\skt{भवन्ति~। तम्}}\Afootnote{A H; \skt{भवन्ति~।}\ K}}\skt{~। }\edtext{\skt{निःस\-र्प्पि\-त\-ङ्गु\-द\-म\-भ्यज्य}}{\lemma{\skt{गुदमभ्यज्य}}\Afootnote{H; \skt{गुदस्य गुदमभ्यज्य}\ A; \skt{गुदमभ्यज्य}\ K}}\skt{ }\edtext{\skt{परि\-स्वे\-द्या\-न्पी\-ड\-येत्~। क्षुद्र\-रो\-ग\-चि\-कि\-त्सिते}}{\lemma{\skt{परि\appbrk{}स्वे\appbrk{}द्या\appbrk{}न्पी\appbrk{}ड\appbrk{}येत्~। क्षुद्र\appbrk{}रो\appbrk{}ग\appbrk{}चि\appbrk{}कि\appbrk{}त्सिते}}\Afootnote{H; \skt{परिस्वेद्यान्तः पीड\appbrk{}ये\appbrk{}त्क्षु\appbrk{}द्र\appbrk{}रोग}\ A; \skt{परिस्वेद्यान्तः पीडयेत्~। क्षुद्र\appbrk{}रो\appbrk{}ग\appbrk{}चि\appbrk{}कि\appbrk{}त्सिते}\ K}}\skt{ }\edtext{\edtext{\skt{गुद\-स्रं\-स\-चि\-कि\-त्सि\-त\-ञ्चा\-वे\-क्ष्येत}}{\lemma{\skt{गुद\appbrk{}स्रं\appbrk{}स\appbrk{}चि\appbrk{}कि\appbrk{}त्सि\appbrk{}तञ्}}\Afootnote{K H; \skt{चिकित्सितं}\ A}}}{\lemma{\skt{चावेक्ष्येत}}\Afootnote{K H; \skt{वा वीक्षेत,}\ A}}\skt{~। वेपथौ }\edtext{\skt{वात\-व्या\-धि\-वि\-धा\-न\-ङ्कु\-र्व्वीत}}{\lemma{\skt{वात\appbrk{}व्या\appbrk{}धि\appbrk{}वि\appbrk{}धा\appbrk{}नङ्}}\Afootnote{A H; \skt{वातव्याधिपानं}\ K}}\skt{~। }\edtext{\skt{जिह्वा\-निः\-स\-र्प्प\-णा\-दि\-षूक्तः}}{\lemma{\skt{जिह्वा\appbrk{}निः\appbrk{}स\appbrk{}र्प्प\appbrk{}णा\appbrk{}दि\appbrk{}षूक्तः}}\Afootnote{K H; \skt{जिह्वा\appbrk{}निः\appbrk{}स\appbrk{}र\appbrk{}णा\appbrk{}दि\appbrk{}षूक्तः}\ A}}\skt{ प्रतीकारः~। अतिप्रवृत्ते वा }\edtext{\skt{न्यग्रोधादिकषा यक्षी\-रे\-क्षु\-र\-स\-घृ\-त\-सं\-सृ\-ष्टै\-श्चैनं}}{\lemma{\skt{न्यग्रोधादिकषा यक्षी\appbrk{}रे\appbrk{}क्षु\appbrk{}र\appbrk{}स\appbrk{}घृ\appbrk{}त\appbrk{}सं\appbrk{}सृ\appbrk{}ष्टै\appbrk{}श्चैनं}}\Afootnote{\skt{जीवशोणिते काश्म\appbrk{}री\appbrk{}फ\appbrk{}ल\appbrk{}ब\appbrk{}द\appbrk{}री\appbrk{}दू\appbrk{}र्वो\appbrk{}शीरैः शृतेन पयसा घृत\appbrk{}म\appbrk{}ण्डा\appbrk{}ञ्ज\appbrk{}न\appbrk{}यु\appbrk{}क्तेन सुशी\appbrk{}ते\appbrk{}ना\appbrk{}स्था\appbrk{}प\appbrk{}येत्, न्यग्रो\appbrk{}धा\appbrk{}दि\appbrk{}क\appbrk{}षा\appbrk{}ये\appbrk{}क्षु\appbrk{}र\appbrk{}स\appbrk{}घृ\appbrk{}त\appbrk{}शो\appbrk{}णि\appbrk{}त\appbrk{}सं\appbrk{}सृ\appbrk{}ष्टै\appbrk{}श्चैनं}\ A; \skt{न्यग्रो\appbrk{}धा\appbrk{}दि\appbrk{}क\appbrk{}षा\appbrk{}य\appbrk{}क्षी\appbrk{}रे\appbrk{}क्षु\appbrk{}र\appbrk{}स\appbrk{}घृ\appbrk{}त\appbrk{}शो\appbrk{}णि\appbrk{}त\appbrk{}सं\appbrk{}सृ\appbrk{}ष्टै\appbrk{}श्चैनं}\ K; \skt{न्यग्रोधादिकषा यक्षी\appbrk{}रे\appbrk{}क्षु\appbrk{}र\appbrk{}स\appbrk{}घृ\appbrk{}त\appbrk{}सं\appbrk{}सृ\appbrk{}ष्टै\appbrk{}श्चैनं}\ H}}\skt{ }\edtext{\skt{बस्ति\-रु\-पा\-च\-रेत्~। शोणि\-त\-नि\-ष्ठी\-वने}}{\lemma{\skt{बस्ति\appbrk{}रु\appbrk{}पा\appbrk{}च\appbrk{}रेत्~। शोणि\appbrk{}त\appbrk{}नि\appbrk{}ष्ठी\appbrk{}वने}}\Afootnote{H; \skt{बस्ति\appbrk{}भि\appbrk{}रु\appbrk{}पा\appbrk{}च\appbrk{}रेत्, शोणितष्ठीवने}\ A; \skt{बस्ति\appbrk{}रु\appbrk{}पा\appbrk{}च\appbrk{}रेत्~। शोणि\appbrk{}त\appbrk{}निः\appbrk{}ष्ठी\appbrk{}वने}\ K}}\skt{ }\edtext{\edtext{\skt{रक्ता\-ति\-सा\-र\-र\-क्त\-पि\-त्त\-क्रि\-या\-श्चा\-न्यस्य}}{\lemma{\skt{रक्ता\appbrk{}ति\appbrk{}सा\appbrk{}र\appbrk{}र\appbrk{}क्त\appbrk{}पि\appbrk{}त्त\appbrk{}क्रि\appbrk{}याश्}}\Afootnote{K H; \skt{रक्त\appbrk{}पि\appbrk{}त्त\appbrk{}र\appbrk{}क्ता\appbrk{}ती\appbrk{}सा\appbrk{}र\appbrk{}क्रि\appbrk{}याश्}\ A}}}{\lemma{\skt{चान्यस्य}}\Afootnote{H; \skt{चास्य}\ A K}}\skt{ }\edtext{\skt{विद\-ध्या\-न्न्य\-ग्रो\-धा\-दि\-ञ्चास्य}}{\lemma{\skt{विद\appbrk{}ध्या\appbrk{}न्न्य\appbrk{}ग्रो\appbrk{}धा\appbrk{}दिञ्}}\Afootnote{H; \skt{विधध्यात्, न्यग्रोधादिं}\ A; \skt{विद\appbrk{}ध्या\appbrk{}त्न्य\appbrk{}ग्रो\appbrk{}धादिं}\ K}}\skt{ }\edtext{\skt{पान\-भो\-ज\-ने\-षू\-प\-यु\-ञ्जीत}}{\lemma{\skt{पान\appbrk{}भो\appbrk{}ज\appbrk{}ने\appbrk{}षू\appbrk{}प\appbrk{}यु\appbrk{}ञ्जीत}}\Afootnote{K H; \skt{विद\appbrk{}ध्या\appbrk{}त्पा\appbrk{}न\appbrk{}भो\appbrk{}ज\appbrk{}नेषु~॥}\ A}}\skt{~।}\edtext{}{\lemma{\skt{अति\appbrk{}स्रु\appbrk{}त\appbrk{}वि\appbrk{}धा\appbrk{}ने\appbrk{}नै\appbrk{}वो\appbrk{}प\appbrk{}च\appbrk{}रेत्~। निःसर्प्पितङ्}}\Afootnote{H; \skt{अपि निःस्रु\appbrk{}त\appbrk{}शो\appbrk{}णि\appbrk{}त\appbrk{}वि\appbrk{}धा\appbrk{}ने\appbrk{}नो\appbrk{}प\appbrk{}च\appbrk{}रेत्, निःसर्पित}\ A; \skt{तम\appbrk{}ति\appbrk{}स्रु\appbrk{}त\appbrk{}वि\appbrk{}धा\appbrk{}ने\appbrk{}नै\appbrk{}वो\appbrk{}च\appbrk{}रेत्~। निःसर्पितं}\ K\textsuperscript{ac}; \skt{तम\appbrk{}ति\appbrk{}स्रु\appbrk{}त\appbrk{}वि\appbrk{}धा\appbrk{}ने\appbrk{}नै\appbrk{}वो\appbrk{}प\appbrk{}च\appbrk{}रेत्~। निःसर्पितं}\ K\textsuperscript{pc}}}\pend

\pstart \ledsidenote{\passref{4.34.14}}\edtext{\skt{जीव\-शो\-णि\-त\-र\-क्त\-पि\-त्तयो}}{\lemma{\skt{जीव\appbrk{}शो\appbrk{}णि\appbrk{}त\appbrk{}र\appbrk{}क्त\appbrk{}पि\appbrk{}त्तयो}}\Afootnote{H; \skt{जीव\appbrk{}शो\appbrk{}णि\appbrk{}त\appbrk{}र\appbrk{}क्त\appbrk{}पि\appbrk{}त्त\appbrk{}योश्च}\ A; \skt{जीव\appbrk{}शो\appbrk{}णि\appbrk{}त\appbrk{}र\appbrk{}क्त\appbrk{}पि\appbrk{}त्त\appbrk{}योर्}\ K}}\skt{ }\edtext{\skt{जिज्ञासार्थम्पिचु}}{\lemma{\skt{पिचु}}\Afootnote{K H; \skt{तस्मिन्पिचुं}\ A}}\skt{ }\edtext{\skt{प्लोतम्वात्र}}{\lemma{\skt{वात्र}}\Afootnote{K H; \skt{वा}\ A}}\skt{ }\edtext{\edtext{\skt{प्रक्षि\-पे\-त्त\-दु\-ष्णो\-द\-क\-प्र\-क्षा\-लि\-त\-म\-व\-र\-ज्येत}}{\lemma{\skt{प्रक्षि\appbrk{}पे\appbrk{}त्त\appbrk{}दु\appbrk{}ष्णो\appbrk{}द\appbrk{}क\appbrk{}प्र\appbrk{}क्षा\appbrk{}लि\appbrk{}तम्}}\Afootnote{H; \skt{क्षिपेत्, यद्यु\appbrk{}ष्णो\appbrk{}द\appbrk{}क\appbrk{}प्र\appbrk{}क्षा\appbrk{}लि\appbrk{}तम्}\ A; \skt{प्रक्षि\appbrk{}पे\appbrk{}त्त\appbrk{}दु\appbrk{}ष्णो\appbrk{}द\appbrk{}क\appbrk{}प्र\appbrk{}क्षा\appbrk{}लि\appbrk{}तम्}\ K}}}{\lemma{\skt{अवरज्येत}}\Afootnote{K H; \skt{अपि वस्त्रं रञ्जयति}\ A}}\skt{ }\edtext{\skt{तज्जीवशोणित }\edtext{\edtext{\edtext{\skt{मित्य\-व\-ग\-न्त\-व्य\-म\-मृ\-ग्भो\-क्तुम्वा}}{\lemma{\skt{इत्य्}}\Afootnote{H; \skt{त}\ A; \skt{इत्य}\ K}}}{\lemma{\skt{अवगन्तव्यम्}}\Afootnote{H; \skt{मवगन्तव्यं;}\ A; \skt{वगन्तव्यम्}\ K}}}{\lemma{\skt{अमृग्भोक्तुम्वा}}\Afootnote{K H; \skt{सभक्तं च}\ A}}}{\lemma{\skt{तज्जीवशोणित म्}}\Afootnote{H; \skt{तज्जीवशोणि}\ A; \skt{तज्जी\appbrk{}व\appbrk{}शो\appbrk{}णि\appbrk{}तम्}\ K}}\skt{ शुने }\edtext{\skt{दद्यात्स}}{\lemma{\skt{दद्यात्स}}\Afootnote{K H; \skt{दद्याच्छ}\ A}}\skt{ }\edtext{\edtext{\edtext{\edtext{\edtext{\skt{यद्यु\-प\-यु\-ञ्ज्या\-त्त\-ज्जी\-व\-शो\-णि\-त\-मि\-त्य\-व\-ग\-न्त\-व्य\-मिति}}{\lemma{\skt{यद्य्}}\Afootnote{H; \skt{क्तुसंमिश्रं}\ A; \emph{om.}\ K}}}{\lemma{\skt{उप\appbrk{}यु\appbrk{}ञ्ज्या\appbrk{}त्त\appbrk{}ज्जी\appbrk{}व\appbrk{}शो\appbrk{}णि\appbrk{}तम्}}\Afootnote{H; \skt{वा, स यद्युपभुञ्जीत तज्जीवशोणि}\ A; \skt{यद्यु\appbrk{}प\appbrk{}यु\appbrk{}ञ्ज्या\appbrk{}त्त\appbrk{}ज्जी\appbrk{}व\appbrk{}शो\appbrk{}णि\appbrk{}तम्}\ K}}}{\lemma{\skt{इत्य्}}\Afootnote{H; \skt{त}\ A; \skt{इ}\ K}}}{\lemma{\skt{अवगन्तव्यम्}}\Afootnote{H; \skt{मवगन्तव्यम्;}\ A; \emph{om.}\ K}}}{\lemma{\skt{इति}}\Afootnote{H; \skt{अन्यथा रक्तपित्तमिति~॥}\ A; \skt{त्यव\appbrk{}ग\appbrk{}न्त\appbrk{}व्य\appbrk{}मिति~॥}\ K}}\skt{~। }\edtext{\skt{तयोश्च}}{\lemma{\skt{तयोश्च}}\Afootnote{K H; \emph{om.}\ A}}\skt{ }\edtext{\edtext{\skt{रक्ति\-पि\-त्ता\-ती\-सा\-र\-क्रि\-या\-म्वि\-द\-धीत}}{\lemma{\skt{रक्ति\appbrk{}पि\appbrk{}त्ता\appbrk{}ती\appbrk{}सा\appbrk{}र\appbrk{}क्रि\appbrk{}याम्}}\Afootnote{H; \emph{om.}\ A; \skt{रक्त\appbrk{}पि\appbrk{}त्ता\appbrk{}ती\appbrk{}सा\appbrk{}र\appbrk{}क्रियां}\ K}}}{\lemma{\skt{विदधीत}}\Afootnote{K H; \emph{om.}\ A}}\skt{~।}\pend

\pstart \ledsidenote{\passref{4.34.15}}\edtext{\skt{स शेषा\-न्ने\-ना\-नि\-ल\-प्रा\-य\-को\-ष्ठे\-ना\-स्नि\-ग्धेन}}{\lemma{\skt{स शेषा\appbrk{}न्ने\appbrk{}ना\appbrk{}नि\appbrk{}ल\appbrk{}प्रा\appbrk{}य\appbrk{}को\appbrk{}ष्ठे\appbrk{}ना\appbrk{}स्नि\appbrk{}ग्धेन}}\Afootnote{H; \skt{सशेषान्नेन बहुदोषेण रूक्षे\appbrk{}णा\appbrk{}नि\appbrk{}ल\appbrk{}प्रा\appbrk{}य\appbrk{}को\appbrk{}ष्ठे\appbrk{}ना\appbrk{}नु\appbrk{}ष्णम}\ A; \skt{स शेषा\appbrk{}न्ना\appbrk{}नि\appbrk{}ल\appbrk{}प्रा\appbrk{}य\appbrk{}को\appbrk{}ष्ठे\appbrk{}ना\appbrk{}स्नि\appbrk{}ग्धेन}\ K}}\skt{ }\edtext{\skt{वा}}{\lemma{\skt{वा}}\Afootnote{K H; \skt{स्निग्धं वा}\ A}}\skt{ }\edtext{\edtext{\skt{पीत\-मौ\-ष\-ध\-मा\-ध्मा\-प\-यति}}{\lemma{\skt{पीतमौषधम्}}\Afootnote{K H; \skt{पीतमौषध}\ A}}}{\lemma{\skt{आध्मापयति}}\Afootnote{K H; \skt{माध्मापयति,}\ A}}\skt{~। }\edtext{\skt{तत्रा\-नि\-ल\-मू\-त्र\-पु\-री\-ष\-सङ्गः}}{\lemma{\skt{तत्रा\appbrk{}नि\appbrk{}ल\appbrk{}मू\appbrk{}त्र\appbrk{}पु\appbrk{}री\appbrk{}ष\appbrk{}सङ्गः}}\Afootnote{A H; \skt{तत्रा\appbrk{}नि\appbrk{}ल\appbrk{}मू\appbrk{}त्र\appbrk{}पु\appbrk{}री\appbrk{}श\appbrk{}संगः~।}\ K}}\skt{~। समु\-न्न\-द्धो\-द\-रता }\edtext{\skt{पार्श्वभङ्गो}}{\lemma{\skt{पार्श्वभङ्गो}}\Afootnote{A H; \skt{हृद\appbrk{}य\appbrk{}पा\appbrk{}र्श्व\appbrk{}भङ्गो}\ K}}\skt{ }\edtext{\skt{गुद\-ब\-स्ति\-तु\-द\-नञ्च}}{\lemma{\skt{गुद\appbrk{}ब\appbrk{}स्ति\appbrk{}तु\appbrk{}द\appbrk{}नञ्च}}\Afootnote{K H; \skt{गुद\appbrk{}ब\appbrk{}स्ति\appbrk{}नि\appbrk{}स्तो\appbrk{}दनं भक्तारुचिश्च}\ A}}\skt{ भवति~। }\edtext{\edtext{\skt{तमा\-ध्मा\-न\-मि\-त्या\-च\-क्षते}}{\lemma{\skt{तमाध्मानम्}}\Afootnote{K H; \skt{तं चाध्मान}\ A}}}{\lemma{\skt{इत्य्}}\Afootnote{K H; \skt{मित्य्}\ A}}\skt{~। तमु\-प\-स्वे\-द्या\-नाह }\edtext{\edtext{\skt{वर्त्ति\-दी\-प\-न\-ब\-स्ति\-क्रि\-या\-भि\-र्योज्या}}{\lemma{\skt{वर्त्ति\appbrk{}दी\appbrk{}प\appbrk{}न\appbrk{}ब\appbrk{}स्ति\appbrk{}क्रि\appbrk{}या\appbrk{}भिर्}}\Afootnote{A H; \skt{वर्त्ति\appbrk{}दी\appbrk{}प\appbrk{}ब\appbrk{}स्ति\appbrk{}क्रि\appbrk{}या\appbrk{}भिर्}\ K}}}{\lemma{\skt{योज्या}}\Afootnote{K H; \emph{om.}\ A}}\skt{ }\edtext{\skt{इति}}{\lemma{\skt{इति}}\Afootnote{K H; \skt{उपचरेत्~॥}\ A}}\skt{~।}\pend

\pstart \ledsidenote{\passref{4.34.16}}\skt{क्षामे\-णा\-ति\-मृ\-दु\-को\-ष्ठेन }\edtext{\edtext{\skt{तीक्ष्ण\-म\-त्यु\-ष्ण\-म\-ति\-ल\-व\-ण\-ति\-रू\-क्षम्वा}}{\lemma{\skt{तीक्ष्णमत्युष्णम्}}\Afootnote{H; \skt{मन्दाग्निना रूक्षेण वाऽतितीक्ष्णो}\ A; \skt{तीक्ष्णमत्युष्णम्}\ K}}}{\lemma{\skt{अति\appbrk{}ल\appbrk{}व\appbrk{}ण\appbrk{}ति\appbrk{}रू\appbrk{}क्षम्वा}}\Afootnote{H; \skt{ष्णाति\appbrk{}ल\appbrk{}व\appbrk{}ण\appbrk{}म\appbrk{}ति\appbrk{}रूक्षं वा}\ A; \skt{अति\appbrk{}ल\appbrk{}व\appbrk{}ण\appbrk{}म\appbrk{}ति\appbrk{}रू\appbrk{}क्षम्वा}\ K}}\skt{ }\edtext{\edtext{\skt{भेष\-ज\-मु\-प\-यु\-क्त\-म्पि\-त्ता\-निलौ}}{\lemma{\skt{भेषजम्}}\Afootnote{K H; \skt{पीतमौषधं}\ A}}}{\lemma{\skt{उपयुक्तम्}}\Afootnote{K H; \emph{om.}\ A}}\skt{ प्रदूष्य }\edtext{\skt{नाति\-ब\-स्ति\-गु\-द\-प\-रि\-क\-र्त्त\-न\-मा\-पा\-द\-यति}}{\lemma{\skt{नाति\appbrk{}ब\appbrk{}स्ति\appbrk{}गु\appbrk{}द\appbrk{}प\appbrk{}रि\appbrk{}क\appbrk{}र्त्त\appbrk{}नम्}}\Afootnote{K H; \skt{परिकर्तिकाम्}\ A}}\skt{~। }\edtext{\edtext{\skt{ताम्प\-रि\-क\-र्त्ति\-के\-त्या\-च\-क्षते}}{\lemma{\skt{ताम्}}\Afootnote{K H; \skt{तत्र गुद\appbrk{}ना\appbrk{}भि\appbrk{}मे\appbrk{}ढ्रे\appbrk{}ब\appbrk{}स्ति\appbrk{}शि\appbrk{}रःसु सदाहं}\ A}}\skt{~॥ तं}}{\lemma{\skt{परि\appbrk{}क\appbrk{}र्त्ति\appbrk{}के\appbrk{}त्या\appbrk{}च\appbrk{}क्षते~॥ तं}}\Afootnote{K H; \skt{परि\appbrk{}क\appbrk{}र्त\appbrk{}न\appbrk{}म\appbrk{}नि\appbrk{}ल\appbrk{}सङ्गो वायुविष्टम्भो भक्तारुचिश्च भवति; तत्र पिच्छा}\ A}}\skt{ }\edtext{\skt{यष्टी\-म\-धु\-क\-कृ\-ष्ण\-ति\-ल\-म\-धु\-घृ\-त\-युक्तैः}}{\lemma{\skt{यष्टी\appbrk{}म\appbrk{}धु\appbrk{}क\appbrk{}कृ\appbrk{}ष्ण\appbrk{}ति\appbrk{}ल\appbrk{}म\appbrk{}धु\appbrk{}घृ\appbrk{}त\appbrk{}युक्तैः}}\Afootnote{K H; \skt{बस्ति\appbrk{}र्य\appbrk{}ष्टी\appbrk{}म\appbrk{}धु\appbrk{}क\appbrk{}कृ\appbrk{}ष्ण\appbrk{}ति\appbrk{}ल\appbrk{}क\appbrk{}ल्क\appbrk{}म\appbrk{}धु\appbrk{}घृ\appbrk{}त\appbrk{}युक्तः,}\ A}}\skt{ }\edtext{\skt{पिच्छा\-ब\-स्ति\-भि\-रा\-स्था\-प\-येत्~। शीता\-म्बु\-प\-रि\-षि\-क्त\-ञ्चैनं}}{\lemma{\skt{पिच्छाबस्तिभिर्… शी}}\Afootnote{\emph{om.}\ A}}\skt{ पयसा भुक्तवन्तं }\edtext{\edtext{\skt{घृत\-म\-ण्ड\-न\-य\-ष्टी\-म\-धु\-क\-सि\-द्धेन}}{\lemma{\skt{घृतम्}}\Afootnote{K H; \skt{घृत}\ A}}\skt{ वा}}{\lemma{\skt{अण्ड\appbrk{}न\appbrk{}य\appbrk{}ष्टी\appbrk{}म\appbrk{}धु\appbrk{}क\appbrk{}सि\appbrk{}द्धेन वा}}\Afootnote{H; \skt{मण्डेन यष्टी\appbrk{}म\appbrk{}धु\appbrk{}क\appbrk{}सि\appbrk{}द्धेन}\ A; \skt{अण्डेन यष्टी\appbrk{}म\appbrk{}धु\appbrk{}क\appbrk{}सि\appbrk{}द्धेन वा}\ K}}\skt{ }\edtext{\skt{तैले\-ना\-नु\-वा\-स\-ये\-दिति}}{\lemma{\skt{तैले\appbrk{}ना\appbrk{}नु\appbrk{}वा\appbrk{}स\appbrk{}ये\appbrk{}दिति}}\Afootnote{K H; \skt{तैलेन वाऽनुवासयेत्~॥}\ A}}\skt{~।}\pend

\pstart \ledsidenote{\passref{4.34.17}}\edtext{\skt{क्रूर\-को\-ष्ठ\-स्या\-ति\-प्र\-भू\-त\-दो\-षस्य वा}}{\lemma{\skt{क्रूर\appbrk{}को\appbrk{}ष्ठ\appbrk{}स्या\appbrk{}ति\appbrk{}प्र\appbrk{}भू\appbrk{}त\appbrk{}दो\appbrk{}षस्य वा}}\Afootnote{K H; \skt{क्रूर\appbrk{}को\appbrk{}ष्ठ\appbrk{}स्या\appbrk{}ति\appbrk{}प्र\appbrk{}भू\appbrk{}त\appbrk{}दो\appbrk{}पस्य}\ A}}\skt{ }\edtext{\skt{मृद्व\-ल्प\-मौ\-ष\-ध\-म\-व\-चा\-रितं}}{\lemma{\skt{मृद्व\appbrk{}ल्प\appbrk{}मौ\appbrk{}ष\appbrk{}ध\appbrk{}म\appbrk{}व\appbrk{}चा\appbrk{}रितं}}\Afootnote{K H; \skt{मृद्वौ\appbrk{}ष\appbrk{}ध\appbrk{}म\appbrk{}व\appbrk{}चा\appbrk{}रितं}\ A}}\skt{ }\edtext{\skt{समुत्क्लेश्य}}{\lemma{\skt{समुत्क्लेश्य}}\Afootnote{H; \skt{समुत्क्लिश्य}\ A; \skt{समुत्क्लेष्य}\ K}}\skt{ दोषान्न निःशे\-षा\-न\-प\-ह\-रति~। }\edtext{\skt{पाकञ्चोपैति}}{\lemma{\skt{पाकञ्चोपैति}}\Afootnote{H; \skt{ततस्ते दोषाः परि\appbrk{}स्रा\appbrk{}व\appbrk{}मा\appbrk{}पा\appbrk{}द\appbrk{}यन्ति,}\ A; \skt{पाकञ्चौपैति~।}\ K}}\skt{~। तत्र }\edtext{\skt{दौर्ब्ब\-ल्यो\-द\-र\-वि\-ष्ट\-म्भा\-रु\-चि\-र्ग्गा\-त्र\-स\-द\-नानि}}{\lemma{\skt{दौर्ब्ब\appbrk{}ल्यो\appbrk{}द\appbrk{}र\appbrk{}वि\appbrk{}ष्ट\appbrk{}म्भा\appbrk{}रु\appbrk{}चि\appbrk{}र्ग्गा\appbrk{}त्र\appbrk{}स\appbrk{}द\appbrk{}नानि}}\Afootnote{H; \skt{दौर्ब\appbrk{}ल्यो\appbrk{}द\appbrk{}र\appbrk{}वि\appbrk{}ष्ट\appbrk{}म्भा\appbrk{}रु\appbrk{}चि\appbrk{}गा\appbrk{}त्र\appbrk{}स\appbrk{}द\appbrk{}नानि}\ A; \skt{दौर्ब\appbrk{}ल्यो\appbrk{}द\appbrk{}र\appbrk{}वि\appbrk{}ष्ट\appbrk{}म्भा\appbrk{}रु\appbrk{}चि\appbrk{}र्गा\appbrk{}त्र\appbrk{}स\appbrk{}द\appbrk{}नानि}\ K}}\skt{ भवन्ति~। स वेदनौ चास्य\-पि\-त्त\-श्ले\-ष्माणौ }\edtext{\edtext{\edtext{\skt{स्रव\-त\-स्त\-म्प\-रि\-स्रा\-व\-मिति}}{\lemma{\skt{स्रवतस्तम्}}\Afootnote{K H; \skt{परिस्रवतः, तं}\ A}}}{\lemma{\skt{परिस्रावम्}}\Afootnote{K H; \skt{परिस्राव}\ A}}}{\lemma{\skt{इति}}\Afootnote{K H; \skt{मित्य्}\ A}}\skt{ आचक्षते~। तम\-ज\-क\-र्ण्ण\-ध\-व\-ति\-निश }\edtext{\edtext{\edtext{\skt{पला\-स\-क\-षा\-यै\-र्म्म\-धु\-यु\-क्तै\-रा\-स्था\-प\-येत्}}{\lemma{\skt{पला\appbrk{}स\appbrk{}क\appbrk{}षा\appbrk{}यैर्}}\Afootnote{H; \skt{पला\appbrk{}श\appbrk{}ब\appbrk{}ला\appbrk{}क\appbrk{}षा\appbrk{}यैर्}\ A; \skt{पलाशकषायै र्}\ K}}}{\lemma{\skt{म्मधुयुक्तैर्}}\Afootnote{H; \skt{मधुसंयुक्तैर्}\ A; \skt{मधुक्तै}\ K}}\skt{~। }\edtext{\skt{उप\-शा\-न्त\-दो\-ष\-म्भूयः}}{\lemma{\skt{भूयः}}\Afootnote{K H; \skt{स्निग्धं च भूयः}\ A}}}{\lemma{\skt{आस्थापयेत्~। उप\appbrk{}शा\appbrk{}न्त\appbrk{}दो\appbrk{}षम्}}\Afootnote{A H; \skt{रास्थापयेत्~। उपशान्तदोषं}\ K}}\skt{ संशोधयेद्}\pend

\pstart \ledsidenote{\passref{4.34.18}}\edtext{\skt{अति\-रू\-क्षे\-ति\-स्निग्धे वा}}{\lemma{\skt{अति\appbrk{}रू\appbrk{}क्षे\appbrk{}ति\appbrk{}स्निग्धे वा}}\Afootnote{K H; \skt{अति\appbrk{}रू\appbrk{}क्षेऽ\appbrk{}ति\appbrk{}स्निग्धे वा}\ A}}\skt{ }\edtext{\edtext{\edtext{\skt{भेष\-ज\-म\-व\-चा\-रि\-त\-म\-प्राप्तं}}{\lemma{\skt{भेषजम्}}\Afootnote{K H; \skt{भेषज}\ A}}}{\lemma{\skt{अवचारितम्}}\Afootnote{K H; \skt{मवचारित}\ A}}}{\lemma{\skt{अप्राप्तं}}\Afootnote{K H; \skt{मप्राप्तं}\ A}}\skt{ }\edtext{\edtext{\skt{पुरी\-ष\-मु\-दी\-र\-यतो}}{\lemma{\skt{पुरीषम्}}\Afootnote{H; \skt{वातवर्च}\ A; \skt{परीषम्}\ K}}}{\lemma{\skt{उदीरयतो}}\Afootnote{K H; \skt{उदीरयति}\ A}}\skt{ }\edtext{\skt{विगेविघातेन}}{\lemma{\skt{विगेविघातेन}}\Afootnote{H; \skt{वेगाघातेन वा,}\ A; \skt{वेगविघातेन}\ K}}\skt{ }\edtext{\skt{प्रवाहिका}}{\lemma{\skt{प्रवाहिका}}\Afootnote{K H; \skt{तदा प्रवाहिका}\ A}}\skt{ भवति~। तत्र }\edtext{\skt{सदाह}}{\lemma{\skt{सदाह}}\Afootnote{H\textsuperscript{pc}; \skt{सवातं सदाहं सशूलं}\ A; \skt{सदाहशू}\ K; \skt{सदा}\ H\textsuperscript{ac}}}\skt{ }\edtext{\skt{शूलं}}{\lemma{\skt{शूलं}}\Afootnote{H; \skt{गुरु पिच्छिलं}\ A; \skt{लं}\ K}}\skt{ }\edtext{\skt{स्वेतं}}{\lemma{\skt{स्वेतं}}\Afootnote{H; \skt{श्वेतं}\ A K}}\skt{ कृष्णं }\edtext{\skt{रक्तम्वा}}{\lemma{\skt{रक्तम्वा}}\Afootnote{K H; \skt{सरक्तं वा}\ A}}\skt{ }\edtext{\skt{भृशं}}{\lemma{\skt{भृशं}}\Afootnote{A H; \skt{भृशां}\ K}}\skt{ प्रवाहमाणः }\edtext{\skt{कफ\-मु\-प\-वि\-शति~। }\edtext{\skt{तम्प्र\-ति\-स्रा\-व\-वि\-धा\-ने\-नो\-प\-च\-रेत्}}{\lemma{\skt{प्रति\appbrk{}स्रा\appbrk{}व\appbrk{}वि\appbrk{}धा\appbrk{}ने\appbrk{}नो\appbrk{}प\appbrk{}च\appbrk{}रेत्}}\Afootnote{K H; \skt{परि\appbrk{}स्रा\appbrk{}व\appbrk{}वि\appbrk{}धा\appbrk{}ने\appbrk{}नो\appbrk{}प\appbrk{}च\appbrk{}रेत्~॥}\ A}}}{\lemma{\skt{उपविशति~। तम्}}\Afootnote{K H; \skt{उपविशति; तां}\ A}}\pend

\pstart \ledsidenote{\passref{4.34.19}}\edtext{\edtext{\skt{यस्तूर्द्ध्वमधो}}{\lemma{\skt{यस्तूर्द्ध्वम्}}\Afootnote{K; \skt{यस्तूर्ध्व}\ A; \skt{यस्तूर्द्ध}\ H}}\skt{ वा}}{\lemma{\skt{अधो वा}}\Afootnote{K; \skt{मधो वा}\ A; \skt{मधोवा}\ H}}\skt{ }\edtext{\skt{भेषजं}}{\lemma{\skt{भेषजं}}\Afootnote{K H; \skt{भेषज}\ A}}\skt{ }\edtext{\skt{वेग\-प्र\-वृ\-त्त\-म\-ज्ञा\-ना\-द्वि\-नि\-हन्ति}}{\lemma{\skt{वेग\appbrk{}प्र\appbrk{}वृ\appbrk{}त्त\appbrk{}म\appbrk{}ज्ञा\appbrk{}ना\appbrk{}द्वि\appbrk{}नि\appbrk{}हन्ति}}\Afootnote{\skt{वेगं प्रवृ\appbrk{}त्त\appbrk{}म\appbrk{}ज्ञ\appbrk{}त्वा\appbrk{}द्वि\appbrk{}नि\appbrk{}हन्ति}\ A; \skt{वेग\appbrk{}प्र\appbrk{}वृ\appbrk{}त्त\appbrk{}म\appbrk{}ज्ञा\appbrk{}ना\appbrk{}द्वि\appbrk{}नि\appbrk{}हन्ति}\ K H}}\skt{ तस्योपसरणं हृदि कुर्व्वन्ति }\edtext{\skt{दोषस्तत्र}}{\lemma{\skt{दोषस्}}\Afootnote{K H; \skt{दोषाः,}\ A}}\skt{ }\edtext{\edtext{\skt{प्रधा\-न\-म\-र्म्म\-स\-न्ता\-पा\-द्वे\-द\-ना\-भि\-र\-त्य\-र्थ\-म्पी\-ड्य\-मा\-दन्तां}}{\lemma{\skt{प्रधा\appbrk{}न\appbrk{}म\appbrk{}र्म्म\appbrk{}स\appbrk{}न्ता\appbrk{}पा\appbrk{}द्वे\appbrk{}द\appbrk{}ना\appbrk{}भिर्}}\Afootnote{H; \skt{प्रधा\appbrk{}न\appbrk{}म\appbrk{}र्म\appbrk{}स\appbrk{}न्ता\appbrk{}पा\appbrk{}द्वे\appbrk{}द\appbrk{}ना\appbrk{}भिर्}\ A; \skt{प्रधा\appbrk{}ण\appbrk{}म\appbrk{}र्म\appbrk{}स\appbrk{}न्ता\appbrk{}पा\appbrk{}द्वे\appbrk{}द\appbrk{}ना\appbrk{}भिर्}\ K}}}{\lemma{\skt{पीड्यमादन्तां}}\Afootnote{K H; \skt{पीड्यमानो दन्ता}\ A}}\skt{ }\edtext{\skt{कट\-क\-टा\-य\-त्यु\-द्वृ\-त्ताक्षो}}{\lemma{\skt{कट\appbrk{}क\appbrk{}टा\appbrk{}य\appbrk{}त्यु\appbrk{}द्वृ\appbrk{}त्ताक्षो}}\Afootnote{K\textsuperscript{pc} H; \skt{न्किटकिटायते उद्गताक्षो}\ A; \skt{कट\appbrk{}क\appbrk{}टा\appbrk{}य\appbrk{}त्यु\appbrk{}द्वृक्षो}\ K\textsuperscript{ac}}}\skt{ जिह्वां }\edtext{\skt{खादति}}{\lemma{\skt{खादति}}\Afootnote{K H; \skt{स्वादति}\ A}}\skt{ }\edtext{\skt{प्रता\-म्य\-मा\-न\-चेता च}}{\lemma{\skt{प्रता\appbrk{}म्य\appbrk{}मा\appbrk{}न\appbrk{}चेता च}}\Afootnote{K H; \skt{प्रता\appbrk{}म्य\appbrk{}त्य\appbrk{}चे\appbrk{}ताश्च}\ A}}\skt{ }\edtext{\skt{भवति~। }\edtext{\skt{तमभ्यज्य}}{\lemma{\skt{अभ्यज्य}}\Afootnote{\skt{परिवर्जयन्ति मूर्खाः; तमभ्यज्य}\ A; \skt{तमभ्यज्य}\ K H}}}{\lemma{\skt{भवति~। तम्}}\Afootnote{A; \skt{भवति}\ K H}}\skt{ }\edtext{\skt{धानस्वेदेन}}{\lemma{\skt{धानस्वेदेन}}\Afootnote{H; \skt{धान्यस्वेदेन}\ A K}}\skt{ }\edtext{\skt{स्वेदयित्वा}}{\lemma{\skt{स्वेदयित्वा}}\Afootnote{K H; \skt{स्वेदयेत्, य}\ A}}\skt{ }\edtext{\skt{तीक्ष्णेन}}{\lemma{\skt{तीक्ष्णेन}}\Afootnote{K H; \skt{ष्टिम\appbrk{}धु\appbrk{}क\appbrk{}सि\appbrk{}द्धेन च तैलेना}\ A}}\skt{ }\edtext{\skt{द्विशिरो विरे\-च\-ने\-नो\-प\-च\-रे\-द्य\-ष्टी\-म\-धु\-क\-मि\-श्रेण}}{\lemma{\skt{द्विशिरो विरे\appbrk{}च\appbrk{}ने\appbrk{}नो\appbrk{}प\appbrk{}च\appbrk{}रे\appbrk{}द्य\appbrk{}ष्टी\appbrk{}म\appbrk{}धु\appbrk{}क\appbrk{}मि\appbrk{}श्रेण}}\Afootnote{H\textsuperscript{pc}; \skt{नुवासयेत्, शिरोविरेचनं चास्मै तीक्ष्णं विदध्यात्, ततो यष्टि\appbrk{}म\appbrk{}धु\appbrk{}क\appbrk{}मि\appbrk{}श्रेण}\ A; \skt{शिरो\appbrk{}वि\appbrk{}रे\appbrk{}च\appbrk{}ने\appbrk{}नो\appbrk{}प\appbrk{}च\appbrk{}रे\appbrk{}द्य\appbrk{}ष्टी\appbrk{}म\appbrk{}धु\appbrk{}क\appbrk{}मि\appbrk{}श्रेण}\ K; \skt{शिरो विरे\appbrk{}च\appbrk{}ने\appbrk{}नो\appbrk{}प\appbrk{}च\appbrk{}रे\appbrk{}द्य\appbrk{}ष्टी\appbrk{}म\appbrk{}धु\appbrk{}क\appbrk{}मि\appbrk{}श्रेण}\ H\textsuperscript{ac}}}\skt{ }\edtext{\edtext{\skt{चैन\-न्त\-ण्डु\-ला\-म्बुना}}{\lemma{\skt{चैनन्}}\Afootnote{K H; \emph{om.}\ A}}}{\lemma{\skt{तण्डुलाम्बुना}}\Afootnote{A H; \skt{ताण्डुलांबना}\ K}}\skt{ }\edtext{\skt{वामयेद्यथा}}{\lemma{\skt{वामयेद्यथा}}\Afootnote{K H; \skt{छर्दयेत्, यथा}\ A}}\skt{ }\edtext{\edtext{\edtext{\skt{दोष\-गु\-णो\-च्छ्रे\-य\-श्चै\-न\-म्ब\-स्ति\-भि\-रु\-पा\-च\-रे\-दिति}}{\lemma{\skt{दोष\appbrk{}गु\appbrk{}णो\appbrk{}च्छ्रे\appbrk{}यश्}}\Afootnote{\skt{दोषोच्छ्रायेण}\ A; \skt{दोषगुणोच्छ्रयं}\ K; \skt{दोषगुणोच्छ्रेय}\ H}}}{\lemma{\skt{चैनम्}}\Afootnote{A K; \skt{श्चैनम्}\ H}}}{\lemma{\skt{उपाचरेदिति}}\Afootnote{K H; \skt{उपाचरेत्~॥}\ A}}\skt{~।}\pend

\pstart \ledsidenote{\passref{4.34.20}}\edtext{\skt{यस्तूर्द्ध्वमधो}}{\lemma{\skt{यस्तूर्द्ध्वम्}}\Afootnote{A K; \skt{यस्तूर्द्धम्}\ H}}\skt{ वा प्रवृत्तदोषः }\edtext{\edtext{\edtext{\skt{शीता\-ङ्गा\-र\-मु\-द\-क\-म\-नि\-ल\-म\-न्यद्वा}}{\lemma{\skt{शीताङ्गारम्}}\Afootnote{H; \skt{शीतागारम्}\ A; \skt{शीतागार}\ K}}}{\lemma{\skt{उदकम्}}\Afootnote{A H; \skt{मुदक}\ K}}}{\lemma{\skt{अनिलम्}}\Afootnote{A H; \skt{मनिलम्}\ K}}\skt{ सेवेत तस्य दोषाः स्रोतः\-स्व\-व\-ली\-य\-माना }\edtext{\skt{घनी\-भा\-व\-मा\-पन्ना}}{\lemma{\skt{घनीभावम्}}\Afootnote{K H; \skt{घानीभावम्}\ A}}\skt{ }\edtext{\skt{वात\-मू\-त्र\-पु\-री\-ष\-ग्र\-ह\-मा\-पाद्य}}{\lemma{\skt{वात\appbrk{}मू\appbrk{}त्र\appbrk{}पु\appbrk{}री\appbrk{}ष\appbrk{}ग्र\appbrk{}हम्}}\Afootnote{H; \skt{वात\appbrk{}मू\appbrk{}त्र\appbrk{}श\appbrk{}कृ\appbrk{}द्ग्र\appbrk{}हम्}\ A; \skt{वात\appbrk{}मू\appbrk{}त्र\appbrk{}पु\appbrk{}रीश ग्रहम्}\ K}}\skt{ विबध्यन्ते~॥ तस्याटोपो }\edtext{\skt{ज्वरो}}{\lemma{\skt{ज्वरो}}\Afootnote{K H; \skt{दाहो}\ A}}\skt{ }\edtext{\skt{दाहो}}{\lemma{\skt{दाहो}}\Afootnote{K H; \skt{ज्वरो}\ A}}\skt{ }\edtext{\skt{वेदना च}}{\lemma{\skt{वेदना च}}\Afootnote{K H; \skt{वेदनाश्च}\ A}}\skt{ तीव्रा }\edtext{\skt{भवति~। }\edtext{\skt{तमाशु}}{\lemma{\skt{आशु}}\Afootnote{A; \skt{माशु}\ K H}}}{\lemma{\skt{भवति~। तम्}}\Afootnote{\skt{भवन्ति; तम्}\ A; \skt{भवति~। त}\ K H}}\skt{ }\edtext{\skt{वामयित्वा}}{\lemma{\skt{वामयित्वा}}\Afootnote{H; \skt{वाम}\ A K}}\skt{ }\edtext{\skt{प्राप्तयित्वा}}{\lemma{\skt{प्राप्तयित्वा}}\Afootnote{H; \skt{यित्वा}\ A K}}\skt{ }\edtext{\skt{प्राप्तकालं}}{\lemma{\skt{प्राप्तकालं}}\Afootnote{H; \skt{प्राप्तकालां}\ A K}}\skt{ क्रियां }\edtext{\skt{योजयेदिति}}{\lemma{\skt{योजयेदिति}}\Afootnote{K H; \skt{कुर्वीत;}\ A}}\skt{~॥ }\edtext{\skt{अधो}}{\lemma{\skt{अधो}}\Afootnote{K H; \skt{अधो}\ A}}\skt{ भागे }\edtext{\skt{त्वभा\-ग\-ह\-र\-द्रव्ये}}{\lemma{\skt{त्वभा\appbrk{}ग\appbrk{}ह\appbrk{}र\appbrk{}द्रव्ये}}\Afootnote{H; \skt{त्वधो\appbrk{}भा\appbrk{}ग\appbrk{}दो\appbrk{}ष\appbrk{}ह\appbrk{}र\appbrk{}द्रव्यं}\ A; \skt{त्वधो\appbrk{}भा\appbrk{}ग\appbrk{}ह\appbrk{}र\appbrk{}द्रव्ये}\ K}}\skt{ }\edtext{\skt{सैन्ध\-वा\-म्ल\-ल\-वण}}{\lemma{\skt{सैन्ध\appbrk{}वा\appbrk{}म्ल\appbrk{}ल\appbrk{}वण}}\Afootnote{H; \skt{सैन्धवाम्ल}\ A; \skt{सैन्ध\appbrk{}वा\appbrk{}म्व्ल\appbrk{}ल\appbrk{}वण}\ K}}\skt{ मूत्रसंसृष्टं }\edtext{\skt{पाय\-ये\-दा\-स्था\-प\-न\-म\-नु\-वा\-स\-नञ्च}}{\lemma{\skt{पाय\appbrk{}ये\appbrk{}दा\appbrk{}स्था\appbrk{}प\appbrk{}नम्}}\Afootnote{K H; \skt{विरेचनाय पाययेत्, आस्थापनम्}\ A}}\skt{ }\edtext{\skt{यथा\-दो\-ष\-म्वि\-द\-ध्या\-द्य\-था\-दो\-ष\-मा\-हा\-र\-क्र\-मश्च}}{\lemma{\skt{आहारक्रमश्च}}\Afootnote{K H; \skt{आहारक्रमं च;}\ A}}\skt{~। }\edtext{\skt{ऊर्द्ध्वभागे तु}}{\lemma{\skt{ऊर्द्ध्वभागे तु}}\Afootnote{H; \skt{उभयतो भागे तू}\ A; \skt{ऊर्ध्वभागे तू}\ K}}\skt{ }\edtext{\skt{उप\-द्र\-व\-वि\-शे\-षा\-त्य\-थास्वं}}{\lemma{\skt{उप\appbrk{}द्र\appbrk{}व\appbrk{}वि\appbrk{}शे\appbrk{}षा\appbrk{}त्य\appbrk{}थास्वं}}\Afootnote{H; \skt{पद्र\appbrk{}व\appbrk{}वि\appbrk{}शे\appbrk{}षा\appbrk{}न्य\appbrk{}थास्वं}\ A; \skt{पद्र\appbrk{}व\appbrk{}वि\appbrk{}शे\appbrk{}षा\appbrk{}द्य\appbrk{}थास्वं}\ K}}\skt{ प्रतिकुर्व्वीत~।}\pend

\pstart \ledsidenote{\passref{4.34.21}}\edtext{\skt{या तु}}{\lemma{\skt{या तु}}\Afootnote{A; \skt{या तू}\ K H\textsuperscript{ac}; \skt{या तूतु}\ H\textsuperscript{pc}}}\skt{ विरेचने }\edtext{\skt{गुद\-प\-रि\-की\-र्त्तिता}}{\lemma{\skt{गुद\appbrk{}प\appbrk{}रि\appbrk{}की\appbrk{}र्त्तिता}}\Afootnote{K H; \skt{गुद\appbrk{}प\appbrk{}रि\appbrk{}क\appbrk{}र्तिका}\ A}}\skt{ }\edtext{\skt{तद्वमाने}}{\lemma{\skt{तद्वमाने}}\Afootnote{\skt{तद्वमने}\ A; \skt{तद्वमाने}\ K H}}\skt{ }\edtext{\skt{कञ्चकषणं}}{\lemma{\skt{कञ्चकषणं}}\Afootnote{K H; \skt{कण्ठक्षणनं,}\ A}}\skt{ }\edtext{\skt{यदधः}}{\lemma{\skt{यदधः}}\Afootnote{K H; \skt{यदधः}\ A}}\skt{ परिस्रवणं स }\edtext{\edtext{\skt{ऊर्द्ध\-ङ्क\-फ\-प्र\-सेकः}}{\lemma{\skt{ऊर्द्धङ्}}\Afootnote{H; \skt{ऊर्ध्व}\ A K}}\skt{~। या}}{\lemma{\skt{कफप्रसेकः~। या}}\Afootnote{K H; \skt{भागे श्लेष्मप्रसेकः, या}\ A}}\skt{ }\edtext{\skt{त्वधः}}{\lemma{\skt{त्वधः}}\Afootnote{K H; \skt{त्वधाः}\ A}}\skt{ }\edtext{\skt{प्रवाहिका स}}{\lemma{\skt{प्रवाहिका स}}\Afootnote{H; \skt{प्रवाहिका सा}\ A; \skt{प्रवाहिकाम्}\ K}}\skt{ }\edtext{\edtext{\skt{ऊर्द्ध\-ङ्शु\-ष्को\-ङ्का\-द्गार}}{\lemma{\skt{ऊर्द्धङ्}}\Afootnote{H; \skt{तूर्ध्वं}\ A; \skt{उर्ध्वं}\ K}}}{\lemma{\skt{शुष्कोङ्काद्गार}}\Afootnote{H\textsuperscript{pc}; \skt{शुष्कोद्गारा}\ A; \skt{शुष्कोद्गार}\ K; \skt{शुष्कोङ्का}\ H\textsuperscript{ac}}}\skt{ इति~॥}\pend

\pstart\versepar \ledsidenote{\passref{4.34.22}}\edtext{\skt{भवति}}{\lemma{\skt{भवति}}\Afootnote{A; \emph{om.}\ K H}}\skt{ चात्र~॥ }\pend
\pstart\versepar \skt{यास्त्वेता व्यापदः प्रोक्ता दशपञ्च च तत्वतः~। }\pend
\pstart\versepar \skt{एता विरेकाति }\edtext{\skt{योगा}}{\lemma{\skt{योगा}}\Afootnote{K H; \skt{योग}\ A}}\skt{ दुर्योगा योगजाः स्मृताः~॥}\edtext{}{\lemma{\skt{चात्र~॥ यास्}}\Afootnote{A; \skt{भवति चात्र~॥ यास्}\ K H}}\pend

\pstart \ledsidenote{\passref{4.34.trailer}}\edtext{\skt{॥ चिकित्सा}}{\lemma{\skt{॥ चिकित्सा}}\Afootnote{H; \skt{इति सुश्रु\appbrk{}त\appbrk{}सं\appbrk{}हि\appbrk{}तायां चिकित्सा}\ A; \skt{चि ल {\symfont ❈}~॥}\ K}}\skt{ }\edtext{\skt{लह्य}}{\lemma{\skt{लह्य}}\Afootnote{H; \skt{स्थाने वम\appbrk{}न\appbrk{}वि\appbrk{}रे\appbrk{}च\appbrk{}न\appbrk{}व्या\appbrk{}प\appbrk{}च्चि\appbrk{}कि\appbrk{}त्सितं नाम चतुस्त्रिंशोऽध्यायः~॥३४~॥}\ A; \emph{om.}\ K}}\skt{~॥}\pend

\endnumbering
\beginnumbering
\eledsection*{Adhyāya 35}
\pstart \ledsidenote{\passref{4.35.1}}\skt{॥ अथातो }\edtext{\edtext{\skt{नेए\-त्र\-ब\-स्ति\-प्र\-मा\-ण\-वि\-भा\-ग\-चि\-कि\-त्सि\-त\-म्व्या\-ख्या\-स्यामः}}{\lemma{\skt{नेए\appbrk{}त्र\appbrk{}ब\appbrk{}स्ति\appbrk{}प्र\appbrk{}मा\appbrk{}ण\appbrk{}वि\appbrk{}भा\appbrk{}ग\appbrk{}चि\appbrk{}कि\appbrk{}त्सि\appbrk{}तम्}}\Afootnote{H; \skt{नेत्र\appbrk{}ब\appbrk{}स्ति\appbrk{}प्र\appbrk{}मा\appbrk{}ण\appbrk{}प्र\appbrk{}वि\appbrk{}भा\appbrk{}ग\appbrk{}चि\appbrk{}कि\appbrk{}त्सितं}\ A; \skt{नेत्र\appbrk{}ब\appbrk{}स्ति\appbrk{}प्र\appbrk{}मा\appbrk{}ण\appbrk{}वि\appbrk{}भा\appbrk{}ग\appbrk{}चि\appbrk{}कि\appbrk{}त्सितं व्या ख्यास्यामः~॥}\ K}}}{\lemma{\skt{व्याख्यास्यामः}}\Afootnote{A H; \emph{om.}\ K}}\skt{~॥}\pend

\pstart \ledsidenote{\passref{4.35.3}}\edtext{\skt{अत्र}}{\lemma{\skt{अत्र}}\Afootnote{K H; \skt{तत्र}\ A}}\skt{ }\edtext{\edtext{\skt{स्नेहा\-दी\-ना\-ङ्क\-र्म्म\-णा\-म्ब\-स्ति\-क\-र्म्म\-प्र\-धा\-न\-त\-म\-मा\-च\-क्षते}}{\lemma{\skt{बस्ति\appbrk{}क\appbrk{}र्म्म\appbrk{}प्र\appbrk{}धा\appbrk{}न\appbrk{}त\appbrk{}मम्}}\Afootnote{H; \skt{बस्तिकर्म प्रधानतम}\ A; \skt{बस्ति\appbrk{}क\appbrk{}र्म\appbrk{}प्र\appbrk{}धा\appbrk{}ण\appbrk{}त\appbrk{}मम्}\ K}}}{\lemma{\skt{आचक्षते}}\Afootnote{K H; \skt{माहुराचार्याः~।}\ A}}\skt{~। }\edtext{\skt{कस्मा\-द\-ने\-क\-क\-र्म्म\-प्र\-का\-र\-त्वात्~। ब स्तिरिहखलु}}{\lemma{\skt{कस्मा\appbrk{}द\appbrk{}ने\appbrk{}क\appbrk{}क\appbrk{}र्म्म\appbrk{}प्र\appbrk{}का\appbrk{}र\appbrk{}त्वात्~। ब स्तिरिहखलु}}\Afootnote{K H; \skt{कस्मा\appbrk{}त\appbrk{}ने\appbrk{}क\appbrk{}क\appbrk{}र्म\appbrk{}क\appbrk{}र\appbrk{}त्वा\appbrk{}द्बस्तेः; इह खलु}\ A}}\skt{ }\edtext{\skt{नाना\-वि\-ध\-द्र\-व्य\-सं\-यो\-गा\-द्दो\-षाणां}}{\lemma{\skt{नाना\appbrk{}वि\appbrk{}ध\appbrk{}द्र\appbrk{}व्य\appbrk{}सं\appbrk{}यो\appbrk{}गा\appbrk{}द्दो\appbrk{}षाणां}}\Afootnote{K H; \skt{बस्ति\appbrk{}र्ना\appbrk{}ना\appbrk{}वि\appbrk{}ध\appbrk{}द्र\appbrk{}व्य\appbrk{}सं\appbrk{}यो\appbrk{}गा\appbrk{}द्दो\appbrk{}षाणां}\ A}}\skt{ }\edtext{\skt{संशोधनं}}{\lemma{\skt{संशोधनं}}\Afootnote{H; \skt{संशोधन}\ A; \skt{संशोधण}\ K}}\skt{ संश\-म\-न\-स\-ङ्ग्र\-ह\-णानि करोति~। }\edtext{\edtext{\skt{क्षीण\-म्व्या\-जी\-क\-रोति}}{\lemma{\skt{क्षीणम्}}\Afootnote{K H; \skt{क्षीण}\ A}}}{\lemma{\skt{व्याजीकरोति}}\Afootnote{K H; \skt{शुक्रं वाजीकरोति,}\ A}}\skt{~। कृशं बृंहयति~। स्थूलं }\edtext{\skt{कर्षयति}}{\lemma{\skt{कर्षयति}}\Afootnote{H; \skt{कर्शयति,}\ A; \skt{कर्मयति~।}\ K}}\skt{~। चक्षुः प्रीणयति~। }\edtext{\edtext{\skt{बली\-प\-लि\-त\-मु\-प\-हन्ति}}{\lemma{\skt{बलीपलितम्}}\Afootnote{K\textsuperscript{pc} H; \skt{वलीपलित}\ A; \skt{पलितम्}\ K\textsuperscript{ac}}}}{\lemma{\skt{उपहन्ति}}\Afootnote{K H; \skt{मपहन्ति,}\ A}}\skt{~। वयः स्थापयति }\edtext{\skt{इति}}{\lemma{\skt{इति}}\Afootnote{K H; \emph{om.}\ A}}\skt{~।}\pend

\pstart\versepar \ledsidenote{\passref{4.35.4}}\skt{शरीरोपचयं }\edtext{\skt{वर्ण्ण\-ब\-ल\-मा\-रो\-ग्य\-मा\-युषः}}{\lemma{\skt{वर्ण्ण\appbrk{}ब\appbrk{}ल\appbrk{}मा\appbrk{}रो\appbrk{}ग्य\appbrk{}मा\appbrk{}युषः}}\Afootnote{H; \skt{वर्णं बल\appbrk{}मा\appbrk{}रो\appbrk{}ग्य\appbrk{}मा\appbrk{}युषः~।}\ A; \skt{वर्ण्ण\appbrk{}प\appbrk{}ल\appbrk{}मा\appbrk{}रो\appbrk{}ग्य\appbrk{}मा\appbrk{}युषः}\ K}}\skt{ }\pend
\pstart\versepar \skt{}\edtext{\skt{परिवृद्धिञ्च}}{\lemma{\skt{परिवृद्धिञ्च}}\Afootnote{H; \skt{कुरुते परिवृद्धिं च}\ A; \skt{परिवृद्धिश्च}\ K}}\skt{ }\edtext{\skt{करोति}}{\lemma{\skt{करोति}}\Afootnote{K H; \emph{om.}\ A}}\skt{ }\edtext{\skt{बस्ति}}{\lemma{\skt{बस्ति}}\Afootnote{K H; \skt{बस्तिः}\ A}}\skt{ }\edtext{\skt{सम्यगुपासिताः}}{\lemma{\skt{सम्यगुपासिताः}}\Afootnote{H; \skt{सम्यगुपासितः}\ A K}}\skt{~।}\pend

\pstart \ledsidenote{\passref{4.35.5}}\skt{तथा ज्वरा\-ती\-सा\-र\-ति\-मि\-र\-प्र\-ति\-श्या\-य\-शि\-रो\-रो\-गा\-धि\-म\-न्था\-र्द्दिता }\edtext{\skt{क्षेप\-का\-घा\-तै\-का\-ङ्ग\-स\-र्व्वा\-ङ्ग\-रो\-गा\-ध्मा\-नो\-द\-र\-शू\-ल\-वृ\-द्ध्यु\-प\-दं\-शा\-ना\-ह\-मू\-त्र\-कृ\-च्छ्र\-गु\-ल्म\-वा\-त\-शो\-णि\-त\-वा\-त\-मू\-त्र\-पु\-री\-ष\-स\-ङ्ग\-शु\-क्रो\-दा\-व\-र्त्त\-शु\-क्रा\-र्त्त\-ब\-स्त\-न्य\-ना\-श\-हृ\-द्ध\-नु\-म\-न्या\-ग्र\-ह\-ण\-श\-र्क्क\-रा\-श्म\-रि\-मू\-ढ\-ग\-र्ब्भ\-प्र\-भृ\-तिषु च}}{\lemma{\skt{क्षेप\appbrk{}का\appbrk{}घा\appbrk{}तै\appbrk{}का\appbrk{}ङ्ग\appbrk{}स\appbrk{}र्व्वा\appbrk{}ङ्ग\appbrk{}रो\appbrk{}गा\appbrk{}ध्मा\appbrk{}नो\appbrk{}द\appbrk{}र\appbrk{}शू\appbrk{}ल\appbrk{}वृ\appbrk{}द्ध्यु\appbrk{}प\appbrk{}दं\appbrk{}शा\appbrk{}ना\appbrk{}ह\appbrk{}मू\appbrk{}त्र\appbrk{}कृ\appbrk{}च्छ्र\appbrk{}गु\appbrk{}ल्म\appbrk{}वा\appbrk{}त\appbrk{}शो\appbrk{}णि\appbrk{}त\appbrk{}वा\appbrk{}त\appbrk{}मू\appbrk{}त्र\appbrk{}पु\appbrk{}री\appbrk{}ष\appbrk{}स\appbrk{}ङ्ग\appbrk{}शु\appbrk{}क्रो\appbrk{}दा\appbrk{}व\appbrk{}र्त्त\appbrk{}शु\appbrk{}क्रा\appbrk{}र्त्त\appbrk{}ब\appbrk{}स्त\appbrk{}न्य\appbrk{}ना\appbrk{}श\appbrk{}हृ\appbrk{}द्ध\appbrk{}नु\appbrk{}म\appbrk{}न्या\appbrk{}ग्र\appbrk{}ह\appbrk{}ण\appbrk{}श\appbrk{}र्क्क\appbrk{}रा\appbrk{}श्म\appbrk{}रि\appbrk{}मू\appbrk{}ढ\appbrk{}ग\appbrk{}र्ब्भ\appbrk{}प्र\appbrk{}भृ\appbrk{}तिषु च}}\Afootnote{\skt{क्षेप\appbrk{}क\appbrk{}प\appbrk{}क्षा\appbrk{}घा\appbrk{}तै\appbrk{}का\appbrk{}ङ्ग\appbrk{}स\appbrk{}र्वा\appbrk{}ङ्ग\appbrk{}रो\appbrk{}गा\appbrk{}ध्मा\appbrk{}नो\appbrk{}द\appbrk{}र\appbrk{}यो\appbrk{}नि\appbrk{}शू\appbrk{}ल\appbrk{}श\appbrk{}र्क\appbrk{}रा\appbrk{}शू\appbrk{}ल\appbrk{}वृ\appbrk{}द्ध्यु\appbrk{}प\appbrk{}दं\appbrk{}शा\appbrk{}ना\appbrk{}ह\appbrk{}मू\appbrk{}त्र\appbrk{}कृ\appbrk{}च्छ्र\appbrk{}गु\appbrk{}ल्म\appbrk{}वा\appbrk{}त\appbrk{}शो\appbrk{}णि\appbrk{}त\appbrk{}वा\appbrk{}त\appbrk{}मू\appbrk{}त्र\appbrk{}पु\appbrk{}री\appbrk{}षो\appbrk{}दा\appbrk{}व\appbrk{}र्त\appbrk{}शु\appbrk{}क्रा\appbrk{}र्त\appbrk{}व\appbrk{}स्त\appbrk{}न्य\appbrk{}ना\appbrk{}श\appbrk{}हृ\appbrk{}द्ध\appbrk{}नु\appbrk{}म\appbrk{}न्या\appbrk{}ग्र\appbrk{}ह\appbrk{}श\appbrk{}र्क\appbrk{}रा\appbrk{}श्म\appbrk{}री\appbrk{}मू\appbrk{}ढ\appbrk{}ग\appbrk{}र्भ\appbrk{}प्र\appbrk{}भृ\appbrk{}तिषु}\ A; \skt{क्षेप\appbrk{}का\appbrk{}घा\appbrk{}तै\appbrk{}कांग सर्वा\appbrk{}ङ्ग\appbrk{}रो\appbrk{}गा\appbrk{}ध्मा\appbrk{}नो\appbrk{}द\appbrk{}र\appbrk{}शू\appbrk{}ल\appbrk{}वृ\appbrk{}द्ध्युप दम्शा\appbrk{}ना\appbrk{}ह\appbrk{}मू\appbrk{}त्र\appbrk{}कृ\appbrk{}च्छ्र\appbrk{}गु\appbrk{}ल्म\appbrk{}वा\appbrk{}त\appbrk{}शो\appbrk{}णि\appbrk{}त\appbrk{}वा\appbrk{}त\appbrk{}मू\appbrk{}त्र\appbrk{}पु\appbrk{}री\appbrk{}श\appbrk{}सं\appbrk{}ग\appbrk{}शु\appbrk{}क्रो\appbrk{}दा\appbrk{}व\appbrk{}र्त्त\appbrk{}शु\appbrk{}क्रा\appbrk{}र्त्त\appbrk{}ब\appbrk{}स्त\appbrk{}न्य\appbrk{}ना\appbrk{}श\appbrk{}हृ\appbrk{}द्ध\appbrk{}नु\appbrk{}म\appbrk{}न्या\appbrk{}ग्र\appbrk{}ह\appbrk{}श\appbrk{}र्क\appbrk{}रा\appbrk{}र्श्म\appbrk{}री\appbrk{}मू\appbrk{}ढ\appbrk{}ग\appbrk{}र्भ\appbrk{}प्र\appbrk{}भृ\appbrk{}तिषु च}\ K; \skt{क्षेप\appbrk{}का\appbrk{}घा\appbrk{}तै\appbrk{}काङ्ग सर्व्वा\appbrk{}ङ्ग\appbrk{}रो\appbrk{}गा\appbrk{}ध्मा\appbrk{}नो\appbrk{}द\appbrk{}र\appbrk{}शू\appbrk{}लवृ द्ध्युप\appbrk{}दं\appbrk{}शा\appbrk{}ना\appbrk{}ह\appbrk{}मू\appbrk{}त्र\appbrk{}कृ\appbrk{}च्छ्र\appbrk{}गु\appbrk{}ल्म\appbrk{}वा\appbrk{}त\appbrk{}शो\appbrk{}णि\appbrk{}त\appbrk{}वा\appbrk{}त\appbrk{}मू\appbrk{}त्र\appbrk{}पु\appbrk{}री\appbrk{}ष\appbrk{}स\appbrk{}ङ्ग\appbrk{}शु\appbrk{}क्रो\appbrk{}दा\appbrk{}व\appbrk{}र्त्त\appbrk{}शु\appbrk{}क्रा\appbrk{}र्त्त\appbrk{}ब\appbrk{}स्त\appbrk{}न्य\appbrk{}ना\appbrk{}श\appbrk{}हृ\appbrk{}द्ध\appbrk{}नु\appbrk{}म\appbrk{}न्या\appbrk{}ग्र\appbrk{}ह\appbrk{}ण\appbrk{}श\appbrk{}र्क्क\appbrk{}रा\appbrk{}श्म\appbrk{}रि\appbrk{}मू\appbrk{}ढ\appbrk{}ग\appbrk{}र्ब्भ\appbrk{}प्र\appbrk{}भृ\appbrk{}तिषु च}\ H}}\skt{ }\edtext{\edtext{\edtext{\skt{विका\-रे\-ष्व\-त्य\-र्थ\-मु\-प\-यु\-ज्यत}}{\lemma{\skt{विकारेष्व्}}\Afootnote{K H; \skt{चा}\ A}}}{\lemma{\skt{अत्यर्थम्}}\Afootnote{K H; \skt{स्यर्थम्}\ A}}}{\lemma{\skt{उपयुज्यत}}\Afootnote{K H; \skt{उपयुज्यते~॥}\ A}}\skt{ }\edtext{\skt{इति}}{\lemma{\skt{इति}}\Afootnote{K H; \emph{om.}\ A}}\skt{~॥}\pend

\pstart\versepar \ledsidenote{\passref{4.35.6}}\edtext{\skt{बस्तिर्व्वाते च}}{\lemma{\skt{बस्तिर्व्वाते च}}\Afootnote{K H; \skt{भवति चात्र बस्तिर्वाते च}\ A}}\skt{ पित्ते च कफे रक्ते च }\edtext{\skt{पूर्यते}}{\lemma{\skt{पूर्यते}}\Afootnote{H; \skt{शस्यते~।}\ A; \skt{पूज्यते~।}\ K}}\skt{~। }\pend
\pstart\versepar \skt{संसर्ग्गे सन्निपाते च }\edtext{\edtext{\skt{बस्तिरेव}}{\lemma{\skt{बस्तिर्}}\Afootnote{A; \skt{बस्ति}\ K H}}}{\lemma{\skt{एव}}\Afootnote{A; \skt{रेव}\ K H}}\skt{ }\edtext{\skt{सदा}}{\lemma{\skt{सदा}}\Afootnote{K H; \emph{om.}\ A}}\skt{ }\edtext{\skt{हितः}}{\lemma{\skt{हितः}}\Afootnote{K H; \skt{हितः सदा~॥}\ A}}\skt{~॥}\pend

\pstart \ledsidenote{\passref{4.35.7}}\skt{तत्र }\edtext{\skt{सम्व\-त्स\-रि\-का\-ष्ट\-द्वि\-र\-ष्ट\-व\-र्षाणां}}{\lemma{\skt{सम्व\appbrk{}त्स\appbrk{}रि\appbrk{}का\appbrk{}ष्ट\appbrk{}द्वि\appbrk{}र\appbrk{}ष्ट\appbrk{}व\appbrk{}र्षाणां}}\Afootnote{K H; \skt{सांव\appbrk{}त्स\appbrk{}रि\appbrk{}का\appbrk{}ष्ट\appbrk{}द्वि\appbrk{}र\appbrk{}ष्ट\appbrk{}व\appbrk{}र्षाणां}\ A}}\skt{ षड\-ष्ट\-द\-शा\-ङ्गु\-ल\-प्र\-मा\-णानि }\edtext{\skt{कनी\-नि\-का\-ना\-मि\-का\-म\-ध्य\-मा\-ङ्गु\-लि\-प\-रि\-णा\-हानि}}{\lemma{\skt{कनी\appbrk{}नि\appbrk{}का\appbrk{}ना\appbrk{}मि\appbrk{}का\appbrk{}म\appbrk{}ध्य\appbrk{}मा\appbrk{}ङ्गु\appbrk{}लि\appbrk{}प\appbrk{}रि\appbrk{}णा\appbrk{}हानि}}\Afootnote{K H; \skt{कनि\appbrk{}ष्ठि\appbrk{}का\appbrk{}ना\appbrk{}मि\appbrk{}का\appbrk{}म\appbrk{}ध्य\appbrk{}मा\appbrk{}ङ्गु\appbrk{}लि\appbrk{}प\appbrk{}रि\appbrk{}णाहा}\ A}}\skt{ }\edtext{\edtext{\skt{अग्रे\-ङ्गु\-ला\-ध्य\-र्द्धा\-ङ्गु\-ल\-ध्य\-र्द्धा\-ङ्गु\-ल\-त्र्य\-ङ्गु\-ल\-स\-न्नि\-वि\-ष्ट\-क\-र्ण्णि\-कानि}}{\lemma{\skt{गुल\appbrk{}ध्य\appbrk{}र्द्धा\appbrk{}ङ्गु\appbrk{}ल\appbrk{}त्र्यङ्गु}}\Afootnote{\emph{om.}\ K\textsuperscript{ac}}}}{\lemma{\skt{अग्रे\appbrk{}ङ्गु\appbrk{}ला\appbrk{}ध्य\appbrk{}र्द्धा\appbrk{}ङ्गु\appbrk{}ल\appbrk{}ध्य\appbrk{}र्द्धा\appbrk{}ङ्गु\appbrk{}ल\appbrk{}त्र्य\appbrk{}ङ्गु\appbrk{}ल\appbrk{}स\appbrk{}न्नि\appbrk{}वि\appbrk{}ष्ट\appbrk{}क\appbrk{}र्ण्णि\appbrk{}कानि}}\Afootnote{H; \skt{न्यग्रेऽ\appbrk{}ध्य\appbrk{}र्धा\appbrk{}ङ्गु\appbrk{}ल\appbrk{}द्व्य\appbrk{}ङ्गु\appbrk{}ला\appbrk{}र्ध\appbrk{}तृ\appbrk{}ती\appbrk{}या\appbrk{}ङ्गु\appbrk{}ल\appbrk{}स\appbrk{}न्नि\appbrk{}वि\appbrk{}ष्ट\appbrk{}क\appbrk{}र्णि\appbrk{}कानि}\ A; \skt{अग्रे\appbrk{}ङ्गु\appbrk{}ला\appbrk{}ध्य\appbrk{}र्द्धां\appbrk{}गु\appbrk{}ल\appbrk{}स\appbrk{}न्नि\appbrk{}वि\appbrk{}ष्ट\appbrk{}क\appbrk{}र्ण्णि\appbrk{}कानि}\ K\textsuperscript{ac}; \skt{अग्रे\appbrk{}ङ्गु\appbrk{}ला\appbrk{}ध्य\appbrk{}र्द्धां\appbrk{}गु\appbrk{}ल\appbrk{}द्व्य\appbrk{}ङ्गु\appbrk{}ल\appbrk{}स\appbrk{}न्नि\appbrk{}वि\appbrk{}ष्ट\appbrk{}क\appbrk{}र्ण्णि\appbrk{}कानि}\ K\textsuperscript{pc}}}\skt{ }\edtext{\skt{श्येन\-ब\-र्हि\-क\-ङ्क\-गृ\-ध्र\-प\-त्र\-ना\-डी\-तुल्य}}{\lemma{\skt{श्येन\appbrk{}ब\appbrk{}र्हि\appbrk{}क\appbrk{}ङ्क\appbrk{}गृ\appbrk{}ध्र\appbrk{}प\appbrk{}त्र\appbrk{}ना\appbrk{}डी\appbrk{}तुल्य}}\Afootnote{H\textsuperscript{ac}; \skt{कङ्क\appbrk{}श्ये\appbrk{}न\appbrk{}ब\appbrk{}र्हि\appbrk{}ण\appbrk{}प\appbrk{}क्ष\appbrk{}ना\appbrk{}डी\appbrk{}तुल्य}\ A; \skt{श्येन\appbrk{}ब\appbrk{}र्हि\appbrk{}कं\appbrk{}क\appbrk{}प\appbrk{}त्र\appbrk{}ना\appbrk{}डी\appbrk{}तुल्य}\ K H\textsuperscript{pc}}}\skt{ }\edtext{\skt{प्रवेशानानि}}{\lemma{\skt{प्रवेशानानि}}\Afootnote{H; \skt{प्रवेशानि}\ A K}}\skt{ मुद्गमाषक }\edtext{\skt{लाय\-मा\-त्रा\-स्रो\-तांसि}}{\lemma{\skt{लाय\appbrk{}मा\appbrk{}त्रा\appbrk{}स्रो\appbrk{}तांसि}}\Afootnote{H; \skt{लाय\appbrk{}मा\appbrk{}त्र\appbrk{}स्रो\appbrk{}तांसि}\ A K}}\skt{ }\edtext{\skt{विद\-ध्या\-न्ने\-त्रा\-णीति}}{\lemma{\skt{नेत्राणीति}}\Afootnote{K H; \skt{नेत्राणि~।}\ A}}\skt{~। तेषु }\edtext{\edtext{\skt{चास्था\-प\-न\-द्र\-व्य\-प्र\-मा\-ण\-मा\-तु\-र\-ह\-स्त\-स\-म्मि\-तौ\-स\-न्मिश्रौ}}{\lemma{\skt{चास्था\appbrk{}प\appbrk{}न\appbrk{}द्र\appbrk{}व्य\appbrk{}प्र\appbrk{}मा\appbrk{}णम्}}\Afootnote{K; \skt{चास्था\appbrk{}प\appbrk{}न\appbrk{}द्र\appbrk{}व्य\appbrk{}प्र\appbrk{}माण}\ A H}}}{\lemma{\skt{आतु\appbrk{}र\appbrk{}ह\appbrk{}स्त\appbrk{}स\appbrk{}म्मि\appbrk{}तौ\appbrk{}स\appbrk{}न्मिश्रौ}}\Afootnote{\skt{मातु\appbrk{}र\appbrk{}ह\appbrk{}स्त\appbrk{}सं\appbrk{}मि\appbrk{}तेन प्रसृतेन संमितौ}\ A; \skt{आतु\appbrk{}र\appbrk{}ह\appbrk{}स्त\appbrk{}स\appbrk{}म्मितौ}\ K; \skt{मातु\appbrk{}र\appbrk{}ह\appbrk{}स्त\appbrk{}स\appbrk{}म्मितौ}\ H\textsuperscript{ac}; \skt{मातु\appbrk{}र\appbrk{}ह\appbrk{}स्त\appbrk{}स\appbrk{}म्मि\appbrk{}तौ\appbrk{}स\appbrk{}न्मिश्रौ}\ H\textsuperscript{pc}}}\skt{ प्रसृतौ द्वौ }\edtext{\skt{चत्वाराष्टौ}}{\lemma{\skt{चत्वाराष्टौ}}\Afootnote{K H; \skt{चत्वारोऽष्टौ च}\ A}}\skt{ }\edtext{\skt{विधेया}}{\lemma{\skt{विधेया}}\Afootnote{K H; \skt{विधेयाः~॥}\ A}}\skt{~।}\pend

\pstart\versepar \ledsidenote{\passref{4.35.8}}\edtext{\skt{वर्षान्तरेषु}}{\lemma{\skt{वर्षान्तरेषु}}\Afootnote{A K H\textsuperscript{pc}; \skt{वर्षान्तषु}\ H\textsuperscript{ac}}}\skt{ नेत्राणां बस्तिमानस्य }\edtext{\edtext{\skt{चाप्यथ}}{\lemma{\skt{चाप्य्}}\Afootnote{K; \skt{चैव}\ A; \skt{चा}\ H}}}{\lemma{\skt{अथ}}\Afootnote{K; \skt{हि~।}\ A; \skt{प्यथ~।}\ H}}\skt{~। }\pend
\pstart\versepar \skt{}\edtext{\skt{वयो}}{\lemma{\skt{वयो}}\Afootnote{H; \skt{वयो}\ A K}}\skt{ बलशरीराणि }\edtext{\skt{समी\-क्ष्यो\-त्क\-र्ष\-ये\-द्विधिं}}{\lemma{\skt{समी\appbrk{}क्ष्यो\appbrk{}त्क\appbrk{}र्ष\appbrk{}ये\appbrk{}द्विधिं}}\Afootnote{A H; \skt{समी\appbrk{}क्यो\appbrk{}त्क\appbrk{}र्ष\appbrk{}ये\appbrk{}द्विधिं~॥}\ K}}\skt{~॥}\pend

\pstart \ledsidenote{\passref{4.35.9}}\edtext{\edtext{\edtext{\edtext{\skt{पञ्च\-विं\-श\-ति\-व\-र्षा\-दू\-र्द्ध्व\-न्द्वा\-द\-शा\-ङ्गु\-ल\-ल\-प्र\-मा\-ण\-मं\-गु\-ष्ठो\-द\-र\-प\-रि\-णा\-ह\-मग्रे}}{\lemma{\skt{गुल\appbrk{}ल\appbrk{}प्र\appbrk{}मा\appbrk{}ण\appbrk{}मं\appbrk{}गुष्ठो}}\Afootnote{\emph{om.}\ K}}}{\lemma{\skt{पञ्च\appbrk{}विं\appbrk{}श\appbrk{}ति\appbrk{}व\appbrk{}र्षा\appbrk{}दू\appbrk{}र्द्ध्वन्}}\Afootnote{H; \skt{पञ्च\appbrk{}विं\appbrk{}श\appbrk{}ते\appbrk{}रूर्ध्वं}\ A; \skt{पञ्च\appbrk{}विं\appbrk{}श\appbrk{}ति\appbrk{}व\appbrk{}र्षा\appbrk{}द\appbrk{}दूर्ध्वं}\ K}}}{\lemma{\skt{द्वाद\appbrk{}शा\appbrk{}ङ्गु\appbrk{}ल\appbrk{}ल\appbrk{}प्र\appbrk{}मा\appbrk{}णम्}}\Afootnote{H; \skt{द्वादशाङ्गुलं, मूले}\ A; \skt{द्वादशां}\ K}}}{\lemma{\skt{अंगु\appbrk{}ष्ठो\appbrk{}द\appbrk{}र\appbrk{}प\appbrk{}रि\appbrk{}णा\appbrk{}हम्}}\Afootnote{H; \skt{ऽङ्गु\appbrk{}ष्ठो\appbrk{}द\appbrk{}र\appbrk{}प\appbrk{}री\appbrk{}णा\appbrk{}हम्, अग्रे कनि\appbrk{}ष्ठि\appbrk{}को\appbrk{}द\appbrk{}र\appbrk{}प\appbrk{}री\appbrk{}णा\appbrk{}हम्,}\ A}}\skt{ त्र्यङ्गु\-ल\-स\-न्नि\-वि\-ष्ट\-क\-र्ण्णिकं गृध्र }\edtext{\skt{पत्र\-ना\-डी\-तु\-ल्य\-प्र\-वेशं}}{\lemma{\skt{पत्र\appbrk{}ना\appbrk{}डी\appbrk{}तु\appbrk{}ल्य\appbrk{}प्र\appbrk{}वेशं}}\Afootnote{K H; \skt{पक्ष\appbrk{}ना\appbrk{}डीं\appbrk{}तु\appbrk{}ल्य\appbrk{}प्र\appbrk{}वेशं,}\ A}}\skt{ }\edtext{\skt{कोलास्थिमात्रं}}{\lemma{\skt{कोलास्थिमात्रं}}\Afootnote{H; \skt{कोलास्थिमात्र}\ A K}}\skt{ }\edtext{\skt{च्छिद्रं}}{\lemma{\skt{च्छिद्रं}}\Afootnote{K H; \skt{छिद्रं,}\ A}}\skt{ }\edtext{\edtext{\skt{क्लिन्न\-क\-ला\-य\-मा\-त्र\-च्छि\-द्र\-मि\-त्येके}}{\lemma{\skt{क्लिन्न\appbrk{}क\appbrk{}ला\appbrk{}य\appbrk{}मा\appbrk{}त्र\appbrk{}च्छि\appbrk{}द्रम्}}\Afootnote{K H; \skt{क्लिन्न\appbrk{}क\appbrk{}ला\appbrk{}य\appbrk{}मा\appbrk{}त्र\appbrk{}छिद्र}\ A}}}{\lemma{\skt{इत्य्}}\Afootnote{K H; \skt{मित्य्}\ A}}\skt{~॥ }\edtext{\skt{आस्था\-प\-न\-मा\-त्र\-प्र\-माणं}}{\lemma{\skt{आस्था\appbrk{}प\appbrk{}न\appbrk{}मा\appbrk{}त्र\appbrk{}प्र\appbrk{}माणं}}\Afootnote{H; \skt{सर्वाणि मूले बस्ति\appbrk{}नि\appbrk{}ब\appbrk{}न्ध\appbrk{}नार्थं द्विकर्णिकानि~। आस्था\appbrk{}प\appbrk{}न\appbrk{}द्र\appbrk{}व्य\appbrk{}प्र\appbrk{}णाणं तु}\ A; \skt{आस्था\appbrk{}प\appbrk{}ण\appbrk{}मा\appbrk{}त्रा\appbrk{}प्रमा}\ K}}\skt{ }\edtext{\skt{द्वादशप्रसृता}}{\lemma{\skt{द्वादशप्रसृता}}\Afootnote{H; \skt{विहितं द्वादशप्रसृताः~।}\ A; \skt{नं द्वादशप्रसृता}\ K}}\skt{ }\edtext{\skt{इति}}{\lemma{\skt{इति}}\Afootnote{K H; \emph{om.}\ A}}\skt{~। }\edtext{\skt{सप्तसप्त्यान्यास्तु}}{\lemma{\skt{सप्तसप्त्यान्यास्तु}}\Afootnote{H\textsuperscript{pc}; \skt{सप्तते}\ A; \skt{सप्तत्या}\ K; \skt{सप्तन्यास्तु}\ H\textsuperscript{ac}}}\skt{ }\edtext{\skt{ऊर्द्ध\-न्ने\-त्र\-प्र\-मा\-ण\-मे\-त\-देव}}{\lemma{\skt{ऊर्द्ध\appbrk{}न्ने\appbrk{}त्र\appbrk{}प्र\appbrk{}मा\appbrk{}ण\appbrk{}मे\appbrk{}त\appbrk{}देव}}\Afootnote{\skt{स्तूर्ध्वं नेत्र\appbrk{}प्र\appbrk{}मा\appbrk{}ण\appbrk{}मे\appbrk{}त\appbrk{}देव,}\ A; \skt{स्तू र्ध्वं नेत्र\appbrk{}प्र\appbrk{}मा\appbrk{}ण\appbrk{}मे\appbrk{}त\appbrk{}देव~।}\ K; \skt{ऊर्द्ध\appbrk{}न्ने\appbrk{}त्र\appbrk{}प्र\appbrk{}मा\appbrk{}ण\appbrk{}मे\appbrk{}त\appbrk{}देव~।}\ H}}\skt{~। }\edtext{\skt{आस्था\-प\-न\-मा\-त्रा\-प्र\-मा\-णन्तु}}{\lemma{\skt{आस्था\appbrk{}प\appbrk{}न\appbrk{}मा\appbrk{}त्रा\appbrk{}प्र\appbrk{}मा\appbrk{}णन्तु}}\Afootnote{K H; \skt{द्रव्यप्रमाणं तु}\ A}}\skt{ }\edtext{\skt{द्विर\-ष्ट\-व\-र्ष\-व\-दिति}}{\lemma{\skt{द्विर\appbrk{}ष्ट\appbrk{}व\appbrk{}र्ष\appbrk{}व\appbrk{}दिति}}\Afootnote{K H; \skt{द्विर\appbrk{}ष्ट\appbrk{}व\appbrk{}र्ष\appbrk{}वत्~॥}\ A}}\skt{~॥}\pend

\pstart\versepar \ledsidenote{\passref{4.35.10}}\edtext{\skt{मृदु\-र्ब्ब\-स्ति\-प्र\-यो\-क्तव्या}}{\lemma{\skt{मृदु\appbrk{}र्ब्ब\appbrk{}स्ति\appbrk{}प्र\appbrk{}यो\appbrk{}क्तव्या}}\Afootnote{H; \skt{मृदुर्बस्तिः प्रयोक्तव्यो}\ A; \skt{मृदुर्बस्ति प्रयोक्तव्या}\ K}}\skt{ }\edtext{\skt{विशे\-षा\-द्बा\-ल\-वृ\-द्धयोः}}{\lemma{\skt{विशे\appbrk{}षा\appbrk{}द्बा\appbrk{}ल\appbrk{}वृ\appbrk{}द्धयोः}}\Afootnote{K H; \skt{विशे\appbrk{}षा\appbrk{}द्बा\appbrk{}ल\appbrk{}वृ\appbrk{}द्धयोः~।}\ A}}\skt{~। }\pend
\pstart\versepar \skt{}\edtext{\skt{तयोस्तीक्ष्णः}}{\lemma{\skt{तयोस्तीक्ष्णः}}\Afootnote{A H; \skt{तयोस्तीःक्ष्णः}\ K}}\skt{ प्रयुक्तस्तु }\edtext{\skt{बस्ति\-हिं\-स्या\-द्ब\-लौ\-जसी}}{\lemma{\skt{बस्ति\appbrk{}हिं\appbrk{}स्या\appbrk{}द्ब\appbrk{}लौ\appbrk{}जसी}}\Afootnote{K H; \skt{बस्ति\appbrk{}र्हिं\appbrk{}स्या\appbrk{}द्ब\appbrk{}ला\appbrk{}युषी~॥}\ A}}\skt{~॥}\pend

\pstart \ledsidenote{\passref{4.35.11}}\edtext{\skt{व्रणनेत्राङ्गुलं}}{\lemma{\skt{व्रणनेत्राङ्गुलं}}\Afootnote{H; \skt{(व्र\appbrk{}ण\appbrk{}ने\appbrk{}त्र\appbrk{}म\appbrk{}ष्टा\appbrk{}ङ्गुलं}\ A; \skt{व्रणे नेत्रांगुलं}\ K}}\skt{ }\edtext{\skt{सुङ्गतुल्यस्रोतो वा हि}}{\lemma{\skt{सुङ्गतुल्यस्रोतो वा हि}}\Afootnote{K H; \skt{मुद्गवाहिस्रोतः;}\ A}}\skt{ }\edtext{\edtext{\skt{व्रणमवेक्ष्य}}{\lemma{\skt{व्रणम्}}\Afootnote{K H; \skt{व्रण}\ A}}\skt{ च}}{\lemma{\skt{अवेक्ष्य च}}\Afootnote{K H; \skt{मवेक्ष्य य}\ A}}\skt{ }\edtext{\skt{स्नेहकषायो}}{\lemma{\skt{स्नेहकषायो}}\Afootnote{H; \skt{थास्वं स्नेहकषाये}\ A; \skt{स्नेहकषायौ}\ K}}\skt{ विदधीत }\edtext{\skt{इति}}{\lemma{\skt{इति}}\Afootnote{K H; \emph{om.}\ A}}\skt{~।}\pend

\pstart \ledsidenote{\passref{4.35.12}}\skt{तत्र नेत्राणि }\edtext{\skt{सुव\-र्ण्ण\-र\-ज\-त\-ता\-म्रा\-य\-स\-द\-न्त\-शृं\-ग\-म\-णि\-सा\-र\-म\-याणि}}{\lemma{\skt{सुव\appbrk{}र्ण्ण\appbrk{}र\appbrk{}ज\appbrk{}त\appbrk{}ता\appbrk{}म्रा\appbrk{}य\appbrk{}स\appbrk{}द\appbrk{}न्त\appbrk{}शृं\appbrk{}ग\appbrk{}म\appbrk{}णि\appbrk{}सा\appbrk{}र\appbrk{}म\appbrk{}याणि}}\Afootnote{H; \skt{सुव\appbrk{}र्ण\appbrk{}र\appbrk{}ज\appbrk{}त\appbrk{}ता\appbrk{}म्रा\appbrk{}यो\appbrk{}री\appbrk{}ति\appbrk{}द\appbrk{}न्त\appbrk{}शृ\appbrk{}ङ्ग\appbrk{}म\appbrk{}णि\appbrk{}त\appbrk{}रु\appbrk{}सा\appbrk{}र\appbrk{}म\appbrk{}यानि}\ A; \skt{सुव\appbrk{}र्ण्ण\appbrk{}र\appbrk{}ज\appbrk{}त\appbrk{}ता\appbrk{}म्रा\appbrk{}य\appbrk{}स\appbrk{}द\appbrk{}न्त\appbrk{}शृ\appbrk{}ङ्ग\appbrk{}म\appbrk{}णि\appbrk{}त\appbrk{}रु\appbrk{}सा\appbrk{}र\appbrk{}म\appbrk{}याणि}\ K}}\skt{ श्लक्ष्णानि दृढानि }\edtext{\skt{गोपुच्छाकृतीनि}}{\lemma{\skt{गोपुच्छाकृतीनि}}\Afootnote{H; \skt{गोपु\appbrk{}च्छा\appbrk{}कृ\appbrk{}ती\appbrk{}न्यृ\appbrk{}जूनि}\ A; \skt{गुपुच्छाकृतिनि~।}\ K}}\skt{~। गुटिकामुखानि }\edtext{\skt{चेति}}{\lemma{\skt{चेति}}\Afootnote{K H; \skt{च~॥}\ A}}\skt{~।}\pend

\pstart \ledsidenote{\passref{4.35.13}}\edtext{\skt{बस्त\-य\-श्चा\-त्रा\-व\-बन्ध्या}}{\lemma{\skt{चात्रावबन्ध्या}}\Afootnote{K H; \skt{च बन्ध्या}\ A}}\skt{ मृदवो }\edtext{\skt{नातिबहला}}{\lemma{\skt{नातिबहला}}\Afootnote{K H; \skt{नातिवहला}\ A}}\skt{ दृढाः प्रमाणवन्तो }\edtext{\edtext{\skt{महि\-ष\-गो\-व\-रा\-हा\-डो\-र\-भ्रा\-णा\-म\-न्य\-त\-म\-मा\-सा\-द्येति}}{\lemma{\skt{महि\appbrk{}ष\appbrk{}गो\appbrk{}व\appbrk{}रा\appbrk{}हा\appbrk{}डो\appbrk{}र\appbrk{}भ्रा\appbrk{}णा\appbrk{}म\appbrk{}न्य\appbrk{}त\appbrk{}मम्}}\Afootnote{H; \skt{गोम\appbrk{}हि\appbrk{}ष\appbrk{}व\appbrk{}रा\appbrk{}हा\appbrk{}जो\appbrk{}र\appbrk{}भ्रा\appbrk{}णाम्~॥}\ A; \skt{महि\appbrk{}ष\appbrk{}गो\appbrk{}रा\appbrk{}हा\appbrk{}डो\appbrk{}र\appbrk{}भ्रा\appbrk{}णा\appbrk{}म\appbrk{}न्य\appbrk{}त\appbrk{}मम्}\ K}}}{\lemma{\skt{आसाद्येति}}\Afootnote{K H; \emph{om.}\ A}}\skt{~॥}\pend

\pstart\versepar \ledsidenote{\passref{4.35.14}}\skt{नेत्रालाभे हिता नाडी }\edtext{\skt{नल\-वे\-ल्व\-स्थि\-स\-म्भवा}}{\lemma{\skt{नल\appbrk{}वे\appbrk{}ल्व\appbrk{}स्थि\appbrk{}स\appbrk{}म्भवा}}\Afootnote{K H; \skt{नल\appbrk{}वं\appbrk{}शा\appbrk{}स्थि\appbrk{}सं\appbrk{}भवा~।}\ A}}\skt{~। }\pend
\pstart\versepar \skt{}\edtext{\skt{बस्त्यालाभे}}{\lemma{\skt{बस्त्यालाभे}}\Afootnote{K H; \skt{बस्त्यलाभे}\ A}}\skt{ }\edtext{\skt{हितश्चर्म्म}}{\lemma{\skt{हितश्चर्म्म}}\Afootnote{K H; \skt{हितं चर्म}\ A}}\skt{ }\edtext{\skt{शूक्ष्मम्वा}}{\lemma{\skt{शूक्ष्मम्वा}}\Afootnote{H; \skt{सूक्ष्मं वा}\ A K}}\skt{ तान्तवं घनं~॥}\pend

\pstart \ledsidenote{\passref{4.35.18}}\skt{तत्र द्विविधो }\edtext{\skt{बस्तिर्न्नैरूहिकः}}{\lemma{\skt{न्नैरूहिकः}}\Afootnote{A H; \skt{नैरूहिक}\ K}}\skt{ }\edtext{\skt{स्नेहिकश्च}}{\lemma{\skt{स्नेहिकश्च}}\Afootnote{H; \skt{स्नैहिकश्च~।}\ A K}}\skt{~॥ }\edtext{\skt{आस्थापनो}}{\lemma{\skt{आस्थापनो}}\Afootnote{K H; \skt{आस्थापनं,}\ A}}\skt{ निरूह }\edtext{\skt{इत्यनर्थान्तरं}}{\lemma{\skt{इत्यनर्थान्तरं}}\Afootnote{H; \skt{इत्यनर्थान्तरं;}\ A; \skt{इत्यर्थान्तरं~।}\ K}}\skt{~। तस्य विकल्पो माधुतैलिकः~। तस्य }\edtext{\skt{पर्यायशब्दो}}{\lemma{\skt{पर्यायशब्दो}}\Afootnote{K H; \skt{प्रयायशब्दो}\ A}}\skt{ }\edtext{\skt{युक्तरथः}}{\lemma{\skt{युक्तरथः}}\Afootnote{K H; \skt{यापनो, युक्तरथः,}\ A}}\skt{ }\edtext{\skt{सिद्धबस्तिरिति}}{\lemma{\skt{सिद्धबस्तिरिति}}\Afootnote{H; \skt{सिद्धबस्तिरिति~।}\ A; \skt{सिद्धस्तिरिति~।}\ K}}\skt{~। }\edtext{\skt{सर्व्व\-दो\-ष\-नि\-र्ह\-र\-णा\-च्छ\-री\-र\-रो\-ह\-णाद्वा}}{\lemma{\skt{सर्व्व\appbrk{}दो\appbrk{}ष\appbrk{}नि\appbrk{}र्ह\appbrk{}र\appbrk{}णा\appbrk{}च्छ\appbrk{}री\appbrk{}र\appbrk{}रो\appbrk{}ह\appbrk{}णाद्वा}}\Afootnote{H; \skt{स दोष\appbrk{}नि\appbrk{}र्ह\appbrk{}र\appbrk{}णा\appbrk{}च्छ\appbrk{}री\appbrk{}र\appbrk{}नी\appbrk{}रो\appbrk{}ह\appbrk{}णाद्वा}\ A; \skt{सर्व\appbrk{}दो\appbrk{}ष\appbrk{}नि\appbrk{}र्ह\appbrk{}र\appbrk{}णा\appbrk{}च्छ\appbrk{}री\appbrk{}रो\appbrk{}ह\appbrk{}णाद्वा}\ K}}\skt{ निरूहः }\edtext{\skt{स्थापनादायुः स्थानाद्वा}}{\lemma{\skt{स्थापनादायुः स्थानाद्वा}}\Afootnote{H; \skt{बयः\appbrk{}स्था\appbrk{}प\appbrk{}ना\appbrk{}दा\appbrk{}युः\appbrk{}स्था\appbrk{}प\appbrk{}नाद्वा}\ A; \skt{वय\appbrk{}स्था\appbrk{}प\appbrk{}ना\appbrk{}दा\appbrk{}यु\appbrk{}स्था\appbrk{}प\appbrk{}नाद्वा आस्थानाद्वा}\ K}}\skt{ }\edtext{\skt{आस्थापन}}{\lemma{\skt{आस्थापन}}\Afootnote{K H; \skt{आस्थापनम्~।}\ A}}\skt{ }\edtext{\skt{इति~। }\edtext{\skt{पश्चा\-त्मा\-धु\-तै\-लि\-क\-वि\-क\-ल्पा\-नु\-प\-दे\-क्ष्यामः}}{\lemma{\skt{चात्मा\appbrk{}धु\appbrk{}तै\appbrk{}लि\appbrk{}क\appbrk{}वि\appbrk{}क\appbrk{}ल्पा\appbrk{}नु\appbrk{}प\appbrk{}दे\appbrk{}क्ष्यामः}}\Afootnote{H; \skt{माधु\appbrk{}तै\appbrk{}लि\appbrk{}क\appbrk{}वि\appbrk{}धानं च निरू\appbrk{}हो\appbrk{}प\appbrk{}क्र\appbrk{}म\appbrk{}चि\appbrk{}कि\appbrk{}त्सिते वक्ष्यामः~।}\ A; \skt{श्चात्मा\appbrk{}धु\appbrk{}तै\appbrk{}लि\appbrk{}क\appbrk{}वि\appbrk{}क\appbrk{}ल्पा\appbrk{}नु\appbrk{}प\appbrk{}दे\appbrk{}क्ष्यामः~।}\ K}}}{\lemma{\skt{इति~। पश्}}\Afootnote{H; \emph{om.}\ A; \skt{इति~। प}\ K}}\skt{~॥ }\edtext{\skt{तत्र}}{\lemma{\skt{तत्र}}\Afootnote{K H; \emph{om.}\ A}}\skt{ यथा प्रमा\-ण\-गु\-ण\-वि\-हितः }\edtext{\skt{स्नेहबस्तिः}}{\lemma{\skt{स्नेहबस्तिः}}\Afootnote{K H; \skt{स्नेह}\ A}}\skt{ पादापकृष्टः~। }\edtext{\skt{स्नेहबस्ति}}{\lemma{\skt{स्नेहबस्ति}}\Afootnote{K\textsuperscript{pc} H; \skt{स्नेबस्ति}\ K\textsuperscript{ac}}}\skt{ }\edtext{\edtext{\skt{विक\-ल्पो\-र्द्ध\-मा\-त्रा\-प\-कृ\-ष्टो\-नु\-वा\-सनः।}}{\lemma{\skt{विक\appbrk{}ल्पो\appbrk{}र्द्ध\appbrk{}मा\appbrk{}त्रा\appbrk{}प\appbrk{}कृ\appbrk{}ष्टो\appbrk{}नु\appbrk{}वा\appbrk{}सनः}}\Afootnote{K H; \skt{विक\appbrk{}ल्पोऽ\appbrk{}नु\appbrk{}वा\appbrk{}सनः}\ A}}\skt{ अनवसन्नपि}}{\lemma{\skt{बस्तिः… स्नेहब}}\Afootnote{\emph{om.}\ A}}\skt{ न }\edtext{\skt{दुष्यत}}{\lemma{\skt{दुष्यत}}\Afootnote{K H; \skt{दुष्य}\ A}}\skt{ }\edtext{\skt{इत्य\-नु\-दि\-व\-सम्वा}}{\lemma{\skt{इत्य\appbrk{}नु\appbrk{}दि\appbrk{}व\appbrk{}सम्वा}}\Afootnote{K H; \skt{त्यनुदिवसं वा}\ A}}\skt{ दीयत }\edtext{\skt{इत्यनुवासनः}}{\lemma{\skt{इत्यनुवासनः}}\Afootnote{K H; \skt{इत्यनुवासनः~।}\ A}}\skt{~। }\edtext{\skt{तस्यातिविकल्पोऽ}}{\lemma{\skt{तस्यातिविकल्पोऽ}}\Afootnote{H; \skt{तस्यापि विकल्पोऽ}\ A; \skt{तस्यापि विकल्पो}\ K}}\skt{ }\edtext{\edtext{\skt{र्द्धर्थ\-मा\-त्रा\-प\-कृ\-ष्टोऽ\-प\-हार्यो}}{\lemma{\skt{र्द्धर्थ\appbrk{}मा\appbrk{}त्रा\appbrk{}प\appbrk{}कृष्टो}}\Afootnote{H\textsuperscript{pc}; \skt{र्धार्ध\appbrk{}मा\appbrk{}त्रा\appbrk{}व\appbrk{}कृष्टो}\ A; \skt{र्धमात्रापकृष्टो}\ K; \skt{र्थमात्रापकृष्टो}\ H\textsuperscript{ac}}}}{\lemma{\skt{ऽपहार्यो}}\Afootnote{H; \skt{ऽपरिहार्यो}\ A; \skt{पहार्यो}\ K}}\skt{ }\edtext{\skt{मात्राबस्तिरिति}}{\lemma{\skt{मात्राबस्तिरिति}}\Afootnote{K H; \skt{मात्राबस्तिरिति~॥}\ A}}\skt{~॥}\edtext{}{\lemma{\skt{अनवसन्नपि न}}\Afootnote{H; \skt{पादा\appbrk{}व(\appbrk{}पा\appbrk{}दा\appbrk{}प)\appbrk{}कृष्टः~। अनुवसन्नपि न}\ A; \skt{अनवसन्नपि न}\ K}}\pend

\pstart\versepar \ledsidenote{\passref{4.35.19}}\skt{निरूहः शोधनो }\edtext{\skt{लेखः}}{\lemma{\skt{लेखः}}\Afootnote{K H; \skt{लेखी}\ A}}\skt{ स्नैहिको }\edtext{\skt{वृष्यबृंहणः}}{\lemma{\skt{वृष्यबृंहणः}}\Afootnote{K H; \skt{बृंहणो}\ A}}\skt{~। }\pend
\pstart\versepar \skt{}\edtext{\edtext{\skt{निरू\-ह\-शो\-धि\-तै\-र्म्मार्ग्गैः}}{\lemma{\skt{निरू\appbrk{}ह\appbrk{}शो\appbrk{}धि\appbrk{}तैर्}}\Afootnote{H; \skt{मतः~। निरू\appbrk{}ह\appbrk{}शो\appbrk{}धि\appbrk{}तान्}\ A; \skt{निरूहः शोधितैर्}\ K}}}{\lemma{\skt{म्मार्ग्गैः}}\Afootnote{K H; \skt{मार्गा}\ A}}\skt{ }\edtext{\skt{सम्यक्स्नेहो}}{\lemma{\skt{सम्यक्स्नेहो}}\Afootnote{K H; \skt{न्सम्यक्स्नेहो}\ A}}\skt{ }\edtext{\skt{विसर्प्पति}}{\lemma{\skt{विसर्प्पति}}\Afootnote{K H; \skt{ऽनुगच्छति~। अपे\appbrk{}त\appbrk{}स\appbrk{}र्व\appbrk{}दो\appbrk{}षासु नाडीष्विव वहज्जलम्~॥}\ A}}\skt{~॥}\pend

\pstart \ledsidenote{\passref{4.35.21}}\edtext{\skt{तत्रो\-न्मा\-द\-भ\-य\-शो\-क\-पि\-पा\-सा\-रो\-च\-का\-जी\-र्ण्ण\-पा\-ण्डु\-रो\-ग\-भ्र\-म\-म\-द\-मू\-र्च्छा\-च्छ\-र्द्दि\-कु\-ष्ठ\-मे\-हो\-द\-र\-स्थौ\-ल्य\-का\-स\-श्वा\-स\-शो\-ष\-शो\-फो\-प\-सृ\-ष्ट\-क्ष\-त\-क्षी\-ण\-न्यू\-न\-त्रि\-मा\-स\-ग\-र्ब्भिणी}}{\lemma{\skt{तत्रो\appbrk{}न्मा\appbrk{}द\appbrk{}भ\appbrk{}य\appbrk{}शो\appbrk{}क\appbrk{}पि\appbrk{}पा\appbrk{}सा\appbrk{}रो\appbrk{}च\appbrk{}का\appbrk{}जी\appbrk{}र्ण्ण\appbrk{}पा\appbrk{}ण्डु\appbrk{}रो\appbrk{}ग\appbrk{}भ्र\appbrk{}म\appbrk{}म\appbrk{}द\appbrk{}मू\appbrk{}र्च्छा\appbrk{}च्छ\appbrk{}र्द्दि\appbrk{}कु\appbrk{}ष्ठ\appbrk{}मे\appbrk{}हो\appbrk{}द\appbrk{}र\appbrk{}स्थौ\appbrk{}ल्य\appbrk{}का\appbrk{}स\appbrk{}श्वा\appbrk{}स\appbrk{}शो\appbrk{}ष\appbrk{}शो\appbrk{}फो\appbrk{}प\appbrk{}सृ\appbrk{}ष्ट\appbrk{}क्ष\appbrk{}त\appbrk{}क्षी\appbrk{}ण\appbrk{}न्यू\appbrk{}न\appbrk{}त्रि\appbrk{}मा\appbrk{}स\appbrk{}ग\appbrk{}र्ब्भिणी}}\Afootnote{K H\textsuperscript{pc}; \skt{तत्रो\appbrk{}न्मा\appbrk{}द\appbrk{}भ\appbrk{}य\appbrk{}शो\appbrk{}क\appbrk{}पि\appbrk{}पा\appbrk{}सा\appbrk{}रो\appbrk{}च\appbrk{}का\appbrk{}जी\appbrk{}र्णा\appbrk{}र्शः\appbrk{}पा\appbrk{}ण्डु\appbrk{}रो\appbrk{}ग\appbrk{}भ्र\appbrk{}म\appbrk{}म\appbrk{}द\appbrk{}मू\appbrk{}र्च्छा\appbrk{}च्छ\appbrk{}र्दि\appbrk{}कु\appbrk{}ष्ठ\appbrk{}मे\appbrk{}हो\appbrk{}द\appbrk{}र\appbrk{}स्थौ\appbrk{}ल्य\appbrk{}श्वा\appbrk{}स\appbrk{}का\appbrk{}स\appbrk{}क\appbrk{}ण्ठ\appbrk{}शो\appbrk{}ष\appbrk{}शो\appbrk{}फो\appbrk{}प\appbrk{}सृ\appbrk{}ष्ट\appbrk{}क्ष\appbrk{}त\appbrk{}क्षी\appbrk{}ण\appbrk{}च\appbrk{}तु\appbrk{}स्त्रि\appbrk{}मा\appbrk{}स\appbrk{}ग\appbrk{}र्भिणी}\ A; \skt{तत्रो\appbrk{}न्मा\appbrk{}द\appbrk{}भ\appbrk{}य\appbrk{}शो\appbrk{}क\appbrk{}पि\appbrk{}पा\appbrk{}सा\appbrk{}रो\appbrk{}च\appbrk{}का\appbrk{}जी\appbrk{}र्ण्ण\appbrk{}पा\appbrk{}ण्डु\appbrk{}रो\appbrk{}ग\appbrk{}भ्र\appbrk{}म\appbrk{}द\appbrk{}मू\appbrk{}र्च्छा\appbrk{}च्छ\appbrk{}र्द्दि\appbrk{}कु\appbrk{}ष्ठ\appbrk{}मे\appbrk{}हो\appbrk{}द\appbrk{}र\appbrk{}स्थौ\appbrk{}ल्य\appbrk{}का\appbrk{}स\appbrk{}श्वा\appbrk{}स\appbrk{}शो\appbrk{}ष\appbrk{}शो\appbrk{}फो\appbrk{}प\appbrk{}सृ\appbrk{}ष्ट\appbrk{}क्ष\appbrk{}त\appbrk{}क्षी\appbrk{}ण\appbrk{}न्यू\appbrk{}न\appbrk{}त्रि\appbrk{}मा\appbrk{}स\appbrk{}ग\appbrk{}र्ब्भिणी}\ H\textsuperscript{ac}}}\skt{ }\edtext{\skt{दुर्व्व\-ला\-ग्न्य\-स\-ह\-वा\-त\-रो\-ग\-दृते च}}{\lemma{\skt{दुर्व्व\appbrk{}ला\appbrk{}ग्न्य\appbrk{}स\appbrk{}ह\appbrk{}वा\appbrk{}त\appbrk{}रो\appbrk{}ग\appbrk{}दृते च}}\Afootnote{H; \skt{दुर्बलाग्न्यसहा बालवृद्धौ च वातरोगादृते}\ A; \skt{दुर्व\appbrk{}ला\appbrk{}ग\appbrk{}न्य\appbrk{}स\appbrk{}ह\appbrk{}वा\appbrk{}त\appbrk{}रो\appbrk{}गा\appbrk{}दृते च}\ K}}\skt{ क्षीणा }\edtext{\skt{नानु\-वा\-स्या\-ना\-स्था\-प्याश्च}}{\lemma{\skt{नानु\appbrk{}वा\appbrk{}स्या\appbrk{}ना\appbrk{}स्था\appbrk{}प्याश्च}}\Afootnote{K H\textsuperscript{pc}; \skt{नानुवास्या नास्थापयि}\ A; \skt{नुवा\appbrk{}स्या\appbrk{}ना\appbrk{}स्था\appbrk{}प्याश्च}\ H\textsuperscript{ac}}}\skt{ }\edtext{\skt{भवन्ति}}{\lemma{\skt{भवन्ति}}\Afootnote{H; \skt{तव्याः~॥}\ A; \skt{भवन्ती}\ K}}\skt{ }\edtext{\skt{इति}}{\lemma{\skt{इति}}\Afootnote{H; \emph{om.}\ A; \skt{ति~॥}\ K}}\skt{~॥}\pend

\pstart\versepar \ledsidenote{\passref{4.35.22}}\skt{उदरी च प्रमेही च कुष्ठी स्थूलश्च मानवः~। }\pend
\pstart\versepar \skt{}\edtext{\skt{अव\-स्था\-स्था\-प\-नी\-याश्च स्तेनानुवास्याः}}{\lemma{\skt{अव\appbrk{}स्था\appbrk{}स्था\appbrk{}प\appbrk{}नी\appbrk{}याश्च स्तेनानुवास्याः}}\Afootnote{\skt{अवश्यं स्थापनीयास्ते नानुवास्याः}\ A; \skt{अव\appbrk{}स्था\appbrk{}स्था\appbrk{}प\appbrk{}नी\appbrk{}यास्ते नानुवास्याः}\ K H\textsuperscript{pc}; \skt{अव\appbrk{}स्था\appbrk{}स्था\appbrk{}प\appbrk{}नी\appbrk{}या\appbrk{}श्चस्ते नानुवास्याः}\ H\textsuperscript{ac}}}\skt{ }\edtext{\skt{कथञ्चन}}{\lemma{\skt{कथञ्चन}}\Afootnote{A H; \skt{कथञ्चनः~।}\ K}}\skt{~॥}\pend

\pstart\versepar \ledsidenote{\passref{4.35.23}}\edtext{\skt{अनु\-वा\-स\-ना\-द्भ\-व\-न्त्ये\-षा\-म्वि\-का\-रा\-णा\-म\-सा\-ध्यता}}{\lemma{\skt{अनु\appbrk{}वा\appbrk{}स\appbrk{}ना\appbrk{}द्भ\appbrk{}वन्त्य्}}\Afootnote{K H; \skt{असाध्यता विकाराणां स्यादे\appbrk{}षा\appbrk{}म\appbrk{}नु\appbrk{}वा\appbrk{}सना}\ A}}\skt{~। }\pend
\pstart\versepar \skt{}\edtext{\skt{असाध्यत्वे}}{\lemma{\skt{असाध्यत्वे}}\Afootnote{K H; \skt{साध्यत्वे}\ A}}\skt{ }\edtext{\skt{पि}}{\lemma{\skt{पि}}\Afootnote{K H; \skt{ऽपि}\ A}}\skt{ भूयिष्ठं }\edtext{\skt{गात्रा\-णा\-म\-व\-सा\-दनं}}{\lemma{\skt{अवसादनं}}\Afootnote{K H; \skt{सदनं भवेत्~॥}\ A}}\skt{~॥}\edtext{}{\lemma{\skt{द्… असाध्यता}}\Afootnote{\emph{om.}\ A}}\pend

\pstart\versepar \ledsidenote{\passref{4.35.25}}\edtext{\skt{पक्वा\-स\-या\-द्ब\-स्ति\-वीर्यं}}{\lemma{\skt{पक्वा\appbrk{}स\appbrk{}या\appbrk{}द्ब\appbrk{}स्ति\appbrk{}वीर्यं}}\Afootnote{H; \skt{पक्वा\appbrk{}श\appbrk{}या\appbrk{}द्ब\appbrk{}स्ति\appbrk{}वीर्यं}\ A; \skt{पक्वा\appbrk{}श\appbrk{}या\appbrk{}द्ब\appbrk{}स्ति\appbrk{}वीर्यं}\ K}}\skt{ }\edtext{\skt{खैर्द्दे\-ह\-म\-नु\-स\-र्प्पति}}{\lemma{\skt{खैर्द्दे\appbrk{}ह\appbrk{}म\appbrk{}नु\appbrk{}स\appbrk{}र्प्पति}}\Afootnote{H; \skt{खैर्दे\appbrk{}ह\appbrk{}म\appbrk{}नु\appbrk{}स\appbrk{}र्पति~।}\ A; \skt{खैर्दे\appbrk{}ह\appbrk{}म\appbrk{}नु\appbrk{}स\appbrk{}र्पति~।}\ K}}\skt{~॥ }\pend
\pstart\versepar \skt{वृक्षमूले }\edtext{\edtext{\edtext{\edtext{\skt{निषि\-क्ता\-ना\-म\-पा\-म्वी\-र्य\-मिव}}{\lemma{\skt{निषिक्तानाम्}}\Afootnote{K H; \skt{निषिक्ताना}\ A}}}{\lemma{\skt{अपाम्}}\Afootnote{K H; \skt{मपां}\ A}}}{\lemma{\skt{वीर्यम्}}\Afootnote{K H; \skt{वीर्य}\ A}}}{\lemma{\skt{इव}}\Afootnote{K H; \skt{मिव}\ A}}\skt{ }\edtext{\skt{द्रुमः}}{\lemma{\skt{द्रुमः}}\Afootnote{K H; \skt{द्रुमम्~॥}\ A}}\skt{~॥}\pend

\pstart\versepar \ledsidenote{\passref{4.35.26}}\edtext{\skt{सर्व्वापि}}{\lemma{\skt{सर्व्वापि}}\Afootnote{K H; \skt{स चापि}\ A}}\skt{ }\edtext{\skt{बस्तिः}}{\lemma{\skt{बस्तिः}}\Afootnote{K H; \skt{सह}\ A}}\skt{ }\edtext{\skt{सहसा}}{\lemma{\skt{सहसा}}\Afootnote{K H; \skt{सा बस्तिः}\ A}}\skt{ केवलः समलो पि वा~॥ }\pend
\pstart\versepar \skt{प्रत्येति वीर्य\-न्त्व\-नि\-लै\-र\-पा\-नाद्यैः }\edtext{\skt{प्रणीयते}}{\lemma{\skt{प्रणीयते}}\Afootnote{K H; \skt{विनीयते~॥}\ A}}\skt{~॥}\pend

\pstart\versepar \ledsidenote{\passref{4.35.27}}\skt{वीर्येण बस्तिरादत्ते दोषा\-ना\-पा\-द\-म\-स्त\-कान्~। }\pend
\pstart\versepar \skt{}\edtext{\skt{पक्वा\-स\-य\-स्थः\-ख\-स्थोऽर्क्क}}{\lemma{\skt{ऽर्क्क}}\Afootnote{H; \skt{निव~॥}\ A; \emph{om.}\ K}}\skt{ }\edtext{\skt{अपो}}{\lemma{\skt{अपो}}\Afootnote{K H; \emph{om.}\ A}}\skt{ }\edtext{\skt{यद्व\-त्म\-ही\-त\-लात्}}{\lemma{\skt{यद्व\appbrk{}त्म\appbrk{}ही\appbrk{}त\appbrk{}लात्}}\Afootnote{K H; \emph{om.}\ A}}\skt{~॥}\edtext{}{\lemma{\skt{दोषा\appbrk{}ना\appbrk{}पा\appbrk{}द\appbrk{}म\appbrk{}स्त\appbrk{}कान्~। पक्वा\appbrk{}स\appbrk{}य\appbrk{}स्थः\appbrk{}खस्थो}}\Afootnote{H; \skt{दोषा\appbrk{}ना\appbrk{}पा\appbrk{}द\appbrk{}म\appbrk{}स्त\appbrk{}का\appbrk{}त्(न्)~। पक्वा\appbrk{}श\appbrk{}य\appbrk{}स्थोऽ\appbrk{}म्ब\appbrk{}रगो भूमेरर्को रसा}\ A; \skt{दोषा\appbrk{}ना\appbrk{}पा\appbrk{}द\appbrk{}म\appbrk{}स्त\appbrk{}कात्~। पक्वाशयस्थः खस्थोर्क}\ K}}\pend

\pstart\versepar \ledsidenote{\passref{4.35.28}}\edtext{\skt{स कटी\-पृ\-ष्ठ\-को\-ष्ठ\-स्था\-न्वी\-र्ये\-णा\-लोड्य}}{\lemma{\skt{स कटी\appbrk{}पृ\appbrk{}ष्ठ\appbrk{}को\appbrk{}ष्ठ\appbrk{}स्था\appbrk{}न्वी\appbrk{}र्ये\appbrk{}णा\appbrk{}लोड्य}}\Afootnote{A H\textsuperscript{pc}; \skt{स कटी\appbrk{}पृ\appbrk{}ष्ठ\appbrk{}को\appbrk{}ष्ठस्थं वीर्येणालोड्य}\ K; \skt{स कटी\appbrk{}पृ\appbrk{}ष्ठ\appbrk{}को\appbrk{}ष्ठ\appbrk{}स्था\appbrk{}न्वी\appbrk{}र्ये\appbrk{}णाड्य}\ H\textsuperscript{ac}}}\skt{ सञ्चयान्~। }\pend
\pstart\versepar \skt{उत्खा\-त\-मू\-ला\-न्ह\-रति }\edtext{\skt{दोषान्वै}}{\lemma{\skt{दोषान्वै}}\Afootnote{K H; \skt{दोषाणां}\ A}}\skt{ साधुयोजितः~॥}\edtext{}{\lemma{\skt{सञ्चयान्~। उत्खा\appbrk{}त\appbrk{}मू\appbrk{}ला\appbrk{}न्ह\appbrk{}रति}}\Afootnote{A H; \skt{सञ्चयान्~। उत्खातमूलां हरति}\ K}}\pend

\pstart\versepar \ledsidenote{\passref{4.35.30}}\skt{वायोर्व्विषहते }\edtext{\skt{वेगान्नान्या}}{\lemma{\skt{वेगान्}}\Afootnote{H; \skt{वेगं}\ A K}}\skt{ }\edtext{\skt{बस्तिमृते}}{\lemma{\skt{बस्तिमृते}}\Afootnote{K H; \skt{बस्तेरृते}\ A}}\skt{ क्रिया~। }\pend
\pstart\versepar \skt{}\edtext{\skt{पव\-ना\-वि\-द्ध\-तो\-यस्य}}{\lemma{\skt{पव\appbrk{}ना\appbrk{}वि\appbrk{}द्ध\appbrk{}तो\appbrk{}यस्य}}\Afootnote{A H; \skt{पव\appbrk{}ना\appbrk{}द्वि\appbrk{}द्ध\appbrk{}तो\appbrk{}यस्य}\ K}}\skt{ }\edtext{\edtext{\skt{वेलामिव}}{\lemma{\skt{वेलाम्}}\Afootnote{K H; \skt{वेला}\ A}}}{\lemma{\skt{इव}}\Afootnote{K H; \skt{वेग}\ A}}\skt{ }\edtext{\skt{महोदधेः}}{\lemma{\skt{महोदधेः}}\Afootnote{K H; \skt{मिवोदधेः~॥}\ A}}\skt{~॥}\pend

\pstart \ledsidenote{\passref{4.35.trailer}}\edtext{\skt{॥ चिकित्सा}}{\lemma{\skt{॥ चिकित्सा}}\Afootnote{H; \skt{इति सुश्रु\appbrk{}त\appbrk{}सं\appbrk{}हि\appbrk{}तायां चिकित्सा}\ A; \skt{॥ चिलहृ~॥ {\symfont ❈}~॥}\ K}}\skt{ }\edtext{\skt{लहृ}}{\lemma{\skt{लहृ}}\Afootnote{H; \skt{स्थाने नेत्र\appbrk{}ब\appbrk{}स्ति\appbrk{}प्र\appbrk{}मा\appbrk{}ण\appbrk{}प्र\appbrk{}वि\appbrk{}भा\appbrk{}ग\appbrk{}चि\appbrk{}कि\appbrk{}त्स्तितं नाम पञ्चत्रिंशोऽध्यायः~॥३५~॥}\ A; \emph{om.}\ K}}\pend

\endnumbering
\beginnumbering
\eledsection*{Adhyāya 36}
\pstart \ledsidenote{\passref{4.36.1}}\edtext{\skt{अथातो}}{\lemma{\skt{अथातो}}\Afootnote{A H; \skt{अथातो}\ K}}\skt{ }\edtext{\edtext{\skt{नेत्र\-ब\-स्ति\-व्या\-प\-च्चि\-कि\-त्सि\-त\-म्व्या\-ख्या\-स्यामः}}{\lemma{\skt{नेत्र\appbrk{}ब\appbrk{}स्ति\appbrk{}व्या\appbrk{}प\appbrk{}च्चि\appbrk{}कि\appbrk{}त्सि\appbrk{}तम्}}\Afootnote{A H; \emph{om.}\ K}}}{\lemma{\skt{व्याख्यास्यामः}}\Afootnote{A H; \emph{om.}\ K}}\skt{~॥}\pend

\pstart \ledsidenote{\passref{4.36.2.add}}\edtext{\edtext{\skt{तत्र नेत्र\-ञ्च\-लि\-त\-म्वि\-व\-र्त्ति\-त\-म्पा\-र्श्वा\-व\-पी\-डि\-त\-म\-त्यु\-त्त्क्षि\-प्त\-म\-व\-स\-न्न\-न्ति\-र्य\-क्क्षि\-प्त\-मिति}}{\lemma{\skt{… अत्युत्त्क्षिप्तम्}}\Afootnote{\emph{om.}\ K\textsuperscript{ac}}}}{\lemma{\skt{… अत्युत्त्क्षिप्तम्}}\Afootnote{\emph{om.}\ K\textsuperscript{pc}}}\skt{ }\edtext{\skt{षघ्र\-णि\-धा\-न\-दोषाः}}{\lemma{\skt{षघ्र\appbrk{}णि\appbrk{}धा\appbrk{}न\appbrk{}दोषाः}}\Afootnote{H; \skt{षट्घ्र\appbrk{}णि\appbrk{}धा\appbrk{}ण\appbrk{}दोषाः~॥}\ K}}\skt{~। }\edtext{\edtext{\skt{अति\-स्थू\-ल\-ङ्क\-र्क्क\-शा\-ग्र\-म\-व\-न\-त\-म\-नु\-भि\-न्न\-स\-न्नि\-कृ\-ष्ट\-क\-र्ण्णिकं}}{\lemma{\skt{अवनतम्}}\Afootnote{H; \emph{om.}\ K}}}{\lemma{\skt{अनु\appbrk{}भि\appbrk{}न्न\appbrk{}स\appbrk{}न्नि\appbrk{}कृ\appbrk{}ष्ट\appbrk{}क\appbrk{}र्ण्णिकं}}\Afootnote{H; \skt{अणु\appbrk{}भि\appbrk{}न्न\appbrk{}स\appbrk{}न्नि\appbrk{}कृ\appbrk{}ष्ट\appbrk{}वि\appbrk{}प्र\appbrk{}कृ\appbrk{}ष्ट\appbrk{}क\appbrk{}र्ण्णिकं}\ K}}\skt{ }\edtext{\edtext{\skt{शूक्ष्मा\-ति\-म\-ज\-च्छि\-द्र\-म\-ल्प\-म\-ति\-दी\-र्घ\-म\-ति\-ह्र\-स्व\-मि\-त्ये\-का\-द\-श\-ने\-त्र\-दोषाः}}{\lemma{\skt{शूक्ष्मा\appbrk{}ति\appbrk{}म\appbrk{}ज\appbrk{}च्छि\appbrk{}द्रम्}}\Afootnote{H; \skt{सूक्ष्मातिच्छिद्र}\ K}}}{\lemma{\skt{अल्पम्}}\Afootnote{H; \skt{मल्पम्}\ K}}\skt{~॥ }\edtext{\skt{बहलताल्पता}}{\lemma{\skt{बहलताल्पता}}\Afootnote{H; \skt{बहलता}\ K\textsuperscript{ac}; \skt{बहलताल्प}\ K\textsuperscript{pc}}}\skt{ }\edtext{\skt{च्छिद्रता}}{\lemma{\skt{च्छिद्रता}}\Afootnote{H; \skt{च्छिद्रता}\ K}}\skt{ }\edtext{\skt{प्रस्त्धानता}}{\lemma{\skt{प्रस्त्धानता}}\Afootnote{H; \skt{प्रस्त्यानता}\ K}}\skt{ }\edtext{\skt{दुर्ब्ब\-ल\-दु\-र्ब्ब\-द्धेति त्वेति}}{\lemma{\skt{दुर्ब्ब\appbrk{}ल\appbrk{}दु\appbrk{}र्ब्ब\appbrk{}द्धेति त्वेति}}\Afootnote{H\textsuperscript{pc}; \skt{दुर्बद्धेति}\ K; \skt{दुर्ब्बल त्वेति}\ H\textsuperscript{ac}}}\skt{ पञ्च बस्तिदोषाः~। अतिपीडितता शिथि\-ल\-पी\-डि\-तता भूयो }\edtext{\skt{वपीडितता}}{\lemma{\skt{वपीडितता}}\Afootnote{H; \skt{भूयो\appbrk{}व\appbrk{}पी\appbrk{}डि\appbrk{}तता}\ K}}\skt{ काला\-ति\-क्र\-म\-तेति }\edtext{\skt{चत्वारः}}{\lemma{\skt{चत्वारः}}\Afootnote{H; \skt{चत्वा}\ K}}\skt{ पीडनदोषाः~। आमं }\edtext{\edtext{\edtext{\skt{हीन\-म\-ति\-मा\-त्र\-म\-ति\-शी\-त\-म\-त्यु\-ष्ण\-म\-ति\-ती\-क्ष्ण\-ण\-ति\-वृ\-द्ध\-म\-ति\-स्नि\-ग्ध\-म\-ति\-रू\-क्ष्म\-ति\-सा\-न्द्र\-म\-ति\-द्र\-व\-मि\-त्ये\-का\-द\-श\-द्र\-व्य\-दोषाः}}{\lemma{\skt{अति\appbrk{}ती\appbrk{}क्ष्ण\appbrk{}ण\appbrk{}ति\appbrk{}वृ\appbrk{}द्धम्}}\Afootnote{H; \skt{अति\appbrk{}ती\appbrk{}क्ष्ण\appbrk{}म\appbrk{}तिमृ}\ K}}}{\lemma{\skt{अतिस्निग्धम्}}\Afootnote{H; \skt{द्वति स्निग्धम्}\ K}}}{\lemma{\skt{अतिरूक्ष्म्}}\Afootnote{H; \skt{अतिरूक्षम्}\ K}}\skt{~। }\edtext{\skt{अवाङ्चीर्षो}}{\lemma{\skt{अवाङ्चीर्षो}}\Afootnote{H; \skt{अवाङ्छीर्षो}\ K}}\skt{ }\edtext{\edtext{\skt{च्छीर्षो\-न्यु\-ब्जौ\-त्ता\-न\-स\-ङ्कु\-चि\-तो\-त्थि\-त\-स\-त्थि\-ता\-द\-द\-क्षि\-ण\-पा\-र्श्व\-शा\-यि\-न\-म\-व\-न\-त\-प्र\-ता\-न\-मि\-त्यष्टौ}}{\lemma{\skt{च्छीर्षोन्य्}}\Afootnote{H; \skt{च्छीर्षो}\ K}}}{\lemma{\skt{उब्जौ\appbrk{}त्ता\appbrk{}न\appbrk{}स\appbrk{}ङ्कु\appbrk{}चि\appbrk{}तो\appbrk{}त्थि\appbrk{}त\appbrk{}स\appbrk{}त्थि\appbrk{}ता\appbrk{}द\appbrk{}द\appbrk{}क्षि\appbrk{}ण\appbrk{}पा\appbrk{}र्श्व\appbrk{}शा\appbrk{}यि\appbrk{}नम्}}\Afootnote{H; \skt{न्युब्जौ\appbrk{}त्ता\appbrk{}न\appbrk{}सं\appbrk{}कु\appbrk{}चि\appbrk{}तो\appbrk{}त्थित सत्थिताद क्षिण\appbrk{}पा\appbrk{}र्श्व\appbrk{}शा\appbrk{}यि\appbrk{}नम्}\ K}}\skt{ शय्यादोषाः~। एवमेताः पञ्च\-च\-त्वा\-रिं\-श\-द्व्या\-पदो }\edtext{\skt{वैद्यनिमत्ता}}{\lemma{\skt{वैद्यनिमत्ता}}\Afootnote{H; \skt{वैद्यनिमित्ता}\ K}}\skt{ भवन्ति~।आ\-तु\-र\-नि\-मि\-त्तास्तु पञ्चदश }\edtext{\skt{आतु\-रो\-प\-द्र\-वीये}}{\lemma{\skt{आतु\appbrk{}रो\appbrk{}प\appbrk{}द्र\appbrk{}वीये}}\Afootnote{H; \skt{आतु\appbrk{}रो\appbrk{}प\appbrk{}द्र\appbrk{}वीय}\ K}}\skt{ वक्ष्यन्ते~। }\edtext{\skt{स्नेहस्त्वष्टाभिः}}{\lemma{\skt{स्नेहस्त्वष्टाभिः}}\Afootnote{H; \skt{स्नेहस्त्वष्टाभिः}\ K}}\skt{ कारणैः प्रणिहितो न प्रत्यागच्छति~। तद्यथा त्रिभिर्दोषै रसनाभिभूतो मलव्यामिश्रो दूरानुप्रविष्टो स्विन्नस्यानुष्णोल्पं भुक्तवतो ल्पांशश्चेत्येताः~। वैद्या\-तु\-र\-नि\-मित्ता भवन्ति~। उभ\-यो\-र्ब्ब\-स्त्यो\-र\-योगः~। }\edtext{\skt{आध्मानं}}{\lemma{\skt{आध्मानं}}\Afootnote{H; \skt{आध्मार्न}\ K}}\skt{ परिकर्त्तः परि\-स्रा\-वः\-प्र\-वा\-हिका }\edtext{\skt{हृद\-यो\-प\-स\-र\-ण\-म\-ङ्ग\-प्र\-ग्रहो}}{\lemma{\skt{अङ्गप्रग्रहो}}\Afootnote{H; \skt{अङ्गप्रग्रहो}\ K}}\skt{ }\edtext{\skt{तियोगो}}{\lemma{\skt{तियोगो}}\Afootnote{H; \skt{तियोगो}\ K}}\skt{ जीवा\-दा\-न\-मि\-त्ये\-ता\-न\-व\-व्या\-पदो वैद्यनिमित्ता एव भवन्ति~॥ षट्सप्तभिः समासेन व्यापदः परिकीर्त्तिताः~।तासां वक्ष्यामि विज्ञानं सिद्धिञ्च तदनन्तरं~॥}\pend

\pstart\versepar \ledsidenote{\passref{4.36.3}}\edtext{\skt{नेत्रे}}{\lemma{\skt{नेत्रे}}\Afootnote{K H; \skt{अथ नेत्रे}\ A}}\skt{ विचलिते }\edtext{\skt{चापि}}{\lemma{\skt{चापि}}\Afootnote{K H; \skt{तथा चैव वि}\ A}}\skt{ वर्त्तिते }\edtext{\skt{गुदवेष्टनं}}{\lemma{\skt{गुदवेष्टनं}}\Afootnote{K H; \skt{गुदे}\ A}}\skt{~। }\pend
\pstart\versepar \skt{}\edtext{\skt{रुक्क्षतम्वा}}{\lemma{\skt{रुक्क्षतम्वा}}\Afootnote{K H; \skt{क्षतं रु}\ A}}\skt{ }\edtext{\skt{भवेत्तत्र}}{\lemma{\skt{भवेत्तत्र}}\Afootnote{K H; \skt{जा वा स्यात्तत्र}\ A}}\skt{ }\edtext{\skt{विधिः}}{\lemma{\skt{विधिः}}\Afootnote{K H; \skt{सद्यः}\ A}}\skt{ }\edtext{\skt{पित्तक्षतापहः}}{\lemma{\skt{पित्तक्षतापहः}}\Afootnote{K H; \skt{क्षतक्रियाः~॥}\ A}}\skt{~॥}\pend

\pstart\versepar \ledsidenote{\passref{4.36.4}}\edtext{\edtext{\skt{अत्यु\-त्क्षि\-प्तेऽ\-व\-सन्ने}}{\lemma{\skt{अत्युत्क्षिप्ते}}\Afootnote{A H; \skt{अत्युक्षिप्ते}\ K}}\skt{ च}}{\lemma{\skt{ऽवसन्ने च}}\Afootnote{A H; \skt{वसन्ने च}\ K}}\skt{ नेत्रे }\edtext{\skt{पायौ च}}{\lemma{\skt{पायौ च}}\Afootnote{K H; \skt{पायौ भ}\ A}}\skt{ }\edtext{\skt{वेदना}}{\lemma{\skt{वेदना}}\Afootnote{K H; \skt{वेद्रुजा~।}\ A}}\skt{~। }\pend
\pstart\versepar \skt{}\edtext{\skt{भवत्यत्रापि}}{\lemma{\skt{भवत्यत्रापि}}\Afootnote{H; \skt{विधिरत्रापि}\ A; \skt{भवत्यत्रापि}\ K}}\skt{ }\edtext{\skt{पित्तघ्नो}}{\lemma{\skt{पित्तघ्नो}}\Afootnote{A H; \skt{पिघ्नौ}\ K\textsuperscript{ac}; \skt{पित्तघ्नौ}\ K\textsuperscript{pc}}}\skt{ }\edtext{\skt{विधिः}}{\lemma{\skt{विधिः}}\Afootnote{K H; \skt{कार्यः}\ A}}\skt{ स्नेहैश्च }\edtext{\skt{सेवनं}}{\lemma{\skt{सेवनं}}\Afootnote{K H; \skt{सेचनम्~॥}\ A}}\skt{~॥}\pend

\pstart\versepar \ledsidenote{\passref{4.36.5}}\edtext{\skt{पार्श्वावपीडे}}{\lemma{\skt{पार्श्वावपीडे}}\Afootnote{K H; \skt{तिर्यक्प्रणिहिते नेत्रे तथा पार्श्वावपी}\ A}}\skt{ }\edtext{\skt{तिर्यक्च}}{\lemma{\skt{तिर्यक्च}}\Afootnote{K H; \skt{डि}\ A}}\skt{ }\edtext{\skt{क्षिप्ते}}{\lemma{\skt{क्षिप्ते}}\Afootnote{K H; \emph{om.}\ A}}\skt{ }\edtext{\skt{नेत्रे}}{\lemma{\skt{नेत्रे}}\Afootnote{K H; \skt{ते~।}\ A}}\skt{ }\edtext{\skt{मुखावृते}}{\lemma{\skt{मुखावृते}}\Afootnote{K; \skt{मुख\appbrk{}स्या\appbrk{}व\appbrk{}र\appbrk{}णा\appbrk{}द्ब\appbrk{}स्तिर्न सम्य\appbrk{}क्प्र\appbrk{}ति\appbrk{}प\appbrk{}द्यते~। ऋजु नेत्रं विधे}\ A; \skt{मुखावृते~। बस्तिर्न्न}\ H}}\skt{~। }\pend
\pstart\versepar \skt{}\edtext{\edtext{\skt{बस्तिर्न्न}}{\lemma{\skt{बस्तिर्}}\Afootnote{\skt{यं स्यात्तत्रर}\ A; \skt{बस्तिर्न}\ K; \emph{om.}\ H}}\skt{ क्रमते}}{\lemma{\skt{न्न क्रमते}}\Afootnote{\skt{सम्यग्विजानता~॥}\ A; \skt{क्रमते}\ K H}}\skt{ }\edtext{\skt{रुक्चगुदे}}{\lemma{\skt{रुक्चगुदे}}\Afootnote{K H; \emph{om.}\ A}}\skt{ }\edtext{\skt{सिद्धिश्च}}{\lemma{\skt{सिद्धिश्च}}\Afootnote{K H; \emph{om.}\ A}}\skt{ }\edtext{\skt{पूर्व्ववत्}}{\lemma{\skt{पूर्व्ववत्}}\Afootnote{K H; \emph{om.}\ A}}\skt{~॥}\pend

\pstart\versepar \ledsidenote{\passref{4.36.6}}\skt{अतिस्थूले कर्क्कशे च }\edtext{\skt{तथा\-श्रि\-म\-ति\-घ\-र्षणं}}{\lemma{\skt{तथा\appbrk{}श्रि\appbrk{}म\appbrk{}ति\appbrk{}घ\appbrk{}र्षणं}}\Afootnote{K H; \skt{नेत्रेऽस्त्रिमति घर्षणा}\ A}}\skt{~। }\pend
\pstart\versepar \skt{}\edtext{\skt{पायौ}}{\lemma{\skt{पायौ}}\Afootnote{K H; \skt{त्~। गुदे भवे}\ A}}\skt{ }\edtext{\skt{ततो}}{\lemma{\skt{ततो}}\Afootnote{K H; \skt{त्क्षतं}\ A}}\skt{ }\edtext{\skt{रुक्क्षतञ्च}}{\lemma{\skt{रुक्क्षतञ्च}}\Afootnote{H; \skt{रुक्च}\ A; \skt{रुक्क्षतश्च}\ K}}\skt{ }\edtext{\edtext{\skt{सिद्धिस्तत्रापि}}{\lemma{\skt{सिद्धिस्}}\Afootnote{K H; \skt{साध}\ A}}}{\lemma{\skt{तत्रापि}}\Afootnote{K H; \skt{नं तस्य}\ A}}\skt{ पूर्ववत्~॥}\pend

\pstart\versepar \ledsidenote{\passref{4.36.7}}\edtext{\skt{निकृष्टकर्ण्णिके}}{\lemma{\skt{निकृष्टकर्ण्णिके}}\Afootnote{K H; \skt{आसन्नकर्णिके}\ A}}\skt{ नेत्रे }\edtext{\skt{भिन्नेनौ}}{\lemma{\skt{भिन्नेनौ}}\Afootnote{H; \skt{भिन्नेऽणौ वाऽ}\ A; \skt{भिन्नेणौ वा}\ K}}\skt{ }\edtext{\skt{वाप्यपार्थकः}}{\lemma{\skt{वाप्यपार्थकः}}\Afootnote{H; \skt{प्यपार्थकः~।}\ A; \skt{प्यपार्थकः~।}\ K}}\skt{~। }\pend
\pstart\versepar \skt{अवसेको }\edtext{\edtext{\skt{भवे\-द्ब\-स्ते\-स्ता\-न्दो\-षा\-न्प\-रि\-व\-र्ज्ज\-येत्}}{\lemma{\skt{भवेद्बस्तेस्तान्}}\Afootnote{H; \skt{भवेद्बस्तेस्त}\ A; \skt{भवेद्बस्तेस्तान्}\ K}}}{\lemma{\skt{दोषा\appbrk{}न्प\appbrk{}रि\appbrk{}व\appbrk{}र्ज्ज\appbrk{}येत्}}\Afootnote{H; \skt{स्माद्दो\appbrk{}षा\appbrk{}न्वि\appbrk{}व\appbrk{}र्ज\appbrk{}येत्~॥}\ A; \skt{दोषां परिवर्जयेत्~॥}\ K}}\skt{~॥}\pend

\pstart\versepar \ledsidenote{\passref{4.36.8}}\skt{प्रकृष्टकर्ण्णिके रक्तं गुद\-म\-र्म्म\-प्र\-पी\-ड\-नात्~। }\pend
\pstart\versepar \skt{क्षरत्यत्रापि पित्तघ्नो }\edtext{\edtext{\skt{विधिर्ब्बस्तिश्च}}{\lemma{\skt{विधिर्}}\Afootnote{K H; \skt{विधि}\ A}}}{\lemma{\skt{ब्बस्तिश्च}}\Afootnote{K H; \skt{र्बस्तिश्च}\ A}}\skt{ पिच्छिलः~॥}\pend

\pstart\versepar \ledsidenote{\passref{4.36.9}}\skt{ह्रस्वे त्वणुस्रोतसि च क्लेशो बस्तिश्च }\edtext{\skt{पूर्व्वशः}}{\lemma{\skt{पूर्व्वशः}}\Afootnote{K H; \skt{पूर्वव}\ A}}\skt{~। }\pend
\pstart\versepar \skt{}\edtext{\skt{प्रत्यागच्छन्नतः}}{\lemma{\skt{प्रत्यागच्छन्नतः}}\Afootnote{K H; \skt{त्~। प्रत्यागच्छंस्ततः}\ A}}\skt{ }\edtext{\skt{कुर्या\-द्घो\-रा\-न्रो\-गा\-न्वि\-घा\-त\-जान्}}{\lemma{\skt{कुर्या\appbrk{}द्घो\appbrk{}रा\appbrk{}न्रो\appbrk{}गा\appbrk{}न्वि\appbrk{}घा\appbrk{}त\appbrk{}जान्}}\Afootnote{H; \skt{कुर्या\appbrk{}द्रो\appbrk{}गा\appbrk{}न्ब\appbrk{}स्ति\appbrk{}वि\appbrk{}घा\appbrk{}त\appbrk{}जान्~॥}\ A; \skt{कुर्याद्घोरान्रोगां विघातजां~॥}\ K}}\skt{~॥}\pend

\pstart\versepar \ledsidenote{\passref{4.36.10}}\skt{दीर्घे महास्रोतसि च }\edtext{\skt{ज्ञेय\-म\-त्य\-व\-पी\-ड\-वत्}}{\lemma{\skt{ज्ञेयम्}}\Afootnote{K H; \skt{ज्ञेय}\ A}}\skt{~। }\pend
\pstart\versepar \skt{प्रस्त्याने बहले चापि बस्तौ }\edtext{\skt{दुर्व्व\-द्ध\-दो\-ष\-वत्}}{\lemma{\skt{दुर्व्व\appbrk{}द्ध\appbrk{}दो\appbrk{}ष\appbrk{}वत्}}\Afootnote{H; \skt{दुर्ब\appbrk{}द्ध\appbrk{}दो\appbrk{}ष\appbrk{}वत्}\ A K}}\skt{~॥}\edtext{}{\lemma{\skt{अत्य\appbrk{}व\appbrk{}पी\appbrk{}ड\appbrk{}वत्~। प्रस्त्याने}}\Afootnote{K H; \skt{मत्य\appbrk{}व\appbrk{}पी\appbrk{}ड\appbrk{}वत्~। प्रस्तीर्णे}\ A}}\pend

\pstart\versepar \ledsidenote{\passref{4.36.11}}\edtext{\edtext{\skt{ज्ञेयमल्पोल्पता}}{\lemma{\skt{ज्ञेयम्}}\Afootnote{K H; \skt{बस्ताव्}\ A}}}{\lemma{\skt{अल्पोल्पता}}\Afootnote{K H; \skt{अल्पेऽल्पता}\ A}}\skt{ }\edtext{\skt{चापि}}{\lemma{\skt{चापि}}\Afootnote{K H; \skt{वाऽपि}\ A}}\skt{ }\edtext{\skt{द्रव्यस्यात्मगुणा}}{\lemma{\skt{द्रव्यस्यात्मगुणा}}\Afootnote{K H; \skt{द्रव्यस्याल्पा गुणा}\ A}}\skt{ मता~। }\pend
\pstart\versepar \skt{}\edtext{\skt{दुर्व्वद्धे}}{\lemma{\skt{दुर्व्वद्धे}}\Afootnote{H; \skt{दुर्बद्धे}\ A K}}\skt{ }\edtext{\skt{चैव}}{\lemma{\skt{चैव}}\Afootnote{K H; \skt{चाणु}\ A}}\skt{ }\edtext{\skt{च्छिद्रे च}}{\lemma{\skt{च्छिद्रे च}}\Afootnote{K H; \skt{भिन्ने च}\ A}}\skt{ विज्ञे\-य\-म्भि\-न्न\-ने\-त्र\-वत्~॥}\pend

\pstart\versepar \ledsidenote{\passref{4.36.12}}\edtext{\skt{अतिप्रपीडितो}}{\lemma{\skt{अतिप्रपीडितो}}\Afootnote{A K; \skt{अतिप्रपी}\ H}}\skt{ }\edtext{\skt{बस्तिः}}{\lemma{\skt{बस्तिः}}\Afootnote{A K; \emph{om.}\ H}}\skt{ }\edtext{\skt{प्रयात्यामाशयं}}{\lemma{\skt{प्रयात्यामाशयं}}\Afootnote{A; \skt{प्रयात्यामाशयं}\ K; \emph{om.}\ H}}\skt{ }\edtext{\skt{ततः}}{\lemma{\skt{ततः}}\Afootnote{A K; \emph{om.}\ H}}\skt{~॥ }\pend
\pstart\versepar \skt{}\edtext{\skt{वातेरितो}}{\lemma{\skt{वातेरितो}}\Afootnote{A K; \emph{om.}\ H}}\skt{ }\edtext{\skt{नासिकाभ्यां}}{\lemma{\skt{नासिकाभ्यां}}\Afootnote{A K; \emph{om.}\ H}}\skt{ }\edtext{\skt{मुखतो वा}}{\lemma{\skt{मुखतो वा}}\Afootnote{A K; \emph{om.}\ H}}\skt{ }\edtext{\skt{प्रवर्त्तते}}{\lemma{\skt{प्रवर्त्तते}}\Afootnote{K; \skt{प्रपद्यते~॥}\ A; \emph{om.}\ H}}\skt{~॥}\pend

\pstart \ledsidenote{\passref{4.36.12ef}}\edtext{\edtext{\skt{च्छर्दिर्मूर्च्छा}}{\lemma{\skt{च्छर्दिर्}}\Afootnote{K\textsuperscript{pc}; \emph{om.}\ K\textsuperscript{ac}}}}{\lemma{\skt{मूर्च्छा}}\Afootnote{K\textsuperscript{pc}; \emph{om.}\ K\textsuperscript{ac}}}\skt{ }\edtext{\skt{सहृल्लास}}{\lemma{\skt{सहृल्लास}}\Afootnote{K\textsuperscript{pc}; \emph{om.}\ K\textsuperscript{ac}}}\skt{ }\edtext{\skt{पिपिसा}}{\lemma{\skt{पिपिसा}}\Afootnote{K\textsuperscript{pc}; \emph{om.}\ K\textsuperscript{ac}}}\skt{ }\edtext{\skt{दाह}}{\lemma{\skt{दाह}}\Afootnote{K\textsuperscript{pc}; \emph{om.}\ K\textsuperscript{ac}}}\skt{ }\edtext{\skt{शोषकृत्}}{\lemma{\skt{शोषकृत्}}\Afootnote{K\textsuperscript{pc}; \emph{om.}\ K\textsuperscript{ac}}}\pend

\pstart\versepar \ledsidenote{\passref{4.36.13}}\skt{तत्र }\edtext{\skt{तूर्णगलापीडं}}{\lemma{\skt{तूर्णगलापीडं}}\Afootnote{K; \skt{तूर्णं गलापीडं}\ A}}\skt{ कुर्याच्चापि }\edtext{\skt{विधूननं}}{\lemma{\skt{विधूननं}}\Afootnote{K; \skt{अवधूननम्~।}\ A}}\skt{~। }\pend
\pstart\versepar \skt{शिरः कायविरेकौ च तीक्ष्णौ सेकांश्च }\edtext{\skt{शीतलां}}{\lemma{\skt{शीतलां}}\Afootnote{K; \skt{शीतलान्~॥}\ A}}\skt{~।}\pend

\pstart\versepar \ledsidenote{\passref{4.36.14}}\edtext{\skt{मृदुना}}{\lemma{\skt{मृदुना}}\Afootnote{K; \emph{om.}\ A}}\skt{ }\edtext{\skt{चावपीडेन}}{\lemma{\skt{चावपीडेन}}\Afootnote{K; \skt{शनैः प्रपीडितो बस्तिः}\ A}}\skt{ पक्वाधानं च }\edtext{\skt{गच्छति}}{\lemma{\skt{गच्छति}}\Afootnote{K; \skt{गच्छति~। न}\ A}}\skt{~। }\pend
\pstart\versepar \skt{}\edtext{\edtext{\skt{बस्तिर्नार्थकरः}}{\lemma{\skt{बस्तिर्}}\Afootnote{K; \skt{च सं}\ A}}}{\lemma{\skt{नार्थकरः}}\Afootnote{K; \skt{पादयत्य्}\ A}}\skt{ }\edtext{\skt{प्रोक्त}}{\lemma{\skt{प्रोक्त}}\Afootnote{K; \skt{अर्थं}\ A}}\skt{ }\edtext{\skt{तस्माद्युक्त्या}}{\lemma{\skt{तस्माद्युक्त्या}}\Afootnote{K; \skt{तस्माद्युक्तं}\ A}}\skt{ प्रपीडयेत्~॥}\pend

\pstart\versepar \ledsidenote{\passref{4.36.15}}\skt{भूयो }\edtext{\skt{भूयोवपीडेन}}{\lemma{\skt{भूयोवपीडेन}}\Afootnote{\skt{भूयोऽवपीडेन}\ A; \skt{भूयो वपीडेन}\ K}}\skt{ वायुरन्तः }\edtext{\skt{प्रपद्यते}}{\lemma{\skt{प्रपद्यते}}\Afootnote{K; \skt{प्रपीड्यते~।}\ A}}\skt{~। }\pend
\pstart\versepar \skt{तेनाध्मानं रुजश्चोग्रा यथास्वं तत्र बस्तयः~।}\pend

\pstart\versepar \ledsidenote{\passref{4.36.16}}\skt{काला\-ति\-क्र\-म\-णा\-त्क्लेशो }\edtext{\skt{व्याधि\-श्चा\-भि\-प्र\-व\-र्त्तते}}{\lemma{\skt{चाभिप्रवर्त्तते}}\Afootnote{K; \skt{चाभिप्रवर्धते~।}\ A}}\skt{~। }\pend
\pstart\versepar \skt{तत्र व्याधि बलघ्नं तु भूयो बस्तिं निधापयेत्~।}\pend

\pstart\versepar \ledsidenote{\passref{4.36.17}}\edtext{\skt{गुदो\-प\-दे\-ह\-शोफौ तु}}{\lemma{\skt{गुदो\appbrk{}प\appbrk{}दे\appbrk{}ह\appbrk{}शोफौ तु}}\Afootnote{K; \skt{गोदो\appbrk{}प\appbrk{}दे\appbrk{}ह\appbrk{}शोफौ तु}\ A}}\skt{ }\edtext{\skt{स्नेहोपक्वः}}{\lemma{\skt{स्नेहोपक्वः}}\Afootnote{K; \skt{स्नेहोऽपक्वः}\ A}}\skt{ करोति हि~। }\pend
\pstart\versepar \skt{तत्र संशोधनो बस्तिर्हितञ्चापि विरेचनं~।}\pend

\pstart\versepar \ledsidenote{\passref{4.36.18}}\edtext{\edtext{\skt{हीनमात्रावुभौ}}{\lemma{\skt{हीनमात्राव्}}\Afootnote{K; \skt{हीनमात्रा}\ A}}}{\lemma{\skt{उभौ}}\Afootnote{K; \skt{वुभौ}\ A}}\skt{ बस्ती }\edtext{\skt{हीनकार्यकरौ}}{\lemma{\skt{हीनकार्यकरौ}}\Afootnote{K; \skt{नातिकार्यकरौ}\ A}}\skt{ मतौ~। }\pend
\pstart\versepar \skt{अतिमात्रौ }\edtext{\skt{समु\-न्ना\-ह\-क्ल\-मा\-ती\-सा\-र\-का\-रकौ}}{\lemma{\skt{समु\appbrk{}न्ना\appbrk{}ह\appbrk{}क्ल\appbrk{}मा\appbrk{}ती\appbrk{}सा\appbrk{}र\appbrk{}का\appbrk{}रकौ}}\Afootnote{K; \skt{तथाऽऽ\appbrk{}आ\appbrk{}ना\appbrk{}ह\appbrk{}क्ल\appbrk{}मा\appbrk{}ती\appbrk{}सा\appbrk{}र\appbrk{}का\appbrk{}रकौ~॥}\ A}}\skt{~।}\pend

\pstart\versepar \ledsidenote{\passref{4.36.19}}\edtext{\skt{तीक्ष्णाष्णौ}}{\lemma{\skt{तीक्ष्णाष्णौ}}\Afootnote{K; \emph{om.}\ A}}\skt{ }\edtext{\skt{पित्तलौ}}{\lemma{\skt{पित्तलौ}}\Afootnote{K; \emph{om.}\ A}}\skt{ }\edtext{\skt{मूर्च्छा}}{\lemma{\skt{मूर्च्छा}}\Afootnote{K; \skt{मूर्च्छां}\ A}}\skt{ }\edtext{\skt{दाहा\-ती\-सा\-र\-का\-रकौ}}{\lemma{\skt{दाहा\appbrk{}ती\appbrk{}सा\appbrk{}र\appbrk{}का\appbrk{}रकौ}}\Afootnote{K; \skt{दाहमतीसारं पित्तं चात्युष्णतीक्ष्णकौ~।}\ A}}\skt{~। }\pend
\pstart\versepar \skt{मृदुशीतावुभौ वात\-वि\-ब\-न्धा\-ध्मा\-न\-का\-रकौ~।}\pend

\pstart \ledsidenote{\passref{4.36.20}}\skt{तत्रहीनादिषु हितः }\edtext{\skt{प्रत्यनीक}}{\lemma{\skt{प्रत्यनीक}}\Afootnote{K; \skt{प्रत्यनीकः}\ A}}\skt{ }\edtext{\skt{क्रियाविधिः}}{\lemma{\skt{क्रियाविधिः}}\Afootnote{K; \skt{क्रियाविधिः~। गुद\appbrk{}ब\appbrk{}स्त्यु\appbrk{}प\appbrk{}देहं तु कुर्यात्सान्द्रो निरूहणः~॥}\ A}}\skt{~।}\pend

\pstart\versepar \ledsidenote{\passref{4.36.22}}\edtext{\skt{स्निग्धो\-ति\-जा\-ड्य\-कृ\-द्रुक्षः}}{\lemma{\skt{स्निग्धो\appbrk{}ति\appbrk{}जा\appbrk{}ड्य\appbrk{}कृ\appbrk{}द्रुक्षः}}\Afootnote{K; \skt{स्निग्धोऽ\appbrk{}ति\appbrk{}जा\appbrk{}ड्य\appbrk{}कृ\appbrk{}द्रूक्षः}\ A}}\skt{ }\edtext{\skt{स्तम्भा\-ध्मा\-न\-कृ\-दु\-च्यते}}{\lemma{\skt{स्तम्भा\appbrk{}ध्मा\appbrk{}न\appbrk{}कृ\appbrk{}दु\appbrk{}च्यते}}\Afootnote{\skt{स्तम्भा\appbrk{}ध्मा\appbrk{}न\appbrk{}कृ\appbrk{}दु\appbrk{}च्यते}\ A K}}\skt{~॥ }\pend
\pstart\versepar \skt{}\edtext{\skt{बस्ति}}{\lemma{\skt{बस्ति}}\Afootnote{\skt{बस्तिं}\ A; \skt{बस्तिर्}\ K}}\skt{ }\edtext{\skt{रुक्षमतिस्निग्धे}}{\lemma{\skt{रुक्षमतिस्निग्धे}}\Afootnote{\skt{रूक्षमतिस्निग्धे}\ A; \skt{उक्षमतिस्निग्धे}\ K}}\skt{ स्निग्धं रूक्षे }\edtext{\skt{प्रयोजयेत्}}{\lemma{\skt{प्रयोजयेत्}}\Afootnote{K; \skt{च दापयेत्~॥}\ A}}\skt{~।}\pend

\pstart\versepar \ledsidenote{\passref{4.36.22.1}}\edtext{\skt{बस्ति\-र्गु\-दो\-प\-लेपं तु}}{\lemma{\skt{गुदोपलेपं तु}}\Afootnote{\skt{गुदो पलेपं तु}\ K}}\skt{ कुर्यात्सान्द्रो निरूहणः~। }\pend
\pstart\versepar \skt{प्रवाहिकाम्वा जन\-ये\-त्त\-नु\-र\-ल्प\-गु\-णा\-वहः~। }\pend
\pstart\versepar \skt{तत्र सान्द्रे तनुम्बस्तिं तनौ सान्द्रं च दापयेत्~।}\pend

\pstart\versepar \ledsidenote{\passref{4.36.23}}\edtext{\skt{अति\-पी\-ड\-न\-व\-द्दोषां}}{\lemma{\skt{अति\appbrk{}पी\appbrk{}ड\appbrk{}न\appbrk{}व\appbrk{}द्दोषां}}\Afootnote{\skt{अति\appbrk{}पी\appbrk{}डि\appbrk{}त\appbrk{}व\appbrk{}द्दोषा}\ A; \skt{अति\appbrk{}पी\appbrk{}ड\appbrk{}न\appbrk{}व\appbrk{}द्दोषां}\ K}}\skt{ }\edtext{\skt{विद्धि}}{\lemma{\skt{विद्धि}}\Afootnote{K; \skt{न्विद्धि}\ A}}\skt{ चाप्यवशीर्षके~। }\pend
\pstart\versepar \skt{अतिपीडित सिद्धिञ्च यथाप्राप्तां प्रयोजयेत्~। }\pend
\pstart\versepar \skt{उच्छी\-र्ष\-क\-स\-मु\-न्नाहो }\edtext{\skt{बस्तः}}{\lemma{\skt{बस्तः}}\Afootnote{K; \skt{बस्तिः}\ A}}\skt{ }\edtext{\skt{कृच्छ्राच्च}}{\lemma{\skt{कृच्छ्राच्च}}\Afootnote{K; \skt{कुर्याच्च}\ A}}\skt{ मेहनं~।}\edtext{}{\lemma{\skt{अतिपीडित… उ}}\Afootnote{\emph{om.}\ A}}\edtext{}{\lemma{\skt{अतिपीडित… उ}}\Afootnote{\emph{om.}\ K\textsuperscript{ac}}}\pend

\pstart\versepar \ledsidenote{\passref{4.36.24}}\skt{तत्रोत्तरो हितो }\edtext{\skt{बस्ति}}{\lemma{\skt{बस्ति}}\Afootnote{K; \skt{बस्तिः}\ A}}\skt{ सुस्विन्नस्य सुखावहः~। }\pend
\pstart\versepar \skt{न्युब्जस्य बस्तिर्नाप्नोति }\edtext{\skt{पक्वाधाणं}}{\lemma{\skt{पक्वाधाणं}}\Afootnote{K; \skt{पक्वाधानं}\ A}}\skt{ }\edtext{\skt{विर्निर्गतः}}{\lemma{\skt{विर्निर्गतः}}\Afootnote{K; \skt{विमार्गगः~॥}\ A}}\skt{~।}\pend

\pstart\versepar \ledsidenote{\passref{4.36.25}}\edtext{\skt{हृद्गुदे}}{\lemma{\skt{हृद्गुदे}}\Afootnote{K; \skt{हृद्गुदं}\ A}}\skt{ बाधते चात्र वायुः }\edtext{\edtext{\skt{पृष्ठमथापि}}{\lemma{\skt{पृष्ठम्}}\Afootnote{K; \skt{कोष्ठ}\ A}}\skt{ च}}{\lemma{\skt{अथापि च}}\Afootnote{K; \skt{मथापि च~।}\ A}}\skt{~। }\pend
\pstart\versepar \skt{उत्तानस्यावृते मार्गे बस्तिः }\edtext{\skt{प्रस्थेति}}{\lemma{\skt{प्रस्थेति}}\Afootnote{K; \skt{नान्तः प्रपद्यते~॥}\ A}}\skt{ }\edtext{\skt{देहिनः}}{\lemma{\skt{देहिनः}}\Afootnote{K; \emph{om.}\ A}}\skt{~।}\pend

\pstart\versepar \ledsidenote{\passref{4.36.26}}\edtext{\skt{कुञ्चितः}}{\lemma{\skt{कुञ्चितः}}\Afootnote{K; \skt{नेत्र\appbrk{}सं\appbrk{}वे\appbrk{}ज\appbrk{}न\appbrk{}भ्रान्तो वायुश्चान्तः प्रकुप्यति~। देहे सङ्कुचिते दत्तः}\ A}}\skt{ }\edtext{\skt{सक्थदेहस्य}}{\lemma{\skt{सक्थदेहस्य}}\Afootnote{K; \skt{सक्थ्नोरप्यु}\ A}}\skt{ }\edtext{\edtext{\skt{यत्क\-म\-य\-म\-सं\-प्राप्य}}{\lemma{\skt{यत्कमयम्}}\Afootnote{K; \skt{भयोस्तथा~॥}\ A}}}{\lemma{\skt{असंप्राप्य}}\Afootnote{K; \emph{om.}\ A}}\skt{~। }\pend
\pstart\versepar \skt{}\edtext{\edtext{\skt{बस्तिर्निवर्तते}}{\lemma{\skt{बस्तिर्}}\Afootnote{K; \emph{om.}\ A}}}{\lemma{\skt{निवर्तते}}\Afootnote{K; \emph{om.}\ A}}\skt{ }\edtext{\skt{कृच्छ्रादन्य}}{\lemma{\skt{कृच्छ्रादन्य}}\Afootnote{K; \emph{om.}\ A}}\skt{ }\edtext{\skt{तत्र}}{\lemma{\skt{तत्र}}\Afootnote{K; \emph{om.}\ A}}\skt{ }\edtext{\skt{नियोजयेत्}}{\lemma{\skt{नियोजयेत्}}\Afootnote{K; \emph{om.}\ A}}\skt{~।}\pend

\pstart \ledsidenote{\passref{4.36.27cd}}\skt{स्थितस्य बस्तिर्दत्तस्तु }\edtext{\edtext{\edtext{\skt{क्षिप्र\-मा\-या\-त्य\-वा\-ङ्मुखः}}{\lemma{\skt{क्षिप्रम्}}\Afootnote{K; \skt{क्षिप्र}\ A}}}{\lemma{\skt{आयात्य्}}\Afootnote{K; \skt{माया}\ A}}}{\lemma{\skt{अवाङ्मुखः}}\Afootnote{K; \skt{त्यवाङ्मुखः~॥}\ A}}\skt{~।}\pend

\pstart\versepar \ledsidenote{\passref{4.36.28}}\skt{न चाशयं तर्पयति तस्मान्नार्थकरो }\edtext{\skt{मतः}}{\lemma{\skt{मतः}}\Afootnote{K; \skt{हि सः~।}\ A}}\skt{~। }\pend
\pstart\versepar \skt{नाप्नो\-ति\-ब\-स्ति\-र्द\-त्तस्तु कृत्स्नं पक्वाशयं पुनः~।}\pend

\pstart\versepar \ledsidenote{\passref{4.36.29}}\skt{दक्षि\-णा\-श्रि\-त\-पा\-र्श्वस्य }\edtext{\skt{वामपार्श्वानुगो हि सः}}{\lemma{\skt{वामपार्श्वानुगो हि सः}}\Afootnote{K; \skt{वामपार्श्वानुगो यतः~।}\ A}}\skt{~। }\pend
\pstart\versepar \skt{न्युब्जादीनां }\edtext{\skt{यथादानं}}{\lemma{\skt{यथादानं}}\Afootnote{K; \skt{प्रदानं}\ A}}\skt{ }\edtext{\skt{तद्वद्बस्तिर्न}}{\lemma{\skt{तद्वद्बस्तिर्न}}\Afootnote{\skt{च बस्तेर्नैव}\ A; \skt{तद्वद्बस्तिर्न}\ K}}\skt{ }\edtext{\skt{शस्यते}}{\lemma{\skt{शस्यते}}\Afootnote{K; \skt{प्रशस्यते~॥}\ A}}\skt{~॥}\pend

\pstart\versepar \ledsidenote{\passref{4.36.30}}\edtext{\skt{यश्चा\-प्य\-नि\-ल\-कोपो}}{\lemma{\skt{यश्चा\appbrk{}प्य\appbrk{}नि\appbrk{}ल\appbrk{}कोपो}}\Afootnote{K; \skt{पश्चा\appbrk{}द\appbrk{}नि\appbrk{}ल\appbrk{}कोपो}\ A}}\skt{ त्र यथास्वं तत्र }\edtext{\skt{सिद्धयः}}{\lemma{\skt{सिद्धयः}}\Afootnote{K; \skt{कारयेत्~। व्यापदः}\ A}}\skt{~। }\pend
\pstart\versepar \skt{}\edtext{\skt{स्नेह\-ब\-स्ति\-नि\-मि\-त्तास्तु}}{\lemma{\skt{स्नेह\appbrk{}ब\appbrk{}स्ति\appbrk{}नि\appbrk{}मि\appbrk{}त्तास्तु}}\Afootnote{K H; \skt{स्नेहबस्तेस्तु}\ A}}\skt{ वक्ष्यन्ते }\edtext{\skt{व्यापदः}}{\lemma{\skt{व्यापदः}}\Afootnote{K H; \emph{om.}\ A}}\skt{ }\edtext{\skt{पुरा}}{\lemma{\skt{पुरा}}\Afootnote{K; \skt{तच्चिकित्सिते~॥}\ A; \skt{पराः~।}\ H\textsuperscript{ac}; \skt{परा~।}\ H\textsuperscript{pc}}}\skt{~।}\edtext{}{\lemma{\skt{… सिद्धयः}}\Afootnote{\emph{om.}\ H\textsuperscript{ac}}}\edtext{}{\lemma{\skt{… सिद्धयः}}\Afootnote{\emph{om.}\ H\textsuperscript{pc}}}\pend

\pstart\versepar \ledsidenote{\passref{4.36.31}}\edtext{\edtext{\skt{स्नेह\-ब\-स्ति\-वि\-धा\-वेव}}{\lemma{\skt{स्नेह\appbrk{}ब\appbrk{}स्ति\appbrk{}वि\appbrk{}धाव्}}\Afootnote{K H; \emph{om.}\ A}}}{\lemma{\skt{एव}}\Afootnote{K H; \skt{अयोगाद्यास्तु व}\ A}}\skt{ }\edtext{\skt{सल\-क्ष\-ण\-चि\-कि\-त्सिताः}}{\lemma{\skt{सल\appbrk{}क्ष\appbrk{}ण\appbrk{}चि\appbrk{}कि\appbrk{}त्सिताः}}\Afootnote{K H; \skt{क्ष्यामि व्यापदः सचिकित्सिताः~।}\ A}}\skt{~। }\pend
\pstart\versepar \skt{}\edtext{\skt{अनुष्णोल्पौषधी}}{\lemma{\skt{अनुष्णोल्पौषधी}}\Afootnote{H\textsuperscript{pc}; \skt{अनुष्णोऽल्पौषधो}\ A; \skt{अनुष्णो ल्पौषध}\ K; \skt{अनुष्णोषधी}\ H\textsuperscript{ac}}}\skt{ हीनो }\edtext{\skt{बस्तिर्न्नेति}}{\lemma{\skt{न्नेति}}\Afootnote{K H; \skt{नैति}\ A}}\skt{ प्रयोजितः~॥}\pend

\pstart\versepar \ledsidenote{\passref{4.36.32}}\edtext{\skt{विष्ट\-म्भा\-ध्मा\-न\-शू\-लैश्च तम\-यो\-ग\-म्प्र\-च\-क्षते}}{\lemma{\skt{विष्ट\appbrk{}म्भा\appbrk{}ध्मा\appbrk{}न\appbrk{}शू\appbrk{}लैश्च तम्}}\Afootnote{A H; \skt{विष्टं\appbrk{}भा\appbrk{}ध्मा\appbrk{}न\appbrk{}शू\appbrk{}लाश्च तम्}\ K}}\skt{~। }\pend
\pstart\versepar \skt{तत्र तीक्ष्णो हितो }\edtext{\edtext{\skt{बस्तिस्तीक्ष्णञ्चापि}}{\lemma{\skt{बस्तिस्}}\Afootnote{A H; \skt{बस्ति}\ K}}}{\lemma{\skt{तीक्ष्णञ्}}\Afootnote{A H; \skt{स्तीक्ष्णं}\ K}}\skt{ विरेचनं~॥}\pend

\pstart\versepar \ledsidenote{\passref{4.36.33}}\edtext{\skt{हृत\-दो\-ष\-प्र\-मा\-णेन}}{\lemma{\skt{हृत\appbrk{}दो\appbrk{}ष\appbrk{}प्र\appbrk{}मा\appbrk{}णेन}}\Afootnote{K H; \emph{om.}\ A}}\skt{ }\edtext{\skt{संसर्ग्गश्च}}{\lemma{\skt{संसर्ग्गश्च}}\Afootnote{K H; \skt{सशेषान्नेऽथवा भुक्ते बहुदोषे}\ A}}\skt{ }\edtext{\skt{विधीयते}}{\lemma{\skt{विधीयते}}\Afootnote{K H; \skt{च योजितः~।}\ A}}\skt{~॥ }\pend
\pstart\versepar \skt{}\edtext{\edtext{\edtext{\skt{अत्या\-सि\-त\-ञ्चा\-ति\-ब\-हु\-र्ब्ब\-स्ति\-र्म्म\-न्दोष्ण}}{\lemma{\skt{अत्यासितञ्}}\Afootnote{K H; \skt{अत्याशित}\ A}}}{\lemma{\skt{चातिबहुर्}}\Afootnote{H; \skt{स्यातिबहु}\ A; \skt{चाति ब}\ K}}}{\lemma{\skt{ब्बस्तिर्}}\Afootnote{H; \skt{र्बस्तिर्}\ A; \skt{स्तिर्}\ K}}\skt{ एव च~॥}\pend

\pstart \ledsidenote{\passref{4.36.34cd}}\skt{तथा }\edtext{\skt{बहुपुरीशञ्च}}{\lemma{\skt{बहुपुरीशञ्च}}\Afootnote{K; \skt{बहुपुरीषं च}\ A}}\skt{ }\edtext{\edtext{\skt{क्षिप्र\-मा\-ध्मा\-प\-ये\-न्नरं}}{\lemma{\skt{क्षिप्रम्}}\Afootnote{K; \skt{क्षिप्र}\ A}}}{\lemma{\skt{आध्मापयेन्}}\Afootnote{K; \skt{माध्मापयेन्}\ A}}\skt{~।}\pend

\pstart\versepar \ledsidenote{\passref{4.36.35}}\edtext{\skt{हृत्क\-टी\-पृ\-ष्ठ\-पा\-र्श्वेषु}}{\lemma{\skt{हृत्क\appbrk{}टी\appbrk{}पृ\appbrk{}ष्ठ\appbrk{}पा\appbrk{}र्श्वेषु}}\Afootnote{K H; \skt{हृत्क\appbrk{}टी\appbrk{}पा\appbrk{}र्श्व\appbrk{}पृ\appbrk{}ष्ठेषु}\ A}}\skt{ }\edtext{\edtext{\skt{शूला\-स्त\-त्रा\-ति\-दा\-रुणाः}}{\lemma{\skt{शूलास्}}\Afootnote{H; \skt{शूलं}\ A; \skt{शूला}\ K}}}{\lemma{\skt{तत्रातिदारुणाः}}\Afootnote{H; \skt{तत्रा\appbrk{}ति\appbrk{}दा\appbrk{}रु\appbrk{}णम्~।}\ A; \skt{स्तत्रातिदारुणाः~।}\ K}}\skt{~॥ }\pend
\pstart\versepar \skt{तत्र तीक्ष्णतरो बस्ति\-र्हि\-त\-ञ्चा\-प्य\-नु\-वा\-सनं~।}\pend

\pstart\versepar \ledsidenote{\passref{4.36.36}}\edtext{\skt{अति\-ती\-क्ष्णो\-ष्ण\-ल\-वणो}}{\lemma{\skt{अति\appbrk{}ती\appbrk{}क्ष्णो\appbrk{}ष्ण\appbrk{}ल\appbrk{}वणो}}\Afootnote{K H; \skt{अति\appbrk{}ती\appbrk{}क्ष्णो\appbrk{}ति\appbrk{}ल\appbrk{}वणो}\ A}}\skt{ रूक्षो बस्तिः प्रयोजितः~॥ }\pend
\pstart\versepar \skt{सपित्तं }\edtext{\skt{कोप\-ये\-द्वा\-यु\-ङ्कु\-र्या\-त्स\-प\-रि\-क\-र्त्ति\-कान्}}{\lemma{\skt{कुर्या\appbrk{}त्स\appbrk{}प\appbrk{}रि\appbrk{}क\appbrk{}र्त्ति\appbrk{}कान्}}\Afootnote{H; \skt{कुर्याच्च परिकर्तिकाम्~॥}\ A; \skt{कुर्या\appbrk{}त्स\appbrk{}प\appbrk{}रि\appbrk{}क\appbrk{}र्त्तिकां~।}\ K}}\skt{~।}\pend

\pstart\versepar \ledsidenote{\passref{4.36.37}}\edtext{\skt{नाभि\-ब\-स्ति\-र्ग्गु\-द\-न्तत्र}}{\lemma{\skt{नाभिबस्तिर्}}\Afootnote{H; \skt{नाभिबस्ति}\ A K}}\skt{ }\edtext{\skt{कृत्यन्त}}{\lemma{\skt{कृत्यन्त}}\Afootnote{K H; \skt{छिनत्ती}\ A}}\skt{ }\edtext{\skt{इति}}{\lemma{\skt{इति}}\Afootnote{K H; \skt{वाति}\ A}}\skt{ देहिनः~॥ }\pend
\pstart\versepar \skt{}\edtext{\skt{पिच्छा\-ब\-स्ति\-र्हि\-त\-स्तत्र}}{\lemma{\skt{तत्र}}\Afootnote{K H; \skt{तस्य}\ A}}\skt{ स्नेहश्च मधुरैः }\edtext{\skt{श्रितः}}{\lemma{\skt{श्रितः}}\Afootnote{K H; \skt{शृतः~॥}\ A}}\skt{~।}\pend

\pstart\versepar \ledsidenote{\passref{4.36.38}}\edtext{\skt{अत्यु\-ष्ण\-ती\-क्ष्ण\-ल\-वणः}}{\lemma{\skt{अत्यु\appbrk{}ष्ण\appbrk{}ती\appbrk{}क्ष्ण\appbrk{}ल\appbrk{}वणः}}\Afootnote{K H; \skt{अत्य\appbrk{}म्ल\appbrk{}ल\appbrk{}व\appbrk{}ण\appbrk{}स्तीक्ष्णः}\ A}}\skt{ परिस्रावाय कल्पते~। }\pend
\pstart\versepar \skt{}\edtext{\skt{दौर्ब्ब\-ल्य\-म\-ङ्ग\-सा\-दञ्च}}{\lemma{\skt{दौर्ब्ब\appbrk{}ल्य\appbrk{}म\appbrk{}ङ्ग\appbrk{}सा\appbrk{}दञ्च}}\Afootnote{H; \skt{दौर्ब\appbrk{}ल्य\appbrk{}म\appbrk{}ङ्ग\appbrk{}सा\appbrk{}दश्च}\ A K\textsuperscript{pc}; \skt{दौर्ब\appbrk{}म\appbrk{}ङ्ग\appbrk{}सा\appbrk{}दश्च}\ K\textsuperscript{ac}}}\skt{ जायते तत्र देहिनः~॥}\pend

\pstart\versepar \ledsidenote{\passref{4.36.39}}\edtext{\skt{परिस्रवत्यधः}}{\lemma{\skt{परिस्रवत्यधः}}\Afootnote{K; \skt{परिस्रवेत्ततः}\ A; \skt{परिस्रवत्यधः}\ H}}\skt{ पित्तं }\edtext{\skt{दाहः}}{\lemma{\skt{दाहः}}\Afootnote{K H; \skt{दाहं}\ A}}\skt{ सञ्जनयेद्गुदे~॥ }\pend
\pstart\versepar \skt{}\edtext{\skt{पिच्छा\-ब\-स्ति\-र्हि\-त\-स्तत्र}}{\lemma{\skt{पिच्छाबस्तिर्}}\Afootnote{A H; \skt{पिच्छर्बस्तिर्}\ K}}\skt{ बस्तिः }\edtext{\skt{क्षीरघृतस्य}}{\lemma{\skt{क्षीरघृतस्य}}\Afootnote{K H; \skt{क्षीरघृतेन}\ A}}\skt{ च~।}\pend

\pstart\versepar \ledsidenote{\passref{4.36.40}}\skt{प्रवाहिका }\edtext{\skt{भवेत्तीक्ष्णा}}{\lemma{\skt{भवेत्तीक्ष्णा}}\Afootnote{K H; \skt{भवेत्तीक्ष्णान्}\ A}}\skt{ निरू\-हा\-त्सा\-नु\-वा\-स\-नात्~॥ }\pend
\pstart\versepar \skt{सदाहशूलंं कृच्छ्रेण }\edtext{\skt{कफ\-न्त\-त्रो\-प\-वे\-श्यते}}{\lemma{\skt{कफ\appbrk{}न्त\appbrk{}त्रो\appbrk{}प\appbrk{}वे\appbrk{}श्यते}}\Afootnote{K H; \skt{कफा\appbrk{}सृ\appbrk{}गु\appbrk{}प\appbrk{}वे\appbrk{}श्यते~॥}\ A}}\skt{~।}\edtext{}{\lemma{\skt{निरू\appbrk{}हा\appbrk{}त्सा\appbrk{}नु\appbrk{}वा\appbrk{}स\appbrk{}नात्~॥ सदाहशूलंं}}\Afootnote{H; \skt{निरू\appbrk{}हा\appbrk{}त्सा\appbrk{}नु\appbrk{}वा\appbrk{}स\appbrk{}ना\appbrk{}त्स\appbrk{}दा\appbrk{}ह\appbrk{}शूलं}\ A K}}\pend

\pstart\versepar \ledsidenote{\passref{4.36.41}}\edtext{\edtext{\skt{पिच्छा\-ब\-स्ति\-र्हि\-त\-स्तत्र}}{\lemma{\skt{पिच्छाबस्तिर्}}\Afootnote{A H; \skt{पिच्छाबस्ति}\ K}}}{\lemma{\skt{हितस्}}\Afootnote{A H; \skt{र्हितस्}\ K}}\skt{ पयसा चैव भोजनं~॥ }\pend
\pstart\versepar \skt{}\edtext{\edtext{\skt{सर्प्पि\-र्म\-धु\-क\-सि\-द्धञ्च}}{\lemma{\skt{सर्प्पिर्}}\Afootnote{K H; \skt{सर्पि}\ A}}}{\lemma{\skt{मधुकसिद्धञ्च}}\Afootnote{K H; \skt{र्मधुरकैः सिद्धं}\ A}}\skt{ }\edtext{\skt{तैल\-म्वा\-प्य\-नु\-वा\-सनं}}{\lemma{\skt{वाप्य्}}\Afootnote{K H; \skt{चाप्य्}\ A}}\skt{~।}\pend

\pstart\versepar \ledsidenote{\passref{4.36.42}}\skt{अतितीक्ष्णो निरूहो वा }\edtext{\edtext{\skt{सत\-त\-म्वा\-नु\-वा\-सनं}}{\lemma{\skt{सततम्}}\Afootnote{K H; \skt{सवाते}\ A}}}{\lemma{\skt{वानुवासनं}}\Afootnote{K H; \skt{चानुवासनः~।}\ A}}\skt{~॥ }\pend
\pstart\versepar \skt{हृद\-य\-स्यो\-प\-स\-रणं }\edtext{\skt{कुरु\-तेऽ\-था\-ङ्ग\-सा\-दनं}}{\lemma{\skt{ऽथाङ्गसादनं}}\Afootnote{H; \skt{वा सवाते चानुवासनः~॥}\ A; \skt{थाङ्गसादनं~।}\ K}}\skt{~।}\pend

\pstart\versepar \ledsidenote{\passref{4.36.43}}\skt{दोषैस्तत्र }\edtext{\skt{रुजास्तीव्रा}}{\lemma{\skt{रुजास्तीव्रा}}\Afootnote{K H; \skt{रुजस्तास्ता}\ A}}\skt{ मदो }\edtext{\skt{मूर्च्छां}}{\lemma{\skt{मूर्च्छां}}\Afootnote{K H; \skt{मूर्च्छा}\ A}}\skt{ }\edtext{\skt{सगौरवं}}{\lemma{\skt{सगौरवं}}\Afootnote{K H; \skt{ऽङ्गगौरवम्~।}\ A}}\skt{~॥ }\pend
\pstart\versepar \skt{सर्व्वदोषहरं बस्तिं शोधनं तत्र दापयेत्~।}\pend

\pstart\versepar \ledsidenote{\passref{4.36.44}}\edtext{\skt{रूक्षस्य}}{\lemma{\skt{रूक्षस्य}}\Afootnote{K H; \skt{रूक्सस्य}\ A}}\skt{ बहुवातस्य }\edtext{\skt{स्याद्दुःखं}}{\lemma{\skt{स्याद्दुःखं}}\Afootnote{K H; \skt{तथा दुः}\ A}}\skt{ }\edtext{\skt{शयितस्य वा}}{\lemma{\skt{शयितस्य वा}}\Afootnote{K H; \skt{शायितस्य च~।}\ A}}\skt{~॥ }\pend
\pstart\versepar \skt{}\edtext{\edtext{\skt{बस्ति\-प्र\-ग्र\-ह\-म\-ङ्गानां}}{\lemma{\skt{बस्तिप्रग्रहम्}}\Afootnote{K H; \skt{बस्तिरङ्गग्र}\ A}}}{\lemma{\skt{अङ्गानां}}\Afootnote{K H; \skt{हं}\ A}}\skt{ }\edtext{\skt{कुर्याद्रुक्षो}}{\lemma{\skt{कुर्याद्रुक्षो}}\Afootnote{H; \skt{कुर्याद्रूक्सो}\ A; \skt{कुर्याद्रु}\ K}}\skt{ }\edtext{\skt{ल्पभेषजाः}}{\lemma{\skt{ल्पभेषजाः}}\Afootnote{H; \skt{मृद्वल्पभेषजः~॥}\ A; \skt{क्षोल्पभेषजाः~।}\ K}}\skt{~।}\pend

\pstart\versepar \ledsidenote{\passref{4.36.45}}\edtext{\skt{तत्रा\-ङ्ग\-सा\-द\-प्र\-स्तम्भ}}{\lemma{\skt{तत्रा\appbrk{}ङ्ग\appbrk{}सा\appbrk{}द\appbrk{}प्र\appbrk{}स्तम्भ}}\Afootnote{K H; \skt{तत्राङ्गसादः प्रस्तम्भो}\ A}}\skt{ }\edtext{\skt{जृम्भो\-द्वे\-ष्ट\-न\-वे\-पकाः}}{\lemma{\skt{जृम्भो\appbrk{}द्वे\appbrk{}ष्ट\appbrk{}न\appbrk{}वे\appbrk{}पकाः}}\Afootnote{A H; \skt{जृम्भो\appbrk{}द्वे\appbrk{}ष्ट\appbrk{}वे\appbrk{}पकाः~।}\ K}}\skt{~॥ }\pend
\pstart\versepar \skt{}\edtext{\skt{सर्व्वभेदश्च}}{\lemma{\skt{सर्व्वभेदश्च}}\Afootnote{K H; \skt{पर्वभेदश्च}\ A}}\skt{ तत्रेष्टाः स्वेदा\-भ्य\-ञ्ज\-न\-ब\-स्तयः~।}\pend

\pstart\versepar \ledsidenote{\passref{4.36.46}}\edtext{\edtext{\edtext{\skt{अत्यु\-ष्ण\-ती\-क्ष्णोऽ\-ति\-ब\-हु\-र्द्द\-त्ता\-ति\-स्वे\-दि\-तस्य}}{\lemma{\skt{अत्युष्णतीक्ष्णो}}\Afootnote{A H; \skt{अत्युष्णतीक्ष्णो}\ K}}}{\lemma{\skt{ऽतिबहुर्}}\Afootnote{H; \skt{ऽतिबहु}\ A; \skt{तिबर्हु}\ K}}\skt{ वा}}{\lemma{\skt{द्दत्ता\appbrk{}ति\appbrk{}स्वे\appbrk{}दि\appbrk{}तस्य वा}}\Afootnote{K H; \skt{र्दत्तोऽ\appbrk{}ति\appbrk{}स्वे\appbrk{}दि\appbrk{}तस्य च~।}\ A}}\skt{~॥ }\pend
\pstart\versepar \skt{अल्पदोषस्य वा बस्ति\-र\-ति\-यो\-गाय कल्पते~।}\pend

\pstart\versepar \ledsidenote{\passref{4.36.47}}\skt{विरे\-च\-ना\-ति\-यो\-गेन }\edtext{\skt{समानन्तत्र}}{\lemma{\skt{तत्र}}\Afootnote{K H; \skt{तस्य}\ A}}\skt{ लक्षणं~॥ }\pend
\pstart\versepar \skt{}\edtext{\skt{पिच्छाबस्तिः}}{\lemma{\skt{पिच्छाबस्तिः}}\Afootnote{K H; \skt{पिच्छाबस्ति}\ A}}\skt{ }\edtext{\edtext{\skt{समाभ्यासस्तस्य}}{\lemma{\skt{समाभ्यासस्}}\Afootnote{K H; \skt{प्रयोगश्च}\ A}}}{\lemma{\skt{तस्य}}\Afootnote{K H; \skt{तत्र}\ A}}\skt{ शीतः सुखावहः~।}\pend

\pstart\versepar \ledsidenote{\passref{4.36.48}}\skt{अतियोगात्परं }\edtext{\skt{हन्ति}}{\lemma{\skt{हन्ति}}\Afootnote{K H; \skt{यत्र}\ A}}\skt{ जीवा\-दा\-न\-म्वि\-रि\-क्त\-वत्~॥ }\pend
\pstart\versepar \skt{देहं }\edtext{\skt{देयांस्तत्र}}{\lemma{\skt{देयांस्}}\Afootnote{H\textsuperscript{pc}; \emph{om.}\ A; \skt{यांस्}\ K H\textsuperscript{ac}}}\skt{ }\edtext{\skt{हितानाहुः}}{\lemma{\skt{हितानाहुः}}\Afootnote{K H; \skt{हितश्चाप्सु}\ A}}\skt{ }\edtext{\skt{पिच्छाबस्तिं}}{\lemma{\skt{पिच्छाबस्तिं}}\Afootnote{H; \skt{पिच्छाबस्तिः}\ A; \skt{पिच्छाबस्ति}\ K}}\skt{ }\edtext{\skt{सशोणितं}}{\lemma{\skt{सशोणितं}}\Afootnote{K H; \skt{सशोणितः~॥}\ A}}\skt{~।}\edtext{}{\lemma{\skt{विरिक्तवत्~॥ देहं}}\Afootnote{H\textsuperscript{pc}; \skt{विरिक्तवत्~। देयस्}\ A; \skt{विरिक्तवत्दे}\ K H\textsuperscript{ac}}}\pend

\pstart\versepar \ledsidenote{\passref{4.36.51}}\edtext{\skt{पक्षाद्विरेको}}{\lemma{\skt{पक्षाद्विरेको}}\Afootnote{H; \skt{पक्साद्विरेको}\ A; \skt{पक्षाद्विरेको}\ K}}\skt{ वान्तस्य }\edtext{\edtext{\skt{निरू\-ह\-ह्य\-श्चा\-नु\-वा\-सितः}}{\lemma{\skt{निरूहह्यश्}}\Afootnote{H\textsuperscript{pc}; \skt{ततश्चापि निरूह}\ A; \skt{निरूह्यश्}\ K; \skt{निरूहश्}\ H\textsuperscript{ac}}}}{\lemma{\skt{चानुवासितः}}\Afootnote{K H; \skt{णम्~।}\ A}}\skt{~॥ }\pend
\pstart\versepar \skt{सद्यो }\edtext{\edtext{\skt{निरु\-होऽ\-नु\-वा\-सस्य}}{\lemma{\skt{निरुहो}}\Afootnote{H; \skt{निरूढो}\ A; \skt{निरूहो}\ K}}}{\lemma{\skt{ऽनुवासस्य}}\Afootnote{H\textsuperscript{pc}; \skt{ऽनुवास्यः}\ A; \skt{नुवास्य}\ K; \skt{ऽनुवास्य}\ H\textsuperscript{ac}}}\skt{ }\edtext{\skt{सप्त\-रा\-त्रा\-द्वि\-रे\-चित}}{\lemma{\skt{सप्त\appbrk{}रा\appbrk{}त्रा\appbrk{}द्वि\appbrk{}रे\appbrk{}चित}}\Afootnote{H; \skt{सप्त\appbrk{}रा\appbrk{}त्रा\appbrk{}द्वि\appbrk{}रे\appbrk{}चितः~॥}\ A; \skt{सप्त\appbrk{}रा\appbrk{}त्रा\appbrk{}द्वि\appbrk{}रे\appbrk{}चित}\ K}}\skt{ }\edtext{\skt{इति}}{\lemma{\skt{इति}}\Afootnote{H; \emph{om.}\ A K}}\skt{~॥}\pend

\pstart \ledsidenote{\passref{4.36.trailer}}\edtext{\skt{॥ ३६~॥}}{\lemma{\skt{॥ ३६~॥}}\Afootnote{\skt{इति सुश्रु\appbrk{}त\appbrk{}सं\appbrk{}हि\appbrk{}तायां चिकित्सास्थाने नेत्र\appbrk{}ब\appbrk{}स्ति\appbrk{}व्या\appbrk{}प\appbrk{}च्चि\appbrk{}कि\appbrk{}त्सितं नाम षट्त्रिंशोऽध्यायः~॥३६~॥}\ A; \skt{इति~॥ चि लह्र॥ {\symfont ❈}~॥}\ K; \skt{ओ~॥ चिकित्सा लम्त्र॥}\ H}}\pend

\endnumbering
\beginnumbering
\eledsection*{Adhyāya 37}
\pstart \ledsidenote{\passref{4.37.1}}\edtext{\skt{अथा\-तोऽ\-नु\-वा\-स\-नो\-त्त\-र\-ब\-स्ति\-चि\-कि\-त्सितं}}{\lemma{\skt{ऽनु\appbrk{}वा\appbrk{}स\appbrk{}नो\appbrk{}त्त\appbrk{}र\appbrk{}ब\appbrk{}स्ति\appbrk{}चि\appbrk{}कि\appbrk{}त्सितं}}\Afootnote{A H; \skt{नुवा\appbrk{}स\appbrk{}नो\appbrk{}त्त\appbrk{}र\appbrk{}ब\appbrk{}स्ति\appbrk{}चि\appbrk{}कि\appbrk{}त्सितं}\ K}}\skt{ व्याख्यास्यामः~॥}\pend

\pstart\versepar \ledsidenote{\passref{4.37.3}}\skt{विरे\-च\-ना\-त्स\-प्त\-रात्रे गते जातबलाय वै~। }\pend
\pstart\versepar \skt{}\edtext{\skt{कृताहाराय सा}}{\lemma{\skt{कृताहाराय सा}}\Afootnote{K H; \skt{कृतान्नायानुवा}\ A}}\skt{ }\edtext{\skt{याह्ने}}{\lemma{\skt{याह्ने}}\Afootnote{K H; \skt{स्याय}\ A}}\skt{ }\edtext{\edtext{\skt{बस्तिर्द्देयो}}{\lemma{\skt{बस्तिर्}}\Afootnote{K H; \skt{सम्य}\ A}}}{\lemma{\skt{द्देयो}}\Afootnote{H; \skt{ग्देयो}\ A; \skt{दे}\ K}}\skt{ }\edtext{\skt{नुवासनः}}{\lemma{\skt{नुवासनः}}\Afootnote{H; \skt{ऽनुवासनः~॥}\ A; \skt{योनुवासनः~।}\ K}}\skt{~॥}\pend

\pstart\versepar \ledsidenote{\passref{4.37.4}}\skt{यथावयो निरूहाणां या }\edtext{\skt{मात्रा\-स\-म्प्र\-की\-र्त्तिताः}}{\lemma{\skt{मात्रा\appbrk{}स\appbrk{}म्प्र\appbrk{}की\appbrk{}र्त्तिताः}}\Afootnote{K H; \skt{मात्राः परिकीर्तिताः~।}\ A}}\skt{~॥ }\pend
\pstart\versepar \skt{}\edtext{\skt{पादापकृष्टास्ताः}}{\lemma{\skt{पादापकृष्टास्ताः}}\Afootnote{K H; \skt{पादावकृष्टास्ताः}\ A}}\skt{ }\edtext{\skt{कार्याः}}{\lemma{\skt{कार्याः}}\Afootnote{A H; \skt{कार्या}\ K}}\skt{ स्नेहबस्तिषु }\edtext{\skt{देदेहिनां}}{\lemma{\skt{देदेहिनां}}\Afootnote{H; \skt{देहिनाम्}\ A K}}\skt{~॥}\pend

\pstart\versepar \ledsidenote{\passref{4.37.5}}\skt{उत्सृ\-ष्टा\-नि\-ल\-वि\-ण्मूत्रे नरे बस्ति\-न्नि\-धा\-प\-येत्~। }\pend
\pstart\versepar \skt{एतैर्हि }\edtext{\skt{विहितो}}{\lemma{\skt{विहितो}}\Afootnote{H; \skt{विहतः}\ A; \skt{विर्हितो}\ K}}\skt{ }\edtext{\edtext{\skt{बस्तिर्न्नैव}}{\lemma{\skt{बस्तिर्}}\Afootnote{H; \skt{स्नेहो}\ A; \skt{बस्ति}\ K}}}{\lemma{\skt{न्नैव}}\Afootnote{K H; \skt{नै}\ A}}\skt{ }\edtext{\skt{चान्तः}}{\lemma{\skt{चान्तः}}\Afootnote{K H; \skt{वान्तः}\ A}}\skt{ }\edtext{\skt{प्रपद्यते}}{\lemma{\skt{प्रपद्यते}}\Afootnote{K H; \skt{प्रतिपद्यते~॥}\ A}}\skt{~॥}\edtext{}{\lemma{\skt{निधापयेत्~। एतैर्हि}}\Afootnote{K H; \skt{विधापयेत्~। एतैर्हि}\ A}}\pend

\pstart\versepar \ledsidenote{\passref{4.37.6}}\edtext{\edtext{\skt{स्नेह\-ब\-स्ति\-र्वि\-धे\-यश्च}}{\lemma{\skt{स्नेहबस्तिर्}}\Afootnote{K H; \skt{स्नेहवस्तिर्}\ A}}}{\lemma{\skt{विधेयश्च}}\Afootnote{K H; \skt{विधेयस्तु}\ A}}\skt{ नाविशुद्धस्य देहिनः~। }\pend
\pstart\versepar \skt{स्नेहवीर्यं }\edtext{\skt{यथा}}{\lemma{\skt{यथा}}\Afootnote{K H; \skt{तथा}\ A}}\skt{ दत्ते }\edtext{\skt{देहे\-ना\-नु\-वि\-स\-र्प्पति}}{\lemma{\skt{देहे\appbrk{}ना\appbrk{}नु\appbrk{}वि\appbrk{}स\appbrk{}र्प्पति}}\Afootnote{K H; \skt{देहं चानुविसर्पति~॥}\ A}}\skt{~॥}\pend

\pstart\versepar \ledsidenote{\passref{4.37.43}}\skt{अशुद्धमपि वातेन केव\-ले\-ना\-ति\-पी\-डितं~। }\pend
\pstart\versepar \skt{अहोरात्रस्य कालेषु }\edtext{\edtext{\skt{सर्व्वे\-ष्वे\-वा\-नु\-वा\-स\-येत्}}{\lemma{\skt{सर्व्वेष्व्}}\Afootnote{H; \skt{सर्वे}\ A K}}}{\lemma{\skt{एवा\appbrk{}नु\appbrk{}वा\appbrk{}स\appbrk{}येत्}}\Afootnote{H; \skt{ष्वेवा\appbrk{}नु\appbrk{}वा\appbrk{}स\appbrk{}येत्~॥}\ A K}}\skt{~॥}\pend

\pstart\versepar \ledsidenote{\passref{4.37.44}}\skt{रूक्षस्य }\edtext{\skt{बहुवातस्यद्वौ}}{\lemma{\skt{बहुवातस्यद्वौ}}\Afootnote{K H; \skt{बहूवातस्य द्वौ}\ A}}\skt{ }\edtext{\skt{त्रीन्वा\-प्य\-नु\-वा\-स\-नात्}}{\lemma{\skt{त्रीन्वाप्य्}}\Afootnote{K H; \skt{त्रीन}\ A}}\skt{~। }\pend
\pstart\versepar \skt{दत्वा }\edtext{\skt{स्निग्धतनुं}}{\lemma{\skt{स्निग्धतनुं}}\Afootnote{A H; \skt{स्निगधतनुं}\ K}}\skt{ ज्ञात्वा ततः }\edtext{\skt{पश्चा\-न्नि\-रू\-ह\-येत्}}{\lemma{\skt{पश्चान्}}\Afootnote{H; \skt{पश्चान्}\ A K}}\skt{~॥}\edtext{}{\lemma{\skt{अनुवासनात्~। दत्वा}}\Afootnote{K H; \skt{प्यनुवासनान्~। दत्त्वा}\ A}}\pend

\pstart\versepar \ledsidenote{\passref{4.37.45}}\skt{अस्निग्धमपि वातेन केव\-ले\-ना\-ति\-पी\-डितं~। }\pend
\pstart\versepar \skt{स्नेह\-प्र\-गा\-ढै\-र्म्म\-ति\-मा\-न्नि\-रूहैः समुपाचरेत्~॥}\pend

\pstart\versepar \ledsidenote{\passref{4.37.46}}\skt{अथ }\edtext{\skt{सम्यङ्निरूढञ्च}}{\lemma{\skt{निरूढञ्च}}\Afootnote{H; \skt{गिरूढं तु}\ A; \skt{निरूढश्च}\ K}}\skt{ }\edtext{\skt{वाता\-दि\-ष्व\-नु\-वा\-स\-नात्}}{\lemma{\skt{वातादिष्व्}}\Afootnote{H; \skt{वातादि}\ A K}}\skt{~। }\pend
\pstart\versepar \skt{}\edtext{\skt{बिल्व\-य\-ष्ट्या\-ह्व\-म\-द\-न\-फ\-ल\-तै\-लै\-र्य\-था\-क्रमं}}{\lemma{\skt{यथाक्रमं}}\Afootnote{A H; \skt{यथाक्रमां~।}\ K}}\skt{~॥}\edtext{}{\lemma{\skt{अनुवासनात्~। बिल्व\appbrk{}य\appbrk{}ष्ट्या\appbrk{}ह्व\appbrk{}म\appbrk{}द\appbrk{}न\appbrk{}फ\appbrk{}ल\appbrk{}तै\appbrk{}लैर्}}\Afootnote{H; \skt{ष्वनुवासयेत्~। बिल्व\appbrk{}य\appbrk{}ष्ट्या\appbrk{}ह्व\appbrk{}म\appbrk{}द\appbrk{}न\appbrk{}फ\appbrk{}ल\appbrk{}तै\appbrk{}लैर्}\ A; \skt{ष्वनुवासनात्~। बिल्व\appbrk{}य\appbrk{}ष्ट्या\appbrk{}ह्व\appbrk{}म\appbrk{}द\appbrk{}न\appbrk{}फ\appbrk{}ल\appbrk{}तै\appbrk{}लैर्}\ K}}\pend

\pstart\versepar \ledsidenote{\passref{4.37.47}}\skt{रात्रौ }\edtext{\skt{बस्तिर्न्न}}{\lemma{\skt{बस्तिर्न्न}}\Afootnote{K H; \skt{बस्तिं न}\ A}}\skt{ }\edtext{\skt{दद्याच्च}}{\lemma{\skt{दद्याच्च}}\Afootnote{K H; \skt{दद्यात्तु}\ A}}\skt{ दोषोत्क्लेशो हि }\edtext{\skt{रात्रिजः}}{\lemma{\skt{रात्रिजः}}\Afootnote{A H; \skt{रात्रितजः~।}\ K}}\skt{~। }\pend
\pstart\versepar \skt{स्नेहवीर्ययुतः कुर्यादाध्मानं गौरवं ज्वरं~॥}\pend

\pstart\versepar \ledsidenote{\passref{4.37.48}}\edtext{\skt{अग्निस्थानगते}}{\lemma{\skt{अग्निस्थानगते}}\Afootnote{H; \skt{अह्नि स्थानस्थिते}\ A; \skt{अह्निस्थानगते}\ K}}\skt{ दोषे वह्नौ }\edtext{\skt{चानुरसेन्विते}}{\lemma{\skt{चानुरसेन्विते}}\Afootnote{H; \skt{चान्नरसान्विते~।}\ A; \skt{च्चानुरसेन्विते~।}\ K}}\skt{~। }\pend
\pstart\versepar \skt{}\edtext{\skt{स्फुट\-स्रो\-तो\-मु\-ख\-न्देहं}}{\lemma{\skt{स्फुट\appbrk{}स्रो\appbrk{}तो\appbrk{}मु\appbrk{}ख\appbrk{}न्देहं}}\Afootnote{K H; \skt{स्फुटस्रोतोमुखे देहे}\ A}}\skt{ स्नेहौजः परिसर्प्पति~॥}\pend

\pstart\versepar \ledsidenote{\passref{4.37.49}}\edtext{\skt{पित्तेऽधिके}}{\lemma{\skt{ऽधिके}}\Afootnote{A H; \skt{धिके}\ K}}\skt{ कफे क्षीणे रूक्षे }\edtext{\skt{वातरुजार्दिते}}{\lemma{\skt{वातरुजार्दिते}}\Afootnote{K H; \skt{वातरुगर्दिते~।}\ A}}\skt{~। }\pend
\pstart\versepar \skt{}\edtext{\edtext{\skt{रात्रावपि}}{\lemma{\skt{रात्राव्}}\Afootnote{H; \skt{नरे रात्रौ}\ A; \skt{रात्रा}\ K}}\skt{ तु}}{\lemma{\skt{अपि तु}}\Afootnote{H; \skt{तु}\ A; \skt{वपि तु}\ K}}\skt{ दातव्यं काले }\edtext{\skt{चोष्णेऽनुवासनं}}{\lemma{\skt{ऽनुवासनं}}\Afootnote{A H; \skt{नुवासनं~।}\ K}}\skt{~॥}\pend

\pstart\versepar \ledsidenote{\passref{4.37.50}}\skt{उष्णे }\edtext{\skt{पित्तेधिकेवापि}}{\lemma{\skt{पित्तेधिकेवापि}}\Afootnote{K H; \skt{पित्ताधिके वाऽपि}\ A}}\skt{ दिवादाहादयो गदाः~। }\pend
\pstart\versepar \skt{}\edtext{\skt{सम्भवन्ति}}{\lemma{\skt{सम्भवन्ति}}\Afootnote{A H; \skt{सम्भवत्ति}\ K}}\skt{ }\edtext{\skt{तदा}}{\lemma{\skt{तदा}}\Afootnote{K H; \skt{यतस्}\ A}}\skt{ }\edtext{\edtext{\skt{स्त्वेनम्प्रदोषे}}{\lemma{\skt{स्त्वेनम्}}\Afootnote{K H; \skt{तस्मा}\ A}}}{\lemma{\skt{प्रदोषे}}\Afootnote{K H; \skt{त्प्रदोषे}\ A}}\skt{ योजयेद्भिषक्~॥}\pend

\pstart\versepar \ledsidenote{\passref{4.37.52}}\skt{अहोरात्रस्य कालेषु }\edtext{\skt{सर्व्वे\-ष्वे\-वा\-नि\-लो\-च्छ्र\-यात्}}{\lemma{\skt{सर्व्वेष्व्}}\Afootnote{H; \skt{सर्वे}\ A K}}\skt{~। }\pend
\pstart\versepar \skt{तीव्रायां }\edtext{\edtext{\skt{रुचि\-जी\-र्ण्णा\-न्न\-म्भो\-ज\-यि\-त्वा\-नु\-वा\-स\-येत्}}{\lemma{\skt{रुचिजीर्ण्णान्नम्}}\Afootnote{H; \skt{रुजि जीर्णान्नं}\ A K}}}{\lemma{\skt{भोज\appbrk{}यि\appbrk{}त्वा\appbrk{}नु\appbrk{}वा\appbrk{}स\appbrk{}येत्}}\Afootnote{K H; \skt{भोज\appbrk{}यि\appbrk{}त्वाऽ\appbrk{}नु\appbrk{}वा\appbrk{}स\appbrk{}येत्~॥}\ A}}\skt{~॥}\edtext{}{\lemma{\skt{एवा\appbrk{}नि\appbrk{}लो\appbrk{}च्छ्र\appbrk{}यात्~। तीव्रायां}}\Afootnote{H; \skt{ष्वेवा\appbrk{}नि\appbrk{}ला\appbrk{}धि\appbrk{}कम्~। तीव्रायां}\ A; \skt{ष्वेवा\appbrk{}नि\appbrk{}लो\appbrk{}च्छ्र\appbrk{}यात्~। तीव्रायां}\ K}}\pend

\pstart\versepar \ledsidenote{\passref{4.37.53}}\skt{न चाभुक्तवतः स्नेहः प्रणिधेयः }\edtext{\skt{कथञ्चन}}{\lemma{\skt{कथञ्चन}}\Afootnote{A H; \skt{कथञ्चनः~।}\ K}}\skt{~। }\pend
\pstart\versepar \skt{}\edtext{\skt{सूक्ष्म\-त्वा\-च्छू\-न्य\-को\-ष्ठस्य}}{\lemma{\skt{सूक्ष्म\appbrk{}त्वा\appbrk{}च्छू\appbrk{}न्य\appbrk{}को\appbrk{}ष्ठस्य}}\Afootnote{K H; \skt{शुद्ध\appbrk{}त्वा\appbrk{}च्छू\appbrk{}न्य\appbrk{}को\appbrk{}ष्ठस्य}\ A}}\skt{ }\edtext{\edtext{\edtext{\skt{क्षिप्र\-मू\-र्द्ध\-म\-थो\-त्प\-तेत्}}{\lemma{\skt{क्षिप्रम्}}\Afootnote{H; \skt{स्नेह}\ A; \skt{क्षिप्र}\ K}}}{\lemma{\skt{ऊर्द्धम्}}\Afootnote{H; \skt{ऊर्ध्वं}\ A; \skt{मूर्ध्व}\ K}}}{\lemma{\skt{अथोत्पतेत्}}\Afootnote{H; \skt{समुत्पतेत्~॥}\ A; \skt{मथोत्पतेत्~।}\ K}}\skt{~॥}\pend

\pstart\versepar \ledsidenote{\passref{4.37.54}}\edtext{\edtext{\edtext{\skt{सदा\-नु\-वा\-स\-ये\-द्भु\-क्त\-मा\-र्द्र\-पा\-णि\-न्न\-र\-म्भि\-षक्}}{\lemma{\skt{सदा\appbrk{}नु\appbrk{}वा\appbrk{}स\appbrk{}ये\appbrk{}द्भु\appbrk{}क्तम्}}\Afootnote{K H; \skt{सदाऽ\appbrk{}नु\appbrk{}वा\appbrk{}स\appbrk{}ये\appbrk{}च्चापि भोजयित्वाऽऽ}\ A}}}{\lemma{\skt{आर्द्रपाणिन्}}\Afootnote{K H; \skt{आर्द्रपाणि}\ A}}}{\lemma{\skt{नरम्}}\Afootnote{K H; \skt{नम्~।}\ A}}\skt{~। }\pend
\pstart\versepar \skt{ज्वर\-म्वि\-द\-ग्ध\-भु\-क्तस्य कुर्यात्स्नेहः प्रयोजितः~॥}\edtext{}{\lemma{\skt{भिषक्~। ज्वरम्}}\Afootnote{K H; \skt{ज्वरं}\ A}}\pend

\pstart\versepar \ledsidenote{\passref{4.37.55}}\edtext{\skt{न }\edtext{\skt{चाति\-स्नि\-ग्ध\-म\-शनं}}{\lemma{\skt{अशनं}}\Afootnote{H; \skt{मशनं}\ A K}}}{\lemma{\skt{न चातिस्निग्धम्}}\Afootnote{H; \skt{न चातिस्निग्ध}\ A K}}\skt{ भोज\-यि\-त्वा\-नु\-वा\-स\-येत्~। }\pend
\pstart\versepar \skt{मदं मूर्च्छाञ्च जनयेद्द्विधा स्नेहः }\edtext{\skt{प्रयोजितः}}{\lemma{\skt{प्रयोजितः}}\Afootnote{A H; \skt{प्रयोर्जितः~।}\ K}}\skt{~॥}\edtext{}{\lemma{\skt{भोज\appbrk{}यि\appbrk{}त्वा\appbrk{}नु\appbrk{}वा\appbrk{}स\appbrk{}येत्~। मदं}}\Afootnote{K H; \skt{भोज\appbrk{}यि\appbrk{}त्वाऽ\appbrk{}नु\appbrk{}वा\appbrk{}स\appbrk{}येत्~। मदं}\ A}}\pend

\pstart\versepar \ledsidenote{\passref{4.37.56}}\skt{रूक्षम्भुक्तवतो ह्यन्न\-म्ब\-ल\-म्व\-र्ण्णञ्च }\edtext{\skt{हीयते}}{\lemma{\skt{हीयते}}\Afootnote{K H; \skt{हापये}\ A}}\skt{~। }\pend
\pstart\versepar \skt{}\edtext{\skt{युक्तस्नेहवतो}}{\lemma{\skt{युक्तस्नेहवतो}}\Afootnote{K H; \skt{त्~। युक्तस्नेहमतो}\ A}}\skt{ }\edtext{\skt{जन्तु\-म्भो\-ज\-यि\-त्वा\-नु\-वा\-स\-येत्}}{\lemma{\skt{भोज\appbrk{}यि\appbrk{}त्वा\appbrk{}नु\appbrk{}वा\appbrk{}स\appbrk{}येत्}}\Afootnote{K H; \skt{भोज\appbrk{}यि\appbrk{}त्वाऽ\appbrk{}नु\appbrk{}वा\appbrk{}स\appbrk{}येत्~॥}\ A}}\skt{~॥}\pend

\pstart\versepar \ledsidenote{\passref{4.37.57}}\edtext{\edtext{\skt{यूष\-क्षी\-र\-र\-सै\-स्त\-स्मा\-द्य\-था\-व्या\-धि\-म\-वेक्ष्य}}{\lemma{\skt{तस्मा\appbrk{}द्य\appbrk{}था\appbrk{}व्या\appbrk{}धिम्}}\Afootnote{K H; \skt{तस्माद्यथाव्याधि}\ A}}\skt{ हि}}{\lemma{\skt{अवेक्ष्य हि}}\Afootnote{K H; \skt{समीक्ष्य वा~।}\ A}}\skt{~। }\pend
\pstart\versepar \skt{}\edtext{\skt{यथो\-चि\-ता\-त्पा\-द\-ही\-न\-म्भो\-ज\-यि\-त्वा\-नु\-वा\-स\-येत्}}{\lemma{\skt{भोज\appbrk{}यि\appbrk{}त्वा\appbrk{}नु\appbrk{}वा\appbrk{}स\appbrk{}येत्}}\Afootnote{K H; \skt{भोज\appbrk{}यि\appbrk{}त्वाऽ\appbrk{}नु\appbrk{}वा\appbrk{}स\appbrk{}येत्~॥}\ A}}\skt{~॥}\pend

\pstart\versepar \ledsidenote{\passref{4.37.57.1}}\skt{न तु }\edtext{\skt{भुक्तवतो}}{\lemma{\skt{भुक्तवतो}}\Afootnote{H; \skt{भुक्तवते}\ K}}\skt{ }\edtext{\skt{देय\-मा\-स्था\-प\-न\-मिति}}{\lemma{\skt{आस्थापनम्}}\Afootnote{H; \skt{आस्थापणम्}\ K}}\skt{ स्थितिः~। }\pend
\pstart\versepar \skt{विशूचिकां }\edtext{\skt{सञ्ज\-न\-ये\-च्च्छ\-र्द्दि\-म्वापि}}{\lemma{\skt{सञ्ज\appbrk{}न\appbrk{}ये\appbrk{}च्च्छ\appbrk{}र्द्दिम्}}\Afootnote{H; \skt{सञ्जनये च्छर्दिं}\ K}}\skt{ सुदारुणां~॥}\pend

\pstart\versepar \ledsidenote{\passref{4.37.57.2}}\skt{निःशेषाः सुखमायान्ति भोजनेना प्रपीडिताः~॥ }\pend
\pstart\versepar \skt{न चास्था\-प\-न\-वि\-क्षि\-प्त\-म\-न्न\-मग्निः प्रधावति~। }\pend
\pstart\versepar \skt{तस्मा\-दा\-स्था\-प\-न\-न्दे\-य\-न्नि\-रा\-हा\-राय जानता~॥}\pend

\pstart\versepar \ledsidenote{\passref{4.37.58}}\skt{अथानुवास्यं }\edtext{\skt{स्वभ्यक्तं}}{\lemma{\skt{स्वभ्यक्तं}}\Afootnote{K H; \skt{स्वभ्यक्त}\ A}}\skt{ }\edtext{\edtext{\skt{स्नानमुष्णाम्बुभिः}}{\lemma{\skt{स्नानम्}}\Afootnote{K H; \emph{om.}\ A}}}{\lemma{\skt{उष्णाम्बुभिः}}\Afootnote{K H; \skt{मुष्णाम्बुस्वेदितं}\ A}}\skt{ शनैः~। }\pend
\pstart\versepar \skt{भोजयित्वा }\edtext{\skt{यथोद्दिष्टं}}{\lemma{\skt{यथोद्दिष्टं}}\Afootnote{K H; \skt{यथाशास्त्रं}\ A}}\skt{ }\edtext{\edtext{\skt{कृत\-ञ्च\-ङ्क्र\-म\-ण\-न्ततः}}{\lemma{\skt{कृतञ्}}\Afootnote{H; \skt{कृत}\ A K}}}{\lemma{\skt{चङ्क्रमणन्}}\Afootnote{K H; \skt{चङ्क्रम्णं}\ A}}\skt{~॥}\pend

\pstart\versepar \ledsidenote{\passref{4.37.59}}\skt{विसृज्य च }\edtext{\skt{शकृन्मूत्रे}}{\lemma{\skt{शकृन्मूत्रे}}\Afootnote{K H; \skt{शकृन्मूत्रं}\ A}}\skt{ योज\-ये\-त्स्ने\-ह\-ब\-स्तिना~। }\pend
\pstart\versepar \skt{}\edtext{\skt{प्रणि\-धा\-न\-वि\-धा\-नन्तु}}{\lemma{\skt{प्रणि\appbrk{}धा\appbrk{}न\appbrk{}वि\appbrk{}धा\appbrk{}नन्तु}}\Afootnote{A H; \skt{प्रणि\appbrk{}धा\appbrk{}न\appbrk{}वि\appbrk{}धा\appbrk{}णन्तु}\ K}}\skt{ निरूहे सम्प्रवक्ष्यते~॥}\pend

\pstart\versepar \ledsidenote{\passref{4.37.60}}\skt{ततः }\edtext{\skt{प्रणिहितः}}{\lemma{\skt{प्रणिहितः}}\Afootnote{K H; \skt{प्रणिहित}\ A}}\skt{ }\edtext{\edtext{\skt{स्नेहस्तूत्तानो}}{\lemma{\skt{स्नेहस्}}\Afootnote{K H; \skt{स्नेह}\ A}}}{\lemma{\skt{तूत्तानो}}\Afootnote{K H; \skt{उत्तानो}\ A}}\skt{ वाक्छतम्भवेत्~। }\pend
\pstart\versepar \skt{प्रसारितैः सर्व्वगात्रैस्तथा वीर्यम्विसर्प्पति~॥}\pend

\pstart\versepar \ledsidenote{\passref{4.37.61}}\edtext{\skt{ताड\-ये\-त्त\-ल\-यो\-रेवं}}{\lemma{\skt{एवं}}\Afootnote{K H; \skt{एनं}\ A}}\skt{ }\edtext{\skt{त्रीस्त्रीन्वाराश्छनैः}}{\lemma{\skt{त्रीस्त्रीन्वाराश्छनैः}}\Afootnote{H; \skt{त्रींस्त्रीन्वाराञ्छनैः}\ A K}}\skt{ शनैः~। }\pend
\pstart\versepar \skt{}\edtext{\edtext{\skt{स्फिजोश्चैनन्ततः}}{\lemma{\skt{स्फिजोश्}}\Afootnote{H; \skt{स्फिचोश्}\ A; \skt{स्फिजो}\ K}}}{\lemma{\skt{चैनन्}}\Afootnote{A H; \skt{श्चनैनं}\ K}}\skt{ }\edtext{\skt{शय्या\-न्त्री\-न्वा\-रा\-नु\-त्क्षि\-पे\-त्पुनः}}{\lemma{\skt{त्रीन्वा\appbrk{}रा\appbrk{}नु\appbrk{}त्क्षि\appbrk{}पे\appbrk{}त्पुनः}}\Afootnote{H; \skt{त्रीन्वा\appbrk{}रा\appbrk{}नु\appbrk{}त्क्षि\appbrk{}पे\appbrk{}त्ततः~॥}\ A; \skt{त्रीन्वा\appbrk{}रा\appbrk{}नु\appbrk{}त्षि\appbrk{}पे\appbrk{}त्पुनः~।}\ K}}\skt{~॥}\pend

\pstart\versepar \ledsidenote{\passref{4.37.62}}\skt{एवम्प्रणिहिते बस्तौ }\edtext{\edtext{\skt{मन्दायामोऽथ}}{\lemma{\skt{मन्दायामो}}\Afootnote{H; \skt{मन्दायासो}\ A; \skt{मन्दायामो}\ K}}}{\lemma{\skt{ऽथ}}\Afootnote{A H\textsuperscript{pc}; \skt{थ}\ K; \skt{ऽ}\ H\textsuperscript{ac}}}\skt{ मन्दवाक्~। }\pend
\pstart\versepar \skt{स्वास्तीर्ण्णो शयने }\edtext{\edtext{\skt{काम\-मा\-सी\-ता\-चा\-रिके}}{\lemma{\skt{कामम्}}\Afootnote{K H; \skt{काम}\ A}}}{\lemma{\skt{आसीताचारिके}}\Afootnote{K H; \skt{मासीताचारिके}\ A}}\skt{ }\edtext{\skt{ततः}}{\lemma{\skt{ततः}}\Afootnote{K H; \skt{रतः~॥}\ A}}\skt{~॥}\edtext{}{\lemma{\skt{मन्दवाक्~। स्वास्तीर्ण्णो}}\Afootnote{K H; \skt{मन्दवाक्~। स्वास्तीर्णे}\ A}}\pend

\pstart\versepar \ledsidenote{\passref{4.37.63}}\edtext{\skt{तत्र}}{\lemma{\skt{तत्र}}\Afootnote{K H; \skt{स तु}\ A}}\skt{ }\edtext{\skt{सैन्धवचूर्ण्णेन}}{\lemma{\skt{सैन्धवचूर्ण्णेन}}\Afootnote{A H; \skt{सेन्धवचूर्ण्णेन}\ K}}\skt{ शताह्वेन च }\edtext{\skt{संयुतः}}{\lemma{\skt{संयुतः}}\Afootnote{K H; \skt{योजितः~।}\ A}}\skt{~। }\pend
\pstart\versepar \skt{}\edtext{\skt{भवेत्सुखोष्णश्च}}{\lemma{\skt{भवेत्सुखोष्णश्च}}\Afootnote{K H; \skt{देयः सुखोष्णश्च}\ A}}\skt{ तथा }\edtext{\skt{निरति}}{\lemma{\skt{निरति}}\Afootnote{K H; \skt{निरेति}\ A}}\skt{ सहसा सुखं~॥}\pend

\pstart\versepar \ledsidenote{\passref{4.37.64}}\skt{यस्यानुवासनो दत्तः सकृ\-द\-न्व\-क्ष\-मा\-व्र\-जेत्~। }\pend
\pstart\versepar \skt{अत्यौ\-ष्णा\-द\-ति\-तै\-क्ष्णाद्वा }\edtext{\skt{वायुना वा}}{\lemma{\skt{वायुना वा}}\Afootnote{A H; \skt{वायुना वाऽ}\ K}}\skt{ }\edtext{\skt{प्रपीडिताः~॥ऽ}}{\lemma{\skt{प्रपीडिताः~॥ऽ}}\Afootnote{H; \skt{प्रपीडितः~॥}\ A; \skt{प्रपीडिताः~।}\ K}}\skt{~।}\edtext{}{\lemma{\skt{सकृ\appbrk{}द\appbrk{}न्व\appbrk{}क्ष\appbrk{}मा\appbrk{}व्र\appbrk{}जेत्~। अत्यौ\appbrk{}ष्णा\appbrk{}द\appbrk{}ति\appbrk{}तै\appbrk{}क्ष्णाद्वा}}\Afootnote{H; \skt{सकृ\appbrk{}द\appbrk{}न्व\appbrk{}क्ष\appbrk{}मा\appbrk{}व्र\appbrk{}जेत्~। अत्यौ\appbrk{}ष्ण्या\appbrk{}द\appbrk{}ति\appbrk{}तै\appbrk{}क्ष्ण्याद्वा}\ A; \skt{सकृदन्वक्षमा व्रजेत्~। अत्यौ\appbrk{}ष्णा\appbrk{}द\appbrk{}ति\appbrk{}तै\appbrk{}क्ष्ण्याद्वा}\ K}}\pend

\pstart\versepar \ledsidenote{\passref{4.37.65}}\edtext{\skt{सवा\-तो\-धि\-क\-मात्रो वा}}{\lemma{\skt{सवा\appbrk{}तो\appbrk{}धि\appbrk{}क\appbrk{}मात्रो वा}}\Afootnote{K H; \skt{सवा\appbrk{}तोऽ\appbrk{}धि\appbrk{}क\appbrk{}मात्रो वा}\ A}}\skt{ }\edtext{\skt{गुरु\-त्वा\-द्वाऽ\-ति\-भे\-षजः}}{\lemma{\skt{गुरु\appbrk{}त्वा\appbrk{}द्वाऽ\appbrk{}ति\appbrk{}भे\appbrk{}षजः}}\Afootnote{K H; \skt{गुरुत्वाद्वा सभेषजः~।}\ A}}\skt{~। }\pend
\pstart\versepar \skt{}\edtext{\skt{तस्यात्यल्पतरो}}{\lemma{\skt{तस्यात्यल्पतरो}}\Afootnote{H; \skt{तस्यान्योऽल्पतरो}\ A; \skt{तस्यात्यल्पतरो देयो न हि स्निह्य तिष्ठति~। तस्यात्यल्पतरो}\ K}}\skt{ देयो न हि }\edtext{\skt{स्निह्यति}}{\lemma{\skt{स्निह्यति}}\Afootnote{K H; \skt{स्निह्यत्य}\ A}}\skt{ तिष्ठति~॥}\pend

\pstart\versepar \ledsidenote{\passref{4.37.66}}\edtext{\skt{विष्ट\-ब्धा\-नि\-ल\-वि\-ण्मू\-त्र\-स्ने\-ह\-ही\-ने\-नु\-वा\-सने}}{\lemma{\skt{विष्ट\appbrk{}ब्धा\appbrk{}नि\appbrk{}ल\appbrk{}वि\appbrk{}ण्मू\appbrk{}त्र\appbrk{}स्ने\appbrk{}ह\appbrk{}ही\appbrk{}ने\appbrk{}नु\appbrk{}वा\appbrk{}सने}}\Afootnote{K H; \skt{विष्ट\appbrk{}ब्धा\appbrk{}नि\appbrk{}ल\appbrk{}वि\appbrk{}ण्मूत्रः स्नेह\appbrk{}ही\appbrk{}नेऽ\appbrk{}नु\appbrk{}वा\appbrk{}सने~।}\ A}}\skt{~। }\pend
\pstart\versepar \skt{}\edtext{\skt{दाहः}}{\lemma{\skt{दाहः}}\Afootnote{K H; \skt{दाह}\ A}}\skt{ }\edtext{\skt{क्लमः}}{\lemma{\skt{क्लमः}}\Afootnote{K H; \skt{क्लम}\ A}}\skt{ प्रवाहार्त्ति }\edtext{\skt{कर\-श्चा\-त्य\-नु\-वा\-सने}}{\lemma{\skt{चात्यनुवासने}}\Afootnote{K H; \skt{चात्यनुवासनः~॥}\ A}}\skt{~॥}\pend

\pstart\versepar \ledsidenote{\passref{4.37.67}}\edtext{\skt{सानिलं}}{\lemma{\skt{सानिलं}}\Afootnote{K H; \skt{सानिलः}\ A}}\skt{ }\edtext{\skt{सपुरीषञ्च}}{\lemma{\skt{सपुरीषञ्च}}\Afootnote{H; \skt{सपुरीषश्च}\ A; \skt{सपुरीशञ्च}\ K}}\skt{ }\edtext{\skt{स्नेहप्रत्येति}}{\lemma{\skt{स्नेहप्रत्येति}}\Afootnote{K H; \skt{स्नेहः प्रत्येति}\ A}}\skt{ यस्य तु~। }\pend
\pstart\versepar \skt{}\edtext{\skt{ओषाचोषौ}}{\lemma{\skt{ओषाचोषौ}}\Afootnote{K H; \skt{ओषचोषौ}\ A}}\skt{ }\edtext{\skt{विनाशीघ्रन्न}}{\lemma{\skt{विनाशीघ्रन्न}}\Afootnote{K H; \skt{विना शीघ्रं स}\ A}}\skt{ }\edtext{\skt{सम्य\-ग\-नु\-वा\-सितं}}{\lemma{\skt{सम्य\appbrk{}ग\appbrk{}नु\appbrk{}वा\appbrk{}सितं}}\Afootnote{K H; \skt{सम्य\appbrk{}ग\appbrk{}नु\appbrk{}वा\appbrk{}सितः~॥}\ A}}\skt{~॥}\pend

\pstart\versepar \ledsidenote{\passref{4.37.68}}\edtext{\edtext{\skt{जीर्ण्णान्नमथ}}{\lemma{\skt{जीर्ण्णान्नम्}}\Afootnote{K H; \skt{जीर्णान्न}\ A}}}{\lemma{\skt{अथ}}\Afootnote{K H; \skt{मथ}\ A}}\skt{ }\edtext{\skt{सायांह्ने}}{\lemma{\skt{सायांह्ने}}\Afootnote{H; \skt{सायाह्ने}\ A K}}\skt{ स्नेहे प्रत्यागते पुनः~। }\pend
\pstart\versepar \skt{}\edtext{\skt{लघ्व\-न्न\-म्भो\-ज\-ये\-त्का\-म\-न्दी\-प्ता\-ग्निन्तु}}{\lemma{\skt{दीप्ताग्निन्तु}}\Afootnote{H; \skt{दीप्ताग्निस्तु}\ A K}}\skt{ नरो यदि~॥}\pend

\pstart\versepar \ledsidenote{\passref{4.37.70}}\edtext{\edtext{\edtext{\edtext{\skt{स्नेह\-ब\-स्ति\-क्र\-म\-न्त्वे\-त\-दे\-व\-मा\-हु\-र्म्म\-नी\-षिणः}}{\lemma{\skt{स्नेह\appbrk{}ब\appbrk{}स्ति\appbrk{}क्र\appbrk{}मन्}}\Afootnote{K H; \skt{स्नेहबस्तिक्रमे}\ A}}}{\lemma{\skt{त्वेतदेवम्}}\Afootnote{K H; \skt{ष्वेवं विधि}\ A}}}{\lemma{\skt{आहुर्}}\Afootnote{K H; \skt{माहु}\ A}}}{\lemma{\skt{म्मनीषिणः}}\Afootnote{K H; \skt{र्मनीषिणः~।}\ A}}\skt{~। }\pend
\pstart\versepar \skt{अनेन विधिना }\edtext{\skt{यद्वा}}{\lemma{\skt{यद्वा}}\Afootnote{K H; \skt{षड्वा}\ A}}\skt{ सप्त }\edtext{\skt{वाष्टौ}}{\lemma{\skt{वाष्टौ}}\Afootnote{K H; \skt{वाऽष्टौ}\ A}}\skt{ }\edtext{\skt{परेण}}{\lemma{\skt{परेण}}\Afootnote{K H; \skt{नवैव}\ A}}\skt{ वा~॥}\pend

\pstart\versepar \ledsidenote{\passref{4.37.71}}\skt{विधेया }\edtext{\skt{बस्तयो}}{\lemma{\skt{बस्तयो}}\Afootnote{K H; \skt{बस्तय}\ A}}\skt{ }\edtext{\edtext{\skt{नृणामन्तरा}}{\lemma{\skt{नृणाम्}}\Afootnote{K H; \skt{स्तेषा}\ A}}\skt{ तु}}{\lemma{\skt{अन्तरा तु}}\Afootnote{K H; \skt{मन्तरा तु}\ A}}\skt{ निरूहणात्~। }\pend
\pstart\versepar \skt{दत्तस्तु प्रथमो }\edtext{\skt{बस्तिः}}{\lemma{\skt{बस्तिः}}\Afootnote{A H; \skt{ब}\ K\textsuperscript{ac}; \skt{बस्ति}\ K\textsuperscript{pc}}}\skt{ }\edtext{\skt{स्नेह\-ये\-द्ब\-स्ति\-वं\-क्षणौ}}{\lemma{\skt{स्नेह\appbrk{}ये\appbrk{}द्ब\appbrk{}स्ति\appbrk{}वं\appbrk{}क्षणौ}}\Afootnote{A K\textsuperscript{pc} H; \skt{स्ति वंक्षणौ~॥}\ K\textsuperscript{ac}}}\skt{~॥}\edtext{}{\lemma{\skt{निरूहणात्~। दत्तस्तु}}\Afootnote{K H; \skt{निरूहणम्~। दत्तस्तु}\ A}}\pend

\pstart\versepar \ledsidenote{\passref{4.37.72}}\edtext{\skt{सम्यग्दत्ते}}{\lemma{\skt{सम्यग्दत्ते}}\Afootnote{K H; \skt{सम्यग्दत्तो}\ A}}\skt{ द्वितीयस्तु }\edtext{\edtext{\skt{कोष्ठस्थमनिलं}}{\lemma{\skt{कोष्ठस्थम्}}\Afootnote{K H; \skt{मूर्धस्थ}\ A}}}{\lemma{\skt{अनिलं}}\Afootnote{K H; \skt{मनिलं}\ A}}\skt{ जयेत्~। }\pend
\pstart\versepar \skt{}\edtext{\skt{बलम्वर्ण्णञ्च}}{\lemma{\skt{वर्ण्णञ्च}}\Afootnote{K H; \skt{वर्णौ च}\ A}}\skt{ }\edtext{\skt{जन\-ये\-त्तृ\-ती\-यस्तु}}{\lemma{\skt{जन\appbrk{}ये\appbrk{}त्तृ\appbrk{}ती\appbrk{}यस्तु}}\Afootnote{K H; \skt{तृतीयस्तु}\ A}}\skt{ }\edtext{\skt{निषेवितः}}{\lemma{\skt{निषेवितः}}\Afootnote{K H; \skt{प्रयोजितः~॥}\ A}}\skt{~॥}\edtext{}{\lemma{\skt{जयेत्~। बलम्}}\Afootnote{K H; \skt{जयेत्~। जनयेद्बल}\ A}}\pend

\pstart\versepar \ledsidenote{\passref{4.37.73}}\edtext{\skt{चतुर्थः}}{\lemma{\skt{चतुर्थः}}\Afootnote{K H; \skt{रसं चतुर्थो}\ A}}\skt{ }\edtext{\skt{स्नेहनिरसं}}{\lemma{\skt{स्नेहनिरसं}}\Afootnote{H; \emph{om.}\ A; \skt{स्नेहतिरसं}\ K}}\skt{ रक्तं }\edtext{\skt{स्नेहति}}{\lemma{\skt{स्नेहति}}\Afootnote{K H; \skt{तु}\ A}}\skt{ पञ्चमः~। }\pend
\pstart\versepar \skt{}\edtext{\skt{षष्ठस्तु}}{\lemma{\skt{षष्ठस्तु}}\Afootnote{K H; \skt{स्नेहयेत्तथा~। षष्ठस्तु}\ A}}\skt{ }\edtext{\skt{स्नेहयेत्मांसं}}{\lemma{\skt{स्नेहयेत्मांसं}}\Afootnote{K H; \skt{स्नेहयेन्मांसं}\ A}}\skt{ मेदः }\edtext{\skt{स्नेहति}}{\lemma{\skt{स्नेहति}}\Afootnote{K H; \emph{om.}\ A}}\skt{ }\edtext{\skt{सप्तमः}}{\lemma{\skt{सप्तमः}}\Afootnote{K H; \skt{सप्तम एव च~॥}\ A}}\skt{~॥}\pend

\pstart\versepar \ledsidenote{\passref{4.37.74}}\skt{अष्टमो }\edtext{\skt{नवमश्चैव}}{\lemma{\skt{चैव}}\Afootnote{K H; \emph{om.}\ A}}\skt{ }\edtext{\skt{सार\-म्म\-ज्जा\-न\-मेव च}}{\lemma{\skt{सार\appbrk{}म्म\appbrk{}ज्जा\appbrk{}न\appbrk{}मेव च}}\Afootnote{K; \skt{चास्थि मज्जानं च यथाक्रमम्~।}\ A; \skt{सार\appbrk{}म्म\appbrk{}ज्जा\appbrk{}मे\appbrk{}वमे च}\ H\textsuperscript{ac}; \skt{सार\appbrk{}म्म\appbrk{}ज्जा\appbrk{}र\appbrk{}सा\appbrk{}म\appbrk{}ज्ज\appbrk{}न\appbrk{}मेव च}\ H\textsuperscript{pc}}}\skt{ }\pend
\pstart\versepar \skt{एवं }\edtext{\skt{शुक्र\-ग\-ता\-न्दो\-षा\-न्द्वि\-गु\-णे\-नाथ}}{\lemma{\skt{द्विगुणेनाथ}}\Afootnote{K H; \skt{द्विगुणः साधु}\ A}}\skt{ }\edtext{\skt{साधयेत्}}{\lemma{\skt{साधयेत्}}\Afootnote{A H; \skt{शोधयेत्~॥}\ K}}\skt{~॥}\pend

\pstart\versepar \ledsidenote{\passref{4.37.75}}\edtext{\skt{अष्टा\-द\-शा\-ष्टा\-द\-श\-का\-न्ब\-स्तीस्तु यो}}{\lemma{\skt{अष्टा\appbrk{}द\appbrk{}शा\appbrk{}ष्टा\appbrk{}द\appbrk{}श\appbrk{}का\appbrk{}न्ब\appbrk{}स्तीस्तु यो}}\Afootnote{H; \skt{अष्टा\appbrk{}द\appbrk{}शा\appbrk{}ष्टा\appbrk{}द\appbrk{}श\appbrk{}का\appbrk{}न्ब\appbrk{}स्तीनां यो}\ A; \skt{अष्टा\appbrk{}द\appbrk{}शा\appbrk{}ष्टा\appbrk{}द\appbrk{}शकां बस्तींस्तु यो}\ K}}\skt{ }\edtext{\skt{निसेवते}}{\lemma{\skt{निसेवते}}\Afootnote{K H; \skt{निषेवते~।}\ A}}\skt{~। }\pend
\pstart\versepar \skt{}\edtext{\skt{यथोक्तेन}}{\lemma{\skt{यथोक्तेन}}\Afootnote{K H; \skt{यथोकेन}\ A}}\skt{ }\edtext{\skt{विधानेन}}{\lemma{\skt{विधानेन}}\Afootnote{A H; \skt{विधाणेन}\ K}}\skt{ परि\-हा\-र\-क्र\-मेण च~॥}\pend

\pstart\versepar \ledsidenote{\passref{4.37.76}}\edtext{\edtext{\skt{सकु\-ञ्ज\-र\-ब\-लोऽ\-श्वस्य}}{\lemma{\skt{सकुञ्जरबलो}}\Afootnote{A H; \skt{सकुञ्जरबलो}\ K}}}{\lemma{\skt{ऽश्वस्य}}\Afootnote{A H; \skt{श्वस्य}\ K}}\skt{ }\edtext{\skt{जवे}}{\lemma{\skt{जवे}}\Afootnote{K H; \skt{जवैस्}\ A}}\skt{ }\edtext{\edtext{\skt{तुल्योऽमरप्रभः}}{\lemma{\skt{तुल्यो}}\Afootnote{A H; \skt{तुल्यो}\ K}}}{\lemma{\skt{ऽमरप्रभः}}\Afootnote{A H; \skt{मरप्रभः~।}\ K}}\skt{~। }\pend
\pstart\versepar \skt{वीत\-पा\-प्मा\-श्रु\-त\-धरः }\edtext{\edtext{\skt{सहस्रायुर्न्नरो}}{\lemma{\skt{सहस्रायुर्}}\Afootnote{H; \skt{सहस्रायु}\ A K}}}{\lemma{\skt{न्नरो}}\Afootnote{K H; \skt{र्नरो}\ A}}\skt{ भवेत्~॥}\pend

\pstart\versepar \ledsidenote{\passref{4.37.77}}\skt{स्नेह\-ब\-स्ति\-न्नि\-रू\-हम्वा }\edtext{\skt{नाभ्यसेदेक}}{\lemma{\skt{नाभ्यसेदेक}}\Afootnote{K H; \skt{नैक}\ A}}\skt{ }\edtext{\skt{एव तु}}{\lemma{\skt{एव तु}}\Afootnote{K H; \skt{मेवाति}\ A}}\skt{~। }\pend
\pstart\versepar \skt{}\edtext{\skt{स्नेहा\-त्पि\-त्त\-क\-फो\-त्क्लेशो}}{\lemma{\skt{स्नेहा\appbrk{}त्पि\appbrk{}त्त\appbrk{}क\appbrk{}फो\appbrk{}त्क्लेशो}}\Afootnote{H; \skt{शीलयेत्~। स्नेहा\appbrk{}द\appbrk{}ग्नि\appbrk{}व\appbrk{}धो\appbrk{}त्क्लेशौ}\ A; \skt{स्नेहा\appbrk{}त्पि\appbrk{}त्त\appbrk{}क\appbrk{}फो\appbrk{}त्केशो}\ K}}\skt{ निरू\-हा\-त्प\-व\-ना\-द्भयं~॥}\pend

\pstart\versepar \ledsidenote{\passref{4.37.78}}\edtext{\edtext{\skt{तस्मा\-न्नि\-रू\-ढोऽ\-नु\-वास्यो}}{\lemma{\skt{तस्मान्निरूढो}}\Afootnote{A H; \skt{तस्मान्निरूढो}\ K}}}{\lemma{\skt{ऽनुवास्यो}}\Afootnote{A H; \skt{नुवास्यो}\ K}}\skt{ }\edtext{\skt{निरू\-ढ\-श्चा\-नु\-वा\-सितः}}{\lemma{\skt{निरूढश्}}\Afootnote{K H; \skt{निरूह्यश्}\ A}}\skt{~। }\pend
\pstart\versepar \skt{}\edtext{\skt{नैव}}{\lemma{\skt{नैव}}\Afootnote{K H; \skt{नैवं}\ A}}\skt{ }\edtext{\skt{पित्तकफोत्क्लेशो}}{\lemma{\skt{पित्तकफोत्क्लेशो}}\Afootnote{H; \skt{पित्तकफोत्क्लेशौ}\ A K}}\skt{ स्यातान्न }\edtext{\skt{पव\-ना\-व\-ना\-द्भयं}}{\lemma{\skt{पव\appbrk{}ना\appbrk{}व\appbrk{}ना\appbrk{}द्भयं}}\Afootnote{H; \skt{पवनाद्भयम्}\ A K}}\skt{~॥}\pend

\pstart\versepar \ledsidenote{\passref{4.37.79}}\edtext{\skt{रूक्षस्य}}{\lemma{\skt{रूक्षस्य}}\Afootnote{K H; \skt{रूक्षाय}\ A}}\skt{ }\edtext{\skt{बहुवातस्य}}{\lemma{\skt{बहुवातस्य}}\Afootnote{K H; \skt{बहुवाताय}\ A}}\skt{ }\edtext{\skt{स्नेहबस्तिन्दिने}}{\lemma{\skt{स्नेहबस्तिन्}}\Afootnote{K H; \skt{स्नेहवस्तिं}\ A}}\skt{ दिने~। }\pend
\pstart\versepar \skt{}\edtext{\edtext{\edtext{\skt{दद्या\-द्वै\-द्य\-स्त\-तोऽ\-न्येषां}}{\lemma{\skt{दद्याद्वैद्यस्}}\Afootnote{A H; \skt{दद्याद्वैद्यस्}\ K}}}{\lemma{\skt{ततो}}\Afootnote{A H; \skt{ततो}\ K}}}{\lemma{\skt{ऽन्येषां}}\Afootnote{H; \skt{ऽन्येषा}\ A; \skt{न्येषाम्}\ K}}\skt{ }\edtext{\skt{मग्न्या\-बा\-ध\-भ\-या\-त्र्य\-हात्}}{\lemma{\skt{मग्न्या\appbrk{}बा\appbrk{}ध\appbrk{}भ\appbrk{}या\appbrk{}त्र्य\appbrk{}हात्}}\Afootnote{A H; \skt{अग\appbrk{}न्या\appbrk{}बा\appbrk{}ध\appbrk{}भ\appbrk{}या\appbrk{}त्र्य\appbrk{}हात्~॥}\ K}}\skt{~॥}\pend

\pstart\versepar \ledsidenote{\passref{4.37.80}}\edtext{\skt{स्नेहोल्पमात्रो}}{\lemma{\skt{स्नेहोल्पमात्रो}}\Afootnote{K H; \skt{स्नेहोऽल्पमात्रो}\ A}}\skt{ }\edtext{\skt{रूक्षा\-णा\-न्दी\-र्घ\-का\-ल\-म\-न\-न्ययः}}{\lemma{\skt{दीर्घ\appbrk{}का\appbrk{}ल\appbrk{}म\appbrk{}न\appbrk{}न्ययः}}\Afootnote{K H; \skt{दीर्घ\appbrk{}का\appbrk{}ल\appbrk{}म\appbrk{}न\appbrk{}त्ययः~।}\ A}}\skt{~। }\pend
\pstart\versepar \skt{तथा निरूहः }\edtext{\edtext{\skt{स्निग्धा\-ना\-म\-ल्प\-मात्रः}}{\lemma{\skt{स्निग्धानाम्}}\Afootnote{K H; \skt{स्निग्धाना}\ A}}}{\lemma{\skt{अल्पमात्रः}}\Afootnote{K H; \skt{मल्पमात्रः}\ A}}\skt{ प्रशस्यते~॥}\pend

\pstart\versepar \ledsidenote{\passref{4.37.81}}\skt{अत ऊर्द्ध्वम्प्रवक्ष्यामि }\edtext{\skt{व्यापदः}}{\lemma{\skt{व्यापदः}}\Afootnote{A H; \skt{व्यपदः}\ K}}\skt{ }\edtext{\skt{स्नेहबस्तिजः}}{\lemma{\skt{स्नेहबस्तिजः}}\Afootnote{K H; \skt{स्नेहबस्तिजाः~।}\ A}}\skt{~। }\pend
\pstart\versepar \skt{बलवन्तो यदा दोषाः }\edtext{\skt{कोष्ठाः}}{\lemma{\skt{कोष्ठाः}}\Afootnote{H; \skt{कोष्ठे}\ A; \skt{कोष्ठ}\ K}}\skt{ }\edtext{\skt{स्युरनिलादयः}}{\lemma{\skt{स्युरनिलादयः}}\Afootnote{K H; \skt{स्युरनिलादयः~॥}\ A}}\skt{~॥}\pend

\pstart\versepar \ledsidenote{\passref{4.37.82}}\skt{अल्पवीर्यन्तदा }\edtext{\skt{स्नेहामभिभूय}}{\lemma{\skt{स्नेहाम्}}\Afootnote{K H; \skt{स्नेहम्}\ A}}\skt{ }\edtext{\skt{पृथग्विधाः}}{\lemma{\skt{पृथग्विधाः}}\Afootnote{H; \skt{पृथग्विधा}\ A K}}\skt{~। }\pend
\pstart\versepar \skt{}\edtext{\skt{कुर्व्व\-न्त्यु\-प\-द्र\-वा\-न्स्नेहः स}}{\lemma{\skt{कुर्व्व\appbrk{}न्त्यु\appbrk{}प\appbrk{}द्र\appbrk{}वा\appbrk{}न्स्नेहः स}}\Afootnote{H; \skt{न्~। कुर्व\appbrk{}न्त्यु\appbrk{}प\appbrk{}द्र\appbrk{}वा\appbrk{}न्स्नेहः स}\ A; \skt{कुर्वन्त्युपद्रवां स्नेहः स}\ K}}\skt{ चापि न }\edtext{\skt{निवर्त्तते}}{\lemma{\skt{निवर्त्तते}}\Afootnote{A H; \skt{निवर्त्ते~॥}\ K}}\skt{~॥}\pend

\pstart\versepar \ledsidenote{\passref{4.37.83}}\skt{तत्र }\edtext{\skt{वाताभिभूते तु}}{\lemma{\skt{वाताभिभूते तु}}\Afootnote{A H; \skt{वाताविभूते तु}\ K}}\skt{ }\edtext{\skt{स्नेह\-व\-क्त्र\-क\-षा\-यता}}{\lemma{\skt{स्नेह\appbrk{}व\appbrk{}क्त्र\appbrk{}क\appbrk{}षा\appbrk{}यता}}\Afootnote{K H; \skt{स्नेहे मुखकषायता~।}\ A}}\skt{~। }\pend
\pstart\versepar \skt{}\edtext{\skt{जृम्भो}}{\lemma{\skt{जृम्भो}}\Afootnote{H; \skt{जृम्भा}\ A; \skt{जृम्भो}\ K}}\skt{ }\edtext{\skt{वातरुजस्तास्ता}}{\lemma{\skt{वातरुजस्तास्ता}}\Afootnote{A H; \skt{वातरुजस्तास्ता}\ K}}\skt{ }\edtext{\edtext{\skt{वेप\-थु\-र्व्वि\-ष\-म\-ज्वरः}}{\lemma{\skt{वेपथुर्}}\Afootnote{K H; \skt{वेपथु}\ A}}}{\lemma{\skt{व्विषमज्वरः}}\Afootnote{K H; \skt{र्विषमज्वरः~॥}\ A}}\skt{~॥}\pend

\pstart\versepar \ledsidenote{\passref{4.37.84}}\edtext{\skt{पित्ताभिभूते}}{\lemma{\skt{पित्ताभिभूते}}\Afootnote{A H; \skt{पित्ताविभूते}\ K}}\skt{ स्नेहे तु }\edtext{\skt{मुख्यस्य}}{\lemma{\skt{मुख्यस्य}}\Afootnote{H; \skt{मुखस्य}\ A K}}\skt{ कटुता भवेत्~। }\pend
\pstart\versepar \skt{ज्वरदाहः }\edtext{\skt{स्तथा}}{\lemma{\skt{तथा}}\Afootnote{H; \skt{ज्वरः}\ A; \skt{ज्वरस्तथा}\ K}}\skt{ स्वेदो नेत्र\-मू\-त्रा\-ङ्ग\-पी\-तता~॥}\edtext{}{\lemma{\skt{भवेत्~। ज्वरदाहः स्}}\Afootnote{H; \skt{भवेत्~। दाहस्तृष्णा}\ A; \skt{भवेत्~। दाह}\ K}}\pend

\pstart\versepar \ledsidenote{\passref{4.37.85}}\skt{श्लेष्माभिभूते स्नेहे तु प्रसेको मधुरास्यता~। }\pend
\pstart\versepar \skt{}\edtext{\skt{गौर\-व\-च्छ\-र्द्दि\-रु\-च्छ्वासः}}{\lemma{\skt{गौर\appbrk{}व\appbrk{}च्छ\appbrk{}र्द्दि\appbrk{}रु\appbrk{}च्छ्वासः}}\Afootnote{K H; \skt{गौरवं छर्दिरुच्छ्वासः}\ A}}\skt{ कृच्छ्राच्छीतज्वरो }\edtext{\skt{रुचिः}}{\lemma{\skt{रुचिः}}\Afootnote{K H; \skt{ऽरुचिः~॥}\ A}}\skt{~॥}\pend

\pstart\versepar \ledsidenote{\passref{4.37.86}}\skt{तत्र दोषाभिभूते तु }\edtext{\skt{स्नेह\-ब\-स्ति\-न्नि\-धा\-प\-येत्}}{\lemma{\skt{स्नेहबस्तिन्}}\Afootnote{K H; \skt{स्नेहे बस्तिं}\ A}}\skt{~। }\pend
\pstart\versepar \skt{}\edtext{\skt{यथा\-स्व\-न्दो\-ष\-श\-म\-ना\-न्यु\-प\-यो\-ज्यानि}}{\lemma{\skt{दोष\appbrk{}श\appbrk{}म\appbrk{}ना\appbrk{}न्यु\appbrk{}प\appbrk{}यो\appbrk{}ज्यानि}}\Afootnote{H; \skt{दोष\appbrk{}श\appbrk{}म\appbrk{}ना\appbrk{}न्यु\appbrk{}प\appbrk{}यो\appbrk{}ज्यानि}\ A K\textsuperscript{pc}; \skt{दोष\appbrk{}श\appbrk{}ना\appbrk{}न्यु\appbrk{}प\appbrk{}यो\appbrk{}ज्यानि}\ K\textsuperscript{ac}}}\skt{ यानि च~॥}\pend

\pstart\versepar \ledsidenote{\passref{4.37.87}}\edtext{\skt{अत्या\-सि\-ते\-त्वा\-ना\-भि\-भ\-वा\-त्स्नेहो}}{\lemma{\skt{अत्या\appbrk{}सि\appbrk{}ते\appbrk{}त्वा\appbrk{}ना\appbrk{}भि\appbrk{}भ\appbrk{}वा\appbrk{}त्स्नेहो}}\Afootnote{H\textsuperscript{pc}; \skt{अत्या\appbrk{}शि\appbrk{}तेऽ\appbrk{}न्ना\appbrk{}भि\appbrk{}भ\appbrk{}वा\appbrk{}त्स्नेहो}\ A; \skt{अत्या\appbrk{}सि\appbrk{}ते\appbrk{}न्ना\appbrk{}भि\appbrk{}भ\appbrk{}वा\appbrk{}त्स्नेहो}\ K; \skt{अत्या\appbrk{}सि\appbrk{}ते\appbrk{}त्वा\appbrk{}भि\appbrk{}भ\appbrk{}वा\appbrk{}त्स्नेहो}\ H\textsuperscript{ac}}}\skt{ नैति यदा तदा~। }\pend
\pstart\versepar \skt{गुरुरामाशयः }\edtext{\skt{शूलो}}{\lemma{\skt{शूलो}}\Afootnote{K H; \skt{शूलं}\ A}}\skt{ वायो\-श्चा\-प्र\-ति\-स\-ञ्चरः~॥}\pend

\pstart\versepar \ledsidenote{\passref{4.37.88}}\edtext{\skt{हृत्पी\-डा\-मु\-ख\-वै\-रस्यं}}{\lemma{\skt{हृत्पी\appbrk{}डा\appbrk{}मु\appbrk{}ख\appbrk{}वै\appbrk{}रस्यं}}\Afootnote{K H; \skt{हृत्पीडा मुखवैरस्य्ंअ}\ A}}\skt{ श्वासो }\edtext{\edtext{\skt{मूर्च्छाज्वरोऽरुचिः}}{\lemma{\skt{मूर्च्छाज्वरो}}\Afootnote{H; \skt{मूर्च्छा भ्रमो}\ A; \skt{मूर्च्छा ज्वरो}\ K}}}{\lemma{\skt{ऽरुचिः}}\Afootnote{A H; \skt{रुचिः}\ K}}\skt{~। }\pend
\pstart\versepar \skt{}\edtext{\skt{तत्रा\-प\-त\-र्प्पि\-त\-स्यान्ते}}{\lemma{\skt{तत्रा\appbrk{}प\appbrk{}त\appbrk{}र्प्पि\appbrk{}त\appbrk{}स्यान्ते}}\Afootnote{K H; \skt{तत्रा\appbrk{}प\appbrk{}त\appbrk{}र्प\appbrk{}ण\appbrk{}स्यान्ते}\ A}}\skt{ दीपनो विधिरिष्यते~॥}\pend

\pstart\versepar \ledsidenote{\passref{4.37.89}}\skt{अशुद्धस्य मलोन्मिश्रः स्नेहो नैति यदा पुनः~। }\pend
\pstart\versepar \skt{}\edtext{\skt{तदा\-ङ्ग\-स\-द\-ना\-ध्माने}}{\lemma{\skt{तदा\appbrk{}ङ्ग\appbrk{}स\appbrk{}द\appbrk{}ना\appbrk{}ध्माने}}\Afootnote{K H; \skt{तदाऽ\appbrk{}ङ्ग\appbrk{}स\appbrk{}द\appbrk{}ना\appbrk{}ध्माने}\ A}}\skt{ }\edtext{\skt{शूलः}}{\lemma{\skt{शूलः}}\Afootnote{K H; \skt{श्वासः}\ A}}\skt{ }\edtext{\skt{श्वाशश्च}}{\lemma{\skt{श्वाशश्च}}\Afootnote{H; \skt{शूलं च}\ A; \skt{श्वासश्च}\ K}}\skt{ जायते~॥}\pend

\pstart\versepar \ledsidenote{\passref{4.37.90}}\skt{पक्वा\-श\-य\-गु\-रु\-त्वञ्च तत्र दद्यान्निरूहणं~। }\pend
\pstart\versepar \skt{}\edtext{\skt{तीक्ष्णा\-ती\-क्ष्णौ\-ष\-धै\-रेव}}{\lemma{\skt{तीक्ष्णा\appbrk{}ती\appbrk{}क्ष्णौ\appbrk{}ष\appbrk{}धैर्}}\Afootnote{H; \skt{तीक्ष्णं तीक्ष्णौषधैर्}\ A; \skt{तीक्ष्ण\appbrk{}ती\appbrk{}क्ष्णौ\appbrk{}ष\appbrk{}धैर्}\ K}}\skt{ सिद्ध\-ञ्चा\-प्य\-नु\-वा\-सनं~॥}\pend

\pstart\versepar \ledsidenote{\passref{4.37.91}}\skt{शुद्धस्य }\edtext{\skt{दूरानुसृते}}{\lemma{\skt{दूरानुसृते}}\Afootnote{A K\textsuperscript{pc} H; \skt{दूनुसृते}\ K\textsuperscript{ac}}}\skt{ स्नेहे }\edtext{\skt{स्नेहोनुदर्शनं}}{\lemma{\skt{स्नेहोनुदर्शनं}}\Afootnote{K H; \skt{स्नेहस्य दर्शनम्~।}\ A}}\skt{~। }\pend
\pstart\versepar \skt{}\edtext{\skt{गात्रे}}{\lemma{\skt{गात्रे}}\Afootnote{K H; \skt{गात्रेषु}\ A}}\skt{ }\edtext{\edtext{\edtext{\edtext{\skt{सर्व्वे\-न्द्रि\-या\-णा\-ञ्चा\-प्य\-नु\-ले\-पोऽ\-व\-पी\-डनं}}{\lemma{\skt{सर्व्वेन्द्रियाणाञ्}}\Afootnote{A H; \skt{सर्वे\appbrk{}न्द्रि\appbrk{}या\appbrk{}ण\appbrk{}आञ्}\ K}}}{\lemma{\skt{चाप्य्}}\Afootnote{K H; \skt{उप}\ A}}}{\lemma{\skt{अनुलेपो}}\Afootnote{H; \skt{लेपो}\ A; \skt{अनुलेपो}\ K}}}{\lemma{\skt{ऽवपीडनं}}\Afootnote{H; \skt{ऽवसादनम्~॥}\ A; \skt{वपीडनं~।}\ K}}\skt{~॥}\pend

\pstart\versepar \ledsidenote{\passref{4.37.92}}\edtext{\skt{स्नेह\-स\-न्धि\-मु\-ख\-ञ्चापि}}{\lemma{\skt{स्नेह\appbrk{}स\appbrk{}न्धि\appbrk{}मु\appbrk{}खञ्}}\Afootnote{H; \skt{स्नेहगन्धि मुखं}\ A K}}\skt{ कास\-श्वा\-सा\-व\-रो\-चकः~। }\pend
\pstart\versepar \skt{अति\-पी\-डि\-त\-व\-त्तत्र सिद्धि\-रा\-स्था\-प\-न\-न्तथा~।}\pend

\pstart\versepar \ledsidenote{\passref{4.37.93}}\skt{अस्वि\-न्न\-स्या\-वि\-शु\-द्धस्य }\edtext{\edtext{\skt{स्नेहोऽल्पस्तु}}{\lemma{\skt{स्नेहो}}\Afootnote{A H; \skt{स्नेहो}\ K}}}{\lemma{\skt{ऽल्पस्तु}}\Afootnote{H; \skt{ऽल्पः सं}\ A; \skt{ल्पस्तु}\ K}}\skt{ प्रयोजितः~। }\pend
\pstart\versepar \skt{शीतो }\edtext{\skt{मृदुश्च}}{\lemma{\skt{मृदुश्च}}\Afootnote{K H; \skt{मृदुश्च}\ A}}\skt{ नाभ्येति ततो }\edtext{\skt{मन्दम्प्रवाहति}}{\lemma{\skt{प्रवाहति}}\Afootnote{K H; \skt{प्रवाहते~॥}\ A}}\skt{~॥}\pend

\pstart\versepar \ledsidenote{\passref{4.37.94}}\edtext{\skt{विब\-न्ध\-गौ\-र\-वा\-ध्मा\-न\-शूलैः}}{\lemma{\skt{विब\appbrk{}न्ध\appbrk{}गौ\appbrk{}र\appbrk{}वा\appbrk{}ध्मा\appbrk{}न\appbrk{}शूलैः}}\Afootnote{K H; \skt{विब\appbrk{}न्ध\appbrk{}गौ\appbrk{}र\appbrk{}वा\appbrk{}ध्मा\appbrk{}न\appbrk{}शूलाः}\ A}}\skt{ पक्वाशयम्प्रति~। }\pend
\pstart\versepar \skt{}\edtext{\skt{तत्रा\-स्था\-प\-न\-मे\-वाशु}}{\lemma{\skt{एवाशु}}\Afootnote{A H; \skt{एवाशुः}\ K}}\skt{ प्रयोज्यं सानुवासनं~॥}\pend

\pstart\versepar \ledsidenote{\passref{4.37.95}}\edtext{\skt{अल्प\-म्भु\-क्त\-व\-तो\-ल्पाशः}}{\lemma{\skt{भुक्तवतोल्पाशः}}\Afootnote{K H; \skt{भुक्तवतोऽल्पो हि}\ A}}\skt{ स्नेहो मन्दगुणस्तथा~। }\pend
\pstart\versepar \skt{दत्तो नैति }\edtext{\skt{क्लमोत्क्लेशो}}{\lemma{\skt{क्लमोत्क्लेशो}}\Afootnote{H; \skt{क्लमोत्क्लेशौ}\ A; \skt{क्लमोत्लेशो}\ K}}\skt{ }\edtext{\edtext{\skt{भृश\-ञ्चा\-रि\-त\-मा\-व\-हेत्}}{\lemma{\skt{चारितम्}}\Afootnote{K H; \skt{चारति}\ A}}}{\lemma{\skt{आवहेत्}}\Afootnote{K H; \skt{मावहेत्~॥}\ A}}\skt{~॥}\pend

\pstart\versepar \ledsidenote{\passref{4.37.96}}\skt{तत्राप्यास्थापनं कार्यं }\edtext{\skt{शोधनीयेन}}{\lemma{\skt{शोधनीयेन}}\Afootnote{A H; \skt{शोधणीयेन}\ K}}\skt{ बस्तिना~। }\pend
\pstart\versepar \skt{}\edtext{\skt{अनुवासनञ्च}}{\lemma{\skt{अनुवासनञ्च}}\Afootnote{K H; \skt{(अन्वासनं च}\ A}}\skt{ स्नेहेन }\edtext{\skt{शोधने}}{\lemma{\skt{शोधने}}\Afootnote{K H; \skt{शोधनी}\ A}}\skt{ }\edtext{\skt{नैव}}{\lemma{\skt{नैव}}\Afootnote{K H; \skt{येन}\ A}}\skt{ शस्यते~॥}\pend

\pstart\versepar \ledsidenote{\passref{4.37.97}}\edtext{\skt{अहोरात्रादपि}}{\lemma{\skt{अहोरात्रादपि}}\Afootnote{K H; \skt{अहोरात्रादपि}\ A}}\skt{ स्नेहः प्रत्यागच्छन्न दुष्यति~। }\pend
\pstart\versepar \skt{कुर्या\-द्ब\-स्ति\-गु\-णा\-श्चापि }\edtext{\skt{जीर्ण्णस्वल्पगुणो}}{\lemma{\skt{जीर्ण्णस्वल्पगुणो}}\Afootnote{H; \skt{जीर्णस्त्वल्पगुणो}\ A; \skt{जीर्ण्णस्वल्पगोणो}\ K}}\skt{ भवेत्~॥}\pend

\pstart\versepar \ledsidenote{\passref{4.37.99}}\edtext{\skt{अना\-या\-न्त\-म\-हो\-रा\-त्रा\-त्स्नेहां}}{\lemma{\skt{अहोरात्रात्स्नेहां}}\Afootnote{K H; \skt{त्वहोरात्रात्स्नेहं}\ A}}\skt{ }\edtext{\skt{शंशो\-ध\-नै\-र्ह\-रेत्}}{\lemma{\skt{शंशोधनैर्}}\Afootnote{H; \skt{संशोधनैर्}\ A; \skt{संशोधनै}\ K}}\skt{~। }\pend
\pstart\versepar \skt{स्नेह\-ब\-स्ता\-व\-ना\-याते }\edtext{\skt{नान्यस्नेहो}}{\lemma{\skt{नान्यस्नेहो}}\Afootnote{K H; \skt{नान्यः स्नेहो}\ A}}\skt{ विधीयते~॥}\pend

\pstart\versepar \ledsidenote{\passref{4.37.100}}\edtext{\skt{इत्युक्तो}}{\lemma{\skt{उक्तो}}\Afootnote{K H; \skt{उक्ता}\ A}}\skt{ व्यापदः }\edtext{\skt{सर्व्वाः}}{\lemma{\skt{सर्व्वाः}}\Afootnote{K H; \skt{सर्वा}\ A}}\skt{ }\edtext{\skt{सल\-क्ष\-ण\-चि\-कि\-त्सिताः~॥ ओ}}{\lemma{\skt{सल\appbrk{}क्ष\appbrk{}ण\appbrk{}चि\appbrk{}कि\appbrk{}त्सिताः~॥ ओ}}\Afootnote{H; \skt{सल\appbrk{}क्ष\appbrk{}ण\appbrk{}चि\appbrk{}कि\appbrk{}त्सिताः~।}\ A; \skt{सल\appbrk{}क्ष\appbrk{}ण\appbrk{}चि\appbrk{}कि\appbrk{}त्सिताः~॥{\symfont ❈}॥}\ K}}\skt{~॥ }\pend
\pstart\versepar \skt{बस्ते\-रु\-त्त\-र\-सं\-ज्ञस्य }\edtext{\skt{शृणु}}{\lemma{\skt{शृणु}}\Afootnote{K H; \skt{विधिं}\ A}}\skt{ वक्ष्याम्यतः परं~।}\pend

\pstart\versepar \ledsidenote{\passref{4.37.101}}\edtext{\skt{चतु\-र्द्द\-शा\-ङ्गु\-ल\-न्नेत्रं}}{\lemma{\skt{नेत्रं}}\Afootnote{K H; \skt{नेत्र}\ A}}\skt{ }\edtext{\edtext{\skt{कार्यम्वैद्येन}}{\lemma{\skt{कार्यम्}}\Afootnote{K H; \emph{om.}\ A}}}{\lemma{\skt{वैद्येन}}\Afootnote{K H; \emph{om.}\ A}}\skt{ }\edtext{\skt{जानता}}{\lemma{\skt{जानता}}\Afootnote{K H; \skt{मातुराङ्}\ A}}\skt{~॥ }\pend
\pstart\versepar \skt{}\edtext{\edtext{\edtext{\skt{माल\-ती\-पु\-ष्प\-वृ\-न्ता\-भ\-न्त\-च्च्छि\-द्र\-ञ्च्छि\-द्र\-मेव}}{\lemma{\skt{माल\appbrk{}ती\appbrk{}पु\appbrk{}ष्प\appbrk{}वृ\appbrk{}न्ता\appbrk{}भन्}}\Afootnote{K H; \skt{गुलसंमितम्~। माल\appbrk{}ती\appbrk{}पु\appbrk{}ष्प\appbrk{}वृन्ता}\ A}}}{\lemma{\skt{तच्च्छिद्रञ्}}\Afootnote{K H; \skt{ग्रं}\ A}}}{\lemma{\skt{एव}}\Afootnote{K H; \skt{सर्ष}\ A}}\skt{ }\edtext{\skt{च}}{\lemma{\skt{च}}\Afootnote{K H; \skt{पनिर्गमम्~॥}\ A}}\skt{~।}\pend

\pstart\versepar \ledsidenote{\passref{4.37.104}}\edtext{\skt{मेढ्रयामसमं}}{\lemma{\skt{मेढ्रयामसमं}}\Afootnote{K H; \skt{मेढ्रायामसमं}\ A}}\skt{ }\edtext{\skt{केचिदिच्छन्ति}}{\lemma{\skt{केचिदिच्छन्ति}}\Afootnote{K H; \skt{केचिदिच्छन्ति}\ A}}\skt{ खलु तद्विदः~॥ }\pend
\pstart\versepar \skt{}\edtext{\skt{स्नेहप्रमाणं}}{\lemma{\skt{स्नेहप्रमाणं}}\Afootnote{K\textsuperscript{pc} H\textsuperscript{pc}; \emph{om.}\ K\textsuperscript{ac}}}\skt{ }\edtext{\edtext{\edtext{\skt{पर\-म\-म्प्र\-सृ\-त\-श्चात्र}}{\lemma{\skt{परमम्}}\Afootnote{K\textsuperscript{pc} H\textsuperscript{pc}; \emph{om.}\ K\textsuperscript{ac}}}}{\lemma{\skt{प्रसृतश्}}\Afootnote{K\textsuperscript{pc} H\textsuperscript{pc}; \emph{om.}\ K\textsuperscript{ac}}}}{\lemma{\skt{चात्र}}\Afootnote{K\textsuperscript{pc} H\textsuperscript{pc}; \emph{om.}\ K\textsuperscript{ac}}}\skt{ }\edtext{\skt{कीर्त्तितः}}{\lemma{\skt{कीर्त्तितः}}\Afootnote{K\textsuperscript{pc} H\textsuperscript{pc}; \emph{om.}\ K\textsuperscript{ac}}}\skt{~। }\pend
\pstart\versepar \skt{}\edtext{\skt{पञ्चविंशति वर्षाणां}}{\lemma{\skt{पञ्चविंशति वर्षाणां}}\Afootnote{K\textsuperscript{pc} H\textsuperscript{pc}; \emph{om.}\ K\textsuperscript{ac}}}\skt{ }\edtext{\skt{स वाग्बु\-द्धि\-वि\-क\-ल्पितं}}{\lemma{\skt{स वाग्बु\appbrk{}द्धि\appbrk{}वि\appbrk{}क\appbrk{}ल्पितं}}\Afootnote{H\textsuperscript{pc}; \emph{om.}\ K\textsuperscript{ac}; \skt{सत्वा\appbrk{}ग्बु\appbrk{}द्धि\appbrk{}वि\appbrk{}क\appbrk{}ल्पितं~।}\ K\textsuperscript{pc}}}\skt{ }\pend
\pstart\versepar \skt{}\edtext{\skt{निकृष्टकर्ण्णिकां}}{\lemma{\skt{निकृष्टकर्ण्णिकां}}\Afootnote{H\textsuperscript{pc}; \emph{om.}\ K\textsuperscript{ac}; \skt{निकृष्टकर्ण्णिकं}\ K\textsuperscript{pc}}}\skt{ }\edtext{\skt{मध्ये}}{\lemma{\skt{मध्ये}}\Afootnote{K\textsuperscript{pc} H\textsuperscript{pc}; \emph{om.}\ K\textsuperscript{ac}}}\skt{ }\edtext{\edtext{\skt{नारी\-णा\-ञ्च\-तु\-र\-ङ्गुलं}}{\lemma{\skt{नारीणाञ्}}\Afootnote{K\textsuperscript{pc} H\textsuperscript{pc}; \emph{om.}\ K\textsuperscript{ac}}}}{\lemma{\skt{चतुरङ्गुलं}}\Afootnote{K\textsuperscript{pc} H\textsuperscript{pc}; \emph{om.}\ K\textsuperscript{ac}}}\skt{~। }\pend
\pstart\versepar \skt{}\edtext{\skt{मूत्र\-स्रो\-त\-प\-री\-णाहं}}{\lemma{\skt{मूत्र\appbrk{}स्रो\appbrk{}त\appbrk{}प\appbrk{}री\appbrk{}णाहं}}\Afootnote{K\textsuperscript{pc} H\textsuperscript{pc}; \emph{om.}\ K\textsuperscript{ac}}}\skt{ }\edtext{\skt{मुङ्गवाहि}}{\lemma{\skt{मुङ्गवाहि}}\Afootnote{H\textsuperscript{pc}; \emph{om.}\ K\textsuperscript{ac}; \skt{मुद्गवाहि}\ K\textsuperscript{pc}}}\skt{ }\edtext{\skt{दशाङ्गुलं}}{\lemma{\skt{दशाङ्गुलं}}\Afootnote{K\textsuperscript{pc} H\textsuperscript{pc}; \emph{om.}\ K\textsuperscript{ac}}}\skt{~॥ }\pend
\pstart\versepar \skt{}\edtext{\skt{तासा\-म\-प\-त्य\-मा\-र्ग्गेषु}}{\lemma{\skt{तासा\appbrk{}म\appbrk{}प\appbrk{}त्य\appbrk{}मा\appbrk{}र्ग्गेषु}}\Afootnote{H; \skt{तासा\appbrk{}म\appbrk{}प\appbrk{}त्य\appbrk{}मार्गे तु}\ A K\textsuperscript{pc}; \emph{om.}\ K\textsuperscript{ac}}}\skt{ }\edtext{\skt{निद\-ध्या\-च्च\-तु\-र\-ङ्गुलं}}{\lemma{\skt{निद\appbrk{}ध्या\appbrk{}च्च\appbrk{}तु\appbrk{}र\appbrk{}ङ्गुलं}}\Afootnote{K\textsuperscript{pc} H; \skt{निद\appbrk{}ध्या\appbrk{}च्च\appbrk{}तु\appbrk{}र\appbrk{}ङ्गु\appbrk{}लम्~॥}\ A; \emph{om.}\ K\textsuperscript{ac}}}\skt{~।}\edtext{}{\lemma{\skt{स्नेहप्रमाणं… ता}}\Afootnote{\emph{om.}\ A}}\edtext{}{\lemma{\skt{स्नेहप्रमाणं… ता}}\Afootnote{\emph{om.}\ H\textsuperscript{ac}}}\pend

\pstart \ledsidenote{\passref{4.37.105}}\edtext{\edtext{\skt{द्व्यङ्गु\-ल\-म्मू\-त्र\-मार्ग्गे}}{\lemma{\skt{द्व्यङ्गुलम्}}\Afootnote{A K\textsuperscript{pc} H; \emph{om.}\ K\textsuperscript{ac}}}\skt{ तु}}{\lemma{\skt{मूत्रमार्ग्गे तु}}\Afootnote{A K\textsuperscript{pc} H; \emph{om.}\ K\textsuperscript{ac}}}\skt{ }\edtext{\skt{कन्यानां}}{\lemma{\skt{कन्यानां}}\Afootnote{A K\textsuperscript{pc} H; \emph{om.}\ K\textsuperscript{ac}}}\skt{ }\edtext{\skt{कर्ण्णिकाङ्गुले}}{\lemma{\skt{कर्ण्णिकाङ्गुले}}\Afootnote{K\textsuperscript{pc} H; \skt{त्वेकमङ्गुलम्~। विधेयं चाङ्गुलं तासां विधिवद्वक्ष्यते यथा~॥}\ A; \emph{om.}\ K\textsuperscript{ac}}}\skt{~॥}\pend

\pstart\versepar \ledsidenote{\passref{4.37.106}}\skt{स्नेहस्य }\edtext{\skt{प्रसृतश्चात्र}}{\lemma{\skt{प्रसृतश्}}\Afootnote{K H; \skt{प्रसृतं}\ A}}\skt{ }\edtext{\skt{स्वांगु\-ली\-मू\-ल\-स\-म्मितः}}{\lemma{\skt{स्वांगु\appbrk{}ली\appbrk{}मू\appbrk{}ल\appbrk{}स\appbrk{}म्मितः}}\Afootnote{K H; \skt{स्वाङ्गु\appbrk{}ली\appbrk{}मू\appbrk{}ल\appbrk{}सं\appbrk{}मि\appbrk{}तम्~।}\ A}}\skt{~। }\pend
\pstart\versepar \skt{}\edtext{\edtext{\edtext{\skt{पर\-म्प्र\-मा\-ण\-म्वि\-हि\-त\-म\-वा\-ग्बु\-द्धि\-वि\-क\-ल्पितं}}{\lemma{\skt{परम्}}\Afootnote{K H; \skt{देयं}\ A}}}{\lemma{\skt{विहितम्}}\Afootnote{K H; \skt{परम}\ A}}}{\lemma{\skt{अवा\appbrk{}ग्बु\appbrk{}द्धि\appbrk{}वि\appbrk{}क\appbrk{}ल्पितं}}\Afootnote{K H; \skt{मर्वा\appbrk{}ग्बु\appbrk{}द्धि\appbrk{}वि\appbrk{}क\appbrk{}ल्पि\appbrk{}तम्~॥}\ A}}\skt{~॥}\pend

\pstart \ledsidenote{\passref{4.37.107}}\skt{औरभ्रः }\edtext{\skt{सौकरोवापि}}{\lemma{\skt{सौकरोवापि}}\Afootnote{H; \skt{शौकरो वाऽपि}\ A; \skt{सौकरो वापि}\ K}}\skt{ बस्तिराजश्च }\edtext{\skt{पूजतिः}}{\lemma{\skt{पूजतिः}}\Afootnote{H; \skt{पूजितः~। तदलाभे प्रयुञ्जीत गलचर्म तु पक्षिणाम्~॥}\ A; \skt{पूजितः~।}\ K}}\skt{~।}\pend

\pstart\versepar \ledsidenote{\passref{4.37.108}}\skt{तस्यालाभे दृतेः पादो मृदुचर्म्म }\edtext{\skt{कृतो}}{\lemma{\skt{कृतो}}\Afootnote{K H; \skt{ततो}\ A}}\skt{ पि वा~॥ }\pend
\pstart\versepar \skt{}\edtext{\skt{अथा\-तु\-र\-मु\-प\-स्नि\-ग्ध\-स्वि\-न्न\-म्प्र\-शि\-थि\-ला\-शयं}}{\lemma{\skt{उप\appbrk{}स्नि\appbrk{}ग्ध\appbrk{}स्वि\appbrk{}न्नम्}}\Afootnote{K H; \skt{उपस्निग्धं स्विन्नं}\ A}}\skt{~।}\pend

\pstart\versepar \ledsidenote{\passref{4.37.109}}\skt{यवागूं सघृ\-त\-क्षी\-रा\-म्पी\-त\-वन्तं यथाबलं~॥ }\pend
\pstart\versepar \skt{निष\-ण्ण\-मा\-जा\-नु\-समे पीठे }\edtext{\skt{सापाश्रये}}{\lemma{\skt{सापाश्रये}}\Afootnote{K H; \skt{सोपाश्रये}\ A}}\skt{ समं~।}\pend

\pstart\versepar \ledsidenote{\passref{4.37.110}}\edtext{\skt{स्वभ्य\-क्त\-ब\-स्ति\-शि\-रसं}}{\lemma{\skt{स्वभ्य\appbrk{}क्त\appbrk{}ब\appbrk{}स्ति\appbrk{}शि\appbrk{}रसं}}\Afootnote{K H; \skt{स्वभ्य\appbrk{}क्त\appbrk{}ब\appbrk{}स्ति\appbrk{}मू\appbrk{}र्धानं}\ A}}\skt{ }\edtext{\skt{तैलनोष्णेन}}{\lemma{\skt{तैलनोष्णेन}}\Afootnote{H; \skt{तैलेनोष्णेन}\ A K}}\skt{ }\edtext{\skt{पूजितः}}{\lemma{\skt{पूजितः}}\Afootnote{K H; \skt{मानवम्~।}\ A}}\skt{~॥ }\pend
\pstart\versepar \skt{ततः समं स्थापयित्वा }\edtext{\skt{नालमस्यमस्य}}{\lemma{\skt{नालमस्यमस्य}}\Afootnote{H\textsuperscript{ac}; \skt{नालमस्य}\ A K H\textsuperscript{pc}}}\skt{ प्रहर्षितं~।}\pend

\pstart\versepar \ledsidenote{\passref{4.37.111}}\skt{पूर्व्वं }\edtext{\skt{सलाकयान्विस्य}}{\lemma{\skt{सलाकयान्विस्य}}\Afootnote{H; \skt{शलाकयाऽन्विष्य}\ A; \skt{शलाकयान्विष्य}\ K}}\skt{ }\edtext{\edtext{\skt{मार्ग्ग\-न्ने\-त्र\-म\-न\-न्तरं}}{\lemma{\skt{मार्ग्गन्नेत्रम्}}\Afootnote{K H; \skt{ततो नेत्र}\ A}}}{\lemma{\skt{अनन्तरं}}\Afootnote{K H; \skt{मनन्तरम्~।}\ A}}\skt{~॥ }\pend
\pstart\versepar \skt{शनैः शनैर्घृताभ्यक्तं }\edtext{\skt{निद\-ध्या\-द\-ङ्गु\-लानि}}{\lemma{\skt{निद\appbrk{}ध्या\appbrk{}द\appbrk{}ङ्गु\appbrk{}लानि}}\Afootnote{H; \skt{विद\appbrk{}ध्या\appbrk{}द\appbrk{}ङ्गु\appbrk{}लानि}\ A; \skt{निद\appbrk{}ध्या\appbrk{}द\appbrk{}ङ्गु\appbrk{}लानि}\ K}}\skt{ षट्~।}\pend

\pstart \ledsidenote{\passref{4.37.112}}\edtext{\edtext{\skt{ततोऽ\-व\-पी\-ड\-ये\-द्बस्तिः}}{\lemma{\skt{ततो}}\Afootnote{H; \skt{मेढ्रयामसमं केचिदिच्छन्ति प्रणिधानकम्~। ततो}\ A; \skt{ततो}\ K}}}{\lemma{\skt{ऽव\appbrk{}पी\appbrk{}ड\appbrk{}ये\appbrk{}द्बस्तिः}}\Afootnote{H; \skt{ऽव\appbrk{}पी\appbrk{}ड\appbrk{}ये\appbrk{}द्बस्तिं}\ A; \skt{वपीडयेद्बस्ति}\ K}}\skt{ शनैर्न्नेत्रञ्च निर्हरेत्~॥}\pend

\pstart\versepar \ledsidenote{\passref{4.37.113}}\skt{ततः }\edtext{\edtext{\skt{प्रत्या\-ग\-त\-स्ने\-ह\-म\-प\-राह्ने}}{\lemma{\skt{प्रत्या\appbrk{}ग\appbrk{}त\appbrk{}स्ने\appbrk{}हम्}}\Afootnote{K H; \skt{प्रत्यागतस्नेह}\ A}}}{\lemma{\skt{अपराह्ने}}\Afootnote{H; \skt{मपराह्णे}\ A; \skt{अपरांह्णे}\ K}}\skt{ विचक्षणः~। }\pend
\pstart\versepar \skt{}\edtext{\skt{भोज\-ये\-त्प\-य\-सा\-मात्रा}}{\lemma{\skt{भोज\appbrk{}ये\appbrk{}त्प\appbrk{}य\appbrk{}सा\appbrk{}मात्रा}}\Afootnote{H; \skt{भोजयेत्पयसा मात्रां}\ A; \skt{भोजयेत्पयसां मात्रां}\ K}}\skt{ यूषे\-णा\-थ\-र\-सेन वा~।}\pend

\pstart\versepar \ledsidenote{\passref{4.37.114}}\skt{अनेन विधिना }\edtext{\skt{दद्या\-द्ब\-स्ती\-स्त्री\-न्च\-तु\-रोऽपि}}{\lemma{\skt{दद्या\appbrk{}द्ब\appbrk{}स्ती\appbrk{}स्त्री\appbrk{}न्च\appbrk{}तुरो}}\Afootnote{H; \skt{दद्या\appbrk{}द्ब\appbrk{}स्तीं\appbrk{}स्त्रीं\appbrk{}श्च\appbrk{}तुरो}\ A; \skt{दध्या\appbrk{}द्ब\appbrk{}स्तीं\appbrk{}स्त्री\appbrk{}श्च\appbrk{}तुरो}\ K}}\skt{ वा~। }\pend
\pstart\versepar \skt{}\edtext{\skt{अनु\-वा\-स\-न\-सि\-द्धिञ्च}}{\lemma{\skt{अनु\appbrk{}वा\appbrk{}स\appbrk{}न\appbrk{}सि\appbrk{}द्धिञ्च}}\Afootnote{K\textsuperscript{pc} H\textsuperscript{pc}; \skt{ऊ}\ A}}\skt{ }\edtext{\skt{वीक्ष्य}}{\lemma{\skt{वीक्ष्य}}\Afootnote{K\textsuperscript{pc} H\textsuperscript{pc}; \skt{र्ध्व}\ A}}\skt{ }\edtext{\skt{कर्म्म}}{\lemma{\skt{कर्म्म}}\Afootnote{K\textsuperscript{pc} H\textsuperscript{pc}; \skt{जान्वै}\ A}}\skt{ प्रयोजयेत्~॥ }\pend
\pstart\versepar \skt{उत्तानायै }\edtext{\skt{स्त्रियै}}{\lemma{\skt{स्त्रियै}}\Afootnote{K H; \skt{विच}\ A}}\skt{ }\edtext{\skt{दद्यादूर्ध्वजात्वै}}{\lemma{\skt{दद्यादूर्ध्वजात्वै}}\Afootnote{K H; \skt{क्षणः~॥}\ A}}\skt{ }\edtext{\skt{समाहितः}}{\lemma{\skt{समाहितः}}\Afootnote{K H; \emph{om.}\ A}}\skt{~।}\edtext{}{\lemma{\skt{अनु\appbrk{}वा\appbrk{}स\appbrk{}न\appbrk{}सि\appbrk{}द्धिञ्… उ}}\Afootnote{\emph{om.}\ H\textsuperscript{ac}}}\edtext{}{\lemma{\skt{अनु\appbrk{}वा\appbrk{}स\appbrk{}न\appbrk{}सि\appbrk{}द्धिञ्… उ}}\Afootnote{\emph{om.}\ K\textsuperscript{ac}}}\edtext{}{\lemma{\skt{प्रयोजयेत्~॥ उत्तानायै}}\Afootnote{K\textsuperscript{pc} H\textsuperscript{pc}; \skt{स्त्रियै दद्यादुत्तानायै}\ A}}\pend

\pstart\versepar \ledsidenote{\passref{4.37.115}}\edtext{\skt{कन्येतरस्यै}}{\lemma{\skt{कन्येतरस्यै}}\Afootnote{K H; \skt{सम्यक्प्र}\ A}}\skt{ }\edtext{\skt{कन्यायै}}{\lemma{\skt{कन्यायै}}\Afootnote{K H; \skt{पीडयेद्योनिं}\ A}}\skt{ दद्या\-त्सु\-मृ\-दु\-पी\-डितं~॥ }\pend
\pstart\versepar \skt{}\edtext{\skt{विकर्ण्णिकेन}}{\lemma{\skt{विकर्ण्णिकेन}}\Afootnote{K H; \skt{त्रिकर्णिकेन}\ A}}\skt{ नेत्रेण }\edtext{\skt{दद्या\-द्यो\-नि\-मु\-ख\-म्प्रति}}{\lemma{\skt{दद्या\appbrk{}द्यो\appbrk{}नि\appbrk{}मु\appbrk{}खम्}}\Afootnote{A H; \skt{दद्याद्योनिमुखं}\ K}}\skt{~।}\pend

\pstart \ledsidenote{\passref{4.37.116}}\skt{गर्ब्भा\-श\-य\-वि\-शु\-द्ध्यर्थं }\edtext{\skt{स्नेहने}}{\lemma{\skt{स्नेहने}}\Afootnote{H; \skt{स्नेहेन}\ A; \skt{स्नेहेने}\ K}}\skt{ द्विगुणेन }\edtext{\skt{तु}}{\lemma{\skt{तु}}\Afootnote{K H; \skt{तु~। क्वाथप्रमाणं प्रसृतं स्त्रिया द्विप्रसृतं भवेत्~॥}\ A}}\skt{~॥}\pend

\pstart \ledsidenote{\passref{4.37.117}}\edtext{\skt{अप्रत्यागच्छति}}{\lemma{\skt{अप्रत्यागच्छति}}\Afootnote{K H; \skt{कन्येतरस्याः कन्या\appbrk{}या\appbrk{}स्त\appbrk{}द्व\appbrk{}द्ब\appbrk{}स्ति\appbrk{}प्र\appbrk{}मा\appbrk{}ण\appbrk{}कम्~। अप्रत्यागच्छति}\ A}}\skt{ भिष\-ग्ब\-स्ता\-वु\-त्त\-र\-सं\-ज्ञिते~।}\pend

\pstart\versepar \ledsidenote{\passref{4.37.118}}\skt{भूयो }\edtext{\skt{बस्तिन्निदध्याच्च}}{\lemma{\skt{निदध्याच्च}}\Afootnote{K H; \skt{निदध्यात्तु}\ A}}\skt{ }\edtext{\skt{संयुतं}}{\lemma{\skt{संयुतं}}\Afootnote{K H; \skt{संयुक्तं}\ A}}\skt{ }\edtext{\skt{शोधनौषधैः}}{\lemma{\skt{शोधनौषधैः}}\Afootnote{K H; \skt{शोधनैर्गणैः~।}\ A}}\skt{~॥ }\pend
\pstart\versepar \skt{}\edtext{\skt{पायौ}}{\lemma{\skt{पायौ}}\Afootnote{K H; \skt{गुदे}\ A}}\skt{ }\edtext{\skt{वर्त्तिः}}{\lemma{\skt{वर्त्तिः}}\Afootnote{H; \skt{वर्तिं}\ A K}}\skt{ }\edtext{\skt{निदध्याद्वा}}{\lemma{\skt{निदध्याद्वा}}\Afootnote{K H; \skt{निदध्याद्वा}\ A}}\skt{ प्रोक्तङ्गुल्म }\edtext{\skt{चिकित्सिते~।}}{\lemma{\skt{प्रोक्तङ्गुल्म चिकित्सिते~।}}\Afootnote{\emph{om.}\ A}}\pend

\pstart\versepar \ledsidenote{\passref{4.37.119}}\skt{प्रवेशयेद्वा }\edtext{\edtext{\skt{मति\-मा\-न्ना\-ब\-स्ति\-द्वा\-र\-मे\-षणीं}}{\lemma{\skt{मतिमान्}}\Afootnote{K H; \skt{मति}\ A}}}{\lemma{\skt{नाब\appbrk{}स्ति\appbrk{}द्वा\appbrk{}र\appbrk{}मे\appbrk{}षणीं}}\Afootnote{K\textsuperscript{pc} H; \skt{मान्ब\appbrk{}स्ति\appbrk{}द्वा\appbrk{}र\appbrk{}म\appbrk{}थै\appbrk{}ष\appbrk{}णीम्~।}\ A; \skt{नाब\appbrk{}स्ति\appbrk{}र\appbrk{}द्वा\appbrk{}र\appbrk{}मे\appbrk{}षणीं~।}\ K\textsuperscript{ac}}}\skt{~॥ }\pend
\pstart\versepar \skt{}\edtext{\skt{पीडयेद्वाप्यधो}}{\lemma{\skt{पीडयेद्वाप्यधो}}\Afootnote{H; \skt{पीडयेद्वाऽप्यधो}\ A; \skt{पीडयेद्वाप्यध}\ K}}\skt{ }\edtext{\edtext{\skt{नाभे\-र्ब्ब\-स्ते\-रु\-परि}}{\lemma{\skt{नाभेर्}}\Afootnote{H; \skt{नाभे}\ A; \skt{र्नाभेर्}\ K}}}{\lemma{\skt{ब्बस्तेरुपरि}}\Afootnote{H; \skt{र्बलेनोत्तर}\ A; \skt{बस्तेरुपरि}\ K}}\skt{ मुष्टिना~।}\pend

\pstart\versepar \ledsidenote{\passref{4.37.120}}\skt{आरग्वधस्य }\edtext{\skt{पत्रेषु}}{\lemma{\skt{पत्रेषु}}\Afootnote{K H; \skt{पत्रैस्तु}\ A}}\skt{ }\edtext{\skt{निर्ग्गुण्ड्याः}}{\lemma{\skt{निर्ग्गुण्ड्याः}}\Afootnote{A H; \skt{निर्गुणण्ड्या}\ K}}\skt{ }\edtext{\skt{सुरसस्य च}}{\lemma{\skt{सुरसस्य च}}\Afootnote{K H; \skt{स्वरसेन च~।}\ A}}\skt{~॥ }\pend
\pstart\versepar \skt{}\edtext{\skt{कार्या}}{\lemma{\skt{कार्या}}\Afootnote{K H; \skt{कुर्या}\ A}}\skt{ }\edtext{\skt{गोमूत्रपिष्टेषु}}{\lemma{\skt{गोमूत्रपिष्टेषु}}\Afootnote{K H; \skt{द्गोमूत्रपिष्टेषु}\ A}}\skt{ }\edtext{\skt{वर्त्तिर्व्वा}}{\lemma{\skt{वर्त्तिर्व्वा}}\Afootnote{K H; \skt{वर्तीर्वा}\ A}}\skt{ }\edtext{\skt{सैन्धवेन}}{\lemma{\skt{सैन्धवेन}}\Afootnote{K H; \skt{ऽपि ससैन्धवाः~॥}\ A}}\skt{ }\edtext{\skt{वा}}{\lemma{\skt{वा}}\Afootnote{K H; \emph{om.}\ A}}\skt{~।}\pend

\pstart\versepar \ledsidenote{\passref{4.37.121}}\edtext{\skt{मुद्गै\-ला\-स\-र्ष\-प\-समा}}{\lemma{\skt{मुद्गै\appbrk{}ला\appbrk{}स\appbrk{}र्ष\appbrk{}प\appbrk{}समा}}\Afootnote{H; \skt{मुद्गै\appbrk{}ला\appbrk{}स\appbrk{}र्ष\appbrk{}प\appbrk{}समाः}\ A; \skt{मुद्गै\appbrk{}ला\appbrk{}स\appbrk{}र्ष\appbrk{}प\appbrk{}समां}\ K}}\skt{ प्रविभज्य }\edtext{\skt{वयांसि च}}{\lemma{\skt{वयांसि च}}\Afootnote{K H; \skt{वयांसि तु~।}\ A}}\skt{~॥ }\pend
\pstart\versepar \skt{}\edtext{\skt{बस्ते\-र्न्ना\-ग\-म\-ना\-र्थाय तान्नि\-द\-ध्या\-च्छ\-ला\-कया}}{\lemma{\skt{बस्ते\appbrk{}र्न्ना\appbrk{}ग\appbrk{}म\appbrk{}ना\appbrk{}र्थाय तान्}}\Afootnote{H; \skt{बस्ते\appbrk{}रा\appbrk{}ग\appbrk{}म\appbrk{}ना\appbrk{}र्थाय ता}\ A; \skt{बस्ते\appbrk{}रा\appbrk{}ग\appbrk{}म\appbrk{}ना\appbrk{}र्थाय तान्}\ K}}\skt{~।}\pend

\pstart\versepar \ledsidenote{\passref{4.37.122}}\edtext{\skt{अगा\-र\-धू\-म\-बृ\-हती}}{\lemma{\skt{अगा\appbrk{}र\appbrk{}धू\appbrk{}म\appbrk{}बृ\appbrk{}हती}}\Afootnote{K H; \skt{आगा\appbrk{}र\appbrk{}धू\appbrk{}म\appbrk{}बृ\appbrk{}हती}\ A}}\skt{ }\edtext{\skt{फल\-पि\-प्प\-ली\-सै\-न्धवैः}}{\lemma{\skt{फल\appbrk{}पि\appbrk{}प्प\appbrk{}ली\appbrk{}सै\appbrk{}न्धवैः}}\Afootnote{H; \skt{पिप्प\appbrk{}ली\appbrk{}फ\appbrk{}ल\appbrk{}सै\appbrk{}न्धवैः~।}\ A; \skt{फल\appbrk{}पि\appbrk{}प्प\appbrk{}लि\appbrk{}सै\appbrk{}न्धवै~।}\ K}}\skt{~॥ }\pend
\pstart\versepar \skt{कृता वा शुक्त\-गो\-मू\-त्र\-सु\-रा\-पिष्टैः सनागरैः~।}\pend

\pstart \ledsidenote{\passref{4.37.123}}\edtext{\skt{शर्क्क\-रा\-म\-धु\-मि\-श्रेण}}{\lemma{\skt{शर्क्क\appbrk{}रा\appbrk{}म\appbrk{}धु\appbrk{}मि\appbrk{}श्रेण}}\Afootnote{K H; \skt{अनु\appbrk{}वा\appbrk{}स\appbrk{}न\appbrk{}सिद्धिं च वीक्ष्य कर्म प्रयोजयेत्~। शर्क\appbrk{}रा\appbrk{}म\appbrk{}धु\appbrk{}मि\appbrk{}श्रेण}\ A}}\skt{ शीतेन मधुकाम्बुना~॥}\pend

\pstart \ledsidenote{\passref{4.37.124}}\edtext{\skt{बस्तौ}}{\lemma{\skt{बस्तौ}}\Afootnote{K H; \skt{दह्यमाने तदा बस्तौ}\ A}}\skt{ }\edtext{\edtext{\skt{दह्यति वा}}{\lemma{\skt{दह्यति वा}}\Afootnote{K H; \skt{द}\ A}}\skt{ बाढं पायौ बस्ति\-म्प्र\-दा\-प\-येत्~॥}}{\lemma{\skt{ह्यति… प्रदा}}\Afootnote{\emph{om.}\ A}}\pend

\pstart\versepar \ledsidenote{\passref{4.37.125}}\skt{शुक्रन्दुष्टं }\edtext{\skt{शोणि\-त\-ञ्चा\-ङ्ग\-ना\-ना\-म\-सृ\-ग्द\-र\-न्तस्य}}{\lemma{\skt{असृग्दरन्तस्य}}\Afootnote{K H; \skt{पुष्पोद्रेकं तस्य}\ A}}\skt{ }\edtext{\skt{नाशञ्च}}{\lemma{\skt{नाशञ्च}}\Afootnote{K H; \skt{नाणं च}\ A}}\skt{ कष्टं~। }\pend
\pstart\versepar \skt{}\edtext{\edtext{\skt{मूत्रा\-घा\-त\-म्मू\-त्र\-दो\-षा\-न्प्र\-वृ\-द्धा\-न्योनौ}}{\lemma{\skt{मूत्राघातम्}}\Afootnote{K H; \skt{मूत्राघातान्}\ A}}}{\lemma{\skt{मूत्र\appbrk{}दो\appbrk{}षा\appbrk{}न्प्र\appbrk{}वृ\appbrk{}द्धा\appbrk{}न्योनौ}}\Afootnote{H; \skt{मूत्र\appbrk{}दो\appbrk{}षा\appbrk{}न्प्र\appbrk{}वृ\appbrk{}द्धा\appbrk{}न्योनि}\ A; \skt{मूत्रदोषां प्रवृद्धान्योनौ}\ K}}\skt{ }\edtext{\skt{रोगा\-न\-प\-रा\-याश्च}}{\lemma{\skt{रोगा\appbrk{}न\appbrk{}प\appbrk{}रा\appbrk{}याश्च}}\Afootnote{K H; \skt{व्याधिं संस्थितिं चापरायाः~॥}\ A}}\skt{ }\edtext{\skt{सङ्गां}}{\lemma{\skt{सङ्गां}}\Afootnote{K H; \emph{om.}\ A}}\skt{~॥}\pend

\pstart\versepar \ledsidenote{\passref{4.37.126}}\edtext{\skt{शुक्राघातं}}{\lemma{\skt{शुक्राघातं}}\Afootnote{K H; \skt{शुक्रोत्सेकं}\ A}}\skt{ }\edtext{\edtext{\skt{शर्क्क\-रा\-म\-श्म\-रीञ्च}}{\lemma{\skt{शर्क्कराम्}}\Afootnote{K H; \skt{शर्करा}\ A}}}{\lemma{\skt{अश्मरीञ्च}}\Afootnote{H; \skt{मश्मरीं च}\ A; \skt{अर्श्मरीञ्च}\ K}}\skt{ }\edtext{\skt{शुक्रं}}{\lemma{\skt{शुक्रं}}\Afootnote{H; \skt{शूलं}\ A K}}\skt{ बस्तौ }\edtext{\skt{वंक्षणे}}{\lemma{\skt{वंक्षणे}}\Afootnote{A H; \skt{वंक्षणो}\ K}}\skt{ मेहने च~। }\pend
\pstart\versepar \skt{}\edtext{\skt{घोरा\-न\-न्या\-न्ब\-स्ति\-जा\-श्चा\-पि\-रो\-गा\-न्मे\-हो\-दृते}}{\lemma{\skt{चापि\appbrk{}रो\appbrk{}गा\appbrk{}न्मे\appbrk{}हो\appbrk{}दृते}}\Afootnote{H; \skt{चापि रोगान्हित्वा मेहा}\ A; \skt{चापि रोगां मेहादृते}\ K}}\skt{ }\edtext{\skt{तूत्तरौ}}{\lemma{\skt{तूत्तरौ}}\Afootnote{H; \skt{नुत्तरो}\ A; \skt{तूत्तरो}\ K}}\skt{ हन्ति बस्तिः~॥}\pend

\pstart\versepar \ledsidenote{\passref{4.37.127}}\skt{सम्यग्दत्तस्य }\edtext{\edtext{\skt{विज्ञानम्व्यापदः}}{\lemma{\skt{विज्ञानम्}}\Afootnote{K H; \skt{लिङ्गानि}\ A}}}{\lemma{\skt{व्यापदः}}\Afootnote{A H; \skt{व्यापद}\ K}}\skt{ क्रम }\edtext{\skt{एव च}}{\lemma{\skt{एव च}}\Afootnote{A H; \skt{एव छ~।}\ K}}\skt{~। }\pend
\pstart\versepar \skt{}\edtext{\edtext{\skt{बस्ते\-रु\-त्त\-र\-सं\-ज्ञस्य}}{\lemma{\skt{बस्तेर्}}\Afootnote{A H; \skt{बस्ते}\ K}}}{\lemma{\skt{उत्तरसंज्ञस्य}}\Afootnote{A H; \skt{रुरसंज्ञस्य}\ K}}\skt{ समानं }\edtext{\skt{स्नेहबस्तिनेति}}{\lemma{\skt{स्नेहबस्तिनेति}}\Afootnote{K H; \skt{स्नेहबस्तिना~॥}\ A}}\skt{~॥ }\edtext{\skt{ओ}}{\lemma{\skt{ओ}}\Afootnote{H; \emph{om.}\ A K}}\skt{~॥}\pend

\pstart \ledsidenote{\passref{4.37.trailer}}\edtext{\skt{चिकित्सालञ्}}{\lemma{\skt{चिकित्सालञ्}}\Afootnote{H; \skt{इति सुश्रु\appbrk{}त\appbrk{}सं\appbrk{}हि\appbrk{}तायां चिकि\appbrk{}त्सा\appbrk{}स्था\appbrk{}नेऽ\appbrk{}नु\appbrk{}वा\appbrk{}स\appbrk{}नो\appbrk{}त्त\appbrk{}र\appbrk{}ब\appbrk{}स्ति\appbrk{}चि\appbrk{}कि\appbrk{}त्सितं नाम सप्तविंशोऽध्यायः~॥३७~॥}\ A; \skt{चिलग्र~॥ {\symfont ❈}~॥}\ K}}\skt{~॥}\pend

\endnumbering
\beginnumbering
\eledsection*{Adhyāya 38}
\pstart \ledsidenote{\passref{4.38.1}}\skt{अथातो निरू\-ह\-क्र\-म\-चि\-कि\-त्सितं }\edtext{\skt{व्याख्यास्यामः}}{\lemma{\skt{व्याख्यास्यामः}}\Afootnote{A H; \emph{om.}\ K}}\skt{~॥}\pend

\pstart\versepar \ledsidenote{\passref{4.38.1.add1}}\edtext{\skt{धन्व\-न्त\-रि\-र्द्ध\-र्म्म\-भृ\-ता\-म्व\-रि\-ष्ठ\-म\-मृ\-तो\-द्भवं}}{\lemma{\skt{धन्व\appbrk{}न्त\appbrk{}रि\appbrk{}र्द्ध\appbrk{}र्म्म\appbrk{}भृ\appbrk{}ताम्}}\Afootnote{\skt{धन्व\appbrk{}न्त\appbrk{}रि\appbrk{}ध\appbrk{}र्म\appbrk{}भृतां}\ K; \skt{धन्व\appbrk{}न्त\appbrk{}रि\appbrk{}र्द्ध\appbrk{}र्म्म\appbrk{}भृ\appbrk{}ताम्}\ H}}\skt{~। }\pend
\pstart\versepar \skt{}\edtext{\skt{चर\-णा\-वु\-प\-सं\-ङ्गृह्य}}{\lemma{\skt{चर\appbrk{}णा\appbrk{}वु\appbrk{}प\appbrk{}सं\appbrk{}ङ्गृह्य}}\Afootnote{\skt{चर\appbrk{}णा\appbrk{}वु\appbrk{}प\appbrk{}सं\appbrk{}गृह्य}\ K; \skt{चर\appbrk{}णा\appbrk{}वु\appbrk{}प\appbrk{}सं\appbrk{}ङ्गृह्य}\ H}}\skt{ सुश्रुतः परिपृच्छति~॥}\pend

\pstart\versepar \ledsidenote{\passref{4.38.1.add2}}\skt{भगवं }\edtext{\skt{स्नेहवस्तेस्तु}}{\lemma{\skt{स्नेहवस्तेस्तु}}\Afootnote{H; \skt{वस्तेस्तु}\ K}}\skt{ तथै\-वा\-स्था\-प\-नस्य च~। }\pend
\pstart\versepar \skt{}\edtext{\skt{प्रमा\-ण\-ङ्क\-ल्प\-न\-ङ्का\-ल\-म्मा\-त्रा\-ञ्चापि}}{\lemma{\skt{प्रमा\appbrk{}ण\appbrk{}ङ्क\appbrk{}ल्प\appbrk{}नङ्}}\Afootnote{H; \skt{प्रमानं कल्पनां}\ K}}\skt{ ब्रवीहि मे~॥ }\pend
\pstart\versepar \skt{तस्य }\edtext{\skt{तद्वचनं}}{\lemma{\skt{तद्वचनं}}\Afootnote{K; \skt{तद्वचनं}\ H}}\skt{ श्रुत्वा प्राब्र\-वी\-द्भि\-ष\-जा\-म्वरः~॥}\pend

\pstart\versepar \ledsidenote{\passref{4.38.1.add3}}\skt{शुद्धे मुहूर्त्ते करणे दिने च नक्षत्रयोगे च तथानुकूले~। }\pend
\pstart\versepar \skt{विवि\-क्त\-सं\-सृ\-ष्ट\-मृ\-दा\-व\-लिप्ते विचि\-त्र\-पु\-ष्पा\-स्तृ\-त\-भू\-प्र\-देशे~॥}\pend

\pstart\versepar \ledsidenote{\passref{4.38.1.add4}}\skt{सुवर्ण्णरूप्याय सताम्रजो वा }\edtext{\skt{मृदूद्भवो वा}}{\lemma{\skt{मृदूद्भवो वा}}\Afootnote{H; \skt{मृदुद्भवो वा}\ K}}\skt{ }\edtext{\edtext{\skt{कलशोऽम्बुपूर्ण्णः}}{\lemma{\skt{कलशो}}\Afootnote{H; \skt{कलशोम्}\ K}}}{\lemma{\skt{ऽम्बुपूर्ण्णः}}\Afootnote{H; \skt{बुपूर्ण्णः~।}\ K}}\skt{~। }\pend
\pstart\versepar \skt{सधान्यरत्नादि }\edtext{\skt{समा\-ल्य\-ग\-न्धै\-र्व\-स्त्रा\-दि\-भि\-श्चा\-प्यथ}}{\lemma{\skt{वस्त्रादिभिश्}}\Afootnote{H; \skt{वस्त्रादिभिश्च}\ K}}\skt{ स्वस्तिकैश्च~॥}\pend

\pstart\versepar \ledsidenote{\passref{4.38.1.add5}}\edtext{\edtext{\edtext{\skt{तच्चा\-मृ\-त\-स्था\-न\-म\-तोऽ\-र्च्च\-नीयं}}{\lemma{\skt{तच्चा\appbrk{}मृ\appbrk{}त\appbrk{}स्था\appbrk{}नम्}}\Afootnote{H; \skt{तच्चामृतस्थान}\ K}}}{\lemma{\skt{अतो}}\Afootnote{H; \skt{मतो}\ K}}}{\lemma{\skt{ऽर्च्चनीयं}}\Afootnote{H; \skt{र्चनीयं}\ K}}\skt{ }\edtext{\skt{पूर्व्वा\-त्त\-तो\-न्य\-द्धृ\-षयः}}{\lemma{\skt{पूर्व्वा\appbrk{}त्त\appbrk{}तो\appbrk{}न्य\appbrk{}द्धृ\appbrk{}षयः}}\Afootnote{H; \skt{पूर्व\appbrk{}त्त\appbrk{}तो\appbrk{}न्ये\appbrk{}द्दृ\appbrk{}षयः}\ K}}\skt{ सदेवाः~। }\pend
\pstart\versepar \skt{संस्पृश्य रत्नानि तथौषधीश्च }\edtext{\skt{देवान्द्विजान्पूज्य च}}{\lemma{\skt{द्विजान्पूज्य च}}\Afootnote{H; \skt{द्विजांन्पूज्य च}\ K}}\skt{ दक्षिणाभिः~॥}\pend

\pstart\versepar \ledsidenote{\passref{4.38.1.add6}}\skt{द्रवाणि सर्व्वाणि तु वामपार्श्वे चूर्ण्णानि कल्कानथ दक्षिणे तु~। }\pend
\pstart\versepar \skt{तन्मध्य }\edtext{\edtext{\edtext{\skt{संस्था\-य\-भि\-ष\-क्प्र\-सन्न}}{\lemma{\skt{संस्था\appbrk{}य\appbrk{}भि\appbrk{}ष\appbrk{}क्प्र\appbrk{}सन्न}}\Afootnote{H\textsuperscript{pc}; \skt{संस्था}\ K H\textsuperscript{ac}}}\skt{ इति पाठः~।भि\-ष\-ग\-प्र\-मत्तः सुनिर्म्मले}}{\lemma{\skt{यभिषक्…~।भि}}\Afootnote{\emph{om.}\ K}}}{\lemma{\skt{यभिषक्…~।भि}}\Afootnote{\emph{om.}\ H\textsuperscript{ac}}}\skt{ }\edtext{\skt{स्यात्सुदृढो}}{\lemma{\skt{स्यात्सुदृढो}}\Afootnote{H; \skt{स्यात्सुदृढे}\ K}}\skt{ विशाले~॥ }\pend
\pstart\versepar \skt{सुवर्ण्णलोहादि }\edtext{\skt{मृदूद्भवो वा}}{\lemma{\skt{मृदूद्भवो वा}}\Afootnote{H; \skt{मृदुद्भवो वा}\ K}}\skt{ }\edtext{\skt{प्रगृ\-ह्य\-कां\-सेऽ\-प्यथ}}{\lemma{\skt{ऽप्यथ}}\Afootnote{\skt{प्यथ}\ K; \skt{ऽप्यथ}\ H}}\skt{ राजते वा~॥}\pend

\pstart\versepar \ledsidenote{\passref{4.38.1.add7}}\skt{दत्वादौ }\edtext{\skt{सैन्धवस्याक्षं}}{\lemma{\skt{सैन्धवस्याक्षं}}\Afootnote{H; \skt{सेन्धस्याक्षं}\ K}}\skt{ मधुनः प्रसृतम्मितं~। }\pend
\pstart\versepar \skt{पात्रे तलेन मथ्नी\-या\-त्त\-द्व\-त्स्नेहं शनैः शनैः~॥}\pend

\pstart\versepar \ledsidenote{\passref{4.38.1.add8}}\skt{गूढ\-स्ने\-ह\-म्वि\-दित्वा तु विनीतः कल्कमावपेत्~। }\pend
\pstart\versepar \skt{दत्वा सम्यक्सुमथिते फलकल्कमतः परं~॥}\pend

\pstart\versepar \ledsidenote{\passref{4.38.1.add9}}\skt{ततो यथो\-चि\-ता\-ञ्चे\-षा\-न्का\-र्षि\-का\-न्श्ल\-क्ष्ण\-पे\-षि\-मान्~। }\pend
\pstart\versepar \skt{गम्भीरे भोजने न्यस्य खजे\-ना\-भि\-म\-थे\-त्ततः~॥}\edtext{}{\lemma{\skt{चेषा\appbrk{}न्का\appbrk{}र्षि\appbrk{}का\appbrk{}न्श्ल\appbrk{}क्ष्ण\appbrk{}पे\appbrk{}षि\appbrk{}मान्~। गम्भीरे}}\Afootnote{H; \skt{केषां कार्षि\appbrk{}का\appbrk{}न्श्ल\appbrk{}क्ष्ण\appbrk{}पे\appbrk{}षिमां~। गम्भीरे}\ K}}\pend

\pstart\versepar \ledsidenote{\passref{4.38.1.add10}}\skt{यथा च साधु मन्येत न सान्द्रो न तनुः समः~। }\pend
\pstart\versepar \skt{कषायः प्रसृतात्पञ्च सुपूर्ण्णात्तत्र दापयेत्~॥}\pend

\pstart\versepar \ledsidenote{\passref{4.38.1.add11}}\edtext{\edtext{\skt{रस\-क्षी\-रा\-म्ल\-मू\-त्रा\-णा\-न्दो\-षा\-व\-स्था\-म\-वेक्ष्य}}{\lemma{\skt{दोषावस्थाम्}}\Afootnote{H; \skt{दोषवस्था}\ K}}\skt{ च}}{\lemma{\skt{अवेक्ष्य च}}\Afootnote{H; \skt{नवेक्ष्य च~।}\ K}}\skt{~। }\pend
\pstart\versepar \skt{एवमञ्जलयः }\edtext{\skt{षड्तैः}}{\lemma{\skt{षड्तैः}}\Afootnote{H; \skt{षड्वै}\ K}}\skt{ कुर्या\-द्व्या\-धि\-त\-पा\-णिना~॥}\pend

\pstart\versepar \ledsidenote{\passref{4.38.1.add12}}\skt{द्वादशप्रसृतं केचिद्विङ्शतिः पलिकन्तथा~। }\pend
\pstart\versepar \skt{}\edtext{\edtext{\skt{द्वात्रि\-ङ्श\-त्प\-लि\-कां\-श्चैव}}{\lemma{\skt{द्वात्रि\appbrk{}ङ्श\appbrk{}त्प\appbrk{}लि\appbrk{}कांश्}}\Afootnote{\skt{द्वात्रिं\appbrk{}श\appbrk{}त्प\appbrk{}लि\appbrk{}कञ्}\ K; \skt{द्वात्रिङ्शत्पलिका}\ H}}}{\lemma{\skt{चैव}}\Afootnote{K; \skt{श्चैव}\ H}}\skt{ त्रिंशदष्टौ तथैव च~॥ }\pend
\pstart\versepar \skt{एवम्भिन्नेषु शास्त्रेषु संशयः }\edtext{\skt{सुमहानभूत्~॥ ओ}}{\lemma{\skt{सुमहानभूत्~॥ ओ}}\Afootnote{H; \skt{सुमहानभूत्~॥ {\symfont ❈}~॥}\ K}}\skt{~॥}\pend

\pstart \ledsidenote{\passref{4.38.3}}\edtext{\edtext{\edtext{\skt{आतु\-र\-म\-नु\-वा\-सि\-त\-मा\-स्था\-प\-येत्}}{\lemma{\skt{आतुरम्}}\Afootnote{H; \skt{अथा}\ A; \skt{अथातुर}\ K}}}{\lemma{\skt{अनुवासितम्}}\Afootnote{H; \skt{नुवासितम्}\ A; \skt{मनुवासितम्}\ K}}\skt{~। स्वभ्यक्तं}}{\lemma{\skt{आस्थापयेत्~। स्वभ्यक्तं}}\Afootnote{K H; \skt{आस्थापयेत्; स्वभ्यक्त}\ A}}\skt{ }\edtext{\edtext{\skt{स्विन्न\-श\-री\-र\-मु\-त्सृष्ट}}{\lemma{\skt{स्विन्नशरीरम्}}\Afootnote{K H; \skt{स्विन्नशरीर}\ A}}}{\lemma{\skt{उत्सृष्ट}}\Afootnote{K H; \skt{मुत्सृष्ट}\ A}}\skt{ }\edtext{\skt{बहि\-र्ब्भा\-ग\-म\-प्र\-वाते}}{\lemma{\skt{बहि\appbrk{}र्ब्भा\appbrk{}ग\appbrk{}म\appbrk{}प्र\appbrk{}वाते}}\Afootnote{K H; \skt{बहि\appbrk{}र्वे\appbrk{}ग\appbrk{}म\appbrk{}वाते}\ A}}\skt{ शुचौ }\edtext{\skt{देशे}}{\lemma{\skt{देशे}}\Afootnote{K H; \skt{वेश्म}\ A}}\skt{ }\edtext{\skt{समविन्यस्तायां}}{\lemma{\skt{समविन्यस्तायां}}\Afootnote{K H; \skt{नि मध्याह्ने प्रततायां}\ A}}\skt{ }\edtext{\edtext{\skt{शय्यायामधः}}{\lemma{\skt{शय्यायाम्}}\Afootnote{K H; \skt{शय्याया}\ A}}}{\lemma{\skt{अधः}}\Afootnote{K H; \skt{मधः}\ A}}\skt{~। }\edtext{\edtext{\skt{सप्र\-ति\-ग्र\-हा\-या\-मा\-श्रो\-णी\-प्र\-देशे}}{\lemma{\skt{सप्र\appbrk{}ति\appbrk{}ग्र\appbrk{}हा\appbrk{}याम्}}\Afootnote{K H; \skt{सुपरिग्रहायां}\ A}}}{\lemma{\skt{आश्रोणीप्रदेशे}}\Afootnote{K H; \skt{श्रोणिप्रदेश}\ A}}\skt{ }\edtext{\skt{सुव्यूहायां}}{\lemma{\skt{सुव्यूहायां}}\Afootnote{K H; \skt{प्रतिव्यूढाया}\ A}}\skt{ }\edtext{\edtext{\skt{शया\-न\-म्वा\-म\-पार्श्वे}}{\lemma{\skt{शयानम्}}\Afootnote{K H; \skt{मनुपधानायां}\ A}}}{\lemma{\skt{वामपार्श्वे}}\Afootnote{K H; \skt{वामपार्श्व}\ A}}\skt{ }\edtext{\skt{नेति}}{\lemma{\skt{नेति}}\Afootnote{H; \skt{शायि}\ A; \emph{om.}\ K}}\skt{~॥ }\edtext{\skt{भवति}}{\lemma{\skt{भवति}}\Afootnote{H; \skt{नमा\appbrk{}कु\appbrk{}ञ्चि\appbrk{}त\appbrk{}दक्षि}\ A; \skt{नेति~॥}\ K}}\skt{ }\edtext{\skt{चात्र}}{\lemma{\skt{चात्र}}\Afootnote{H; \skt{णस}\ A; \skt{भ~॥}\ K}}\skt{ }\edtext{\skt{श्लोकः}}{\lemma{\skt{श्लोकः}}\Afootnote{H; \skt{क्थिमि}\ A; \emph{om.}\ K}}\skt{~॥ }\edtext{\edtext{\skt{सव्य\-म्प्र\-सा\-र\-ये\-त्सक्थि}}{\lemma{\skt{सव्यम्}}\Afootnote{K H; \skt{तर}\ A}}}{\lemma{\skt{प्रसारयेत्सक्थि}}\Afootnote{H; \skt{प्रसारितसक्थिं}\ A; \skt{प्रसारयेत्सक्थिं}\ K}}\skt{ दक्षिणं चापसंहरेत्~। }\edtext{\skt{मध्याह्ने }\edtext{\skt{सुमनाजीर्ण्णे}}{\lemma{\skt{सुमनाजीर्ण्णे}}\Afootnote{K H; \skt{सुमनसं जी}\ A}}\skt{ }\edtext{\skt{निरन्नो}}{\lemma{\skt{निरन्नो}}\Afootnote{K H; \skt{र्णान्नं}\ A}}\skt{ }\edtext{\skt{वाग्य\-तो\-न\-रः\-यतः}}{\lemma{\skt{वाग्य\appbrk{}तो\appbrk{}न\appbrk{}रः\appbrk{}यतः}}\Afootnote{H\textsuperscript{pc}; \skt{वाग्यतं सुनिषण्णदेहं विदित्वा,}\ A; \skt{वाग्यतो यतः~॥}\ K; \skt{वाग्यतोयतः~॥}\ H\textsuperscript{ac}}}}{\lemma{\skt{दक्षिणं… सु}}\Afootnote{\emph{om.}\ A}}\skt{~॥ ततो }\edtext{\skt{वामपादोपरि}}{\lemma{\skt{वामपादोपरि}}\Afootnote{K H; \skt{वाम\appbrk{}पा\appbrk{}द\appbrk{}स्यो\appbrk{}परि}\ A}}\skt{ नेत्रं }\edtext{\edtext{\skt{कृत्वे\-त\-र\-पा\-दा\-ङ्गु\-ष्ठा\-ङ्गु\-ली\-भ्या\-ङ्क\-र्ण्णि\-को\-परि}}{\lemma{\skt{कृत्वे\appbrk{}त\appbrk{}र\appbrk{}पा\appbrk{}दा\appbrk{}ङ्गु\appbrk{}ष्ठा\appbrk{}ङ्गु\appbrk{}ली\appbrk{}भ्याङ्}}\Afootnote{K H; \skt{कृत्वे\appbrk{}त\appbrk{}र\appbrk{}पा\appbrk{}दा\appbrk{}ङ्गु\appbrk{}ष्ठा\appbrk{}ङ्गु\appbrk{}लिभ्यां}\ A}}}{\lemma{\skt{कर्ण्णिकोपरि}}\Afootnote{H\textsuperscript{pc}; \skt{कर्णिकामुपरि निष्पीड्य, सव्य\appbrk{}पा\appbrk{}णि\appbrk{}क\appbrk{}नि\appbrk{}ष्ठि\appbrk{}का\appbrk{}ना\appbrk{}मि\appbrk{}काभ्यां बस्तेर्मुखार्धं सङ्कोच्य, मध्य\appbrk{}मा\appbrk{}प्र\appbrk{}दे\appbrk{}शि\appbrk{}न्य\appbrk{}ङ्गु\appbrk{}ष्ठै\appbrk{}रर्धं तु विवृतास्यं कृत्वा, बस्तावौषधं प्रक्षिप्य, दक्षि\appbrk{}ण\appbrk{}ह\appbrk{}स्ता\appbrk{}ङ्गु\appbrk{}ष्ठेन प्रदे\appbrk{}शि\appbrk{}नी\appbrk{}म\appbrk{}ध्य\appbrk{}माभ्यां चानु\appbrk{}सि\appbrk{}क्त\appbrk{}म\appbrk{}ना\appbrk{}य\appbrk{}त\appbrk{}म\appbrk{}बु\appbrk{}द्बु\appbrk{}द\appbrk{}म\appbrk{}स\appbrk{}ङ्कु\appbrk{}चि\appbrk{}त\appbrk{}म\appbrk{}वा\appbrk{}त\appbrk{}मौ\appbrk{}ष\appbrk{}धा\appbrk{}स\appbrk{}न्न\appbrk{}मु\appbrk{}प\appbrk{}सं\appbrk{}गृह्य, पुनरुपरि}\ A; \skt{कर्ण्णि\appbrk{}को\appbrk{}प\appbrk{}रि\appbrk{}पीड्य बस्तावौष}\ K; \skt{कर्ण्णिकोपरि निपीड्य}\ H\textsuperscript{ac}}}\skt{ }\edtext{\skt{निपीड्य }\edtext{\skt{बस्ता\-वौ\-ष\-ध\-मा\-सि\-च्यो\-च्छ्रा\-स्यौ\-ष\-धान्ते}}{\lemma{\skt{औष\appbrk{}ध\appbrk{}मा\appbrk{}सि\appbrk{}च्यो\appbrk{}च्छ्रा\appbrk{}स्यौ\appbrk{}ष\appbrk{}धान्ते}}\Afootnote{H; \skt{णेनावसिञ्चेत्, ततः सूत्रे\appbrk{}णै\appbrk{}वौ\appbrk{}ष\appbrk{}धान्ते}\ A}}}{\lemma{\skt{निपीड्य बस्ताव्}}\Afootnote{H\textsuperscript{pc}; \skt{तदितरेण गृहीत्वा दक्षि}\ A; \skt{बस्ताव्}\ H\textsuperscript{ac}}}\skt{ }\edtext{\skt{चाववध्य मूत्रेण}}{\lemma{\skt{चाववध्य मूत्रेण}}\Afootnote{K H; \skt{द्विस्त्रिर्वाऽऽआवेष्ट्य बध्नीयात्, अथ दक्षि\appbrk{}णे\appbrk{}नो\appbrk{}त्ता\appbrk{}नेन पाणिना बस्तिं गृहीत्वा वाम\appbrk{}ह\appbrk{}स्त\appbrk{}म\appbrk{}ध्य\appbrk{}मा\appbrk{}ङ्गु\appbrk{}लि\appbrk{}प्र\appbrk{}दे\appbrk{}शि\appbrk{}नीभ्यां नेत्र\appbrk{}मु\appbrk{}प\appbrk{}सं\appbrk{}गृ\appbrk{}ह्या\appbrk{}ङ्गु\appbrk{}ष्टेन}\ A}}\skt{ }\edtext{\edtext{\skt{सम्य\-ग्ब\-द्ध\-ञ्चा\-तुरा}}{\lemma{\skt{सम्यग्बद्धञ्}}\Afootnote{K H; \skt{नेत्रद्वारं}\ A}}\skt{ य}}{\lemma{\skt{चातुरा य}}\Afootnote{K H; \skt{पिधाय,}\ A}}\skt{ }\edtext{\skt{घृता\-भ्य\-क्त\-ने\-त्राग्रं}}{\lemma{\skt{घृता\appbrk{}भ्य\appbrk{}क्त\appbrk{}ने\appbrk{}त्राग्रं}}\Afootnote{H; \skt{घृता\appbrk{}भ्य\appbrk{}क्ता\appbrk{}ग्र\appbrk{}नेत्रं}\ A; \skt{घृताभ्यक्तं नेत्राग्रं}\ K}}\skt{ }\edtext{\skt{घृताभ्यक्त}}{\lemma{\skt{घृताभ्यक्त}}\Afootnote{K H; \skt{घृताक्तमु}\ A}}\skt{ }\edtext{\skt{गुदाय}}{\lemma{\skt{गुदाय}}\Afootnote{K H; \skt{पादाय}\ A}}\skt{ }\edtext{\edtext{\edtext{\edtext{\skt{प्रण\-ये\-द\-नु\-त्सृ\-ज\-न्न\-नु\-पृ\-ष्ठ\-वं\-श\-म\-नु\-सु\-ख\-मा\-क\-र्ण्णि\-क\-मिति}}{\lemma{\skt{प्रण\appbrk{}ये\appbrk{}द\appbrk{}नु\appbrk{}त्सृ\appbrk{}ज\appbrk{}न्न\appbrk{}नु\appbrk{}पृ\appbrk{}ष्ठ\appbrk{}वं\appbrk{}शम्}}\Afootnote{K H; \skt{प्रय\appbrk{}च्छे\appbrk{}द\appbrk{}नु\appbrk{}पृ\appbrk{}ष्ठ\appbrk{}वंशं}\ A}}}{\lemma{\skt{अनुसुखम्}}\Afootnote{K H; \skt{सममुन्मुख}\ A}}}{\lemma{\skt{आकर्ण्णिकम्}}\Afootnote{K H; \skt{माकर्णिकं}\ A}}}{\lemma{\skt{इति}}\Afootnote{K H; \skt{नेत्रं प्रणिधत्स्वेति ब्रूयात्~॥}\ A}}\skt{~॥}\edtext{}{\lemma{\skt{परि… औष}}\Afootnote{\emph{om.}\ K}}\pend

\pstart\versepar \ledsidenote{\passref{4.38.4}}\skt{बस्तिं }\edtext{\skt{सव्ये}}{\lemma{\skt{सव्ये}}\Afootnote{A H; \skt{सव्य}\ K}}\skt{ करे कृत्वा }\edtext{\skt{दक्षिणेन}}{\lemma{\skt{दक्षिणेन}}\Afootnote{H; \skt{दक्षिणेना}\ A K}}\skt{ अवपीडयेत्~। }\pend
\pstart\versepar \skt{एकेनैव च पीडेन न द्रुतन्न विलम्बितं~॥}\edtext{}{\lemma{\skt{अवपीडयेत्~। एकेनैव च}}\Afootnote{H; \skt{वपीडयेत्~। एकेनैवाव}\ A; \skt{वपीडयेत्~। एकेनैव च}\ K}}\pend

\pstart \ledsidenote{\passref{4.38.5}}\skt{ततो }\edtext{\edtext{\skt{नेत्रमपनीय}}{\lemma{\skt{नेत्रम्}}\Afootnote{K H; \skt{नेत्र}\ A}}}{\lemma{\skt{अपनीय}}\Afootnote{K H; \skt{मपनीय}\ A}}\skt{ }\edtext{\edtext{\skt{त्रिंश\-न्मा\-त्रा\-पी\-ड\-न\-का\-ला\-दु\-पे\-त्यो\-त्ति\-ष्ठे\-त्या\-तु\-र\-म्ब्रू\-यात्}}{\lemma{\skt{त्रिंश\appbrk{}न्मा\appbrk{}त्रा\appbrk{}पी\appbrk{}ड\appbrk{}न\appbrk{}का\appbrk{}ला\appbrk{}दु\appbrk{}पे\appbrk{}त्यो\appbrk{}त्ति\appbrk{}ष्ठे\appbrk{}त्या\appbrk{}तु\appbrk{}रम्}}\Afootnote{K H; \skt{त्रिंशन्मात्राः पीड\appbrk{}न\appbrk{}का\appbrk{}ला\appbrk{}दु\appbrk{}पे\appbrk{}क्ष्यो\appbrk{}त्ति\appbrk{}ष्ठे\appbrk{}त्या\appbrk{}तुरं}\ A}}\skt{~। उत्कु\-टु\-को\-भ\-व\-त\-स्ते\-रा\-ग\-म\-ना\-येति}}{\lemma{\skt{ब्रूयात्~। उत्कु\appbrk{}टु\appbrk{}को\appbrk{}भ\appbrk{}व\appbrk{}त\appbrk{}स्ते\appbrk{}रा\appbrk{}ग\appbrk{}म\appbrk{}ना\appbrk{}येति}}\Afootnote{K H; \skt{ब्रूयात्~। अथा\appbrk{}तु\appbrk{}र\appbrk{}मु\appbrk{}प\appbrk{}वे\appbrk{}श\appbrk{}ये\appbrk{}दु\appbrk{}त्कु\appbrk{}टुकं बस्त्या\appbrk{}ग\appbrk{}म\appbrk{}ना\appbrk{}र्थम्~। निरू\appbrk{}ह\appbrk{}प्र\appbrk{}त्या\appbrk{}ग\appbrk{}म\appbrk{}न\appbrk{}का\appbrk{}लस्तु मुहूर्तो भवति~॥}\ A}}\skt{~।}\pend

\pstart\versepar \ledsidenote{\passref{4.38.6}}\skt{अनेन विधिना }\edtext{\skt{दद्याद्बस्तिं}}{\lemma{\skt{दद्याद्बस्तिं}}\Afootnote{K H; \skt{बस्तिं दद्याद्ब}\ A}}\skt{ }\edtext{\skt{बस्तिविशारदः}}{\lemma{\skt{बस्तिविशारदः}}\Afootnote{K H; \skt{स्तिविशारदः~।}\ A}}\skt{~॥ }\pend
\pstart\versepar \skt{द्वितीयम्वा तृतीयम्वा चतुर्थम्वा यथार्थतः~॥}\pend

\pstart \ledsidenote{\passref{4.38.7}}\skt{सम्य\-ङ्नि\-रू\-ढ\-लिङ्गे तु }\edtext{\edtext{\skt{प्राप्त\-ब\-स्ति\-न्नि\-वा\-र\-येत्}}{\lemma{\skt{प्राप्तबस्तिन्}}\Afootnote{K H; \skt{प्राओते बस्तिं}\ A}}}{\lemma{\skt{निवारयेत्}}\Afootnote{K H; \skt{निवारयेत्~। विशे\appbrk{}षा\appbrk{}त्सु\appbrk{}कु\appbrk{}मा\appbrk{}राणां हीन एव क्रमो हितः~॥}\ A}}\skt{~।}\pend

\pstart\versepar \ledsidenote{\passref{4.38.8}}\skt{अपि }\edtext{\skt{हीनं}}{\lemma{\skt{हीनं}}\Afootnote{K H; \skt{हीन}\ A}}\skt{ }\edtext{\skt{क्रमङ्कुर्यान्न च}}{\lemma{\skt{कुर्यान्न च}}\Afootnote{K H; \skt{कुर्यान्न तु}\ A}}\skt{ कुर्यादतिक्रमं~॥ }\pend
\pstart\versepar \skt{विशे\-षा\-त्सु\-कु\-मा\-राणां हीन एव }\edtext{\skt{क्रमो}}{\lemma{\skt{क्रमो}}\Afootnote{K\textsuperscript{pc} H\textsuperscript{pc}; \skt{स्य स्या}\ A}}\skt{ }\edtext{\skt{हितः}}{\lemma{\skt{हितः}}\Afootnote{K\textsuperscript{pc} H\textsuperscript{pc}; \skt{द्बस्तिर्}\ A}}\skt{ }\pend
\pstart\versepar \skt{}\edtext{\skt{अल्पाल्पवेगो}}{\lemma{\skt{अल्पाल्पवेगो}}\Afootnote{K H; \skt{अल्पोऽल्पवेगो}\ A}}\skt{ }\edtext{\skt{विन्वात}}{\lemma{\skt{विन्वात}}\Afootnote{K H; \emph{om.}\ A}}\skt{ }\edtext{\skt{हीनोहीनो}}{\lemma{\skt{हीनोहीनो}}\Afootnote{K H; \skt{हीन}\ A}}\skt{ }\edtext{\skt{निरूहणः}}{\lemma{\skt{निरूहणः}}\Afootnote{H; \skt{मलानिलः~॥}\ A; \skt{निरूहनः~।}\ K}}\skt{~॥}\edtext{}{\lemma{\skt{विशेषात्… अ}}\Afootnote{\emph{om.}\ H\textsuperscript{ac}}}\edtext{}{\lemma{\skt{विशेषात्… अ}}\Afootnote{\emph{om.}\ K\textsuperscript{ac}}}\edtext{}{\lemma{\skt{विशेषात्… एव}}\Afootnote{\emph{om.}\ A}}\pend

\pstart\versepar \ledsidenote{\passref{4.38.8.1}}\edtext{\skt{मूर्च्छा\-शू\-ल\-क\-फ\-प्रायो}}{\lemma{\skt{मूर्च्छा\appbrk{}शू\appbrk{}ल\appbrk{}क\appbrk{}फ\appbrk{}प्रायो}}\Afootnote{K\textsuperscript{pc} H\textsuperscript{pc}; \emph{om.}\ K\textsuperscript{ac}}}\skt{ }\edtext{\skt{महा\-वे\-गा\-ति\-श\-ब्दितः}}{\lemma{\skt{महा\appbrk{}वे\appbrk{}गा\appbrk{}ति\appbrk{}श\appbrk{}ब्दितः}}\Afootnote{K\textsuperscript{pc} H\textsuperscript{pc}; \emph{om.}\ K\textsuperscript{ac}}}\skt{~। }\pend
\pstart\versepar \skt{}\edtext{\skt{यस्य}}{\lemma{\skt{यस्य}}\Afootnote{K\textsuperscript{pc} H; \emph{om.}\ K\textsuperscript{ac}}}\skt{ }\edtext{\skt{मूत्रपुरीषञ्च}}{\lemma{\skt{मूत्रपुरीषञ्च}}\Afootnote{K\textsuperscript{pc} H; \emph{om.}\ K\textsuperscript{ac}}}\skt{ }\edtext{\skt{कफवायुश्च}}{\lemma{\skt{कफवायुश्च}}\Afootnote{K\textsuperscript{pc} H; \emph{om.}\ K\textsuperscript{ac}}}\skt{ }\edtext{\skt{गच्छति}}{\lemma{\skt{गच्छति}}\Afootnote{K\textsuperscript{pc} H; \emph{om.}\ K\textsuperscript{ac}}}\skt{~॥}\edtext{}{\lemma{\skt{… य}}\Afootnote{\emph{om.}\ H\textsuperscript{ac}}}\pend

\pstart\versepar \ledsidenote{\passref{4.38.8.2}}\skt{क्रमेण लघुता चैव }\edtext{\skt{सुनिरूहः}}{\lemma{\skt{सुनिरूहः}}\Afootnote{H\textsuperscript{pc}; \skt{सुनिरढः}\ K\textsuperscript{pc}}}\skt{ समानवः~। }\pend
\pstart\versepar \skt{सुनिरूढन्ततो जन्तुं स्नातं भुक्तरसोदनं~॥}\edtext{}{\lemma{\skt{… सु}}\Afootnote{\emph{om.}\ K\textsuperscript{ac}}}\edtext{}{\lemma{\skt{… सु}}\Afootnote{\emph{om.}\ H\textsuperscript{ac}}}\pend

\pstart\versepar \ledsidenote{\passref{4.38.8.3}}\skt{यथोक्तेन विधानेन योज\-ये\-त्स्ने\-ह\-ब\-स्तिना~। }\pend
\pstart\versepar \skt{तदहस्तस्य पव\-ना\-द्भ\-य\-म्ब\-ल\-व\-दि\-ष्यते~॥}\edtext{}{\lemma{\skt{… त}}\Afootnote{\emph{om.}\ K\textsuperscript{ac}}}\edtext{}{\lemma{\skt{… त}}\Afootnote{\emph{om.}\ H\textsuperscript{ac}}}\pend

\pstart\versepar \ledsidenote{\passref{4.38.8.4}}\edtext{\skt{रसोदनस्तेन}}{\lemma{\skt{रसोदनस्तेन}}\Afootnote{H; \skt{रसोदनंस्तेन}\ K}}\skt{ शस्त\-स्त\-द\-ह\-श्चा\-नु\-वा\-सनं~। }\pend
\pstart\versepar \skt{पश्चा\-द\-ग्नि\-ब\-ल\-म्मत्वा पवनस्य च }\edtext{\skt{वेष्टितं}}{\lemma{\skt{वेष्टितं}}\Afootnote{H; \skt{चेष्टितं~।}\ K}}\skt{~॥}\pend

\pstart \ledsidenote{\passref{4.38.8.5}}\skt{अन्नोपस्तम्भिते काष्ठे }\edtext{\skt{स्नेह\-ब\-स्ति\-र्व्वि\-धी\-यते}}{\lemma{\skt{स्नेह\appbrk{}ब\appbrk{}स्ति\appbrk{}र्व्वि\appbrk{}धी\appbrk{}यते}}\Afootnote{H; \skt{स्नेह\appbrk{}ब\appbrk{}स्ति\appbrk{}र्वि\appbrk{}धी\appbrk{}यते~॥ {\symfont ❈}~॥}\ K}}\skt{~॥ }\edtext{\skt{ओ}}{\lemma{\skt{ओ}}\Afootnote{H; \emph{om.}\ K}}\skt{~॥}\pend

\pstart\versepar \ledsidenote{\passref{4.38.8.6}}\edtext{\skt{विवि\-क्त\-ता\-म\-न\-स्तुष्टि}}{\lemma{\skt{विवि\appbrk{}क्त\appbrk{}ता\appbrk{}म\appbrk{}न\appbrk{}स्तुष्टि}}\Afootnote{H; \skt{विवि\appbrk{}क्ता\appbrk{}म\appbrk{}न\appbrk{}स्तुष्टि}\ K}}\skt{ स्निग्धता व्याधिनिग्रहः~॥ }\pend
\pstart\versepar \skt{आस्था\-प\-न\-स्ने\-ह\-बस्त्योः सम्यग्दाने तु लक्षणं~॥}\pend

\pstart \ledsidenote{\passref{4.38.17}}\edtext{\skt{अना\-या\-न्त\-म्मु\-हू\-र्त्तात्तु}}{\lemma{\skt{अनायान्तम्}}\Afootnote{K H; \skt{अन्नोपस्तम्भिते कोष्ठे स्नेह\appbrk{}ब\appbrk{}स्ति\appbrk{}र्वि\appbrk{}धी\appbrk{}यते~॥ अनायान्तं}\ A}}\skt{ निरूहं शोधनैर्हरेत्~।}\pend

\pstart\versepar \ledsidenote{\passref{4.38.18}}\edtext{\edtext{\skt{निरूहैरेव}}{\lemma{\skt{निरूहैर्}}\Afootnote{K H; \skt{तीक्ष्णैर्निरूहैर्}\ A}}}{\lemma{\skt{एव}}\Afootnote{K H; \emph{om.}\ A}}\skt{ }\edtext{\skt{मति\-मा\-न्क्षा\-र\-मू\-त्रा\-म्ल\-सं\-युतैः}}{\lemma{\skt{मति\appbrk{}मा\appbrk{}न्क्षा\appbrk{}र\appbrk{}मू\appbrk{}त्रा\appbrk{}म्ल\appbrk{}सं\appbrk{}युतैः}}\Afootnote{A H; \skt{मतिमां क्षार\appbrk{}मू\appbrk{}त्रां\appbrk{}म्ल\appbrk{}सं\appbrk{}युतैः~।}\ K}}\skt{~॥ }\pend
\pstart\versepar \skt{}\edtext{\edtext{\skt{विगु\-णा\-नि\-ल\-वि\-ष्ट\-ब्ध\-श्चि\-र\-न्ति\-ष्ठ\-न्नि\-रू\-हणः}}{\lemma{\skt{विगु\appbrk{}णा\appbrk{}नि\appbrk{}ल\appbrk{}वि\appbrk{}ष्ट\appbrk{}ब्धश्}}\Afootnote{K H; \skt{विगु\appbrk{}णा\appbrk{}नि\appbrk{}ल\appbrk{}वि\appbrk{}ष्टब्धं}\ A}}}{\lemma{\skt{निरूहणः}}\Afootnote{K H; \skt{निरूहणम्~॥}\ A}}\skt{~।}\pend

\pstart\versepar \ledsidenote{\passref{4.38.19}}\edtext{\skt{शूला\-र\-ति\-ज्व\-रा\-टो\-पा\-त्म\-र\-णम्वा}}{\lemma{\skt{शूला\appbrk{}र\appbrk{}ति\appbrk{}ज्व\appbrk{}रा\appbrk{}टो\appbrk{}पा\appbrk{}त्म\appbrk{}र\appbrk{}णम्वा}}\Afootnote{H; \skt{शूला\appbrk{}र\appbrk{}ति\appbrk{}ज्व\appbrk{}रा\appbrk{}ना\appbrk{}हा\appbrk{}न्म\appbrk{}रणं वा}\ A; \skt{शूला\appbrk{}र\appbrk{}ति\appbrk{}र्ज्व\appbrk{}रा\appbrk{}टो\appbrk{}पा\appbrk{}त्म\appbrk{}र\appbrk{}णम्वा}\ K}}\skt{ प्रवर्त्तयेत्~॥ }\pend
\pstart\versepar \skt{न तु }\edtext{\skt{भुक्तवते}}{\lemma{\skt{भुक्तवते}}\Afootnote{K H; \skt{भुक्तवतो}\ A}}\skt{ देय\-मा\-स्था\-प\-न\-मिति स्थितिः~।}\pend

\pstart\versepar \ledsidenote{\passref{4.38.20}}\edtext{\skt{आमं हि}}{\lemma{\skt{आमं हि}}\Afootnote{K H; \emph{om.}\ A}}\skt{ }\edtext{\skt{तधरेद्भुक्तं}}{\lemma{\skt{तधरेद्भुक्तं}}\Afootnote{H; \skt{वि}\ A; \skt{तद्धरेद्भुक्तं}\ K}}\skt{ }\edtext{\skt{च्छर्द्दिम्वा}}{\lemma{\skt{च्छर्द्दिम्वा}}\Afootnote{K H; \skt{सूचिकां वा}\ A}}\skt{ }\edtext{\skt{जनयेद्भृशं}}{\lemma{\skt{जनयेद्भृशं}}\Afootnote{H; \skt{जनयेच्छर्दिं वाऽपि सुदारुणाम्~॥}\ A; \skt{जनयेद्भृषं~।}\ K}}\skt{~॥ }\pend
\pstart\versepar \skt{}\edtext{\skt{कोप\-ये\-त्स\-र्व्व\-दो\-षान्वा}}{\lemma{\skt{कोप\appbrk{}ये\appbrk{}त्स\appbrk{}र्व्व\appbrk{}दो\appbrk{}षान्वा}}\Afootnote{A H; \skt{कोप\appbrk{}ये\appbrk{}द्स\appbrk{}र्व\appbrk{}दो\appbrk{}षान्वा}\ K}}\skt{ }\edtext{\skt{तस्मा\-द्द\-द्या\-द\-ना\-सिते}}{\lemma{\skt{तस्मा\appbrk{}द्द\appbrk{}द्या\appbrk{}द\appbrk{}ना\appbrk{}सिते}}\Afootnote{K H; \skt{तस्मा\appbrk{}द्द\appbrk{}द्या\appbrk{}द\appbrk{}भो\appbrk{}जिने~॥}\ A}}\skt{~।}\pend

\pstart\versepar \ledsidenote{\passref{4.38.23}}\edtext{\skt{आव\-स्थि\-का\-न्क्र\-मा\-न्वापि}}{\lemma{\skt{आव\appbrk{}स्थि\appbrk{}का\appbrk{}न्क्र\appbrk{}मा\appbrk{}न्वापि}}\Afootnote{H; \skt{आवस्थिकं क्रमं चापि}\ A; \skt{आवस्थिकं क्रमां वापि}\ K}}\skt{ }\edtext{\skt{मत्वा}}{\lemma{\skt{मत्वा}}\Afootnote{K H; \skt{बुद्ध्वा}\ A}}\skt{ कार्यन्निरूहणं~॥ }\pend
\pstart\versepar \skt{}\edtext{\edtext{\skt{मले\-प\-हृ\-त\-दो\-षा\-णा\-म्ब\-ल\-व\-त्त्वन्न}}{\lemma{\skt{मले\appbrk{}प\appbrk{}हृ\appbrk{}त\appbrk{}दो\appbrk{}षा\appbrk{}णाम्}}\Afootnote{H; \skt{मलेऽपकृष्टे दोषाणां}\ A; \skt{मलेपहृते दोषाणां}\ K}}}{\lemma{\skt{बलवत्त्वन्न}}\Afootnote{A H; \skt{बलवत्व न}\ K}}\skt{ विद्यते~।}\pend

\pstart\versepar \ledsidenote{\passref{4.38.24}}\skt{क्षीराण्यम्लानि मूत्राणि स्नेहाः क्वाथा रसास्तथा~॥ }\pend
\pstart\versepar \skt{लवणानि }\edtext{\skt{तथा}}{\lemma{\skt{तथा}}\Afootnote{K H; \skt{फलं}\ A}}\skt{ क्षौद्रं शताह्वा }\edtext{\skt{सर्षपा}}{\lemma{\skt{सर्षपा}}\Afootnote{K H; \skt{सर्षपं}\ A}}\skt{ वचा~।}\pend

\pstart\versepar \ledsidenote{\passref{4.38.25}}\edtext{\edtext{\skt{कुष्टन्त्रिकटकं}}{\lemma{\skt{कुष्टन्}}\Afootnote{H; \skt{एला}\ A; \skt{कुष्ठं}\ K}}}{\lemma{\skt{त्रिकटकं}}\Afootnote{H; \skt{त्रिकटुकं}\ A K}}\skt{ रास्ना }\edtext{\skt{सरला}}{\lemma{\skt{सरला}}\Afootnote{K H; \skt{सरलो}\ A}}\skt{ }\edtext{\skt{सुरदारु च}}{\lemma{\skt{सुरदारु च}}\Afootnote{K H; \skt{देवदारु च~॥}\ A}}\skt{~॥ }\pend
\pstart\versepar \skt{}\edtext{\skt{रजने}}{\lemma{\skt{रजने}}\Afootnote{K H; \skt{रजनी}\ A}}\skt{ मधुकं हिङ्गु }\edtext{\skt{एला}}{\lemma{\skt{एला}}\Afootnote{K H; \skt{कुष्ठं}\ A}}\skt{ }\edtext{\skt{संशोधनानि}}{\lemma{\skt{संशोधनानि}}\Afootnote{A H; \skt{संशोधाणानि}\ K}}\skt{ च~।}\pend

\pstart\versepar \ledsidenote{\passref{4.38.26}}\edtext{\skt{कटुकां}}{\lemma{\skt{कटुकां}}\Afootnote{K H; \skt{कटुका}\ A}}\skt{ }\edtext{\skt{शर्क्करां}}{\lemma{\skt{शर्क्करां}}\Afootnote{K H; \skt{शर्करा}\ A}}\skt{ }\edtext{\skt{मुस्ता\-मु\-शी\-र\-ञ्च\-न्दनं}}{\lemma{\skt{मुस्ताम्}}\Afootnote{K H; \skt{मुस्तम्}\ A}}\skt{ }\edtext{\skt{शठीं}}{\lemma{\skt{शठीं}}\Afootnote{K H; \skt{शटी~॥}\ A}}\skt{~॥ }\pend
\pstart\versepar \skt{}\edtext{\edtext{\skt{मञ्जि\-ष्ठा\-म्म\-द\-न\-ञ्चण्डां}}{\lemma{\skt{मञ्जिष्ठाम्}}\Afootnote{K H; \skt{मञ्जिष्ठा}\ A}}}{\lemma{\skt{चण्डां}}\Afootnote{K H; \skt{चण्डा}\ A}}\skt{ }\edtext{\skt{त्रायमाणां}}{\lemma{\skt{त्रायमाणां}}\Afootnote{K H; \skt{त्रायमाणा}\ A}}\skt{ रसाञ्जनं~।}\pend

\pstart\versepar \ledsidenote{\passref{4.38.27}}\edtext{\skt{बिल्वमध्यं}}{\lemma{\skt{बिल्वमध्यं}}\Afootnote{A H; \skt{बिल्वमध्य}\ K}}\skt{ }\edtext{\edtext{\skt{यवा\-नी\-ञ्च\-फ\-लिनी}}{\lemma{\skt{यवानीञ्}}\Afootnote{K H; \skt{यवानी}\ A}}}{\lemma{\skt{चफलिनी}}\Afootnote{A H; \skt{च फलिनीं}\ K}}\skt{ शक्रजा यवान्~॥ }\pend
\pstart\versepar \skt{काकोली क्षीरकाकोली }\edtext{\skt{जीवकर्षभकौ}}{\lemma{\skt{जीवकर्षभकौ}}\Afootnote{K H; \skt{जीवकर्षभका}\ A}}\skt{ }\edtext{\skt{तथा}}{\lemma{\skt{तथा}}\Afootnote{K H; \skt{वुभौ~॥}\ A}}\skt{~।}\edtext{}{\lemma{\skt{यवान्~॥ काकोली}}\Afootnote{K H; \skt{यवाः~॥ काकोली}\ A}}\pend

\pstart\versepar \ledsidenote{\passref{4.38.28}}\edtext{\skt{मेदा}}{\lemma{\skt{मेदा}}\Afootnote{K H; \skt{तथा भेदा}\ A}}\skt{ }\edtext{\skt{चैव}}{\lemma{\skt{चैव}}\Afootnote{H; \emph{om.}\ A; \skt{जैव}\ K}}\skt{ }\edtext{\skt{महामेदा}}{\lemma{\skt{महामेदा}}\Afootnote{K H; \skt{महाभेदा}\ A}}\skt{ }\edtext{\skt{ऋद्धि}}{\lemma{\skt{ऋद्धि}}\Afootnote{K H; \skt{ऋद्धिर्}\ A}}\skt{ वृद्धिर्मधूलिका~॥ }\pend
\pstart\versepar \skt{निरूहेषु यथालाभमेष वर्ग्गो विधीयते~।}\pend

\pstart\versepar \ledsidenote{\passref{4.38.29}}\edtext{\skt{क्वाथस्य}}{\lemma{\skt{क्वाथस्य}}\Afootnote{K H; \skt{स्वस्थे क्वाथस्य}\ A}}\skt{ }\edtext{\skt{भागाश्चत्वारः}}{\lemma{\skt{भागाश्}}\Afootnote{K H; \emph{om.}\ A}}\skt{ }\edtext{\skt{स्नेहभागश्च}}{\lemma{\skt{स्नेहभागश्च}}\Afootnote{K H; \skt{भागाः स्नेहस्य}\ A}}\skt{ पञ्चमः~॥ }\pend
\pstart\versepar \skt{}\edtext{\edtext{\skt{स्वस्थेऽनिले}}{\lemma{\skt{स्वस्थे}}\Afootnote{K H; \skt{क्रुद्धे}\ A}}}{\lemma{\skt{ऽनिले}}\Afootnote{A H; \skt{निले}\ K}}\skt{ }\edtext{\skt{चतुर्थस्तु}}{\lemma{\skt{चतुर्थस्तु}}\Afootnote{A H; \skt{चतुर्थस्तु}\ K}}\skt{ }\edtext{\skt{कुपि\-ते\-न्य\-त\-रो\-ष्टमः}}{\lemma{\skt{कुपि\appbrk{}ते\appbrk{}न्य\appbrk{}त\appbrk{}रो\appbrk{}ष्टमः}}\Afootnote{K H; \skt{षष्ठः पित्ते कफेऽष्टमः~॥}\ A}}\skt{~।}\pend

\pstart\versepar \ledsidenote{\passref{4.38.30}}\skt{सर्व्वेषु चाष्टमो भागः कल्का\-नां\-ल\-व\-ण\-म्पुनः~॥ }\pend
\pstart\versepar \skt{}\edtext{\edtext{\skt{क्षौद्र\-मू\-ल\-घृ\-त\-क्षी\-र\-म\-म्ल\-मां\-स\-र\-स\-स्तथा}}{\lemma{\skt{क्षौद्र\appbrk{}मू\appbrk{}ल\appbrk{}घृ\appbrk{}त\appbrk{}क्षी\appbrk{}रम्}}\Afootnote{H; \skt{क्षौद्रं मूत्रं फलं क्षीर}\ A; \skt{क्षौद्र\appbrk{}मू\appbrk{}त्र\appbrk{}घृ\appbrk{}त\appbrk{}क्षीर}\ K}}}{\lemma{\skt{अम्ल\appbrk{}मां\appbrk{}स\appbrk{}र\appbrk{}स\appbrk{}स्तथा}}\Afootnote{H; \skt{मम्लं मांसरसं तथा~॥}\ A; \skt{मम्व्ल\appbrk{}मां\appbrk{}स\appbrk{}र\appbrk{}स\appbrk{}न्तथा~।}\ K}}\skt{~।}\pend

\pstart \ledsidenote{\passref{4.38.31}}\skt{युक्त्या }\edtext{\edtext{\skt{प्रदा\-प\-ये\-द्धी\-मा\-न्नि\-रूहे}}{\lemma{\skt{प्रदा\appbrk{}प\appbrk{}ये\appbrk{}द्धी\appbrk{}मान्}}\Afootnote{K H; \skt{प्रक\appbrk{}ल्प\appbrk{}ये\appbrk{}द्धी\appbrk{}मान्}\ A}}}{\lemma{\skt{निरूहे}}\Afootnote{K H; \skt{निरू}\ A}}\skt{ }\edtext{\skt{कल्पनात्वियं}}{\lemma{\skt{कल्पनात्वियं}}\Afootnote{K H; \skt{ह्कल्पना त्वियम्~॥}\ A}}\skt{~॥}\pend

\pstart\versepar \ledsidenote{\passref{4.38.32}}\edtext{\skt{क्लक\-स्ने\-ह\-क\-षा\-याणां}}{\lemma{\skt{क्लक\appbrk{}स्ने\appbrk{}ह\appbrk{}क\appbrk{}षा\appbrk{}याणां}}\Afootnote{H; \skt{कल्क\appbrk{}स्ने\appbrk{}ह\appbrk{}क\appbrk{}षा\appbrk{}याणा}\ A; \skt{कल्क\appbrk{}स्ने\appbrk{}ह\appbrk{}क\appbrk{}षा\appbrk{}या\appbrk{}णाम्}\ K}}\skt{ }\edtext{\skt{मवि\-वे\-का\-द्भि\-ष\-ग्वरैः}}{\lemma{\skt{मवि\appbrk{}वे\appbrk{}का\appbrk{}द्भि\appbrk{}ष\appbrk{}ग्वरैः}}\Afootnote{H; \skt{मवि\appbrk{}वे\appbrk{}का\appbrk{}द्भि\appbrk{}ष\appbrk{}ग्वरैः~।}\ A; \skt{अवि\appbrk{}वे\appbrk{}का\appbrk{}द्भि\appbrk{}ष\appbrk{}ग्वरैः~।}\ K}}\skt{~। }\pend
\pstart\versepar \skt{}\edtext{\skt{बस्तिस्तु}}{\lemma{\skt{बस्तिस्तु}}\Afootnote{K H; \skt{बस्तेः सु}\ A}}\skt{ }\edtext{\skt{कल्पितः}}{\lemma{\skt{कल्पितः}}\Afootnote{K H; \skt{कल्पना}\ A}}\skt{ }\edtext{\skt{प्रोक्तस्तस्य}}{\lemma{\skt{प्रोक्तस्}}\Afootnote{K H; \skt{प्रोक्ता}\ A}}\skt{ दानं }\edtext{\skt{यथार्थवकृत्}}{\lemma{\skt{यथार्थवकृत्}}\Afootnote{H\textsuperscript{pc}; \skt{यथार्थकृत्}\ A K H\textsuperscript{ac}}}\pend

\pstart\versepar \ledsidenote{\passref{4.38.77}}\edtext{\skt{वात\-घ्नौ\-ष\-ध\-निः\-क्वाथाः}}{\lemma{\skt{वात\appbrk{}घ्नौ\appbrk{}ष\appbrk{}ध\appbrk{}निः\appbrk{}क्वाथाः}}\Afootnote{K H; \skt{वात\appbrk{}घ्नौ\appbrk{}ष\appbrk{}ध\appbrk{}नि\appbrk{}ष्क्वाथाः}\ A}}\skt{ सैन्धवत्रिवृता युताः~। }\pend
\pstart\versepar \skt{}\edtext{\skt{साम्लाः}}{\lemma{\skt{साम्लाः}}\Afootnote{A H; \skt{साम्व्लाः}\ K}}\skt{ सुखोष्णा }\edtext{\skt{देयाः स्युर्ब्बस्तयः}}{\lemma{\skt{देयाः स्युर्}}\Afootnote{H; \skt{योज्याः स्युर्}\ A; \skt{देया स्युर्}\ K}}\skt{ }\edtext{\skt{कुपितेनिले}}{\lemma{\skt{कुपितेनिले}}\Afootnote{K H; \skt{कुपितेऽनिले~॥}\ A}}\skt{~॥}\pend

\pstart\versepar \ledsidenote{\passref{4.38.78}}\edtext{\skt{न्यग्रोधादिगणः}}{\lemma{\skt{न्यग्रोधादिगणः}}\Afootnote{K H; \skt{न्यग्रोधादिगण}\ A}}\skt{ क्वाथाः काकोल्यादि समायुता~। }\pend
\pstart\versepar \skt{विधेया बस्तयः पित्ते }\edtext{\skt{सक्षौ\-द्र\-घृ\-त\-श\-र्क्कराः}}{\lemma{\skt{सक्षौ\appbrk{}द्र\appbrk{}घृ\appbrk{}त\appbrk{}श\appbrk{}र्क्कराः}}\Afootnote{K H; \skt{ससर्पिष्काः सशर्कराः~॥}\ A}}\skt{~॥}\pend

\pstart\versepar \ledsidenote{\passref{4.38.79}}\edtext{\skt{आर\-ग्व\-ध\-ग\-ण\-क्वाथाः}}{\lemma{\skt{आर\appbrk{}ग्व\appbrk{}ध\appbrk{}ग\appbrk{}ण\appbrk{}क्वाथाः}}\Afootnote{K H; \skt{आर\appbrk{}ग्व\appbrk{}धा\appbrk{}दि\appbrk{}नि\appbrk{}ष्क्वाथाः}\ A}}\skt{ पिप्पल्यादि समायुताः~। }\pend
\pstart\versepar \skt{}\edtext{\skt{सक्षौद्रमूत्रा}}{\lemma{\skt{सक्षौद्रमूत्रा}}\Afootnote{A H; \skt{सक्षौद्रमूत्र}\ K}}\skt{ }\edtext{\skt{देयाः स्युर्ब्बस्तयः}}{\lemma{\skt{देयाः स्युर्}}\Afootnote{A H; \skt{देया स्यु}\ K}}\skt{ कुपिते कफे~॥}\pend

\pstart\versepar \ledsidenote{\passref{4.38.80}}\skt{शर्क्क\-रे\-क्षु\-र\-स\-क्षी\-र\-घृ\-त\-युक्ताः सुशीतलाः~। }\pend
\pstart\versepar \skt{}\edtext{\skt{क्षीर\-वृ\-क्ष\-क\-षा\-याद्याः}}{\lemma{\skt{क्षीर\appbrk{}वृ\appbrk{}क्ष\appbrk{}क\appbrk{}षा\appbrk{}याद्याः}}\Afootnote{H; \skt{क्षीर\appbrk{}वृ\appbrk{}क्ष\appbrk{}क\appbrk{}षा\appbrk{}याढ्या}\ A; \skt{क्षार\appbrk{}वृ\appbrk{}क्ष\appbrk{}क\appbrk{}षा\appbrk{}याद्या}\ K}}\skt{ बस्तयः शोणिते हिताः~॥}\pend

\pstart\versepar \ledsidenote{\passref{4.38.81}}\edtext{\edtext{\skt{शोध\-न\-द्र\-व्य\-निः\-क्वा\-था\-स्त\-त्क\-ल्क\-स्ने\-ह\-सं\-युताः}}{\lemma{\skt{शोध\appbrk{}न\appbrk{}द्र\appbrk{}व्य\appbrk{}निः\appbrk{}क्वा\appbrk{}थास्}}\Afootnote{K H; \skt{शोध\appbrk{}न\appbrk{}द्र\appbrk{}व्य\appbrk{}नि\appbrk{}ष्क्वाथा}\ A}}}{\lemma{\skt{तत्क\appbrk{}ल्क\appbrk{}स्ने\appbrk{}ह\appbrk{}सं\appbrk{}युताः}}\Afootnote{K H; \skt{स्तत्क\appbrk{}ल्क\appbrk{}स्ने\appbrk{}ह\appbrk{}सै\appbrk{}न्धवैः~। युक्ताः}\ A}}\skt{~। }\pend
\pstart\versepar \skt{खजेन }\edtext{\skt{मथिता}}{\lemma{\skt{मथिता}}\Afootnote{K H; \skt{मथित}\ A}}\skt{ }\edtext{\skt{देया}}{\lemma{\skt{देया}}\Afootnote{K H; \emph{om.}\ A}}\skt{ बस्तयः }\edtext{\skt{शोधने}}{\lemma{\skt{शोधने}}\Afootnote{K H; \skt{शोधनाः}\ A}}\skt{ }\edtext{\skt{मताः}}{\lemma{\skt{मताः}}\Afootnote{K H; \skt{स्मृताः~॥}\ A}}\skt{~॥}\pend

\pstart\versepar \ledsidenote{\passref{4.38.87}}\edtext{\skt{प्रियं\-ग्वा\-दि\-ग\-ण\-क्वाथाः}}{\lemma{\skt{प्रियं\appbrk{}ग्वा\appbrk{}दि\appbrk{}ग\appbrk{}ण\appbrk{}क्वाथाः}}\Afootnote{K H; \skt{प्रिय\appbrk{}ङ्ग्वा\appbrk{}दि\appbrk{}ग\appbrk{}ण\appbrk{}क्वाथा}\ A}}\skt{ }\edtext{\skt{पिष्टा\-म्ब\-स्ष्टा\-दि\-सं\-युताः}}{\lemma{\skt{पिष्टा\appbrk{}म्ब\appbrk{}स्ष्टा\appbrk{}दि\appbrk{}सं\appbrk{}युताः}}\Afootnote{H; \skt{अम्बष्ठाद्येन संयुताः~।}\ A; \skt{पिष्टाम्बष्टादि संयुताः~।}\ K}}\skt{~। }\pend
\pstart\versepar \skt{सक्षौद्राः }\edtext{\skt{सघृताश्चापि}}{\lemma{\skt{चापि}}\Afootnote{K H; \skt{चै}\ A}}\skt{ }\edtext{\skt{बस्तयो}}{\lemma{\skt{बस्तयो}}\Afootnote{K H; \skt{व}\ A}}\skt{ ग्राहिणो }\edtext{\skt{मताः}}{\lemma{\skt{मताः}}\Afootnote{K H; \skt{बस्तयः स्मृताः~॥}\ A}}\skt{~॥}\pend

\pstart\versepar \ledsidenote{\passref{4.38.87.1}}\skt{त्रिफ\-ला\-क\-ल्क\-गो\-मू\-त्र\-क्षौ\-द्र\-क्षा\-र\-स\-म\-न्विताः~। }\pend
\pstart\versepar \skt{ऊष\-का\-दि\-स\-मा\-वापा बस्तयो लेखना मताः~॥}\pend

\pstart\versepar \ledsidenote{\passref{4.38.87.2}}\edtext{\skt{बृंह\-ण\-द्र\-व्य\-निः\-क्वा\-थाः\-क\-ल्कै\-र्म्म\-धु\-र\-सं\-युताः}}{\lemma{\skt{बृंह\appbrk{}ण\appbrk{}द्र\appbrk{}व्य\appbrk{}निः\appbrk{}क्वा\appbrk{}थाः\appbrk{}क\appbrk{}ल्कैर्}}\Afootnote{H; \skt{बृङ्ह\appbrk{}ण\appbrk{}द्र\appbrk{}व्य\appbrk{}निः\appbrk{}क्वाथाः कल्कैर्}\ K}}\skt{~। }\pend
\pstart\versepar \skt{सर्प्पि\-र्म्मां\-स\-र\-सो\-पेता बस्तयो बृंहणा मताः~॥}\pend

\pstart\versepar \ledsidenote{\passref{4.38.87.3}}\skt{चट\-का\-ण्डो\-च्च\-ट\-क्वाथाः सक्षी\-र\-घृ\-त\-श\-र्क्कराः~। }\pend
\pstart\versepar \skt{आत्म\-गु\-प्त\-फ\-ला\-वापाः स्मृता व्याजीकरा नृणां~॥}\pend

\pstart\versepar \ledsidenote{\passref{4.38.87.4}}\skt{बद\-र्यै\-रा\-व\-ती\-शेलू शाल्मली धन्वनाङ्कुराः~। }\pend
\pstart\versepar \skt{क्षीरसिद्धास्तु शीतासृक्संयुक्तः पिच्छिला मताः~॥}\pend

\pstart\versepar \ledsidenote{\passref{4.38.88}}\edtext{\skt{एतेष्वेव च}}{\lemma{\skt{एतेष्वेव च}}\Afootnote{K H; \skt{एतेष्वेव च}\ A}}\skt{ योगेषु }\edtext{\skt{स्नेहः}}{\lemma{\skt{स्नेहः}}\Afootnote{K H; \skt{स्नेहाः}\ A}}\skt{ सिद्धाः पृथक्पृथक्~। }\pend
\pstart\versepar \skt{तत्कर्म्मसाधना }\edtext{\edtext{\skt{वैद्यैर्व्विधेया}}{\lemma{\skt{वैद्यैर्}}\Afootnote{K H; \skt{सम्य}\ A}}}{\lemma{\skt{व्विधेया}}\Afootnote{K H; \skt{ग्विधेयाः}\ A}}\skt{ }\edtext{\skt{स्नेहबस्तिषु}}{\lemma{\skt{स्नेहबस्तिषु}}\Afootnote{K H; \skt{स्नेहबस्तयः~॥}\ A}}\skt{~॥}\edtext{}{\lemma{\skt{पृथक्पृथक्~। तत्कर्म्मसाधना}}\Afootnote{K H; \skt{पृथक्पृथक्~। समस्तेष्वथवा}\ A}}\pend

\pstart\versepar \ledsidenote{\passref{4.38.89}}\skt{वन्ध्यानां }\edtext{\skt{शतपाकेशो}}{\lemma{\skt{शतपाकेशो}}\Afootnote{K H; \skt{शतपाकेन}\ A}}\skt{ शोधितानां यथाक्रमं~। }\pend
\pstart\versepar \skt{}\edtext{\skt{बलातेलैन}}{\lemma{\skt{बलातेलैन}}\Afootnote{H; \skt{बलातैलेन}\ A K}}\skt{ }\edtext{\skt{देयाः }\edtext{\skt{स्युर्ब्ब\-स्त\-य\-स्त्रै\-वृ\-तेन}}{\lemma{\skt{ब्बस्तयस्}}\Afootnote{K H; \skt{बस्तय}\ A}}}{\lemma{\skt{देयाः स्युर्}}\Afootnote{A H; \skt{देया स्युर्}\ K}}\skt{ वा~॥ }\edtext{\skt{ओ}}{\lemma{\skt{ओ}}\Afootnote{H; \emph{om.}\ A K}}\skt{~॥}\edtext{}{\lemma{\skt{त्रैवृतेन वा}}\Afootnote{K H; \skt{स्त्रैवृतेन च~॥}\ A}}\pend

\pstart\versepar \ledsidenote{\passref{4.38.89.1}}\edtext{\skt{सिद्धानतः}}{\lemma{\skt{सिद्धानतः}}\Afootnote{H\textsuperscript{pc}; \emph{om.}\ H\textsuperscript{ac}}}\skt{ }\edtext{\skt{प्रवक्ष्यामि}}{\lemma{\skt{प्रवक्ष्यामि}}\Afootnote{H\textsuperscript{pc}; \emph{om.}\ H\textsuperscript{ac}}}\skt{ }\edtext{\edtext{\skt{बस्ती\-न्ना\-ना\-ग\-दा\-प\-हान्}}{\lemma{\skt{बस्तीन्}}\Afootnote{H\textsuperscript{pc}; \emph{om.}\ H\textsuperscript{ac}}}}{\lemma{\skt{नाना\appbrk{}ग\appbrk{}दा\appbrk{}प\appbrk{}हान्}}\Afootnote{H\textsuperscript{pc}; \emph{om.}\ H\textsuperscript{ac}}}\skt{~। }\pend
\pstart\versepar \skt{}\edtext{\skt{दश\-मू\-ला\-ज\-मां\-सस्य}}{\lemma{\skt{दश\appbrk{}मू\appbrk{}ला\appbrk{}ज\appbrk{}मां\appbrk{}सस्य}}\Afootnote{H\textsuperscript{pc}; \emph{om.}\ H\textsuperscript{ac}}}\skt{ }\edtext{\skt{क्वाथ\-म\-ष्टां\-श\-शे\-षितं}}{\lemma{\skt{क्वाथ\appbrk{}म\appbrk{}ष्टां\appbrk{}श\appbrk{}शे\appbrk{}षितं}}\Afootnote{H\textsuperscript{pc}; \emph{om.}\ H\textsuperscript{ac}}}\skt{~॥}\pend

\pstart\versepar \ledsidenote{\passref{4.38.89.2}}\edtext{\skt{सिद्धानतः}}{\lemma{\skt{सिद्धानतः}}\Afootnote{H\textsuperscript{pc}; \emph{om.}\ H\textsuperscript{ac}}}\skt{ }\edtext{\skt{प्रवक्ष्यामि}}{\lemma{\skt{प्रवक्ष्यामि}}\Afootnote{H\textsuperscript{pc}; \emph{om.}\ H\textsuperscript{ac}}}\skt{ }\edtext{\edtext{\skt{बस्ती\-न्ना\-ना\-ग\-दा\-प\-हान्}}{\lemma{\skt{बस्तीन्}}\Afootnote{H\textsuperscript{pc}; \emph{om.}\ H\textsuperscript{ac}}}}{\lemma{\skt{नाना\appbrk{}ग\appbrk{}दा\appbrk{}प\appbrk{}हान्}}\Afootnote{H\textsuperscript{pc}; \emph{om.}\ H\textsuperscript{ac}}}\skt{~। }\pend
\pstart\versepar \skt{}\edtext{\skt{दश\-मू\-ला\-ज\-मां\-सस्य}}{\lemma{\skt{दश\appbrk{}मू\appbrk{}ला\appbrk{}ज\appbrk{}मां\appbrk{}सस्य}}\Afootnote{H\textsuperscript{pc}; \emph{om.}\ H\textsuperscript{ac}}}\skt{ }\edtext{\skt{क्वाथ\-म\-ष्टां\-श\-शे\-षितं}}{\lemma{\skt{क्वाथ\appbrk{}म\appbrk{}ष्टां\appbrk{}श\appbrk{}शे\appbrk{}षितं}}\Afootnote{H\textsuperscript{pc}; \emph{om.}\ H\textsuperscript{ac}}}\skt{~॥}\pend

\pstart\versepar \ledsidenote{\passref{4.38.89.3}}\edtext{\skt{फलं}}{\lemma{\skt{फलं}}\Afootnote{H\textsuperscript{pc}; \emph{om.}\ H\textsuperscript{ac}}}\skt{ }\edtext{\edtext{\skt{सैन्धवमम्लञ्च}}{\lemma{\skt{सैन्धवम्}}\Afootnote{H\textsuperscript{pc}; \emph{om.}\ H\textsuperscript{ac}}}}{\lemma{\skt{अम्लञ्च}}\Afootnote{H\textsuperscript{pc}; \emph{om.}\ H\textsuperscript{ac}}}\skt{ }\edtext{\skt{संयोज्य}}{\lemma{\skt{संयोज्य}}\Afootnote{H\textsuperscript{pc}; \emph{om.}\ H\textsuperscript{ac}}}\skt{ }\edtext{\skt{खजमूर्च्छितं}}{\lemma{\skt{खजमूर्च्छितं}}\Afootnote{H\textsuperscript{pc}; \emph{om.}\ H\textsuperscript{ac}}}\skt{~। }\pend
\pstart\versepar \skt{}\edtext{\edtext{\edtext{\skt{सुखो\-ष्ण\-म्वा\-त\-रो\-ग\-घ्न\-म्ब\-ल\-व\-र्ण्ण\-क\-र\-म्परं}}{\lemma{\skt{सुखोष्णम्}}\Afootnote{H\textsuperscript{pc}; \emph{om.}\ H\textsuperscript{ac}}}}{\lemma{\skt{वातरोगघ्नम्}}\Afootnote{H\textsuperscript{pc}; \emph{om.}\ H\textsuperscript{ac}}}}{\lemma{\skt{बल\appbrk{}व\appbrk{}र्ण्ण\appbrk{}क\appbrk{}र\appbrk{}म्परं}}\Afootnote{H\textsuperscript{pc}; \emph{om.}\ H\textsuperscript{ac}}}\skt{~॥}\pend

\pstart\versepar \ledsidenote{\passref{4.38.89.4}}\edtext{\skt{बस्ति\-र्न्नि\-र\-त्य\-य\-मिमं}}{\lemma{\skt{बस्ति\appbrk{}र्न्नि\appbrk{}र\appbrk{}त्य\appbrk{}य\appbrk{}मिमं}}\Afootnote{H\textsuperscript{pc}; \emph{om.}\ H\textsuperscript{ac}}}\skt{ }\edtext{\skt{विद\-ध्या\-द्द\-श\-मू\-लिकं}}{\lemma{\skt{विद\appbrk{}ध्या\appbrk{}द्द\appbrk{}श\appbrk{}मू\appbrk{}लिकं}}\Afootnote{H\textsuperscript{pc}; \emph{om.}\ H\textsuperscript{ac}}}\skt{~। }\pend
\pstart\versepar \skt{}\edtext{\skt{त्राय\-मा\-णा\-प\-टो\-लञ्च}}{\lemma{\skt{त्राय\appbrk{}मा\appbrk{}णा\appbrk{}प\appbrk{}टो\appbrk{}लञ्च}}\Afootnote{H\textsuperscript{pc}; \emph{om.}\ H\textsuperscript{ac}}}\skt{ }\edtext{\skt{पञ्च\-मू\-ल\-त्रि\-क\-ण्टके}}{\lemma{\skt{पञ्च\appbrk{}मू\appbrk{}ल\appbrk{}त्रि\appbrk{}क\appbrk{}ण्टके}}\Afootnote{H\textsuperscript{pc}; \emph{om.}\ H\textsuperscript{ac}}}\skt{~॥}\pend

\pstart\versepar \ledsidenote{\passref{4.38.89.5}}\edtext{\skt{बलायवैश्च}}{\lemma{\skt{बलायवैश्च}}\Afootnote{H\textsuperscript{pc}; \emph{om.}\ H\textsuperscript{ac}}}\skt{ }\edtext{\skt{निःक्वाथ्य}}{\lemma{\skt{निःक्वाथ्य}}\Afootnote{H\textsuperscript{pc}; \emph{om.}\ H\textsuperscript{ac}}}\skt{ }\edtext{\skt{पाद\-शे\-ष\-ञ्ज\-ला\-ढकं}}{\lemma{\skt{पाद\appbrk{}शे\appbrk{}ष\appbrk{}ञ्ज\appbrk{}ला\appbrk{}ढकं}}\Afootnote{H\textsuperscript{pc}; \emph{om.}\ H\textsuperscript{ac}}}\skt{~। }\pend
\pstart\versepar \skt{}\edtext{\skt{छाग\-मां\-स\-र\-सा\-र्द्धेन}}{\lemma{\skt{छाग\appbrk{}मां\appbrk{}स\appbrk{}र\appbrk{}सा\appbrk{}र्द्धेन}}\Afootnote{H\textsuperscript{pc}; \emph{om.}\ H\textsuperscript{ac}}}\skt{ }\edtext{\edtext{\skt{योजि\-त\-म्वि\-प\-चे\-त्पुनः}}{\lemma{\skt{योजितम्}}\Afootnote{H\textsuperscript{pc}; \emph{om.}\ H\textsuperscript{ac}}}}{\lemma{\skt{विपचेत्पुनः}}\Afootnote{H\textsuperscript{pc}; \emph{om.}\ H\textsuperscript{ac}}}\skt{~॥}\pend

\pstart\versepar \ledsidenote{\passref{4.38.89.6}}\edtext{\edtext{\skt{प्रस्थशेषन्ततः}}{\lemma{\skt{प्रस्थशेषन्}}\Afootnote{H\textsuperscript{pc}; \emph{om.}\ H\textsuperscript{ac}}}}{\lemma{\skt{ततः}}\Afootnote{H\textsuperscript{pc}; \emph{om.}\ H\textsuperscript{ac}}}\skt{ }\edtext{\skt{पिष्टैः}}{\lemma{\skt{पिष्टैः}}\Afootnote{H\textsuperscript{pc}; \emph{om.}\ H\textsuperscript{ac}}}\skt{ }\edtext{\edtext{\skt{कान्ता\-कृ\-ष्णा\-द्य\-नै\-र्युतं}}{\lemma{\skt{कान्ता\appbrk{}कृ\appbrk{}ष्णा\appbrk{}द्य\appbrk{}नैर्}}\Afootnote{H\textsuperscript{pc}; \emph{om.}\ H\textsuperscript{ac}}}}{\lemma{\skt{युतं}}\Afootnote{H\textsuperscript{pc}; \emph{om.}\ H\textsuperscript{ac}}}\skt{~। }\pend
\pstart\versepar \skt{}\edtext{\edtext{\skt{सर्प्पि\-स्तै\-ल\-स\-मा\-क्षी\-र\-सै\-न्ध\-वै\-र्म्म\-द\-नेन}}{\lemma{\skt{सर्प्पि\appbrk{}स्तै\appbrk{}ल\appbrk{}स\appbrk{}मा\appbrk{}क्षी\appbrk{}र\appbrk{}सै\appbrk{}न्ध\appbrk{}वैर्}}\Afootnote{H\textsuperscript{pc}; \emph{om.}\ H\textsuperscript{ac}}}}{\lemma{\skt{म्मदनेन}}\Afootnote{H\textsuperscript{pc}; \emph{om.}\ H\textsuperscript{ac}}}\skt{ }\edtext{\skt{च}}{\lemma{\skt{च}}\Afootnote{H\textsuperscript{pc}; \emph{om.}\ H\textsuperscript{ac}}}\skt{~॥}\pend

\pstart\versepar \ledsidenote{\passref{4.38.89.7}}\edtext{\edtext{\edtext{\edtext{\skt{चक्षु\-ष्य\-न्दी\-प\-न\-म्ब\-स्ति\-न्द\-द्या\-त्मां\-स\-ब\-ल\-प्रदं}}{\lemma{\skt{चक्षुष्यन्}}\Afootnote{H\textsuperscript{pc}; \emph{om.}\ H\textsuperscript{ac}}}}{\lemma{\skt{दीपनम्}}\Afootnote{H\textsuperscript{pc}; \emph{om.}\ H\textsuperscript{ac}}}}{\lemma{\skt{बस्तिन्}}\Afootnote{H\textsuperscript{pc}; \emph{om.}\ H\textsuperscript{ac}}}}{\lemma{\skt{दद्या\appbrk{}त्मां\appbrk{}स\appbrk{}ब\appbrk{}ल\appbrk{}प्रदं}}\Afootnote{H\textsuperscript{pc}; \emph{om.}\ H\textsuperscript{ac}}}\skt{~। }\pend
\pstart\versepar \skt{}\edtext{\edtext{\skt{एर\-ण्ड\-न्त्रि\-फ\-लां\-मू\-ला\-द्बृ\-हत्यौ}}{\lemma{\skt{एरण्डन्}}\Afootnote{H\textsuperscript{pc}; \emph{om.}\ H\textsuperscript{ac}}}}{\lemma{\skt{त्रिफ\appbrk{}लां\appbrk{}मू\appbrk{}ला\appbrk{}द्बृ\appbrk{}हत्यौ}}\Afootnote{H\textsuperscript{pc}; \emph{om.}\ H\textsuperscript{ac}}}\skt{ }\edtext{\skt{प्रग्रहं}}{\lemma{\skt{प्रग्रहं}}\Afootnote{H\textsuperscript{pc}; \emph{om.}\ H\textsuperscript{ac}}}\skt{ }\edtext{\skt{स्थिरे}}{\lemma{\skt{स्थिरे}}\Afootnote{H\textsuperscript{pc}; \emph{om.}\ H\textsuperscript{ac}}}\skt{~॥}\pend

\pstart\versepar \ledsidenote{\passref{4.38.89.8}}\edtext{\skt{अश्वगन्धामृता}}{\lemma{\skt{अश्वगन्धामृता}}\Afootnote{H\textsuperscript{pc}; \emph{om.}\ H\textsuperscript{ac}}}\skt{ }\edtext{\skt{गन्धा}}{\lemma{\skt{गन्धा}}\Afootnote{H\textsuperscript{pc}; \emph{om.}\ H\textsuperscript{ac}}}\skt{ }\edtext{\skt{पला\-शा\-सु\-र\-दारु च}}{\lemma{\skt{पला\appbrk{}शा\appbrk{}सु\appbrk{}र\appbrk{}दारु च}}\Afootnote{H\textsuperscript{pc}; \emph{om.}\ H\textsuperscript{ac}}}\skt{~। }\pend
\pstart\versepar \skt{}\edtext{\skt{रास्ना}}{\lemma{\skt{रास्ना}}\Afootnote{H\textsuperscript{pc}; \emph{om.}\ H\textsuperscript{ac}}}\skt{ }\edtext{\skt{पुन\-र्न्न\-व\-ब\-ला\-श्व\-दंष्ट्रा}}{\lemma{\skt{पुन\appbrk{}र्न्न\appbrk{}व\appbrk{}ब\appbrk{}ला\appbrk{}श्व\appbrk{}दंष्ट्रा}}\Afootnote{H\textsuperscript{pc}; \emph{om.}\ H\textsuperscript{ac}}}\skt{ }\edtext{\skt{चपलोन्मिताः}}{\lemma{\skt{चपलोन्मिताः}}\Afootnote{H\textsuperscript{pc}; \emph{om.}\ H\textsuperscript{ac}}}\skt{~॥}\pend

\pstart\versepar \ledsidenote{\passref{4.38.89.9}}\edtext{\skt{पक्त्वा}}{\lemma{\skt{पक्त्वा}}\Afootnote{H\textsuperscript{pc}; \emph{om.}\ H\textsuperscript{ac}}}\skt{ }\edtext{\skt{जलाढकं}}{\lemma{\skt{जलाढकं}}\Afootnote{H\textsuperscript{pc}; \emph{om.}\ H\textsuperscript{ac}}}\skt{ }\edtext{\skt{चाम्लद्व्याढके}}{\lemma{\skt{चाम्लद्व्याढके}}\Afootnote{H\textsuperscript{pc}; \emph{om.}\ H\textsuperscript{ac}}}\skt{ }\edtext{\skt{पादशोषिते}}{\lemma{\skt{पादशोषिते}}\Afootnote{H\textsuperscript{pc}; \emph{om.}\ H\textsuperscript{ac}}}\skt{~। }\pend
\pstart\versepar \skt{}\edtext{\skt{पिष्टैः}}{\lemma{\skt{पिष्टैः}}\Afootnote{H\textsuperscript{pc}; \emph{om.}\ H\textsuperscript{ac}}}\skt{ }\edtext{\skt{शताह्वाहपुषा}}{\lemma{\skt{शताह्वाहपुषा}}\Afootnote{H\textsuperscript{pc}; \emph{om.}\ H\textsuperscript{ac}}}\skt{ }\edtext{\skt{कान्ता}}{\lemma{\skt{कान्ता}}\Afootnote{H\textsuperscript{pc}; \emph{om.}\ H\textsuperscript{ac}}}\skt{ }\edtext{\skt{कृष्णा}}{\lemma{\skt{कृष्णा}}\Afootnote{H\textsuperscript{pc}; \emph{om.}\ H\textsuperscript{ac}}}\skt{ }\edtext{\skt{रसाञ्जनैः}}{\lemma{\skt{रसाञ्जनैः}}\Afootnote{H\textsuperscript{pc}; \emph{om.}\ H\textsuperscript{ac}}}\skt{~॥}\pend

\pstart\versepar \ledsidenote{\passref{4.38.89.10}}\edtext{\skt{मुस्ते\-न्द्र\-य\-व\-ष\-ड्ग्रन्था}}{\lemma{\skt{मुस्ते\appbrk{}न्द्र\appbrk{}य\appbrk{}व\appbrk{}ष\appbrk{}ड्ग्रन्था}}\Afootnote{H\textsuperscript{pc}; \emph{om.}\ H\textsuperscript{ac}}}\skt{ }\edtext{\edtext{\skt{मधु\-कै\-श्चा\-क्ष\-भा\-गिकैः}}{\lemma{\skt{मधुकैश्}}\Afootnote{H\textsuperscript{pc}; \emph{om.}\ H\textsuperscript{ac}}}}{\lemma{\skt{चाक्षभागिकैः}}\Afootnote{H\textsuperscript{pc}; \emph{om.}\ H\textsuperscript{ac}}}\skt{~। }\pend
\pstart\versepar \skt{}\edtext{\skt{ससै\-न्ध\-व\-घृ\-त\-क्षौ\-द्र\-तै\-ल\-मू\-त्रैश्च}}{\lemma{\skt{ससै\appbrk{}न्ध\appbrk{}व\appbrk{}घृ\appbrk{}त\appbrk{}क्षौ\appbrk{}द्र\appbrk{}तै\appbrk{}ल\appbrk{}मू\appbrk{}त्रैश्च}}\Afootnote{H\textsuperscript{pc}; \emph{om.}\ H\textsuperscript{ac}}}\skt{ }\edtext{\skt{योजितः}}{\lemma{\skt{योजितः}}\Afootnote{H\textsuperscript{pc}; \emph{om.}\ H\textsuperscript{ac}}}\skt{~॥}\pend

\pstart\versepar \ledsidenote{\passref{4.38.89.11}}\edtext{\skt{सन्ध्य\-स्थि\-पा\-र्श्व\-जं\-घोरु}}{\lemma{\skt{सन्ध्य\appbrk{}स्थि\appbrk{}पा\appbrk{}र्श्व\appbrk{}जं\appbrk{}घोरु}}\Afootnote{H\textsuperscript{pc}; \emph{om.}\ H\textsuperscript{ac}}}\skt{ }\edtext{\skt{पाद\-पृ\-ष्ठ\-त्रि\-का\-श्रितं}}{\lemma{\skt{पाद\appbrk{}पृ\appbrk{}ष्ठ\appbrk{}त्रि\appbrk{}का\appbrk{}श्रितं}}\Afootnote{H\textsuperscript{pc}; \emph{om.}\ H\textsuperscript{ac}}}\skt{~। }\pend
\pstart\versepar \skt{}\edtext{\skt{शूलं}}{\lemma{\skt{शूलं}}\Afootnote{H\textsuperscript{pc}; \emph{om.}\ H\textsuperscript{ac}}}\skt{ }\edtext{\edtext{\skt{कफा\-वृ\-त\-म्वा\-त\-वि\-ण्मू\-त्र\-ग्र\-हणं}}{\lemma{\skt{कफावृतम्}}\Afootnote{H\textsuperscript{pc}; \emph{om.}\ H\textsuperscript{ac}}}}{\lemma{\skt{वात\appbrk{}वि\appbrk{}ण्मू\appbrk{}त्र\appbrk{}ग्र\appbrk{}हणं}}\Afootnote{H\textsuperscript{pc}; \emph{om.}\ H\textsuperscript{ac}}}\skt{ }\edtext{\skt{सरुक्}}{\lemma{\skt{सरुक्}}\Afootnote{H\textsuperscript{pc}; \emph{om.}\ H\textsuperscript{ac}}}\skt{~॥}\pend

\pstart\versepar \ledsidenote{\passref{4.38.89.12}}\edtext{\skt{अश्म\-री\-श\-र्क्क\-रा\-ध्मा\-न\-ग्र\-ह\-ण्य\-र्शो\-ग\-दा\-नपि}}{\lemma{\skt{अश्म\appbrk{}री\appbrk{}श\appbrk{}र्क्क\appbrk{}रा\appbrk{}ध्मा\appbrk{}न\appbrk{}ग्र\appbrk{}ह\appbrk{}ण्य\appbrk{}र्शो\appbrk{}ग\appbrk{}दा\appbrk{}नपि}}\Afootnote{H\textsuperscript{pc}; \emph{om.}\ H\textsuperscript{ac}}}\skt{~। }\pend
\pstart\versepar \skt{}\edtext{\skt{एरण्डादिरयं}}{\lemma{\skt{एरण्डादिरयं}}\Afootnote{H\textsuperscript{pc}; \emph{om.}\ H\textsuperscript{ac}}}\skt{ }\edtext{\skt{हन्या\-द्ब\-स्ति\-ले\-ख\-न\-दी\-पनः}}{\lemma{\skt{हन्या\appbrk{}द्ब\appbrk{}स्ति\appbrk{}ले\appbrk{}ख\appbrk{}न\appbrk{}दी\appbrk{}पनः}}\Afootnote{H\textsuperscript{pc}; \emph{om.}\ H\textsuperscript{ac}}}\skt{~॥}\pend

\pstart\versepar \ledsidenote{\passref{4.38.89.13}}\edtext{\skt{मुस्ता}}{\lemma{\skt{मुस्ता}}\Afootnote{H\textsuperscript{pc}; \emph{om.}\ H\textsuperscript{ac}}}\skt{ }\edtext{\skt{पाठामृतातिक्ता}}{\lemma{\skt{पाठामृतातिक्ता}}\Afootnote{H\textsuperscript{pc}; \emph{om.}\ H\textsuperscript{ac}}}\skt{ }\edtext{\skt{बला}}{\lemma{\skt{बला}}\Afootnote{H\textsuperscript{pc}; \emph{om.}\ H\textsuperscript{ac}}}\skt{ }\edtext{\skt{रास्ना}}{\lemma{\skt{रास्ना}}\Afootnote{H\textsuperscript{pc}; \emph{om.}\ H\textsuperscript{ac}}}\skt{ }\edtext{\skt{पुनर्न्नवा}}{\lemma{\skt{पुनर्न्नवा}}\Afootnote{H\textsuperscript{pc}; \emph{om.}\ H\textsuperscript{ac}}}\skt{~। }\pend
\pstart\versepar \skt{}\edtext{\skt{मञ्जि\-ष्ठा\-र\-ग्व\-धो\-शीर}}{\lemma{\skt{मञ्जि\appbrk{}ष्ठा\appbrk{}र\appbrk{}ग्व\appbrk{}धो\appbrk{}शीर}}\Afootnote{H\textsuperscript{pc}; \emph{om.}\ H\textsuperscript{ac}}}\skt{ }\edtext{\skt{त्राय\-मा\-णा\-क्ष\-गो\-क्षु\-रान्}}{\lemma{\skt{त्राय\appbrk{}मा\appbrk{}णा\appbrk{}क्ष\appbrk{}गो\appbrk{}क्षु\appbrk{}रान्}}\Afootnote{H\textsuperscript{pc}; \emph{om.}\ H\textsuperscript{ac}}}\skt{~॥}\pend

\pstart\versepar \ledsidenote{\passref{4.38.89.14}}\edtext{\skt{सहा\-ल्प\-प\-ञ्च\-मू\-लेन}}{\lemma{\skt{सहा\appbrk{}ल्प\appbrk{}प\appbrk{}ञ्च\appbrk{}मू\appbrk{}लेन}}\Afootnote{H\textsuperscript{pc}; \emph{om.}\ H\textsuperscript{ac}}}\skt{ }\edtext{\skt{पलिकां}}{\lemma{\skt{पलिकां}}\Afootnote{H\textsuperscript{pc}; \emph{om.}\ H\textsuperscript{ac}}}\skt{ }\edtext{\skt{मदनाष्टकां}}{\lemma{\skt{मदनाष्टकां}}\Afootnote{H\textsuperscript{pc}; \emph{om.}\ H\textsuperscript{ac}}}\skt{~। }\pend
\pstart\versepar \skt{}\edtext{\skt{जलद्रोणे}}{\lemma{\skt{जलद्रोणे}}\Afootnote{H\textsuperscript{pc}; \emph{om.}\ H\textsuperscript{ac}}}\skt{ }\edtext{\edtext{\edtext{\skt{पचे\-त्क्वा\-थ\-म्पा\-द\-शे\-ष\-म्पुनः}}{\lemma{\skt{पचेत्क्वाथम्}}\Afootnote{H\textsuperscript{pc}; \emph{om.}\ H\textsuperscript{ac}}}}{\lemma{\skt{पादशेषम्}}\Afootnote{H\textsuperscript{pc}; \emph{om.}\ H\textsuperscript{ac}}}}{\lemma{\skt{पुनः}}\Afootnote{H\textsuperscript{pc}; \emph{om.}\ H\textsuperscript{ac}}}\skt{ }\edtext{\skt{पचेत्}}{\lemma{\skt{पचेत्}}\Afootnote{H\textsuperscript{pc}; \emph{om.}\ H\textsuperscript{ac}}}\skt{~॥}\pend

\pstart\versepar \ledsidenote{\passref{4.38.89.15}}\edtext{\skt{क्षीरप्रस्थेन}}{\lemma{\skt{क्षीरप्रस्थेन}}\Afootnote{H\textsuperscript{pc}; \emph{om.}\ H\textsuperscript{ac}}}\skt{ }\edtext{\skt{संयुक्तं}}{\lemma{\skt{संयुक्तं}}\Afootnote{H\textsuperscript{pc}; \emph{om.}\ H\textsuperscript{ac}}}\skt{ }\edtext{\edtext{\skt{क्षीर\-पि\-ष्ट\-म्प\-रि\-श्रुतं}}{\lemma{\skt{क्षीरपिष्टम्}}\Afootnote{H\textsuperscript{pc}; \emph{om.}\ H\textsuperscript{ac}}}}{\lemma{\skt{परिश्रुतं}}\Afootnote{H\textsuperscript{pc}; \emph{om.}\ H\textsuperscript{ac}}}\skt{~। }\pend
\pstart\versepar \skt{}\edtext{\skt{जङ्ग\-ला\-नु\-र\-सो\-त्पादः}}{\lemma{\skt{जङ्ग\appbrk{}ला\appbrk{}नु\appbrk{}र\appbrk{}सो\appbrk{}त्पादः}}\Afootnote{H\textsuperscript{pc}; \emph{om.}\ H\textsuperscript{ac}}}\skt{ }\edtext{\skt{सर्प्पिषो}}{\lemma{\skt{सर्प्पिषो}}\Afootnote{H\textsuperscript{pc}; \emph{om.}\ H\textsuperscript{ac}}}\skt{ }\edtext{\edtext{\skt{मदनस्तथा}}{\lemma{\skt{मदनस्}}\Afootnote{H\textsuperscript{pc}; \emph{om.}\ H\textsuperscript{ac}}}}{\lemma{\skt{तथा}}\Afootnote{H\textsuperscript{pc}; \emph{om.}\ H\textsuperscript{ac}}}\skt{~॥}\pend

\pstart\versepar \ledsidenote{\passref{4.38.89.16}}\edtext{\skt{शता\-ह्वा\-फ\-लि\-नी\-य\-ष्टी\-व\-त्सकैः}}{\lemma{\skt{शता\appbrk{}ह्वा\appbrk{}फ\appbrk{}लि\appbrk{}नी\appbrk{}य\appbrk{}ष्टी\appbrk{}व\appbrk{}त्सकैः}}\Afootnote{H\textsuperscript{pc}; \emph{om.}\ H\textsuperscript{ac}}}\skt{ }\edtext{\skt{सरसाञ्जनैः}}{\lemma{\skt{सरसाञ्जनैः}}\Afootnote{H\textsuperscript{pc}; \emph{om.}\ H\textsuperscript{ac}}}\skt{~। }\pend
\pstart\versepar \skt{}\edtext{\skt{कार्षिकैः}}{\lemma{\skt{कार्षिकैः}}\Afootnote{H\textsuperscript{pc}; \emph{om.}\ H\textsuperscript{ac}}}\skt{ }\edtext{\skt{सैन्धवयुतैः}}{\lemma{\skt{सैन्धवयुतैः}}\Afootnote{H\textsuperscript{pc}; \emph{om.}\ H\textsuperscript{ac}}}\skt{ }\edtext{\edtext{\skt{कल्कैर्ब्बस्तिः}}{\lemma{\skt{कल्कैर्}}\Afootnote{H\textsuperscript{pc}; \emph{om.}\ H\textsuperscript{ac}}}}{\lemma{\skt{ब्बस्तिः}}\Afootnote{H\textsuperscript{pc}; \emph{om.}\ H\textsuperscript{ac}}}\skt{ }\edtext{\skt{सुयोजितः}}{\lemma{\skt{सुयोजितः}}\Afootnote{H\textsuperscript{pc}; \emph{om.}\ H\textsuperscript{ac}}}\skt{~॥}\pend

\pstart\versepar \ledsidenote{\passref{4.38.89.17}}\edtext{\skt{वाता\-सृ\-क्मो\-ह\-शो\-फार्शो}}{\lemma{\skt{वाता\appbrk{}सृ\appbrk{}क्मो\appbrk{}ह\appbrk{}शो\appbrk{}फार्शो}}\Afootnote{H\textsuperscript{pc}; \emph{om.}\ H\textsuperscript{ac}}}\skt{ गुल्म\-मू\-त्र\-वि\-ब\-न्ध\-नुत्~। }\pend
\pstart\versepar \skt{विस\-र्प्प\-ज्व\-र\-वि\-ड्भङ्ग }\edtext{\skt{रक्त\-पि\-त्त\-वि\-ना\-शनः}}{\lemma{\skt{रक्त\appbrk{}पि\appbrk{}त्त\appbrk{}वि\appbrk{}ना\appbrk{}शनः}}\Afootnote{H\textsuperscript{pc}; \emph{om.}\ H\textsuperscript{ac}}}\skt{~॥}\edtext{}{\lemma{\skt{गुल्म\appbrk{}मू\appbrk{}त्र\appbrk{}वि\appbrk{}ब\appbrk{}न्ध\appbrk{}नुत्~। विस\appbrk{}र्प्प\appbrk{}ज्व\appbrk{}र\appbrk{}वि\appbrk{}ड्भङ्ग}}\Afootnote{H\textsuperscript{pc}; \emph{om.}\ H\textsuperscript{ac}}}\pend

\pstart\versepar \ledsidenote{\passref{4.38.89.20}}\edtext{\skt{वल्यः}}{\lemma{\skt{वल्यः}}\Afootnote{H\textsuperscript{pc}; \emph{om.}\ H\textsuperscript{ac}}}\skt{ }\edtext{\edtext{\skt{सन्दी\-प\-नो\-वृ\-ष्य\-श्च\-क्षुष्यः}}{\lemma{\skt{सन्दी\appbrk{}प\appbrk{}नो\appbrk{}वृ\appbrk{}ष्यश्}}\Afootnote{H\textsuperscript{pc}; \emph{om.}\ H\textsuperscript{ac}}}}{\lemma{\skt{चक्षुष्यः}}\Afootnote{H\textsuperscript{pc}; \emph{om.}\ H\textsuperscript{ac}}}\skt{ }\edtext{\skt{शूलनाशनः}}{\lemma{\skt{शूलनाशनः}}\Afootnote{H\textsuperscript{pc}; \emph{om.}\ H\textsuperscript{ac}}}\skt{~। }\pend
\pstart\versepar \skt{}\edtext{\skt{यायनानानयं}}{\lemma{\skt{यायनानानयं}}\Afootnote{H\textsuperscript{pc}; \emph{om.}\ H\textsuperscript{ac}}}\skt{ }\edtext{\skt{राजा}}{\lemma{\skt{राजा}}\Afootnote{H\textsuperscript{pc}; \emph{om.}\ H\textsuperscript{ac}}}\skt{ }\edtext{\skt{बस्तिमुस्तादिको}}{\lemma{\skt{बस्तिमुस्तादिको}}\Afootnote{H\textsuperscript{pc}; \emph{om.}\ H\textsuperscript{ac}}}\skt{ }\edtext{\skt{मतः}}{\lemma{\skt{मतः}}\Afootnote{H\textsuperscript{pc}; \emph{om.}\ H\textsuperscript{ac}}}\skt{~॥}\pend

\pstart\versepar \ledsidenote{\passref{4.38.90}}\skt{नर\-स्यो\-त्त\-म\-स\-त्वस्य तीक्ष्णं बस्ती\-न्प्र\-दा\-प\-येत्~। }\pend
\pstart\versepar \skt{}\edtext{\skt{मध्य\-म्म\-ध्य\-म\-स\-त्वस्य}}{\lemma{\skt{मध्यमसत्वस्य}}\Afootnote{K H; \skt{मं मध्यसत्त्वस्य}\ A}}\skt{ }\edtext{\skt{हीनसत्वस्य वा}}{\lemma{\skt{हीनसत्वस्य वा}}\Afootnote{K H; \skt{विपरीतस्य वै}\ A}}\skt{ मृदुं~॥}\edtext{}{\lemma{\skt{बस्ती\appbrk{}न्प्र\appbrk{}दा\appbrk{}प\appbrk{}येत्~। मध्यम्}}\Afootnote{H; \skt{बस्तिं निधापयेत्~। मध्य}\ A; \skt{बस्तीं प्रदापयेत्~। मध्यं}\ K}}\pend

\pstart\versepar \ledsidenote{\passref{4.38.91}}\skt{एव\-ङ्का\-ल\-म्ब\-ल\-न्दो\-ष\-म्वि\-का\-रञ्च विकारवित्~। }\pend
\pstart\versepar \skt{बस्ति\-द्र\-व्य\-ब\-ल\-ञ्चैव वीक्ष्य }\edtext{\edtext{\skt{बस्ति\-म्प्र\-यो\-ज\-येत्}}{\lemma{\skt{बस्तिम्}}\Afootnote{K H; \skt{बस्ती}\ A}}}{\lemma{\skt{प्रयोजयेत्}}\Afootnote{K H; \skt{न्प्रयोजयेत्~॥}\ A}}\skt{~॥}\pend

\pstart\versepar \ledsidenote{\passref{4.38.92}}\edtext{\skt{दद्यादुत्क्लेशनं}}{\lemma{\skt{दद्यादुत्क्लेशनं}}\Afootnote{A H; \skt{दद्यादुत्क्लेशनं}\ K}}\skt{ पूर्व्वम्मध्ये दोषहरं पुनः~। }\pend
\pstart\versepar \skt{पश्चा\-त्सं\-श\-म\-नी\-यञ्च }\edtext{\skt{बस्तिं}}{\lemma{\skt{बस्तिं}}\Afootnote{K H; \skt{दद्या}\ A}}\skt{ }\edtext{\skt{बस्तिविशारदः}}{\lemma{\skt{बस्तिविशारदः}}\Afootnote{K H; \skt{द्बस्तिं विचक्षणः~॥}\ A}}\skt{~॥}\pend

\pstart\versepar \ledsidenote{\passref{4.38.96}}\skt{॥ नृपाणां }\edtext{\skt{तत्समानानां}}{\lemma{\skt{तत्समानानां}}\Afootnote{A H; \skt{तत्समाणानां}\ K}}\skt{ तथा }\edtext{\edtext{\skt{वसुमतामपि}}{\lemma{\skt{वसुमताम्}}\Afootnote{H; \skt{सुमहताम्}\ A; \skt{वसुमता}\ K}}}{\lemma{\skt{अपि}}\Afootnote{A H; \skt{मपिः~।}\ K}}\skt{~। }\pend
\pstart\versepar \skt{नारीणां }\edtext{\skt{सुकु\-मा\-रा\-णा\-न्त\-थैव}}{\lemma{\skt{तथैव}}\Afootnote{K H; \emph{om.}\ A}}\skt{ }\edtext{\skt{शिशुवृद्धयोः}}{\lemma{\skt{शिशुवृद्धयोः}}\Afootnote{K H; \skt{शिशु\appbrk{}स्थ\appbrk{}वि\appbrk{}र\appbrk{}यो\appbrk{}रपि~॥}\ A}}\skt{~॥}\pend

\pstart\versepar \ledsidenote{\passref{4.38.97}}\skt{दोष\-नि\-र्ह\-र\-णा\-र्थाय बल\-व\-र्ण्णो\-द\-याय च~। }\pend
\pstart\versepar \skt{समा\-से\-नो\-प\-दे\-क्ष्यामि }\edtext{\skt{विधा\-न\-म्मा\-धु\-तै\-लिकं}}{\lemma{\skt{विधानम्}}\Afootnote{A H; \skt{विधाणं}\ K}}\skt{~॥}\pend

\pstart\versepar \ledsidenote{\passref{4.38.98}}\edtext{\skt{यान\-स्त्री\-पा\-न\-भो\-ज्येषु}}{\lemma{\skt{यान\appbrk{}स्त्री\appbrk{}पा\appbrk{}न\appbrk{}भो\appbrk{}ज्येषु}}\Afootnote{H; \skt{यान\appbrk{}स्त्री\appbrk{}भो\appbrk{}ज्य\appbrk{}पा\appbrk{}नेषु}\ A; \skt{यान\appbrk{}स्त्री\appbrk{}पा\appbrk{}ण\appbrk{}भो\appbrk{}ज्येषु}\ K}}\skt{ }\edtext{\skt{परिहारो न}}{\lemma{\skt{परिहारो न}}\Afootnote{K H; \skt{नियमश्चात्र}\ A}}\skt{ }\edtext{\skt{रोच्यते}}{\lemma{\skt{रोच्यते}}\Afootnote{K H; \skt{नोच्यते~।}\ A}}\skt{~। }\pend
\pstart\versepar \skt{फलञ्च विपुलन्दृष्टं }\edtext{\skt{व्याप\-द\-ञ्चा\-प्य\-स\-म्भवः}}{\lemma{\skt{व्यापदञ्}}\Afootnote{K H; \skt{व्यापदां}\ A}}\skt{~॥}\pend

\pstart \ledsidenote{\passref{4.38.99}}\edtext{\skt{यदिच्छति}}{\lemma{\skt{यदिच्छति}}\Afootnote{K H; \skt{योज्यस्त्वतः सुखेनैव निरू\appbrk{}ह\appbrk{}क्र\appbrk{}म\appbrk{}मि\appbrk{}च्छता~। यदेच्छति}\ A}}\skt{ }\edtext{\skt{तदैवेष}}{\lemma{\skt{तदैवेष}}\Afootnote{K H; \skt{तदैवैष}\ A}}\skt{ }\edtext{\edtext{\skt{निरू\-ह\-क्र\-म\-म\-र्हति}}{\lemma{\skt{निरूहक्रमम्}}\Afootnote{K H; \skt{प्रयो}\ A}}}{\lemma{\skt{अर्हति}}\Afootnote{K H; \skt{क्तव्यो विपश्चिता~॥}\ A}}\skt{~।}\pend

\pstart\versepar \ledsidenote{\passref{4.38.100}}\skt{मधुतैले समे स्यातां }\edtext{\skt{क्वाथ\-श्चै\-र\-ण्ड\-मू\-लजः}}{\lemma{\skt{क्वाथ\appbrk{}श्चै\appbrk{}र\appbrk{}ण्ड\appbrk{}मू\appbrk{}लजः}}\Afootnote{H; \skt{क्वाथ\appbrk{}श्चै\appbrk{}र\appbrk{}ण्ड\appbrk{}मू\appbrk{}लजः~।}\ A; \skt{क्वाथ\appbrk{}श्चै\appbrk{}र\appbrk{}ण्ड\appbrk{}मू\appbrk{}लजः~।}\ K}}\skt{~॥ }\pend
\pstart\versepar \skt{}\edtext{\skt{कर्षश्च}}{\lemma{\skt{कर्षश्च}}\Afootnote{K H; \skt{पलार्धं}\ A}}\skt{ }\edtext{\edtext{\edtext{\skt{शत\-पु\-ष्पा\-णा\-न्त\-तोऽर्द्धं}}{\lemma{\skt{शतपुष्पाणान्}}\Afootnote{K H; \skt{शतपुष्पाया}\ A}}}{\lemma{\skt{ततो}}\Afootnote{H; \skt{स्ततो}\ A; \skt{ततो}\ K}}}{\lemma{\skt{ऽर्द्धं}}\Afootnote{A H; \skt{र्धं}\ K}}\skt{ सैन्धवस्य च~।}\pend

\pstart\versepar \ledsidenote{\passref{4.38.101}}\edtext{\skt{फलनैकेन}}{\lemma{\skt{फलनैकेन}}\Afootnote{H; \skt{फलेनैकेन}\ A K}}\skt{ }\edtext{\skt{संयुक्तः}}{\lemma{\skt{संयुक्तः}}\Afootnote{A H; \skt{संयुक्तऋ}\ K}}\skt{ खजेन }\edtext{\skt{सुविलोडितः}}{\lemma{\skt{सुविलोडितः}}\Afootnote{K H; \skt{च विलोडितः~।}\ A}}\skt{~॥ }\pend
\pstart\versepar \skt{देयः सुखोष्णो भिषजा }\edtext{\skt{माधु\-तै\-लि\-क\-सं\-युतः}}{\lemma{\skt{माधु\appbrk{}तै\appbrk{}लि\appbrk{}क\appbrk{}सं\appbrk{}युतः}}\Afootnote{K H; \skt{माधु\appbrk{}तै\appbrk{}लि\appbrk{}क\appbrk{}सं\appbrk{}ज्ञितः~॥}\ A}}\skt{~।}\pend

\pstart\versepar \ledsidenote{\passref{4.38.102}}\skt{वचा }\edtext{\skt{मधु च}}{\lemma{\skt{मधु च}}\Afootnote{K H; \skt{मधुक}\ A}}\skt{ तैलञ्च क्वाथः }\edtext{\skt{सचससैन्धवः}}{\lemma{\skt{सचससैन्धवः}}\Afootnote{K H; \skt{सरससैन्धवः~।}\ A}}\skt{~॥ }\pend
\pstart\versepar \skt{}\edtext{\skt{फल\-पि\-प्प\-लि\-सं\-युक्ते}}{\lemma{\skt{फल\appbrk{}पि\appbrk{}प्प\appbrk{}लि\appbrk{}सं\appbrk{}युक्ते}}\Afootnote{K H; \skt{पिप्प\appbrk{}ली\appbrk{}फ\appbrk{}ल\appbrk{}सं\appbrk{}युक्तो}\ A}}\skt{ }\edtext{\skt{वस्तिर्युक्तरथः}}{\lemma{\skt{वस्तिर्}}\Afootnote{H; \skt{बस्तिर्}\ A K}}\skt{ स्मृतः~।}\pend

\pstart\versepar \ledsidenote{\passref{4.38.103}}\skt{सुरदारु }\edtext{\skt{वचा}}{\lemma{\skt{वचा}}\Afootnote{K H; \skt{वरा}\ A}}\skt{ रास्ना शतपुष्पा }\edtext{\skt{तथैव च}}{\lemma{\skt{तथैव च}}\Afootnote{K H; \skt{वचा मधु~।}\ A}}\skt{~॥ }\pend
\pstart\versepar \skt{}\edtext{\skt{हिंगु\-सै\-न्ध\-व\-सं\-यु\-क्तो\-ब\-स्ति\-र्द्दो\-ष\-हरः}}{\lemma{\skt{हिंगु\appbrk{}सै\appbrk{}न्ध\appbrk{}व\appbrk{}सं\appbrk{}यु\appbrk{}क्तो\appbrk{}ब\appbrk{}स्ति\appbrk{}र्द्दो\appbrk{}ष\appbrk{}हरः}}\Afootnote{H; \skt{हिङ्गु\appbrk{}सै\appbrk{}न्ध\appbrk{}व\appbrk{}सं\appbrk{}युक्तो बस्तिर्दोषहरः}\ A K}}\skt{ स्मृतः~।}\pend

\pstart\versepar \ledsidenote{\passref{4.38.104}}\edtext{\skt{पञ्च\-मू\-ली\-क\-षा\-यन्तु}}{\lemma{\skt{पञ्च\appbrk{}मू\appbrk{}ली\appbrk{}क\appbrk{}षा\appbrk{}यन्तु}}\Afootnote{H; \skt{पञ्च\appbrk{}मू\appbrk{}ली\appbrk{}क\appbrk{}षायं च}\ A; \skt{पञ्च\appbrk{}मू\appbrk{}ली\appbrk{}क\appbrk{}षा\appbrk{}यान्तु}\ K}}\skt{ तैलम्मागधिका }\edtext{\skt{मधुः}}{\lemma{\skt{मधुः}}\Afootnote{K H; \skt{मधु~।}\ A}}\skt{~॥ }\pend
\pstart\versepar \skt{}\edtext{\skt{बस्तिरेष}}{\lemma{\skt{बस्तिरेष}}\Afootnote{K H; \skt{बस्तिरेष}\ A}}\skt{ विधातव्यः सशताह्वः ससैन्धवः~।}\pend

\pstart\versepar \ledsidenote{\passref{4.38.105}}\skt{यव\-को\-ल\-कु\-ल\-त्थानां क्वाथो }\edtext{\skt{मागधिकायुतः}}{\lemma{\skt{मागधिकायुतः}}\Afootnote{K H; \skt{मागधिका मधु~।}\ A}}\skt{~॥ }\pend
\pstart\versepar \skt{ससैन्धवः }\edtext{\skt{समधुकः}}{\lemma{\skt{समधुकः}}\Afootnote{K H; \skt{सयष्ट्याह्वः}\ A}}\skt{ }\edtext{\edtext{\skt{सिद्धबस्तिरिति}}{\lemma{\skt{सिद्धबस्तिर्}}\Afootnote{A H; \skt{सिद्धबस्ति}\ K}}}{\lemma{\skt{इति}}\Afootnote{A H; \skt{रिति}\ K}}\skt{ स्मृतः~।}\pend

\pstart\versepar \ledsidenote{\passref{4.38.112}}\skt{अवेक्ष्य भेषजम्बुद्ध्या विकारञ्च विकारवित्~॥ }\pend
\pstart\versepar \skt{बीजेनानेन }\edtext{\skt{मति\-मा\-न्कु\-र्या\-द्ब\-स्ति\-श\-ता\-न्यपि}}{\lemma{\skt{मति\appbrk{}मा\appbrk{}न्कु\appbrk{}र्या\appbrk{}द्ब\appbrk{}स्ति\appbrk{}श\appbrk{}ता\appbrk{}न्यपि}}\Afootnote{H; \skt{शास्त्रज्ञः कुर्या\appbrk{}द्ब\appbrk{}स्ति\appbrk{}श\appbrk{}ता\appbrk{}न्यपि~॥}\ A; \skt{मतिमां कुर्या\appbrk{}द्ब\appbrk{}स्ति\appbrk{}श\appbrk{}ता\appbrk{}न्यपि~।}\ K}}\skt{~।}\pend

\pstart\versepar \ledsidenote{\passref{4.38.113}}\skt{अजीर्ण्णेन प्रयुञ्जीत दिवास्वप्नञ्च वर्ज्जयेत्~॥ }\pend
\pstart\versepar \skt{आहारञ्चारिकं }\edtext{\edtext{\skt{शेषमन्यत्कामं}}{\lemma{\skt{शेषम्}}\Afootnote{K H; \skt{शेष}\ A}}}{\lemma{\skt{अन्यत्कामं}}\Afootnote{K H; \skt{मन्यत्कामं}\ A}}\skt{ समाचरेत्~।}\edtext{}{\lemma{\skt{वर्ज्जयेत्~॥ आहारञ्}}\Afootnote{K H; \skt{वर्जयेत्~। आहारा}\ A}}\pend

\pstart\versepar \ledsidenote{\passref{4.38.114}}\edtext{\skt{यस्मात्मधु च}}{\lemma{\skt{यस्मात्मधु च}}\Afootnote{K H; \skt{यस्मान्मधु च}\ A}}\skt{ तैलञ्च }\edtext{\skt{प्राधान्येनात्र}}{\lemma{\skt{प्राधान्येनात्र}}\Afootnote{K H; \skt{प्राधान्येन प्र}\ A}}\skt{ }\edtext{\skt{वर्त्तते}}{\lemma{\skt{वर्त्तते}}\Afootnote{K H; \skt{दीयते~।}\ A}}\skt{~॥ }\pend
\pstart\versepar \skt{माधुतैलिक }\edtext{\edtext{\skt{इत्येष}}{\lemma{\skt{इत्य्}}\Afootnote{A K; \skt{इ}\ H}}}{\lemma{\skt{एष}}\Afootnote{K; \skt{एवं भिष\appbrk{}ग्भि\appbrk{}र्ब\appbrk{}स्ति\appbrk{}रु\appbrk{}च्यते~॥}\ A; \skt{त्येष}\ H}}\skt{ }\edtext{\skt{संज्ञितो}}{\lemma{\skt{संज्ञितो}}\Afootnote{K H; \emph{om.}\ A}}\skt{ }\edtext{\skt{बस्तिचिन्तकैः}}{\lemma{\skt{बस्तिचिन्तकैः}}\Afootnote{K H; \emph{om.}\ A}}\skt{~।}\pend

\pstart\versepar \ledsidenote{\passref{4.38.115}}\skt{रथेष्वपि च युक्तेषु हस्त्यश्वे }\edtext{\skt{वापि}}{\lemma{\skt{वापि}}\Afootnote{H; \skt{चापि}\ A K}}\skt{ कल्पिते~॥ }\pend
\pstart\versepar \skt{यस्मान्न }\edtext{\skt{प्रति\-षि\-द्धे\-य\-मतो}}{\lemma{\skt{प्रति\appbrk{}षि\appbrk{}द्धे\appbrk{}य\appbrk{}मतो}}\Afootnote{K H; \skt{प्रति\appbrk{}षि\appbrk{}द्धोऽ\appbrk{}य\appbrk{}मतो}\ A}}\skt{ युक्तरथः स्मृतः~।}\pend

\pstart\versepar \ledsidenote{\passref{4.38.116}}\edtext{\skt{बस्तौ}}{\lemma{\skt{बस्तौ}}\Afootnote{K H; \skt{बलोपच}\ A}}\skt{ }\edtext{\skt{युक्तरथे}}{\lemma{\skt{युक्तरथे}}\Afootnote{H; \skt{यवर्णानां}\ A; \skt{युक्तरथै}\ K}}\skt{ }\edtext{\skt{यस्माद्ध्रुवासिद्धिः}}{\lemma{\skt{यस्माद्ध्रुवासिद्धिः}}\Afootnote{K; \skt{यस्मा\appbrk{}द्व्या\appbrk{}धि\appbrk{}श\appbrk{}तस्य च~। भवत्येतेन सिद्धिस्}\ A; \skt{यस्माद्ध्रुवासिद्धिः}\ H}}\skt{ }\edtext{\skt{प्रकीर्त्तिताः}}{\lemma{\skt{प्रकीर्त्तिताः}}\Afootnote{K H; \skt{तु}\ A}}\skt{~॥ }\pend
\pstart\versepar \skt{}\edtext{\skt{सिद्धबस्तिरिति}}{\lemma{\skt{सिद्धबस्तिरिति}}\Afootnote{K H; \skt{सिद्धबस्तिर्}\ A}}\skt{ }\edtext{\skt{प्रोक्ता मुनि\-भि\-र्न्नै\-ग\-मै\-षिभिः~।}}{\lemma{\skt{रिति… न्नैगमैषिभिः}}\Afootnote{\emph{om.}\ A}}\pend

\pstart\versepar \ledsidenote{\passref{4.38.117}}\edtext{\edtext{\edtext{\skt{सुखि\-ना\-म\-ल्प\-दो\-षा\-णा\-न्नि\-त्य\-स्नि\-ग्धाश्च}}{\lemma{\skt{सुखिनाम्}}\Afootnote{K H; \skt{सुखिना}\ A}}}{\lemma{\skt{अल्पदोषाणान्}}\Afootnote{K H; \skt{मल्पदोषाणां}\ A}}\skt{ ये}}{\lemma{\skt{नित्यस्निग्धाश्च ये}}\Afootnote{H; \skt{नित्यं स्निग्धाश्च ये}\ A; \skt{नित्यस्निग्धाश्च ये}\ K}}\skt{ नराः~॥ }\pend
\pstart\versepar \skt{मृदुकोष्ठाश्च ये तेषां विधेया }\edtext{\skt{माधुतैलिकः}}{\lemma{\skt{माधुतैलिकः}}\Afootnote{K H; \skt{माधुतैलिकाः~॥}\ A}}\skt{~।}\pend

\pstart\versepar \ledsidenote{\passref{4.38.118}}\skt{मृदु\-त्वा\-त्पा\-द\-ही\-न\-त्वा\-द\-कृ\-त्स्न\-वि\-धि\-से\-व\-नात्~॥ }\pend
\pstart\versepar \skt{एक\-ब\-स्ति\-प्र\-दा\-नाच्च }\edtext{\edtext{\skt{सिद्ध\-ब\-स्ति\-ष्व\-यन्त्र}}{\lemma{\skt{सिद्धबस्तिष्व्}}\Afootnote{K H; \skt{सिद्धबस्ति}\ A}}}{\lemma{\skt{अयन्त्र}}\Afootnote{K H; \skt{ष्वयन्त्र}\ A}}\skt{ }\edtext{\skt{णेति}}{\lemma{\skt{णेति}}\Afootnote{K H; \skt{णा~॥}\ A}}\skt{~॥}\pend

\pstart \ledsidenote{\passref{4.38.trailer}}\edtext{\skt{॥ }\edtext{\edtext{\skt{चिकि\-त्सा\-स्था\-न\-म\-ष्ट\-त्रिं\-श\-त्त\-मोऽ\-ध्यायः}}{\lemma{\skt{अष्टत्रिंशत्तमो}}\Afootnote{H; \skt{माष्टत्रिंशो}\ A; \emph{om.}\ K}}}{\lemma{\skt{ऽध्यायः}}\Afootnote{A H; \emph{om.}\ K}}}{\lemma{\skt{॥ चिकित्सास्थानम्}}\Afootnote{H; \skt{इति सुश्रु\appbrk{}त\appbrk{}सं\appbrk{}हि\appbrk{}तायां चिकित्सास्थाने निरू\appbrk{}ह\appbrk{}क्र\appbrk{}म\appbrk{}चि\appbrk{}कि\appbrk{}त्सितं ना}\ A; \skt{चिलह्रा~॥ {\symfont ❈}~॥}\ K}}\skt{~॥}\pend

\endnumbering
\beginnumbering
\eledsection*{Adhyāya 39}
\pstart \ledsidenote{\passref{4.39.1}}\skt{अथात }\edtext{\skt{आतु\-रो\-प\-द्र\-व\-कि\-त्सितं}}{\lemma{\skt{आतु\appbrk{}रो\appbrk{}प\appbrk{}द्र\appbrk{}व\appbrk{}कि\appbrk{}त्सितं}}\Afootnote{\skt{आतु\appbrk{}रो\appbrk{}प\appbrk{}द्र\appbrk{}व\appbrk{}चि\appbrk{}कि\appbrk{}त्सितं}\ A K H}}\skt{ }\edtext{\skt{व्याख्यास्यामः}}{\lemma{\skt{व्याख्यास्यामः}}\Afootnote{A H; \skt{व्या व्याख्यास्यामः~॥}\ K}}\skt{~॥}\pend

\pstart\versepar \ledsidenote{\passref{4.39.3}}\edtext{\skt{स्नेहपीस्य}}{\lemma{\skt{स्नेहपीस्य}}\Afootnote{\skt{स्नेहपीतस्य}\ A; \skt{स्नेहपीडस्य}\ K H}}\skt{ वान्तस्य }\edtext{\skt{विरि\-क्त\-स्या\-श्रु\-ता\-सृजः}}{\lemma{\skt{विरि\appbrk{}क्त\appbrk{}स्या\appbrk{}श्रु\appbrk{}ता\appbrk{}सृजः}}\Afootnote{K H; \skt{विरिक्तस्य स्रुतासृजः~।}\ A}}\skt{~। }\pend
\pstart\versepar \skt{निरूढस्य च }\edtext{\skt{कायाग्निर्मन्दो}}{\lemma{\skt{कायाग्निर्}}\Afootnote{A K; \skt{कायाग्निं}\ H}}\skt{ भवति देहिनः~।}\pend

\pstart\versepar \ledsidenote{\passref{4.39.4}}\edtext{\skt{सो न्नैर\-त्य\-र्थ\-गु\-रु\-भि\-रु\-प\-युक्तैः}}{\lemma{\skt{सो न्नैर्}}\Afootnote{K H; \skt{सोऽग्नैर्}\ A}}\skt{ प्रशाम्यति~। }\pend
\pstart\versepar \skt{अल्पो }\edtext{\skt{महद्भिर्बहुधा}}{\lemma{\skt{बहुधा}}\Afootnote{K; \skt{बहुभिश्}\ A; \skt{व्वसुधा}\ H}}\skt{ }\edtext{\skt{छादिते}}{\lemma{\skt{छादिते}}\Afootnote{K; \skt{छादितो}\ A; \skt{च्छादिते}\ H}}\skt{ }\edtext{\skt{ग्निरिवेन्धनैः}}{\lemma{\skt{ग्निर्}}\Afootnote{K; \skt{ऽग्निर्}\ A H}}\skt{~।}\pend

\pstart\versepar \ledsidenote{\passref{4.39.5}}\skt{स चाल्पै\-र्ल\-घु\-भि\-श्चा\-न्नै\-रु\-प\-यु\-क्तै\-र्वि\-व\-र्धते~। }\pend
\pstart\versepar \skt{काष्ठै\-र\-णु\-भि\-र\-ल्पैश्च सन्धुक्षित इवानलः~।}\pend

\pstart\versepar \ledsidenote{\passref{4.39.6}}\skt{हृत\-दो\-ष\-प्र\-मा\-णेन स }\edtext{\skt{चाहारविधिः}}{\lemma{\skt{चाहारविधिः}}\Afootnote{K H; \skt{दाहारविधिः}\ A}}\skt{ स्मृतः~। }\pend
\pstart\versepar \skt{त्रीणि चात्र प्रमाणानि प्रस्थो }\edtext{\skt{र्धाढ\-क\-मा\-ढ\-कम्}}{\lemma{\skt{र्धाढकम्}}\Afootnote{K; \skt{ऽर्धाढकम्}\ A H}}\skt{~॥}\pend

\pstart\versepar \ledsidenote{\passref{4.39.7}}\skt{तत्रावरं प्रस्थमात्रं द्वे शेषे मध्यमोत्तमे~। }\pend
\pstart\versepar \skt{प्रस्थे }\edtext{\skt{त्वेके}}{\lemma{\skt{त्वेके}}\Afootnote{K\textsuperscript{pc} H; \skt{परि}\ A; \skt{त्वेको}\ K\textsuperscript{ac}}}\skt{ }\edtext{\skt{गते}}{\lemma{\skt{गते}}\Afootnote{K H; \skt{स्रुते}\ A}}\skt{ देया }\edtext{\skt{यवा\-गू\-स्व\-ल्प\-ता\-ण्डुला}}{\lemma{\skt{यवा\appbrk{}गू\appbrk{}स्व\appbrk{}ल्प\appbrk{}ता\appbrk{}ण्डुला}}\Afootnote{K; \skt{यवागूः खल्पतण्डुला~॥}\ A; \skt{यवा\appbrk{}गू\appbrk{}स्व\appbrk{}ल्प\appbrk{}त\appbrk{}ण्डुला~॥}\ H}}\skt{~।}\pend

\pstart\versepar \ledsidenote{\passref{4.39.8}}\skt{द्वे चैवार्धाढके देये तिस्रश्चाप्याढके गते~। }\pend
\pstart\versepar \skt{}\edtext{\edtext{\skt{विले\-पी\-मु\-चि\-ता\-द्भ\-क्ता\-च्च\-तु\-र्थां\-श\-कृ\-ता\-म्भि\-षक्}}{\lemma{\skt{विले\appbrk{}पी\appbrk{}मु\appbrk{}चि\appbrk{}ता\appbrk{}द्भ\appbrk{}क्ता\appbrk{}च्च\appbrk{}तु\appbrk{}र्थां\appbrk{}श\appbrk{}कृ\appbrk{}ताम्}}\Afootnote{A K H\textsuperscript{pc}; \skt{विलेपीमुचिरा द्भक्ता\appbrk{}च्च\appbrk{}तु\appbrk{}र्थां\appbrk{}श\appbrk{}कृ\appbrk{}ताम्}\ H\textsuperscript{ac}}}}{\lemma{\skt{भिषक्}}\Afootnote{K H; \skt{ततः~॥}\ A}}\skt{~।}\pend

\pstart\versepar \ledsidenote{\passref{4.39.9}}\edtext{\edtext{\skt{दद्यात्क्लिन्नन्ततः}}{\lemma{\skt{दद्यात्क्लिन्नन्}}\Afootnote{K H; \skt{दद्यादुक्ते}\ A}}}{\lemma{\skt{ततः}}\Afootnote{K H; \skt{न वि}\ A}}\skt{ }\edtext{\skt{पश्चा\-त्मु\-द्ग\-यू\-षेण}}{\lemma{\skt{पश्चा\appbrk{}त्मु\appbrk{}द्ग\appbrk{}यू\appbrk{}षेण}}\Afootnote{\skt{धिना क्लिन्नसिक्थाम}\ A; \skt{पश्चा\appbrk{}त्मु\appbrk{}द्ग\appbrk{}यू\appbrk{}षेतो}\ K; \skt{पश्चा\appbrk{}त्मु\appbrk{}द्ग\appbrk{}यू\appbrk{}षेण}\ H}}\skt{ }\edtext{\skt{भोजनं}}{\lemma{\skt{भोजनं}}\Afootnote{K H; \skt{पिच्छिलाम्~।}\ A}}\skt{~। }\pend
\pstart\versepar \skt{}\edtext{\skt{त्रिभागकलितं}}{\lemma{\skt{त्रिभागकलितं}}\Afootnote{K\textsuperscript{pc} H; \skt{अग्निग्धलवणं स्वच्छ\appbrk{}मु\appbrk{}द्ग\appbrk{}यू\appbrk{}ष\appbrk{}युतं}\ A; \skt{अर्द्धां\appbrk{}श\appbrk{}मु\appbrk{}चि\appbrk{}ता\appbrk{}द्भक्ता}\ K\textsuperscript{ac}}}\skt{ }\edtext{\skt{सम्य\-ग\-कृ\-ते\-ना\-र्द्ध\-स\-म्मितं}}{\lemma{\skt{सम्य\appbrk{}ग\appbrk{}कृ\appbrk{}ते\appbrk{}ना\appbrk{}र्द्ध\appbrk{}स\appbrk{}म्मितं}}\Afootnote{K\textsuperscript{pc} H; \skt{ततः~॥}\ A; \skt{दकृताख्येन बुद्धिमां~।}\ K\textsuperscript{ac}}}\skt{~।}\pend

\pstart \ledsidenote{\passref{4.39.10b}}\skt{ततस्तु कृतसंज्ञेन हृद्ये\-ने\-न्द्रि\-य\-बो\-धिना~।}\pend

\pstart\versepar \ledsidenote{\passref{4.39.11}}\edtext{\skt{त्रीनङ्शाद्वितरे }\edtext{\skt{भोक्तु\-मा\-तु\-रा\-योनं}}{\lemma{\skt{आतुरायोनं}}\Afootnote{\skt{आतुरायौदनं}\ A; \skt{आतुरायोदनं}\ K H}}}{\lemma{\skt{त्रीनङ्शाद्वितरे भोक्तुम्}}\Afootnote{K; \skt{त्रीनं\appbrk{}शा\appbrk{}न्वि\appbrk{}त\appbrk{}रे\appbrk{}द्भो\appbrk{}क्तुम्}\ A; \skt{त्रीनं\appbrk{}शा\appbrk{}द्वि\appbrk{}त\appbrk{}रे\appbrk{}द्भो\appbrk{}क्तुंम्}\ H\textsuperscript{ac} H\textsuperscript{pc}}}\skt{ }\edtext{\skt{प्रति}}{\lemma{\skt{प्रति}}\Afootnote{K H; \skt{मृदु~।}\ A}}\skt{~। }\pend
\pstart\versepar \skt{ततो यथोचितं भक्तं भोक्तुमस्मै विचक्षणः~।}\pend

\pstart\versepar \ledsidenote{\passref{4.39.12}}\edtext{\skt{लावै\-न\-ह\-रि\-णा\-दीनां}}{\lemma{\skt{लावै\appbrk{}न\appbrk{}ह\appbrk{}रि\appbrk{}णा\appbrk{}दीनां}}\Afootnote{K; \skt{लावै\appbrk{}ण\appbrk{}ह\appbrk{}रि\appbrk{}णा\appbrk{}दीनां}\ A H}}\skt{ रसै\-र्द\-द्या\-त्सु\-सं\-स्कृतैः~। }\pend
\pstart\versepar \skt{हीन\-म\-ध्यो\-त्त\-मे\-ष्वेषु विरेकेषु प्रकीर्त्तितः~।}\pend

\pstart \ledsidenote{\passref{4.39.13a}}\skt{एकद्वित्रिगुणः सम्यगाहारस्य क्रमो }\edtext{\skt{हितः}}{\lemma{\skt{हितः}}\Afootnote{K H; \skt{त्वयम्~।}\ A}}\skt{~।}\pend

\pstart \ledsidenote{\passref{4.39.14b}}\edtext{\skt{वेद\-ना\-ला\-भ\-नि\-य\-म\-शो\-क\-वै\-चि\-त्य\-हे\-तुभिः}}{\lemma{\skt{वेद\appbrk{}ना\appbrk{}ला\appbrk{}भ\appbrk{}नि\appbrk{}य\appbrk{}म\appbrk{}शो\appbrk{}क\appbrk{}वै\appbrk{}चि\appbrk{}त्य\appbrk{}हे\appbrk{}तुभिः}}\Afootnote{\skt{वेद\appbrk{}ना\appbrk{}ला\appbrk{}भ\appbrk{}नि\appbrk{}य\appbrk{}म\appbrk{}शो\appbrk{}क\appbrk{}वै\appbrk{}चि\appbrk{}त्त्य\appbrk{}हे\appbrk{}तुभिः}\ A K H}}\skt{~।}\pend

\pstart\versepar \ledsidenote{\passref{4.39.15}}\edtext{\skt{नरा\-नु\-पो\-षि\-तं\-श्चापि}}{\lemma{\skt{नरा\appbrk{}नु\appbrk{}पो\appbrk{}षि\appbrk{}तंश्}}\Afootnote{K; \skt{नरा\appbrk{}नु\appbrk{}पो\appbrk{}षि\appbrk{}तांश्}\ A H}}\skt{ विरि\-क्त\-व\-दु\-पा\-च\-रेत्~। }\pend
\pstart\versepar \skt{आढ\-का\-र्धा\-ढ\-क\-प्र\-स्थ\-संख्या ह्येषा विरेचने~।}\pend

\pstart\versepar \ledsidenote{\passref{4.39.16}}\edtext{\skt{श्लेष्मान्तो हि}}{\lemma{\skt{श्लेष्मान्तो हि}}\Afootnote{K H; \skt{श्लेष्मान्तत्वा}\ A}}\skt{ }\edtext{\skt{विरेकः}}{\lemma{\skt{विरेकः}}\Afootnote{K H; \skt{द्विरेक}\ A}}\skt{ }\edtext{\skt{स्यादेतदिच्छन्ति}}{\lemma{\skt{स्यादेतदिच्छन्ति}}\Afootnote{K H; \skt{स्य न तामिच्छति}\ A}}\skt{ }\edtext{\skt{केचनः}}{\lemma{\skt{केचनः}}\Afootnote{K H; \skt{तद्विदः~।}\ A}}\skt{~। }\pend
\pstart\versepar \skt{एको विरेकः श्लेष्मान्तो न द्वितीयः }\edtext{\skt{कदाचनः}}{\lemma{\skt{कदाचनः}}\Afootnote{K; \skt{ऽस्ति कश्चन~॥}\ A; \skt{कदाचन~॥}\ H}}\skt{~।}\pend

\pstart\versepar \ledsidenote{\passref{4.39.17}}\skt{बलं यत्त्रिविधं }\edtext{\skt{प्रोक्तमतस्तत्र}}{\lemma{\skt{अतस्}}\Afootnote{A K; \skt{मतस्}\ H}}\skt{ क्रमस्त्रिधा~। }\pend
\pstart\versepar \skt{}\edtext{\edtext{\skt{तत्रा\-नु\-क्र\-म\-मे\-कन्तु}}{\lemma{\skt{तत्रानुक्रमम्}}\Afootnote{K H; \skt{तत्रानुक्रम}\ A}}}{\lemma{\skt{एकन्तु}}\Afootnote{K H; \skt{मेकं तु}\ A}}\skt{ }\edtext{\skt{बलस्थे}}{\lemma{\skt{बलस्थे}}\Afootnote{K H; \skt{बलस्थः}\ A}}\skt{ }\edtext{\skt{सकृदाचरेत्}}{\lemma{\skt{सकृदाचरेत्}}\Afootnote{K H; \skt{सकृदाचरेत्~॥}\ A}}\skt{~।}\pend

\pstart\versepar \ledsidenote{\passref{4.39.18}}\edtext{\skt{मध्यमे द्वौ}}{\lemma{\skt{मध्यमे द्वौ}}\Afootnote{K H; \skt{द्विराचरेन्मध्य}\ A}}\skt{ }\edtext{\skt{क्रमौ}}{\lemma{\skt{क्रमौ}}\Afootnote{K H; \emph{om.}\ A}}\skt{ }\edtext{\edtext{\skt{कुर्या\-त्त्री\-न्वा\-रा\-न्दु\-र्बले}}{\lemma{\skt{कुर्यात्त्रीन्वारान्}}\Afootnote{K H; \skt{बलस्त्रीन्वारान्}\ A}}}{\lemma{\skt{दुर्बले}}\Afootnote{K H; \skt{दुर्बलस्}\ A}}\skt{ तथा~। }\pend
\pstart\versepar \skt{केचिदेवं क्रमं प्राहु\-र्म\-न्द\-म\-ध्यो\-त्त\-मा\-ग्निषु}\pend

\pstart\versepar \ledsidenote{\passref{4.39.21}}\skt{केवलं }\edtext{\skt{स्नेहपीत्तो वा}}{\lemma{\skt{स्नेहपीत्तो वा}}\Afootnote{K; \skt{स्नेहपीतो वा}\ A H}}\skt{ वान्तो यश्चापि केवलं~। }\pend
\pstart\versepar \skt{स सप्तरात्रं मनुजो भुञ्जीत लघु भोजनं~।}\pend

\pstart\versepar \ledsidenote{\passref{4.39.22}}\edtext{\skt{सिरावेधकृतो यश्च}}{\lemma{\skt{सिरावेधकृतो यश्च}}\Afootnote{K; \skt{कृतः सिराव्यधो यस्य कृतं यस्य}\ A; \skt{शिरावेधकृतो यश्च}\ H}}\skt{ }\edtext{\skt{विरिक्तो यश्च}}{\lemma{\skt{विरिक्तो यश्च}}\Afootnote{K; \skt{च}\ A; \skt{विविक्तो यश्च}\ H}}\skt{ }\edtext{\skt{मानवः}}{\lemma{\skt{मानवः}}\Afootnote{K H; \skt{शोधनम्~।}\ A}}\skt{~। }\pend
\pstart\versepar \skt{}\edtext{\edtext{\edtext{\skt{मास\-म्प\-रि\-ह\-रे\-द\-ग्नि\-र्या\-वन्न}}{\lemma{\skt{मासम्}}\Afootnote{K H; \skt{स ना}\ A}}}{\lemma{\skt{परि\appbrk{}ह\appbrk{}रे\appbrk{}द\appbrk{}ग्निर्}}\Afootnote{K H; \skt{परिहरेन्मासं}\ A}}}{\lemma{\skt{यावन्न}}\Afootnote{K H; \skt{यावद्वा}\ A}}\skt{ }\edtext{\skt{बलवां}}{\lemma{\skt{बलवां}}\Afootnote{K; \skt{बलवा}\ A H}}\skt{ }\edtext{\skt{भवेत्}}{\lemma{\skt{भवेत्}}\Afootnote{K; \skt{न्भवेत्~॥}\ A H}}\skt{~।}\pend

\pstart\versepar \ledsidenote{\passref{4.39.23}}\skt{त्र्यहं त्र्यहं परि\-ह\-रे\-दे\-कै\-क\-म्ब\-स्ति\-मा\-तुरः~। }\pend
\pstart\versepar \skt{}\edtext{\skt{तृतीये}}{\lemma{\skt{तृतीये}}\Afootnote{K H; \skt{तृतीये तु}\ A}}\skt{ }\edtext{\skt{परिहारे तु}}{\lemma{\skt{परिहारे तु}}\Afootnote{K H; \skt{परीहारे}\ A}}\skt{ यथायोगं समाचरेत्~।}\pend

\pstart\versepar \ledsidenote{\passref{4.39.24}}\edtext{\skt{तैल\-पा\-त्रा\-म\-मृ\-द्भा\-ण्ड\-स\-ध\-र्माणो}}{\lemma{\skt{तैल\appbrk{}पा\appbrk{}त्रा\appbrk{}म\appbrk{}मृ\appbrk{}द्भा\appbrk{}ण्ड\appbrk{}स\appbrk{}ध\appbrk{}र्माणो}}\Afootnote{K H; \skt{तैल\appbrk{}पू\appbrk{}र्णा\appbrk{}मृ\appbrk{}द्भा\appbrk{}ण्ड\appbrk{}स\appbrk{}ध\appbrk{}र्माणो}\ A}}\skt{ व्रणातुराः~। }\pend
\pstart\versepar \skt{}\edtext{\skt{स्निग्ध\-शु\-द्धा\-ग्नि\-रो\-गार्त्ता}}{\lemma{\skt{स्निग्ध\appbrk{}शु\appbrk{}द्धा\appbrk{}ग्नि\appbrk{}रो\appbrk{}गार्त्ता}}\Afootnote{K H; \skt{स्निग्ध\appbrk{}शु\appbrk{}द्धा\appbrk{}क्षि\appbrk{}रो\appbrk{}गार्ता}\ A}}\skt{ ज्वरा\-ती\-सा\-रि\-णश्च ये~।}\pend

\pstart\versepar \ledsidenote{\passref{4.39.25}}\skt{क्रुध्यतः }\edtext{\skt{कुपितः}}{\lemma{\skt{कुपितः}}\Afootnote{K H; \skt{कुपितं}\ A}}\skt{ पित्तं }\edtext{\skt{कुर्या\-न्स्ता\-न्स्ता\-नु\-प\-द्र\-वान्}}{\lemma{\skt{कुर्यान्स्तान्स्}}\Afootnote{K; \skt{कुर्यात्तांस्}\ A H}}\skt{~। }\pend
\pstart\versepar \skt{आयास्यतः शोचतो वा चित्तं }\edtext{\skt{विभ्रममर्च्छति}}{\lemma{\skt{अर्च्छति}}\Afootnote{K H; \skt{ऋच्छति~॥}\ A}}\skt{~।}\pend

\pstart\versepar \ledsidenote{\passref{4.39.26}}\edtext{\skt{मैथु\-नो\-प\-ग\-मा\-द्घोरां}}{\lemma{\skt{मैथु\appbrk{}नो\appbrk{}प\appbrk{}ग\appbrk{}मा\appbrk{}द्घोरां}}\Afootnote{K; \skt{मैथु\appbrk{}नो\appbrk{}प\appbrk{}ग\appbrk{}मा\appbrk{}द्घोरा}\ A H}}\skt{ }\edtext{\skt{व्याधीनर्च्छति}}{\lemma{\skt{व्याधीनर्च्छति}}\Afootnote{K; \skt{न्व्याधीनाप्नोति}\ A; \skt{न्व्याधीनर्च्छति}\ H}}\skt{ दुर्मतिः~। }\pend
\pstart\versepar \skt{आक्षेपकं }\edtext{\skt{पक्षा\-घा\-त\-म\-ङ्ग\-प्र\-ग्र\-ह\-मेव}}{\lemma{\skt{पक्षाघातम्}}\Afootnote{K; \skt{पक्षघातम्}\ A H}}\skt{ च~।}\pend

\pstart\versepar \ledsidenote{\passref{4.39.27}}\skt{गुह्यप्रदेशे श्वयथुं }\edtext{\skt{काशश्वासौ च}}{\lemma{\skt{काशश्वासौ च}}\Afootnote{K; \skt{कासश्वासौ च}\ A H}}\skt{ दारुणौ~। }\pend
\pstart\versepar \skt{}\edtext{\skt{शुक्रवच्चास्य}}{\lemma{\skt{शुक्रवच्चास्य}}\Afootnote{K H; \skt{रुधिरं शुक्रवच्चा}\ A}}\skt{ }\edtext{\skt{रुधिरं}}{\lemma{\skt{रुधिरं}}\Afootnote{K; \skt{पि}\ A; \skt{रुधिरां}\ H}}\skt{ }\edtext{\skt{सरुजं}}{\lemma{\skt{सरुजं}}\Afootnote{K H; \skt{सरज}\ A}}\skt{ }\edtext{\skt{सम्प्रवर्त्तते}}{\lemma{\skt{सम्प्रवर्त्तते}}\Afootnote{K H; \skt{स्कं प्रवर्तते~॥}\ A}}\skt{~।}\pend

\pstart\versepar \ledsidenote{\passref{4.39.28}}\edtext{\skt{दिवास्वप्नात्तु}}{\lemma{\skt{दिवास्वप्नात्तु}}\Afootnote{K H; \skt{लभते च दिवास्वप्ना}\ A}}\skt{ }\edtext{\skt{लभते तांस्ता}}{\lemma{\skt{लभते तांस्ता}}\Afootnote{K H; \skt{त्तांस्ता}\ A}}\skt{ }\edtext{\skt{व्याधीन्कफात्मकां}}{\lemma{\skt{व्याधीन्कफात्मकां}}\Afootnote{K; \skt{न्व्याधी\appbrk{}न्क\appbrk{}फा\appbrk{}त्म\appbrk{}कान्~। प्लीहोदरं}\ A; \skt{न्व्याधी\appbrk{}न्क\appbrk{}फा\appbrk{}त्म\appbrk{}कान्~। प्र}\ H}}\skt{~। }\pend
\pstart\versepar \skt{}\edtext{\skt{प्रतिश्यायं}}{\lemma{\skt{प्रतिश्यायं}}\Afootnote{A K; \emph{om.}\ H}}\skt{ }\edtext{\skt{ज्वरं}}{\lemma{\skt{ज्वरं}}\Afootnote{K; \skt{पाण्डुतां श्वयथुं ज्वरम्~॥}\ A; \skt{तिश्यायञ्ज्वरं}\ H}}\skt{ }\edtext{\skt{शोफं}}{\lemma{\skt{शोफं}}\Afootnote{K H; \emph{om.}\ A}}\skt{ }\edtext{\skt{पाण्डुरोगं}}{\lemma{\skt{पाण्डुरोगं}}\Afootnote{K H; \emph{om.}\ A}}\skt{ }\edtext{\skt{प्लिहोदरं}}{\lemma{\skt{प्लिहोदरं}}\Afootnote{K H; \emph{om.}\ A}}\skt{~।}\pend

\pstart\versepar \ledsidenote{\passref{4.39.29}}\skt{मोहं सदनमङ्गानां }\edtext{\edtext{\skt{सावि\-पा\-क\-म\-रो\-चकं}}{\lemma{\skt{साविपाकम्}}\Afootnote{K H; \skt{अविपाकं}\ A}}}{\lemma{\skt{अरोचकं}}\Afootnote{K H; \skt{तथाऽरुचिम्~।}\ A}}\skt{~। }\pend
\pstart\versepar \skt{तमसा }\edtext{\skt{चाभिभूतश्च}}{\lemma{\skt{चाभिभूतश्च}}\Afootnote{K; \skt{चाभिभूतस्तु}\ A; \skt{चाभिभूताश्च}\ H}}\skt{ स्वप्न\-मे\-वा\-भि\-न\-न्दति~।}\pend

\pstart\versepar \ledsidenote{\passref{4.39.30}}\skt{उच्चैः सम्भाषणाद्वायुः शिर\-स्या\-पा\-द\-ये\-द्रुजं~। }\pend
\pstart\versepar \skt{}\edtext{\skt{आन्ध्य\-म्वा\-धि\-र्य\-म\-घ्रत्वं}}{\lemma{\skt{वाधिर्यमघ्रत्वं}}\Afootnote{K H; \skt{जाड्यमजिघ्रत्वं}\ A}}\skt{ }\edtext{\skt{मूकता}}{\lemma{\skt{मूकता}}\Afootnote{K H; \skt{बाधिर्यं मूकतां}\ A}}\skt{ }\edtext{\edtext{\skt{जडतामपि}}{\lemma{\skt{जडताम्}}\Afootnote{K H; \skt{तथा~॥}\ A}}}{\lemma{\skt{अपि}}\Afootnote{K H; \emph{om.}\ A}}\skt{~।}\pend

\pstart \ledsidenote{\passref{4.39.31a}}\edtext{\skt{अर्दितं}}{\lemma{\skt{अर्दितं}}\Afootnote{K H; \emph{om.}\ A}}\skt{ }\edtext{\edtext{\edtext{\skt{सह\-नु\-म्मो\-क्ष\-म\-धी\-मन्थ}}{\lemma{\skt{सहनुम्}}\Afootnote{K; \skt{हनु}\ A; \skt{सहनु}\ H}}}{\lemma{\skt{मोक्षम्}}\Afootnote{K; \skt{मोक्ष}\ A; \skt{र्म्मोक्षम्}\ H}}\skt{ दारुणं}}{\lemma{\skt{अधीमन्थ दारुणं}}\Afootnote{\skt{मधी\appbrk{}म\appbrk{}न्थ\appbrk{}म\appbrk{}र्दितं च सुदारुणम्~।}\ A; \skt{अधीमन्थञ्च दारुणं}\ K H}}\skt{~।}\pend

\pstart\versepar \ledsidenote{\passref{4.39.32}}\skt{लभते दन्तचालञ्च तांस्तां\-श्चा\-न्या\-नु\-प\-द्र\-वान्~। }\pend
\pstart\versepar \skt{यानयानात्तु लभते }\edtext{\skt{च्छर्दिं}}{\lemma{\skt{च्छर्दिं}}\Afootnote{K H; \skt{छर्दि}\ A}}\skt{ }\edtext{\skt{मूर्च्छा}}{\lemma{\skt{मूर्च्छा}}\Afootnote{A K; \skt{मूर्च्छाम्}\ H}}\skt{ }\edtext{\skt{भ्रमं}}{\lemma{\skt{भ्रमं}}\Afootnote{K H; \skt{भ्रम}\ A}}\skt{ }\edtext{\skt{क्लमं}}{\lemma{\skt{क्लमं}}\Afootnote{K H; \skt{क्लमान्~॥}\ A}}\skt{~।}\edtext{}{\lemma{\skt{चान्या\appbrk{}नु\appbrk{}प\appbrk{}द्र\appbrk{}वान्~। यानयानात्तु}}\Afootnote{K H; \skt{चान्या\appbrk{}नु\appbrk{}प\appbrk{}द्र\appbrk{}वान्~। यानयानेन}\ A}}\pend

\pstart\versepar \ledsidenote{\passref{4.39.33}}\edtext{\skt{तथाङ्गप्रग्रहं}}{\lemma{\skt{तथाङ्गप्रग्रहं}}\Afootnote{K H; \skt{तथैवाङ्गग्रहं}\ A}}\skt{ घोर\-मि\-न्द्रि\-या\-णाञ्च विभ्रमं~। }\pend
\pstart\versepar \skt{}\edtext{\skt{चिरासनात्तथा}}{\lemma{\skt{चिरासनात्तथा}}\Afootnote{A K; \skt{चिरासनान्तथा}\ H}}\skt{ स्थानाच्छ्रोण्यां भवति वेदना~॥}\pend

\pstart\versepar \ledsidenote{\passref{4.39.34}}\edtext{\skt{अति\-च\-ङ्क्र\-म\-णा\-द्वा\-यु\-र्जं\-घोर्व्वोः}}{\lemma{\skt{जंघोर्व्वोः}}\Afootnote{K H; \skt{जङ्घयोः}\ A}}\skt{ कुरुते }\edtext{\skt{रुजः}}{\lemma{\skt{रुजः}}\Afootnote{A K; \skt{रुजाः~॥}\ H}}\skt{~। }\pend
\pstart\versepar \skt{}\edtext{\skt{सक्थ्नोः}}{\lemma{\skt{सक्थ्नोः}}\Afootnote{K H; \skt{सक्थि}\ A}}\skt{ प्रशोषं शोफम्वा पाद\-ह\-र्ष\-म\-थापि वा~।}\pend

\pstart \ledsidenote{\passref{4.39.35a}}\skt{शीत\-स\-म्भो\-ग\-तो\-यानां सेवा मारुतवृद्धये~।}\pend

\pstart\versepar \ledsidenote{\passref{4.39.36}}\skt{वाता\-त\-पा\-भ्या\-म्वै\-वर्ण्यं ज्वरं }\edtext{\skt{वापि}}{\lemma{\skt{वापि}}\Afootnote{K H; \skt{चापि}\ A}}\skt{ }\edtext{\skt{निगच्छति}}{\lemma{\skt{निगच्छति}}\Afootnote{K; \skt{समाप्नुया}\ A; \skt{निग च्छति~।}\ H\textsuperscript{ac}; \skt{नि यच्छति~।}\ H\textsuperscript{pc}}}\skt{~॥ }\pend
\pstart\versepar \skt{}\edtext{\skt{विरु\-द्धा\-ध्य\-श\-ना\-न्मृत्युं}}{\lemma{\skt{विरु\appbrk{}द्धा\appbrk{}ध्य\appbrk{}श\appbrk{}ना\appbrk{}न्मृत्युं}}\Afootnote{K; \skt{त्~। विरु\appbrk{}द्धा\appbrk{}ध्य\appbrk{}श\appbrk{}ना\appbrk{}न्मृत्युं}\ A; \skt{निरु\appbrk{}द्धा\appbrk{}ध्य\appbrk{}श\appbrk{}ना\appbrk{}त्मृ\appbrk{}त्युम्}\ H}}\skt{ व्याधिं वा }\edtext{\skt{घोरमर्च्छति}}{\lemma{\skt{अर्च्छति}}\Afootnote{K H; \skt{ऋच्छति~॥}\ A}}\skt{~।}\pend

\pstart\versepar \ledsidenote{\passref{4.39.37}}\edtext{\skt{असा\-त्म्य\-भो\-ज\-ना\-द्ध\-न्या\-द्ब\-ल\-व\-र्ण्णा\-व\-सं\-शयं}}{\lemma{\skt{असा\appbrk{}त्म्य\appbrk{}भो\appbrk{}ज\appbrk{}ना\appbrk{}द्ध\appbrk{}न्या\appbrk{}द्ब\appbrk{}ल\appbrk{}व\appbrk{}र्ण्णा\appbrk{}व\appbrk{}सं\appbrk{}शयं}}\Afootnote{K H; \skt{असात्म्यभोजनं हन्या\appbrk{}द्ब\appbrk{}ल\appbrk{}व\appbrk{}र्ण\appbrk{}म\appbrk{}सं\appbrk{}श\appbrk{}यम्~।}\ A}}\skt{~। }\pend
\pstart\versepar \skt{}\edtext{\skt{अनात्मवन्तः}}{\lemma{\skt{अनात्मवन्तः}}\Afootnote{A K; \skt{अनात्मवन्तम्}\ H}}\skt{ पशुवद्भुञ्जते ये }\edtext{\skt{प्रमाणतः}}{\lemma{\skt{प्रमाणतः}}\Afootnote{K H; \skt{ऽप्रमाणतः~।}\ A}}\skt{~। }\pend
\pstart\versepar \skt{रोगानीकस्य ते }\edtext{\skt{मूल\-म\-जी\-र्ण्णा\-त्प्रा\-प्नु\-वन्ति}}{\lemma{\skt{अजी\appbrk{}र्ण्णा\appbrk{}त्प्रा\appbrk{}प्नु\appbrk{}वन्ति}}\Afootnote{K H; \skt{अजीर्णं प्राप्नुवन्ति}\ A}}\skt{ }\edtext{\skt{हः}}{\lemma{\skt{हः}}\Afootnote{K; \skt{हि~॥}\ A; \skt{ह~॥}\ H}}\skt{~।}\pend

\pstart\versepar \ledsidenote{\passref{4.39.38}}\edtext{\skt{व्यापदां}}{\lemma{\skt{व्यापदां}}\Afootnote{A K; \skt{व्यापदा}\ H}}\skt{ }\edtext{\skt{कारणं}}{\lemma{\skt{कारणं}}\Afootnote{A K; \skt{त्कारणम्}\ H}}\skt{ वीक्ष्य }\edtext{\skt{व्यापत्स्वेतेषु}}{\lemma{\skt{एतेषु}}\Afootnote{K H; \skt{एतासु}\ A}}\skt{ बुद्धिमान्~। }\pend
\pstart\versepar \skt{प्रय\-ते\-ता\-तु\-रा\-रोग्यै प्रत्यनीकेन हेतुना~।}\edtext{}{\lemma{\skt{बुद्धिमान्~। प्रय\appbrk{}ते\appbrk{}ता\appbrk{}तु\appbrk{}रा\appbrk{}रोग्यै}}\Afootnote{K H; \skt{बुद्धिमान्~। प्रय\appbrk{}ते\appbrk{}ता\appbrk{}तु\appbrk{}रा\appbrk{}रोग्ये}\ A}}\pend

\pstart\versepar \ledsidenote{\passref{4.39.39}}\skt{विरि\-क्त\-वा\-न्तै\-र्ह\-रि\-णै\-ण\-ला\-वकाः }\edtext{\skt{शशाश्च}}{\lemma{\skt{शशाश्च}}\Afootnote{K H; \skt{शशश्च}\ A}}\skt{ }\edtext{\skt{सेव्याः}}{\lemma{\skt{सेव्याः}}\Afootnote{K H; \skt{सेव्यः}\ A}}\skt{ }\edtext{\skt{सम\-यू\-र\-ति\-त्तिराः}}{\lemma{\skt{सम\appbrk{}यू\appbrk{}र\appbrk{}ति\appbrk{}त्तिराः}}\Afootnote{K H; \skt{सम\appbrk{}यू\appbrk{}र\appbrk{}ति\appbrk{}त्तिरिः~।}\ A}}\skt{~। }\pend
\pstart\versepar \skt{}\edtext{\skt{सषष्टिकाश्चापि}}{\lemma{\skt{चापि}}\Afootnote{K H; \skt{चैव}\ A}}\skt{ पुरा\-ण\-शा\-ल\-य\-स्त\-थैव मुद्गा लघु यच्च }\edtext{\skt{कीर्त्तितमिति}}{\lemma{\skt{इति}}\Afootnote{K H; \emph{om.}\ A}}\skt{~॥}\pend

\pstart \ledsidenote{\passref{4.39.trailer}}\edtext{\skt{॥ चि ३९~॥ {\symfont ❈}}}{\lemma{\skt{॥ चि ३९~॥ {\symfont ❈}}}\Afootnote{K; \skt{इति सुश्रु\appbrk{}त\appbrk{}सं\appbrk{}हि\appbrk{}तायां चिकित्सास्थाने आतु\appbrk{}रो\appbrk{}प\appbrk{}द्र\appbrk{}व\appbrk{}चि\appbrk{}कि\appbrk{}त्सितं नामै\appbrk{}को\appbrk{}न\appbrk{}च\appbrk{}त्वा\appbrk{}रिं\appbrk{}शोऽ\appbrk{}ध्यायः~॥३९~॥}\ A; \skt{॥ चिकित्सास्थाने एको\appbrk{}न\appbrk{}च\appbrk{}त्वा\appbrk{}रिं\appbrk{}श\appbrk{}त्तमो ध्यायः~॥}\ H}}\skt{~॥}\pend

\endnumbering
\beginnumbering
\eledsection*{Adhyāya 40}
\pstart \ledsidenote{\passref{4.40.1}}\skt{अथातो }\edtext{\skt{धूमनस्तः}}{\lemma{\skt{धूमनस्तः}}\Afootnote{K; \skt{धूमनस्य}\ A; \skt{धूमनस्त}\ H}}\skt{ }\edtext{\skt{कव\-ल\-चि\-कि\-त्सितं}}{\lemma{\skt{कव\appbrk{}ल\appbrk{}चि\appbrk{}कि\appbrk{}त्सितं}}\Afootnote{K H; \skt{कव\appbrk{}ल\appbrk{}ग्र\appbrk{}ह\appbrk{}चि\appbrk{}कि\appbrk{}त्सितं}\ A}}\skt{ }\edtext{\skt{व्याख्यास्यामः}}{\lemma{\skt{व्याख्यास्यामः}}\Afootnote{A H; \skt{व्या व्याख्यास्यामः~॥}\ K}}\skt{~॥}\pend

\pstart \ledsidenote{\passref{4.40.3}}\edtext{\skt{पञ्चविधो}}{\lemma{\skt{पञ्चविधो}}\Afootnote{K H; \skt{धूमः पञ्चविधो}\ A}}\skt{ }\edtext{\skt{धूमो}}{\lemma{\skt{धूमो}}\Afootnote{K H; \emph{om.}\ A}}\skt{ भवति~। तद्यथा~। प्रायोगिकः स्नैहिको }\edtext{\skt{वैरेचनिकः}}{\lemma{\skt{वैरेचनिकः}}\Afootnote{K H; \skt{वैरेनिकः,}\ A}}\skt{ कासघ्नो }\edtext{\skt{वामनीय}}{\lemma{\skt{वामनीय}}\Afootnote{K H; \skt{वामनीयश्}\ A}}\skt{ }\edtext{\skt{इति}}{\lemma{\skt{इति}}\Afootnote{K H; \skt{चेति~॥}\ A}}\skt{~॥}\pend

\pstart \ledsidenote{\passref{4.40.4}}\skt{तत्रैलादिना कुष्ठ\-त\-ग\-र\-व\-र्ज्येन श्लक्ष्णपिष्टेन द्वादशाङ्गुलं }\edtext{\skt{शरकाण्डं}}{\lemma{\skt{शरकाण्डं}}\Afootnote{K H; \skt{शरकाण्}\ A}}\skt{ }\edtext{\edtext{\edtext{\skt{दशा\-ङ्गु\-ल\-ग\-र्भ\-म\-ङ्गुल}}{\lemma{\skt{दशा\appbrk{}ङ्गु\appbrk{}ल\appbrk{}ग\appbrk{}र्भम्}}\Afootnote{K\textsuperscript{ac} H; \skt{ड}\ A}}\skt{ मङ्गुष्ठोदर}}{\lemma{\skt{दशा\appbrk{}ङ्गु\appbrk{}ल\appbrk{}ग\appbrk{}र्भम्… अङ्}}\Afootnote{\emph{om.}\ K\textsuperscript{pc}}}\skt{ परिणाहं}}{\lemma{\skt{अङ्गुल मङ्गुष्ठोदर परिणाहं}}\Afootnote{\skt{मङ्गु\appbrk{}लि\appbrk{}प\appbrk{}रि\appbrk{}णाहं}\ A; \skt{अङ्गुल परिणाहं}\ K\textsuperscript{ac}; \skt{अंगु\appbrk{}ष्ठो\appbrk{}द\appbrk{}र\appbrk{}प\appbrk{}रि\appbrk{}णाहं}\ H}}\skt{ }\edtext{\skt{क्षौमे\-णा\-ष्टा\-ङ्गु\-ल\-म्प\-रि\-वे\-ष्ट्या\-ले\-प\-ये\-देषा}}{\lemma{\skt{परि\appbrk{}वे\appbrk{}ष्ट्या\appbrk{}ले\appbrk{}प\appbrk{}ये\appbrk{}देषा}}\Afootnote{K H; \skt{वेष्टयित्वा लेपयेदेषा}\ A}}\skt{ वर्त्तिः }\edtext{\skt{प्रायोगिकः}}{\lemma{\skt{प्रायोगिकः}}\Afootnote{K H; \skt{प्रायोगिके,}\ A}}\skt{~॥ }\edtext{\skt{स्नेह\-फ\-ल\-सा\-र\-म\-धू\-च्छि\-ष्ट\-स\-र्ज\-र\-स\-गु\-ग्गु\-लु\-प्र\-भृ\-ति\-भिश्च}}{\lemma{\skt{स्नेह\appbrk{}फ\appbrk{}ल\appbrk{}सा\appbrk{}र\appbrk{}म\appbrk{}धू\appbrk{}च्छि\appbrk{}ष्ट\appbrk{}स\appbrk{}र्ज\appbrk{}र\appbrk{}स\appbrk{}गु\appbrk{}ग्गु\appbrk{}लु\appbrk{}प्र\appbrk{}भृ\appbrk{}ति\appbrk{}भिश्च}}\Afootnote{K H; \skt{स्नेह\appbrk{}फ\appbrk{}ल\appbrk{}सा\appbrk{}र\appbrk{}म\appbrk{}धू\appbrk{}च्छि\appbrk{}ष्ट\appbrk{}स\appbrk{}र्ज\appbrk{}र\appbrk{}स\appbrk{}गु\appbrk{}ग्गु\appbrk{}लु\appbrk{}प्र\appbrk{}भृ\appbrk{}तिभिः}\ A}}\skt{ स्नेहमिश्रैः }\edtext{\skt{स्नैहिकैः}}{\lemma{\skt{स्नैहिकैः}}\Afootnote{K H; \skt{स्नैहिके,}\ A}}\skt{~॥ }\edtext{\skt{शिरो\-वि\-रे\-च\-न\-द्र\-व्यै\-र्वै\-रे\-च\-निकैः}}{\lemma{\skt{वैरेचनिकैः}}\Afootnote{K H; \skt{वैरेचने,}\ A}}\skt{~॥ }\edtext{\skt{बृह\-ती\-क\-ण्ट\-का\-रि\-का\-त्रि\-क\-टु\-क\-का\-स\-म\-र्द\-हिं\-ग्वि\-ङ्गु\-दी\-त्व\-ङ्म\-नः\-शि\-ला\-च्छि\-न्न\-रु\-हा\-क\-र्क\-ट\-क\-सृ\-ङ्गी\-प्र\-भृ\-तिभिः}}{\lemma{\skt{बृह\appbrk{}ती\appbrk{}क\appbrk{}ण्ट\appbrk{}का\appbrk{}रि\appbrk{}का\appbrk{}त्रि\appbrk{}क\appbrk{}टु\appbrk{}क\appbrk{}का\appbrk{}स\appbrk{}म\appbrk{}र्द\appbrk{}हिं\appbrk{}ग्वि\appbrk{}ङ्गु\appbrk{}दी\appbrk{}त्व\appbrk{}ङ्म\appbrk{}नः\appbrk{}शि\appbrk{}ला\appbrk{}च्छि\appbrk{}न्न\appbrk{}रु\appbrk{}हा\appbrk{}क\appbrk{}र्क\appbrk{}ट\appbrk{}क\appbrk{}सृ\appbrk{}ङ्गी\appbrk{}प्र\appbrk{}भृ\appbrk{}तिभिः}}\Afootnote{K; \skt{बृह\appbrk{}ती\appbrk{}क\appbrk{}ण्ट\appbrk{}का\appbrk{}रि\appbrk{}का\appbrk{}त्रि\appbrk{}क\appbrk{}टु\appbrk{}का\appbrk{}स\appbrk{}म\appbrk{}र्द\appbrk{}हि\appbrk{}ङ्ग्वि\appbrk{}ङ्गु\appbrk{}दी\appbrk{}त्व\appbrk{}ङ्म\appbrk{}नः\appbrk{}शि\appbrk{}ला\appbrk{}च्छि\appbrk{}न्न\appbrk{}रु\appbrk{}हा\appbrk{}क\appbrk{}र्क\appbrk{}ट\appbrk{}शृ\appbrk{}ङ्गी\appbrk{}प्र\appbrk{}भृ\appbrk{}तिभिः}\ A; \skt{बृह\appbrk{}ती\appbrk{}क\appbrk{}ण्ट\appbrk{}का\appbrk{}रि\appbrk{}का\appbrk{}त्रि\appbrk{}क\appbrk{}टु\appbrk{}क\appbrk{}का\appbrk{}स\appbrk{}म\appbrk{}र्द्द\appbrk{}हिं\appbrk{}ग्वि\appbrk{}ङ्गु\appbrk{}दी\appbrk{}त्व\appbrk{}ङ्म\appbrk{}नः\appbrk{}शि\appbrk{}ला\appbrk{}छि\appbrk{}न्न\appbrk{}रु\appbrk{}हा\appbrk{}क\appbrk{}र्क्क\appbrk{}ट\appbrk{}क\appbrk{}शृ\appbrk{}ङ्गी\appbrk{}प्र\appbrk{}भृ\appbrk{}तिभिः}\ H}}\skt{ }\edtext{\skt{कासहरैः}}{\lemma{\skt{कासहरैः}}\Afootnote{K; \skt{कासहरैश्च}\ A H}}\skt{ }\edtext{\skt{कासघ्नैः}}{\lemma{\skt{कासघ्नैः}}\Afootnote{K H; \skt{कासघ्ने,}\ A}}\skt{~। }\edtext{\skt{स्नायु\-च\-र्म\-खु\-र\-शृ\-ङ्ग\-क\-र्क\-टा\-स्थि\-शु\-ष्क\-म\-त्स्य\-व\-ल्लू\-र\-क्रि\-मि\-भि\-र्वा\-म\-नी\-यैश्च}}{\lemma{\skt{स्नायु\appbrk{}च\appbrk{}र्म\appbrk{}खु\appbrk{}र\appbrk{}शृ\appbrk{}ङ्ग\appbrk{}क\appbrk{}र्क\appbrk{}टा\appbrk{}स्थि\appbrk{}शु\appbrk{}ष्क\appbrk{}म\appbrk{}त्स्य\appbrk{}व\appbrk{}ल्लू\appbrk{}र\appbrk{}क्रि\appbrk{}मि\appbrk{}भिर्}}\Afootnote{K; \skt{स्नायु\appbrk{}च\appbrk{}र्म\appbrk{}खु\appbrk{}र\appbrk{}शृ\appbrk{}ङ्ग\appbrk{}क\appbrk{}र्क\appbrk{}ट\appbrk{}का\appbrk{}स्थि\appbrk{}शु\appbrk{}ष्क\appbrk{}म\appbrk{}त्स्य\appbrk{}व\appbrk{}ल्लू\appbrk{}र\appbrk{}कृ\appbrk{}मि\appbrk{}प्र\appbrk{}भृ\appbrk{}ति\appbrk{}भिर्}\ A; \skt{स्नायु\appbrk{}च\appbrk{}र्म्म\appbrk{}खु\appbrk{}र\appbrk{}शृ\appbrk{}ङ्ग\appbrk{}क\appbrk{}र्क्क\appbrk{}ट\appbrk{}का\appbrk{}स्थि\appbrk{}शु\appbrk{}ष्क\appbrk{}म\appbrk{}त्स्य\appbrk{}व\appbrk{}ल्लू\appbrk{}र\appbrk{}क्रि\appbrk{}मि\appbrk{}भिर्}\ H}}\skt{ }\edtext{\skt{वामनीय}}{\lemma{\skt{वामनीय}}\Afootnote{K H; \skt{वामनीये~॥}\ A}}\skt{ }\edtext{\skt{इति}}{\lemma{\skt{इति}}\Afootnote{K H; \emph{om.}\ A}}\skt{~॥}\pend

\pstart \ledsidenote{\passref{4.40.5}}\skt{तत्र }\edtext{\edtext{\skt{बस्ति\-ने\-त्र\-व्य\-व्यै\-र्दू\-म\-ने\-त्र\-द्र\-व्यानि}}{\lemma{\skt{बस्ति\appbrk{}ने\appbrk{}त्र\appbrk{}व्य\appbrk{}व्यैर्}}\Afootnote{K; \skt{बस्ति\appbrk{}ने\appbrk{}त्र\appbrk{}द्र\appbrk{}व्यैर्}\ A H}}}{\lemma{\skt{दूम\appbrk{}ने\appbrk{}त्र\appbrk{}द्र\appbrk{}व्यानि}}\Afootnote{K; \skt{धूम\appbrk{}ने\appbrk{}त्र\appbrk{}द्र\appbrk{}व्याणि}\ A H}}\skt{ व्याख्यातानि भवन्ति~॥ धूमनेत्रं तु }\edtext{\edtext{\skt{कनी\-नि\-का\-प\-रि\-णा\-ह\-मग्रे}}{\lemma{\skt{कनी\appbrk{}नि\appbrk{}का\appbrk{}प\appbrk{}रि\appbrk{}णा\appbrk{}हम्}}\Afootnote{K H; \skt{कनि\appbrk{}ष्ठि\appbrk{}का\appbrk{}प\appbrk{}रि\appbrk{}णाह}\ A}}}{\lemma{\skt{अग्रे}}\Afootnote{K H; \skt{मग्रे}\ A}}\skt{ कला\-य\-मा\-त्र\-स्रोतो }\edtext{\skt{मूले ङ्गुष्ठपरिणाहं}}{\lemma{\skt{मूले ङ्गुष्ठपरिणाहं}}\Afootnote{K H; \skt{मूलेऽ\appbrk{}ङ्गु\appbrk{}ष्ठ\appbrk{}प\appbrk{}रि\appbrk{}णाहं}\ A}}\skt{ धूम\-व\-र्त्ति\-प्र\-वे\-श\-स्रोतो }\edtext{\skt{ङ्गुलानि}}{\lemma{\skt{ङ्गुलानि}}\Afootnote{K H; \skt{ऽङ्गुलान्य}\ A}}\skt{ }\edtext{\skt{चतु\-श्च\-त्वा\-रिं\-श\-त्प्रा\-यो\-गिके}}{\lemma{\skt{चतु\appbrk{}श्च\appbrk{}त्वा\appbrk{}रिं\appbrk{}श\appbrk{}त्प्रा\appbrk{}यो\appbrk{}गिके}}\Afootnote{K H; \skt{ष्टच\appbrk{}त्वा\appbrk{}रिं\appbrk{}श\appbrk{}त्प्रा\appbrk{}यो\appbrk{}गिके,}\ A}}\skt{~। }\edtext{\skt{द्वात्रिंशत्स्नैहिके}}{\lemma{\skt{द्वात्रिंशत्स्नैहिके}}\Afootnote{K H; \skt{द्वात्रिंशत्स्नेहने,}\ A}}\skt{~। चतुर्विंशति }\edtext{\skt{वैरेचनिके}}{\lemma{\skt{वैरेचनिके}}\Afootnote{K H; \skt{र्वैरेचने,}\ A}}\skt{~। }\edtext{\skt{षोडशकासघ्ने}}{\lemma{\skt{षोडशकासघ्ने}}\Afootnote{K\textsuperscript{pc} H; \skt{षोडशाङ्गुलं कासघ्ने}\ A; \skt{षोडश सघ्ने~।}\ K\textsuperscript{ac}}}\skt{~। }\edtext{\skt{दश}}{\lemma{\skt{दश}}\Afootnote{K H; \emph{om.}\ A}}\skt{ }\edtext{\skt{वामनीय}}{\lemma{\skt{वामनीय}}\Afootnote{K; \skt{वामनीये च~। एतेऽपि कोला\appbrk{}स्थि\appbrk{}मा\appbrk{}त्र\appbrk{}च्छिद्रे भवतः~। व्रण\appbrk{}ने\appbrk{}त्र\appbrk{}म\appbrk{}ष्टा\appbrk{}ङ्गुलं व्रणधूपनार्थं कला\appbrk{}य\appbrk{}प\appbrk{}रि\appbrk{}म\appbrk{}ण्डलं कुल\appbrk{}त्थ\appbrk{}वा\appbrk{}हि\appbrk{}स्रोत}\ A; \skt{वामनीय इति~।}\ H}}\skt{ }\edtext{\skt{इति}}{\lemma{\skt{इति}}\Afootnote{A K; \emph{om.}\ H}}\skt{~॥}\pend

\pstart \ledsidenote{\passref{4.40.6}}\skt{अथ सुखोपविष्टः सुमना }\edtext{\edtext{\skt{ऋज्व\-धो\-दृ\-ष्टि\-र\-त\-न्द्रितः}}{\lemma{\skt{ऋज्वधोदृष्टिर्}}\Afootnote{K H; \skt{ऋज्वधोदृष्टि}\ A}}}{\lemma{\skt{अतन्द्रितः}}\Afootnote{K H; \skt{रतन्द्रितः}\ A}}\skt{~। }\edtext{\edtext{\skt{स्नेह\-प्र\-दी\-प्ता\-ग्रा\-म्व\-र्त्ति\-ने\-त्र\-स्रो\-तसि}}{\lemma{\skt{स्नेह\appbrk{}प्र\appbrk{}दी\appbrk{}प्ता\appbrk{}ग्राम्}}\Afootnote{K H; \skt{स्थिरे समाहिते शरावे प्रक्षिप्य}\ A}}}{\lemma{\skt{वर्त्ति\appbrk{}ने\appbrk{}त्र\appbrk{}स्रो\appbrk{}तसि}}\Afootnote{K; \skt{वर्तिं मूल\appbrk{}च्छि\appbrk{}द्रे\appbrk{}णा\appbrk{}न्येन शरावेण}\ A; \skt{वर्त्ति\appbrk{}ने\appbrk{}त्र\appbrk{}स्रो\appbrk{}त्रसि}\ H}}\skt{ }\edtext{\skt{प्रणिधाय}}{\lemma{\skt{प्रणिधाय}}\Afootnote{K H; \skt{पिधाय}\ A}}\skt{ }\edtext{\edtext{\skt{धूमम्पिबेदिति}}{\lemma{\skt{धूमम्}}\Afootnote{K H; \skt{तस्मिन्छिद्रे नेत्रमूलं संयोज्य धूममासेवेत, प्रशान्ते धूमे वर्तिमवशिष्टां प्रक्षिप्य पुनरपि धूमं}\ A}}}{\lemma{\skt{पिबेदिति}}\Afootnote{K H; \skt{पाय\appbrk{}ये\appbrk{}दा\appbrk{}दो\appbrk{}ष\appbrk{}वि\appbrk{}शुद्धेः; एष धूम\appbrk{}पा\appbrk{}नो\appbrk{}पा\appbrk{}य\appbrk{}विधिः~॥}\ A}}\skt{~॥}\pend

\pstart\versepar \ledsidenote{\passref{4.40.7}}\skt{मुखेन }\edtext{\skt{तम्पिबेत्पूर्वं}}{\lemma{\skt{पिबेत्पूर्वं}}\Afootnote{A K; \skt{पिबेम्पूर्व्वन्}\ H}}\skt{ नासिकाभ्यान्ततः पिबेत्~। }\pend
\pstart\versepar \skt{मुखपीतं मुखेनैव वमेत्पीतञ्च }\edtext{\skt{नस्ततः}}{\lemma{\skt{नस्ततः}}\Afootnote{\skt{नासया॥}\ A; \skt{नस्ततः}\ K H}}\skt{~।}\pend

\pstart\versepar \ledsidenote{\passref{4.40.8}}\edtext{\skt{यो वमेन्नस्ततो धूमन्नस्तपीतं मुखेन वा}}{\lemma{\skt{… मु}}\Afootnote{\emph{om.}\ A}}\skt{~। }\pend
\pstart\versepar \skt{स }\edtext{\skt{नेत्र\-क\-र्ण्ण\-ना\-सा\-स्य\-सं\-श्र\-या\-ल्ल\-भते गदान्~॥}}{\lemma{\skt{वा… गदान्}}\Afootnote{\emph{om.}\ A}}\edtext{}{\lemma{\skt{मुखेन वा~। स}}\Afootnote{K H; \skt{मुखेन}\ A}}\pend

\pstart \ledsidenote{\passref{4.40.9}}\skt{विशेषतस्तु प्रायोगिकं }\edtext{\skt{घ्राणेनादद्यात्~। स्नैहिकं}}{\lemma{\skt{घ्राणेनादद्यात्~। स्नैहिकं}}\Afootnote{K H; \skt{घ्राणेनाददीत, स्नैहिकं}\ A}}\skt{ मुखनासाभ्यां~। }\edtext{\skt{नासिकाभ्यां}}{\lemma{\skt{नासिकाभ्यां}}\Afootnote{K H; \skt{नासिकया}\ A}}\skt{ वैरेचनिकं~। }\edtext{\skt{मुखे\-नै\-वे\-त\-रा\-विति}}{\lemma{\skt{इति}}\Afootnote{K H; \emph{om.}\ A}}\skt{~॥}\pend

\pstart \ledsidenote{\passref{4.40.11}}\edtext{\skt{तत्र श्रम\-भ\-या\-म\-र्षो\-ष्ण\-वि\-ष\-र\-क्त\-पि\-त्त\-म\-द\-मू\-र्च्छा\-पि\-पा\-सा\-ता\-लु\-शो\-ष\-मू\-र्धा\-भि\-घा\-तो\-द्गा\-र\-ति\-मि\-र\-मे\-हो\-द\-रा\-ध्मा\-नो\-र्ध्व\-वा\-तो\-न्त्र\-वृ\-द्धि\-बा\-ल\-वृ\-द्ध\-दु\-र्ब\-ल\-वि\-र\-क्ता\-स्था\-पि\-त\-जा\-ग\-रि\-त\-ग\-र्भि\-णी\-रू\-क्ष\-क्ष\-त\-क्षी\-णो\-रस्काः}}{\lemma{\skt{तत्र श्रम\appbrk{}भ\appbrk{}या\appbrk{}म\appbrk{}र्षो\appbrk{}ष्ण\appbrk{}वि\appbrk{}ष\appbrk{}र\appbrk{}क्त\appbrk{}पि\appbrk{}त्त\appbrk{}म\appbrk{}द\appbrk{}मू\appbrk{}र्च्छा\appbrk{}पि\appbrk{}पा\appbrk{}सा\appbrk{}ता\appbrk{}लु\appbrk{}शो\appbrk{}ष\appbrk{}मू\appbrk{}र्धा\appbrk{}भि\appbrk{}घा\appbrk{}तो\appbrk{}द्गा\appbrk{}र\appbrk{}ति\appbrk{}मि\appbrk{}र\appbrk{}मे\appbrk{}हो\appbrk{}द\appbrk{}रा\appbrk{}ध्मा\appbrk{}नो\appbrk{}र्ध्व\appbrk{}वा\appbrk{}तो\appbrk{}न्त्र\appbrk{}वृ\appbrk{}द्धि\appbrk{}बा\appbrk{}ल\appbrk{}वृ\appbrk{}द्ध\appbrk{}दु\appbrk{}र्ब\appbrk{}ल\appbrk{}वि\appbrk{}र\appbrk{}क्ता\appbrk{}स्था\appbrk{}पि\appbrk{}त\appbrk{}जा\appbrk{}ग\appbrk{}रि\appbrk{}त\appbrk{}ग\appbrk{}र्भि\appbrk{}णी\appbrk{}रू\appbrk{}क्ष\appbrk{}क्ष\appbrk{}त\appbrk{}क्षी\appbrk{}णो\appbrk{}रस्काः}}\Afootnote{K\textsuperscript{pc}; \skt{तत्र शोक\appbrk{}श्र\appbrk{}म\appbrk{}भ\appbrk{}या\appbrk{}म\appbrk{}र्षौ\appbrk{}ष्ण्य\appbrk{}वि\appbrk{}ष\appbrk{}र\appbrk{}क्त\appbrk{}पि\appbrk{}त्त\appbrk{}म\appbrk{}द\appbrk{}मू\appbrk{}र्च्छा\appbrk{}दा\appbrk{}ह\appbrk{}पि\appbrk{}पा\appbrk{}सा\appbrk{}पा\appbrk{}ण्डु\appbrk{}रो\appbrk{}ग\appbrk{}ता\appbrk{}लु\appbrk{}शो\appbrk{}ष\appbrk{}च्छ\appbrk{}र्दि\appbrk{}शि\appbrk{}रोऽ\appbrk{}भि\appbrk{}घा\appbrk{}तो\appbrk{}द्गा\appbrk{}रा\appbrk{}प\appbrk{}त\appbrk{}र्पि\appbrk{}त\appbrk{}ति\appbrk{}मि\appbrk{}र\appbrk{}प्र\appbrk{}मे\appbrk{}हो\appbrk{}द\appbrk{}रा\appbrk{}ध्मा\appbrk{}नो\appbrk{}र्ध्व\appbrk{}वा\appbrk{}तार्ता बाल\appbrk{}वृ\appbrk{}द्ध\appbrk{}दु\appbrk{}र्ब\appbrk{}ल\appbrk{}वि\appbrk{}रि\appbrk{}क्ता\appbrk{}स्था\appbrk{}पि\appbrk{}त\appbrk{}जा\appbrk{}ग\appbrk{}रि\appbrk{}त\appbrk{}ग\appbrk{}र्भि\appbrk{}णी\appbrk{}रू\appbrk{}क्ष\appbrk{}क्षी\appbrk{}ण\appbrk{}क्ष\appbrk{}तो\appbrk{}रस्क}\ A; \skt{तत्र श्रम\appbrk{}भ\appbrk{}या\appbrk{}म\appbrk{}र्षो\appbrk{}ष्ण\appbrk{}वि\appbrk{}ष\appbrk{}र\appbrk{}क्त\appbrk{}पि\appbrk{}त्त\appbrk{}म\appbrk{}द\appbrk{}मू\appbrk{}र्च्छा\appbrk{}पि\appbrk{}पा\appbrk{}सा\appbrk{}ता\appbrk{}लु\appbrk{}शो\appbrk{}ष\appbrk{}मू\appbrk{}र्धा\appbrk{}भि\appbrk{}घा\appbrk{}तो\appbrk{}द्गार रमे\appbrk{}हो\appbrk{}द\appbrk{}रा\appbrk{}ध्मा\appbrk{}नो\appbrk{}र्ध्व\appbrk{}वा\appbrk{}तो\appbrk{}न्त्र\appbrk{}वृ\appbrk{}द्धि\appbrk{}बा\appbrk{}ल\appbrk{}वृ\appbrk{}द्ध\appbrk{}दु\appbrk{}र्ब\appbrk{}ल\appbrk{}वि\appbrk{}र\appbrk{}क्ता\appbrk{}स्था\appbrk{}पि\appbrk{}त\appbrk{}जा\appbrk{}ग\appbrk{}रि\appbrk{}त\appbrk{}ग\appbrk{}र्भि\appbrk{}णी\appbrk{}रू\appbrk{}क्ष\appbrk{}क्ष\appbrk{}त\appbrk{}क्षी\appbrk{}णो\appbrk{}रस्काः}\ K\textsuperscript{ac}; \skt{तत्र श्रम\appbrk{}भ\appbrk{}या\appbrk{}म\appbrk{}र्षो\appbrk{}ष्ण\appbrk{}वि\appbrk{}ष\appbrk{}र\appbrk{}क्त\appbrk{}पि\appbrk{}त्त\appbrk{}म\appbrk{}द\appbrk{}मू\appbrk{}र्च्छा\appbrk{}पि\appbrk{}पा\appbrk{}सा\appbrk{}ता\appbrk{}लु\appbrk{}शो\appbrk{}ष\appbrk{}मू\appbrk{}र्द्धा\appbrk{}भि\appbrk{}घा\appbrk{}तो\appbrk{}द्गा\appbrk{}र\appbrk{}ति\appbrk{}मि\appbrk{}र\appbrk{}मे\appbrk{}हो\appbrk{}द\appbrk{}रा\appbrk{}ध्मा\appbrk{}नो\appbrk{}र्द्ध\appbrk{}वा\appbrk{}तो\appbrk{}न्त्र\appbrk{}वृ\appbrk{}द्धि\appbrk{}बा\appbrk{}ल\appbrk{}वृ\appbrk{}द्ध\appbrk{}दु\appbrk{}र्ब्ब\appbrk{}ल\appbrk{}वि\appbrk{}र\appbrk{}क्ता\appbrk{}स्था\appbrk{}पि\appbrk{}त\appbrk{}जा\appbrk{}ग\appbrk{}रि\appbrk{}त\appbrk{}ग\appbrk{}र्ब्भि\appbrk{}णी\appbrk{}रू\appbrk{}क्ष\appbrk{}क्ष\appbrk{}त\appbrk{}क्षी\appbrk{}णो\appbrk{}रस्काः}\ H\textsuperscript{ac} H\textsuperscript{pc}}}\skt{ }\edtext{\skt{मधु\-घृ\-त\-द\-धि\-दु\-ग्ध\-म\-द्य\-य\-वा\-गू\-पीताः}}{\lemma{\skt{मधु\appbrk{}घृ\appbrk{}त\appbrk{}द\appbrk{}धि\appbrk{}दु\appbrk{}ग्ध\appbrk{}म\appbrk{}द्य\appbrk{}य\appbrk{}वा\appbrk{}गू\appbrk{}पीताः}}\Afootnote{K; \skt{मधु\appbrk{}घृ\appbrk{}त\appbrk{}द\appbrk{}धि\appbrk{}दु\appbrk{}ग्ध\appbrk{}म\appbrk{}त्स्य\appbrk{}म\appbrk{}द्य\appbrk{}य\appbrk{}वा\appbrk{}गू\appbrk{}पीता}\ A; \skt{मधु\appbrk{}घृ\appbrk{}त\appbrk{}द\appbrk{}धि\appbrk{}दु\appbrk{}ग्ध\appbrk{}म\appbrk{}द्य\appbrk{}य\appbrk{}वा\appbrk{}गू\appbrk{}पीताः}\ H}}\skt{ }\edtext{\skt{स्वल्पकफाश्च}}{\lemma{\skt{स्वल्पकफाश्च}}\Afootnote{K H; \skt{ल्पकफाश्च न}\ A}}\skt{ }\edtext{\skt{धूमन्न}}{\lemma{\skt{धूमन्न}}\Afootnote{K H; \skt{धूममा}\ A}}\skt{ }\edtext{\skt{सेवेरन्निति}}{\lemma{\skt{सेवेरन्निति}}\Afootnote{K H; \skt{सेवेरन्~॥}\ A}}\skt{~॥}\pend

\pstart\versepar \ledsidenote{\passref{4.40.12}}\skt{अकालपीतो }\edtext{\skt{भ्रमरी}}{\lemma{\skt{भ्रमरी}}\Afootnote{K H; \skt{कुरुते भ्रमं मूर्च्छां}\ A}}\skt{ }\edtext{\skt{शिरोरोगं}}{\lemma{\skt{शिरोरोगं}}\Afootnote{K H; \skt{शिरोरुजम्~।}\ A}}\skt{ }\edtext{\skt{करोति च}}{\lemma{\skt{करोति च}}\Afootnote{K; \emph{om.}\ A; \skt{करोति हि~।}\ H}}\skt{~। }\pend
\pstart\versepar \skt{}\edtext{\edtext{\skt{नासा\-स्रो\-तो\-क्षि\-जि\-ह्वा\-नां\-मु\-प\-घा\-तञ्च}}{\lemma{\skt{नासा\appbrk{}स्रो\appbrk{}तो\appbrk{}क्षि\appbrk{}जि\appbrk{}ह्वा\appbrk{}नांम्}}\Afootnote{K; \skt{घ्राण\appbrk{}श्रो\appbrk{}त्रा\appbrk{}क्षि\appbrk{}जि\appbrk{}ह्वा\appbrk{}नाम्}\ A; \skt{नासा\appbrk{}स्रो\appbrk{}त्रोऽ\appbrk{}क्षि\appbrk{}जि\appbrk{}ह्वा\appbrk{}नाम्}\ H}}}{\lemma{\skt{उपघातञ्च}}\Afootnote{A K\textsuperscript{pc} H; \skt{उपा घातञ्च}\ K\textsuperscript{ac}}}\skt{ दारुणं~॥}\pend

\pstart \ledsidenote{\passref{4.40.13}}\edtext{\edtext{\skt{धूमः}}{\lemma{\skt{धूमः}}\Afootnote{K H; \skt{आद्यास्तु त्रयो धू}\ A}}\skt{ स्नैहिको वैरेचनिकः }\edtext{\skt{प्रायोगिको}}{\lemma{\skt{प्रायोगिको}}\Afootnote{K H; \skt{मा}\ A}}}{\lemma{\skt{मः… प्रा}}\Afootnote{\emph{om.}\ A}}\skt{ द्वादशसु }\edtext{\skt{कालेषूपयोगो}}{\lemma{\skt{कालेषूपयोगो}}\Afootnote{K H; \skt{कालेषूपादेयाः~।}\ A}}\skt{ }\edtext{\skt{भवति}}{\lemma{\skt{भवति}}\Afootnote{K H; \emph{om.}\ A}}\skt{~॥ तद्यथा }\edtext{\edtext{\skt{क्षुत\-द\-न्त\-प्र\-क्षा\-ल\-न\-न\-स्य\-स्ना\-न\-भो\-ज\-ना\-दि\-वा\-स्व\-प्न\-मै\-थु\-न\-च्छ\-र्द\-न\-मू\-त्रो\-च्चा\-रा\-भि\-रु\-षि\-त\-श\-स्त्र\-क\-र्मा\-न्ते\-ष्विति}}{\lemma{\skt{क्षुत\appbrk{}द\appbrk{}न्त\appbrk{}प्र\appbrk{}क्षा\appbrk{}ल\appbrk{}न\appbrk{}न\appbrk{}स्य\appbrk{}स्ना\appbrk{}न\appbrk{}भो\appbrk{}ज\appbrk{}ना\appbrk{}दि\appbrk{}वा\appbrk{}स्व\appbrk{}प्न\appbrk{}मै\appbrk{}थु\appbrk{}न\appbrk{}च्छ\appbrk{}र्द\appbrk{}न\appbrk{}मू\appbrk{}त्रो\appbrk{}च्चा\appbrk{}रा\appbrk{}भिर्}}\Afootnote{\skt{क्षुत\appbrk{}द\appbrk{}न्त\appbrk{}प्र\appbrk{}क्षा\appbrk{}ल\appbrk{}न\appbrk{}न\appbrk{}स्य\appbrk{}स्ना\appbrk{}न\appbrk{}भो\appbrk{}ज\appbrk{}न\appbrk{}दि\appbrk{}वा\appbrk{}स्व\appbrk{}प्न\appbrk{}मै\appbrk{}थु\appbrk{}न\appbrk{}च्छ\appbrk{}र्दि\appbrk{}मू\appbrk{}त्रो\appbrk{}च्चा\appbrk{}रह}\ A; \skt{क्षुत\appbrk{}द\appbrk{}न्त\appbrk{}प्र\appbrk{}क्षा\appbrk{}ल\appbrk{}न\appbrk{}न\appbrk{}स्य\appbrk{}स्ना\appbrk{}न\appbrk{}भो\appbrk{}ज\appbrk{}ना\appbrk{}दि\appbrk{}वा\appbrk{}स्व\appbrk{}प्न\appbrk{}मै\appbrk{}थु\appbrk{}न\appbrk{}च्छ\appbrk{}र्द\appbrk{}न\appbrk{}मू\appbrk{}त्रो\appbrk{}च्चारा}\ K; \skt{क्षुत\appbrk{}द\appbrk{}न्त\appbrk{}प्र\appbrk{}क्षा\appbrk{}ल\appbrk{}न\appbrk{}न\appbrk{}स्य\appbrk{}स्ना\appbrk{}न\appbrk{}भो\appbrk{}ज\appbrk{}न\appbrk{}दि\appbrk{}वा\appbrk{}स्व\appbrk{}प्न\appbrk{}मै\appbrk{}थु\appbrk{}न\appbrk{}च्छ\appbrk{}र्द्द\appbrk{}न\appbrk{}मू\appbrk{}त्रो\appbrk{}च्चा\appbrk{}रा\appbrk{}भिर्}\ H}}}{\lemma{\skt{उषि\appbrk{}त\appbrk{}श\appbrk{}स्त्र\appbrk{}क\appbrk{}र्मा\appbrk{}न्तेष्व्}}\Afootnote{H; \skt{सित\appbrk{}रु\appbrk{}षि\appbrk{}त\appbrk{}श\appbrk{}स्त्र\appbrk{}क\appbrk{}र्मा\appbrk{}न्तेष्व्}\ A; \skt{भिरु\appbrk{}षि\appbrk{}त\appbrk{}श\appbrk{}स्त्र\appbrk{}क\appbrk{}र्मा\appbrk{}न्तेष्व्}\ K}}\skt{~॥ }\edtext{\skt{तस्य}}{\lemma{\skt{तस्य}}\Afootnote{K H; \skt{तत्र}\ A}}\skt{ विभागो }\edtext{\skt{मूत्रो\-च्चा\-र\-क्ष\-व\-थु\-रु\-षि\-त\-मै\-थु\-ना\-न्तेषु}}{\lemma{\skt{मूत्रो\appbrk{}च्चा\appbrk{}र\appbrk{}क्ष\appbrk{}व\appbrk{}थु\appbrk{}रु\appbrk{}षि\appbrk{}त\appbrk{}मै\appbrk{}थु\appbrk{}ना\appbrk{}न्तेषु}}\Afootnote{K H; \skt{मूत्रो\appbrk{}च्चा\appbrk{}र\appbrk{}क्ष\appbrk{}व\appbrk{}थु\appbrk{}ह\appbrk{}सि\appbrk{}त\appbrk{}रु\appbrk{}षि\appbrk{}त\appbrk{}मै\appbrk{}थु\appbrk{}ना\appbrk{}न्तेषु}\ A}}\skt{ स्नैहिकः~। }\edtext{\edtext{\skt{दिवा\-स्व\-प्न\-च्छ\-र्द\-ना\-न्त\-यो\-र्वै\-रे\-च\-निकः}}{\lemma{\skt{दिवा\appbrk{}स्व\appbrk{}प्न\appbrk{}च्छ\appbrk{}र्द\appbrk{}ना\appbrk{}न्त\appbrk{}योर्}}\Afootnote{K H; \skt{स्नान\appbrk{}च्छ\appbrk{}र्दि\appbrk{}वा\appbrk{}स्व\appbrk{}प्ना\appbrk{}न्तेषु}\ A}}}{\lemma{\skt{वैरेचनिकः}}\Afootnote{K H; \skt{वैरेचनिकः,}\ A}}\skt{~। दन्त\-प्र\-क्षा\-ल\-नस्य स्नान\-भो\-ज\-न\-श\-स्त्र\-क\-र्मा\-न्तेषु प्रायोगिक इति~॥}\pend

\pstart \ledsidenote{\passref{4.40.14}}\skt{तत्र स्नैहिको वातं शमयति~। }\edtext{\skt{स्नेहा\-दु\-प\-ले\-पि\-त्वाच्च}}{\lemma{\skt{स्नेहा\appbrk{}दु\appbrk{}प\appbrk{}ले\appbrk{}पि\appbrk{}त्वाच्च}}\Afootnote{K; \skt{स्नेहा\appbrk{}दु\appbrk{}प\appbrk{}ले\appbrk{}पाच्च;}\ A; \skt{स्नेहा\appbrk{}दु\appbrk{}प\appbrk{}ले\appbrk{}प\appbrk{}यि\appbrk{}त्वाच्च~।}\ H}}\skt{~। }\edtext{\skt{वैरेचनिकः}}{\lemma{\skt{वैरेचनिकः}}\Afootnote{K H; \skt{वैरेचनः}\ A}}\skt{ }\edtext{\edtext{\skt{श्लेष्मा\-ण\-मु\-त्क्लि\-श्या\-प\-क\-र्षति}}{\lemma{\skt{श्लेष्माणम्}}\Afootnote{K H; \skt{श्लेष्माण}\ A}}}{\lemma{\skt{उत्क्लि\appbrk{}श्या\appbrk{}प\appbrk{}क\appbrk{}र्षति}}\Afootnote{K H; \skt{मुत्क्ले\appbrk{}श्या\appbrk{}प\appbrk{}क\appbrk{}र्षति,}\ A}}\skt{ }\edtext{\skt{रौक्ष्या\-त्तै\-क्ष्ण्या\-दौ\-ष्ण्या\-द्वै\-ष\-द्याच्च}}{\lemma{\skt{रौक्ष्या\appbrk{}त्तै\appbrk{}क्ष्ण्या\appbrk{}दौ\appbrk{}ष्ण्या\appbrk{}द्वै\appbrk{}ष\appbrk{}द्याच्च}}\Afootnote{K H; \skt{रौक्ष्या\appbrk{}त्तै\appbrk{}क्ष्ण्या\appbrk{}दौ\appbrk{}ष्ण्या\appbrk{}द्वै\appbrk{}श\appbrk{}द्याच्च;}\ A}}\skt{~। प्रायोगिकः }\edtext{\edtext{\edtext{\skt{श्लेष्मा\-ण\-मु\-त्क्ले\-श\-य\-त्यु\-त्क्लिष्टं}}{\lemma{\skt{श्लेष्माणम्}}\Afootnote{K H; \skt{श्लेष्माण}\ A}}}{\lemma{\skt{उत्क्लेशयत्य्}}\Afootnote{K H; \skt{मुत्क्लेशय}\ A}}}{\lemma{\skt{उत्क्लिष्टं}}\Afootnote{K H; \skt{त्युत्क्लिष्टं}\ A}}\skt{ }\edtext{\skt{चापकर्षयति}}{\lemma{\skt{चापकर्षयति}}\Afootnote{K; \skt{चापकर्षति}\ A H}}\skt{ }\edtext{\skt{शमयति वा}}{\lemma{\skt{शमयति वा}}\Afootnote{K H; \skt{शमयति वातं}\ A}}\skt{ साधा\-र\-ण\-त्वा\-त्पू\-र्वा\-भ्या\-मिति~॥}\pend

\pstart\versepar \ledsidenote{\passref{4.40.15}}\edtext{\skt{धूम\-प्र\-यो\-गा\-त्पु\-रुषः}}{\lemma{\skt{धूम\appbrk{}प्र\appbrk{}यो\appbrk{}गा\appbrk{}त्पु\appbrk{}रुषः}}\Afootnote{K H; \skt{भवति चात्र नरो धूमोपयोगाच्च}\ A}}\skt{ प्रस\-न्ने\-न्द्रि\-य\-वा\-ङ्मना~। }\pend
\pstart\versepar \skt{}\edtext{\skt{दृढकेशद्विजः}}{\lemma{\skt{दृढकेशद्विजः}}\Afootnote{K; \skt{दृढकेशद्विज}\ A H}}\skt{ }\edtext{\skt{श्मश्रु\-सु\-ग\-न्धि\-वि\-श\-दा\-ननः}}{\lemma{\skt{श्मश्रु\appbrk{}सु\appbrk{}ग\appbrk{}न्धि\appbrk{}वि\appbrk{}श\appbrk{}दा\appbrk{}ननः}}\Afootnote{K H; \skt{श्मश्रुः सुग\appbrk{}न्धि\appbrk{}वि\appbrk{}श\appbrk{}दा\appbrk{}ननः~॥}\ A}}\skt{~।}\pend

\pstart \ledsidenote{\passref{4.40.16}}\skt{तथा }\edtext{\skt{कास\-श्वा\-सा\-रो\-च\-का\-स्य\-प्र\-ले\-प\-स्व\-र\-भे\-द\-मु\-खा\-स्रा\-व\-क्ष\-व\-थु\-व\-म\-थु\-क्रा\-थ\-त\-न्द्री\-ह\-नु\-म\-न्या\-स्त\-म्भ\-पी\-न\-स\-शि\-रो\-रो\-ग\-लो\-च\-न\-शू\-ला\-वा\-त\-क\-फ\-नि\-मि\-त्ता\-श्चास्य}}{\lemma{\skt{कास\appbrk{}श्वा\appbrk{}सा\appbrk{}रो\appbrk{}च\appbrk{}का\appbrk{}स्य\appbrk{}प्र\appbrk{}ले\appbrk{}प\appbrk{}स्व\appbrk{}र\appbrk{}भे\appbrk{}द\appbrk{}मु\appbrk{}खा\appbrk{}स्रा\appbrk{}व\appbrk{}क्ष\appbrk{}व\appbrk{}थु\appbrk{}व\appbrk{}म\appbrk{}थु\appbrk{}क्रा\appbrk{}थ\appbrk{}त\appbrk{}न्द्री\appbrk{}ह\appbrk{}नु\appbrk{}म\appbrk{}न्या\appbrk{}स्त\appbrk{}म्भ\appbrk{}पी\appbrk{}न\appbrk{}स\appbrk{}शि\appbrk{}रो\appbrk{}रो\appbrk{}ग\appbrk{}लो\appbrk{}च\appbrk{}न\appbrk{}शू\appbrk{}ला\appbrk{}वा\appbrk{}त\appbrk{}क\appbrk{}फ\appbrk{}नि\appbrk{}मि\appbrk{}त्ताश्}}\Afootnote{\skt{कास\appbrk{}श्वा\appbrk{}सा\appbrk{}रो\appbrk{}च\appbrk{}का\appbrk{}स्यो\appbrk{}प\appbrk{}ले\appbrk{}प\appbrk{}स्व\appbrk{}र\appbrk{}भे\appbrk{}द\appbrk{}मु\appbrk{}खा\appbrk{}स्रा\appbrk{}व\appbrk{}क्ष\appbrk{}व\appbrk{}थु\appbrk{}व\appbrk{}म\appbrk{}थु\appbrk{}क्र\appbrk{}थ\appbrk{}त\appbrk{}न्द्रा\appbrk{}नि\appbrk{}द्रा\appbrk{}ह\appbrk{}नु\appbrk{}म\appbrk{}न्या\appbrk{}स्तम्भाः पीन\appbrk{}स\appbrk{}शि\appbrk{}रो\appbrk{}रो\appbrk{}ग\appbrk{}क\appbrk{}र्णा\appbrk{}क्षि\appbrk{}शूला वात\appbrk{}क\appbrk{}फ\appbrk{}नि\appbrk{}मि\appbrk{}त्ताश्}\ A; \skt{कास\appbrk{}श्वा\appbrk{}सा\appbrk{}रो\appbrk{}च\appbrk{}का\appbrk{}स्य\appbrk{}प्र\appbrk{}ले\appbrk{}प\appbrk{}स्व\appbrk{}र\appbrk{}भे\appbrk{}द\appbrk{}मु\appbrk{}खा\appbrk{}स्रा\appbrk{}व\appbrk{}क्ष\appbrk{}व\appbrk{}थु\appbrk{}व\appbrk{}म\appbrk{}थु\appbrk{}क्रा\appbrk{}थ\appbrk{}त\appbrk{}न्द्री\appbrk{}ह\appbrk{}नुमा न्यास्त\appbrk{}म्भ\appbrk{}पी\appbrk{}न\appbrk{}स\appbrk{}शि\appbrk{}रो\appbrk{}रो\appbrk{}ग\appbrk{}लो\appbrk{}च\appbrk{}न\appbrk{}शू\appbrk{}ला\appbrk{}वा\appbrk{}त\appbrk{}क\appbrk{}फ\appbrk{}नि\appbrk{}मि\appbrk{}त्ताश्}\ K\textsuperscript{ac}; \skt{कास\appbrk{}श्वा\appbrk{}सा\appbrk{}रो\appbrk{}च\appbrk{}का\appbrk{}स्य\appbrk{}प्र\appbrk{}ले\appbrk{}प\appbrk{}स्व\appbrk{}र\appbrk{}भे\appbrk{}द\appbrk{}मु\appbrk{}खा\appbrk{}स्रा\appbrk{}व\appbrk{}क्ष\appbrk{}व\appbrk{}थु\appbrk{}व\appbrk{}म\appbrk{}थु\appbrk{}क्रा\appbrk{}थ\appbrk{}त\appbrk{}न्द्री\appbrk{}ह\appbrk{}नु\appbrk{}म\appbrk{}न्या\appbrk{}स्त\appbrk{}म्भ\appbrk{}पी\appbrk{}न\appbrk{}स\appbrk{}शि\appbrk{}रो\appbrk{}रो\appbrk{}ग\appbrk{}लो\appbrk{}च\appbrk{}न\appbrk{}शू\appbrk{}ला\appbrk{}वा\appbrk{}त\appbrk{}क\appbrk{}फ\appbrk{}नि\appbrk{}मि\appbrk{}त्ताश्}\ K\textsuperscript{pc}; \skt{कास\appbrk{}श्वा\appbrk{}सा\appbrk{}रो\appbrk{}च\appbrk{}का\appbrk{}स्य\appbrk{}प्र\appbrk{}ले\appbrk{}प\appbrk{}स्व\appbrk{}र\appbrk{}भे\appbrk{}द\appbrk{}मु\appbrk{}खा\appbrk{}स्रा\appbrk{}व\appbrk{}क्ष\appbrk{}व\appbrk{}थु\appbrk{}व\appbrk{}म\appbrk{}थु\appbrk{}क्र\appbrk{}थ\appbrk{}त\appbrk{}न्द्री\appbrk{}ह\appbrk{}नु\appbrk{}म\appbrk{}न्या\appbrk{}स्त\appbrk{}म्भ\appbrk{}पी\appbrk{}न\appbrk{}स\appbrk{}शि\appbrk{}रो\appbrk{}रो\appbrk{}ग\appbrk{}लो\appbrk{}च\appbrk{}न\appbrk{}शू\appbrk{}ला\appbrk{}वा\appbrk{}त\appbrk{}क\appbrk{}फ\appbrk{}नि\appbrk{}मि\appbrk{}त्ताश्}\ H}}\skt{ मुखरोगा न }\edtext{\skt{भवन्ति}}{\lemma{\skt{भवन्ति}}\Afootnote{A K\textsuperscript{pc} H\textsuperscript{pc}; \skt{पहन्ति~।}\ K\textsuperscript{ac} H\textsuperscript{ac}}}\skt{~।}\pend

\pstart \ledsidenote{\passref{4.40.17}}\skt{तस्य योगा\-यो\-गा\-ति\-योगा }\edtext{\skt{विज्ञातव्या}}{\lemma{\skt{विज्ञातव्या}}\Afootnote{K H; \skt{विज्ञातव्याः~।}\ A}}\skt{ }\edtext{\skt{भवन्ति}}{\lemma{\skt{भवन्ति}}\Afootnote{K H; \skt{तत्र}\ A}}\skt{~। }\edtext{\skt{योगस्तु}}{\lemma{\skt{योगस्तु}}\Afootnote{K H; \skt{योगो रो}\ A}}\skt{ }\edtext{\skt{दृष्टफलार्थः}}{\lemma{\skt{दृष्टफलार्थः}}\Afootnote{K H; \skt{गप्रशमनः,}\ A}}\skt{~। अयोगो }\edtext{\skt{रोगानिर्जयः}}{\lemma{\skt{रोगानिर्जयः}}\Afootnote{K; \skt{रोगा}\ A; \skt{योगानिर्ज्जयः~।}\ H}}\skt{~। }\edtext{\edtext{\edtext{\skt{अति\-यो\-ग\-स्ता\-लु\-शो\-ष\-प\-रि\-दा\-ह\-पि\-पा\-सा\-भ्र\-म\-म\-द\-मू\-र्च्छा\-क\-र्ण्ण\-क्ष्वे\-ल\-दृ\-ष्टि\-रो\-गा\-दौ\-र्ब\-ल्य\-यु\-त\-मिति}}{\lemma{\skt{अतियोगस्}}\Afootnote{K H; \skt{प्रशमनः,}\ A}}}{\lemma{\skt{तालु\appbrk{}शो\appbrk{}ष\appbrk{}प\appbrk{}रि\appbrk{}दा\appbrk{}ह\appbrk{}पि\appbrk{}पा\appbrk{}सा\appbrk{}भ्र\appbrk{}म\appbrk{}म\appbrk{}द\appbrk{}मू\appbrk{}र्च्छा\appbrk{}क\appbrk{}र्ण्ण\appbrk{}क्ष्वे\appbrk{}ल\appbrk{}दृ\appbrk{}ष्टि\appbrk{}रो\appbrk{}गा\appbrk{}दौ\appbrk{}र्ब\appbrk{}ल्य\appbrk{}यु\appbrk{}तम्}}\Afootnote{K H\textsuperscript{pc}; \skt{तालु\appbrk{}ग\appbrk{}ल\appbrk{}शो\appbrk{}ष\appbrk{}प\appbrk{}रि\appbrk{}दा\appbrk{}ह\appbrk{}पि\appbrk{}पा\appbrk{}सा\appbrk{}मू\appbrk{}र्च्छा\appbrk{}भ्र\appbrk{}म\appbrk{}म\appbrk{}द\appbrk{}क\appbrk{}र्ण\appbrk{}क्ष्वे\appbrk{}ड\appbrk{}दृ\appbrk{}ष्टि\appbrk{}ना\appbrk{}सा\appbrk{}रो\appbrk{}ग\appbrk{}दौ\appbrk{}र्ब\appbrk{}ल्यान्य}\ A; \skt{तालु\appbrk{}शो\appbrk{}ष\appbrk{}प\appbrk{}रि\appbrk{}दा\appbrk{}ह\appbrk{}पि\appbrk{}पा\appbrk{}सा\appbrk{}भ्र\appbrk{}म\appbrk{}म\appbrk{}द\appbrk{}मू\appbrk{}र्च्छा\appbrk{}क\appbrk{}र्ण्ण\appbrk{}क्ष्वे\appbrk{}ल\appbrk{}दृ\appbrk{}ष्टि\appbrk{}रो\appbrk{}गा\appbrk{}दौ\appbrk{}र्ब्ब\appbrk{}ल्य\appbrk{}युत}\ H\textsuperscript{ac}}}}{\lemma{\skt{इति}}\Afootnote{K H; \skt{तियोगो जनयति~॥}\ A}}\skt{~॥}\pend

\pstart \ledsidenote{\passref{4.40.18}}\edtext{\skt{तत्र}}{\lemma{\skt{तत्र}}\Afootnote{K H; \emph{om.}\ A}}\skt{ }\edtext{\edtext{\skt{प्रायो\-गि\-क\-स्त्रीं\-स्त्री\-नु\-च्छ्वा\-सा\-ना\-द\-द्या\-त्प\-र्या\-ग\-ता\-न्स्त्री\-न्स्त्रीं\-श्च\-तुरो}}{\lemma{\skt{प्रायोगिकस्त्रींस्}}\Afootnote{K; \skt{प्रायोगिकं त्रीं}\ A H}}}{\lemma{\skt{त्रीनु\appbrk{}च्छ्वा\appbrk{}सा\appbrk{}ना\appbrk{}द\appbrk{}द्या\appbrk{}त्प\appbrk{}र्या\appbrk{}ग\appbrk{}ता\appbrk{}न्स्त्री\appbrk{}न्स्त्रींश्}}\Afootnote{K; \skt{स्त्रीनु\appbrk{}च्छ्वा\appbrk{}सा\appbrk{}ना\appbrk{}द\appbrk{}दीत मुखनासिकाभ्यां च पर्यायांस्त्रींश्}\ A; \skt{त्रीनु\appbrk{}च्छ्वा\appbrk{}सा\appbrk{}ना\appbrk{}द\appbrk{}द्या\appbrk{}त्प\appbrk{}र्या\appbrk{}ग\appbrk{}ता\appbrk{}स्त्री\appbrk{}स्त्रीश्}\ H}}\skt{ वेति~॥ स्नैहिकं }\edtext{\skt{याव\-द\-स्रु\-प्र\-वृत्तिः}}{\lemma{\skt{याव\appbrk{}द\appbrk{}स्रु\appbrk{}प्र\appbrk{}वृत्तिः}}\Afootnote{K H; \skt{याव\appbrk{}द\appbrk{}श्रु\appbrk{}प्र\appbrk{}वृत्तिः,}\ A}}\skt{~॥ }\edtext{\edtext{\skt{वैरे\-च\-नि\-क\-मा\-दो\-ष\-द\-र्श\-नात्}}{\lemma{\skt{वैरेचनिकम्}}\Afootnote{K H; \skt{वैरेचनिक}\ A}}\skt{~॥ ग्रासान्तरेषु}}{\lemma{\skt{आदो\appbrk{}ष\appbrk{}द\appbrk{}र्श\appbrk{}नात्~॥ ग्रासान्तरेषु}}\Afootnote{K H; \skt{मादोषदर्शना}\ A}}\skt{ }\edtext{\skt{कासहरं~॥ वैरे\-च\-नि\-क\-व\-त्ति\-ल\-त\-ण्डु\-ल\-य\-वा\-गू\-पीतो }\edtext{\skt{वामनीयं}}{\lemma{\skt{वामनीयं}}\Afootnote{K H; \skt{वामनीयः,}\ A}}}{\lemma{\skt{त्… ति}}\Afootnote{\emph{om.}\ A}}\skt{ }\edtext{\skt{यथायोग}}{\lemma{\skt{यथायोग}}\Afootnote{K H; \skt{ग्रासान्तरेषु कासघ्न}\ A}}\skt{ इति~॥}\pend

\pstart\versepar \ledsidenote{\passref{4.40.19}}\edtext{\skt{तत्र }\edtext{\skt{व्रण\-ब\-स्ति\-ने\-त्रेण}}{\lemma{\skt{व्रण\appbrk{}ब\appbrk{}स्ति\appbrk{}ने\appbrk{}त्रेण}}\Afootnote{K; \skt{व्रण\appbrk{}ने\appbrk{}त्र\appbrk{}ब\appbrk{}स्ति\appbrk{}ने\appbrk{}त्रेण}\ H}}\skt{ }\edtext{\skt{व्रणधूमनेत्रं}}{\lemma{\skt{व्रणधूमनेत्रं}}\Afootnote{K H; \skt{व्रणधू}\ A}}}{\lemma{\skt{… व्र}}\Afootnote{\emph{om.}\ A}}\skt{ }\edtext{\edtext{\skt{व्याख्यातम्भवति}}{\lemma{\skt{व्याख्यातम्}}\Afootnote{K H; \skt{मं}\ A}}}{\lemma{\skt{भवति}}\Afootnote{K H; \emph{om.}\ A}}\skt{~॥ }\pend
\pstart\versepar \skt{}\edtext{\skt{शरा\-व\-स\-म्पु\-ट\-युतं}}{\lemma{\skt{शरा\appbrk{}व\appbrk{}स\appbrk{}म्पु\appbrk{}ट\appbrk{}युतं}}\Afootnote{K H; \skt{शरावसंपुटो}\ A}}\skt{ }\edtext{\edtext{\skt{व्रण\-धू\-म\-न्न\-ये\-द्व्रणं}}{\lemma{\skt{व्रणधूमन्}}\Afootnote{K H; \emph{om.}\ A}}}{\lemma{\skt{नयेद्व्रणं}}\Afootnote{K H; \skt{पनीते}\ A}}\skt{~। }\pend
\pstart\versepar \skt{}\edtext{\skt{धूमनेत्रेण}}{\lemma{\skt{धूमनेत्रेण}}\Afootnote{K H; \skt{न नेत्रेण}\ A}}\skt{ }\edtext{\skt{मति\-मा\-न्वे\-द\-ना\-स्रा\-व\-शा\-न्तये}}{\lemma{\skt{मति\appbrk{}मा\appbrk{}न्वे\appbrk{}द\appbrk{}ना\appbrk{}स्रा\appbrk{}व\appbrk{}शा\appbrk{}न्तये}}\Afootnote{K H; \skt{व्रणमानयेत्, धूम\appbrk{}पा\appbrk{}ना\appbrk{}द्वे\appbrk{}द\appbrk{}नो\appbrk{}प\appbrk{}शमो व्रण\appbrk{}वै\appbrk{}श\appbrk{}द्य\appbrk{}मा\appbrk{}स्रा\appbrk{}वो\appbrk{}प\appbrk{}श\appbrk{}मश्च भवति~॥}\ A}}\skt{~॥}\pend

\pstart\versepar \ledsidenote{\passref{4.40.20}}\edtext{\edtext{\skt{विधिरेष}}{\lemma{\skt{विधिर्}}\Afootnote{K H; \skt{विधि}\ A}}}{\lemma{\skt{एष}}\Afootnote{K H; \skt{रेष}\ A}}\skt{ समासेन }\edtext{\skt{धूमस्याभिहितो}}{\lemma{\skt{धूमस्याभिहितो}}\Afootnote{A K H\textsuperscript{pc}; \skt{धूप स्याभिहितो}\ H\textsuperscript{ac}}}\skt{ मया~॥ }\pend
\pstart\versepar \skt{नस्यस्यातः प्रवक्ष्यामि विधिं }\edtext{\skt{कुशलसम्मितं}}{\lemma{\skt{कुशलसम्मितं}}\Afootnote{K H; \skt{निरवशेषतः~॥}\ A}}\skt{~।}\pend

\pstart \ledsidenote{\passref{4.40.21}}\edtext{\skt{औषधपक्वो वा}}{\lemma{\skt{औषधपक्वो वा}}\Afootnote{K\textsuperscript{ac}; \skt{औष\appbrk{}ध\appbrk{}मौ\appbrk{}ष\appbrk{}ध\appbrk{}सिद्धो वा}\ A; \skt{औषध औषधपक्वो वा}\ K\textsuperscript{pc}; \skt{औष\appbrk{}ध\appbrk{}मौ\appbrk{}ष\appbrk{}ध\appbrk{}पक्वो वा}\ H}}\skt{ स्नेहो नासिकाभ्यां }\edtext{\skt{प्रदीयते}}{\lemma{\skt{प्रदीयते}}\Afootnote{K; \skt{दीयत इति नस्यम्~। तद्द्विविधं शिरोविरे}\ A; \skt{प्रदीयते~। नस्यन्त}\ H}}\skt{~। }\edtext{\edtext{\edtext{\skt{नस्य\-न्त\-त्प\-ञ्च\-वि\-क\-ल्प\-म्भ\-वति}}{\lemma{\skt{नस्यन्}}\Afootnote{K; \skt{चनं, स्ने}\ A; \emph{om.}\ H}}}{\lemma{\skt{तत्प\appbrk{}ञ्च\appbrk{}वि\appbrk{}क\appbrk{}ल्पम्}}\Afootnote{K; \skt{हनं च~। तद्द्विविधमपि}\ A; \skt{त्पञ्चविकल्पम्}\ H}}}{\lemma{\skt{भवति}}\Afootnote{K H; \skt{पञ्चधा~।}\ A}}\skt{~। तद्यथा~। नस्यं }\edtext{\skt{शिरोविरेचनं }\edtext{\skt{प्रतिमर्शो}}{\lemma{\skt{प्रतिमर्शो}}\Afootnote{K H\textsuperscript{pc}; \skt{प्रतिमर्शो,}\ A}}\skt{ }\edtext{\skt{वपीडः}}{\lemma{\skt{वपीडः}}\Afootnote{K H\textsuperscript{pc}; \skt{अवपीडः,}\ A}}\skt{ }\edtext{\skt{प्रधमनमिति}}{\lemma{\skt{इति}}\Afootnote{K H\textsuperscript{pc}; \skt{च~।}\ A}}\skt{~। }\edtext{\edtext{\skt{तेषा\-न्न\-स्य\-म्प्र\-धाणं}}{\lemma{\skt{तेषान्}}\Afootnote{K H\textsuperscript{pc}; \skt{तेषु}\ A}}}{\lemma{\skt{प्रधाणं}}\Afootnote{K; \skt{प्रधानं}\ A H\textsuperscript{pc}}}\skt{ शिरोविरेचनं}}{\lemma{\skt{नस्यं… न}}\Afootnote{\emph{om.}\ H\textsuperscript{ac}}}\skt{ }\edtext{\edtext{\skt{येन}}{\lemma{\skt{येन}}\Afootnote{K H\textsuperscript{pc}; \skt{च;}\ A}}\skt{ }\edtext{\skt{नस्य}}{\lemma{\skt{नस्य}}\Afootnote{K H\textsuperscript{pc}; \skt{न}\ A}}\skt{ }\edtext{\skt{विकल्पः}}{\lemma{\skt{विकल्पः}}\Afootnote{K; \skt{स्यविकल्पः}\ A; \skt{इति विकल्पः}\ H\textsuperscript{pc}}}\skt{ प्रतिमर्शः~। शिरो\-वि\-रे\-च\-न\-वि\-कल्पो }\edtext{\skt{वपीडः}}{\lemma{\skt{वपीडः}}\Afootnote{K H; \skt{ऽवपीडः}\ A}}\skt{ }\edtext{\skt{प्रधमनञ्चेति}}{\lemma{\skt{चेति}}\Afootnote{K H; \skt{च;}\ A}}}{\lemma{\skt{प्रधाणं… शि}}\Afootnote{\emph{om.}\ H\textsuperscript{ac}}}\skt{~। ततो }\edtext{\skt{नस्तशब्दः}}{\lemma{\skt{नस्तशब्दः}}\Afootnote{K H; \skt{नस्यशब्दः}\ A}}\skt{ }\edtext{\skt{पञ्चविधा}}{\lemma{\skt{पञ्चविधा}}\Afootnote{K H; \skt{पञ्चधा}\ A}}\skt{ }\edtext{\skt{णिपतति}}{\lemma{\skt{णिपतति}}\Afootnote{K; \skt{नियमितः~॥}\ A; \skt{निपतति~।}\ H}}\skt{~।}\pend

\pstart \ledsidenote{\passref{4.40.22}}\skt{तत्र यः }\edtext{\skt{स्नेहार्थं}}{\lemma{\skt{स्नेहार्थं}}\Afootnote{K H; \skt{स्नेहनार्थं}\ A}}\skt{ शून्यशिरसां }\edtext{\skt{ग्रीवास्कन्धोरसां}}{\lemma{\skt{ग्रीवास्कन्धोरसां}}\Afootnote{K H; \skt{ग्रीवास्कन्धोरसां च}\ A}}\skt{ बलजननार्थं दृष्टि\-प्र\-सा\-द\-ज\-न\-ना\-र्थम्वा }\edtext{\skt{स्नेह}}{\lemma{\skt{स्नेह}}\Afootnote{K H; \skt{स्नेहो}\ A}}\skt{ }\edtext{\skt{उपादीयते}}{\lemma{\skt{उपादीयते}}\Afootnote{K; \skt{विधीयते}\ A; \skt{उपादीर्यते~।}\ H}}\skt{ }\edtext{\skt{तस्मिन्वैशेषिके}}{\lemma{\skt{तस्मिन्वैशेषिके}}\Afootnote{K H; \skt{तस्मिन्वैशेषिको}\ A}}\skt{ }\edtext{\skt{नस्तशब्दः}}{\lemma{\skt{नस्तशब्दः}}\Afootnote{K H; \skt{नस्यशब्दः~।}\ A}}\skt{~। तत्तु }\edtext{\skt{देय\-म्वा\-ता\-भि\-पन्ने}}{\lemma{\skt{वाताभिपन्ने}}\Afootnote{K H; \skt{वाताभिभूते}\ A}}\skt{ शिरसि }\edtext{\skt{तथा}}{\lemma{\skt{तथा}}\Afootnote{K; \emph{om.}\ A; \skt{तद्यथा~। तथा}\ H}}\skt{ }\edtext{\skt{दन्त\-के\-श\-श्म\-स्रु\-णा\-त\-दा\-रु\-ण\-क\-र्ण्ण\-क्ष्वे\-ल\-ति\-मि\-र\-स्व\-रो\-प\-घा\-त\-ना\-सा\-रो\-ग\-मु\-ख\-रो\-ग\-शो\-षा\-व\-बा\-हु\-का\-का\-ल\-ज\-व\-ली\-प\-लि\-त\-प्र\-बो\-धेषु}}{\lemma{\skt{दन्त\appbrk{}के\appbrk{}श\appbrk{}श्म\appbrk{}स्रु\appbrk{}णा\appbrk{}त\appbrk{}दा\appbrk{}रु\appbrk{}ण\appbrk{}क\appbrk{}र्ण्ण\appbrk{}क्ष्वे\appbrk{}ल\appbrk{}ति\appbrk{}मि\appbrk{}र\appbrk{}स्व\appbrk{}रो\appbrk{}प\appbrk{}घा\appbrk{}त\appbrk{}ना\appbrk{}सा\appbrk{}रो\appbrk{}ग\appbrk{}मु\appbrk{}ख\appbrk{}रो\appbrk{}ग\appbrk{}शो\appbrk{}षा\appbrk{}व\appbrk{}बा\appbrk{}हु\appbrk{}का\appbrk{}का\appbrk{}ल\appbrk{}ज\appbrk{}व\appbrk{}ली\appbrk{}प\appbrk{}लि\appbrk{}त\appbrk{}प्र\appbrk{}बो\appbrk{}धेषु}}\Afootnote{K H; \skt{दन्त\appbrk{}के\appbrk{}श\appbrk{}श्म\appbrk{}श्रु\appbrk{}प्र\appbrk{}पा\appbrk{}त\appbrk{}दा\appbrk{}रु\appbrk{}ण\appbrk{}क\appbrk{}र्ण\appbrk{}शू\appbrk{}ल\appbrk{}क\appbrk{}र्ण\appbrk{}क्ष्वे\appbrk{}ड\appbrk{}ति\appbrk{}मि\appbrk{}र\appbrk{}स्व\appbrk{}रो\appbrk{}प\appbrk{}घा\appbrk{}त\appbrk{}ना\appbrk{}सा\appbrk{}रो\appbrk{}गा\appbrk{}स्य\appbrk{}शो\appbrk{}षा\appbrk{}व\appbrk{}बा\appbrk{}हु\appbrk{}का\appbrk{}का\appbrk{}ल\appbrk{}ज\appbrk{}व\appbrk{}ली\appbrk{}प\appbrk{}लि\appbrk{}त\appbrk{}प्रा\appbrk{}दु\appbrk{}र्भा\appbrk{}व\appbrk{}दा\appbrk{}रु\appbrk{}ण\appbrk{}प्र\appbrk{}बो\appbrk{}धेषु}\ A}}\skt{ }\edtext{\skt{वातपैत्तिकेषु च}}{\lemma{\skt{वातपैत्तिकेषु च}}\Afootnote{K H; \skt{वातपैत्तिकेषु}\ A}}\skt{ }\edtext{\skt{मुखरोगेषु}}{\lemma{\skt{मुखरोगेषु}}\Afootnote{K H; \skt{मुख\appbrk{}रो\appbrk{}गे\appbrk{}ष्व\appbrk{}न्येषु च}\ A}}\skt{ वात\-पि\-त्त\-ह\-र\-द्र\-व्य\-सि\-द्धेन स्नेहेनेति~॥}\pend

\pstart \ledsidenote{\passref{4.40.23}}\edtext{\skt{शिरो\-वि\-रे\-च\-न\-मपि}}{\lemma{\skt{अपि}}\Afootnote{K H; \emph{om.}\ A}}\skt{ }\edtext{\skt{श्लेष्म\-णा\-भि\-व्या\-प्त\-ता\-लु\-ग\-त\-शि\-रसां}}{\lemma{\skt{श्लेष्म\appbrk{}णा\appbrk{}भि\appbrk{}व्या\appbrk{}प्त\appbrk{}ता\appbrk{}लु\appbrk{}ग\appbrk{}त\appbrk{}शि\appbrk{}रसां}}\Afootnote{K H; \skt{श्लेष्म\appbrk{}णाऽ\appbrk{}भि\appbrk{}व्या\appbrk{}प्त\appbrk{}ता\appbrk{}लु\appbrk{}क\appbrk{}ण्ठ\appbrk{}शि\appbrk{}रसा}\ A}}\skt{ }\edtext{\skt{तथा}}{\lemma{\skt{तथा}}\Afootnote{K H; \skt{म}\ A}}\skt{ }\edtext{\skt{रोच\-क\-शि\-रो\-गौ\-र\-व\-शू\-ल\-पी\-न\-सा\-व\-भे\-द\-क\-क्रि\-मि\-प्र\-ति\-श्या\-या\-प\-स्मा\-र\-घ्रा\-णा\-ज्ञा\-ने\-ष्व\-न्येषु}}{\lemma{\skt{रोच\appbrk{}क\appbrk{}शि\appbrk{}रो\appbrk{}गौ\appbrk{}र\appbrk{}व\appbrk{}शू\appbrk{}ल\appbrk{}पी\appbrk{}न\appbrk{}सा\appbrk{}व\appbrk{}भे\appbrk{}द\appbrk{}क\appbrk{}क्रि\appbrk{}मि\appbrk{}प्र\appbrk{}ति\appbrk{}श्या\appbrk{}या\appbrk{}प\appbrk{}स्मा\appbrk{}र\appbrk{}घ्रा\appbrk{}णा\appbrk{}ज्ञा\appbrk{}नेष्व्}}\Afootnote{K H; \skt{रोच\appbrk{}क\appbrk{}शि\appbrk{}रो\appbrk{}गौ\appbrk{}र\appbrk{}व\appbrk{}शू\appbrk{}ल\appbrk{}पी\appbrk{}न\appbrk{}सा\appbrk{}र्धा\appbrk{}व\appbrk{}भे\appbrk{}द\appbrk{}क\appbrk{}कृ\appbrk{}मि\appbrk{}प्र\appbrk{}ति\appbrk{}श्या\appbrk{}या\appbrk{}प\appbrk{}स्मा\appbrk{}र\appbrk{}ग\appbrk{}न्धा\appbrk{}ज्ञा\appbrk{}नेष्व्}\ A}}\skt{ }\edtext{\skt{चोर्ध्व\-ज\-त्रु\-ग\-तेषु}}{\lemma{\skt{चोर्ध्व\appbrk{}ज\appbrk{}त्रु\appbrk{}ग\appbrk{}तेषु}}\Afootnote{K H; \skt{चोर्ध्व\appbrk{}ज\appbrk{}त्रु\appbrk{}ग\appbrk{}तेषु कफजेषु}\ A}}\skt{ विकारेषु }\edtext{\skt{शिरो\-वि\-रे\-च\-न\-द्र\-व्य\-सि\-द्धेन}}{\lemma{\skt{शिरो\appbrk{}वि\appbrk{}रे\appbrk{}च\appbrk{}न\appbrk{}द्र\appbrk{}व्य\appbrk{}सि\appbrk{}द्धेन}}\Afootnote{K H; \skt{शिरो\appbrk{}वि\appbrk{}रे\appbrk{}च\appbrk{}न\appbrk{}द्र\appbrk{}व्यै\appbrk{}स्त\appbrk{}त्सि\appbrk{}द्धेन}\ A}}\skt{ }\edtext{\skt{स्नेहेनेति}}{\lemma{\skt{स्नेहेनेति}}\Afootnote{K; \skt{वा स्नेहेनेति~॥}\ A; \skt{स्नेहेनैवेति~॥}\ H}}\skt{~।}\pend

\pstart \ledsidenote{\passref{4.40.24}}\edtext{\edtext{\edtext{\skt{तत्त्वे\-त\-द्द्वि\-वि\-ध\-म\-प्युक्तं}}{\lemma{\skt{तत्त्वे\appbrk{}त\appbrk{}द्द्वि\appbrk{}वि\appbrk{}धम्}}\Afootnote{K H; \skt{तत्रैतद्द्विविध}\ A}}}{\lemma{\skt{अप्य्}}\Afootnote{K H; \skt{मप्य}\ A}}}{\lemma{\skt{उक्तं}}\Afootnote{K H; \skt{भुक्त\appbrk{}व\appbrk{}तोऽ\appbrk{}न्न\appbrk{}काले पूर्वाह्णे श्लेष्मरोगिणां, मध्याह्ने पित्तरोगिणां, अपराह्णे वातरोगिणाम्~॥}\ A}}\skt{~॥}\pend

\pstart \ledsidenote{\passref{4.40.25}}\edtext{\skt{तत्रो\-च्चा\-रि\-त\-वा\-त\-मू\-त्र\-पु\-री\-षाय}}{\lemma{\skt{तत्रो\appbrk{}च्चा\appbrk{}रि\appbrk{}त\appbrk{}वा\appbrk{}त\appbrk{}मू\appbrk{}त्र\appbrk{}पु\appbrk{}री\appbrk{}षाय}}\Afootnote{K; \skt{अथ पुरुषाय शिरो\appbrk{}वि\appbrk{}रे\appbrk{}च\appbrk{}नी\appbrk{}याय त्यक्त\appbrk{}मू\appbrk{}त्र\appbrk{}पु\appbrk{}री\appbrk{}षाय}\ A; \skt{तत्रोच्चारित तमू\appbrk{}त्र\appbrk{}पु\appbrk{}री\appbrk{}षाया}\ H\textsuperscript{ac}; \skt{तत्रोच्चारित वात\appbrk{}मू\appbrk{}त्र\appbrk{}पु\appbrk{}री\appbrk{}षाया}\ H\textsuperscript{pc}}}\skt{ भुक्तवते }\edtext{\skt{व्यभ्रे}}{\lemma{\skt{व्यभ्रे}}\Afootnote{A K; \emph{om.}\ H}}\skt{ }\edtext{\skt{काले}}{\lemma{\skt{काले}}\Afootnote{A K; \emph{om.}\ H}}\skt{ }\edtext{\skt{पाणि\-ता\-पो\-प\-स्वे\-दि\-त\-ग\-ल\-क\-पो\-ल\-ल\-ला\-ट\-प्र\-दे\-शाय}}{\lemma{\skt{पाणि\appbrk{}ता\appbrk{}पो\appbrk{}प\appbrk{}स्वे\appbrk{}दि\appbrk{}त\appbrk{}ग\appbrk{}ल\appbrk{}क\appbrk{}पो\appbrk{}ल\appbrk{}ल\appbrk{}ला\appbrk{}ट\appbrk{}प्र\appbrk{}दे\appbrk{}शाय}}\Afootnote{K; \skt{दन्त\appbrk{}का\appbrk{}ष्ठ\appbrk{}धू\appbrk{}म\appbrk{}पा\appbrk{}नाभ्यां विशु\appbrk{}द्ध\appbrk{}व\appbrk{}क्त्र\appbrk{}स्रो\appbrk{}तसे पाणि\appbrk{}ता\appbrk{}प\appbrk{}प\appbrk{}रि\appbrk{}स्वि\appbrk{}न्न\appbrk{}मृ\appbrk{}दि\appbrk{}त\appbrk{}ग\appbrk{}ल\appbrk{}क\appbrk{}पो\appbrk{}ल\appbrk{}ल\appbrk{}ला\appbrk{}ट\appbrk{}प्र\appbrk{}दे\appbrk{}शाय}\ A; \skt{द्रव्यत्रे काले पाणि\appbrk{}ता\appbrk{}पो\appbrk{}प\appbrk{}स्वे\appbrk{}दि\appbrk{}त\appbrk{}मृ\appbrk{}दि\appbrk{}त\appbrk{}ग\appbrk{}ल\appbrk{}क\appbrk{}पो\appbrk{}ल\appbrk{}ल\appbrk{}ला\appbrk{}ट\appbrk{}प्र\appbrk{}दे\appbrk{}शाय}\ H}}\skt{ वाता\-त\-प\-र\-जो\-हीने }\edtext{\edtext{\skt{वेश्म\-न्यु\-त्ता\-न\-शा\-यिने}}{\lemma{\skt{वेश्मन्य्}}\Afootnote{K H; \skt{वेश्म}\ A}}}{\lemma{\skt{उत्तानशायिने}}\Afootnote{K H; \skt{न्युत्तानशायिने}\ A}}\skt{ }\edtext{\skt{किञ्चि\-त्प्र\-वि\-ल\-म्बि\-त\-शि\-रसे}}{\lemma{\skt{किञ्चि\appbrk{}त्प्र\appbrk{}वि\appbrk{}ल\appbrk{}म्बि\appbrk{}त\appbrk{}शि\appbrk{}रसे}}\Afootnote{K H; \skt{प्रसा\appbrk{}रि\appbrk{}त\appbrk{}क\appbrk{}र\appbrk{}च\appbrk{}र\appbrk{}णाय किञ्चि\appbrk{}त्प्र\appbrk{}वि\appbrk{}ल\appbrk{}म्बि\appbrk{}त\appbrk{}शि\appbrk{}रसे}\ A}}\skt{ }\edtext{\skt{वस्त्रा\-व\-च्छा\-दि\-त\-ने\-त्राय}}{\lemma{\skt{वस्त्रा\appbrk{}व\appbrk{}च्छा\appbrk{}दि\appbrk{}त\appbrk{}ने\appbrk{}त्राय}}\Afootnote{K H; \skt{वस्त्रा\appbrk{}च्छा\appbrk{}दि\appbrk{}त\appbrk{}ने\appbrk{}त्राय}\ A}}\skt{ वाम\-ह\-स्त\-प्र\-दे\-शि\-न्य\-ग्रो\-न्ना\-मि\-त\-ना\-सा\-ग्राय }\edtext{\skt{विशुद्धस्रोतसे}}{\lemma{\skt{विशुद्धस्रोतसे}}\Afootnote{K H; \skt{विशुद्धस्रोतसि}\ A}}\skt{ दक्षिणहस्तेन }\edtext{\skt{शुक्त्या}}{\lemma{\skt{शुक्त्या}}\Afootnote{K H; \skt{स्नेहमुष्णाम्बुना प्रतप्तं रज\appbrk{}त\appbrk{}सु\appbrk{}व\appbrk{}र्ण\appbrk{}ता\appbrk{}म्र\appbrk{}म\appbrk{}णि\appbrk{}मृ\appbrk{}त्पा\appbrk{}त्र\appbrk{}शु\appbrk{}क्ती\appbrk{}ना\appbrk{}म\appbrk{}न्य\appbrk{}त\appbrk{}मस्थं शुक्त्या}\ A}}\skt{ पिचुना वा }\edtext{\edtext{\edtext{\skt{सुखो\-ष्ण\-म\-द्रु\-त\-म\-वि\-च्छि\-न्न\-धा\-रा\-मा\-सि\-ञ्चे\-दिति}}{\lemma{\skt{अद्रुतम्}}\Afootnote{K; \skt{स्नेहमद्रुतम्}\ A; \emph{om.}\ H}}}{\lemma{\skt{अवि\appbrk{}च्छि\appbrk{}न्न\appbrk{}धा\appbrk{}राम्}}\Afootnote{K; \skt{आसि\appbrk{}ञ्चे\appbrk{}द\appbrk{}व्य\appbrk{}व\appbrk{}च्छि\appbrk{}न्न\appbrk{}धारं}\ A; \skt{अद्रुतमव च्छिन्नधारा}\ H\textsuperscript{ac}; \skt{अद्रुतम विच्छिन्नधारा}\ H\textsuperscript{pc}}}}{\lemma{\skt{आसिञ्चेदिति}}\Afootnote{K; \emph{om.}\ A; \skt{सा सिञ्चेदिति~।}\ H}}\skt{~। }\edtext{\skt{यथा च}}{\lemma{\skt{यथा च}}\Afootnote{K H; \skt{यथा}\ A}}\skt{ }\edtext{\skt{नेत्रेण}}{\lemma{\skt{नेत्रेण}}\Afootnote{K H; \skt{नेत्रे न}\ A}}\skt{ प्राप्नोति~॥}\pend

\pstart\versepar \ledsidenote{\passref{4.40.26}}\skt{स्नेहे }\edtext{\skt{वसिच्यति}}{\lemma{\skt{वसिच्यति}}\Afootnote{K; \skt{ऽवसिच्यमाने तु}\ A; \skt{ऽवसिच्यति}\ H}}\skt{ }\edtext{\skt{शिरो न}}{\lemma{\skt{शिरो न}}\Afootnote{K H; \skt{शिरो नै}\ A}}\skt{ }\edtext{\edtext{\skt{कथञ्चन}}{\lemma{\skt{कथञ्}}\Afootnote{K; \skt{व}\ A; \skt{कञ्}\ H}}}{\lemma{\skt{चन}}\Afootnote{K H; \skt{प्र}\ A}}\skt{ कम्पयेत्~। }\pend
\pstart\versepar \skt{न }\edtext{\skt{भाषेन्न च}}{\lemma{\skt{भाषेन्न च}}\Afootnote{K H; \skt{कुप्येन्न प्र}\ A}}\skt{ }\edtext{\skt{कुप्येच्च न}}{\lemma{\skt{कुप्येच्च न}}\Afootnote{K H; \skt{भाषेच्च न}\ A}}\skt{ }\edtext{\skt{क्षुयान्न}}{\lemma{\skt{क्षुयान्न}}\Afootnote{A K; \skt{क्षुधान्न}\ H}}\skt{ }\edtext{\skt{हसेदपि}}{\lemma{\skt{हसेदपि}}\Afootnote{K; \skt{हसेत्तथा~॥}\ A; \skt{सहेदपि~॥}\ H}}\skt{~।}\pend

\pstart\versepar \ledsidenote{\passref{4.40.27}}\skt{एतैर्हि विहतः स्नेहो न सम्य\-क्प्र\-ति\-प\-द्यते~। }\pend
\pstart\versepar \skt{ततः }\edtext{\skt{कासः}}{\lemma{\skt{कासः}}\Afootnote{K; \skt{कास}\ A H}}\skt{ }\edtext{\skt{प्रतिश्यायः}}{\lemma{\skt{प्रतिश्यायः}}\Afootnote{K H; \skt{प्रतिश्याय}\ A}}\skt{ }\edtext{\skt{शिरो\-क्षि\-ग\-द\-स\-म्भवः}}{\lemma{\skt{शिरो\appbrk{}क्षि\appbrk{}ग\appbrk{}द\appbrk{}स\appbrk{}म्भवः}}\Afootnote{K H; \skt{शिरोऽ\appbrk{}क्षि\appbrk{}ग\appbrk{}द\appbrk{}सं\appbrk{}भवः~॥}\ A}}\skt{~।}\pend

\pstart \ledsidenote{\passref{4.40.28}}\skt{तस्य }\edtext{\edtext{\skt{प्रमाणन्त्वष्टौ}}{\lemma{\skt{प्रमाणन्}}\Afootnote{K H; \skt{प्रमाण}\ A}}}{\lemma{\skt{त्वष्टौ}}\Afootnote{K H; \skt{मष्तौ}\ A}}\skt{ बिन्दवः }\edtext{\skt{प्रदे\-शि\-नी\-प\-र्व\-द्व\-य\-निः\-स्रुताः}}{\lemma{\skt{प्रदे\appbrk{}शि\appbrk{}नी\appbrk{}प\appbrk{}र्व\appbrk{}द्व\appbrk{}य\appbrk{}निः\appbrk{}स्रुताः}}\Afootnote{K H; \skt{प्रदे\appbrk{}शि\appbrk{}नी\appbrk{}प\appbrk{}र्व\appbrk{}द्व\appbrk{}य\appbrk{}निः\appbrk{}सृताः}\ A}}\skt{~। }\edtext{\skt{एषा}}{\lemma{\skt{एषा}}\Afootnote{K H; \emph{om.}\ A}}\skt{ }\edtext{\skt{मात्रा}}{\lemma{\skt{मात्रा}}\Afootnote{K H; \emph{om.}\ A}}\skt{ }\edtext{\skt{प्रथमा}}{\lemma{\skt{प्रथमा}}\Afootnote{K H; \skt{प्रथमा मात्रा,}\ A}}\skt{~। द्वितीया शुक्तिः तृतीया }\edtext{\edtext{\skt{पाणि\-शु\-क्ति\-रि\-त्ये\-ता\-स्तिस्रो}}{\lemma{\skt{एतास्}}\Afootnote{K H; \skt{एता}\ A}}}{\lemma{\skt{तिस्रो}}\Afootnote{K H; \skt{स्तिस्रो}\ A}}\skt{ }\edtext{\skt{मानाः}}{\lemma{\skt{मानाः}}\Afootnote{K H; \skt{मात्रा}\ A}}\skt{ }\edtext{\skt{यथा\-ब\-ल\-म्प्र\-योज्या}}{\lemma{\skt{प्रयोज्या}}\Afootnote{K H; \skt{प्रयोज्याः~॥}\ A}}\skt{ }\edtext{\skt{भवन्तीति}}{\lemma{\skt{भवन्तीति}}\Afootnote{K H; \emph{om.}\ A}}\skt{~॥}\pend

\pstart\versepar \ledsidenote{\passref{4.40.30}}\skt{शृङ्गा\-ट\-क\-म\-भि\-प्लाव्य निरेति वदनाद्यथा~। }\pend
\pstart\versepar \skt{}\edtext{\skt{कफो\-त्क्ले\-श\-भ\-या\-च्चै\-न\-न्नि\-ष्ठी\-वे\-द\-वि\-धा\-र\-यन्}}{\lemma{\skt{निष्ठी\appbrk{}वे\appbrk{}द\appbrk{}वि\appbrk{}धा\appbrk{}र\appbrk{}यन्}}\Afootnote{K; \skt{निष्ठी\appbrk{}वे\appbrk{}द\appbrk{}वि\appbrk{}धा\appbrk{}र\appbrk{}यन्~॥}\ A; \skt{निष्ठी\appbrk{}वे\appbrk{}द\appbrk{}भि\appbrk{}धा\appbrk{}र\appbrk{}यन्~॥}\ H}}\skt{~॥}\pend

\pstart \ledsidenote{\passref{4.40.31}}\edtext{\skt{दन्ते}}{\lemma{\skt{दन्ते}}\Afootnote{K H; \skt{दत्ते च}\ A}}\skt{ }\edtext{\edtext{\skt{पुनरपि}}{\lemma{\skt{पुनर्}}\Afootnote{K H; \skt{पुन}\ A}}}{\lemma{\skt{अपि}}\Afootnote{K H; \skt{रपि}\ A}}\skt{ संस्वेद्य }\edtext{\edtext{\skt{गल\-क\-पो\-ला\-दी\-न्धू\-म\-मा\-से\-वेत}}{\lemma{\skt{धूमम्}}\Afootnote{K H; \skt{धूम}\ A}}}{\lemma{\skt{आसेवेत}}\Afootnote{K H; \skt{मासेवेत,}\ A}}\skt{~। }\edtext{\edtext{\skt{भोज\-ये\-च्चै\-न\-म\-न\-भि\-ष्यन्ति}}{\lemma{\skt{भोजयेच्चैनम्}}\Afootnote{K H; \skt{भोजयेच्चैनम्}\ A}}}{\lemma{\skt{अनभिष्यन्ति}}\Afootnote{K H\textsuperscript{ac}; \skt{अभिष्यन्दि,}\ A; \skt{अनभिष्यन्दि~।}\ H\textsuperscript{pc}}}\skt{~। }\edtext{\edtext{\skt{अथा\-स्या\-चा\-रि\-क\-मु\-प\-दि\-शेत्}}{\lemma{\skt{अथा\appbrk{}स्या\appbrk{}चा\appbrk{}रि\appbrk{}कम्}}\Afootnote{K H\textsuperscript{pc}; \skt{ततोऽ\appbrk{}स्या\appbrk{}चा\appbrk{}रि\appbrk{}कम्}\ A; \skt{य थास्याचारिकम्}\ H\textsuperscript{ac}}}\skt{~। क्रोधा\-दि\-म्प\-रि\-ह\-रेच्च}}{\lemma{\skt{उपदिशेत्~। क्रोधादिम्}}\Afootnote{K H; \skt{आदिशे}\ A}}\skt{~। }\edtext{\skt{रजो\-धू\-म\-म\-न\-स्ता\-प\-द्र\-व\-पा\-न\-शि\-र\-स्ना\-ना\-दीनि}}{\lemma{\skt{रजो\appbrk{}धू\appbrk{}म\appbrk{}म\appbrk{}न\appbrk{}स्ता\appbrk{}प\appbrk{}द्र\appbrk{}व\appbrk{}पा\appbrk{}न\appbrk{}शि\appbrk{}र\appbrk{}स्ना\appbrk{}ना\appbrk{}दीनि}}\Afootnote{K H; \skt{त्; रजो\appbrk{}धू\appbrk{}म\appbrk{}स्ने\appbrk{}हा\appbrk{}त\appbrk{}प\appbrk{}म\appbrk{}द्य\appbrk{}द्र\appbrk{}व\appbrk{}पा\appbrk{}न\appbrk{}शि\appbrk{}रः\appbrk{}स्ना\appbrk{}ना\appbrk{}ति\appbrk{}या\appbrk{}न\appbrk{}क्रो\appbrk{}धा\appbrk{}दीनि च परिहरेत्~॥}\ A}}\skt{~।}\edtext{}{\lemma{\skt{त्… र}}\Afootnote{\emph{om.}\ A}}\pend

\pstart \ledsidenote{\passref{4.40.32}}\edtext{\skt{तस्य योगा\-यो\-गा\-ति\-यो\-ग\-वि\-ज्ञानं}}{\lemma{\skt{तस्य योगा\appbrk{}यो\appbrk{}गा\appbrk{}ति\appbrk{}यो\appbrk{}ग\appbrk{}वि\appbrk{}ज्ञानं}}\Afootnote{K H; \skt{तस्य योगा\appbrk{}ति\appbrk{}यो\appbrk{}गा\appbrk{}यो\appbrk{}गा\appbrk{}ना\appbrk{}मिदं विज्ञानं}\ A}}\skt{ भवति~।}\pend

\pstart\versepar \ledsidenote{\passref{4.40.33}}\skt{लाघवं शिरसो योगे }\edtext{\skt{सुख\-स्व\-प्न\-प्र\-बो\-धणं}}{\lemma{\skt{सुख\appbrk{}स्व\appbrk{}प्न\appbrk{}प्र\appbrk{}बो\appbrk{}धणं}}\Afootnote{K; \skt{सुख\appbrk{}स्व\appbrk{}प्न\appbrk{}प्र\appbrk{}बो\appbrk{}ध\appbrk{}नम्}\ A H}}\skt{~। }\pend
\pstart\versepar \skt{विकारोपशमः }\edtext{\edtext{\edtext{\skt{शुद्धि\-रि\-न्द्रि\-या\-णा\-म्परं}}{\lemma{\skt{शुद्धिर्}}\Afootnote{K H; \skt{शुद्धि}\ A}}}{\lemma{\skt{इन्द्रियाणाम्}}\Afootnote{K H; \skt{रिन्दिर्याणां}\ A}}}{\lemma{\skt{परं}}\Afootnote{K H; \skt{मनः}\ A}}\skt{ सुखं~॥}\pend

\pstart \ledsidenote{\passref{4.40.34}}\edtext{\skt{अयोगे}}{\lemma{\skt{अयोगे}}\Afootnote{K H; \emph{om.}\ A}}\skt{ }\edtext{\skt{रोगानिर्जयः}}{\lemma{\skt{रोगानिर्जयः}}\Afootnote{K H; \emph{om.}\ A}}\skt{ कफप्रसेकः शिरसो }\edtext{\skt{गुरु\-ते\-न्द्रि\-य\-वि\-भ्रमः~। क्>}}{\lemma{\skt{गुरु\appbrk{}ते\appbrk{}न्द्रि\appbrk{}य\appbrk{}वि\appbrk{}भ्रमः~। क्>}}\Afootnote{\skt{गुरु\appbrk{}ते\appbrk{}न्द्रि\appbrk{}य\appbrk{}वि\appbrk{}भ्रमः}\ A K H}}\skt{ }\edtext{\edtext{\skt{अतियोगे तालु\-शो\-ष\-पि\-पा\-सा\-भ्र\-म\-म\-द\-मूर्च्छा चेति }\edtext{\edtext{\skt{लक्ष\-ण\-म्मू\-र्ध्न्य\-ति\-स्निग्धे}}{\lemma{\skt{मूर्ध्न्य्}}\Afootnote{K H; \skt{मूर्ध्न}\ A}}}{\lemma{\skt{अतिस्निग्धे}}\Afootnote{K H; \skt{यतिस्निग्धे}\ A}}}{\lemma{\skt{क्>… ल}}\Afootnote{\emph{om.}\ A}}}{\lemma{\skt{क्>… ल}}\Afootnote{\emph{om.}\ K\textsuperscript{ac}}}\skt{ रूक्षं तत्रा\-व\-चा\-र\-येत्~।}\pend

\pstart \ledsidenote{\passref{4.40.36}}\edtext{\skt{शिरो\-वि\-रे\-च\-न\-स्यापि}}{\lemma{\skt{शिरो\appbrk{}वि\appbrk{}रे\appbrk{}च\appbrk{}न\appbrk{}स्यापि}}\Afootnote{K H; \emph{om.}\ A}}\skt{ चत्वारो बिन्दवः षड्वा }\edtext{\skt{तथाष्टौ वा}}{\lemma{\skt{तथाष्टौ वा}}\Afootnote{K H; \skt{तथाऽष्टौ वा}\ A}}\skt{ }\edtext{\skt{यथा\-ब\-ल\-म्शि\-रो\-वि\-च\-क\-स्ने\-हस्य}}{\lemma{\skt{शिरो\appbrk{}वि\appbrk{}च\appbrk{}क\appbrk{}स्ने\appbrk{}हस्य}}\Afootnote{K; \skt{शिरो\appbrk{}वि\appbrk{}रे\appbrk{}क\appbrk{}स्ने\appbrk{}हस्य}\ A; \skt{शिरो\appbrk{}वि\appbrk{}रे\appbrk{}च\appbrk{}क\appbrk{}स्ने\appbrk{}हस्य}\ H}}\skt{ प्रमा\-ण\-म\-भि\-नि\-र्दि\-शेत्~।}\pend

\pstart\versepar \ledsidenote{\passref{4.40.37}}\edtext{\skt{तस्य}}{\lemma{\skt{तस्य}}\Afootnote{K H; \skt{नस्ये}\ A}}\skt{ }\edtext{\edtext{\skt{त्रीण्यभिवर्ण्ण्यन्ते}}{\lemma{\skt{त्रीण्य्}}\Afootnote{K H; \skt{त्रीण्यु}\ A}}}{\lemma{\skt{अभिवर्ण्ण्यन्ते}}\Afootnote{K H; \skt{पदिष्टानि}\ A}}\skt{ लक्षणानि प्रयोगतः~। }\pend
\pstart\versepar \skt{}\edtext{\skt{शुद्धि\-ही\-ना\-ति\-स\-ञ्ज्ञानि}}{\lemma{\skt{शुद्धि\appbrk{}ही\appbrk{}ना\appbrk{}ति\appbrk{}स\appbrk{}ञ्ज्ञानि}}\Afootnote{K H; \skt{शुद्ध(\appbrk{}शु\appbrk{}द्धि)\appbrk{}ही\appbrk{}ना\appbrk{}ति\appbrk{}सं\appbrk{}ज्ञानि}\ A}}\skt{ विशे\-षा\-च्छा\-स्त्र\-चि\-न्तकैः~।}\pend

\pstart \ledsidenote{\passref{4.40.38}}\skt{लाघवं शिरसः }\edtext{\skt{शुद्धिः}}{\lemma{\skt{शुद्धिः}}\Afootnote{K H; \skt{सुद्धिः}\ A}}\skt{ }\edtext{\skt{स्रोतसां}}{\lemma{\skt{स्रोतसां}}\Afootnote{A K; \skt{स्रोत्रसां}\ H}}\skt{ व्याधिनिर्जयः~। चित्ते\-न्द्रि\-य\-प्र\-सा\-दश्च }\edtext{\skt{शिरसः}}{\lemma{\skt{शिरसः}}\Afootnote{A K H\textsuperscript{ac}; \skt{मनसः}\ H\textsuperscript{pc}}}\skt{ शुद्धिलक्षणं~॥}\pend

\pstart\versepar \ledsidenote{\passref{4.40.39}}\edtext{\skt{कण्डूपदे गुरुता}}{\lemma{\skt{कण्डूपदे गुरुता}}\Afootnote{\skt{कण्डूपदेहौ गुरुता}\ A; \skt{कण्डू\appbrk{}प\appbrk{}दे/\appbrk{}उ\appbrk{}न्च्ले\appbrk{}अ\appbrk{}र्>\appbrk{}गु\appbrk{}रुता}\ K; \skt{कण्डू\appbrk{}प्र\appbrk{}दे\appbrk{}ह\appbrk{}गु\appbrk{}रुता}\ H}}\skt{ }\edtext{\skt{स्रोतसां}}{\lemma{\skt{स्रोतसां}}\Afootnote{A K; \skt{स्रोत्रसाङ्}\ H}}\skt{ }\edtext{\skt{कफसम्भवः}}{\lemma{\skt{कफसम्भवः}}\Afootnote{K H; \skt{कफसंस्रवः~।}\ A}}\skt{~। }\pend
\pstart\versepar \skt{}\edtext{\skt{हीनाविशुद्धे}}{\lemma{\skt{हीनाविशुद्धे}}\Afootnote{K H; \skt{मूर्ध्नि हीनविशुद्धे}\ A}}\skt{ }\edtext{\skt{शिरसि}}{\lemma{\skt{शिरसि}}\Afootnote{K H; \skt{तु}\ A}}\skt{ लक्ष\-ण\-म्प\-रि\-की\-र्त्ति\-तम्~॥}\pend

\pstart\versepar \ledsidenote{\passref{4.40.40}}\skt{मस्तुलुङ्गागमो वा }\edtext{\edtext{\skt{तद्दृ\-ष्टि\-रि\-न्द्रि\-य\-वि\-भ्रमः}}{\lemma{\skt{तद्दृष्टिर्}}\Afootnote{K H; \skt{तवृद्धि}\ A}}}{\lemma{\skt{इन्द्रियविभ्रमः}}\Afootnote{K H; \skt{रिन्द्रियविभ्रमः~।}\ A}}\skt{~। }\pend
\pstart\versepar \skt{शून्यता शिरसश्चापि मूर्ध्नि गाढविरेचिते~॥}\pend

\pstart\versepar \ledsidenote{\passref{4.40.41}}\skt{हीनातिशुद्धे शिरसि कफ\-वा\-त\-घ्न\-मा\-च\-रेत्~। }\pend
\pstart\versepar \skt{सम्यग्विरेचिते }\edtext{\skt{चापि}}{\lemma{\skt{चापि}}\Afootnote{K H; \skt{रसि}\ A}}\skt{ }\edtext{\edtext{\skt{सर्पिर्नस्ते}}{\lemma{\skt{सर्पिर्}}\Afootnote{K H; \skt{सर्पि}\ A}}}{\lemma{\skt{नस्ते}}\Afootnote{K H; \skt{र्न}\ A}}\skt{ }\edtext{\skt{नुषेचयेत्}}{\lemma{\skt{नुषेचयेत्}}\Afootnote{K; \skt{स्यं निषेचयेत्~॥}\ A; \skt{ऽनुषेचयेत्~॥}\ H}}\skt{~॥}\edtext{}{\lemma{\skt{आचरेत्~। सम्यग्विरेचिते}}\Afootnote{K; \skt{आचरेत्~। सम्यग्विशुद्धे शि}\ A; \skt{आचरेत्~॥ सम्यक्विरेचिते}\ H}}\pend

\pstart \ledsidenote{\passref{4.40.42}}\edtext{\skt{नस्ये तु}}{\lemma{\skt{नस्ये तु}}\Afootnote{K H; \skt{(एकान्तरं द्व्यनतरं वा सप्ताहं वा पुनः पुनः~। एक\appbrk{}विं\appbrk{}श\appbrk{}ति\appbrk{}रात्रं वा यावद्वा साधु मन्यते~॥)}\ A}}\skt{ }\edtext{\skt{परिर्हर्त्तव्या}}{\lemma{\skt{परिर्हर्त्तव्या}}\Afootnote{K; \emph{om.}\ A; \skt{परिहर्त्तव्या}\ H}}\skt{ }\edtext{\skt{भवन्ति}}{\lemma{\skt{भवन्ति}}\Afootnote{K H; \emph{om.}\ A}}\skt{~। }\edtext{\skt{तद्यथा}}{\lemma{\skt{तद्यथा}}\Afootnote{K H; \emph{om.}\ A}}\pend

\pstart \ledsidenote{\passref{4.40.47}}\edtext{\skt{भुक्त\-वा\-न\-प\-त\-र्पि\-तो\-त्यर्थं}}{\lemma{\skt{भुक्त\appbrk{}वा\appbrk{}न\appbrk{}प\appbrk{}त\appbrk{}र्पि\appbrk{}तो\appbrk{}त्यर्थं}}\Afootnote{K H; \skt{नस्येन परिहर्तव्यो भुत\appbrk{}वा\appbrk{}न\appbrk{}प\appbrk{}त\appbrk{}र्पि\appbrk{}तोऽ\appbrk{}त्यर्थ}\ A}}\skt{ तरु\-ण\-प्र\-ति\-श्यायी }\edtext{\skt{गर्भिनी मद्य\-द्र\-व\-पी\-डितो}}{\lemma{\skt{गर्भिनी मद्य\appbrk{}द्र\appbrk{}व\appbrk{}पी\appbrk{}डितो}}\Afootnote{K; \skt{गर्भिणी पीत\appbrk{}स्ने\appbrk{}हो\appbrk{}द\appbrk{}क\appbrk{}म\appbrk{}द्य\appbrk{}द्रवोऽ}\ A; \skt{गर्ब्भिणी मद्य\appbrk{}द्र\appbrk{}व\appbrk{}पी\appbrk{}डितो}\ H}}\skt{ }\edtext{\skt{जीर्ण्णीवस्तिदन्तः}}{\lemma{\skt{जीर्ण्णीवस्तिदन्तः}}\Afootnote{K H; \skt{जीर्णी दत्तबस्तिः}\ A}}\skt{ क्रुद्धो }\edtext{\skt{रागार्त्तः}}{\lemma{\skt{रागार्त्तः}}\Afootnote{K H; \skt{गरार्त}\ A}}\skt{ }\edtext{\skt{तृष्णाभिभूतो}}{\lemma{\skt{तृष्णाभिभूतो}}\Afootnote{K H; \skt{स्तृषितः शोकाभिभूतः आन्तो}\ A}}\skt{ बालो वृद्धः }\edtext{\skt{श्रान्तो}}{\lemma{\skt{श्रान्तो}}\Afootnote{K H; \emph{om.}\ A}}\skt{ }\edtext{\skt{वेगा\-व\-रो\-ध\-रतः}}{\lemma{\skt{वेगा\appbrk{}व\appbrk{}रो\appbrk{}ध\appbrk{}रतः}}\Afootnote{K H; \skt{वेगावरोधितः}\ A}}\skt{ }\edtext{\edtext{\skt{शिर\-स्ना\-तु\-का\-म\-श्चेति}}{\lemma{\skt{शिर\appbrk{}स्ना\appbrk{}तु\appbrk{}का\appbrk{}मश्}}\Afootnote{K H; \skt{शिरः\appbrk{}स्ना\appbrk{}तु\appbrk{}का\appbrk{}मश्}\ A}}}{\lemma{\skt{चेति}}\Afootnote{K H; \skt{चेति; अनार्तवे चाभ्रे नस्यधूमौ परिहरेत्~॥}\ A}}\skt{~॥}\pend

\pstart\versepar \ledsidenote{\passref{4.40.49}}\edtext{\skt{नस्ये}}{\lemma{\skt{नस्ये}}\Afootnote{K H; \skt{भवतश्चात्र नस्ये}\ A}}\skt{ शिरोविरेके च व्यापदो द्विविधाः स्मृताः~। }\pend
\pstart\versepar \skt{}\edtext{\skt{दोषो\-त्क्ले\-श\-क्ष\-या\-च्चैव}}{\lemma{\skt{दोषो\appbrk{}त्क्ले\appbrk{}श\appbrk{}क्ष\appbrk{}या\appbrk{}च्चैव}}\Afootnote{K H; \skt{दोषो\appbrk{}त्क्ले\appbrk{}शा\appbrk{}त्क्ष\appbrk{}या\appbrk{}च्चैव}\ A}}\skt{ }\edtext{\skt{विज्ञेया च}}{\lemma{\skt{विज्ञेया च}}\Afootnote{K H; \skt{विज्ञेयास्ता}\ A}}\skt{ यथाक्रमं~।}\pend

\pstart\versepar \ledsidenote{\passref{4.40.50}}\edtext{\skt{तत्रो\-त्क्ले\-श\-नि\-मि\-त्तासु}}{\lemma{\skt{तत्रो\appbrk{}त्क्ले\appbrk{}श\appbrk{}नि\appbrk{}मि\appbrk{}त्तासु}}\Afootnote{K H; \skt{दोषो\appbrk{}क्ले\appbrk{}श\appbrk{}नि\appbrk{}मि\appbrk{}त्तास्तु}\ A}}\skt{ }\edtext{\skt{युञ्ज्या\-च्छ\-म\-न\-शो\-धणे}}{\lemma{\skt{युञ्ज्या\appbrk{}च्छ\appbrk{}म\appbrk{}न\appbrk{}शो\appbrk{}धणे}}\Afootnote{K; \skt{जये\appbrk{}च्छ\appbrk{}म\appbrk{}न\appbrk{}शो\appbrk{}धनैः~।}\ A; \skt{युञ्ज्या\appbrk{}च्छा\appbrk{}म\appbrk{}न\appbrk{}शो\appbrk{}धने~।}\ H\textsuperscript{ac}; \skt{युञ्ज्या च्छमनकोपने~।}\ H\textsuperscript{pc}}}\skt{~। }\pend
\pstart\versepar \skt{अथ क्षयनिमित्तासु यथास्वं }\edtext{\skt{बृंहणं}}{\lemma{\skt{बृंहणं}}\Afootnote{A K; \skt{बृंहनं}\ H}}\skt{ हितं~॥}\pend

\pstart \ledsidenote{\passref{4.40.51}}\edtext{\skt{प्रतिमर्शेस्तु}}{\lemma{\skt{प्रतिमर्शेस्तु}}\Afootnote{K; \skt{प्रतिमर्शश्}\ A; \skt{प्रतिमर्शस्तु}\ H}}\skt{ चतुर्दशसु }\edtext{\skt{कालेषूपदेयो}}{\lemma{\skt{कालेषूपदेयो}}\Afootnote{K H; \skt{कालेषूपादेयः;}\ A}}\skt{ }\edtext{\skt{भवति}}{\lemma{\skt{भवति}}\Afootnote{K H; \emph{om.}\ A}}\skt{~। तद्यथा~। }\edtext{\skt{काल्योत्थितेन}}{\lemma{\skt{काल्योत्थितेन}}\Afootnote{K H; \skt{तल्पोत्थितेन,}\ A}}\skt{~। प्रक्षा\-लि\-त\-द\-न्तेन~। }\edtext{\skt{गृहान्निर्गच्छता}}{\lemma{\skt{गृहान्}}\Afootnote{A K; \skt{गृहा}\ H}}\skt{~। }\edtext{\skt{व्याया\-म\-व्य\-वा\-य\-प\-रि\-श्रा\-तेन}}{\lemma{\skt{व्याया\appbrk{}म\appbrk{}व्य\appbrk{}वा\appbrk{}य\appbrk{}प\appbrk{}रि\appbrk{}श्रा\appbrk{}तेन}}\Afootnote{K H; \skt{व्याया\appbrk{}म\appbrk{}व्य\appbrk{}वा\appbrk{}या\appbrk{}ध्व\appbrk{}प\appbrk{}रि\appbrk{}श्रा\appbrk{}न्तेन,}\ A}}\skt{~। }\edtext{\skt{मूत्रो\-च्चा\-र\-क\-व\-ला\-ञ्ज\-नान्ते}}{\lemma{\skt{मूत्रो\appbrk{}च्चा\appbrk{}र\appbrk{}क\appbrk{}व\appbrk{}ला\appbrk{}ञ्ज\appbrk{}नान्ते}}\Afootnote{A K; \skt{मूत्रो\appbrk{}च्चा\appbrk{}र\appbrk{}क\appbrk{}व\appbrk{}ला\appbrk{}ञ्ज\appbrk{}ना\appbrk{}नान्ते}\ H}}\skt{ भुक्तवता~। छर्दितवता~। दिवा\-स्व\-प्नो\-त्थि\-तेन सायञ्चेति~॥}\pend

\pstart \ledsidenote{\passref{4.40.52}}\skt{तत्र }\edtext{\skt{काल्यो\-त्थि\-ते\-ना\-से\-वितः}}{\lemma{\skt{काल्यो\appbrk{}त्थि\appbrk{}ते\appbrk{}ना\appbrk{}से\appbrk{}वितः}}\Afootnote{K H; \skt{तल्पो\appbrk{}त्थि\appbrk{}ते\appbrk{}ना\appbrk{}से\appbrk{}वितः}\ A}}\skt{ प्रतिमर्शो }\edtext{\edtext{\skt{रात्रावुपचितं}}{\lemma{\skt{रात्राव्}}\Afootnote{K H; \skt{रात्रा}\ A}}}{\lemma{\skt{उपचितं}}\Afootnote{K H; \skt{वुपचितं}\ A}}\skt{ }\edtext{\skt{नासास्रोतोगतं}}{\lemma{\skt{नासास्रोतोगतं}}\Afootnote{A K; \skt{नासा\appbrk{}स्रो\appbrk{}त्रो\appbrk{}ग\appbrk{}तम्}\ H}}\skt{ }\edtext{\skt{मलमपहन्ति}}{\lemma{\skt{अपहन्ति}}\Afootnote{K; \skt{उपहन्ति मनःप्रसादं च करोति,}\ A; \skt{उपहन्ति~।}\ H}}\skt{~। }\edtext{\skt{प्रक्षा\-लि\-त\-द\-न्तेन}}{\lemma{\skt{प्रक्षा\appbrk{}लि\appbrk{}त\appbrk{}द\appbrk{}न्तेन}}\Afootnote{K H; \skt{प्रक्षा\appbrk{}लि\appbrk{}त\appbrk{}द\appbrk{}न्तेना}\ A}}\skt{ }\edtext{\skt{दन्तानां}}{\lemma{\skt{दन्तानां}}\Afootnote{K H; \skt{सेवितो दन्तानां}\ A}}\skt{ }\edtext{\skt{दृढता}}{\lemma{\skt{दृढता}}\Afootnote{K H; \skt{दृढतां}\ A}}\skt{ वद\-न\-सौ\-ग\-न्ध्य\-ञ्चा\-पा\-द\-यति~। गृहान्निर्गच्छता }\edtext{\skt{नासास्रोतसो}}{\lemma{\skt{नासास्रोतसो}}\Afootnote{K; \skt{सेवितो नासासेओतसः}\ A; \skt{नासास्रोत्रसां}\ H}}\skt{ }\edtext{\skt{स्नेहक्लिन्नतया}}{\lemma{\skt{स्नेहक्लिन्नतया}}\Afootnote{K H; \skt{क्लिन्नतया}\ A}}\skt{~। रजोधूमो वा न }\edtext{\skt{प्रबाधते}}{\lemma{\skt{प्रबाधते}}\Afootnote{K H; \skt{बाधते,}\ A}}\skt{~। }\edtext{\skt{व्याया\-म\-व्य\-वा\-य\-प\-रि\-श्रा\-न्तेन}}{\lemma{\skt{व्याया\appbrk{}म\appbrk{}व्य\appbrk{}वा\appbrk{}य\appbrk{}प\appbrk{}रि\appbrk{}श्रा\appbrk{}न्तेन}}\Afootnote{K H; \skt{व्याया\appbrk{}म\appbrk{}मै\appbrk{}थु\appbrk{}ना\appbrk{}ध्व\appbrk{}प\appbrk{}रि\appbrk{}श्रा\appbrk{}न्तेना}\ A}}\skt{ }\edtext{\edtext{\skt{श्रममपहन्ति}}{\lemma{\skt{श्रमम्}}\Afootnote{K H; \skt{सेवितः श्रमम्}\ A}}}{\lemma{\skt{अपहन्ति}}\Afootnote{K H; \skt{उपहन्ति,}\ A}}\skt{~। मूत्रोच्चारान्ते }\edtext{\skt{दृष्टि\-गु\-रु\-त्व\-म\-प\-न\-यति}}{\lemma{\skt{दृष्टिगुरुत्वम्}}\Afootnote{K H; \skt{सेवितो दृष्टेर्गुरुत्वम्}\ A}}\skt{~। कवलाञ्जनान्ते }\edtext{\skt{दृष्टिं}}{\lemma{\skt{दृष्टिं}}\Afootnote{K H; \skt{सेवितो दृष्टिं}\ A}}\skt{ प्रसादयति~। भुक्तवता }\edtext{\skt{स्रोतो}}{\lemma{\skt{स्रोतो}}\Afootnote{K; \skt{सेवितः स्रोतसां}\ A; \skt{स्रोत्रो}\ H}}\skt{ विशुद्धिं }\edtext{\skt{लघुत्वं}}{\lemma{\skt{लघुत्वं}}\Afootnote{K H; \skt{लघुतां}\ A}}\skt{ चापादयति~। }\edtext{\skt{छर्दितवता}}{\lemma{\skt{छर्दितवता}}\Afootnote{K H; \skt{वान्तेनासेवितः}\ A}}\skt{ }\edtext{\skt{स्रोतोविलग्नं}}{\lemma{\skt{स्रोतोविलग्नं}}\Afootnote{A K; \skt{स्रोत्रोविलग्नं}\ H}}\skt{ }\edtext{\edtext{\skt{श्लेष्माणमपोह्य}}{\lemma{\skt{श्लेष्माणम्}}\Afootnote{K H; \skt{श्लेष्माण}\ A}}}{\lemma{\skt{अपोह्य}}\Afootnote{K H; \skt{मपोह्य}\ A}}\skt{ }\edtext{\skt{भक्तच्छंदं}}{\lemma{\skt{भक्तच्छंदं}}\Afootnote{K H; \skt{भक्ताकाङ्क्षाम्}\ A}}\skt{ }\edtext{\skt{कुरुते}}{\lemma{\skt{कुरुते}}\Afootnote{K H; \skt{आपादयति,}\ A}}\skt{~। }\edtext{\skt{दिवा\-स्व\-प्नो\-त्थि\-तेन}}{\lemma{\skt{दिवा\appbrk{}स्व\appbrk{}प्नो\appbrk{}त्थि\appbrk{}तेन}}\Afootnote{K H; \skt{दिवा\appbrk{}स्व\appbrk{}प्नो\appbrk{}त्थि\appbrk{}तेना}\ A}}\skt{ }\edtext{\skt{निद्राशेषं}}{\lemma{\skt{निद्राशेषं}}\Afootnote{K H; \skt{सेवितो निद्राशेषं}\ A}}\skt{ }\edtext{\edtext{\skt{गौर\-व\-ञ्चा\-प\-हत्य}}{\lemma{\skt{गौरवञ्}}\Afootnote{K H; \skt{गुरुत्वं मलं}\ A}}}{\lemma{\skt{चापहत्य}}\Afootnote{K; \skt{चापोह्य}\ A; \skt{चापहृत्य}\ H}}\skt{ चित्तैकाग्र्यं जनयति~। सायं }\edtext{\edtext{\skt{सुख\-नि\-द्रा\-प्र\-बो\-ध\-मा\-पा\-द\-य\-तीति}}{\lemma{\skt{सुख\appbrk{}नि\appbrk{}द्रा\appbrk{}प्र\appbrk{}बो\appbrk{}धम्}}\Afootnote{K H; \skt{चासेवितः सुख\appbrk{}नि\appbrk{}द्रा\appbrk{}प्र\appbrk{}बोधं}\ A}}}{\lemma{\skt{आपादयतीति}}\Afootnote{K H; \skt{चेति~॥}\ A}}\skt{~॥}\pend

\pstart\versepar \ledsidenote{\passref{4.40.53}}\edtext{\skt{ईष\-दु\-च्छिं\-घ\-णा\-त्स्नेहो}}{\lemma{\skt{ईष\appbrk{}दु\appbrk{}च्छिं\appbrk{}घ\appbrk{}णा\appbrk{}त्स्नेहो}}\Afootnote{K; \skt{ईषदुच्छिङ्घतः स्नेहो}\ A; \skt{ईष\appbrk{}दु\appbrk{}च्छि\appbrk{}ङ्घ\appbrk{}ना\appbrk{}त्स्नेहो}\ H}}\skt{ }\edtext{\skt{यदा}}{\lemma{\skt{यदा}}\Afootnote{K H; \skt{याव}\ A}}\skt{ }\edtext{\skt{वक्तुम्प्रपद्यते}}{\lemma{\skt{वक्तुम्}}\Afootnote{K H; \skt{द्वक्त्रं}\ A}}\skt{~। }\pend
\pstart\versepar \skt{}\edtext{\skt{नस्ते}}{\lemma{\skt{नस्ते}}\Afootnote{K H; \skt{नस्ये}\ A}}\skt{ }\edtext{\skt{निशिक्तं तं}}{\lemma{\skt{निशिक्तं तं}}\Afootnote{K; \skt{निषिक्तं तं}\ A H}}\skt{ }\edtext{\skt{विद्यात्प्रतिमर्शं}}{\lemma{\skt{विद्यात्प्रतिमर्शं}}\Afootnote{A K; \skt{विद्या\appbrk{}त्प्र\appbrk{}ति\appbrk{}म\appbrk{}र्षम्}\ H}}\skt{ प्रमाणतः~॥}\pend

\pstart\versepar \ledsidenote{\passref{4.40.54}}\edtext{\skt{तेन}}{\lemma{\skt{तेन}}\Afootnote{K H; \skt{नस्येन}\ A}}\skt{ रोगाः }\edtext{\skt{प्रशाम्यन्ति}}{\lemma{\skt{प्रशाम्यन्ति}}\Afootnote{K H; \skt{शाम्यन्ति}\ A}}\skt{ }\edtext{\edtext{\skt{नरा\-णा\-मू\-र्ध्व\-ज\-त्रुजाः}}{\lemma{\skt{नराणाम्}}\Afootnote{K H; \skt{नराणा}\ A}}}{\lemma{\skt{ऊर्ध्वजत्रुजाः}}\Afootnote{K H; \skt{मूर्ध्वजत्रुजाः~।}\ A}}\skt{~। }\pend
\pstart\versepar \skt{इन्द्रियाणाञ्च वैमल्यं }\edtext{\skt{कुर्या\-दा\-स्य\-सु\-गन्धि}}{\lemma{\skt{कुर्या\appbrk{}दा\appbrk{}स्य\appbrk{}सु\appbrk{}गन्धि}}\Afootnote{K H; \skt{कुर्यादास्यं सुगन्धि}\ A}}\skt{ च~॥}\pend

\pstart\versepar \ledsidenote{\passref{4.40.55}}\skt{हनु\-द\-न्त\-शि\-रो\-ग्री\-वा\-त्रि\-क\-बा\-हू\-रसां बलं~। }\pend
\pstart\versepar \skt{}\edtext{\edtext{\skt{वली\-प\-लि\-त\-खा\-लि\-त्य\-व्यं\-गा\-ना\-ञ्चा\-प्य\-स\-म्भवः}}{\lemma{\skt{वली\appbrk{}प\appbrk{}लि\appbrk{}त\appbrk{}खा\appbrk{}लि\appbrk{}त्य\appbrk{}व्यं\appbrk{}गा\appbrk{}नाञ्}}\Afootnote{K H; \skt{वली\appbrk{}प\appbrk{}लि\appbrk{}त\appbrk{}खा\appbrk{}लि\appbrk{}त्य\appbrk{}व्य\appbrk{}ङ्गानां}\ A}}}{\lemma{\skt{असम्भवः}}\Afootnote{K H; \skt{असंभवम्~॥}\ A}}\skt{~॥}\pend

\pstart\versepar \ledsidenote{\passref{4.40.56}}\skt{तैलं }\edtext{\skt{कफे च}}{\lemma{\skt{कफे च}}\Afootnote{K H; \skt{कफे स}\ A}}\skt{ }\edtext{\skt{वाते च}}{\lemma{\skt{वाते च}}\Afootnote{K H; \skt{वाते स्या}\ A}}\skt{ }\edtext{\skt{केवलं}}{\lemma{\skt{केवलं}}\Afootnote{K H; \skt{त्केवले}\ A}}\skt{ पवने वसां~। }\pend
\pstart\versepar \skt{}\edtext{\edtext{\skt{दद्यान्नस्तः}}{\lemma{\skt{दद्यान्}}\Afootnote{K H; \skt{दद्या}\ A}}}{\lemma{\skt{नस्तः}}\Afootnote{K H; \skt{त्सर्पिः}\ A}}\skt{ सदा पित्ते }\edtext{\skt{सर्पि\-र्म\-ज्जा\-स\-मा\-रुते}}{\lemma{\skt{सर्पि\appbrk{}र्म\appbrk{}ज्जा\appbrk{}स\appbrk{}मा\appbrk{}रुते}}\Afootnote{K H; \skt{मज्जानं च समारुते~॥}\ A}}\skt{~॥}\pend

\pstart\versepar \ledsidenote{\passref{4.40.57}}\skt{चतुर्विधस्य स्नेहस्य }\edtext{\edtext{\skt{विधिरेष}}{\lemma{\skt{विधिर्}}\Afootnote{K H; \skt{विधि}\ A}}}{\lemma{\skt{एष}}\Afootnote{K H; \skt{रेवं}\ A}}\skt{ प्रकीर्त्तितः~। }\pend
\pstart\versepar \skt{श्लेष्म\-स्था\-ना\-वि\-रो\-धि\-त्वा\-त्तेषु तैलं }\edtext{\skt{विशिष्यते}}{\lemma{\skt{विशिष्यते}}\Afootnote{K H; \skt{विधीयते~॥}\ A}}\skt{~॥}\pend

\pstart \ledsidenote{\passref{4.40.44}}\edtext{\skt{अवपीडन्तु}}{\lemma{\skt{अवपीडन्तु}}\Afootnote{K H; \skt{अवपीडस्तु}\ A}}\skt{ }\edtext{\skt{शिरो\-वि\-रे\-च\-न\-व\-द\-भि\-ष्य\-ण्ण\-स\-र्प\-द\-ष्ट\-मू\-र्च्छि\-त\-वि\-सं\-ज्ञेभ्यो}}{\lemma{\skt{शिरो\appbrk{}वि\appbrk{}रे\appbrk{}च\appbrk{}न\appbrk{}व\appbrk{}द\appbrk{}भि\appbrk{}ष्य\appbrk{}ण्ण\appbrk{}स\appbrk{}र्प\appbrk{}द\appbrk{}ष्ट\appbrk{}मू\appbrk{}र्च्छि\appbrk{}त\appbrk{}वि\appbrk{}सं\appbrk{}ज्ञेभ्यो}}\Afootnote{K H; \skt{शिरो\appbrk{}वि\appbrk{}रे\appbrk{}च\appbrk{}न\appbrk{}व\appbrk{}द\appbrk{}भि\appbrk{}ष्य\appbrk{}ण्ण\appbrk{}स\appbrk{}र्प\appbrk{}द\appbrk{}ष्ट\appbrk{}वि\appbrk{}सं\appbrk{}ज्ञेभ्यो}\ A}}\skt{ }\edtext{\edtext{\edtext{\skt{दद्या\-च्छि\-रो\-वि\-रे\-च\-न\-द्र\-व्या\-णा\-म\-न्य\-त\-म\-मा\-सा\-द्या\-व\-पी\-ड\-मि\-ष्य\-तेति}}{\lemma{\skt{दद्या\appbrk{}च्छि\appbrk{}रो\appbrk{}वि\appbrk{}रे\appbrk{}च\appbrk{}न\appbrk{}द्र\appbrk{}व्या\appbrk{}णाम्}}\Afootnote{K H; \skt{दद्या\appbrk{}च्छि\appbrk{}रो\appbrk{}चि\appbrk{}रे\appbrk{}च\appbrk{}न\appbrk{}द्र\appbrk{}व्या\appbrk{}णाम्}\ A}}}{\lemma{\skt{आसा\appbrk{}द्या\appbrk{}व\appbrk{}पी\appbrk{}डम्}}\Afootnote{K H; \skt{अव\appbrk{}पि\appbrk{}ष्या\appbrk{}व\appbrk{}पीड्य}\ A}}}{\lemma{\skt{इष्यतेति}}\Afootnote{K H; \skt{च,}\ A}}\skt{~॥}\pend

\pstart \ledsidenote{\passref{4.40.46}}\edtext{\skt{चेतो\-वि\-का\-र\-क्रि\-मि\-वि\-षा\-दिति}}{\lemma{\skt{चेतो\appbrk{}वि\appbrk{}का\appbrk{}र\appbrk{}क्रि\appbrk{}मि\appbrk{}वि\appbrk{}षा\appbrk{}दिति}}\Afootnote{\skt{चेतो\appbrk{}वि\appbrk{}का\appbrk{}र\appbrk{}कृ\appbrk{}मि\appbrk{}वि\appbrk{}षा\appbrk{}भि\appbrk{}प\appbrk{}न्नानां चूर्णं प्रधमेत्~॥}\ A; \skt{चेतो\appbrk{}वि\appbrk{}का\appbrk{}र\appbrk{}क्रि\appbrk{}मि\appbrk{}वि\appbrk{}षा…\appbrk{}दिति~॥}\ K; \skt{चेतो\appbrk{}वि\appbrk{}का\appbrk{}र\appbrk{}क्रि\appbrk{}मि\appbrk{}वि\appbrk{}षा\appbrk{}भि\appbrk{}प\appbrk{}न्ना\appbrk{}ना\appbrk{}ञ्चूर्ण्णं प्रधमेदिति~॥}\ H}}\skt{~॥}\pend

\pstart \ledsidenote{\passref{4.40.44a}}\edtext{\skt{शर्क\-रे\-क्षु\-र\-स\-क्षीरा}}{\lemma{\skt{शर्क\appbrk{}रे\appbrk{}क्षु\appbrk{}र\appbrk{}स\appbrk{}क्षीरा}}\Afootnote{K H; \skt{शर्क\appbrk{}रे\appbrk{}क्षु\appbrk{}र\appbrk{}स\appbrk{}क्षीर}\ A}}\skt{ }\edtext{\skt{घृतमांसरसां}}{\lemma{\skt{घृतमांसरसां}}\Afootnote{K; \skt{घृत\appbrk{}मां\appbrk{}स\appbrk{}र\appbrk{}सा\appbrk{}नाम्}\ A; \skt{घृत\appbrk{}मां\appbrk{}स\appbrk{}र\appbrk{}सम्}\ H}}\skt{ }\edtext{\skt{पृथक्क्षीणानां}}{\lemma{\skt{पृथक्क्षीणानां}}\Afootnote{K H; \skt{अन्यतमं क्षीणानां}\ A}}\skt{ }\edtext{\skt{नस्ततो}}{\lemma{\skt{नस्ततो}}\Afootnote{K H; \skt{शोणित}\ A}}\skt{ }\edtext{\skt{दद्या\-द्द्र\-क्त\-पि\-त्त\-ग\-देषु}}{\lemma{\skt{दद्या\appbrk{}द्द्र\appbrk{}क्त\appbrk{}पि\appbrk{}त्त\appbrk{}ग\appbrk{}देषु}}\Afootnote{K; \skt{पित्ते च विदध्यात्~॥}\ A; \skt{दद्या\appbrk{}द्र\appbrk{}क्त\appbrk{}पि\appbrk{}त्त\appbrk{}ग\appbrk{}देषु}\ H}}\skt{ }\edtext{\skt{च}}{\lemma{\skt{च}}\Afootnote{K H; \emph{om.}\ A}}\skt{~॥}\pend

\pstart\versepar \ledsidenote{\passref{4.40.58}}\skt{अतः परम्प्रवक्ष्यामि }\edtext{\skt{कवडग्रहणे}}{\lemma{\skt{कवडग्रहणे}}\Afootnote{K H; \skt{कवलग्रहणे}\ A}}\skt{ विधिं~। }\pend
\pstart\versepar \skt{चतुर्धा }\edtext{\skt{कवडस्नेही}}{\lemma{\skt{कवडस्नेही}}\Afootnote{K H; \skt{कवलः स्नेही}\ A}}\skt{ प्रसादी शोधिरोपणौ~।}\pend

\pstart \ledsidenote{\passref{4.40.59}}\skt{स्न्गिधोष्णैः स्नैहिको वाते स्वादुशीतैः प्रसादनः पित्ते }\edtext{\skt{कट्वम्ब्ललवणै}}{\lemma{\skt{कट्वम्ब्ललवणै}}\Afootnote{K; \skt{कट्वम्ललवणै}\ A H}}\skt{ }\edtext{\skt{रूक्षैः}}{\lemma{\skt{रूक्षैः}}\Afootnote{K H; \skt{रूक्षो}\ A}}\skt{ }\edtext{\skt{संशोधणः}}{\lemma{\skt{संशोधणः}}\Afootnote{K; \skt{ष्णैः शोधनः}\ A; \skt{संशोधनः}\ H}}\skt{ कफे~।}\pend

\pstart \ledsidenote{\passref{4.40.60}}\skt{कषा\-य\-ति\-क्त\-म\-धुरैः }\edtext{\skt{कटूष्णो}}{\lemma{\skt{कटूष्णो}}\Afootnote{K H; \skt{कटूष्णै}\ A}}\skt{ रोपणो }\edtext{\skt{व्रणे}}{\lemma{\skt{व्रणे}}\Afootnote{K H; \skt{व्रणे~। चतुर्विधस्य चैवास्य विशेषोऽयं प्रकीर्तितः~॥}\ A}}\skt{~।}\pend

\pstart \ledsidenote{\passref{4.40.61}}\skt{तत्र }\edtext{\edtext{\skt{त्रिक\-टु\-का\-शि\-त\-स\-र्ष\-प\-सि\-द्धा\-र्थ\-क\-क\-ल्क\-मा\-लोड्य}}{\lemma{\skt{त्रिक\appbrk{}टु\appbrk{}का\appbrk{}शि\appbrk{}त\appbrk{}स\appbrk{}र्ष\appbrk{}प\appbrk{}सि\appbrk{}द्धा\appbrk{}र्थ\appbrk{}क\appbrk{}क\appbrk{}ल्कम्}}\Afootnote{K H; \skt{त्रिक\appbrk{}टु\appbrk{}क\appbrk{}व\appbrk{}चा\appbrk{}स\appbrk{}र्ष\appbrk{}प\appbrk{}ह\appbrk{}री\appbrk{}त\appbrk{}की\appbrk{}कल्क}\ A}}}{\lemma{\skt{आलोड्य}}\Afootnote{K H; \skt{मालोड्य}\ A}}\skt{ }\edtext{\edtext{\skt{तैल\-मू\-त्र\-क्षा\-र\-सु\-रा\-शु\-क्त\-म\-धू\-ना\-म\-न्य\-त\-मेन}}{\lemma{\skt{तैल\appbrk{}मू\appbrk{}त्र\appbrk{}क्षा\appbrk{}र\appbrk{}सु\appbrk{}रा\appbrk{}शु\appbrk{}क्त\appbrk{}म\appbrk{}धू\appbrk{}नाम्}}\Afootnote{K; \skt{तैल\appbrk{}शु\appbrk{}क्त\appbrk{}सु\appbrk{}रा\appbrk{}मू\appbrk{}त्र\appbrk{}क्षा\appbrk{}र\appbrk{}म\appbrk{}धूना}\ A; \skt{तैल\appbrk{}मू\appbrk{}त्र\appbrk{}क्षा\appbrk{}र\appbrk{}सु\appbrk{}रा\appbrk{}शु\appbrk{}क्र\appbrk{}म\appbrk{}धू\appbrk{}नाम्}\ H}}}{\lemma{\skt{अन्यतमेन}}\Afootnote{K H; \skt{मन्यतमेन}\ A}}\skt{ }\edtext{\skt{सल\-व\-ण\-म\-भि\-प्र\-त\-प्त\-मु\-प\-स्वि\-न्न\-मृ\-दि\-त\-ग\-ल\-क\-पो\-ल\-ल\-ला\-ट\-प्र\-देशे}}{\lemma{\skt{उप\appbrk{}स्वि\appbrk{}न्न\appbrk{}मृ\appbrk{}दि\appbrk{}त\appbrk{}ग\appbrk{}ल\appbrk{}क\appbrk{}पो\appbrk{}ल\appbrk{}ल\appbrk{}ला\appbrk{}ट\appbrk{}प्र\appbrk{}देशे}}\Afootnote{K H; \skt{उप\appbrk{}स्वि\appbrk{}न्न\appbrk{}मृ\appbrk{}दि\appbrk{}त\appbrk{}ग\appbrk{}ल\appbrk{}क\appbrk{}पो\appbrk{}ल\appbrk{}ल\appbrk{}ला\appbrk{}ट\appbrk{}प्र\appbrk{}देशो}\ A}}\skt{ धारयेत्~॥}\pend

\pstart\versepar \ledsidenote{\passref{4.40.62}}\skt{सुखं सञ्चार्यते या तु }\edtext{\skt{मात्रा सा}}{\lemma{\skt{मात्रा सा}}\Afootnote{K; \skt{मात्रा स (सा)}\ A; \skt{सा मात्रा सा}\ H}}\skt{ }\edtext{\skt{कवलः}}{\lemma{\skt{कवलः}}\Afootnote{A K; \skt{कवलाः}\ H}}\skt{ }\edtext{\skt{स्मृताः}}{\lemma{\skt{स्मृताः}}\Afootnote{K H; \skt{स्मृतः~।}\ A}}\skt{~। }\pend
\pstart\versepar \skt{असञ्चार्या तु या मात्रा }\edtext{\skt{गण्डूषः सा}}{\lemma{\skt{गण्डूषः सा}}\Afootnote{K H; \skt{गण्डूषः स}\ A}}\skt{ }\edtext{\skt{प्रकीर्त्तिताः}}{\lemma{\skt{प्रकीर्त्तिताः}}\Afootnote{K H; \skt{प्रकीर्तितः~॥}\ A}}\skt{~॥}\pend

\pstart \ledsidenote{\passref{4.40.63}}\skt{तावच्च धारयितव्यो }\edtext{\skt{नन्य\-म\-न\-सा\-नु\-न्न\-त\-दे\-हेन}}{\lemma{\skt{नन्य\appbrk{}म\appbrk{}न\appbrk{}सा\appbrk{}नु\appbrk{}न्न\appbrk{}त\appbrk{}दे\appbrk{}हेन}}\Afootnote{K H; \skt{ऽन\appbrk{}न्य\appbrk{}म\appbrk{}न\appbrk{}सो\appbrk{}न्न\appbrk{}त\appbrk{}दे\appbrk{}हेन}\ A}}\skt{ }\edtext{\skt{याव\-द्दो\-ष\-प्र\-ति\-पू\-र्ण्ण\-ग\-ल\-क\-पो\-लत्वं}}{\lemma{\skt{याव\appbrk{}द्दो\appbrk{}ष\appbrk{}प्र\appbrk{}ति\appbrk{}पू\appbrk{}र्ण्ण\appbrk{}ग\appbrk{}ल\appbrk{}क\appbrk{}पो\appbrk{}लत्वं}}\Afootnote{K H; \skt{याव\appbrk{}द्दो\appbrk{}ष\appbrk{}प\appbrk{}रि\appbrk{}पू\appbrk{}र्ण\appbrk{}क\appbrk{}पो\appbrk{}लत्वं}\ A}}\skt{ }\edtext{\skt{नासा\-स्रो\-तो\-न\-य\-न\-प\-रि\-प्ल\-वश्च}}{\lemma{\skt{नासा\appbrk{}स्रो\appbrk{}तो\appbrk{}न\appbrk{}य\appbrk{}न\appbrk{}प\appbrk{}रि\appbrk{}प्ल\appbrk{}वश्च}}\Afootnote{K; \skt{नासा\appbrk{}स्रो\appbrk{}तो\appbrk{}न\appbrk{}य\appbrk{}न\appbrk{}प\appbrk{}रि\appbrk{}प्ला\appbrk{}वश्च}\ A; \skt{नासा\appbrk{}स्रो\appbrk{}त्रो\appbrk{}न\appbrk{}य\appbrk{}न\appbrk{}प\appbrk{}रि\appbrk{}प्ल\appbrk{}वश्च}\ H}}\skt{ भवति तदा विमोक्तव्यः~। पुनश्चान्यो }\edtext{\skt{ग्रहीतव्यः}}{\lemma{\skt{ग्रहीतव्यः}}\Afootnote{K H; \skt{ग्रहीतव्य इति~॥}\ A}}\skt{~॥}\pend

\pstart\versepar \ledsidenote{\passref{4.40.64}}\skt{एवं }\edtext{\skt{स्नेह\-प\-यः\-क्षौ\-द्र\-र\-स\-मू\-त्रा\-म्ब्ल\-स\-म्भृताः}}{\lemma{\skt{स्नेह\appbrk{}प\appbrk{}यः\appbrk{}क्षौ\appbrk{}द्र\appbrk{}र\appbrk{}स\appbrk{}मू\appbrk{}त्रा\appbrk{}म्ब्ल\appbrk{}स\appbrk{}म्भृताः}}\Afootnote{K; \skt{स्नेह\appbrk{}प\appbrk{}यः\appbrk{}क्षौ\appbrk{}द्र\appbrk{}र\appbrk{}स\appbrk{}मू\appbrk{}त्रा\appbrk{}म्ल\appbrk{}सं\appbrk{}भृताः}\ A H}}\skt{~। }\pend
\pstart\versepar \skt{कसा\-यो\-ष्णो\-द\-का\-भ्याञ्च }\edtext{\skt{कवडा}}{\lemma{\skt{कवडा}}\Afootnote{K H; \skt{कवला}\ A}}\skt{ दोषतो हिताः~॥}\pend

\pstart\versepar \ledsidenote{\passref{4.40.65}}\edtext{\edtext{\edtext{\skt{व्याधे\-र\-प\-च\-य\-स्तु\-ष्टि\-र्वि\-षदं}}{\lemma{\skt{अपचयस्}}\Afootnote{K H; \skt{अपचय}\ A}}}{\lemma{\skt{तुष्टिर्}}\Afootnote{K H; \skt{स्तुष्टि}\ A}}}{\lemma{\skt{विषदं}}\Afootnote{K H; \skt{र्वैशद्यं}\ A}}\skt{ वक्त्रलाघवं~। }\pend
\pstart\versepar \skt{}\edtext{\skt{इन्द्रि\-या\-णा\-म्प्र\-सा\-दश्च}}{\lemma{\skt{इन्द्रियाणाम्}}\Afootnote{K H; \skt{इन्दिर्याणां}\ A}}\skt{ }\edtext{\skt{कवडे}}{\lemma{\skt{कवडे}}\Afootnote{K H; \skt{कवले}\ A}}\skt{ शुद्धिलक्षणं~।}\pend

\pstart\versepar \ledsidenote{\passref{4.40.66}}\skt{हीने जाड्यकफोत्क्लेशा वर\-स\-ज्ञा\-न\-मेव च~। }\pend
\pstart\versepar \skt{}\edtext{\skt{अतियोगात्मुखे}}{\lemma{\skt{अतियोगात्मुखे}}\Afootnote{K H; \skt{अतियोगान्मुखे}\ A}}\skt{ पाकः }\edtext{\edtext{\skt{शोष\-स्तृ\-ष्णा\-रु\-चि\-क्लमः}}{\lemma{\skt{शोषस्}}\Afootnote{K H; \skt{शोष}\ A}}}{\lemma{\skt{तृष्णारुचिक्लमः}}\Afootnote{K; \skt{तृष्णारुचिक्लमाः~॥}\ A H}}\skt{~।}\pend

\pstart\versepar \ledsidenote{\passref{4.40.67}}\skt{शोधनीये विशेषेण }\edtext{\edtext{\skt{भवत्येतत्तु}}{\lemma{\skt{भवत्य्}}\Afootnote{K H; \skt{भवन्त्य्}\ A}}}{\lemma{\skt{एतत्तु}}\Afootnote{K H; \skt{एव न}\ A}}\skt{ }\edtext{\skt{लक्षणं}}{\lemma{\skt{लक्षणं}}\Afootnote{K H; \skt{संशयः~।}\ A}}\skt{~॥ }\pend
\pstart\versepar \skt{}\edtext{\skt{तिलां}}{\lemma{\skt{तिलां}}\Afootnote{K H; \skt{तिला}\ A}}\skt{ }\edtext{\skt{नीलोत्पलं}}{\lemma{\skt{नीलोत्पलं}}\Afootnote{A K; \skt{नीलोत्पलां}\ H}}\skt{ सर्पिः शर्करा क्षीरमेव च~।}\pend

\pstart\versepar \ledsidenote{\passref{4.40.68}}\skt{सक्षौद्रो }\edtext{\skt{दग्धवर्क्त्रस्य}}{\lemma{\skt{दग्धवर्क्त्रस्य}}\Afootnote{K; \skt{दग्धवक्त्रस्य}\ A H}}\skt{ गण्डूषो दाहनाशनः~। }\pend
\pstart\versepar \skt{कवलस्य }\edtext{\skt{समासेन}}{\lemma{\skt{समासेन}}\Afootnote{K\textsuperscript{pc} H; \skt{विधिर्ह्येष समासेन}\ A; \emph{om.}\ K\textsuperscript{ac}}}\skt{ }\edtext{\edtext{\skt{विधिरेषः}}{\lemma{\skt{विधिर्}}\Afootnote{K\textsuperscript{pc} H; \emph{om.}\ A K\textsuperscript{ac}}}}{\lemma{\skt{एषः}}\Afootnote{K\textsuperscript{pc}; \emph{om.}\ A K\textsuperscript{ac}; \skt{एष}\ H}}\skt{ प्रकीर्त्तितः~।}\pend

\pstart\versepar \ledsidenote{\passref{4.40.69}}\skt{विभज्य भेषजं बुध्या कुर्वीत प्रतिसारणं~॥ }\pend
\pstart\versepar \skt{कल्को रसक्रिया क्षौद्रं चूर्ण्णञ्चेति चतुर्विधं~॥}\pend

\pstart\versepar \ledsidenote{\passref{4.40.70}}\skt{अङ्गु\-ल्य\-ग्र\-प्र\-णी\-तन्तु यथा\-स्व\-म्मु\-ख\-रो\-गिणां~॥ }\pend
\pstart\versepar \skt{}\edtext{\skt{तस्य}}{\lemma{\skt{तस्य}}\Afootnote{K H; \skt{तस्मि}\ A}}\skt{ }\edtext{\edtext{\skt{योगमयोगञ्च}}{\lemma{\skt{योगम्}}\Afootnote{K H; \skt{न्योग}\ A}}}{\lemma{\skt{अयोगञ्च}}\Afootnote{K H; \skt{मयोगं च}\ A}}\skt{ कवलोक्तं विभावयेत्~।}\pend

\pstart \ledsidenote{\passref{4.40.71}}\edtext{\skt{भोज\-ये\-त्चा\-न\-भि\-ष्यन्दि}}{\lemma{\skt{भोज\appbrk{}ये\appbrk{}त्चा\appbrk{}न\appbrk{}भि\appbrk{}ष्यन्दि}}\Afootnote{K\textsuperscript{pc} H; \skt{तानेव शम\appbrk{}ये\appbrk{}द्व्या\appbrk{}धी\appbrk{}न्क\appbrk{}वलो यानपोहति~। दोष\appbrk{}घ्न\appbrk{}म\appbrk{}न\appbrk{}भि\appbrk{}ष्यन्दि}\ A; \skt{भोज\appbrk{}ये\appbrk{}त\appbrk{}च्चा\appbrk{}न\appbrk{}भि\appbrk{}ष्यन्दि}\ K\textsuperscript{ac}}}\skt{ }\edtext{\edtext{\skt{दोषघ्नमशनं}}{\lemma{\skt{दोषघ्नम्}}\Afootnote{K H; \skt{भोजये}\ A}}}{\lemma{\skt{अशनं}}\Afootnote{K H; \skt{च्च तथा}\ A}}\skt{ }\edtext{\skt{नरमिति}}{\lemma{\skt{इति}}\Afootnote{K H; \emph{om.}\ A}}\skt{~॥}\pend

\pstart \ledsidenote{\passref{4.40.trailer0}}\skt{॥ चि ४०~॥}\pend

\pstart\versepar \ledsidenote{\passref{4.40.72}}\edtext{\skt{स्नेहः}}{\lemma{\skt{स्नेहः}}\Afootnote{\skt{स्नेह}\ K H}}\skt{ }\edtext{\skt{स्वेदो}}{\lemma{\skt{स्वेदो}}\Afootnote{K H\textsuperscript{ac}; \skt{स्वेदोथ}\ H\textsuperscript{pc}}}\skt{ वमनं विरेकं व्यापदं तथा~। }\pend
\pstart\versepar \skt{नेत्रे }\edtext{\skt{प्रमाण}}{\lemma{\skt{प्रमाण}}\Afootnote{K; \skt{प्रमाणम्}\ H}}\skt{ व्याप\-च्चा\-नु\-वा\-स\-न\-स\-हो\-त्तरं~।}\pend

\pstart \ledsidenote{\passref{4.40.73}}\skt{निरू\-हो\-त्त\-र\-पीडा च धूमनस्तेन वै }\edtext{\skt{दश~॥ {\symfont ❈}}}{\lemma{\skt{दश~॥ {\symfont ❈}}}\Afootnote{K; \skt{दश~॥ ०~॥}\ H}}\skt{~॥}\pend

\pstart \ledsidenote{\passref{4.40.trailer1}}\edtext{\skt{यदु\-क्त\-ञ्च\-त्वा\-रिं\-श\-च्चि\-कि\-त्सि\-ता\-नीति तत्सर्वं }\edtext{\skt{समाप्तमिति~॥ {\symfont ❈}}}{\lemma{\skt{इति~॥ {\symfont ❈}}}\Afootnote{K H; \skt{इति सुश्रु\appbrk{}त\appbrk{}सं\appbrk{}हि\appbrk{}तायां चिकित्सास्थाने धूम\appbrk{}न\appbrk{}स्य\appbrk{}क\appbrk{}व\appbrk{}ल\appbrk{}ग्र\appbrk{}ह\appbrk{}चि\appbrk{}कि\appbrk{}त्सितं नाम चत्वारिंशोऽध्यायः~॥}\ A}}\skt{~॥}}{\lemma{\skt{… इ}}\Afootnote{\emph{om.}\ A}}\pend

\endnumbering
